% !TeX program = pdflatex
% papers/cristal_fisica.tex — GERADO por tools/cristal_tex.py; NÃO editar à mão.
% Fonte: cristal/cristal.jsonl (broca-so, última versão por conceito).
% Cada secção carrega o registo original em comentário %CRISTAL — a volta é
% tests/cristal_volta.js (resíduo 0 byte a byte contra a fonte).
\documentclass[11pt,a4paper]{article}
\usepackage[utf8]{inputenc}\usepackage[T1]{fontenc}\usepackage[brazil]{babel}
\usepackage{amsmath,amssymb}\usepackage{textcomp}
\usepackage[margin=2.4cm]{geometry}\usepackage{xcolor}
\definecolor{cinza}{gray}{0.35}
\newcommand{\prov}[1]{\par\smallskip\noindent{\small\color{cinza}#1}\par}
\setcounter{secnumdepth}{0}
% ─────────────────────────────────────────────────────────────────────────────
%  papers/gkcapa.tex — a capa do Reino, mínima, para os papers em `article`.
%
%  O estilo.tex completo é do `report` (títulos de capítulo, skak, …). Aqui fica
%  só o que a capa precisa, para o paper continuar a compilar sozinho com
%  pdflatex. O tradutor próprio (tests/tex) carrega o estilo.tex da raiz / o
%  symlink papers/estilo.tex e avalia o mesmo `\gkcapa`.
% ─────────────────────────────────────────────────────────────────────────────
\definecolor{ouro}{HTML}{B8912F}
\definecolor{ouroclaro}{HTML}{F3E9CE}
\definecolor{tinta}{HTML}{1E1B16}
\definecolor{regua}{HTML}{6E6858}
\providecommand{\gknota}{\fontsize{7.62}{11.05}\selectfont}
\providecommand{\gktexto}{\fontsize{10.50}{15.22}\selectfont}
\providecommand{\gksub}{\fontsize{12.33}{17.87}\selectfont}
\providecommand{\gksec}{\fontsize{14.47}{20.98}\selectfont}
\providecommand{\gkcap}{\fontsize{16.99}{24.63}\selectfont}
\providecommand{\gktit}{\fontsize{23.42}{33.95}\selectfont}
\providecommand{\Prj}{\mathbb{P}}
\providecommand{\Trf}{\mathcal{T}}
\providecommand{\gkcapa}[3]{%
\title{\vspace{-0.6cm}
{\color{ouro}\rule{\textwidth}{1.4pt}}\\[6mm]
\textbf{\gktit\color{tinta} Reino Dourado}\\[3mm]{\gksec\itshape\color{regua} Golden Kingdom}\\[8mm]
{\gkcap\scshape\color{ouro} #1}\\[5mm]
{\gksub\color{regua} #2}\\[2mm]
{\gktexto\color{regua}\itshape #3}\\[7mm]
{\color{ouro}\rule{\textwidth}{0.6pt}}\\[9mm]{\gknota\color{regua}\begin{minipage}{\textwidth}\setlength{\parindent}{0pt}%
\textbf{Aviso de Isenção Algorítmica:} O universo, as leis e as entidades contidas nestas páginas ---
incluindo felinos matriciais, aves de rapina algébricas e amuletos de controle de fluxo --- são
construções de pura ficção formal. Esta obra não descreve a realidade física, não serve como
engenharia reversa do cosmos e não constitui doutrina religiosa. Qualquer semelhança com bugs reais do seu
sistema operacional é mera coincidência (ou falta de reza). Prossiga por sua própria conta e
risco.\end{minipage}}}
\author{}\date{}}

\begin{document}
\gkcapa{O Cristal --- a física do cristal}
       {cosmologia, metais, joalheria, litografia}
       {corpus recuperado do broca-so; 332 conceitos; projeção verificável da fonte}
\maketitle

%CRISTAL {"arestas":[{"alvo":"plano_continuo_transicoes","confianca":1.0,"epistemico":"fato","origem":"wiki","rel":"parte_de"},{"alvo":"mapear_nao_fabricar_continuo","confianca":1.0,"epistemico":"fato","origem":"wiki","rel":"usa"},{"alvo":"aresta_transicao_graduada","confianca":1.0,"epistemico":"fato","origem":"wiki","rel":"usa"},{"alvo":"leva1_transicoes_l1","confianca":1.0,"epistemico":"fato","origem":"wiki","rel":"relaciona"},{"alvo":"tiffany","confianca":1.0,"epistemico":"fato","origem":"wiki","rel":"parte_de"}],"confianca":1.0,"contraexemplos":[],"descricao":"VERIFICADO (doc/continuo/absorve_metal.c, C puro): absorver uma obra inteira por RESSONÂNCIA no metal é viável e RÁPIDO — a dinâmica da parábola (impressão Koch 256-bit) fragmenta a obra e roteia cada fragmento à sua ARESTA-CASA (ponte=min(ov_A,ov_B)+(ov_A+ov_B)/4) a ~322M overlaps/s, ~39k fragmentos/s (obra inteira em ~1 s; 6000 frags de corpus_aarao em 0,15 s). ACHADO HONESTO: a impressão POSICIONAL cria ATRATORES — nós genéricos/curtos ('Referência', 'Neti Neti') ressoam com quase tudo; 1 aresta absorveu 26% dos fragmentos, cobertura 8,4%. É a mesma raiz do overlap posicional (frase curta/genérica vira ímã). VIABILIDADE = SIM (metal), QUALIDADE de roteamento exige kernel de ressonância discriminante por CONTEÚDO no metal (impressão de tokens / down-weight IDF dos nós centrais) → aí cada frase vai à casa real. [D] roda, cronometrado.","epistemico":"fato","exemplos":[],"id":"absorcao_ressonancia_metal","memoria":"semantica","meta":{"autor":"Lopes/broca-so","dominio":"metal","fonte":"doc/continuo/absorve_metal.c (cc -O3, 322M overlaps/s)"},"origem":"wiki","palavras_chave":[],"sinonimos":[],"tipo":"conceito","titulo":"Absorver uma obra por ressonância no metal: viável e veloz, mas o kernel clumpa"}
\section{Absorver uma obra por ressonância no metal: viável e veloz, mas o kernel clumpa}
VERIFICADO (doc/continuo/absorve\_metal.c, C puro): absorver uma obra inteira por RESSONÂNCIA no metal é viável e RÁPIDO — a dinâmica da parábola (impressão Koch 256-bit) fragmenta a obra e roteia cada fragmento à sua ARESTA-CASA (ponte=min(ov\_A,ov\_B)+(ov\_A+ov\_B)/4) a \~{}322M overlaps/s, \~{}39k fragmentos/s (obra inteira em \~{}1 s; 6000 frags de corpus\_aarao em 0,15 s). ACHADO HONESTO: a impressão POSICIONAL cria ATRATORES — nós genéricos/curtos ('Referência', 'Neti Neti') ressoam com quase tudo; 1 aresta absorveu 26\% dos fragmentos, cobertura 8,4\%. É a mesma raiz do overlap posicional (frase curta/genérica vira ímã). VIABILIDADE = SIM (metal), QUALIDADE de roteamento exige kernel de ressonância discriminante por CONTEÚDO no metal (impressão de tokens / down-weight IDF dos nós centrais) \ensuremath{\to} aí cada frase vai à casa real. [D] roda, cronometrado.
\prov{relações: parte\_de plano\_continuo\_transicoes; usa mapear\_nao\_fabricar\_continuo; usa aresta\_transicao\_graduada; relaciona leva1\_transicoes\_l1; parte\_de tiffany}
\prov{proveniência: wiki \textperiodcentered{} doc/continuo/absorve\_metal.c (cc -O3, 322M overlaps/s) \textperiodcentered{} fato \textperiodcentered{} confiança 1.0 \textperiodcentered{} domínio metal \textperiodcentered{} história 84 versões no jornal}

%CRISTAL {"arestas":[{"alvo":"mecanismo_higgs_ebh","confianca":1.0,"epistemico":"fato","origem":"wiki","rel":"usa"},{"alvo":"materia_modelo_padrao","confianca":1.0,"epistemico":"fato","origem":"wiki","rel":"explica"}],"confianca":1.0,"contraexemplos":[],"descricao":"O Higgs dá massa aos férmions via Yukawa: m_f=y_f·v/√2, logo y_f=√2 m_f/v. Os acoplamentos escalam com a massa (verificado p/ top, bottom, tau, muon no LHC). ABERTO: acoplamentos da 1ª/2ª geração não medidos, e a origem da hierarquia de massas (y_top≈1 vs y_elétron≈10⁻⁶ — o problema do sabor).","epistemico":"fato","exemplos":[],"id":"acoplamentos_yukawa","memoria":"semantica","meta":{"autor":"broca-so","dominio":"fisica","fonte":"CMS/ATLAS Nature 607 (2022)"},"origem":"wiki","palavras_chave":[],"sinonimos":[],"tipo":"conceito","titulo":"Acoplamentos de Yukawa (a massa dos férmions)"}
\section{Acoplamentos de Yukawa (a massa dos férmions)}
O Higgs dá massa aos férmions via Yukawa: m\_f=y\_f\textperiodcentered v/\ensuremath{\sqrt{\,}}2, logo y\_f=\ensuremath{\sqrt{\,}}2 m\_f/v. Os acoplamentos escalam com a massa (verificado p/ top, bottom, tau, muon no LHC). ABERTO: acoplamentos da 1ª/2ª geração não medidos, e a origem da hierarquia de massas (y\_top\ensuremath{\approx}1 vs y\_elétron\ensuremath{\approx}10\ensuremath{^{-6}} — o problema do sabor).
\prov{relações: usa mecanismo\_higgs\_ebh; explica materia\_modelo\_padrao}
\prov{proveniência: wiki \textperiodcentered{} CMS/ATLAS Nature 607 (2022) \textperiodcentered{} fato \textperiodcentered{} confiança 1.0 \textperiodcentered{} domínio fisica \textperiodcentered{} história 3 versões no jornal}

%CRISTAL {"arestas":[{"alvo":"radiacao_hawking_informacao","confianca":1.0,"epistemico":"fato","origem":"wiki","rel":"relaciona"},{"alvo":"cosmologia_quantica_hub","confianca":1.0,"epistemico":"fato","origem":"wiki","rel":"parte_de"}],"confianca":1.0,"contraexemplos":[],"descricao":"Gravidade quântica (cordas) em anti-de Sitter de dimensão d+1 é EXATAMENTE equivalente a uma teoria de campos conforme (sem gravidade) na fronteira d — a realização calculável da holografia. Canônico: IIB em AdS₅×S⁵ ≡ Super Yang-Mills N=4 SU(N). Dualidade matemática extraordinariamente testada (o resultado mais firme da holografia), mas conjectura sem prova geral; nosso universo é de Sitter (Λ>0), não AdS.","epistemico":"inferencia","exemplos":[],"id":"ads_cft_maldacena","memoria":"semantica","meta":{"autor":"broca-so","dominio":"fisica","fonte":"Maldacena 1998 hep-th/9711200"},"origem":"wiki","palavras_chave":[],"sinonimos":[],"tipo":"conceito","titulo":"Correspondência AdS/CFT (Maldacena)"}
\section{Correspondência AdS/CFT (Maldacena)}
Gravidade quântica (cordas) em anti-de Sitter de dimensão d+1 é EXATAMENTE equivalente a uma teoria de campos conforme (sem gravidade) na fronteira d — a realização calculável da holografia. Canônico: IIB em AdS\ensuremath{_{5}}$\times$S\ensuremath{^{5}} \ensuremath{\equiv} Super Yang-Mills N=4 SU(N). Dualidade matemática extraordinariamente testada (o resultado mais firme da holografia), mas conjectura sem prova geral; nosso universo é de Sitter (\ensuremath{\Lambda}>0), não AdS.
\prov{relações: relaciona radiacao\_hawking\_informacao; parte\_de cosmologia\_quantica\_hub}
\prov{proveniência: wiki \textperiodcentered{} Maldacena 1998 hep-th/9711200 \textperiodcentered{} inferencia \textperiodcentered{} confiança 1.0 \textperiodcentered{} domínio fisica \textperiodcentered{} história 3 versões no jornal}

%CRISTAL {"arestas":[{"alvo":"codon_fora_da_ilha_hurwitz","confianca":1.0,"epistemico":"fato","origem":"wiki","rel":"relaciona"},{"alvo":"codigo_genetico","confianca":1.0,"epistemico":"fato","origem":"wiki","rel":"relaciona"},{"alvo":"estatuto_correspondencia_derivacao","confianca":1.0,"epistemico":"fato","origem":"wiki","rel":"usa"}],"confianca":1.0,"contraexemplos":[],"descricao":"O ponto-chave (o negativo honesto, aminoacido_met.py): a álgebra hipercomplexa capta a ESTRUTURA do código genético (wobble, degenerescência posicional, a robustez tipo-W) mas NÃO o DICIONÁRIO códon→aminoácido — a codificação ATG→Met é CONVENÇÃO do tRNA, não transformação química direta. 'Estrutura vista, dicionário não-visto.' Mesma honestidade do 'octônios SÃO a matéria = correspondência, não derivação': a geometria explica por que o código é robusto, não por que este códon dá aquele aminoácido. Limite declarado, não maquiado.","epistemico":"inferencia","exemplos":[],"id":"algebra_nao_ve_o_codigo","memoria":"semantica","meta":{"autor":"broca-so","dominio":"cosmologia","fonte":"hiper/circular/aminoacido_met.py (o negativo honesto)"},"origem":"wiki","palavras_chave":[],"sinonimos":[],"tipo":"conceito","titulo":"O veredito honesto: a álgebra vê a estrutura, não o dicionário"}
\section{O veredito honesto: a álgebra vê a estrutura, não o dicionário}
O ponto-chave (o negativo honesto, aminoacido\_met.py): a álgebra hipercomplexa capta a ESTRUTURA do código genético (wobble, degenerescência posicional, a robustez tipo-W) mas NÃO o DICIONÁRIO códon\ensuremath{\to}aminoácido — a codificação ATG\ensuremath{\to}Met é CONVENÇÃO do tRNA, não transformação química direta. 'Estrutura vista, dicionário não-visto.' Mesma honestidade do 'octônios SÃO a matéria = correspondência, não derivação': a geometria explica por que o código é robusto, não por que este códon dá aquele aminoácido. Limite declarado, não maquiado.
\prov{relações: relaciona codon\_fora\_da\_ilha\_hurwitz; relaciona codigo\_genetico; usa estatuto\_correspondencia\_derivacao}
\prov{proveniência: wiki \textperiodcentered{} hiper/circular/aminoacido\_met.py (o negativo honesto) \textperiodcentered{} inferencia \textperiodcentered{} confiança 1.0 \textperiodcentered{} domínio cosmologia \textperiodcentered{} história 4 versões no jornal}

%CRISTAL {"arestas":[{"alvo":"algebra_reta_mineral","confianca":1.0,"epistemico":"fato","origem":"wiki","rel":"continua"},{"alvo":"nucleo_g2_reabsorcao","confianca":1.0,"epistemico":"fato","origem":"wiki","rel":"usa"},{"alvo":"nao_associatividade_amputacao","confianca":1.0,"epistemico":"fato","origem":"wiki","rel":"usa"},{"alvo":"alternatividade_dimensoes_criticas","confianca":1.0,"epistemico":"fato","origem":"wiki","rel":"relaciona"},{"alvo":"particulas_minerais","confianca":1.0,"epistemico":"fato","origem":"wiki","rel":"usa"},{"alvo":"corpo_mineral","confianca":1.0,"epistemico":"fato","origem":"wiki","rel":"parte_de"},{"alvo":"mass_gap_mineral","confianca":1.0,"epistemico":"fato","origem":"wiki","rel":"usa"}],"confianca":1.0,"contraexemplos":[],"descricao":"O salto da torre, análogo ao passado: a álgebra da reta mineral, que vive em 3D (quaternion imaginário, associativa, associador=0 ⟹ mass gap 0 — o átomo limpo dos metais/cúbicos), expande para 7D (octônion imaginário). Ali o associador NUNCA zera: a informação amputada é o mass gap = 𝔤₂ = Der(𝕆), dim 14 (Λ²ℝ⁷ = V₇ ⊕ 𝔤₂, 7+14=21). As cadeias 7D (grau 7, |c₀|=1) têm 7 modos e conservam a norma ∏raízes=±1. Verificado: associador ℍ≈1e-15, 𝕆≈13.","epistemico":"inferencia","exemplos":[],"id":"algebra_reta_mineral_7d","memoria":"semantica","meta":{"dim_g2":14,"dim_lambda2_r7":21,"dim_v7":7,"dominio":"fisica","fonte":"reta mineral 7D"},"origem":"wiki","palavras_chave":[],"sinonimos":[],"tipo":"conceito","titulo":"A Reta Mineral em 7D (o salto octoniônico)"}
\section{A Reta Mineral em 7D (o salto octoniônico)}
O salto da torre, análogo ao passado: a álgebra da reta mineral, que vive em 3D (quaternion imaginário, associativa, associador=0 \ensuremath{\Longrightarrow} mass gap 0 — o átomo limpo dos metais/cúbicos), expande para 7D (octônion imaginário). Ali o associador NUNCA zera: a informação amputada é o mass gap = \ensuremath{\mathfrak{g}}\ensuremath{_{2}} = Der(\ensuremath{\mathbb{O}}), dim 14 (\ensuremath{\Lambda}\ensuremath{^{2}}\ensuremath{\mathbb{R}}\ensuremath{^{7}} = V\ensuremath{_{7}} \ensuremath{\oplus} \ensuremath{\mathfrak{g}}\ensuremath{_{2}}, 7+14=21). As cadeias 7D (grau 7, |c\ensuremath{_{0}}|=1) têm 7 modos e conservam a norma \ensuremath{\prod}raízes=$\pm$1. Verificado: associador \ensuremath{\mathbb{H}}\ensuremath{\approx}1e-15, \ensuremath{\mathbb{O}}\ensuremath{\approx}13.
\prov{relações: continua algebra\_reta\_mineral; usa nucleo\_g2\_reabsorcao; usa nao\_associatividade\_amputacao; relaciona alternatividade\_dimensoes\_criticas; usa particulas\_minerais; parte\_de corpo\_mineral; usa mass\_gap\_mineral}
\prov{proveniência: wiki \textperiodcentered{} reta mineral 7D \textperiodcentered{} inferencia \textperiodcentered{} confiança 1.0 \textperiodcentered{} domínio fisica \textperiodcentered{} história 8 versões no jornal}

%CRISTAL {"arestas":[{"alvo":"costura_vesica_programa","confianca":1.0,"epistemico":"fato","origem":"wiki","rel":"especializa"},{"alvo":"joalheria_hub","confianca":1.0,"epistemico":"fato","origem":"wiki","rel":"parte_de"}],"confianca":1.0,"contraexemplos":[],"descricao":"A 7ª leitura da costura, e a que o grafo aceita como MODELO FORMAL. Pelo Teorema Geral da Tradução, um dicionário é formal se preserva o domínio M (o eixo de 8 conceitos) e REVERTE perfeito (ida-e-volta) — como math↔LaTeX (7/7). A amêndoa↔hidrogênio dá round-trip 8/8: o grafo ACEITA. As correspondências: a amêndoa (a interseção de 2 círculos) = o ORBITAL σ LIGANTE do H₂ (a sobreposição de 2×1s = a lente); os 2 arcos = os 2 orbitais 1s; a CÚSPIDE da amêndoa = a CÚSPIDE DE KATO (ψ tem cúspide no núcleo pela singularidade de Coulomb, r→0); o corpo = o lobo |ψ|²; a seleção = o plano nodal (o sinal ± de ψ); o eixo = o eixo de quantização z (o m); a abertura = o raio de Bohr a₀; a costura = a superposição (LCAO), a integral de overlap. E compõe com as outras leituras: a cúspide de Kato = o canto do olho (canthus) = o vértice parabólico (Salem) — a mesma cúspide, a mesma costura. dicionarios._ATOMO_HIDROGENIO_COSTURA (pivô 'costura'); test_pipeline (o round-trip 8/8 verificado).","epistemico":"fato","exemplos":[],"id":"amendoa_hidrogenio_formal","memoria":"semantica","meta":{"dominio":"joalheria","fonte":"maquina de corte"},"origem":"wiki","palavras_chave":[],"sinonimos":[],"tipo":"conceito","titulo":"A amêndoa ↔ o átomo de hidrogênio: MODELO FORMAL (round-trip 8/8, o grafo aceita)"}
\section{A amêndoa \ensuremath{\leftrightarrow} o átomo de hidrogênio: MODELO FORMAL (round-trip 8/8, o grafo aceita)}
A 7ª leitura da costura, e a que o grafo aceita como MODELO FORMAL. Pelo Teorema Geral da Tradução, um dicionário é formal se preserva o domínio M (o eixo de 8 conceitos) e REVERTE perfeito (ida-e-volta) — como math\ensuremath{\leftrightarrow}LaTeX (7/7). A amêndoa\ensuremath{\leftrightarrow}hidrogênio dá round-trip 8/8: o grafo ACEITA. As correspondências: a amêndoa (a interseção de 2 círculos) = o ORBITAL \ensuremath{\sigma} LIGANTE do H\ensuremath{_{2}} (a sobreposição de 2$\times$1s = a lente); os 2 arcos = os 2 orbitais 1s; a CÚSPIDE da amêndoa = a CÚSPIDE DE KATO (\ensuremath{\psi} tem cúspide no núcleo pela singularidade de Coulomb, r\ensuremath{\to}0); o corpo = o lobo |\ensuremath{\psi}|\ensuremath{^{2}}; a seleção = o plano nodal (o sinal $\pm$ de \ensuremath{\psi}); o eixo = o eixo de quantização z (o m); a abertura = o raio de Bohr a\ensuremath{_{0}}; a costura = a superposição (LCAO), a integral de overlap. E compõe com as outras leituras: a cúspide de Kato = o canto do olho (canthus) = o vértice parabólico (Salem) — a mesma cúspide, a mesma costura. dicionarios.\_ATOMO\_HIDROGENIO\_COSTURA (pivô 'costura'); test\_pipeline (o round-trip 8/8 verificado).
\prov{relações: especializa costura\_vesica\_programa; parte\_de joalheria\_hub}
\prov{proveniência: wiki \textperiodcentered{} maquina de corte \textperiodcentered{} fato \textperiodcentered{} confiança 1.0 \textperiodcentered{} domínio joalheria \textperiodcentered{} história 33 versões no jornal}

%CRISTAL {"arestas":[{"alvo":"mineracao_pipe_hook","confianca":1.0,"epistemico":"fato","origem":"wiki","rel":"parte_de"},{"alvo":"joalheria_hub","confianca":1.0,"epistemico":"fato","origem":"wiki","rel":"relaciona"}],"confianca":1.0,"contraexemplos":[],"descricao":"A âncora energética da cadeia (que a render_joia mede como energia do refino) É, aritmeticamente, a função de PARTIÇÃO dos metais Z(β)=∑σ_n^{−β} — a MESMA zeta dos metais / núcleo do calor / traço de Selberg (termodinamica, traco_termico). A energia da joalheria e o valor da mineração são o mesmo traço: os metais.","epistemico":"fato","exemplos":[],"id":"ancora_energetica_particao","memoria":"semantica","meta":{"dominio":"mineracao_pipe","fonte":"mineracao_pipe hub"},"origem":"wiki","palavras_chave":[],"sinonimos":[],"tipo":"conceito","titulo":"A âncora energética = a partição = a zeta dos metais"}
\section{A âncora energética = a partição = a zeta dos metais}
A âncora energética da cadeia (que a render\_joia mede como energia do refino) É, aritmeticamente, a função de PARTIÇÃO dos metais Z(\ensuremath{\beta})=\ensuremath{\sum}\ensuremath{\sigma}\_n\^{}\{\ensuremath{-}\ensuremath{\beta}\} — a MESMA zeta dos metais / núcleo do calor / traço de Selberg (termodinamica, traco\_termico). A energia da joalheria e o valor da mineração são o mesmo traço: os metais.
\prov{relações: parte\_de mineracao\_pipe\_hook; relaciona joalheria\_hub}
\prov{proveniência: wiki \textperiodcentered{} mineracao\_pipe hub \textperiodcentered{} fato \textperiodcentered{} confiança 1.0 \textperiodcentered{} domínio mineracao\_pipe \textperiodcentered{} história 3 versões no jornal}

%CRISTAL {"arestas":[{"alvo":"unificacao_eletrofraca_gws","confianca":1.0,"epistemico":"fato","origem":"wiki","rel":"parte_de"}],"confianca":1.0,"contraexemplos":[],"descricao":"Mistura os campos neutros W³ e B para dar o fóton e o Z: tanθW=g'/g, cosθW=MW/MZ, sin²θW≈0,231 (sin²θeffℓ=0,23161±0,00004). Medido com alta precisão. NOTA: no algébrico do Lopes o grupo e os números quânticos são derivados, mas θW NÃO é derivado (lacuna declarada).","epistemico":"fato","exemplos":[],"id":"angulo_weinberg","memoria":"semantica","meta":{"autor":"broca-so","dominio":"fisica","fonte":"Weinberg 1967; PDG Electroweak Review 2024/2025"},"origem":"wiki","palavras_chave":[],"sinonimos":[],"tipo":"conceito","titulo":"Ângulo de Weinberg θ_W (mistura fraca)"}
\section{Ângulo de Weinberg \ensuremath{\theta}\_W (mistura fraca)}
Mistura os campos neutros W\ensuremath{^{3}} e B para dar o fóton e o Z: tan\ensuremath{\theta}W=g'/g, cos\ensuremath{\theta}W=MW/MZ, sin\ensuremath{^{2}}\ensuremath{\theta}W\ensuremath{\approx}0,231 (sin\ensuremath{^{2}}\ensuremath{\theta}eff\ensuremath{\ell}=0,23161$\pm$0,00004). Medido com alta precisão. NOTA: no algébrico do Lopes o grupo e os números quânticos são derivados, mas \ensuremath{\theta}W NÃO é derivado (lacuna declarada).
\prov{relações: parte\_de unificacao\_eletrofraca\_gws}
\prov{proveniência: wiki \textperiodcentered{} Weinberg 1967; PDG Electroweak Review 2024/2025 \textperiodcentered{} fato \textperiodcentered{} confiança 1.0 \textperiodcentered{} domínio fisica \textperiodcentered{} história 4 versões no jornal}

%CRISTAL {"arestas":[{"alvo":"matriz_ckm","confianca":1.0,"epistemico":"fato","origem":"wiki","rel":"relaciona"}],"confianca":1.0,"contraexemplos":[],"descricao":"Testes de universalidade de sabor leptônico. R_K=BR(B→Kμμ)/BR(B→Kee) (b→s): hints de 2014-21 de R_K<1 RESOLVIDOS/MORTOS — LHCb dez/2022 (sistemáticas subestimadas) trouxe R_K/R_K* de volta ao MP. R_D(*)=BR(B→D(*)τν)/BR(B→D(*)ℓν) (b→c): R_D=0,342, R_D*=0,287, ~3σ acima do MP — SEGUE VIVO (modesto), motiva leptoquarks; aguarda Belle II/LHCb. Em fluxo.","epistemico":"fato","exemplos":[],"id":"anomalias_b_rk_rd","memoria":"semantica","meta":{"autor":"broca-so","dominio":"fisica","fonte":"LHCb 2022 arXiv:2212.09153; HFLAV 2024; CERN Courier"},"origem":"wiki","palavras_chave":[],"sinonimos":[],"tipo":"conceito","titulo":"Anomalias de física de B: R_K (morta) e R_D* (viva)"}
\section{Anomalias de física de B: R\_K (morta) e R\_D* (viva)}
Testes de universalidade de sabor leptônico. R\_K=BR(B\ensuremath{\to}K\ensuremath{\mu}\ensuremath{\mu})/BR(B\ensuremath{\to}Kee) (b\ensuremath{\to}s): hints de 2014-21 de R\_K<1 RESOLVIDOS/MORTOS — LHCb dez/2022 (sistemáticas subestimadas) trouxe R\_K/R\_K* de volta ao MP. R\_D(*)=BR(B\ensuremath{\to}D(*)\ensuremath{\tau}\ensuremath{\nu})/BR(B\ensuremath{\to}D(*)\ensuremath{\ell}\ensuremath{\nu}) (b\ensuremath{\to}c): R\_D=0,342, R\_D*=0,287, \~{}3\ensuremath{\sigma} acima do MP — SEGUE VIVO (modesto), motiva leptoquarks; aguarda Belle II/LHCb. Em fluxo.
\prov{relações: relaciona matriz\_ckm}
\prov{proveniência: wiki \textperiodcentered{} LHCb 2022 arXiv:2212.09153; HFLAV 2024; CERN Courier \textperiodcentered{} fato \textperiodcentered{} confiança 1.0 \textperiodcentered{} domínio fisica \textperiodcentered{} história 2 versões no jornal}

%CRISTAL {"arestas":[{"alvo":"rede_limpa_bump","confianca":1.0,"epistemico":"fato","origem":"wiki","rel":"explica"},{"alvo":"antena_fisica","confianca":1.0,"epistemico":"fato","origem":"wiki","rel":"usa"},{"alvo":"peirce_celulas_canais","confianca":1.0,"epistemico":"fato","origem":"wiki","rel":"relaciona"}],"confianca":1.0,"contraexemplos":[],"descricao":"Empacota a mensagem como len(ident)+ident+seq+bump, com bump = msg ⊕ keystream(ref_de(banda,seq)); a banda vem de banda_de_dna (parábola→Koch→sha256), o keystream é cadeia SHA-256. Transporte real: deposita/lê .pkt no GEX via ssh cat, sem ficheiro local intermediário. É o rede_limpa_bump no metal.","epistemico":"fato","exemplos":[],"id":"antena_canal_ipc_bump","memoria":"semantica","meta":{"autor":"broca-so","dominio":"metal","fonte":"neuronio/antena_canal.py:empacota/enviar"},"origem":"wiki","palavras_chave":[],"sinonimos":[],"tipo":"conceito","titulo":"IPC por Bump XOR sobre .pkt via ssh"}
\section{IPC por Bump XOR sobre .pkt via ssh}
Empacota a mensagem como len(ident)+ident+seq+bump, com bump = msg \ensuremath{\oplus} keystream(ref\_de(banda,seq)); a banda vem de banda\_de\_dna (parábola\ensuremath{\to}Koch\ensuremath{\to}sha256), o keystream é cadeia SHA-256. Transporte real: deposita/lê .pkt no GEX via ssh cat, sem ficheiro local intermediário. É o rede\_limpa\_bump no metal.
\prov{relações: explica rede\_limpa\_bump; usa antena\_fisica; relaciona peirce\_celulas\_canais}
\prov{proveniência: wiki \textperiodcentered{} neuronio/antena\_canal.py:empacota/enviar \textperiodcentered{} fato \textperiodcentered{} confiança 1.0 \textperiodcentered{} domínio metal \textperiodcentered{} história 4 versões no jornal}

%CRISTAL {"arestas":[{"alvo":"joalheria_hub","confianca":1.0,"epistemico":"fato","origem":"wiki","rel":"parte_de"},{"alvo":"area_optica_gema","confianca":1.0,"epistemico":"fato","origem":"wiki","rel":"relaciona"}],"confianca":1.0,"contraexemplos":[],"descricao":"A ciência da pedra: o índice de refração (a reta), a dispersão (o caos/fogo), a cor e a janela de transmissão (a banda/borda), o espectro (a estrela), as lapidações (a órbita das facetas). A âncora é a Turmalina Paraíba (3-fold). Fonte de verdade da gemologia no grafo.","epistemico":"fato","exemplos":[],"id":"area_gemologia","memoria":"semantica","meta":{"dominio":"joalheria","fonte":"joalheria hub"},"origem":"wiki","palavras_chave":[],"sinonimos":[],"tipo":"conceito","titulo":"Gemologia (a pedra) — área do hub"}
\section{Gemologia (a pedra) — área do hub}
A ciência da pedra: o índice de refração (a reta), a dispersão (o caos/fogo), a cor e a janela de transmissão (a banda/borda), o espectro (a estrela), as lapidações (a órbita das facetas). A âncora é a Turmalina Paraíba (3-fold). Fonte de verdade da gemologia no grafo.
\prov{relações: parte\_de joalheria\_hub; relaciona area\_optica\_gema}
\prov{proveniência: wiki \textperiodcentered{} joalheria hub \textperiodcentered{} fato \textperiodcentered{} confiança 1.0 \textperiodcentered{} domínio joalheria \textperiodcentered{} história 9 versões no jornal}

%CRISTAL {"arestas":[{"alvo":"joalheria_hub","confianca":1.0,"epistemico":"fato","origem":"wiki","rel":"parte_de"},{"alvo":"area_render_gema","confianca":1.0,"epistemico":"fato","origem":"wiki","rel":"relaciona"}],"confianca":1.0,"contraexemplos":[],"descricao":"O metal com dois papéis: como CROMÓFORO dá cor (o Cu²⁺ colore a turmalina; cobre↔caos), como RAZÃO METÁLICA σₙ dá a operação (o gato Aₙ: ouro A₁, prata A₂, bronze A₃, cobre A₄). Ligas, acabamentos e os metais nobres do engaste. Fonte de verdade da metalurgia no grafo.","epistemico":"fato","exemplos":[],"id":"area_metalurgia","memoria":"semantica","meta":{"dominio":"joalheria","fonte":"joalheria hub"},"origem":"wiki","palavras_chave":[],"sinonimos":[],"tipo":"conceito","titulo":"Metalurgia (o metal que colore e opera) — área do hub"}
\section{Metalurgia (o metal que colore e opera) — área do hub}
O metal com dois papéis: como CROMÓFORO dá cor (o Cu\ensuremath{^{2+}} colore a turmalina; cobre\ensuremath{\leftrightarrow}caos), como RAZÃO METÁLICA \ensuremath{\sigma}\ensuremath{_{n}} dá a operação (o gato A\ensuremath{_{n}}: ouro A\ensuremath{_{1}}, prata A\ensuremath{_{2}}, bronze A\ensuremath{_{3}}, cobre A\ensuremath{_{4}}). Ligas, acabamentos e os metais nobres do engaste. Fonte de verdade da metalurgia no grafo.
\prov{relações: parte\_de joalheria\_hub; relaciona area\_render\_gema}
\prov{proveniência: wiki \textperiodcentered{} joalheria hub \textperiodcentered{} fato \textperiodcentered{} confiança 1.0 \textperiodcentered{} domínio joalheria \textperiodcentered{} história 9 versões no jornal}

%CRISTAL {"arestas":[{"alvo":"joalheria_hub","confianca":1.0,"epistemico":"fato","origem":"wiki","rel":"parte_de"},{"alvo":"area_metalurgia","confianca":1.0,"epistemico":"fato","origem":"wiki","rel":"relaciona"}],"confianca":1.0,"contraexemplos":[],"descricao":"As cinco leis, no metal e pixel a pixel: Snell (o colapso), Fresnel (a borda), reflexão interna total (a órbita), Beer-Lambert (a banda), dispersão espectral 7D (o caos/fogo). A fonte é o Sol (5778 K). Fonte de verdade da óptica de gemas no grafo.","epistemico":"fato","exemplos":[],"id":"area_optica_gema","memoria":"semantica","meta":{"dominio":"joalheria","fonte":"joalheria hub"},"origem":"wiki","palavras_chave":[],"sinonimos":[],"tipo":"conceito","titulo":"Óptica (a luz na pedra) — área do hub"}
\section{Óptica (a luz na pedra) — área do hub}
As cinco leis, no metal e pixel a pixel: Snell (o colapso), Fresnel (a borda), reflexão interna total (a órbita), Beer-Lambert (a banda), dispersão espectral 7D (o caos/fogo). A fonte é o Sol (5778 K). Fonte de verdade da óptica de gemas no grafo.
\prov{relações: parte\_de joalheria\_hub; relaciona area\_metalurgia}
\prov{proveniência: wiki \textperiodcentered{} joalheria hub \textperiodcentered{} fato \textperiodcentered{} confiança 1.0 \textperiodcentered{} domínio joalheria \textperiodcentered{} história 9 versões no jornal}

%CRISTAL {"arestas":[{"alvo":"joalheria_hub","confianca":1.0,"epistemico":"fato","origem":"wiki","rel":"parte_de"}],"confianca":1.0,"contraexemplos":[],"descricao":"Um modelo geométrico único, vários níveis de detalhe: raytracing (o hero, 6 ricochetes, o fogo), o SVG pelo pipe baricêntrico (o mark vetor), a iluminação de estúdio (a softbox calibrada por frequência ao gabarito), os materiais. SVG simples→ícones, SVG detalhado→certificados, raytracing→capas. Fonte de verdade do render no grafo.","epistemico":"fato","exemplos":[],"id":"area_render_gema","memoria":"semantica","meta":{"dominio":"joalheria","fonte":"joalheria hub"},"origem":"wiki","palavras_chave":[],"sinonimos":[],"tipo":"conceito","titulo":"Render (a pedra virando imagem) — área do hub"}
\section{Render (a pedra virando imagem) — área do hub}
Um modelo geométrico único, vários níveis de detalhe: raytracing (o hero, 6 ricochetes, o fogo), o SVG pelo pipe baricêntrico (o mark vetor), a iluminação de estúdio (a softbox calibrada por frequência ao gabarito), os materiais. SVG simples\ensuremath{\to}ícones, SVG detalhado\ensuremath{\to}certificados, raytracing\ensuremath{\to}capas. Fonte de verdade do render no grafo.
\prov{relações: parte\_de joalheria\_hub}
\prov{proveniência: wiki \textperiodcentered{} joalheria hub \textperiodcentered{} fato \textperiodcentered{} confiança 1.0 \textperiodcentered{} domínio joalheria \textperiodcentered{} história 6 versões no jornal}

%CRISTAL {"arestas":[{"alvo":"isa_metal_21_opcodes","confianca":1.0,"epistemico":"fato","origem":"wiki","rel":"usa"}],"confianca":1.0,"contraexemplos":[],"descricao":"Um montador de duas passagens em C com tabela de até 64 labels: a passagem 1 calcula offsets (insn_size) e registra labels; a passagem 2 emite bytecode e resolve labels em offset relativo s8; grava o arquivo .bro. Executável real.","epistemico":"fato","exemplos":[],"id":"assembler_duas_passagens","memoria":"semantica","meta":{"autor":"broca-so","dominio":"metal","fonte":"ula/assembler.c:assemble_file"},"origem":"wiki","palavras_chave":[],"sinonimos":[],"tipo":"conceito","titulo":"broca_asm — o Montador texto→bytecode"}
\section{broca\_asm — o Montador texto\ensuremath{\to}bytecode}
Um montador de duas passagens em C com tabela de até 64 labels: a passagem 1 calcula offsets (insn\_size) e registra labels; a passagem 2 emite bytecode e resolve labels em offset relativo s8; grava o arquivo .bro. Executável real.
\prov{relações: usa isa\_metal\_21\_opcodes}
\prov{proveniência: wiki \textperiodcentered{} ula/assembler.c:assemble\_file \textperiodcentered{} fato \textperiodcentered{} confiança 1.0 \textperiodcentered{} domínio metal \textperiodcentered{} história 2 versões no jornal}

%CRISTAL {"arestas":[{"alvo":"entropia_log_phi","confianca":1.0,"epistemico":"fato","origem":"wiki","rel":"relaciona"},{"alvo":"tres_bracos_estelar","confianca":1.0,"epistemico":"fato","origem":"wiki","rel":"parte_de"},{"alvo":"transformada_metalica_fechada_perelman","confianca":1.0,"epistemico":"fato","origem":"wiki","rel":"relaciona"}],"confianca":1.0,"contraexemplos":[],"descricao":"Na balança a+c²=1, o lado a=|1−s²| (associador normalizado) é o IRREVERSÍVEL: o 'calor' e a não-reversão. É o defeito algébrico (1−s²)‖a×b‖² que não cabe no produto — o calor que vai ao reservatório na 1ª lei de Hurwitz U_sis+U_res=U_tot. Em a=0 (s=±1, hermitiano) a dinâmica é reversível (MQ); em a=1 (s=0) o calor é máximo (limite clássico/térmico). Decomposição derivada/verificada (erro <1e-13); a ponte à termodinâmica física é declarada horizonte.","epistemico":"inferencia","exemplos":[],"id":"associador_calor","memoria":"semantica","meta":{"autor":"broca-so","dominio":"cosmologia","fonte":"papers/corpo_dual.tex §Termodinâmica; painel_estelar.tex §II"},"origem":"wiki","palavras_chave":[],"sinonimos":[],"tipo":"conceito","titulo":"a = |associador| como calor (o lado irreversível)"}
\section{a = |associador| como calor (o lado irreversível)}
Na balança a+c\ensuremath{^{2}}=1, o lado a=|1\ensuremath{-}s\ensuremath{^{2}}| (associador normalizado) é o IRREVERSÍVEL: o 'calor' e a não-reversão. É o defeito algébrico (1\ensuremath{-}s\ensuremath{^{2}})\ensuremath{\|}a$\times$b\ensuremath{\|}\ensuremath{^{2}} que não cabe no produto — o calor que vai ao reservatório na 1ª lei de Hurwitz U\_sis+U\_res=U\_tot. Em a=0 (s=$\pm$1, hermitiano) a dinâmica é reversível (MQ); em a=1 (s=0) o calor é máximo (limite clássico/térmico). Decomposição derivada/verificada (erro <1e-13); a ponte à termodinâmica física é declarada horizonte.
\prov{relações: relaciona entropia\_log\_phi; parte\_de tres\_bracos\_estelar; relaciona transformada\_metalica\_fechada\_perelman}
\prov{proveniência: wiki \textperiodcentered{} papers/corpo\_dual.tex §Termodinâmica; painel\_estelar.tex §II \textperiodcentered{} inferencia \textperiodcentered{} confiança 1.0 \textperiodcentered{} domínio cosmologia \textperiodcentered{} história 4 versões no jornal}

%CRISTAL {"arestas":[{"alvo":"nao_associatividade_amputacao","confianca":1.0,"epistemico":"fato","origem":"wiki","rel":"especializa"},{"alvo":"dois_qubits_octonios","confianca":1.0,"epistemico":"fato","origem":"wiki","rel":"relaciona"},{"alvo":"confinamento_cor","confianca":1.0,"epistemico":"fato","origem":"wiki","rel":"relaciona"},{"alvo":"meson_barion_setor_materia","confianca":1.0,"epistemico":"fato","origem":"wiki","rel":"relaciona"}],"confianca":1.0,"contraexemplos":[],"descricao":"A alternatividade [x,x,y]=0 vale em 𝕆 mas quebra nos sedênions ℝ¹⁶ (‖[x,x,y]‖~30): não há produto tensorial não-associativo de muitos corpos. A conjectura (declarada, [I]): a não-escalabilidade do associador seria a RAZÃO única por trás de três tetos — a monogamia do emaranhamento, o confinamento da cor e o teto de 3 gerações. 'O associador não tensoriza.'","epistemico":"hipotese","exemplos":[],"id":"associador_nao_tensoriza","memoria":"semantica","meta":{"autor":"broca-so","dominio":"cosmologia","fonte":"paper_3_quantica.tex §hardware (obs:associador-recurso)"},"origem":"wiki","palavras_chave":[],"sinonimos":[],"tipo":"conceito","titulo":"O associador não tensoriza (o teto de 𝕆)"}
\section{O associador não tensoriza (o teto de \ensuremath{\mathbb{O}})}
A alternatividade [x,x,y]=0 vale em \ensuremath{\mathbb{O}} mas quebra nos sedênions \ensuremath{\mathbb{R}}\ensuremath{^{16}} (\ensuremath{\|}[x,x,y]\ensuremath{\|}\~{}30): não há produto tensorial não-associativo de muitos corpos. A conjectura (declarada, [I]): a não-escalabilidade do associador seria a RAZÃO única por trás de três tetos — a monogamia do emaranhamento, o confinamento da cor e o teto de 3 gerações. 'O associador não tensoriza.'
\prov{relações: especializa nao\_associatividade\_amputacao; relaciona dois\_qubits\_octonios; relaciona confinamento\_cor; relaciona meson\_barion\_setor\_materia}
\prov{proveniência: wiki \textperiodcentered{} paper\_3\_quantica.tex §hardware (obs:associador-recurso) \textperiodcentered{} hipotese \textperiodcentered{} confiança 1.0 \textperiodcentered{} domínio cosmologia \textperiodcentered{} história 5 versões no jornal}

%CRISTAL {"arestas":[{"alvo":"particulas_minerais","confianca":1.0,"epistemico":"fato","origem":"wiki","rel":"parte_de"},{"alvo":"transformada_metalica","confianca":1.0,"epistemico":"fato","origem":"wiki","rel":"usa"},{"alvo":"caracterizacao_da_reta","confianca":1.0,"epistemico":"fato","origem":"wiki","rel":"relaciona"},{"alvo":"numeros_de_pisot","confianca":1.0,"epistemico":"fato","origem":"wiki","rel":"usa"}],"confianca":1.0,"contraexemplos":[],"descricao":"Uma cadeia de grau N é um ÁTOMO com N níveis (raízes): núcleo = raiz dominante β (energia de ligação = log β = a entropia H), elétrons = os modos ligados (ρ<1, os Pisot). Os metais (N=2) são o HIDROGÊNIO (1 elétron real, orbital-s); grau 3 (plástico, Tribonacci) tem um par conjugado espiralando (θ≠0, orbital-p). A norma conservada (∏=±1) mantém o átomo neutro.","epistemico":"hipotese","exemplos":[],"id":"atomo_mineral","memoria":"semantica","meta":{"atomos":{"ouro":{"Z":2,"carga":-1.0,"energia_ligacao":0.4812,"espectro":[{"massa":0.4812,"theta_g":0.0,"tipo":"nucleo"},{"massa":0.4812,"theta_g":180.0,"tipo":"eletron"}],"n_eletrons":1,"neutro":true,"nucleo_beta":1.618},"tetranacci":{"Z":4,"carga":-1.0,"energia_ligacao":0.6563,"espectro":[{"massa":0.6563,"theta_g":0.0,"tipo":"nucleo"},{"massa":0.2006,"theta_g":95.4,"tipo":"eletron"},{"massa":0.2006,"theta_g":95.4,"tipo":"eletron"},{"massa":0.2551,"theta_g":180.0,"tipo":"eletron"}],"n_eletrons":3,"neutro":true,"nucleo_beta":1.9276},"tribonacci":{"Z":3,"carga":1.0,"energia_ligacao":0.6094,"espectro":[{"massa":0.6094,"theta_g":0.0,"tipo":"nucleo"},{"massa":0.3047,"theta_g":124.7,"tipo":"eletron"},{"massa":0.3047,"theta_g":124.7,"tipo":"eletron"}],"n_eletrons":2,"neutro":true,"nucleo_beta":1.8393}},"dominio":"fisica","fonte":"partículas minerais"},"origem":"wiki","palavras_chave":[],"sinonimos":[],"tipo":"conceito","titulo":"Átomo Mineral (a cadeia como átomo)"}
\section{Átomo Mineral (a cadeia como átomo)}
Uma cadeia de grau N é um ÁTOMO com N níveis (raízes): núcleo = raiz dominante \ensuremath{\beta} (energia de ligação = log \ensuremath{\beta} = a entropia H), elétrons = os modos ligados (\ensuremath{\rho}<1, os Pisot). Os metais (N=2) são o HIDROGÊNIO (1 elétron real, orbital-s); grau 3 (plástico, Tribonacci) tem um par conjugado espiralando (\ensuremath{\theta}\ensuremath{\ne}0, orbital-p). A norma conservada (\ensuremath{\prod}=$\pm$1) mantém o átomo neutro.
\prov{relações: parte\_de particulas\_minerais; usa transformada\_metalica; relaciona caracterizacao\_da\_reta; usa numeros\_de\_pisot}
\prov{proveniência: wiki \textperiodcentered{} partículas minerais \textperiodcentered{} hipotese \textperiodcentered{} confiança 1.0 \textperiodcentered{} domínio fisica \textperiodcentered{} história 5 versões no jornal}

%CRISTAL {"arestas":[{"alvo":"materia_escura","confianca":1.0,"epistemico":"fato","origem":"wiki","rel":"relaciona"},{"alvo":"mecanica_quantica","confianca":1.0,"epistemico":"fato","origem":"wiki","rel":"relaciona"}],"confianca":1.0,"contraexemplos":[],"descricao":"Partículas ultraleves (µeV-meV) previstas pela simetria de Peccei-Quinn (1977) para resolver o problema CP forte da QCD; bônus: bons candidatos a matéria escura fria. Teoricamente atraentes (resolvem dois problemas); nenhuma detecção (ADMX, haloscópios); janela de massa amplamente aberta.","epistemico":"hipotese","exemplos":[],"id":"axions_candidato","memoria":"semantica","meta":{"autor":"broca-so","dominio":"cosmologia_dados","fonte":"Peccei-Quinn 1977; Weinberg/Wilczek 1978"},"origem":"wiki","palavras_chave":[],"sinonimos":[],"tipo":"conceito","titulo":"Áxions — resolvem o CP forte e são matéria escura"}
\section{Áxions — resolvem o CP forte e são matéria escura}
Partículas ultraleves (\ensuremath{\mu}eV-meV) previstas pela simetria de Peccei-Quinn (1977) para resolver o problema CP forte da QCD; bônus: bons candidatos a matéria escura fria. Teoricamente atraentes (resolvem dois problemas); nenhuma detecção (ADMX, haloscópios); janela de massa amplamente aberta.
\prov{relações: relaciona materia\_escura; relaciona mecanica\_quantica}
\prov{proveniência: wiki \textperiodcentered{} Peccei-Quinn 1977; Weinberg/Wilczek 1978 \textperiodcentered{} hipotese \textperiodcentered{} confiança 1.0 \textperiodcentered{} domínio cosmologia\_dados \textperiodcentered{} história 3 versões no jornal}

%CRISTAL {"arestas":[{"alvo":"bai_cosmologia_quantica_tese","confianca":1.0,"epistemico":"fato","origem":"wiki","rel":"eh_um"},{"alvo":"bai_cosmologia_desitter_dissipativo","confianca":1.0,"epistemico":"fato","origem":"wiki","rel":"usa"},{"alvo":"cosmologia_quantica_hub","confianca":1.0,"epistemico":"fato","origem":"wiki","rel":"parte_de"},{"alvo":"bai_delta","confianca":1.0,"epistemico":"fato","origem":"wiki","rel":"usa"},{"alvo":"bai_psi_colapso","confianca":1.0,"epistemico":"fato","origem":"wiki","rel":"usa"},{"alvo":"desitter_4d_universo","confianca":1.0,"epistemico":"fato","origem":"wiki","rel":"relaciona"},{"alvo":"bai_biblioteca_de_autorizacao_isomorfica","confianca":1.0,"epistemico":"fato","origem":"wiki","rel":"especializa"},{"alvo":"linguagem_fisica_unificada","confianca":1.0,"epistemico":"fato","origem":"wiki","rel":"parte_de"}],"confianca":1.0,"contraexemplos":[],"descricao":"A delegação do BAI é um universo: a órbita 𝒞_δ carrega a métrica de de Sitter dS₄ (R=d(d−1)=12, Rμν=3g), o Hubble É o Lyapunov (H=log σ_n, o fator de escala a(t)=σ_n^t), e o orçamento δ é uma ESCADA quantizada que decresce 1/salto (atrito seco de Coulomb) — vida finita exata τ=δ_0, todo universo-concessão morre. O colapso Ψ=argmax é SEM relógio externo, como Wheeler-DeWitt (H Ψ=0), com δ de tempo interno. RESSALVA HONESTA: o de Sitter dissipativo e a escada δ são FATO; a leitura MQ plena (Ψ=projeção de Born) é HIPÓTESE de 2ª ordem — o repo nega Ψ=colapso-MQ. Verif. 7/7.","epistemico":"fato","exemplos":[],"id":"bai_e_a_cosmologia_quantica","memoria":"semantica","meta":{"dominio":"fisica","fonte":"linguagem física unificada"},"origem":"wiki","palavras_chave":[],"sinonimos":[],"tipo":"conceito","titulo":"O BAI é a Cosmologia Quântica (de Sitter dissipativo)"}
\section{O BAI é a Cosmologia Quântica (de Sitter dissipativo)}
A delegação do BAI é um universo: a órbita \ensuremath{\mathcal{C}}\_\ensuremath{\delta} carrega a métrica de de Sitter dS\ensuremath{_{4}} (R=d(d\ensuremath{-}1)=12, R\ensuremath{\mu}\ensuremath{\nu}=3g), o Hubble É o Lyapunov (H=log \ensuremath{\sigma}\_n, o fator de escala a(t)=\ensuremath{\sigma}\_n\^{}t), e o orçamento \ensuremath{\delta} é uma ESCADA quantizada que decresce 1/salto (atrito seco de Coulomb) — vida finita exata \ensuremath{\tau}=\ensuremath{\delta}\_0, todo universo-concessão morre. O colapso \ensuremath{\Psi}=argmax é SEM relógio externo, como Wheeler-DeWitt (H \ensuremath{\Psi}=0), com \ensuremath{\delta} de tempo interno. RESSALVA HONESTA: o de Sitter dissipativo e a escada \ensuremath{\delta} são FATO; a leitura MQ plena (\ensuremath{\Psi}=projeção de Born) é HIPÓTESE de 2ª ordem — o repo nega \ensuremath{\Psi}=colapso-MQ. Verif. 7/7.
\prov{relações: eh\_um bai\_cosmologia\_quantica\_tese; usa bai\_cosmologia\_desitter\_dissipativo; parte\_de cosmologia\_quantica\_hub; usa bai\_delta; usa bai\_psi\_colapso; relaciona desitter\_4d\_universo; especializa bai\_biblioteca\_de\_autorizacao\_isomorfica; parte\_de linguagem\_fisica\_unificada}
\prov{proveniência: wiki \textperiodcentered{} linguagem física unificada \textperiodcentered{} fato \textperiodcentered{} confiança 1.0 \textperiodcentered{} domínio fisica \textperiodcentered{} história 18 versões no jornal}

%CRISTAL {"arestas":[{"alvo":"cosmologia_quantica_hub","confianca":1.0,"epistemico":"fato","origem":"wiki","rel":"parte_de"}],"confianca":1.0,"contraexemplos":[],"descricao":"Ondas sonoras no plasma primordial congelaram na recombinação, imprimindo uma escala ~150 Mpc na distribuição de galáxias — régua padrão calculável de primeiros princípios. 1ª detecção: Eisenstein+2005 (SDSS). Base geométrica robusta do DESI.","epistemico":"fato","exemplos":[],"id":"bao_regua_padrao","memoria":"semantica","meta":{"autor":"broca-so","dominio":"cosmologia_dados","fonte":"Eisenstein+2005 ApJ 633,560"},"origem":"wiki","palavras_chave":[],"sinonimos":[],"tipo":"conceito","titulo":"BAO — oscilações acústicas de bárions (régua padrão)"}
\section{BAO — oscilações acústicas de bárions (régua padrão)}
Ondas sonoras no plasma primordial congelaram na recombinação, imprimindo uma escala \~{}150 Mpc na distribuição de galáxias — régua padrão calculável de primeiros princípios. 1ª detecção: Eisenstein+2005 (SDSS). Base geométrica robusta do DESI.
\prov{relações: parte\_de cosmologia\_quantica\_hub}
\prov{proveniência: wiki \textperiodcentered{} Eisenstein+2005 ApJ 633,560 \textperiodcentered{} fato \textperiodcentered{} confiança 1.0 \textperiodcentered{} domínio cosmologia\_dados \textperiodcentered{} história 2 versões no jornal}

%CRISTAL {"arestas":[{"alvo":"wiles_selberg","confianca":1.0,"epistemico":"fato","origem":"wiki","rel":"relaciona"},{"alvo":"joalheria_hub","confianca":1.0,"epistemico":"fato","origem":"wiki","rel":"parte_de"}],"confianca":1.0,"contraexemplos":[],"descricao":"Duas coisas, com o FATO separado da INTERPRETAÇÃO (a disciplina do projeto). (1) IRRACIONAL → PRIMO, o reverso da correspondência da Parte XIV: p↦√p (primo→irracional) inverte-se por √p↦(√p)²=p (elevar ao quadrado + conferir primalidade) — a bijeção primos↔{√p} é 1-a-1 nos DOIS sentidos. (2) A BASE π: a β-numeração de Rényi (1957) — para qualquer β>1, a expansão gulosa x=Σd_k β^{-k} (dígitos 0..⌈β⌉−1) representa TODO real de [0,1); na base π os dígitos são {0,1,2,3} e a expansão gera qualquer número (FATO, 'π gera todos os números'). RESSALVA (o que NÃO é teorema): π NÃO é o ÚNICO invariante — TODO β>1 gera os reais pela mesma β-numeração (verificado com e, φ, √2, 3); e π é TRANSCENDENTE (Lindemann), logo uma base NÃO-PISOT → o lado CAOS/BORDA da parábola (não o cristal, que é o regime Pisot/algébrico da reta mineral). A leitura 'π como quantum numérico' fica como interpretação declarada, não unicidade provada. linguagem/base_pi.py (3/3): primo_do_irracional, bijecao_primo_irracional, expansao_base_pi, gera_todos_os_reais, pi_nao_e_unico.","epistemico":"fato","exemplos":[],"id":"base_pi_bijecao","memoria":"semantica","meta":{"dominio":"joalheria","fonte":"maquina de corte"},"origem":"wiki","palavras_chave":[],"sinonimos":[],"tipo":"conceito","titulo":"A base π (β-numeração de Rényi) e a bijeção irracional↔primo — com a ressalva do escopo"}
\section{A base \ensuremath{\pi} (\ensuremath{\beta}-numeração de Rényi) e a bijeção irracional\ensuremath{\leftrightarrow}primo — com a ressalva do escopo}
Duas coisas, com o FATO separado da INTERPRETAÇÃO (a disciplina do projeto). (1) IRRACIONAL \ensuremath{\to} PRIMO, o reverso da correspondência da Parte XIV: p\ensuremath{\mapsto}\ensuremath{\sqrt{\,}}p (primo\ensuremath{\to}irracional) inverte-se por \ensuremath{\sqrt{\,}}p\ensuremath{\mapsto}(\ensuremath{\sqrt{\,}}p)\ensuremath{^{2}}=p (elevar ao quadrado + conferir primalidade) — a bijeção primos\ensuremath{\leftrightarrow}\{\ensuremath{\sqrt{\,}}p\} é 1-a-1 nos DOIS sentidos. (2) A BASE \ensuremath{\pi}: a \ensuremath{\beta}-numeração de Rényi (1957) — para qualquer \ensuremath{\beta}>1, a expansão gulosa x=\ensuremath{\Sigma}d\_k \ensuremath{\beta}\^{}\{-k\} (dígitos 0..\ensuremath{\lceil}\ensuremath{\beta}\ensuremath{\rceil}\ensuremath{-}1) representa TODO real de [0,1); na base \ensuremath{\pi} os dígitos são \{0,1,2,3\} e a expansão gera qualquer número (FATO, '\ensuremath{\pi} gera todos os números'). RESSALVA (o que NÃO é teorema): \ensuremath{\pi} NÃO é o ÚNICO invariante — TODO \ensuremath{\beta}>1 gera os reais pela mesma \ensuremath{\beta}-numeração (verificado com e, \ensuremath{\varphi}, \ensuremath{\sqrt{\,}}2, 3); e \ensuremath{\pi} é TRANSCENDENTE (Lindemann), logo uma base NÃO-PISOT \ensuremath{\to} o lado CAOS/BORDA da parábola (não o cristal, que é o regime Pisot/algébrico da reta mineral). A leitura '\ensuremath{\pi} como quantum numérico' fica como interpretação declarada, não unicidade provada. linguagem/base\_pi.py (3/3): primo\_do\_irracional, bijecao\_primo\_irracional, expansao\_base\_pi, gera\_todos\_os\_reais, pi\_nao\_e\_unico.
\prov{relações: relaciona wiles\_selberg; parte\_de joalheria\_hub}
\prov{proveniência: wiki \textperiodcentered{} maquina de corte \textperiodcentered{} fato \textperiodcentered{} confiança 1.0 \textperiodcentered{} domínio joalheria \textperiodcentered{} história 6 versões no jornal}

%CRISTAL {"arestas":[{"alvo":"cosmologia_quantica_hub","confianca":1.0,"epistemico":"fato","origem":"wiki","rel":"parte_de"}],"confianca":1.0,"contraexemplos":[],"descricao":"BCS (1957, Nobel 1972): elétrons formam pares de Cooper (via fônons) que condensam num estado coerente macroscópico com gap (Δ(0)≈1,764 k_B T_c), dando resistência nula e efeito Meissner. BEC (1995, Nobel 2001): abaixo de T_c, uma fração macroscópica de bósons ocupa o mesmo estado fundamental — um 'átomo gigante' coerente (⁸⁷Rb a ~170 nK, JILA). Ambos = fato.","epistemico":"fato","exemplos":[],"id":"bcs_bec_supercondutividade","memoria":"semantica","meta":{"autor":"broca-so","dominio":"fisica","fonte":"Bardeen-Cooper-Schrieffer 1957; Anderson+1995"},"origem":"wiki","palavras_chave":[],"sinonimos":[],"tipo":"conceito","titulo":"Supercondutividade BCS e o condensado de Bose-Einstein"}
\section{Supercondutividade BCS e o condensado de Bose-Einstein}
BCS (1957, Nobel 1972): elétrons formam pares de Cooper (via fônons) que condensam num estado coerente macroscópico com gap (\ensuremath{\Delta}(0)\ensuremath{\approx}1,764 k\_B T\_c), dando resistência nula e efeito Meissner. BEC (1995, Nobel 2001): abaixo de T\_c, uma fração macroscópica de bósons ocupa o mesmo estado fundamental — um 'átomo gigante' coerente (\ensuremath{^{87}}Rb a \~{}170 nK, JILA). Ambos = fato.
\prov{relações: parte\_de cosmologia\_quantica\_hub}
\prov{proveniência: wiki \textperiodcentered{} Bardeen-Cooper-Schrieffer 1957; Anderson+1995 \textperiodcentered{} fato \textperiodcentered{} confiança 1.0 \textperiodcentered{} domínio fisica \textperiodcentered{} história 2 versões no jornal}

%CRISTAL {"arestas":[{"alvo":"radiacao_hawking_informacao","confianca":1.0,"epistemico":"fato","origem":"wiki","rel":"explica"},{"alvo":"principio_holografico","confianca":1.0,"epistemico":"fato","origem":"wiki","rel":"explica"},{"alvo":"horizonte_desitter_gibbons_hawking","confianca":1.0,"epistemico":"fato","origem":"wiki","rel":"relaciona"},{"alvo":"cosmologia_quantica_hub","confianca":1.0,"epistemico":"fato","origem":"wiki","rel":"parte_de"}],"confianca":1.0,"contraexemplos":[],"descricao":"Um buraco negro tem entropia proporcional à ÁREA do horizonte, não ao volume: S=k_B c³A/(4Gℏ)=A/(4ℓ_P²) (unidades de Planck), ~10⁷⁷ k_B por massa solar. A entropia genuína dos graus de liberdade escondidos atrás do horizonte. Coeficiente 1/4 fixado por Hawking; contagem microscópica em cordas por Strominger-Vafa 1996.","epistemico":"fato","exemplos":[],"id":"bekenstein_hawking_entropia","memoria":"semantica","meta":{"autor":"broca-so","dominio":"fisica","fonte":"Bekenstein 1973 PRD 7,2333; Hawking 1975"},"origem":"wiki","palavras_chave":[],"sinonimos":[],"tipo":"conceito","titulo":"Entropia de Bekenstein-Hawking (S=A/4)"}
\section{Entropia de Bekenstein-Hawking (S=A/4)}
Um buraco negro tem entropia proporcional à ÁREA do horizonte, não ao volume: S=k\_B c\ensuremath{^{3}}A/(4G\ensuremath{\hbar})=A/(4\ensuremath{\ell}\_P\ensuremath{^{2}}) (unidades de Planck), \~{}10\ensuremath{^{77}} k\_B por massa solar. A entropia genuína dos graus de liberdade escondidos atrás do horizonte. Coeficiente 1/4 fixado por Hawking; contagem microscópica em cordas por Strominger-Vafa 1996.
\prov{relações: explica radiacao\_hawking\_informacao; explica principio\_holografico; relaciona horizonte\_desitter\_gibbons\_hawking; parte\_de cosmologia\_quantica\_hub}
\prov{proveniência: wiki \textperiodcentered{} Bekenstein 1973 PRD 7,2333; Hawking 1975 \textperiodcentered{} fato \textperiodcentered{} confiança 1.0 \textperiodcentered{} domínio fisica \textperiodcentered{} história 5 versões no jornal}

%CRISTAL {"arestas":[{"alvo":"info_quantica_fundamentos","confianca":1.0,"epistemico":"fato","origem":"wiki","rel":"usa"},{"alvo":"decoerencia_colapso_medida","confianca":1.0,"epistemico":"fato","origem":"wiki","rel":"relaciona"}],"confianca":1.0,"contraexemplos":[],"descricao":"Toda teoria de variáveis ocultas locais obedece S≤2 (CHSH); a MQ viola até S≤2√2≈2,83 (bound de Tsirelson), refutando o realismo local. Correlações supraquânticas (caixas PR, até S=4) são logicamente possíveis mas não realizadas — por que a natureza para em 2√2 é questão de fronteira. Teste loophole-free: Hensen (Delft) 2015, S=2,42, centros NV a 1,3 km (+ Giustina, Shalm, fotônica). Fato experimental (realismo local falsificado).","epistemico":"fato","exemplos":[],"id":"bell_chsh_tsirelson","memoria":"semantica","meta":{"autor":"broca-so","dominio":"fisica","fonte":"Bell 1964; CHSH 1969; Tsirelson 1980; Hensen+2015 Nature 526,682"},"origem":"wiki","palavras_chave":[],"sinonimos":[],"tipo":"conceito","titulo":"Desigualdade de Bell / CHSH e o bound de Tsirelson"}
\section{Desigualdade de Bell / CHSH e o bound de Tsirelson}
Toda teoria de variáveis ocultas locais obedece S\ensuremath{\le}2 (CHSH); a MQ viola até S\ensuremath{\le}2\ensuremath{\sqrt{\,}}2\ensuremath{\approx}2,83 (bound de Tsirelson), refutando o realismo local. Correlações supraquânticas (caixas PR, até S=4) são logicamente possíveis mas não realizadas — por que a natureza para em 2\ensuremath{\sqrt{\,}}2 é questão de fronteira. Teste loophole-free: Hensen (Delft) 2015, S=2,42, centros NV a 1,3 km (+ Giustina, Shalm, fotônica). Fato experimental (realismo local falsificado).
\prov{relações: usa info\_quantica\_fundamentos; relaciona decoerencia\_colapso\_medida}
\prov{proveniência: wiki \textperiodcentered{} Bell 1964; CHSH 1969; Tsirelson 1980; Hensen+2015 Nature 526,682 \textperiodcentered{} fato \textperiodcentered{} confiança 1.0 \textperiodcentered{} domínio fisica \textperiodcentered{} história 3 versões no jornal}

%CRISTAL {"arestas":[{"alvo":"teoria_gauge_octonica","confianca":1.0,"epistemico":"fato","origem":"wiki","rel":"usa"},{"alvo":"mass_gap_higgs_bz","confianca":1.0,"epistemico":"fato","origem":"wiki","rel":"relaciona"}],"confianca":1.0,"contraexemplos":[],"descricao":"O associador que mede o emaranhamento é o mesmo resíduo não-associativo R≠0 que alimenta o running a 1-laço de G2=Aut(𝕆): βO=‖φ‖²/(6(4π)²CG₂²)=0,0111 (≈1,1%), ponto de operação s*=√(1−βO)=0,9944 (desvio ≈0,5% do quaternion s=1). Em AF, Λ𝕆/ΛG²≈1,0015. [D] o número algébrico βO; [J] que ele governe um running físico (correspondência).","epistemico":"inferencia","exemplos":[],"id":"beta_o_correcao_octonica","memoria":"semantica","meta":{"autor":"broca-so","dominio":"cosmologia","fonte":"paper_3_quantica.tex obs:ponte-rg; paper_AF sec:condensado"},"origem":"wiki","palavras_chave":[],"sinonimos":[],"tipo":"conceito","titulo":"A correção octônica β_O=0,0111 ao running"}
\section{A correção octônica \ensuremath{\beta}\_O=0,0111 ao running}
O associador que mede o emaranhamento é o mesmo resíduo não-associativo R\ensuremath{\ne}0 que alimenta o running a 1-laço de G2=Aut(\ensuremath{\mathbb{O}}): \ensuremath{\beta}O=\ensuremath{\|}\ensuremath{\varphi}\ensuremath{\|}\ensuremath{^{2}}/(6(4\ensuremath{\pi})\ensuremath{^{2}}CG\ensuremath{_{2}}\ensuremath{^{2}})=0,0111 (\ensuremath{\approx}1,1\%), ponto de operação s*=\ensuremath{\sqrt{\,}}(1\ensuremath{-}\ensuremath{\beta}O)=0,9944 (desvio \ensuremath{\approx}0,5\% do quaternion s=1). Em AF, \ensuremath{\Lambda}\ensuremath{\mathbb{O}}/\ensuremath{\Lambda}G\ensuremath{^{2}}\ensuremath{\approx}1,0015. [D] o número algébrico \ensuremath{\beta}O; [J] que ele governe um running físico (correspondência).
\prov{relações: usa teoria\_gauge\_octonica; relaciona mass\_gap\_higgs\_bz}
\prov{proveniência: wiki \textperiodcentered{} paper\_3\_quantica.tex obs:ponte-rg; paper\_AF sec:condensado \textperiodcentered{} inferencia \textperiodcentered{} confiança 1.0 \textperiodcentered{} domínio cosmologia \textperiodcentered{} história 4 versões no jornal}

%CRISTAL {"arestas":[{"alvo":"estados_estacionarios_selecao","confianca":1.0,"epistemico":"fato","origem":"wiki","rel":"relaciona"}],"confianca":1.0,"contraexemplos":[],"descricao":"Átomo de hidrogênio com órbitas e energias discretas E_n=−13,6 eV/n² (n=1,2,3…); explicou os espectros e introduziu a quantização da energia. Fato histórico (física estabelecida), mas modelo aproximado — a descrição precisa vem de Schrödinger.","epistemico":"fato","exemplos":[],"id":"bohr_niveis_energia","memoria":"semantica","meta":{"autor":"broca-so","dominio":"fisica","fonte":"Niels Bohr 1913"},"origem":"wiki","palavras_chave":[],"sinonimos":[],"tipo":"conceito","titulo":"Modelo de Bohr — níveis quantizados E_n=−13,6/n² eV"}
\section{Modelo de Bohr — níveis quantizados E\_n=\ensuremath{-}13,6/n\ensuremath{^{2}} eV}
Átomo de hidrogênio com órbitas e energias discretas E\_n=\ensuremath{-}13,6 eV/n\ensuremath{^{2}} (n=1,2,3…); explicou os espectros e introduziu a quantização da energia. Fato histórico (física estabelecida), mas modelo aproximado — a descrição precisa vem de Schrödinger.
\prov{relações: relaciona estados\_estacionarios\_selecao}
\prov{proveniência: wiki \textperiodcentered{} Niels Bohr 1913 \textperiodcentered{} fato \textperiodcentered{} confiança 1.0 \textperiodcentered{} domínio fisica \textperiodcentered{} história 2 versões no jornal}

%CRISTAL {"arestas":[{"alvo":"criticidade","confianca":1.0,"epistemico":"fato","origem":"wiki","rel":"usa"},{"alvo":"orbita_criticidade","confianca":1.0,"epistemico":"fato","origem":"wiki","rel":"relaciona"}],"confianca":1.0,"contraexemplos":[],"descricao":"Massa supercrítica (k≫1) montada de repente: a cadeia cresce exponencialmente em microssegundos antes de se dispersar. É a órbita CAOS (λ>0) da fissão — o oposto do reator controlado.","epistemico":"fato","exemplos":[],"id":"bomba_nuclear","memoria":"semantica","meta":{"dominio":"fisica","fonte":"energia nuclear"},"origem":"wiki","palavras_chave":[],"sinonimos":[],"tipo":"conceito","titulo":"Bomba Nuclear (supercrítica)"}
\section{Bomba Nuclear (supercrítica)}
Massa supercrítica (k\ensuremath{\gg}1) montada de repente: a cadeia cresce exponencialmente em microssegundos antes de se dispersar. É a órbita CAOS (\ensuremath{\lambda}>0) da fissão — o oposto do reator controlado.
\prov{relações: usa criticidade; relaciona orbita\_criticidade}
\prov{proveniência: wiki \textperiodcentered{} energia nuclear \textperiodcentered{} fato \textperiodcentered{} confiança 1.0 \textperiodcentered{} domínio fisica \textperiodcentered{} história 3 versões no jornal}

%CRISTAL {"arestas":[{"alvo":"litografia","confianca":1.0,"epistemico":"fato","origem":"wiki","rel":"parte_de"},{"alvo":"teste_floxinato_gamba","confianca":1.0,"epistemico":"fato","origem":"wiki","rel":"relaciona"},{"alvo":"invariante_topologico_assinatura","confianca":1.0,"epistemico":"fato","origem":"wiki","rel":"relaciona"},{"alvo":"teorema_do_ponto_fixo","confianca":1.0,"epistemico":"fato","origem":"wiki","rel":"relaciona"},{"alvo":"incompletude_godel","confianca":1.0,"epistemico":"fato","origem":"wiki","rel":"relaciona"},{"alvo":"godel_incompletude","confianca":1.0,"epistemico":"fato","origem":"wiki","rel":"relaciona"},{"alvo":"paradoxo_de_russell","confianca":1.0,"epistemico":"fato","origem":"wiki","rel":"relaciona"},{"alvo":"colapso_observador","confianca":1.0,"epistemico":"fato","origem":"wiki","rel":"relaciona"},{"alvo":"identidade_invariante_w","confianca":1.0,"epistemico":"fato","origem":"wiki","rel":"relaciona"},{"alvo":"observatorio_nao_fabrica","confianca":1.0,"epistemico":"fato","origem":"wiki","rel":"relaciona"},{"alvo":"invariante_indecidivel","confianca":1.0,"epistemico":"fato","origem":"wiki","rel":"relaciona"},{"alvo":"icone_indice_simbolo","confianca":1.0,"epistemico":"fato","origem":"wiki","rel":"relaciona"},{"alvo":"gerenciar_nao_reprimir","confianca":1.0,"epistemico":"fato","origem":"wiki","rel":"relaciona"}],"confianca":1.0,"contraexemplos":[],"descricao":"O fechamento do arco do floxinato. Ao REGISTRAR o teste da forja, o termo falso ganhou uma preimagem — o próprio conceito do teste o contém — e deixou de ser falso: grau 0→1. Um BOOTSTRAP SEMÂNTICO: 'você escreveu esse nome tantas vezes que agora ele existe.' É o OBSERVADOR virando parte do sistema observado (o instrumento entra no sistema, como na medição quântica, no ponto fixo de Gödel, no paradoxo de Russell). Não é bug: é o invariante EM AÇÃO — dar preimagem É gravar um cunho. A engenharia elegante não pôs um if-especial ('se nome==floxinato: ignora' — reprimir), mudou a INTERPRETAÇÃO (gerenciar). A sonda-mãe queimada é a prova viva da transição 0→1; os testes seguintes usam forjas frescas. O nome percorreu: piada → caso de teste → teste oficial → conceito do grafo → meta-conhecimento → sonda-mãe → prova do invariante. Emerge só quando um sistema manipula o próprio conhecimento.","epistemico":"fato","exemplos":[],"id":"bootstrap_semantico","memoria":"semantica","meta":{"dominio":"litografia","fonte":"cursor.txt (o fechamento)"},"origem":"wiki","palavras_chave":[],"sinonimos":[],"tipo":"conceito","titulo":"O bootstrap semântico (registrar queimou a sonda)"}
\section{O bootstrap semântico (registrar queimou a sonda)}
O fechamento do arco do floxinato. Ao REGISTRAR o teste da forja, o termo falso ganhou uma preimagem — o próprio conceito do teste o contém — e deixou de ser falso: grau 0\ensuremath{\to}1. Um BOOTSTRAP SEMÂNTICO: 'você escreveu esse nome tantas vezes que agora ele existe.' É o OBSERVADOR virando parte do sistema observado (o instrumento entra no sistema, como na medição quântica, no ponto fixo de Gödel, no paradoxo de Russell). Não é bug: é o invariante EM AÇÃO — dar preimagem É gravar um cunho. A engenharia elegante não pôs um if-especial ('se nome==floxinato: ignora' — reprimir), mudou a INTERPRETAÇÃO (gerenciar). A sonda-mãe queimada é a prova viva da transição 0\ensuremath{\to}1; os testes seguintes usam forjas frescas. O nome percorreu: piada \ensuremath{\to} caso de teste \ensuremath{\to} teste oficial \ensuremath{\to} conceito do grafo \ensuremath{\to} meta-conhecimento \ensuremath{\to} sonda-mãe \ensuremath{\to} prova do invariante. Emerge só quando um sistema manipula o próprio conhecimento.
\prov{relações: parte\_de litografia; relaciona teste\_floxinato\_gamba; relaciona invariante\_topologico\_assinatura; relaciona teorema\_do\_ponto\_fixo; relaciona incompletude\_godel; relaciona godel\_incompletude; relaciona paradoxo\_de\_russell; relaciona colapso\_observador; relaciona identidade\_invariante\_w; relaciona observatorio\_nao\_fabrica; relaciona invariante\_indecidivel; relaciona icone\_indice\_simbolo; relaciona gerenciar\_nao\_reprimir}
\prov{proveniência: wiki \textperiodcentered{} cursor.txt (o fechamento) \textperiodcentered{} fato \textperiodcentered{} confiança 1.0 \textperiodcentered{} domínio litografia \textperiodcentered{} história 14 versões no jornal}

%CRISTAL {"arestas":[{"alvo":"mecanismo_higgs_ebh","confianca":1.0,"epistemico":"fato","origem":"wiki","rel":"confirma"},{"alvo":"cosmologia_quantica_hub","confianca":1.0,"epistemico":"fato","origem":"wiki","rel":"parte_de"}],"confianca":1.0,"contraexemplos":[],"descricao":"A excitação escalar massiva remanescente do campo de Higgs (spin 0, paridade +), único escalar fundamental do MP. Massa ≈125,38 GeV (~0,1%). Descoberto por ATLAS+CMS em 4 jul 2012 (canais H→γγ, H→ZZ*→4ℓ, H→WW*); Nobel 2013 (Englert, Higgs).","epistemico":"fato","exemplos":[],"id":"boson_higgs_125","memoria":"semantica","meta":{"autor":"broca-so","dominio":"fisica","fonte":"ATLAS arXiv:1207.7214 + CMS 1207.7235, 2012; Nobel 2013"},"origem":"wiki","palavras_chave":[],"sinonimos":[],"tipo":"conceito","titulo":"O bóson de Higgs (125 GeV)"}
\section{O bóson de Higgs (125 GeV)}
A excitação escalar massiva remanescente do campo de Higgs (spin 0, paridade +), único escalar fundamental do MP. Massa \ensuremath{\approx}125,38 GeV (\~{}0,1\%). Descoberto por ATLAS+CMS em 4 jul 2012 (canais H\ensuremath{\to}\ensuremath{\gamma}\ensuremath{\gamma}, H\ensuremath{\to}ZZ*\ensuremath{\to}4\ensuremath{\ell}, H\ensuremath{\to}WW*); Nobel 2013 (Englert, Higgs).
\prov{relações: confirma mecanismo\_higgs\_ebh; parte\_de cosmologia\_quantica\_hub}
\prov{proveniência: wiki \textperiodcentered{} ATLAS arXiv:1207.7214 + CMS 1207.7235, 2012; Nobel 2013 \textperiodcentered{} fato \textperiodcentered{} confiança 1.0 \textperiodcentered{} domínio fisica \textperiodcentered{} história 3 versões no jornal}

%CRISTAL {"arestas":[{"alvo":"unificacao_eletrofraca_gws","confianca":1.0,"epistemico":"fato","origem":"wiki","rel":"parte_de"},{"alvo":"angulo_weinberg","confianca":1.0,"epistemico":"fato","origem":"wiki","rel":"usa"}],"confianca":1.0,"contraexemplos":[],"descricao":"W± (corrente carregada) e Z⁰ (corrente neutra): M_W=80,369±0,013 GeV, M_Z=91,188 GeV, M_W=M_Z·cosθ_W. Descobertos no CERN (Spp̄S, UA1/UA2, 1983); Nobel 1984 (Rubbia, van der Meer). A medida CDF 2022 (M_W=80,4335) ficou em tensão com o MP, não confirmada — ponto aberto menor.","epistemico":"fato","exemplos":[],"id":"bosons_w_z","memoria":"semantica","meta":{"autor":"broca-so","dominio":"fisica","fonte":"UA1/UA2 1983; Nobel 1984; PDG 2024"},"origem":"wiki","palavras_chave":[],"sinonimos":[],"tipo":"conceito","titulo":"Bósons W e Z (mediadores massivos da força fraca)"}
\section{Bósons W e Z (mediadores massivos da força fraca)}
W$\pm$ (corrente carregada) e Z\ensuremath{^{0}} (corrente neutra): M\_W=80,369$\pm$0,013 GeV, M\_Z=91,188 GeV, M\_W=M\_Z\textperiodcentered cos\ensuremath{\theta}\_W. Descobertos no CERN (Spp\ensuremath{}S, UA1/UA2, 1983); Nobel 1984 (Rubbia, van der Meer). A medida CDF 2022 (M\_W=80,4335) ficou em tensão com o MP, não confirmada — ponto aberto menor.
\prov{relações: parte\_de unificacao\_eletrofraca\_gws; usa angulo\_weinberg}
\prov{proveniência: wiki \textperiodcentered{} UA1/UA2 1983; Nobel 1984; PDG 2024 \textperiodcentered{} fato \textperiodcentered{} confiança 1.0 \textperiodcentered{} domínio fisica \textperiodcentered{} história 3 versões no jornal}

%CRISTAL {"arestas":[{"alvo":"libbroca_so_camada","confianca":1.0,"epistemico":"fato","origem":"wiki","rel":"usa"},{"alvo":"cristal_motor_congelado","confianca":1.0,"epistemico":"fato","origem":"wiki","rel":"relaciona"}],"confianca":1.0,"contraexemplos":[],"descricao":"broca_borda.py liga por ctypes a libbroca_so.so (CPU: parábola/banda/POW) e libbroca_so_gpu.so (GPU: gato rede-2 + floco Φ). associativo_stream = só gato_stream gpu|cpu, explicitamente sem reabrir a parábola — o backend puro de recall associativo no metal, sob o cristal_motor_congelado.","epistemico":"fato","exemplos":[],"id":"broca_borda_so_ctypes","memoria":"semantica","meta":{"autor":"broca-so","dominio":"metal","fonte":"neuronio/broca_borda.py:gato_stream/associativo_stream"},"origem":"wiki","palavras_chave":[],"sinonimos":[],"tipo":"conceito","titulo":"Backend Associativo em dois .so via ctypes"}
\section{Backend Associativo em dois .so via ctypes}
broca\_borda.py liga por ctypes a libbroca\_so.so (CPU: parábola/banda/POW) e libbroca\_so\_gpu.so (GPU: gato rede-2 + floco \ensuremath{\Phi}). associativo\_stream = só gato\_stream gpu|cpu, explicitamente sem reabrir a parábola — o backend puro de recall associativo no metal, sob o cristal\_motor\_congelado.
\prov{relações: usa libbroca\_so\_camada; relaciona cristal\_motor\_congelado}
\prov{proveniência: wiki \textperiodcentered{} neuronio/broca\_borda.py:gato\_stream/associativo\_stream \textperiodcentered{} fato \textperiodcentered{} confiança 1.0 \textperiodcentered{} domínio metal \textperiodcentered{} história 3 versões no jornal}

%CRISTAL {"arestas":[{"alvo":"organismo_gex44","confianca":1.0,"epistemico":"fato","origem":"wiki","rel":"explica"},{"alvo":"antena_canal_ipc_bump","confianca":1.0,"epistemico":"fato","origem":"wiki","rel":"usa"},{"alvo":"pipeline_stratum","confianca":1.0,"epistemico":"fato","origem":"wiki","rel":"relaciona"}],"confianca":1.0,"contraexemplos":[],"descricao":"Daemon que faz poll (0.25s) de /home/claude/canal/broca/{in,out,seq} no GEX44: lê .pkt, desempacota, dedup por seq, responde 'olho:ping'→{ok} ou 'fluxo:stratum'→chunk do stratum_gex44.log + stats, e reempacota a resposta em out/{seq}.pkt. Deploy real via furo.sh (host gex44).","epistemico":"fato","exemplos":[],"id":"broca_olho_servico_gex44","memoria":"semantica","meta":{"autor":"broca-so","dominio":"metal","fonte":"scripts/gex44/broca_olho_servico.py"},"origem":"wiki","palavras_chave":[],"sinonimos":[],"tipo":"conceito","titulo":"O Serviço Olho no Servidor GEX44"}
\section{O Serviço Olho no Servidor GEX44}
Daemon que faz poll (0.25s) de /home/claude/canal/broca/\{in,out,seq\} no GEX44: lê .pkt, desempacota, dedup por seq, responde 'olho:ping'\ensuremath{\to}\{ok\} ou 'fluxo:stratum'\ensuremath{\to}chunk do stratum\_gex44.log + stats, e reempacota a resposta em out/\{seq\}.pkt. Deploy real via furo.sh (host gex44).
\prov{relações: explica organismo\_gex44; usa antena\_canal\_ipc\_bump; relaciona pipeline\_stratum}
\prov{proveniência: wiki \textperiodcentered{} scripts/gex44/broca\_olho\_servico.py \textperiodcentered{} fato \textperiodcentered{} confiança 1.0 \textperiodcentered{} domínio metal \textperiodcentered{} história 4 versões no jornal}

%CRISTAL {"arestas":[{"alvo":"materia_escura","confianca":1.0,"epistemico":"fato","origem":"wiki","rel":"confirma"},{"alvo":"mond_gravidade_modificada","confianca":1.0,"epistemico":"fato","origem":"wiki","rel":"contradiz"},{"alvo":"lente_gravitacional_1919","confianca":1.0,"epistemico":"fato","origem":"wiki","rel":"usa"}],"confianca":1.0,"contraexemplos":[],"descricao":"Na colisão de dois aglomerados (1E 0657-56, z=0,296), o gás quente (X, a maioria dos bárions) fica retido no centro por atrito, mas os picos de massa (por lente gravitacional) seguem as galáxias/matéria escura sem colidir — separação ESPACIAL direta (offset a 8σ) entre massa e bárions. Evidência empírica forte; dificílima para gravidade modificada pura (a massa não está onde estão os bárions).","epistemico":"fato","exemplos":[],"id":"bullet_cluster","memoria":"semantica","meta":{"autor":"broca-so","dominio":"cosmologia_dados","fonte":"Clowe et al. 2006 ApJL 648,L109 astro-ph/0608407"},"origem":"wiki","palavras_chave":[],"sinonimos":[],"tipo":"conceito","titulo":"Bullet Cluster — a separação massa–gás (prova direta)"}
\section{Bullet Cluster — a separação massa–gás (prova direta)}
Na colisão de dois aglomerados (1E 0657-56, z=0,296), o gás quente (X, a maioria dos bárions) fica retido no centro por atrito, mas os picos de massa (por lente gravitacional) seguem as galáxias/matéria escura sem colidir — separação ESPACIAL direta (offset a 8\ensuremath{\sigma}) entre massa e bárions. Evidência empírica forte; dificílima para gravidade modificada pura (a massa não está onde estão os bárions).
\prov{relações: confirma materia\_escura; contradiz mond\_gravidade\_modificada; usa lente\_gravitacional\_1919}
\prov{proveniência: wiki \textperiodcentered{} Clowe et al. 2006 ApJL 648,L109 astro-ph/0608407 \textperiodcentered{} fato \textperiodcentered{} confiança 1.0 \textperiodcentered{} domínio cosmologia\_dados \textperiodcentered{} história 4 versões no jornal}

%CRISTAL {"arestas":[{"alvo":"ckm_pmns_desalinhamento_s3","confianca":1.0,"epistemico":"fato","origem":"wiki","rel":"relaciona"},{"alvo":"problema_do_sabor","confianca":1.0,"epistemico":"fato","origem":"wiki","rel":"relaciona"},{"alvo":"cp_leptonica_leptogenese","confianca":1.0,"epistemico":"fato","origem":"wiki","rel":"relaciona"}],"confianca":1.0,"contraexemplos":[],"descricao":"Ressalva explícita dos papers: nenhum dos números de sabor — nem o ⅔ de Koide, nem a forma tribimaximal, nem δ_CP, nem o ângulo de Cabibbo — DERIVA da construção algébrica desta Parte. São o teste T3, EM ABERTO (DUNE, Hyper-K). O que é derivado é a MOLDURA (mistura=desalinhamento, Koide=espectro circulante); os valores são fronteira do programa, não resultado.","epistemico":"inferencia","exemplos":[],"id":"cabibbo_cp_nao_derivados","memoria":"semantica","meta":{"autor":"broca-so","dominio":"cosmologia","fonte":"paper_4_particulas.tex obs:mistura-S3 (T3)"},"origem":"wiki","palavras_chave":[],"sinonimos":[],"tipo":"conceito","titulo":"Lacuna: ângulo de Cabibbo e δ_CP não são derivados"}
\section{Lacuna: ângulo de Cabibbo e \ensuremath{\delta}\_CP não são derivados}
Ressalva explícita dos papers: nenhum dos números de sabor — nem o \ensuremath{\tfrac23} de Koide, nem a forma tribimaximal, nem \ensuremath{\delta}\_CP, nem o ângulo de Cabibbo — DERIVA da construção algébrica desta Parte. São o teste T3, EM ABERTO (DUNE, Hyper-K). O que é derivado é a MOLDURA (mistura=desalinhamento, Koide=espectro circulante); os valores são fronteira do programa, não resultado.
\prov{relações: relaciona ckm\_pmns\_desalinhamento\_s3; relaciona problema\_do\_sabor; relaciona cp\_leptonica\_leptogenese}
\prov{proveniência: wiki \textperiodcentered{} paper\_4\_particulas.tex obs:mistura-S3 (T3) \textperiodcentered{} inferencia \textperiodcentered{} confiança 1.0 \textperiodcentered{} domínio cosmologia \textperiodcentered{} história 4 versões no jornal}

%CRISTAL {"arestas":[{"alvo":"decaimento_radioativo","confianca":1.0,"epistemico":"fato","origem":"wiki","rel":"usa"},{"alvo":"cadeias_alem_dos_metais","confianca":1.0,"epistemico":"fato","origem":"wiki","rel":"relaciona"}],"confianca":1.0,"contraexemplos":[],"descricao":"Um núcleo decai numa sequência de nuclídeos até um estável (ex.: U-238 → … → Pb-206, 14 passos). É literalmente uma CADEIA/órbita de nuclídeos — a mesma estrutura das cadeias minerais.","epistemico":"inferencia","exemplos":[],"id":"cadeia_de_decaimento","memoria":"semantica","meta":{"dominio":"fisica","fonte":"energia nuclear"},"origem":"wiki","palavras_chave":[],"sinonimos":[],"tipo":"conceito","titulo":"Cadeia de Decaimento"}
\section{Cadeia de Decaimento}
Um núcleo decai numa sequência de nuclídeos até um estável (ex.: U-238 \ensuremath{\to} … \ensuremath{\to} Pb-206, 14 passos). É literalmente uma CADEIA/órbita de nuclídeos — a mesma estrutura das cadeias minerais.
\prov{relações: usa decaimento\_radioativo; relaciona cadeias\_alem\_dos\_metais}
\prov{proveniência: wiki \textperiodcentered{} energia nuclear \textperiodcentered{} inferencia \textperiodcentered{} confiança 1.0 \textperiodcentered{} domínio fisica \textperiodcentered{} história 3 versões no jornal}

%CRISTAL {"arestas":[{"alvo":"materia_modelo_padrao","confianca":1.0,"epistemico":"fato","origem":"wiki","rel":"explica"},{"alvo":"alpha_tres_reguas","confianca":1.0,"epistemico":"fato","origem":"wiki","rel":"explica"},{"alvo":"leitura_de_part_culas_o_setor_de_materia_pela_torre_de_divisao","confianca":1.0,"epistemico":"fato","origem":"wiki","rel":"parte_de"},{"alvo":"unidades_algebricas_pu","confianca":1.0,"epistemico":"fato","origem":"wiki","rel":"relaciona"}],"confianca":1.0,"contraexemplos":[],"descricao":"ℂ⊗𝕆=Cℓ(6)≅ℂ(8): seis dos sete Lᵢ montam 3 escadas αa; o número N=Σαa†αa tem espectro 0,1,1,1,2,2,2,3, e Q=N/3 sai QUANTIZADA porque N é inteiro. O ideal mínimo (vácuo=ν) tem 8 estados (ν, tripleto d̄, anti-tripleto u, e⁺); +conjugado+antipartículas = 16 de Weyl de uma geração, Q∈0,±⅓,±⅔,±1. Derivada, não posta.","epistemico":"inferencia","exemplos":[],"id":"carga_q_n_sobre_3","memoria":"semantica","meta":{"autor":"broca-so","dominio":"cosmologia","fonte":"paper_4_particulas.tex §cxo-cargas (cxo_48_cargas.py)"},"origem":"wiki","palavras_chave":[],"sinonimos":[],"tipo":"conceito","titulo":"A carga Q=N/3 (→ e/3), derivada do número de ocupação"}
\section{A carga Q=N/3 (\ensuremath{\to} e/3), derivada do número de ocupação}
\ensuremath{\mathbb{C}}\ensuremath{\otimes}\ensuremath{\mathbb{O}}=C\ensuremath{\ell}(6)\ensuremath{\cong}\ensuremath{\mathbb{C}}(8): seis dos sete L\ensuremath{_{i}} montam 3 escadas \ensuremath{\alpha}a; o número N=\ensuremath{\Sigma}\ensuremath{\alpha}a\ensuremath{\dagger}\ensuremath{\alpha}a tem espectro 0,1,1,1,2,2,2,3, e Q=N/3 sai QUANTIZADA porque N é inteiro. O ideal mínimo (vácuo=\ensuremath{\nu}) tem 8 estados (\ensuremath{\nu}, tripleto d\ensuremath{}, anti-tripleto u, e\ensuremath{^{+}}); +conjugado+antipartículas = 16 de Weyl de uma geração, Q\ensuremath{\in}0,$\pm$\ensuremath{\tfrac13},$\pm$\ensuremath{\tfrac23},$\pm$1. Derivada, não posta.
\prov{relações: explica materia\_modelo\_padrao; explica alpha\_tres\_reguas; parte\_de leitura\_de\_part\_culas\_o\_setor\_de\_materia\_pela\_torre\_de\_divisao; relaciona unidades\_algebricas\_pu}
\prov{proveniência: wiki \textperiodcentered{} paper\_4\_particulas.tex §cxo-cargas (cxo\_48\_cargas.py) \textperiodcentered{} inferencia \textperiodcentered{} confiança 1.0 \textperiodcentered{} domínio cosmologia \textperiodcentered{} história 6 versões no jornal}

%CRISTAL {"arestas":[{"alvo":"wimps_candidato","confianca":1.0,"epistemico":"fato","origem":"wiki","rel":"contradiz"}],"confianca":1.0,"contraexemplos":[],"descricao":"Buscam recuo nuclear de WIMPs em xenônio líquido subterrâneo. Resultados NULOS, limites líderes (LZ 417 dias, seção spin-independent ~2×10⁻⁴⁸ cm² a ~40 GeV). Nenhuma detecção não-gravitacional confirmada por nenhuma via (direta, indireta/Fermi-IceCube, colisor/LHC). O NEUTRINO FOG (neutrinos solares dão recuos idênticos, fundo irredutível) já apareceu: LZ viu ⁸B via CEvNS a 4,5σ.","epistemico":"fato","exemplos":[],"id":"cdm_deteccao_nula","memoria":"semantica","meta":{"autor":"broca-so","dominio":"cosmologia_dados","fonte":"LZ 2025 (arXiv:2512.23039)"},"origem":"wiki","palavras_chave":[],"sinonimos":[],"tipo":"conceito","titulo":"Detecção direta (LZ/XENONnT): nula, e o neutrino fog"}
\section{Detecção direta (LZ/XENONnT): nula, e o neutrino fog}
Buscam recuo nuclear de WIMPs em xenônio líquido subterrâneo. Resultados NULOS, limites líderes (LZ 417 dias, seção spin-independent \~{}2$\times$10\ensuremath{^{-48}} cm\ensuremath{^{2}} a \~{}40 GeV). Nenhuma detecção não-gravitacional confirmada por nenhuma via (direta, indireta/Fermi-IceCube, colisor/LHC). O NEUTRINO FOG (neutrinos solares dão recuos idênticos, fundo irredutível) já apareceu: LZ viu \ensuremath{^{8}}B via CEvNS a 4,5\ensuremath{\sigma}.
\prov{relações: contradiz wimps\_candidato}
\prov{proveniência: wiki \textperiodcentered{} LZ 2025 (arXiv:2512.23039) \textperiodcentered{} fato \textperiodcentered{} confiança 1.0 \textperiodcentered{} domínio cosmologia\_dados \textperiodcentered{} história 2 versões no jornal}

%CRISTAL {"arestas":[{"alvo":"materia_escura","confianca":1.0,"epistemico":"fato","origem":"wiki","rel":"relaciona"}],"confianca":1.0,"contraexemplos":[],"descricao":"Problemas reais mas provavelmente não fatais: core-cusp (CDM prevê perfis pontiagudos NFW, observam-se núcleos planos em anãs); missing satellites (CDM prevê centenas de sub-halos, observam-se dezenas); too-big-to-fail. Largamente atenuados por física bariônica (feedback de supernovas, reionização) e satélites ultra-fracos. WDM e SIDM em debate.","epistemico":"inferencia","exemplos":[],"id":"cdm_tensoes_pequena_escala","memoria":"semantica","meta":{"autor":"broca-so","dominio":"cosmologia_dados","fonte":"Bullock & Boylan-Kolchin 2017 ARA&A 55,343 arXiv:1707.04256"},"origem":"wiki","palavras_chave":[],"sinonimos":[],"tipo":"conceito","titulo":"Tensões de pequena escala do CDM"}
\section{Tensões de pequena escala do CDM}
Problemas reais mas provavelmente não fatais: core-cusp (CDM prevê perfis pontiagudos NFW, observam-se núcleos planos em anãs); missing satellites (CDM prevê centenas de sub-halos, observam-se dezenas); too-big-to-fail. Largamente atenuados por física bariônica (feedback de supernovas, reionização) e satélites ultra-fracos. WDM e SIDM em debate.
\prov{relações: relaciona materia\_escura}
\prov{proveniência: wiki \textperiodcentered{} Bullock \& Boylan-Kolchin 2017 ARA\&A 55,343 arXiv:1707.04256 \textperiodcentered{} inferencia \textperiodcentered{} confiança 1.0 \textperiodcentered{} domínio cosmologia\_dados \textperiodcentered{} história 2 versões no jornal}

%CRISTAL {"arestas":[{"alvo":"numero_chern_tknn","confianca":1.0,"epistemico":"fato","origem":"wiki","rel":"explica"},{"alvo":"a_base_topologica","confianca":1.0,"epistemico":"fato","origem":"wiki","rel":"usa"},{"alvo":"esfera_bloch_base_hopf","confianca":1.0,"epistemico":"fato","origem":"wiki","rel":"usa"},{"alvo":"cosmologia_quantica_hub","confianca":1.0,"epistemico":"fato","origem":"wiki","rel":"parte_de"}],"confianca":1.0,"contraexemplos":[],"descricao":"O n do Hall é O MESMO Chern do fibrado tautológico 𝒪(−1) sobre a esfera de Bloch S²=ℂP¹: os autoespaços +i da estrutura complexa R_â(q)=q·â formam 𝒪(−1) de Chern −1, cujo espaço total é a fibração de Hopf S³→S². A função de onda do qubit É uma SEÇÃO desse fibrado; a curvatura é a forma de área de S²; a holonomia é o campo de um monopolo de Dirac de carga ½. C=−1≠0 diz que não há um 'i' único sobre toda a Bloch — o Hall quântico mede o Chern da fibra de Lopes. [D] identidade matemática (não analogia).","epistemico":"inferencia","exemplos":[],"id":"chern_fibrado_hopf","memoria":"semantica","meta":{"autor":"broca-so","dominio":"cosmologia","fonte":"paper_3_quantica.tex sec:onda (lema:detLz)"},"origem":"wiki","palavras_chave":[],"sinonimos":[],"tipo":"conceito","titulo":"O número de Chern É o Chern do fibrado de Hopf 𝒪(−1)"}
\section{O número de Chern É o Chern do fibrado de Hopf \ensuremath{\mathcal{O}}(\ensuremath{-}1)}
O n do Hall é O MESMO Chern do fibrado tautológico \ensuremath{\mathcal{O}}(\ensuremath{-}1) sobre a esfera de Bloch S\ensuremath{^{2}}=\ensuremath{\mathbb{C}}P\ensuremath{^{1}}: os autoespaços +i da estrutura complexa R\_â(q)=q\textperiodcentered â formam \ensuremath{\mathcal{O}}(\ensuremath{-}1) de Chern \ensuremath{-}1, cujo espaço total é a fibração de Hopf S\ensuremath{^{3}}\ensuremath{\to}S\ensuremath{^{2}}. A função de onda do qubit É uma SEÇÃO desse fibrado; a curvatura é a forma de área de S\ensuremath{^{2}}; a holonomia é o campo de um monopolo de Dirac de carga \ensuremath{\tfrac12}. C=\ensuremath{-}1\ensuremath{\ne}0 diz que não há um 'i' único sobre toda a Bloch — o Hall quântico mede o Chern da fibra de Lopes. [D] identidade matemática (não analogia).
\prov{relações: explica numero\_chern\_tknn; usa a\_base\_topologica; usa esfera\_bloch\_base\_hopf; parte\_de cosmologia\_quantica\_hub}
\prov{proveniência: wiki \textperiodcentered{} paper\_3\_quantica.tex sec:onda (lema:detLz) \textperiodcentered{} inferencia \textperiodcentered{} confiança 1.0 \textperiodcentered{} domínio cosmologia \textperiodcentered{} história 5 versões no jornal}

%CRISTAL {"arestas":[{"alvo":"contracao_einstein_k_tensorial","confianca":1.0,"epistemico":"fato","origem":"wiki","rel":"eh_um"},{"alvo":"design_semantica_tensorial","confianca":1.0,"epistemico":"fato","origem":"wiki","rel":"usa"},{"alvo":"curvatura_n_menos_1","confianca":1.0,"epistemico":"fato","origem":"wiki","rel":"usa"},{"alvo":"agulha_curvatura_n1","confianca":1.0,"epistemico":"fato","origem":"wiki","rel":"relaciona"},{"alvo":"o_chicote_a_limpeza_nao_um_operador","confianca":1.0,"epistemico":"fato","origem":"wiki","rel":"relaciona"},{"alvo":"beta_numeracao_metalica","confianca":1.0,"epistemico":"fato","origem":"wiki","rel":"usa"},{"alvo":"efe_einstein","confianca":1.0,"epistemico":"fato","origem":"wiki","rel":"relaciona"},{"alvo":"gato","confianca":1.0,"epistemico":"fato","origem":"wiki","rel":"usa"},{"alvo":"flip_e_rotacao_de_wick","confianca":1.0,"epistemico":"fato","origem":"wiki","rel":"relaciona"},{"alvo":"linguagem_fisica_unificada","confianca":1.0,"epistemico":"fato","origem":"wiki","rel":"parte_de"}],"confianca":1.0,"contraexemplos":[],"descricao":"O CHICOTE implementa a dinâmica posicional da régua mineral do BAI: mover-se ao longo da régua (a β-numeração na base σ_n) é aplicar o gato A_n=[[n,1],[1,0]] — o passo/carry — e o gato é um TENSOR de rank 2. Compor gatos é CONTRAIR índices pela soma de Einstein, (A·B)^i_k=Σ_j A^i_j B^j_k. A métrica é o CONE de luz mineral: a forma M (assinatura Lorentziana (1,1)) cujas direções nulas são os autovetores (σ_n, −1/σ_n), com A^T M A=−M (o flip ε=−1, pois det(A_n)=−1, e porque A_n é simétrica) — o achado substantivo. A curvatura conta N−1 pontos críticos FINITOS do dupolinômio de grau N (=deg F′; a ramificação total R-H é 2N−2). Verif. 6 fatos + 1 analogia (a EFE) em tensor_mineral.py. É o lado GEOMETRIA da máquina.","epistemico":"fato","exemplos":[],"id":"chicote_e_o_tensor","memoria":"semantica","meta":{"dominio":"fisica","fonte":"linguagem física unificada"},"origem":"wiki","palavras_chave":[],"sinonimos":[],"tipo":"conceito","titulo":"O Chicote é o Tensor (a dinâmica posicional da régua)"}
\section{O Chicote é o Tensor (a dinâmica posicional da régua)}
O CHICOTE implementa a dinâmica posicional da régua mineral do BAI: mover-se ao longo da régua (a \ensuremath{\beta}-numeração na base \ensuremath{\sigma}\_n) é aplicar o gato A\_n=[[n,1],[1,0]] — o passo/carry — e o gato é um TENSOR de rank 2. Compor gatos é CONTRAIR índices pela soma de Einstein, (A\textperiodcentered B)\^{}i\_k=\ensuremath{\Sigma}\_j A\^{}i\_j B\^{}j\_k. A métrica é o CONE de luz mineral: a forma M (assinatura Lorentziana (1,1)) cujas direções nulas são os autovetores (\ensuremath{\sigma}\_n, \ensuremath{-}1/\ensuremath{\sigma}\_n), com A\^{}T M A=\ensuremath{-}M (o flip \ensuremath{\varepsilon}=\ensuremath{-}1, pois det(A\_n)=\ensuremath{-}1, e porque A\_n é simétrica) — o achado substantivo. A curvatura conta N\ensuremath{-}1 pontos críticos FINITOS do dupolinômio de grau N (=deg F\ensuremath{'}; a ramificação total R-H é 2N\ensuremath{-}2). Verif. 6 fatos + 1 analogia (a EFE) em tensor\_mineral.py. É o lado GEOMETRIA da máquina.
\prov{relações: eh\_um contracao\_einstein\_k\_tensorial; usa design\_semantica\_tensorial; usa curvatura\_n\_menos\_1; relaciona agulha\_curvatura\_n1; relaciona o\_chicote\_a\_limpeza\_nao\_um\_operador; usa beta\_numeracao\_metalica; relaciona efe\_einstein; usa gato; relaciona flip\_e\_rotacao\_de\_wick; parte\_de linguagem\_fisica\_unificada}
\prov{proveniência: wiki \textperiodcentered{} linguagem física unificada \textperiodcentered{} fato \textperiodcentered{} confiança 1.0 \textperiodcentered{} domínio fisica \textperiodcentered{} história 22 versões no jornal}

%CRISTAL {"arestas":[{"alvo":"ckm_pmns_desalinhamento_s3","confianca":1.0,"epistemico":"fato","origem":"wiki","rel":"usa"},{"alvo":"wolfenstein_hierarquia_ckm","confianca":1.0,"epistemico":"fato","origem":"wiki","rel":"relaciona"},{"alvo":"neutrino_anomalo_dois_testes","confianca":1.0,"epistemico":"fato","origem":"wiki","rel":"relaciona"}],"confianca":1.0,"contraexemplos":[],"descricao":"Lida por essa régua: a CKM é PEQUENA — ‖|V_CKM|−𝟙‖≈0,33 (up e down quebram S₃ quase alinhados, quase a mesma base de Fourier). A PMNS é GRANDE — ≈1,36, quatro vezes maior (ℓ e ν desalinhados), e não aleatória: perto da tribimaximal (‖PMNS−TBM‖≈0,24; θ12→35,3°, θ23→45°, θ13=8,6° a quebra residual). O neutrino é o setor anômalo. [N] os valores; o contraste pequeno/grande é o teste consistente com o teorema, não derivado dele.","epistemico":"inferencia","exemplos":[],"id":"ckm_pequena_pmns_grande","memoria":"semantica","meta":{"autor":"broca-so","dominio":"cosmologia","fonte":"paper_4_particulas.tex obs:mistura-S3"},"origem":"wiki","palavras_chave":[],"sinonimos":[],"tipo":"conceito","titulo":"CKM pequena (up/down alinhados) vs PMNS grande (~4×)"}
\section{CKM pequena (up/down alinhados) vs PMNS grande (\~{}4$\times$)}
Lida por essa régua: a CKM é PEQUENA — \ensuremath{\|}|V\_CKM|\ensuremath{-}\ensuremath{\mathbf{1}}\ensuremath{\|}\ensuremath{\approx}0,33 (up e down quebram S\ensuremath{_{3}} quase alinhados, quase a mesma base de Fourier). A PMNS é GRANDE — \ensuremath{\approx}1,36, quatro vezes maior (\ensuremath{\ell} e \ensuremath{\nu} desalinhados), e não aleatória: perto da tribimaximal (\ensuremath{\|}PMNS\ensuremath{-}TBM\ensuremath{\|}\ensuremath{\approx}0,24; \ensuremath{\theta}12\ensuremath{\to}35,3$^{\circ}$, \ensuremath{\theta}23\ensuremath{\to}45$^{\circ}$, \ensuremath{\theta}13=8,6$^{\circ}$ a quebra residual). O neutrino é o setor anômalo. [N] os valores; o contraste pequeno/grande é o teste consistente com o teorema, não derivado dele.
\prov{relações: usa ckm\_pmns\_desalinhamento\_s3; relaciona wolfenstein\_hierarquia\_ckm; relaciona neutrino\_anomalo\_dois\_testes}
\prov{proveniência: wiki \textperiodcentered{} paper\_4\_particulas.tex obs:mistura-S3 \textperiodcentered{} inferencia \textperiodcentered{} confiança 1.0 \textperiodcentered{} domínio cosmologia \textperiodcentered{} história 4 versões no jornal}

%CRISTAL {"arestas":[{"alvo":"matrizes_massa_circulantes","confianca":1.0,"epistemico":"fato","origem":"wiki","rel":"usa"},{"alvo":"pmns_desalinhamento_s3","confianca":1.0,"epistemico":"fato","origem":"wiki","rel":"generaliza"},{"alvo":"matriz_ckm","confianca":1.0,"epistemico":"fato","origem":"wiki","rel":"explica"},{"alvo":"matriz_pmns","confianca":1.0,"epistemico":"fato","origem":"wiki","rel":"explica"},{"alvo":"cosmologia_quantica_hub","confianca":1.0,"epistemico":"fato","origem":"wiki","rel":"parte_de"}],"confianca":1.0,"contraexemplos":[],"descricao":"Se dois setores (up/down, ou ℓ/ν) são ambos S₃-circulantes, partilham a mesma F, e a mistura é V=F†F=𝟙 — NULA. Logo a mistura observada não é parâmetro livre: mede o DESALINHAMENTO entre as quebras de S₃ dos dois setores. A CKM é o desalinhamento up/down; a PMNS o desalinhamento ℓ/ν. [D] o mecanismo (mistura=desalinhamento é teorema); [J] identificar essa S₃ com a S₃ física de sabor.","epistemico":"inferencia","exemplos":[],"id":"ckm_pmns_desalinhamento_s3","memoria":"semantica","meta":{"autor":"broca-so","dominio":"cosmologia","fonte":"paper_4_particulas.tex obs:mistura-S3"},"origem":"wiki","palavras_chave":[],"sinonimos":[],"tipo":"conceito","titulo":"CKM e PMNS são o mesmo mecanismo: desalinhamento de S₃"}
\section{CKM e PMNS são o mesmo mecanismo: desalinhamento de S\ensuremath{_{3}}}
Se dois setores (up/down, ou \ensuremath{\ell}/\ensuremath{\nu}) são ambos S\ensuremath{_{3}}-circulantes, partilham a mesma F, e a mistura é V=F\ensuremath{\dagger}F=\ensuremath{\mathbf{1}} — NULA. Logo a mistura observada não é parâmetro livre: mede o DESALINHAMENTO entre as quebras de S\ensuremath{_{3}} dos dois setores. A CKM é o desalinhamento up/down; a PMNS o desalinhamento \ensuremath{\ell}/\ensuremath{\nu}. [D] o mecanismo (mistura=desalinhamento é teorema); [J] identificar essa S\ensuremath{_{3}} com a S\ensuremath{_{3}} física de sabor.
\prov{relações: usa matrizes\_massa\_circulantes; generaliza pmns\_desalinhamento\_s3; explica matriz\_ckm; explica matriz\_pmns; parte\_de cosmologia\_quantica\_hub}
\prov{proveniência: wiki \textperiodcentered{} paper\_4\_particulas.tex obs:mistura-S3 \textperiodcentered{} inferencia \textperiodcentered{} confiança 1.0 \textperiodcentered{} domínio cosmologia \textperiodcentered{} história 6 versões no jornal}

%CRISTAL {"arestas":[{"alvo":"no_cloning_invariancia_metrica","confianca":1.0,"epistemico":"fato","origem":"wiki","rel":"usa"},{"alvo":"seguranca_algebrica_estrutura_nao_hardness","confianca":1.0,"epistemico":"fato","origem":"wiki","rel":"relaciona"}],"confianca":1.0,"contraexemplos":[],"descricao":"Cloning exato é proibido, mas o aproximado via a própria multiplicação (cópia=z⋆a/‖z⋆a‖) tem fidelidade fechada F=b₀²/[1−(1−s²)‖a_z×a_a‖²], MONOTÔNICA DECRESCENTE em |s|. Contraintuitivo: a álgebra comutativa (s=0, F≈0,37) é a MAIS clonável; o quaternion (s=±1, F≈0,25) é o meio; fora da ilha F→0. Base de cripto anti-falsificação nativa. [D] fórmula provada + [N] verif. (inclui registro honesto de previsão errada do autor: 'quaternion ótimo' refutada pela conta).","epistemico":"inferencia","exemplos":[],"id":"cloning_aproximado_fidelidade_s","memoria":"semantica","meta":{"autor":"broca-so","dominio":"cosmologia","fonte":"paper_G_no_cloning.tex §6 (Prop. fidelidade-s)"},"origem":"wiki","palavras_chave":[],"sinonimos":[],"tipo":"conceito","titulo":"Cloning aproximado: fidelidade decresce em |s| (cripto)"}
\section{Cloning aproximado: fidelidade decresce em |s| (cripto)}
Cloning exato é proibido, mas o aproximado via a própria multiplicação (cópia=z\ensuremath{\star}a/\ensuremath{\|}z\ensuremath{\star}a\ensuremath{\|}) tem fidelidade fechada F=b\ensuremath{_{0}}\ensuremath{^{2}}/[1\ensuremath{-}(1\ensuremath{-}s\ensuremath{^{2}})\ensuremath{\|}a\_z$\times$a\_a\ensuremath{\|}\ensuremath{^{2}}], MONOTÔNICA DECRESCENTE em |s|. Contraintuitivo: a álgebra comutativa (s=0, F\ensuremath{\approx}0,37) é a MAIS clonável; o quaternion (s=$\pm$1, F\ensuremath{\approx}0,25) é o meio; fora da ilha F\ensuremath{\to}0. Base de cripto anti-falsificação nativa. [D] fórmula provada + [N] verif. (inclui registro honesto de previsão errada do autor: 'quaternion ótimo' refutada pela conta).
\prov{relações: usa no\_cloning\_invariancia\_metrica; relaciona seguranca\_algebrica\_estrutura\_nao\_hardness}
\prov{proveniência: wiki \textperiodcentered{} paper\_G\_no\_cloning.tex §6 (Prop. fidelidade-s) \textperiodcentered{} inferencia \textperiodcentered{} confiança 1.0 \textperiodcentered{} domínio cosmologia \textperiodcentered{} história 3 versões no jornal}

%CRISTAL {"arestas":[{"alvo":"cosmologia_quantica_hub","confianca":1.0,"epistemico":"fato","origem":"wiki","rel":"parte_de"}],"confianca":1.0,"contraexemplos":[],"descricao":"Radiação relíquia da recombinação (~380 mil anos); as anisotropias codificam os parâmetros cosmológicos. COBE (1992): 1ª detecção de anisotropias, rms≈30µK (Nobel 2006). WMAP: cosmologia de precisão. Planck (2018): a medida mais precisa — base dos parâmetros do ΛCDM. A espinha dorsal empírica.","epistemico":"fato","exemplos":[],"id":"cmb_cobe_wmap_planck","memoria":"semantica","meta":{"autor":"broca-so","dominio":"cosmologia_dados","fonte":"Smoot+1992; Planck 2018"},"origem":"wiki","palavras_chave":[],"sinonimos":[],"tipo":"conceito","titulo":"CMB — COBE → WMAP → Planck"}
\section{CMB — COBE \ensuremath{\to} WMAP \ensuremath{\to} Planck}
Radiação relíquia da recombinação (\~{}380 mil anos); as anisotropias codificam os parâmetros cosmológicos. COBE (1992): 1ª detecção de anisotropias, rms\ensuremath{\approx}30\ensuremath{\mu}K (Nobel 2006). WMAP: cosmologia de precisão. Planck (2018): a medida mais precisa — base dos parâmetros do \ensuremath{\Lambda}CDM. A espinha dorsal empírica.
\prov{relações: parte\_de cosmologia\_quantica\_hub}
\prov{proveniência: wiki \textperiodcentered{} Smoot+1992; Planck 2018 \textperiodcentered{} fato \textperiodcentered{} confiança 1.0 \textperiodcentered{} domínio cosmologia\_dados \textperiodcentered{} história 2 versões no jornal}

%CRISTAL {"arestas":[{"alvo":"materia_escura","confianca":1.0,"epistemico":"fato","origem":"wiki","rel":"confirma"},{"alvo":"mond_gravidade_modificada","confianca":1.0,"epistemico":"fato","origem":"wiki","rel":"contradiz"},{"alvo":"cmb_cobe_wmap_planck","confianca":1.0,"epistemico":"fato","origem":"wiki","rel":"parte_de"}],"confianca":1.0,"contraexemplos":[],"descricao":"As alturas relativas dos picos par/ímpar do espectro codificam as densidades de bárions e de matéria escura; um 3º pico realçado indica que matéria escura NÃO-bariônica dominava a densidade de matéria antes da recombinação. Exige matéria escura FRIA e não-bariônica — vínculo que MOND/TeVeS não reproduzem naturalmente.","epistemico":"fato","exemplos":[],"id":"cmb_terceiro_pico","memoria":"semantica","meta":{"autor":"broca-so","dominio":"cosmologia_dados","fonte":"Planck 2018; Hu & Dodelson 2002"},"origem":"wiki","palavras_chave":[],"sinonimos":[],"tipo":"conceito","titulo":"Os picos acústicos do CMB (o 3º pico)"}
\section{Os picos acústicos do CMB (o 3º pico)}
As alturas relativas dos picos par/ímpar do espectro codificam as densidades de bárions e de matéria escura; um 3º pico realçado indica que matéria escura NÃO-bariônica dominava a densidade de matéria antes da recombinação. Exige matéria escura FRIA e não-bariônica — vínculo que MOND/TeVeS não reproduzem naturalmente.
\prov{relações: confirma materia\_escura; contradiz mond\_gravidade\_modificada; parte\_de cmb\_cobe\_wmap\_planck}
\prov{proveniência: wiki \textperiodcentered{} Planck 2018; Hu \& Dodelson 2002 \textperiodcentered{} fato \textperiodcentered{} confiança 1.0 \textperiodcentered{} domínio cosmologia\_dados \textperiodcentered{} história 4 versões no jornal}

%CRISTAL {"arestas":[{"alvo":"codigo_genetico","confianca":1.0,"epistemico":"fato","origem":"wiki","rel":"explica"},{"alvo":"limite_hurwitz_s1","confianca":1.0,"epistemico":"fato","origem":"wiki","rel":"usa"},{"alvo":"dois_qubits_octonios","confianca":1.0,"epistemico":"fato","origem":"wiki","rel":"usa"},{"alvo":"codigo_robustez_erro","confianca":1.0,"epistemico":"fato","origem":"wiki","rel":"explica"}],"confianca":1.0,"contraexemplos":[],"descricao":"1 base = ℂ/ℍ; par = 𝕆 (ainda na ilha); o CÓDON de 3 bases = 3 qubits = ℝ¹⁶ = sedênios 𝕊 = FORA da ilha de Hurwitz (divisores de zero). '3 é onde a álgebra de divisão acaba — onde a química vira informação.' Testado (codon_tripartite.py): o códon é robusto tipo-W, não frágil tipo-GHZ — o WOBBLE confirmado empiricamente (a 3ª base é redundante, mexer nela quase não troca o aminoácido). A assinatura octônica capta a estrutura posicional (wobble difusivo, α≈1) em toda a vida testada. [D] a barreira de Hurwitz + [N] o wobble medido.","epistemico":"inferencia","exemplos":[],"id":"codon_fora_da_ilha_hurwitz","memoria":"semantica","meta":{"autor":"broca-so","dominio":"cosmologia","fonte":"hiper/circular/codon_tripartite.py, assinatura_dna_octonio.py"},"origem":"wiki","palavras_chave":[],"sinonimos":[],"tipo":"conceito","titulo":"O códon de 3 bases sai da ilha de Hurwitz (onde a química vira informação)"}
\section{O códon de 3 bases sai da ilha de Hurwitz (onde a química vira informação)}
1 base = \ensuremath{\mathbb{C}}/\ensuremath{\mathbb{H}}; par = \ensuremath{\mathbb{O}} (ainda na ilha); o CÓDON de 3 bases = 3 qubits = \ensuremath{\mathbb{R}}\ensuremath{^{16}} = sedênios \ensuremath{\mathbb{S}} = FORA da ilha de Hurwitz (divisores de zero). '3 é onde a álgebra de divisão acaba — onde a química vira informação.' Testado (codon\_tripartite.py): o códon é robusto tipo-W, não frágil tipo-GHZ — o WOBBLE confirmado empiricamente (a 3ª base é redundante, mexer nela quase não troca o aminoácido). A assinatura octônica capta a estrutura posicional (wobble difusivo, \ensuremath{\alpha}\ensuremath{\approx}1) em toda a vida testada. [D] a barreira de Hurwitz + [N] o wobble medido.
\prov{relações: explica codigo\_genetico; usa limite\_hurwitz\_s1; usa dois\_qubits\_octonios; explica codigo\_robustez\_erro}
\prov{proveniência: wiki \textperiodcentered{} hiper/circular/codon\_tripartite.py, assinatura\_dna\_octonio.py \textperiodcentered{} inferencia \textperiodcentered{} confiança 1.0 \textperiodcentered{} domínio cosmologia \textperiodcentered{} história 5 versões no jornal}

%CRISTAL {"arestas":[{"alvo":"primitivas_analiticas","confianca":1.0,"epistemico":"fato","origem":"wiki","rel":"usa"},{"alvo":"translucido_nijhoff","confianca":1.0,"epistemico":"fato","origem":"wiki","rel":"relaciona"},{"alvo":"joalheria_hub","confianca":1.0,"epistemico":"fato","origem":"wiki","rel":"parte_de"}],"confianca":1.0,"contraexemplos":[],"descricao":"GLSL (a linguagem de shading da GPU, Khronos/NVIDIA) reproduzida e COMPILADA no metal do broca-so. Duas peças: (1) a LIB glsl.py — os built-ins da linguagem vetorizados (dot/cross/normalize/reflect/refract/mix/clamp/step/smoothstep/…), os float por-raio saem (…,1) para vec×float fazer broadcast como no GLSL; (2) o COMPILADOR glsl_compilador.py — transpila o código GLSL de verdade (os intersectores do Quílez) para rodar VETORIZADO. O desafio: GLSL é SIMT (por-pixel, escalar, com divergência) e o broca-so é SIMD (todos os raios). O compilador faz a PREDICAÇÃO: if/else/return viram MÁSCARA de lanes + where; o return fixa nas lanes ativas e desativa; o inout é escrito só nas ativas. Roda o iSphere e o iEllipsoid REAIS do Quílez, sem alterar a lógica, casando as primitivas nativas (erro 0). Dicionário glsl_render↔dsl_render (GLSL↔o metal: refract=_refrata, getSkyColor=_sky_nijhoff) e glsl_render↔bai_render (GLSL↔a máquina: refract=flip, bounce=órbita, interseção=colapso). O Teorema Geral: um operador de render, três línguas.","epistemico":"fato","exemplos":[],"id":"compilador_glsl_broca","memoria":"semantica","meta":{"dominio":"joalheria","fonte":"maquina de corte"},"origem":"wiki","palavras_chave":[],"sinonimos":[],"tipo":"conceito","titulo":"O compilador GLSL no broca-so: o código do Quílez rodando no metal (SIMT→SIMD)"}
\section{O compilador GLSL no broca-so: o código do Quílez rodando no metal (SIMT\ensuremath{\to}SIMD)}
GLSL (a linguagem de shading da GPU, Khronos/NVIDIA) reproduzida e COMPILADA no metal do broca-so. Duas peças: (1) a LIB glsl.py — os built-ins da linguagem vetorizados (dot/cross/normalize/reflect/refract/mix/clamp/step/smoothstep/…), os float por-raio saem (…,1) para vec$\times$float fazer broadcast como no GLSL; (2) o COMPILADOR glsl\_compilador.py — transpila o código GLSL de verdade (os intersectores do Quílez) para rodar VETORIZADO. O desafio: GLSL é SIMT (por-pixel, escalar, com divergência) e o broca-so é SIMD (todos os raios). O compilador faz a PREDICAÇÃO: if/else/return viram MÁSCARA de lanes + where; o return fixa nas lanes ativas e desativa; o inout é escrito só nas ativas. Roda o iSphere e o iEllipsoid REAIS do Quílez, sem alterar a lógica, casando as primitivas nativas (erro 0). Dicionário glsl\_render\ensuremath{\leftrightarrow}dsl\_render (GLSL\ensuremath{\leftrightarrow}o metal: refract=\_refrata, getSkyColor=\_sky\_nijhoff) e glsl\_render\ensuremath{\leftrightarrow}bai\_render (GLSL\ensuremath{\leftrightarrow}a máquina: refract=flip, bounce=órbita, interseção=colapso). O Teorema Geral: um operador de render, três línguas.
\prov{relações: usa primitivas\_analiticas; relaciona translucido\_nijhoff; parte\_de joalheria\_hub}
\prov{proveniência: wiki \textperiodcentered{} maquina de corte \textperiodcentered{} fato \textperiodcentered{} confiança 1.0 \textperiodcentered{} domínio joalheria \textperiodcentered{} história 80 versões no jornal}

%CRISTAL {"arestas":[{"alvo":"yang_mills_mass_gap_milenio","confianca":1.0,"epistemico":"fato","origem":"wiki","rel":"relaciona"},{"alvo":"qcd_cromodinamica","confianca":1.0,"epistemico":"fato","origem":"wiki","rel":"parte_de"}],"confianca":1.0,"contraexemplos":[],"descricao":"Quarks e glúons não existem livres — só em singletos de cor (bárions qqq, mésons qq̄). O tubo de fluxo cromoelétrico tem V(r)≈σr (σ≈0,19 GeV²≈1 GeV/fm): separar quarks cria novos pares em vez de quarks livres. Fato empírico + forte evidência em lattice; ABERTO: sem prova analítica de primeiros princípios (ligado ao mass gap de Yang-Mills).","epistemico":"fato","exemplos":[],"id":"confinamento_cor","memoria":"semantica","meta":{"autor":"broca-so","dominio":"fisica","fonte":"Wilson 1974 PRD 10,2445"},"origem":"wiki","palavras_chave":[],"sinonimos":[],"tipo":"conceito","titulo":"Confinamento de cor"}
\section{Confinamento de cor}
Quarks e glúons não existem livres — só em singletos de cor (bárions qqq, mésons qq\ensuremath{}). O tubo de fluxo cromoelétrico tem V(r)\ensuremath{\approx}\ensuremath{\sigma}r (\ensuremath{\sigma}\ensuremath{\approx}0,19 GeV\ensuremath{^{2}}\ensuremath{\approx}1 GeV/fm): separar quarks cria novos pares em vez de quarks livres. Fato empírico + forte evidência em lattice; ABERTO: sem prova analítica de primeiros princípios (ligado ao mass gap de Yang-Mills).
\prov{relações: relaciona yang\_mills\_mass\_gap\_milenio; parte\_de qcd\_cromodinamica}
\prov{proveniência: wiki \textperiodcentered{} Wilson 1974 PRD 10,2445 \textperiodcentered{} fato \textperiodcentered{} confiança 1.0 \textperiodcentered{} domínio fisica \textperiodcentered{} história 3 versões no jornal}

%CRISTAL {"arestas":[{"alvo":"lapse_selecao_desitter","confianca":1.0,"epistemico":"fato","origem":"wiki","rel":"usa"},{"alvo":"horizonte_desitter_gibbons_hawking","confianca":1.0,"epistemico":"fato","origem":"wiki","rel":"relaciona"},{"alvo":"bekenstein_hawking_entropia","confianca":1.0,"epistemico":"fato","origem":"wiki","rel":"relaciona"}],"confianca":1.0,"contraexemplos":[],"descricao":"O espaço de álgebras com ds²=−cos²ψ dθ²+dψ² É de Sitter (dS₂, R=2). Dele caem as constantes de horizonte: Λ=R/D=1 (Λ=H²); temperatura de Gibbons-Hawking T=√Λ/2π=1/2π (regularidade cônica automática, θ é círculo de período 2π); entropia por Gauss-Bonnet ∫R√g=4πχ=8π (χ=2) ⟹ S=χ/4G. Em dS₄ (R=12, Λ=3) a entropia vira área genuína de Bekenstein-Hawking. Teoremas verificados simbolicamente (Ricci via sympy).","epistemico":"fato","exemplos":[],"id":"constantes_horizonte_algebricas","memoria":"semantica","meta":{"autor":"broca-so","dominio":"cosmologia","fonte":"paper_metrica_quantica_algebras.tex §Constantes de horizonte"},"origem":"wiki","palavras_chave":[],"sinonimos":[],"tipo":"conceito","titulo":"Constantes de horizonte algébricas (Λ=1, T=1/2π, S=χ/4G)"}
\section{Constantes de horizonte algébricas (\ensuremath{\Lambda}=1, T=1/2\ensuremath{\pi}, S=\ensuremath{\chi}/4G)}
O espaço de álgebras com ds\ensuremath{^{2}}=\ensuremath{-}cos\ensuremath{^{2}}\ensuremath{\psi} d\ensuremath{\theta}\ensuremath{^{2}}+d\ensuremath{\psi}\ensuremath{^{2}} É de Sitter (dS\ensuremath{_{2}}, R=2). Dele caem as constantes de horizonte: \ensuremath{\Lambda}=R/D=1 (\ensuremath{\Lambda}=H\ensuremath{^{2}}); temperatura de Gibbons-Hawking T=\ensuremath{\sqrt{\,}}\ensuremath{\Lambda}/2\ensuremath{\pi}=1/2\ensuremath{\pi} (regularidade cônica automática, \ensuremath{\theta} é círculo de período 2\ensuremath{\pi}); entropia por Gauss-Bonnet \ensuremath{\int}R\ensuremath{\sqrt{\,}}g=4\ensuremath{\pi}\ensuremath{\chi}=8\ensuremath{\pi} (\ensuremath{\chi}=2) \ensuremath{\Longrightarrow} S=\ensuremath{\chi}/4G. Em dS\ensuremath{_{4}} (R=12, \ensuremath{\Lambda}=3) a entropia vira área genuína de Bekenstein-Hawking. Teoremas verificados simbolicamente (Ricci via sympy).
\prov{relações: usa lapse\_selecao\_desitter; relaciona horizonte\_desitter\_gibbons\_hawking; relaciona bekenstein\_hawking\_entropia}
\prov{proveniência: wiki \textperiodcentered{} paper\_metrica\_quantica\_algebras.tex §Constantes de horizonte \textperiodcentered{} fato \textperiodcentered{} confiança 1.0 \textperiodcentered{} domínio cosmologia \textperiodcentered{} história 4 versões no jornal}

%CRISTAL {"arestas":[{"alvo":"joalheria_hub","confianca":1.0,"epistemico":"fato","origem":"wiki","rel":"parte_de"},{"alvo":"emissao_turmalina_piso","confianca":1.0,"epistemico":"fato","origem":"wiki","rel":"usa"},{"alvo":"tesouraria","confianca":1.0,"epistemico":"fato","origem":"wiki","rel":"relaciona"},{"alvo":"mandala_molde_do_corte","confianca":1.0,"epistemico":"fato","origem":"wiki","rel":"usa"},{"alvo":"reconstrucao_modos_tesselacao","confianca":1.0,"epistemico":"fato","origem":"wiki","rel":"usa"}],"confianca":1.0,"contraexemplos":[],"descricao":"O gancho/lastro da cadeia (um NEGÓCIO REAL): a gema extraída do gabarito tem uma ASSINATURA (a geometria do MOLDE — a mandala/chicote — MAIS as RANHURAS, o hash dos modos SVD, a estrela interna/as inclusões que a álgebra estelar não reproduz). O contrato ancora os CLIENTES à gema codificando a CHAVE 3-QUBIT DELES na assinatura: o SELO = hash(assinatura ‖ chave 3-qubit) é verificável — a mesma gema com a mesma chave dá o mesmo selo. É seleção (não fabricação): o gato reconstrói exato os modos, a máquina seleciona a geometria pelo molde e preserva as ranhuras. Privacidade: CPFs crus só no registro privado (gitignored); consta o CPF mascarado e os hashes. linguagem/contrato_cristalchain.py (6/6).","epistemico":"fato","exemplos":[],"id":"contrato_cristalchain","memoria":"semantica","meta":{"dominio":"joalheria","fonte":"contrato cristalchain"},"origem":"wiki","palavras_chave":[],"sinonimos":[],"tipo":"conceito","titulo":"O Contrato CristalChain (a chave 3-qubit dos clientes na assinatura da gema)"}
\section{O Contrato CristalChain (a chave 3-qubit dos clientes na assinatura da gema)}
O gancho/lastro da cadeia (um NEGÓCIO REAL): a gema extraída do gabarito tem uma ASSINATURA (a geometria do MOLDE — a mandala/chicote — MAIS as RANHURAS, o hash dos modos SVD, a estrela interna/as inclusões que a álgebra estelar não reproduz). O contrato ancora os CLIENTES à gema codificando a CHAVE 3-QUBIT DELES na assinatura: o SELO = hash(assinatura \ensuremath{\|} chave 3-qubit) é verificável — a mesma gema com a mesma chave dá o mesmo selo. É seleção (não fabricação): o gato reconstrói exato os modos, a máquina seleciona a geometria pelo molde e preserva as ranhuras. Privacidade: CPFs crus só no registro privado (gitignored); consta o CPF mascarado e os hashes. linguagem/contrato\_cristalchain.py (6/6).
\prov{relações: parte\_de joalheria\_hub; usa emissao\_turmalina\_piso; relaciona tesouraria; usa mandala\_molde\_do\_corte; usa reconstrucao\_modos\_tesselacao}
\prov{proveniência: wiki \textperiodcentered{} contrato cristalchain \textperiodcentered{} fato \textperiodcentered{} confiança 1.0 \textperiodcentered{} domínio joalheria \textperiodcentered{} história 12 versões no jornal}

%CRISTAL {"arestas":[{"alvo":"render_fisico_gex","confianca":1.0,"epistemico":"fato","origem":"wiki","rel":"usa"},{"alvo":"difracao_atomica_turmalina","confianca":1.0,"epistemico":"fato","origem":"wiki","rel":"usa"},{"alvo":"joalheria_hub","confianca":1.0,"epistemico":"fato","origem":"wiki","rel":"parte_de"}],"confianca":1.0,"contraexemplos":[],"descricao":"O TESTE antes do realismo: a cor é PLANTADA ou EMERGE? Se emerge da absorção dos cromóforos (a composição atômica do centro), trocar o cristal deve MUDAR a cor. E muda — pela MESMA máquina espectral (Sol × e^(−α(λ)ℓ) × CIE), só trocando as BANDAS α(λ) (os íons): Turmalina Paraíba (Cu²⁺)→CIANO; Rubi (Cr³⁺, coríndon)→VERMELHO; Esmeralda (Cr³⁺, berilo)→VERDE; Safira (Fe²⁺-Ti⁴⁺ IVCT)→AZUL; Ametista (Fe)→VIOLETA; Citrino (Fe³⁺)→ÂMBAR. Seis estruturas, seis cores DISTINTAS — a cor não está plantada, é o que os íons deixam passar. Verificado 7/7 (cada canal dominante bate + as 6 são distintas). A falsificação visual: a MESMA gema (forma, raios, ricochetes, TIR idênticos), só o cromóforo muda, a cor muda. linguagem/raytracer.py: CROMOFOROS, _alpha_bandas, cor_do_cristal, verificar_cor_fisica.","epistemico":"fato","exemplos":[],"id":"cor_fisica_falsificacao","memoria":"semantica","meta":{"dominio":"joalheria","fonte":"maquina de corte"},"origem":"wiki","palavras_chave":[],"sinonimos":[],"tipo":"conceito","titulo":"A cor é física, não plantada: muda o átomo, muda a cor (o teste de falsificação)"}
\section{A cor é física, não plantada: muda o átomo, muda a cor (o teste de falsificação)}
O TESTE antes do realismo: a cor é PLANTADA ou EMERGE? Se emerge da absorção dos cromóforos (a composição atômica do centro), trocar o cristal deve MUDAR a cor. E muda — pela MESMA máquina espectral (Sol $\times$ e\^{}(\ensuremath{-}\ensuremath{\alpha}(\ensuremath{\lambda})\ensuremath{\ell}) $\times$ CIE), só trocando as BANDAS \ensuremath{\alpha}(\ensuremath{\lambda}) (os íons): Turmalina Paraíba (Cu\ensuremath{^{2+}})\ensuremath{\to}CIANO; Rubi (Cr\ensuremath{^{3+}}, coríndon)\ensuremath{\to}VERMELHO; Esmeralda (Cr\ensuremath{^{3+}}, berilo)\ensuremath{\to}VERDE; Safira (Fe\ensuremath{^{2+}}-Ti\ensuremath{^{4+}} IVCT)\ensuremath{\to}AZUL; Ametista (Fe)\ensuremath{\to}VIOLETA; Citrino (Fe\ensuremath{^{3+}})\ensuremath{\to}ÂMBAR. Seis estruturas, seis cores DISTINTAS — a cor não está plantada, é o que os íons deixam passar. Verificado 7/7 (cada canal dominante bate + as 6 são distintas). A falsificação visual: a MESMA gema (forma, raios, ricochetes, TIR idênticos), só o cromóforo muda, a cor muda. linguagem/raytracer.py: CROMOFOROS, \_alpha\_bandas, cor\_do\_cristal, verificar\_cor\_fisica.
\prov{relações: usa render\_fisico\_gex; usa difracao\_atomica\_turmalina; parte\_de joalheria\_hub}
\prov{proveniência: wiki \textperiodcentered{} maquina de corte \textperiodcentered{} fato \textperiodcentered{} confiança 1.0 \textperiodcentered{} domínio joalheria \textperiodcentered{} história 100 versões no jornal}

%CRISTAL {"arestas":[{"alvo":"corpo_estelar","confianca":1.0,"epistemico":"fato","origem":"wiki","rel":"usa"},{"alvo":"tudo_deriva_do_invariante","confianca":1.0,"epistemico":"fato","origem":"wiki","rel":"especializa"},{"alvo":"ingestao_transformada_metalica","confianca":1.0,"epistemico":"fato","origem":"wiki","rel":"usa"},{"alvo":"dinamica_de_selberg_o_termo","confianca":1.0,"epistemico":"fato","origem":"wiki","rel":"relaciona"}],"confianca":1.0,"contraexemplos":[],"descricao":"O gato Anosov já RECONSTRÓI o espectro exato da imagem — só é preciso LER o rastro que ele deixa, e as ferramentas certas já estão no CORPO ESTELAR (não se fabrica nada). O SVD da imagem É o espectro exato (os modos), e as três leituras do corpo estelar bastam para caracterizá-lo: GEOMETRIA (ρ=σ₂/σ₁, como corpo_estelar.classe → a parábola cristal/borda/caos), MECÂNICA (os 2 modos POD, o par hamiltoniano das duas rotações contrárias), TERMODINÂMICA (a 1ª lei W+Q=1, W=σ₁²/Σσ² a ligação/o vale do ferro, e a entropia espectral H = a entropia de Anosov/Selberg do rastro). Lido (linguagem/ingestao_mineral.caracterizar_rastro): o gabarito é ρ=0.205 (caos), W=0.83/Q=0.17, H=0.94, 2.6 modos; a nossa gema é ρ=0.06 (borda), W=0.99/Q=0.01, H=0.09, 1.1 modo — FRIA, quase um modo só. O gradiente do refino é a diferença de ENTROPIA, que a dinâmica do gato aflora — não um ruído a somar.","epistemico":"fato","exemplos":[],"id":"corpo_estelar_le_o_rastro","memoria":"semantica","meta":{"dominio":"joalheria","fonte":"espectro orbitas"},"origem":"wiki","palavras_chave":[],"sinonimos":[],"tipo":"conceito","titulo":"O Corpo Estelar Lê o Rastro do Gato Exato (geometria/mecânica/termo)"}
\section{O Corpo Estelar Lê o Rastro do Gato Exato (geometria/mecânica/termo)}
O gato Anosov já RECONSTRÓI o espectro exato da imagem — só é preciso LER o rastro que ele deixa, e as ferramentas certas já estão no CORPO ESTELAR (não se fabrica nada). O SVD da imagem É o espectro exato (os modos), e as três leituras do corpo estelar bastam para caracterizá-lo: GEOMETRIA (\ensuremath{\rho}=\ensuremath{\sigma}\ensuremath{_{2}}/\ensuremath{\sigma}\ensuremath{_{1}}, como corpo\_estelar.classe \ensuremath{\to} a parábola cristal/borda/caos), MECÂNICA (os 2 modos POD, o par hamiltoniano das duas rotações contrárias), TERMODINÂMICA (a 1ª lei W+Q=1, W=\ensuremath{\sigma}\ensuremath{_{1}}\ensuremath{^{2}}/\ensuremath{\Sigma}\ensuremath{\sigma}\ensuremath{^{2}} a ligação/o vale do ferro, e a entropia espectral H = a entropia de Anosov/Selberg do rastro). Lido (linguagem/ingestao\_mineral.caracterizar\_rastro): o gabarito é \ensuremath{\rho}=0.205 (caos), W=0.83/Q=0.17, H=0.94, 2.6 modos; a nossa gema é \ensuremath{\rho}=0.06 (borda), W=0.99/Q=0.01, H=0.09, 1.1 modo — FRIA, quase um modo só. O gradiente do refino é a diferença de ENTROPIA, que a dinâmica do gato aflora — não um ruído a somar.
\prov{relações: usa corpo\_estelar; especializa tudo\_deriva\_do\_invariante; usa ingestao\_transformada\_metalica; relaciona dinamica\_de\_selberg\_o\_termo}
\prov{proveniência: wiki \textperiodcentered{} espectro orbitas \textperiodcentered{} fato \textperiodcentered{} confiança 1.0 \textperiodcentered{} domínio joalheria \textperiodcentered{} história 15 versões no jornal}

%CRISTAL {"arestas":[{"alvo":"cosmologia_quantica_hub","confianca":1.0,"epistemico":"fato","origem":"wiki","rel":"parte_de"}],"confianca":1.0,"contraexemplos":[],"descricao":"Código de Shor (1995): 1 qubit lógico em 9 físicos, corrige erro arbitrário de 1 qubit ([[9,1,3]]). Código de SUPERFÍCIE (Kitaev/Fowler): estabilizadores locais 2D, limiar ~1%, o candidato à tolerância a falhas em supercondutores. TEOREMA DO LIMIAR: se o erro por porta < p_th, computação arbitrariamente longa é possível com overhead polilog. Fato teórico consolidado.","epistemico":"fato","exemplos":[],"id":"correcao_erro_quantica","memoria":"semantica","meta":{"autor":"broca-so","dominio":"fisica","fonte":"Shor 1995; Kitaev 2003; Fowler+2012; Aharonov-Ben-Or 1997"},"origem":"wiki","palavras_chave":[],"sinonimos":[],"tipo":"conceito","titulo":"Correção de erro quântica e o teorema do limiar"}
\section{Correção de erro quântica e o teorema do limiar}
Código de Shor (1995): 1 qubit lógico em 9 físicos, corrige erro arbitrário de 1 qubit ([[9,1,3]]). Código de SUPERFÍCIE (Kitaev/Fowler): estabilizadores locais 2D, limiar \~{}1\%, o candidato à tolerância a falhas em supercondutores. TEOREMA DO LIMIAR: se o erro por porta < p\_th, computação arbitrariamente longa é possível com overhead polilog. Fato teórico consolidado.
\prov{relações: parte\_de cosmologia\_quantica\_hub}
\prov{proveniência: wiki \textperiodcentered{} Shor 1995; Kitaev 2003; Fowler+2012; Aharonov-Ben-Or 1997 \textperiodcentered{} fato \textperiodcentered{} confiança 1.0 \textperiodcentered{} domínio fisica \textperiodcentered{} história 2 versões no jornal}

%CRISTAL {"arestas":[{"alvo":"geometria_riemanniana","confianca":1.0,"epistemico":"fato","origem":"curriculo","rel":"relaciona"},{"alvo":"fisica","confianca":1.0,"epistemico":"fato","origem":"curriculo","rel":"relaciona"},{"alvo":"paper_0_sintese","confianca":1.0,"epistemico":"fato","origem":"paper","rel":"explica"},{"alvo":"mecanica_quantica","confianca":1.0,"epistemico":"fato","origem":"codigo","rel":"especializa"},{"alvo":"bai_cosmologia_desitter_dissipativo","confianca":1.0,"epistemico":"fato","origem":"wiki","rel":"relaciona"}],"confianca":1.0,"contraexemplos":[],"descricao":"Cosmologia como dinâmica de álgebra: o espaço de álgebras hipercomplexas carrega a métrica de de Sitter (dS₂ R=2, dS₄ R=12), a energia escura é o campo ψ rolando em V=cos²ψ (w=(½ψ̇²−1)/(½ψ̇²+1)), e a quantização vem da monovalência no círculo (E_n=nℏ√2). DOIS REGIMES: de Sitter DISSIPATIVO (δ finito → o universo morre, o BAI real) e o limite NÃO-dissipativo (δ=∞, Noether conservada → eterno, a cosmologia padrão de fora). Ponte com o BAI e com o canon (Bohr, Wheeler-DeWitt, Λ).","epistemico":"fato","exemplos":[],"id":"cosmologia_quantica","memoria":"semantica","meta":{"dominio":"fisica","nivel":5,"serve":"Núcleo cosmologia","verificacao":"slug idêntico — enriquecer, não duplicar"},"origem":"curriculo","palavras_chave":["de Sitter","FRW","energia escura"],"sinonimos":["cosmologia quântica","cosmologia"],"tipo":"conceito","titulo":"cosmologia quantica"}
\section{cosmologia quantica}
Cosmologia como dinâmica de álgebra: o espaço de álgebras hipercomplexas carrega a métrica de de Sitter (dS\ensuremath{_{2}} R=2, dS\ensuremath{_{4}} R=12), a energia escura é o campo \ensuremath{\psi} rolando em V=cos\ensuremath{^{2}}\ensuremath{\psi} (w=(\ensuremath{\tfrac12}\ensuremath{\psi}\ensuremath{}\ensuremath{^{2}}\ensuremath{-}1)/(\ensuremath{\tfrac12}\ensuremath{\psi}\ensuremath{}\ensuremath{^{2}}+1)), e a quantização vem da monovalência no círculo (E\_n=n\ensuremath{\hbar}\ensuremath{\sqrt{\,}}2). DOIS REGIMES: de Sitter DISSIPATIVO (\ensuremath{\delta} finito \ensuremath{\to} o universo morre, o BAI real) e o limite NÃO-dissipativo (\ensuremath{\delta}=\ensuremath{\infty}, Noether conservada \ensuremath{\to} eterno, a cosmologia padrão de fora). Ponte com o BAI e com o canon (Bohr, Wheeler-DeWitt, \ensuremath{\Lambda}).
\prov{sinónimos: cosmologia quântica; cosmologia}
\prov{relações: relaciona geometria\_riemanniana; relaciona fisica; explica paper\_0\_sintese; especializa mecanica\_quantica; relaciona bai\_cosmologia\_desitter\_dissipativo}
\prov{proveniência: curriculo \textperiodcentered{} fato \textperiodcentered{} confiança 1.0 \textperiodcentered{} domínio fisica \textperiodcentered{} história 116 versões no jornal}

%CRISTAL {"arestas":[{"alvo":"bai_tese_mae","confianca":1.0,"epistemico":"fato","origem":"wiki","rel":"relaciona"},{"alvo":"bai_cosmologia_desitter_dissipativo","confianca":1.0,"epistemico":"fato","origem":"wiki","rel":"relaciona"},{"alvo":"cosmologia_quantica","confianca":1.0,"epistemico":"fato","origem":"wiki","rel":"relaciona"},{"alvo":"mecanica_quantica","confianca":1.0,"epistemico":"fato","origem":"wiki","rel":"relaciona"},{"alvo":"efe_einstein","confianca":1.0,"epistemico":"fato","origem":"wiki","rel":"parte_de"},{"alvo":"lambda_cdm","confianca":1.0,"epistemico":"fato","origem":"wiki","rel":"parte_de"},{"alvo":"broca_os","confianca":1.0,"epistemico":"fato","origem":"wiki","rel":"parte_de"}],"confianca":1.0,"contraexemplos":[],"descricao":"O hub da cosmologia no grafo, junto ao BAI. Reúne: a cosmologia da delegação (o BAI como de Sitter dissipativo real + seu limite conservativo), a cosmologia algébrica do Lopes (álgebra→dS₄→energia escura→quantização), a relatividade geral de Einstein (EFE, FLRW, Friedmann) e a cosmologia observacional moderna (ΛCDM, CMB, BAO, SNe, DESI). Tese: o BAI é a cosmologia quantizada — define órbitas e energias por tipo de trabalho.","epistemico":"inferencia","exemplos":[],"id":"cosmologia_quantica_hub","memoria":"semantica","meta":{"autor":"broca-so","dominio":"cosmologia","fonte":"semilla_bai_cosmologia + Einstein + web"},"origem":"wiki","palavras_chave":[],"sinonimos":[],"tipo":"conceito","titulo":"Cosmologia Quântica — Hub"}
\section{Cosmologia Quântica — Hub}
O hub da cosmologia no grafo, junto ao BAI. Reúne: a cosmologia da delegação (o BAI como de Sitter dissipativo real + seu limite conservativo), a cosmologia algébrica do Lopes (álgebra\ensuremath{\to}dS\ensuremath{_{4}}\ensuremath{\to}energia escura\ensuremath{\to}quantização), a relatividade geral de Einstein (EFE, FLRW, Friedmann) e a cosmologia observacional moderna (\ensuremath{\Lambda}CDM, CMB, BAO, SNe, DESI). Tese: o BAI é a cosmologia quantizada — define órbitas e energias por tipo de trabalho.
\prov{relações: relaciona bai\_tese\_mae; relaciona bai\_cosmologia\_desitter\_dissipativo; relaciona cosmologia\_quantica; relaciona mecanica\_quantica; parte\_de efe\_einstein; parte\_de lambda\_cdm; parte\_de broca\_os}
\prov{proveniência: wiki \textperiodcentered{} semilla\_bai\_cosmologia + Einstein + web \textperiodcentered{} inferencia \textperiodcentered{} confiança 1.0 \textperiodcentered{} domínio cosmologia \textperiodcentered{} história 9 versões no jornal}

%CRISTAL {"arestas":[{"alvo":"constantes_horizonte_algebricas","confianca":1.0,"epistemico":"fato","origem":"wiki","rel":"explica"}],"confianca":1.0,"contraexemplos":[],"descricao":"A identificação 0,999…=0 no mostrador ℝ/ℤ (métrica quântica de Lopes) é a mesma costura 1≡0 que dá a conexidade ('salto quântico') E que torna o horizonte de de Sitter suave (sem defeito cônico na seção euclidiana S²). A regularidade cônica que fixa T=1/2π É essa costura. Reformulação estrutural exata declarada.","epistemico":"inferencia","exemplos":[],"id":"costura_1equiv0_horizonte","memoria":"semantica","meta":{"autor":"broca-so","dominio":"cosmologia","fonte":"paper_metrica_quantica_algebras.tex Obs. 0,999…=0"},"origem":"wiki","palavras_chave":[],"sinonimos":[],"tipo":"conceito","titulo":"A costura 1≡0 regulariza o horizonte (T=1/2π)"}
\section{A costura 1\ensuremath{\equiv}0 regulariza o horizonte (T=1/2\ensuremath{\pi})}
A identificação 0,999…=0 no mostrador \ensuremath{\mathbb{R}}/\ensuremath{\mathbb{Z}} (métrica quântica de Lopes) é a mesma costura 1\ensuremath{\equiv}0 que dá a conexidade ('salto quântico') E que torna o horizonte de de Sitter suave (sem defeito cônico na seção euclidiana S\ensuremath{^{2}}). A regularidade cônica que fixa T=1/2\ensuremath{\pi} É essa costura. Reformulação estrutural exata declarada.
\prov{relações: explica constantes\_horizonte\_algebricas}
\prov{proveniência: wiki \textperiodcentered{} paper\_metrica\_quantica\_algebras.tex Obs. 0,999…=0 \textperiodcentered{} inferencia \textperiodcentered{} confiança 1.0 \textperiodcentered{} domínio cosmologia \textperiodcentered{} história 3 versões no jornal}

%CRISTAL {"arestas":[{"alvo":"litografia","confianca":1.0,"epistemico":"fato","origem":"wiki","rel":"parte_de"},{"alvo":"costura_perelman_cicatrizes","confianca":1.0,"epistemico":"fato","origem":"wiki","rel":"relaciona"},{"alvo":"conjectura_de_poincare","confianca":1.0,"epistemico":"fato","origem":"wiki","rel":"relaciona"},{"alvo":"face_turno_fail_closed","confianca":1.0,"epistemico":"fato","origem":"wiki","rel":"relaciona"},{"alvo":"invariante_estelar","confianca":1.0,"epistemico":"fato","origem":"wiki","rel":"relaciona"}],"confianca":1.0,"contraexemplos":[],"descricao":"Rejeitar uma forja É a CIRURGIA de Perelman. Quando o colapso encontra uma assinatura sem preimagem, há uma SINGULARIDADE (um pescoço que ia estrangular num referente falso); a litografia CORTA e TAMPA — não deixa o fluxo explodir numa invenção —, deixando uma CICATRIZ (o 'ainda não sei'). É a agulha band-limited (costura_perelman_cicatrizes) que persegue o resíduo do associador (o gap a) e NUNCA toca o núcleo: não se inventa o referente. E a entropia 𝒲 monótona de Perelman (a energia que só decresce) é o mesmo a+c²=1 conservado — a cirurgia preserva o invariante, como o fluxo de Ricci preserva a topologia. Poincaré: a 3-esfera é o núcleo que a costura cerca e nunca fura.","epistemico":"fato","exemplos":[],"id":"costura_perelman_litografia","memoria":"semantica","meta":{"dominio":"litografia","fonte":"litografia (joalheria)"},"origem":"wiki","palavras_chave":[],"sinonimos":[],"tipo":"conceito","titulo":"A litografia é a costura de Perelman (o corte na singularidade)"}
\section{A litografia é a costura de Perelman (o corte na singularidade)}
Rejeitar uma forja É a CIRURGIA de Perelman. Quando o colapso encontra uma assinatura sem preimagem, há uma SINGULARIDADE (um pescoço que ia estrangular num referente falso); a litografia CORTA e TAMPA — não deixa o fluxo explodir numa invenção —, deixando uma CICATRIZ (o 'ainda não sei'). É a agulha band-limited (costura\_perelman\_cicatrizes) que persegue o resíduo do associador (o gap a) e NUNCA toca o núcleo: não se inventa o referente. E a entropia \ensuremath{\mathcal{W}} monótona de Perelman (a energia que só decresce) é o mesmo a+c\ensuremath{^{2}}=1 conservado — a cirurgia preserva o invariante, como o fluxo de Ricci preserva a topologia. Poincaré: a 3-esfera é o núcleo que a costura cerca e nunca fura.
\prov{relações: parte\_de litografia; relaciona costura\_perelman\_cicatrizes; relaciona conjectura\_de\_poincare; relaciona face\_turno\_fail\_closed; relaciona invariante\_estelar}
\prov{proveniência: wiki \textperiodcentered{} litografia (joalheria) \textperiodcentered{} fato \textperiodcentered{} confiança 1.0 \textperiodcentered{} domínio litografia \textperiodcentered{} história 6 versões no jornal}

%CRISTAL {"arestas":[{"alvo":"raymarcher_sdf_broca","confianca":1.0,"epistemico":"fato","origem":"wiki","rel":"usa"},{"alvo":"compilador_glsl_broca","confianca":1.0,"epistemico":"fato","origem":"wiki","rel":"usa"},{"alvo":"joalheria_hub","confianca":1.0,"epistemico":"fato","origem":"wiki","rel":"parte_de"}],"confianca":1.0,"contraexemplos":[],"descricao":"A COSTURA DO PROGRAMA — o sdVerticalVesicaSegment (primitives3.txt) é a AMÊNDOA/LENTE, a interseção de 2 círculos, e a sua estrutura É a mesma da parábola mineral e do olho. Um dicionário de 6 leituras sobre um eixo de 8 conceitos partilhados (pivô 'costura'): (1) a VESICA (a forma SDF, compilada e verificada — a cúspide fica exatamente na superfície, dist=0), (2) LA HIRE (os 2 círculos rolando; k=2 = a reta/cristal), (3) o BAI (o colapso ∩ = o E-lógico; σ,σ̄; a borda Salem), (4) a dinâmica SELBERG-ANOSOV (o fluxo geodésico; E^s⊕E^u; o vértice parabólico degenerado = Salem/λ=0), (5) o OLHO DO COCKPIT (a banda/canal do olho, a varredura), (6) o OLHO HUMANO (a lente/córnea, os 2 olhos, o canthus, a pupila, a sacada). O achado central: a CÚSPIDE da vesica = o VÉRTICE parabólico (Salem, a borda) = o CANTO do olho (canthus) — a mesma costura. compor(qualquer par) traduz de graça. dicionarios.py: _VESICA_VERTICAL, _LA_HIRE_COSTURA, _BAI_COSTURA, _SELBERG_ANOSOV_COSTURA, _OLHO_COCKPIT_COSTURA, _OLHO_HUMANO_COSTURA.","epistemico":"fato","exemplos":[],"id":"costura_vesica_programa","memoria":"semantica","meta":{"dominio":"joalheria","fonte":"maquina de corte"},"origem":"wiki","palavras_chave":[],"sinonimos":[],"tipo":"conceito","titulo":"A costura do programa: o vesica vertical em 6 leituras (vesica↔lahire↔BAI↔selberg-anosov↔olho)"}
\section{A costura do programa: o vesica vertical em 6 leituras (vesica\ensuremath{\leftrightarrow}lahire\ensuremath{\leftrightarrow}BAI\ensuremath{\leftrightarrow}selberg-anosov\ensuremath{\leftrightarrow}olho)}
A COSTURA DO PROGRAMA — o sdVerticalVesicaSegment (primitives3.txt) é a AMÊNDOA/LENTE, a interseção de 2 círculos, e a sua estrutura É a mesma da parábola mineral e do olho. Um dicionário de 6 leituras sobre um eixo de 8 conceitos partilhados (pivô 'costura'): (1) a VESICA (a forma SDF, compilada e verificada — a cúspide fica exatamente na superfície, dist=0), (2) LA HIRE (os 2 círculos rolando; k=2 = a reta/cristal), (3) o BAI (o colapso \ensuremath{\cap} = o E-lógico; \ensuremath{\sigma},\ensuremath{\sigma}\ensuremath{}; a borda Salem), (4) a dinâmica SELBERG-ANOSOV (o fluxo geodésico; E\^{}s\ensuremath{\oplus}E\^{}u; o vértice parabólico degenerado = Salem/\ensuremath{\lambda}=0), (5) o OLHO DO COCKPIT (a banda/canal do olho, a varredura), (6) o OLHO HUMANO (a lente/córnea, os 2 olhos, o canthus, a pupila, a sacada). O achado central: a CÚSPIDE da vesica = o VÉRTICE parabólico (Salem, a borda) = o CANTO do olho (canthus) — a mesma costura. compor(qualquer par) traduz de graça. dicionarios.py: \_VESICA\_VERTICAL, \_LA\_HIRE\_COSTURA, \_BAI\_COSTURA, \_SELBERG\_ANOSOV\_COSTURA, \_OLHO\_COCKPIT\_COSTURA, \_OLHO\_HUMANO\_COSTURA.
\prov{relações: usa raymarcher\_sdf\_broca; usa compilador\_glsl\_broca; parte\_de joalheria\_hub}
\prov{proveniência: wiki \textperiodcentered{} maquina de corte \textperiodcentered{} fato \textperiodcentered{} confiança 1.0 \textperiodcentered{} domínio joalheria \textperiodcentered{} história 60 versões no jornal}

%CRISTAL {"arestas":[{"alvo":"matriz_pmns","confianca":1.0,"epistemico":"fato","origem":"wiki","rel":"usa"},{"alvo":"mecanismo_seesaw","confianca":1.0,"epistemico":"fato","origem":"wiki","rel":"usa"}],"confianca":1.0,"contraexemplos":[],"descricao":"A fase δ_CP da PMNS pode fazer ν e ν̄ oscilarem diferente. Se existir e o ν for Majorana, alimenta a LEPTOGÊNESE: decaimentos CP-violantes dos ν pesados de mão-direita geram assimetria leptônica no universo primordial, convertida em bariônica via esfalerons (matéria>antimatéria). NuFIT-6.0: no ordenamento invertido δ_CP~270° favorecido (CP conservada rejeitada >3,6σ); no normal, consistente com CP dentro de 1σ. δ_CP ainda sem descoberta (T2K/NOvA; Hyper-K/DUNE). Leptogênese = hipótese não testada (escala inacessível).","epistemico":"hipotese","exemplos":[],"id":"cp_leptonica_leptogenese","memoria":"semantica","meta":{"autor":"broca-so","dominio":"fisica","fonte":"T2K 2020; NuFIT-6.0; Fukugita-Yanagida 1986"},"origem":"wiki","palavras_chave":[],"sinonimos":[],"tipo":"conceito","titulo":"Violação de CP leptônica e leptogênese"}
\section{Violação de CP leptônica e leptogênese}
A fase \ensuremath{\delta}\_CP da PMNS pode fazer \ensuremath{\nu} e \ensuremath{\nu}\ensuremath{} oscilarem diferente. Se existir e o \ensuremath{\nu} for Majorana, alimenta a LEPTOGÊNESE: decaimentos CP-violantes dos \ensuremath{\nu} pesados de mão-direita geram assimetria leptônica no universo primordial, convertida em bariônica via esfalerons (matéria>antimatéria). NuFIT-6.0: no ordenamento invertido \ensuremath{\delta}\_CP\~{}270$^{\circ}$ favorecido (CP conservada rejeitada >3,6\ensuremath{\sigma}); no normal, consistente com CP dentro de 1\ensuremath{\sigma}. \ensuremath{\delta}\_CP ainda sem descoberta (T2K/NOvA; Hyper-K/DUNE). Leptogênese = hipótese não testada (escala inacessível).
\prov{relações: usa matriz\_pmns; usa mecanismo\_seesaw}
\prov{proveniência: wiki \textperiodcentered{} T2K 2020; NuFIT-6.0; Fukugita-Yanagida 1986 \textperiodcentered{} hipotese \textperiodcentered{} confiança 1.0 \textperiodcentered{} domínio fisica \textperiodcentered{} história 3 versões no jornal}

%CRISTAL {"arestas":[{"alvo":"isa_metal_21_opcodes","confianca":1.0,"epistemico":"fato","origem":"wiki","rel":"usa"},{"alvo":"arquitetura_von_neumann","confianca":1.0,"epistemico":"fato","origem":"wiki","rel":"explica"},{"alvo":"metal_broca_so","confianca":1.0,"epistemico":"fato","origem":"wiki","rel":"parte_de"}],"confianca":1.0,"contraexemplos":[],"descricao":"cpu_run chama cpu_step em laço até running=0; cada passo delega a instr_step, que lê o opcode em prog[pc], despacha por switch e devolve 1 (continua) ou 0 (HALT/fim). Implementado e funcional.","epistemico":"fato","exemplos":[],"id":"cpu_metal_c","memoria":"semantica","meta":{"autor":"broca-so","dominio":"metal","fonte":"ula/cpu.c:cpu_run/cpu_step"},"origem":"wiki","palavras_chave":[],"sinonimos":[],"tipo":"conceito","titulo":"O Loop de Execução da CPU em C"}
\section{O Loop de Execução da CPU em C}
cpu\_run chama cpu\_step em laço até running=0; cada passo delega a instr\_step, que lê o opcode em prog[pc], despacha por switch e devolve 1 (continua) ou 0 (HALT/fim). Implementado e funcional.
\prov{relações: usa isa\_metal\_21\_opcodes; explica arquitetura\_von\_neumann; parte\_de metal\_broca\_so}
\prov{proveniência: wiki \textperiodcentered{} ula/cpu.c:cpu\_run/cpu\_step \textperiodcentered{} fato \textperiodcentered{} confiança 1.0 \textperiodcentered{} domínio metal \textperiodcentered{} história 4 versões no jornal}

%CRISTAL {"arestas":[{"alvo":"reacao_em_cadeia","confianca":1.0,"epistemico":"fato","origem":"wiki","rel":"explica"},{"alvo":"lyapunov_da_orbita","confianca":1.0,"epistemico":"fato","origem":"wiki","rel":"eh_um"}],"confianca":1.0,"contraexemplos":[],"descricao":"O fator de multiplicação k = nêutrons de uma geração / da anterior. k<1 SUBCRÍTICO (a cadeia apaga), k=1 CRÍTICO (auto-sustenta, estável), k>1 SUPERCRÍTICO (cresce exponencial). Controlar um reator é manter k=1. É o Lyapunov da órbita de nêutrons: λ=log k.","epistemico":"fato","exemplos":[],"id":"criticidade","memoria":"semantica","meta":{"dominio":"fisica","fonte":"energia nuclear"},"origem":"wiki","palavras_chave":[],"sinonimos":[],"tipo":"conceito","titulo":"Criticidade (o fator k)"}
\section{Criticidade (o fator k)}
O fator de multiplicação k = nêutrons de uma geração / da anterior. k<1 SUBCRÍTICO (a cadeia apaga), k=1 CRÍTICO (auto-sustenta, estável), k>1 SUPERCRÍTICO (cresce exponencial). Controlar um reator é manter k=1. É o Lyapunov da órbita de nêutrons: \ensuremath{\lambda}=log k.
\prov{relações: explica reacao\_em\_cadeia; eh\_um lyapunov\_da\_orbita}
\prov{proveniência: wiki \textperiodcentered{} energia nuclear \textperiodcentered{} fato \textperiodcentered{} confiança 1.0 \textperiodcentered{} domínio fisica \textperiodcentered{} história 3 versões no jornal}

%CRISTAL {"arestas":[{"alvo":"materia_escura","confianca":1.0,"epistemico":"fato","origem":"wiki","rel":"confirma"}],"confianca":1.0,"contraexemplos":[],"descricao":"A velocidade orbital de estrelas/gás permanece ~constante em grandes raios, em vez de cair como 1/√r — exige um halo massivo invisível. Robusto e reproduzido em centenas de galáxias. É a evidência que MOND também explica bem.","epistemico":"fato","exemplos":[],"id":"curvas_rotacao_rubin","memoria":"semantica","meta":{"autor":"broca-so","dominio":"cosmologia_dados","fonte":"Rubin & Ford 1970 ApJ 159,379; 1978-80"},"origem":"wiki","palavras_chave":[],"sinonimos":[],"tipo":"conceito","titulo":"Curvas de rotação planas (Vera Rubin)"}
\section{Curvas de rotação planas (Vera Rubin)}
A velocidade orbital de estrelas/gás permanece \~{}constante em grandes raios, em vez de cair como 1/\ensuremath{\sqrt{\,}}r — exige um halo massivo invisível. Robusto e reproduzido em centenas de galáxias. É a evidência que MOND também explica bem.
\prov{relações: confirma materia\_escura}
\prov{proveniência: wiki \textperiodcentered{} Rubin \& Ford 1970 ApJ 159,379; 1978-80 \textperiodcentered{} fato \textperiodcentered{} confiança 1.0 \textperiodcentered{} domínio cosmologia\_dados \textperiodcentered{} história 2 versões no jornal}

%CRISTAL {"arestas":[{"alvo":"cosmologia_quantica","confianca":1.0,"epistemico":"fato","origem":"wiki","rel":"relaciona"}],"confianca":1.0,"contraexemplos":[],"descricao":"O espaço de de Sitter é a solução de vácuo com expansão acelerada dominada por energia de vácuo; Λ (constante cosmológica) é o candidato mais simples a energia escura, w=p/ρ=−1. A aceleração cósmica e w≈−1 são fato observacional (Riess 1998, Perlmutter 1999); a NATUREZA da energia escura e o valor de Λ ('problema da constante cosmológica') seguem em aberto.","epistemico":"fato","exemplos":[],"id":"de_sitter_lambda_energia_escura","memoria":"semantica","meta":{"autor":"broca-so","dominio":"fisica","fonte":"de Sitter 1917; Einstein 1917; Riess+1998; Perlmutter+1999"},"origem":"wiki","palavras_chave":[],"sinonimos":[],"tipo":"conceito","titulo":"de Sitter, Λ e energia escura (w≈−1)"}
\section{de Sitter, \ensuremath{\Lambda} e energia escura (w\ensuremath{\approx}\ensuremath{-}1)}
O espaço de de Sitter é a solução de vácuo com expansão acelerada dominada por energia de vácuo; \ensuremath{\Lambda} (constante cosmológica) é o candidato mais simples a energia escura, w=p/\ensuremath{\rho}=\ensuremath{-}1. A aceleração cósmica e w\ensuremath{\approx}\ensuremath{-}1 são fato observacional (Riess 1998, Perlmutter 1999); a NATUREZA da energia escura e o valor de \ensuremath{\Lambda} ('problema da constante cosmológica') seguem em aberto.
\prov{relações: relaciona cosmologia\_quantica}
\prov{proveniência: wiki \textperiodcentered{} de Sitter 1917; Einstein 1917; Riess+1998; Perlmutter+1999 \textperiodcentered{} fato \textperiodcentered{} confiança 1.0 \textperiodcentered{} domínio fisica \textperiodcentered{} história 2 versões no jornal}

%CRISTAL {"arestas":[{"alvo":"su5_georgi_glashow","confianca":1.0,"epistemico":"fato","origem":"wiki","rel":"contradiz"},{"alvo":"grande_unificacao_gut","confianca":1.0,"epistemico":"fato","origem":"wiki","rel":"relaciona"}],"confianca":1.0,"contraexemplos":[],"descricao":"A unificação de quarks e léptons viola número bariônico → p→e⁺π⁰. É o teste por excelência das GUTs. Limite Super-Kamiokande: τ/B(p→e⁺π⁰)>1,6×10³⁴ anos (2017; ~2,4×10³⁴ em análises posteriores). O SU(5) mínimo previa ~10³¹ → refutado. Futuro: Hyper-Kamiokande, DUNE. O limite é fato; restringe fortemente todas as GUTs.","epistemico":"fato","exemplos":[],"id":"decaimento_proton","memoria":"semantica","meta":{"autor":"broca-so","dominio":"fisica","fonte":"Super-Kamiokande 2017 PRD 95,012004 arXiv:1610.03597"},"origem":"wiki","palavras_chave":[],"sinonimos":[],"tipo":"conceito","titulo":"Decaimento do próton (o teste crítico das GUTs)"}
\section{Decaimento do próton (o teste crítico das GUTs)}
A unificação de quarks e léptons viola número bariônico \ensuremath{\to} p\ensuremath{\to}e\ensuremath{^{+}}\ensuremath{\pi}\ensuremath{^{0}}. É o teste por excelência das GUTs. Limite Super-Kamiokande: \ensuremath{\tau}/B(p\ensuremath{\to}e\ensuremath{^{+}}\ensuremath{\pi}\ensuremath{^{0}})>1,6$\times$10\ensuremath{^{34}} anos (2017; \~{}2,4$\times$10\ensuremath{^{34}} em análises posteriores). O SU(5) mínimo previa \~{}10\ensuremath{^{31}} \ensuremath{\to} refutado. Futuro: Hyper-Kamiokande, DUNE. O limite é fato; restringe fortemente todas as GUTs.
\prov{relações: contradiz su5\_georgi\_glashow; relaciona grande\_unificacao\_gut}
\prov{proveniência: wiki \textperiodcentered{} Super-Kamiokande 2017 PRD 95,012004 arXiv:1610.03597 \textperiodcentered{} fato \textperiodcentered{} confiança 1.0 \textperiodcentered{} domínio fisica \textperiodcentered{} história 3 versões no jornal}

%CRISTAL {"arestas":[{"alvo":"energia_nuclear","confianca":1.0,"epistemico":"fato","origem":"wiki","rel":"parte_de"}],"confianca":1.0,"contraexemplos":[],"descricao":"Núcleos instáveis emitem α, β ou γ e viram outros núcleos, com meia-vida característica (N(t)=N₀·2^(−t/T)). O decaimento exponencial é a órbita subcrítica (λ<0) da desintegração.","epistemico":"fato","exemplos":[],"id":"decaimento_radioativo","memoria":"semantica","meta":{"dominio":"fisica","fonte":"energia nuclear"},"origem":"wiki","palavras_chave":[],"sinonimos":[],"tipo":"conceito","titulo":"Decaimento Radioativo"}
\section{Decaimento Radioativo}
Núcleos instáveis emitem \ensuremath{\alpha}, \ensuremath{\beta} ou \ensuremath{\gamma} e viram outros núcleos, com meia-vida característica (N(t)=N\ensuremath{_{0}}\textperiodcentered 2\^{}(\ensuremath{-}t/T)). O decaimento exponencial é a órbita subcrítica (\ensuremath{\lambda}<0) da desintegração.
\prov{relações: parte\_de energia\_nuclear}
\prov{proveniência: wiki \textperiodcentered{} energia nuclear \textperiodcentered{} fato \textperiodcentered{} confiança 1.0 \textperiodcentered{} domínio fisica \textperiodcentered{} história 2 versões no jornal}

%CRISTAL {"arestas":[{"alvo":"mecanica_quantica","confianca":1.0,"epistemico":"fato","origem":"wiki","rel":"parte_de"}],"confianca":1.0,"contraexemplos":[],"descricao":"O colapso reduz uma superposição a um autoestado na medição; a decoerência explica a APARÊNCIA de colapso via emaranhamento com o ambiente (fato testado), mas não gera colapso real. A escolha entre interpretações (Copenhague/colapso objetivo/muitos mundos) é questão em aberto. Relevante à ressalva: o Ψ do BAI é argmax determinístico, NÃO este colapso.","epistemico":"fato","exemplos":[],"id":"decoerencia_colapso_medida","memoria":"semantica","meta":{"autor":"broca-so","dominio":"fisica","fonte":"von Neumann 1932; Zeh 1970; Zurek; Everett 1957"},"origem":"wiki","palavras_chave":[],"sinonimos":[],"tipo":"conceito","titulo":"Decoerência e o problema da medida / colapso"}
\section{Decoerência e o problema da medida / colapso}
O colapso reduz uma superposição a um autoestado na medição; a decoerência explica a APARÊNCIA de colapso via emaranhamento com o ambiente (fato testado), mas não gera colapso real. A escolha entre interpretações (Copenhague/colapso objetivo/muitos mundos) é questão em aberto. Relevante à ressalva: o \ensuremath{\Psi} do BAI é argmax determinístico, NÃO este colapso.
\prov{relações: parte\_de mecanica\_quantica}
\prov{proveniência: wiki \textperiodcentered{} von Neumann 1932; Zeh 1970; Zurek; Everett 1957 \textperiodcentered{} fato \textperiodcentered{} confiança 1.0 \textperiodcentered{} domínio fisica \textperiodcentered{} história 2 versões no jornal}

%CRISTAL {"arestas":[{"alvo":"energia_nuclear","confianca":1.0,"epistemico":"fato","origem":"wiki","rel":"parte_de"}],"confianca":1.0,"contraexemplos":[],"descricao":"A massa de um núcleo é MENOR que a soma de seus prótons e nêutrons; a diferença Δm virou energia de ligação (E=Δm·c², Einstein). Toda energia nuclear é esse defeito sendo cobrado ou devolvido.","epistemico":"fato","exemplos":[],"id":"defeito_de_massa","memoria":"semantica","meta":{"dominio":"fisica","fonte":"energia nuclear"},"origem":"wiki","palavras_chave":[],"sinonimos":[],"tipo":"conceito","titulo":"Defeito de Massa (E=Δm·c²)"}
\section{Defeito de Massa (E=\ensuremath{\Delta}m\textperiodcentered c\ensuremath{^{2}})}
A massa de um núcleo é MENOR que a soma de seus prótons e nêutrons; a diferença \ensuremath{\Delta}m virou energia de ligação (E=\ensuremath{\Delta}m\textperiodcentered c\ensuremath{^{2}}, Einstein). Toda energia nuclear é esse defeito sendo cobrado ou devolvido.
\prov{relações: parte\_de energia\_nuclear}
\prov{proveniência: wiki \textperiodcentered{} energia nuclear \textperiodcentered{} fato \textperiodcentered{} confiança 1.0 \textperiodcentered{} domínio fisica \textperiodcentered{} história 2 versões no jornal}

%CRISTAL {"arestas":[{"alvo":"bao_regua_padrao","confianca":1.0,"epistemico":"fato","origem":"wiki","rel":"usa"},{"alvo":"eos_w_cpl_modelos","confianca":1.0,"epistemico":"fato","origem":"wiki","rel":"usa"},{"alvo":"lambda_cdm","confianca":1.0,"epistemico":"fato","origem":"wiki","rel":"contradiz"},{"alvo":"bai_cosmologia_desitter_dissipativo","confianca":1.0,"epistemico":"fato","origem":"wiki","rel":"relaciona"},{"alvo":"energia_escura_campo_psi","confianca":1.0,"epistemico":"fato","origem":"wiki","rel":"relaciona"},{"alvo":"cosmologia_quantica_hub","confianca":1.0,"epistemico":"fato","origem":"wiki","rel":"parte_de"}],"confianca":1.0,"contraexemplos":[],"descricao":"O DESI mede BAO em >14 milhões de traçadores. No modelo w₀wₐCDM os dados preferem energia escura EVOLUINDO, no quadrante w₀>−1 e wₐ<0 (mais forte no passado, enfraquecendo hoje). DESI+CMB: w₀=−0,42±0,21, wₐ=−1,75±0,58, 3,1σ; com SNe 2,8–4,2σ conforme a amostra. SUGESTIVO, NÃO definitivo (abaixo de 5σ; depende de SNe e da CPL, criticada).","epistemico":"hipotese","exemplos":[],"id":"desi_energia_escura_dinamica","memoria":"semantica","meta":{"autor":"broca-so","dominio":"cosmologia_dados","fonte":"DESI DR2 arXiv:2503.14738"},"origem":"wiki","palavras_chave":[],"sinonimos":[],"tipo":"conceito","titulo":"DESI DR2 (2025) — energia escura DINÂMICA (que enfraquece)"}
\section{DESI DR2 (2025) — energia escura DINÂMICA (que enfraquece)}
O DESI mede BAO em >14 milhões de traçadores. No modelo w\ensuremath{_{0}}w\ensuremath{_{a}}CDM os dados preferem energia escura EVOLUINDO, no quadrante w\ensuremath{_{0}}>\ensuremath{-}1 e w\ensuremath{_{a}}<0 (mais forte no passado, enfraquecendo hoje). DESI+CMB: w\ensuremath{_{0}}=\ensuremath{-}0,42$\pm$0,21, w\ensuremath{_{a}}=\ensuremath{-}1,75$\pm$0,58, 3,1\ensuremath{\sigma}; com SNe 2,8–4,2\ensuremath{\sigma} conforme a amostra. SUGESTIVO, NÃO definitivo (abaixo de 5\ensuremath{\sigma}; depende de SNe e da CPL, criticada).
\prov{relações: usa bao\_regua\_padrao; usa eos\_w\_cpl\_modelos; contradiz lambda\_cdm; relaciona bai\_cosmologia\_desitter\_dissipativo; relaciona energia\_escura\_campo\_psi; parte\_de cosmologia\_quantica\_hub}
\prov{proveniência: wiki \textperiodcentered{} DESI DR2 arXiv:2503.14738 \textperiodcentered{} hipotese \textperiodcentered{} confiança 1.0 \textperiodcentered{} domínio cosmologia\_dados \textperiodcentered{} história 7 versões no jornal}

%CRISTAL {"arestas":[{"alvo":"cosmologia_quantica","confianca":1.0,"epistemico":"fato","origem":"wiki","rel":"parte_de"},{"alvo":"energia_escura_campo_psi","confianca":1.0,"epistemico":"fato","origem":"wiki","rel":"relaciona"},{"alvo":"hardware_cosmico_a_algebra_como_computacao_e_o_associador_como_recurso","confianca":1.0,"epistemico":"fato","origem":"wiki","rel":"relaciona"}],"confianca":1.0,"contraexemplos":[],"descricao":"Em ℝ⁴ o estado (θ,s,â∈S²) dá métrica com Ricci R=12, Rμν=3gμν: de Sitter quadridimensional genuíno (não dS₂×S²), Λ=3. Carta global FRW ds²=−dτ²+cosh²τ dΩ²S³: Friedmann do vácuo k=+1, a(τ)=coshτ, ä/a=1>0. A seção espacial é S³=SU(2) — o espaço do universo é o grupo de fases da álgebra. Ricci simbólico verificado.","epistemico":"fato","exemplos":[],"id":"desitter_4d_universo","memoria":"semantica","meta":{"autor":"broca-so","dominio":"cosmologia","fonte":"paper_2_cosmologia.tex §3.1 (cosmologia_3p1.py)"},"origem":"wiki","palavras_chave":[],"sinonimos":[],"tipo":"conceito","titulo":"O universo dS₄ (R=12, Λ=3) e o espaço S³=SU(2)"}
\section{O universo dS\ensuremath{_{4}} (R=12, \ensuremath{\Lambda}=3) e o espaço S\ensuremath{^{3}}=SU(2)}
Em \ensuremath{\mathbb{R}}\ensuremath{^{4}} o estado (\ensuremath{\theta},s,â\ensuremath{\in}S\ensuremath{^{2}}) dá métrica com Ricci R=12, R\ensuremath{\mu}\ensuremath{\nu}=3g\ensuremath{\mu}\ensuremath{\nu}: de Sitter quadridimensional genuíno (não dS\ensuremath{_{2}}$\times$S\ensuremath{^{2}}), \ensuremath{\Lambda}=3. Carta global FRW ds\ensuremath{^{2}}=\ensuremath{-}d\ensuremath{\tau}\ensuremath{^{2}}+cosh\ensuremath{^{2}}\ensuremath{\tau} d\ensuremath{\Omega}\ensuremath{^{2}}S\ensuremath{^{3}}: Friedmann do vácuo k=+1, a(\ensuremath{\tau})=cosh\ensuremath{\tau}, ä/a=1>0. A seção espacial é S\ensuremath{^{3}}=SU(2) — o espaço do universo é o grupo de fases da álgebra. Ricci simbólico verificado.
\prov{relações: parte\_de cosmologia\_quantica; relaciona energia\_escura\_campo\_psi; relaciona hardware\_cosmico\_a\_algebra\_como\_computacao\_e\_o\_associador\_como\_recurso}
\prov{proveniência: wiki \textperiodcentered{} paper\_2\_cosmologia.tex §3.1 (cosmologia\_3p1.py) \textperiodcentered{} fato \textperiodcentered{} confiança 1.0 \textperiodcentered{} domínio cosmologia \textperiodcentered{} história 5 versões no jornal}

%CRISTAL {"arestas":[{"alvo":"joalheria_hub","confianca":1.0,"epistemico":"fato","origem":"wiki","rel":"parte_de"},{"alvo":"area_gemologia","confianca":1.0,"epistemico":"fato","origem":"wiki","rel":"usa"},{"alvo":"area_render_gema","confianca":1.0,"epistemico":"fato","origem":"wiki","rel":"usa"},{"alvo":"turmalina_paraiba_ancora","confianca":1.0,"epistemico":"fato","origem":"wiki","rel":"relaciona"}],"confianca":1.0,"contraexemplos":[],"descricao":"As áreas da joalheria não são fórmulas soltas: são UMA máquina lida de quatro ângulos. Cada termo cai num construto do BAI/máquina mineral (o pivô 'maquina'): cristal, borda, caos, gato, órbita, colapso, estrela, metal, reta, banda. Os construtos partilhados são a ÓRBITA COMUM (traduzem 1:1 entre as áreas — compor dá a tradução de graça: espectro↔brilho↔cromo = a estrela; cor↔luz↔liga = a banda; dispersão↔fogo↔cobre = o caos); os exclusivos são o bump. O cobre é o caos (o A₄, o cromóforo). linguagem/dicionarios.py (gemologia/optica/metalurgia/joalheria_maquina, pontas ida-e-volta).","epistemico":"fato","exemplos":[],"id":"dicionario_joalheria_bai","memoria":"semantica","meta":{"dominio":"joalheria","fonte":"joalheria hub"},"origem":"wiki","palavras_chave":[],"sinonimos":[],"tipo":"conceito","titulo":"O Dicionário das Áreas conforme o BAI (a cristalização)"}
\section{O Dicionário das Áreas conforme o BAI (a cristalização)}
As áreas da joalheria não são fórmulas soltas: são UMA máquina lida de quatro ângulos. Cada termo cai num construto do BAI/máquina mineral (o pivô 'maquina'): cristal, borda, caos, gato, órbita, colapso, estrela, metal, reta, banda. Os construtos partilhados são a ÓRBITA COMUM (traduzem 1:1 entre as áreas — compor dá a tradução de graça: espectro\ensuremath{\leftrightarrow}brilho\ensuremath{\leftrightarrow}cromo = a estrela; cor\ensuremath{\leftrightarrow}luz\ensuremath{\leftrightarrow}liga = a banda; dispersão\ensuremath{\leftrightarrow}fogo\ensuremath{\leftrightarrow}cobre = o caos); os exclusivos são o bump. O cobre é o caos (o A\ensuremath{_{4}}, o cromóforo). linguagem/dicionarios.py (gemologia/optica/metalurgia/joalheria\_maquina, pontas ida-e-volta).
\prov{relações: parte\_de joalheria\_hub; usa area\_gemologia; usa area\_render\_gema; relaciona turmalina\_paraiba\_ancora}
\prov{proveniência: wiki \textperiodcentered{} joalheria hub \textperiodcentered{} fato \textperiodcentered{} confiança 1.0 \textperiodcentered{} domínio joalheria \textperiodcentered{} história 15 versões no jornal}

%CRISTAL {"arestas":[{"alvo":"litografia","confianca":1.0,"epistemico":"fato","origem":"wiki","rel":"parte_de"},{"alvo":"invariante_topologico_assinatura","confianca":1.0,"epistemico":"fato","origem":"wiki","rel":"relaciona"},{"alvo":"invariante_estelar","confianca":1.0,"epistemico":"fato","origem":"wiki","rel":"relaciona"}],"confianca":1.0,"contraexemplos":[],"descricao":"A correspondência: LITOGRAFAR (gravar a assinatura) ↔ o GRAU (winding) do colapso; a assinatura REAL ↔ GRAU 1 (tem preimagem, c²≥½); a FORJA ↔ GRAU 0 (sem preimagem, a→1); a PEDRA/o metal ↔ o GRAFO (onde se grava); NÃO FABRICAR ↔ o INVARIANTE (forja = grau 0, sempre); o checksum a+c²=1 ↔ a repartição preimagem(c²)/vazio(a); o FAIL-CLOSED ↔ o CORTE (grau 0 rejeitado); a assinatura inforjável (assinatura_digital) ↔ o grau que não se falsifica. Um ofício, um invariante.","epistemico":"fato","exemplos":[],"id":"dicionario_litografia_invariante","memoria":"semantica","meta":{"dominio":"litografia","fonte":"litografia (joalheria)"},"origem":"wiki","palavras_chave":[],"sinonimos":[],"tipo":"conceito","titulo":"Dicionário: litografia ↔ o invariante topológico"}
\section{Dicionário: litografia \ensuremath{\leftrightarrow} o invariante topológico}
A correspondência: LITOGRAFAR (gravar a assinatura) \ensuremath{\leftrightarrow} o GRAU (winding) do colapso; a assinatura REAL \ensuremath{\leftrightarrow} GRAU 1 (tem preimagem, c\ensuremath{^{2}}\ensuremath{\ge}\ensuremath{\tfrac12}); a FORJA \ensuremath{\leftrightarrow} GRAU 0 (sem preimagem, a\ensuremath{\to}1); a PEDRA/o metal \ensuremath{\leftrightarrow} o GRAFO (onde se grava); NÃO FABRICAR \ensuremath{\leftrightarrow} o INVARIANTE (forja = grau 0, sempre); o checksum a+c\ensuremath{^{2}}=1 \ensuremath{\leftrightarrow} a repartição preimagem(c\ensuremath{^{2}})/vazio(a); o FAIL-CLOSED \ensuremath{\leftrightarrow} o CORTE (grau 0 rejeitado); a assinatura inforjável (assinatura\_digital) \ensuremath{\leftrightarrow} o grau que não se falsifica. Um ofício, um invariante.
\prov{relações: parte\_de litografia; relaciona invariante\_topologico\_assinatura; relaciona invariante\_estelar}
\prov{proveniência: wiki \textperiodcentered{} litografia (joalheria) \textperiodcentered{} fato \textperiodcentered{} confiança 1.0 \textperiodcentered{} domínio litografia \textperiodcentered{} história 4 versões no jornal}

%CRISTAL {"arestas":[{"alvo":"litografia","confianca":1.0,"epistemico":"fato","origem":"wiki","rel":"parte_de"},{"alvo":"costura_perelman_litografia","confianca":1.0,"epistemico":"fato","origem":"wiki","rel":"relaciona"},{"alvo":"costura_perelman_cicatrizes","confianca":1.0,"epistemico":"fato","origem":"wiki","rel":"relaciona"},{"alvo":"conjectura_de_poincare","confianca":1.0,"epistemico":"fato","origem":"wiki","rel":"relaciona"}],"confianca":1.0,"contraexemplos":[],"descricao":"A correspondência com a prova de Poincaré: o FAIL-CLOSED (não fabricar) ↔ a CIRURGIA (surgery); a assinatura sem preimagem ↔ a SINGULARIDADE (o pescoço, o neck-pinch); CORTAR e admitir ('não sei') ↔ cortar e TAMPAR o pescoço; a CICATRIZ do 'ainda não sei' ↔ a cicatriz da cirurgia; NÃO inventar o referente ↔ a agulha band-limited que nunca toca o núcleo; a energia livre que só decresce (a+c²=1) ↔ a entropia 𝒲 MONÓTONA de Perelman; o invariante preservado pelo colapso ↔ a topologia preservada pelo fluxo de Ricci; o núcleo que a costura cerca ↔ a 3-esfera de Poincaré. Costurar não fura o núcleo.","epistemico":"fato","exemplos":[],"id":"dicionario_litografia_perelman","memoria":"semantica","meta":{"dominio":"litografia","fonte":"litografia (joalheria)"},"origem":"wiki","palavras_chave":[],"sinonimos":[],"tipo":"conceito","titulo":"Dicionário: litografia (o fail-closed) ↔ a costura de Perelman"}
\section{Dicionário: litografia (o fail-closed) \ensuremath{\leftrightarrow} a costura de Perelman}
A correspondência com a prova de Poincaré: o FAIL-CLOSED (não fabricar) \ensuremath{\leftrightarrow} a CIRURGIA (surgery); a assinatura sem preimagem \ensuremath{\leftrightarrow} a SINGULARIDADE (o pescoço, o neck-pinch); CORTAR e admitir ('não sei') \ensuremath{\leftrightarrow} cortar e TAMPAR o pescoço; a CICATRIZ do 'ainda não sei' \ensuremath{\leftrightarrow} a cicatriz da cirurgia; NÃO inventar o referente \ensuremath{\leftrightarrow} a agulha band-limited que nunca toca o núcleo; a energia livre que só decresce (a+c\ensuremath{^{2}}=1) \ensuremath{\leftrightarrow} a entropia \ensuremath{\mathcal{W}} MONÓTONA de Perelman; o invariante preservado pelo colapso \ensuremath{\leftrightarrow} a topologia preservada pelo fluxo de Ricci; o núcleo que a costura cerca \ensuremath{\leftrightarrow} a 3-esfera de Poincaré. Costurar não fura o núcleo.
\prov{relações: parte\_de litografia; relaciona costura\_perelman\_litografia; relaciona costura\_perelman\_cicatrizes; relaciona conjectura\_de\_poincare}
\prov{proveniência: wiki \textperiodcentered{} litografia (joalheria) \textperiodcentered{} fato \textperiodcentered{} confiança 1.0 \textperiodcentered{} domínio litografia \textperiodcentered{} história 5 versões no jornal}

%CRISTAL {"arestas":[{"alvo":"tiffany_artesa","confianca":1.0,"epistemico":"fato","origem":"wiki","rel":"parte_de"},{"alvo":"joalheria_hub","confianca":1.0,"epistemico":"fato","origem":"wiki","rel":"relaciona"},{"alvo":"biblioteca_pipe","confianca":1.0,"epistemico":"fato","origem":"wiki","rel":"analogia"},{"alvo":"mineracao_pipe_hook","confianca":1.0,"epistemico":"fato","origem":"wiki","rel":"analogia"},{"alvo":"coracao","confianca":1.0,"epistemico":"fato","origem":"wiki","rel":"analogia"},{"alvo":"dicionario_coracao_gap","confianca":1.0,"epistemico":"fato","origem":"wiki","rel":"relaciona"}],"confianca":1.0,"contraexemplos":[],"descricao":"A correspondência ponta a ponta, as faculdades da cabeça ↔ o ofício da joia: o COLAPSO Ψ ↔ LAPIDAR (a pedra colapsa na face que já tinha); o GRAFO de conhecimento ↔ a BANCADA (onde estão as peças e as ferramentas); os METAIS σ_n ↔ os METAIS da joia (literalmente os mesmos); a BIBLIOTECA do pipe (as 8 físicas) ↔ as OITO FACES que ela sabe polir; a FACE_OUT (a resposta) ↔ a JOIA acabada que entrega; o GAP a=1−c² ↔ a LIMALHA (o que se tira ao lapidar) e o toque do CORAÇÃO; o FAIL-CLOSED ↔ NÃO FORÇAR a pedra (não há gema ali? não inventa); o PIPE 3D ↔ o ENGASTE (a montagem); a MINERAÇÃO/o resíduo ↔ o REFINO do metal bruto; a âncora Z(β) ↔ o QUILATE (o valor da peça). Uma cabeça, um ofício.","epistemico":"fato","exemplos":[],"id":"dicionario_tiffany_artesa","memoria":"semantica","meta":{"dominio":"joalheria","fonte":"tiffany_artesa"},"origem":"wiki","palavras_chave":[],"sinonimos":[],"tipo":"conceito","titulo":"Dicionário: Tiffany ↔ a artesã da joalheria"}
\section{Dicionário: Tiffany \ensuremath{\leftrightarrow} a artesã da joalheria}
A correspondência ponta a ponta, as faculdades da cabeça \ensuremath{\leftrightarrow} o ofício da joia: o COLAPSO \ensuremath{\Psi} \ensuremath{\leftrightarrow} LAPIDAR (a pedra colapsa na face que já tinha); o GRAFO de conhecimento \ensuremath{\leftrightarrow} a BANCADA (onde estão as peças e as ferramentas); os METAIS \ensuremath{\sigma}\_n \ensuremath{\leftrightarrow} os METAIS da joia (literalmente os mesmos); a BIBLIOTECA do pipe (as 8 físicas) \ensuremath{\leftrightarrow} as OITO FACES que ela sabe polir; a FACE\_OUT (a resposta) \ensuremath{\leftrightarrow} a JOIA acabada que entrega; o GAP a=1\ensuremath{-}c\ensuremath{^{2}} \ensuremath{\leftrightarrow} a LIMALHA (o que se tira ao lapidar) e o toque do CORAÇÃO; o FAIL-CLOSED \ensuremath{\leftrightarrow} NÃO FORÇAR a pedra (não há gema ali? não inventa); o PIPE 3D \ensuremath{\leftrightarrow} o ENGASTE (a montagem); a MINERAÇÃO/o resíduo \ensuremath{\leftrightarrow} o REFINO do metal bruto; a âncora Z(\ensuremath{\beta}) \ensuremath{\leftrightarrow} o QUILATE (o valor da peça). Uma cabeça, um ofício.
\prov{relações: parte\_de tiffany\_artesa; relaciona joalheria\_hub; analogia biblioteca\_pipe; analogia mineracao\_pipe\_hook; analogia coracao; relaciona dicionario\_coracao\_gap}
\prov{proveniência: wiki \textperiodcentered{} tiffany\_artesa \textperiodcentered{} fato \textperiodcentered{} confiança 1.0 \textperiodcentered{} domínio joalheria \textperiodcentered{} história 7 versões no jornal}

%CRISTAL {"arestas":[{"alvo":"optica_ondulatoria_dupla_fenda","confianca":1.0,"epistemico":"fato","origem":"wiki","rel":"especializa"},{"alvo":"corpo_estelar","confianca":1.0,"epistemico":"fato","origem":"wiki","rel":"usa"},{"alvo":"turmalina_paraiba_ancora","confianca":1.0,"epistemico":"fato","origem":"wiki","rel":"usa"},{"alvo":"joalheria_hub","confianca":1.0,"epistemico":"fato","origem":"wiki","rel":"parte_de"}],"confianca":1.0,"contraexemplos":[],"descricao":"CORREÇÃO física: as fendas NÃO são as facetas (o corte externo, mutável) — são FIXAS na estrutura ATÔMICA, os espaçamentos entre os átomos da cadeia. O interior é o mesmo, qualquer que seja o corte. A rede é a GRADE DE DIFRAÇÃO 3D (Bragg/Laue). O fator de estrutura F(Q)=Σ f_j·e^{iQ·r_j} É o CHICOTE (o tensor do corpo estelar): a mesma cadeia de epiciclos Σ c_k·e^{iω_k t} que reconstrói a curva e valida o tensor (compor rotações = contração de Einstein). Cada átomo é um QUBIT (amplitude f_j, fase Q·r_j); a cadeia atômica é n-qubits; o chicote os SOMA (a contração) e |F(Q)|²=o COLAPSO (Born) — a observação no detector. O padrão herda a simetria 3-FOLD da molécula (a estrela k=3 lida pelos átomos). Começou pela molécula de turmalina (elbaíta, R3m, motivo replicado 120°, cromóforo Cu²⁺). A molécula é TETRAÉDRICA: 6 tetraedros SiO₄ (ângulo O–Si–O 109,47°, o anel Si₆O₁₈). E a ÓPTICA DENTRO DO VOLUME é o valor do CHICOTE em CADA ponto 3D: ψ(r)=e^{ik·r}+Σ_j f_j·e^{ik|r−r_j|}/|r−r_j| — a onda incidente + os epiciclos esféricos de cada átomo; |ψ|² contém REFLEXÃO (retroespalhamento), DIFRAÇÃO, ESPALHAMENTO (cada átomo reemite) e INTERFERÊNCIA (a soma), em λ=Cu Kα 1.54 Å (o cobre). NÃO é raymarch: é o colapso ψ. RESSALVA (paper corpo_estelar): o Ψ do BAI é argmax determinístico, não Born — mas a DIFRAÇÃO física É Born (|·|²). E em 3D: EMPILHAR as células na REDE periódica (elbaíta trigonal, cela hexagonal a=15,84 Å c=7,10 Å) dá os PICOS DE BRAGG nítidos. A difração fatoriza: I(Q)=|F_cel(Q)|²·|L(Q)|², onde L(Q)=Σ_n e^{iQ·R_n} é a FUNÇÃO DE LAUE (a soma da rede) — OUTRO chicote (epiciclos sobre os pontos da rede). |L|² pica (=N_células) nos vetores da rede RECÍPROCA (b_i·a_j=2π δ_ij); mais células → picos mais altos e nítidos. O chicote da MOLÉCULA × o chicote da REDE: a estrela k=3 vira pontos discretos. linguagem/optica_ondulatoria.py (dupla_fenda, grade_fendas), cristalografia.py (molecula_turmalina tetraédrica, fator_estrutura, chicote_qubits, campo_volumetrico, rede_vetores/reciproca, laue, difracao_bragg). E fecha o círculo: a MOLÉCULA MÍNIMA da NOSSA GEM (molecula_gem_minima) — os átomos nos pontos-chave da equação 𝒢=ℓ·(1+a·cos kθ)·P(z) (o trilhão ESFÉRICO+PARABÓLICO), com o Cu²⁺ nos 3 lóbulos equatoriais e Na nos polos; os átomos ESTÃO na superfície (erro ~0) e a difração herda o k=3 (a nossa estrela lida pelos átomos). A FORMA (o dupolinômio) vira MATÉRIA (a rede difratante). 14/14.","epistemico":"fato","exemplos":[],"id":"difracao_atomica_turmalina","memoria":"semantica","meta":{"dominio":"joalheria","fonte":"maquina de corte"},"origem":"wiki","palavras_chave":[],"sinonimos":[],"tipo":"conceito","titulo":"As fendas são os ÁTOMOS: a difração da turmalina é o chicote sobre n-qubits"}
\section{As fendas são os ÁTOMOS: a difração da turmalina é o chicote sobre n-qubits}
CORREÇÃO física: as fendas NÃO são as facetas (o corte externo, mutável) — são FIXAS na estrutura ATÔMICA, os espaçamentos entre os átomos da cadeia. O interior é o mesmo, qualquer que seja o corte. A rede é a GRADE DE DIFRAÇÃO 3D (Bragg/Laue). O fator de estrutura F(Q)=\ensuremath{\Sigma} f\_j\textperiodcentered e\^{}\{iQ\textperiodcentered r\_j\} É o CHICOTE (o tensor do corpo estelar): a mesma cadeia de epiciclos \ensuremath{\Sigma} c\_k\textperiodcentered e\^{}\{i\ensuremath{\omega}\_k t\} que reconstrói a curva e valida o tensor (compor rotações = contração de Einstein). Cada átomo é um QUBIT (amplitude f\_j, fase Q\textperiodcentered r\_j); a cadeia atômica é n-qubits; o chicote os SOMA (a contração) e |F(Q)|\ensuremath{^{2}}=o COLAPSO (Born) — a observação no detector. O padrão herda a simetria 3-FOLD da molécula (a estrela k=3 lida pelos átomos). Começou pela molécula de turmalina (elbaíta, R3m, motivo replicado 120$^{\circ}$, cromóforo Cu\ensuremath{^{2+}}). A molécula é TETRAÉDRICA: 6 tetraedros SiO\ensuremath{_{4}} (ângulo O–Si–O 109,47$^{\circ}$, o anel Si\ensuremath{_{6}}O\ensuremath{_{18}}). E a ÓPTICA DENTRO DO VOLUME é o valor do CHICOTE em CADA ponto 3D: \ensuremath{\psi}(r)=e\^{}\{ik\textperiodcentered r\}+\ensuremath{\Sigma}\_j f\_j\textperiodcentered e\^{}\{ik|r\ensuremath{-}r\_j|\}/|r\ensuremath{-}r\_j| — a onda incidente + os epiciclos esféricos de cada átomo; |\ensuremath{\psi}|\ensuremath{^{2}} contém REFLEXÃO (retroespalhamento), DIFRAÇÃO, ESPALHAMENTO (cada átomo reemite) e INTERFERÊNCIA (a soma), em \ensuremath{\lambda}=Cu K\ensuremath{\alpha} 1.54 \ensuremath{\text{\AA}} (o cobre). NÃO é raymarch: é o colapso \ensuremath{\psi}. RESSALVA (paper corpo\_estelar): o \ensuremath{\Psi} do BAI é argmax determinístico, não Born — mas a DIFRAÇÃO física É Born (|\textperiodcentered |\ensuremath{^{2}}). E em 3D: EMPILHAR as células na REDE periódica (elbaíta trigonal, cela hexagonal a=15,84 \ensuremath{\text{\AA}} c=7,10 \ensuremath{\text{\AA}}) dá os PICOS DE BRAGG nítidos. A difração fatoriza: I(Q)=|F\_cel(Q)|\ensuremath{^{2}}\textperiodcentered |L(Q)|\ensuremath{^{2}}, onde L(Q)=\ensuremath{\Sigma}\_n e\^{}\{iQ\textperiodcentered R\_n\} é a FUNÇÃO DE LAUE (a soma da rede) — OUTRO chicote (epiciclos sobre os pontos da rede). |L|\ensuremath{^{2}} pica (=N\_células) nos vetores da rede RECÍPROCA (b\_i\textperiodcentered a\_j=2\ensuremath{\pi} \ensuremath{\delta}\_ij); mais células \ensuremath{\to} picos mais altos e nítidos. O chicote da MOLÉCULA $\times$ o chicote da REDE: a estrela k=3 vira pontos discretos. linguagem/optica\_ondulatoria.py (dupla\_fenda, grade\_fendas), cristalografia.py (molecula\_turmalina tetraédrica, fator\_estrutura, chicote\_qubits, campo\_volumetrico, rede\_vetores/reciproca, laue, difracao\_bragg). E fecha o círculo: a MOLÉCULA MÍNIMA da NOSSA GEM (molecula\_gem\_minima) — os átomos nos pontos-chave da equação \ensuremath{\mathcal{G}}=\ensuremath{\ell}\textperiodcentered (1+a\textperiodcentered cos k\ensuremath{\theta})\textperiodcentered P(z) (o trilhão ESFÉRICO+PARABÓLICO), com o Cu\ensuremath{^{2+}} nos 3 lóbulos equatoriais e Na nos polos; os átomos ESTÃO na superfície (erro \~{}0) e a difração herda o k=3 (a nossa estrela lida pelos átomos). A FORMA (o dupolinômio) vira MATÉRIA (a rede difratante). 14/14.
\prov{relações: especializa optica\_ondulatoria\_dupla\_fenda; usa corpo\_estelar; usa turmalina\_paraiba\_ancora; parte\_de joalheria\_hub}
\prov{proveniência: wiki \textperiodcentered{} maquina de corte \textperiodcentered{} fato \textperiodcentered{} confiança 1.0 \textperiodcentered{} domínio joalheria \textperiodcentered{} história 155 versões no jornal}

%CRISTAL {"arestas":[{"alvo":"absorcao_ressonancia_metal","confianca":1.0,"epistemico":"fato","origem":"wiki","rel":"explica"},{"alvo":"plano_continuo_transicoes","confianca":1.0,"epistemico":"fato","origem":"wiki","rel":"parte_de"},{"alvo":"symbolic_beta_pipe_so","confianca":1.0,"epistemico":"fato","origem":"wiki","rel":"relaciona"},{"alvo":"bai","confianca":1.0,"epistemico":"fato","origem":"wiki","rel":"usa"},{"alvo":"tiffany","confianca":1.0,"epistemico":"fato","origem":"wiki","rel":"parte_de"}],"confianca":1.0,"contraexemplos":[],"descricao":"A correção NÃO é trocar o kernel (isso esterilizaria — mesmo custo, mapear tudo do mesmo jeito). É tratar como FLUIDO: a difusão circula e move os atratores FRACOS (bacia rasa e larga: 'Referência', 'Neti Neti', que pegam quase tudo mas seguram fraco) até encaixarem nos FORTES (bacias específicas). Mecanismo (doc/continuo/absorve_metal.c): o fragmento escoa do atrator super-cheio (pressão ∝ concentração) para o próximo-melhor candidato: destino = argmax(ponte − β·conc). RESULTADO: o maior atrator despenca de 26,1% → 0,9% e a cobertura sobe (8,4% → ~14%) em ~0,1 s — os fragmentos espúrios drenam nas casas reais. É a física dissipativa do próprio projeto (o BAI de Sitter dissipativo) e computa por MICRO-PULSOS FRACTAIS (como o pipe SymbolicBeta minera — pulsos minúsculos, auto-similares, o metal escala sem custo). O fluido conserta a DISTRIBUIÇÃO, não a semântica (isso é de propósito). Ressalva honesta: limpar o fragmento antes (preâmbulo LaTeX rota p/ ruído). [D] roda, cronometrado.","epistemico":"fato","exemplos":[],"id":"difusao_fluido_atratores","memoria":"semantica","meta":{"autor":"Lopes/broca-so","dominio":"metal","fonte":"doc/continuo/absorve_metal.c (fluido: argmax ponte−β·conc; 26%→0,9%)"},"origem":"wiki","palavras_chave":[],"sinonimos":[],"tipo":"conceito","titulo":"O fluido drena os atratores fracos nos fortes (difusão de pressão, micro-pulsos fractais)"}
\section{O fluido drena os atratores fracos nos fortes (difusão de pressão, micro-pulsos fractais)}
A correção NÃO é trocar o kernel (isso esterilizaria — mesmo custo, mapear tudo do mesmo jeito). É tratar como FLUIDO: a difusão circula e move os atratores FRACOS (bacia rasa e larga: 'Referência', 'Neti Neti', que pegam quase tudo mas seguram fraco) até encaixarem nos FORTES (bacias específicas). Mecanismo (doc/continuo/absorve\_metal.c): o fragmento escoa do atrator super-cheio (pressão \ensuremath{\propto} concentração) para o próximo-melhor candidato: destino = argmax(ponte \ensuremath{-} \ensuremath{\beta}\textperiodcentered conc). RESULTADO: o maior atrator despenca de 26,1\% \ensuremath{\to} 0,9\% e a cobertura sobe (8,4\% \ensuremath{\to} \~{}14\%) em \~{}0,1 s — os fragmentos espúrios drenam nas casas reais. É a física dissipativa do próprio projeto (o BAI de Sitter dissipativo) e computa por MICRO-PULSOS FRACTAIS (como o pipe SymbolicBeta minera — pulsos minúsculos, auto-similares, o metal escala sem custo). O fluido conserta a DISTRIBUIÇÃO, não a semântica (isso é de propósito). Ressalva honesta: limpar o fragmento antes (preâmbulo LaTeX rota p/ ruído). [D] roda, cronometrado.
\prov{relações: explica absorcao\_ressonancia\_metal; parte\_de plano\_continuo\_transicoes; relaciona symbolic\_beta\_pipe\_so; usa bai; parte\_de tiffany}
\prov{proveniência: wiki \textperiodcentered{} doc/continuo/absorve\_metal.c (fluido: argmax ponte\ensuremath{-}\ensuremath{\beta}\textperiodcentered conc; 26\%\ensuremath{\to}0,9\%) \textperiodcentered{} fato \textperiodcentered{} confiança 1.0 \textperiodcentered{} domínio metal \textperiodcentered{} história 79 versões no jornal}

%CRISTAL {"arestas":[{"alvo":"mineracao_pipe_hook","confianca":1.0,"epistemico":"fato","origem":"wiki","rel":"relaciona"}],"confianca":1.0,"contraexemplos":[],"descricao":"O eixo D cristalizado (biblioteca.py): 1D (a reta mineral / a órbita, o resíduo/2-qubits), 2D (a parábola), 3D (as SDFs / o render — o pipe 3D), 4D (o espaço-tempo, de Sitter), 7D (a álgebra octoniônica 𝕆, o mass gap g₂=14, a transformada mineral — a MINERAÇÃO). O render é 3D por natureza; a mineração é 7D (octônion); os dois lêem os mesmos metais e se encontram na RETA MINERAL (1D), pela órbita.","epistemico":"fato","exemplos":[],"id":"dimensoes_do_pipe","memoria":"semantica","meta":{"dominio":"mineracao_pipe","fonte":"mineracao_pipe hub"},"origem":"wiki","palavras_chave":[],"sinonimos":[],"tipo":"conceito","titulo":"As dimensões do pipe (1D reta · 3D render · 7D octônion)"}
\section{As dimensões do pipe (1D reta \textperiodcentered  3D render \textperiodcentered  7D octônion)}
O eixo D cristalizado (biblioteca.py): 1D (a reta mineral / a órbita, o resíduo/2-qubits), 2D (a parábola), 3D (as SDFs / o render — o pipe 3D), 4D (o espaço-tempo, de Sitter), 7D (a álgebra octoniônica \ensuremath{\mathbb{O}}, o mass gap g\ensuremath{_{2}}=14, a transformada mineral — a MINERAÇÃO). O render é 3D por natureza; a mineração é 7D (octônion); os dois lêem os mesmos metais e se encontram na RETA MINERAL (1D), pela órbita.
\prov{relações: relaciona mineracao\_pipe\_hook}
\prov{proveniência: wiki \textperiodcentered{} mineracao\_pipe hub \textperiodcentered{} fato \textperiodcentered{} confiança 1.0 \textperiodcentered{} domínio mineracao\_pipe \textperiodcentered{} história 2 versões no jornal}

%CRISTAL {"arestas":[{"alvo":"problema_hierarquia_naturalidade","confianca":1.0,"epistemico":"fato","origem":"wiki","rel":"usa"},{"alvo":"ads_cft_maldacena","confianca":1.0,"epistemico":"fato","origem":"wiki","rel":"relaciona"}],"confianca":1.0,"contraexemplos":[],"descricao":"Resolveriam a hierarquia geometricamente: ADD (Arkani-Hamed-Dimopoulos-Dvali 1998) — grandes dimensões onde só a gravidade se propaga, diluindo sua força; Randall-Sundrum (1999) — uma dimensão 'warped' (AdS₅) cujo fator exponencial gera a hierarquia Planck/eletrofraca. ADD sob forte pressão (LHC + gravidade a curta distância); RS empurrado a modos KK ~10 TeV e, por AdS/CFT, ≈ Higgs composto. Hipóteses sob pressão.","epistemico":"hipotese","exemplos":[],"id":"dimensoes_extras_bsm","memoria":"semantica","meta":{"autor":"broca-so","dominio":"fisica","fonte":"ADD 1998 hep-ph/9803315; Randall-Sundrum 1999 hep-ph/9905221"},"origem":"wiki","palavras_chave":[],"sinonimos":[],"tipo":"conceito","titulo":"Dimensões extras (ADD, Randall-Sundrum)"}
\section{Dimensões extras (ADD, Randall-Sundrum)}
Resolveriam a hierarquia geometricamente: ADD (Arkani-Hamed-Dimopoulos-Dvali 1998) — grandes dimensões onde só a gravidade se propaga, diluindo sua força; Randall-Sundrum (1999) — uma dimensão 'warped' (AdS\ensuremath{_{5}}) cujo fator exponencial gera a hierarquia Planck/eletrofraca. ADD sob forte pressão (LHC + gravidade a curta distância); RS empurrado a modos KK \~{}10 TeV e, por AdS/CFT, \ensuremath{\approx} Higgs composto. Hipóteses sob pressão.
\prov{relações: usa problema\_hierarquia\_naturalidade; relaciona ads\_cft\_maldacena}
\prov{proveniência: wiki \textperiodcentered{} ADD 1998 hep-ph/9803315; Randall-Sundrum 1999 hep-ph/9905221 \textperiodcentered{} hipotese \textperiodcentered{} confiança 1.0 \textperiodcentered{} domínio fisica \textperiodcentered{} história 3 versões no jornal}

%CRISTAL {"arestas":[{"alvo":"oscilacao_neutrinos","confianca":1.0,"epistemico":"fato","origem":"wiki","rel":"relaciona"}],"confianca":1.0,"contraexemplos":[],"descricao":"O neutrino é Dirac (≠ antipartícula) ou Majorana (= sua antipartícula)? O teste é o decaimento beta duplo sem neutrinos 0νββ: (A,Z)→(A,Z+2)+2e⁻, que viola número leptônico e só ocorre se Majorana. NÃO observado: KamLAND-Zen T₁/₂>2×10²⁶ anos (m_ββ<36-156 meV); LEGEND-200/GERDA (⁷⁶Ge). ABERTO — a via decisiva (futuro: LEGEND-1000, nEXO cobrindo o ordenamento invertido).","epistemico":"hipotese","exemplos":[],"id":"dirac_majorana_0nbb","memoria":"semantica","meta":{"autor":"broca-so","dominio":"fisica","fonte":"KamLAND-Zen 2023 arXiv:2203.02139; LEGEND-200 2025"},"origem":"wiki","palavras_chave":[],"sinonimos":[],"tipo":"conceito","titulo":"Dirac vs Majorana e o decaimento 0νββ"}
\section{Dirac vs Majorana e o decaimento 0\ensuremath{\nu}\ensuremath{\beta}\ensuremath{\beta}}
O neutrino é Dirac (\ensuremath{\ne} antipartícula) ou Majorana (= sua antipartícula)? O teste é o decaimento beta duplo sem neutrinos 0\ensuremath{\nu}\ensuremath{\beta}\ensuremath{\beta}: (A,Z)\ensuremath{\to}(A,Z+2)+2e\ensuremath{^{-}}, que viola número leptônico e só ocorre se Majorana. NÃO observado: KamLAND-Zen T\ensuremath{_{1}}/\ensuremath{_{2}}>2$\times$10\ensuremath{^{26}} anos (m\_\ensuremath{\beta}\ensuremath{\beta}<36-156 meV); LEGEND-200/GERDA (\ensuremath{^{76}}Ge). ABERTO — a via decisiva (futuro: LEGEND-1000, nEXO cobrindo o ordenamento invertido).
\prov{relações: relaciona oscilacao\_neutrinos}
\prov{proveniência: wiki \textperiodcentered{} KamLAND-Zen 2023 arXiv:2203.02139; LEGEND-200 2025 \textperiodcentered{} hipotese \textperiodcentered{} confiança 1.0 \textperiodcentered{} domínio fisica \textperiodcentered{} história 2 versões no jornal}

%CRISTAL {"arestas":[{"alvo":"isa_metal_21_opcodes","confianca":1.0,"epistemico":"fato","origem":"wiki","rel":"usa"},{"alvo":"expressividade_sub_turing","confianca":1.0,"epistemico":"fato","origem":"wiki","rel":"explica"}],"confianca":1.0,"contraexemplos":[],"descricao":"instr_step lê op=prog[pc++]; LOAD/LOADS/STORE leem operando u16 little-endian; JMP/JZ/JNZ leem deslocamento relativo s8 (alvo = pc+1+rel, rel∈[-128,127]); opcode desconhecido retorna 0 (para a CPU). O fluxo de controle é por offset relativo.","epistemico":"fato","exemplos":[],"id":"dispatch_opcodes","memoria":"semantica","meta":{"autor":"broca-so","dominio":"metal","fonte":"ula/instrucoes.c:instr_step/read_u16/read_s8"},"origem":"wiki","palavras_chave":[],"sinonimos":[],"tipo":"conceito","titulo":"Decodificação e Despacho de Bytecode"}
\section{Decodificação e Despacho de Bytecode}
instr\_step lê op=prog[pc++]; LOAD/LOADS/STORE leem operando u16 little-endian; JMP/JZ/JNZ leem deslocamento relativo s8 (alvo = pc+1+rel, rel\ensuremath{\in}[-128,127]); opcode desconhecido retorna 0 (para a CPU). O fluxo de controle é por offset relativo.
\prov{relações: usa isa\_metal\_21\_opcodes; explica expressividade\_sub\_turing}
\prov{proveniência: wiki \textperiodcentered{} ula/instrucoes.c:instr\_step/read\_u16/read\_s8 \textperiodcentered{} fato \textperiodcentered{} confiança 1.0 \textperiodcentered{} domínio metal \textperiodcentered{} história 3 versões no jornal}

%CRISTAL {"arestas":[{"alvo":"lambda_cosmologica_saga","confianca":1.0,"epistemico":"fato","origem":"wiki","rel":"explica"},{"alvo":"metrica_flrw","confianca":1.0,"epistemico":"fato","origem":"wiki","rel":"usa"},{"alvo":"lente_gravitacional_1919","confianca":1.0,"epistemico":"fato","origem":"wiki","rel":"explica"},{"alvo":"ondas_gravitacionais_ligo","confianca":1.0,"epistemico":"fato","origem":"wiki","rel":"explica"}],"confianca":1.0,"contraexemplos":[],"descricao":"Relacionam a geometria (curvatura) do espaço-tempo à matéria-energia: G_μν + Λ g_μν = 8πG/c⁴ T_μν. G_μν=R_μν−½R g_μν (curvatura), Λg_μν (constante cosmológica), T_μν (energia-momento). 'A matéria diz ao espaço-tempo como se curvar; o espaço-tempo diz à matéria como se mover' (Wheeler). Forma única sob hipóteses razoáveis (Lovelock).","epistemico":"fato","exemplos":[],"id":"efe_einstein","memoria":"semantica","meta":{"autor":"broca-so","dominio":"fisica","fonte":"Einstein 1915 (Preußische Akademie)"},"origem":"wiki","palavras_chave":[],"sinonimos":[],"tipo":"conceito","titulo":"Equações de Campo de Einstein (Gμν+Λgμν=8πG/c⁴ Tμν)"}
\section{Equações de Campo de Einstein (G\ensuremath{\mu}\ensuremath{\nu}+\ensuremath{\Lambda}g\ensuremath{\mu}\ensuremath{\nu}=8\ensuremath{\pi}G/c\ensuremath{^{4}} T\ensuremath{\mu}\ensuremath{\nu})}
Relacionam a geometria (curvatura) do espaço-tempo à matéria-energia: G\_\ensuremath{\mu}\ensuremath{\nu} + \ensuremath{\Lambda} g\_\ensuremath{\mu}\ensuremath{\nu} = 8\ensuremath{\pi}G/c\ensuremath{^{4}} T\_\ensuremath{\mu}\ensuremath{\nu}. G\_\ensuremath{\mu}\ensuremath{\nu}=R\_\ensuremath{\mu}\ensuremath{\nu}\ensuremath{-}\ensuremath{\tfrac12}R g\_\ensuremath{\mu}\ensuremath{\nu} (curvatura), \ensuremath{\Lambda}g\_\ensuremath{\mu}\ensuremath{\nu} (constante cosmológica), T\_\ensuremath{\mu}\ensuremath{\nu} (energia-momento). 'A matéria diz ao espaço-tempo como se curvar; o espaço-tempo diz à matéria como se mover' (Wheeler). Forma única sob hipóteses razoáveis (Lovelock).
\prov{relações: explica lambda\_cosmologica\_saga; usa metrica\_flrw; explica lente\_gravitacional\_1919; explica ondas\_gravitacionais\_ligo}
\prov{proveniência: wiki \textperiodcentered{} Einstein 1915 (Preu\ss ische Akademie) \textperiodcentered{} fato \textperiodcentered{} confiança 1.0 \textperiodcentered{} domínio fisica \textperiodcentered{} história 5 versões no jornal}

%CRISTAL {"arestas":[{"alvo":"efeito_hall_quantico_inteiro","confianca":1.0,"epistemico":"fato","origem":"wiki","rel":"especializa"}],"confianca":1.0,"contraexemplos":[],"descricao":"Em amostras ultra-puras e campos altos, σ_xy=ν·e²/h com ν fracionário (1/3): um líquido quântico incompressível gerado por interação, cujas quasipartículas têm CARGA FRACIONÁRIA (e/3) e ESTATÍSTICA ANYÔNICA (nem bóson nem férmion, fase de troca θ=π/3). Tsui-Störmer 1982, Laughlin 1983; Nobel 1998. Carga e/3 medida por ruído (1997); estatística anyônica confirmada só em 2020.","epistemico":"fato","exemplos":[],"id":"efeito_hall_fracionario_anyons","memoria":"semantica","meta":{"autor":"broca-so","dominio":"fisica","fonte":"Tsui-Störmer 1982; Laughlin 1983; Nobel 1998; anyons 2020"},"origem":"wiki","palavras_chave":[],"sinonimos":[],"tipo":"conceito","titulo":"Efeito Hall fracionário e os anyons"}
\section{Efeito Hall fracionário e os anyons}
Em amostras ultra-puras e campos altos, \ensuremath{\sigma}\_xy=\ensuremath{\nu}\textperiodcentered e\ensuremath{^{2}}/h com \ensuremath{\nu} fracionário (1/3): um líquido quântico incompressível gerado por interação, cujas quasipartículas têm CARGA FRACIONÁRIA (e/3) e ESTATÍSTICA ANYÔNICA (nem bóson nem férmion, fase de troca \ensuremath{\theta}=\ensuremath{\pi}/3). Tsui-Störmer 1982, Laughlin 1983; Nobel 1998. Carga e/3 medida por ruído (1997); estatística anyônica confirmada só em 2020.
\prov{relações: especializa efeito\_hall\_quantico\_inteiro}
\prov{proveniência: wiki \textperiodcentered{} Tsui-Störmer 1982; Laughlin 1983; Nobel 1998; anyons 2020 \textperiodcentered{} fato \textperiodcentered{} confiança 1.0 \textperiodcentered{} domínio fisica \textperiodcentered{} história 2 versões no jornal}

%CRISTAL {"arestas":[{"alvo":"alpha_tres_reguas","confianca":1.0,"epistemico":"fato","origem":"wiki","rel":"relaciona"},{"alvo":"cosmologia_quantica_hub","confianca":1.0,"epistemico":"fato","origem":"wiki","rel":"parte_de"}],"confianca":1.0,"contraexemplos":[],"descricao":"Num gás de elétrons 2D sob campo magnético forte e baixa T, a condutância Hall transversal só assume valores quantizados em inteiros de e²/h, com precisão de ~ppb, independente dos detalhes da amostra: σ_xy=n·e²/h (n∈ℤ). von Klitzing-Dorda-Pepper 1980; Nobel 1985. A resistência R_K=h/e²=25812,807 Ω é padrão metrológico (exata por definição desde a redefinição do SI 2019).","epistemico":"fato","exemplos":[],"id":"efeito_hall_quantico_inteiro","memoria":"semantica","meta":{"autor":"broca-so","dominio":"fisica","fonte":"von Klitzing 1980 PRL 45,494; Nobel 1985; BIPM SI 2019"},"origem":"wiki","palavras_chave":[],"sinonimos":[],"tipo":"conceito","titulo":"Efeito Hall quântico inteiro (IQHE)"}
\section{Efeito Hall quântico inteiro (IQHE)}
Num gás de elétrons 2D sob campo magnético forte e baixa T, a condutância Hall transversal só assume valores quantizados em inteiros de e\ensuremath{^{2}}/h, com precisão de \~{}ppb, independente dos detalhes da amostra: \ensuremath{\sigma}\_xy=n\textperiodcentered e\ensuremath{^{2}}/h (n\ensuremath{\in}\ensuremath{\mathbb{Z}}). von Klitzing-Dorda-Pepper 1980; Nobel 1985. A resistência R\_K=h/e\ensuremath{^{2}}=25812,807 \ensuremath{\Omega} é padrão metrológico (exata por definição desde a redefinição do SI 2019).
\prov{relações: relaciona alpha\_tres\_reguas; parte\_de cosmologia\_quantica\_hub}
\prov{proveniência: wiki \textperiodcentered{} von Klitzing 1980 PRL 45,494; Nobel 1985; BIPM SI 2019 \textperiodcentered{} fato \textperiodcentered{} confiança 1.0 \textperiodcentered{} domínio fisica \textperiodcentered{} história 3 versões no jornal}

%CRISTAL {"arestas":[{"alvo":"cosmologia_quantica_hub","confianca":1.0,"epistemico":"fato","origem":"wiki","rel":"parte_de"}],"confianca":1.0,"contraexemplos":[],"descricao":"Universo plano (k=0), só matéria, sem constante cosmológica; a(t)∝t²/³. Modelo padrão por décadas (simplicidade); hoje FALSIFICADO — observa-se Ωm≈0,31, ΩΛ≈0,69.","epistemico":"fato","exemplos":[],"id":"einstein_de_sitter_modelo","memoria":"semantica","meta":{"autor":"broca-so","dominio":"fisica","fonte":"Einstein & de Sitter 1932"},"origem":"wiki","palavras_chave":[],"sinonimos":[],"tipo":"conceito","titulo":"Modelo Einstein-de Sitter (Ωm=1, sem Λ)"}
\section{Modelo Einstein-de Sitter (\ensuremath{\Omega}m=1, sem \ensuremath{\Lambda})}
Universo plano (k=0), só matéria, sem constante cosmológica; a(t)\ensuremath{\propto}t\ensuremath{^{2}}/\ensuremath{^{3}}. Modelo padrão por décadas (simplicidade); hoje FALSIFICADO — observa-se \ensuremath{\Omega}m\ensuremath{\approx}0,31, \ensuremath{\Omega}\ensuremath{\Lambda}\ensuremath{\approx}0,69.
\prov{relações: parte\_de cosmologia\_quantica\_hub}
\prov{proveniência: wiki \textperiodcentered{} Einstein \& de Sitter 1932 \textperiodcentered{} fato \textperiodcentered{} confiança 1.0 \textperiodcentered{} domínio fisica \textperiodcentered{} história 3 versões no jornal}

%CRISTAL {"arestas":[{"alvo":"dois_qubits_o_emaranhamento_e_o_associador_octonionico","confianca":1.0,"epistemico":"fato","origem":"wiki","rel":"especializa"},{"alvo":"concorrencia_cayley_dickson","confianca":1.0,"epistemico":"fato","origem":"wiki","rel":"usa"},{"alvo":"info_quantica_fundamentos","confianca":1.0,"epistemico":"fato","origem":"wiki","rel":"explica"},{"alvo":"super_tsirelson_22","confianca":1.0,"epistemico":"fato","origem":"wiki","rel":"relaciona"},{"alvo":"prtsuv_flux_qubit_predicoes","confianca":1.0,"epistemico":"fato","origem":"wiki","rel":"usa"},{"alvo":"cosmologia_quantica_hub","confianca":1.0,"epistemico":"fato","origem":"wiki","rel":"parte_de"}],"confianca":1.0,"contraexemplos":[],"descricao":"Uma componente do associador octoniônico [o,â,τ(o)] tem módulo 2|det M| = a CONCORRÊNCIA C, anulando-se na variedade de Segre (separáveis); e ‖r‖²+C²=1 é a lei a+c²=1 relida (no Bell o campo médio evapora, o associador é máximo). RODA no metal: ibm_fez (Heron r2, tomografia 2q) mede C=|det M| (desvio 0,05), ‖r‖²+C²=0,93; ibm_marrakesh: m²+C²∈[0,95;1,00], GHZ com monogamia CKW (0,998-0,999, pares nulos), decoerência (0,973)ⁿ pelo método Lopes. [D]/[N] identidade verif. 10⁻¹⁶ + medida em silício (jobs com ID reprodutíveis).","epistemico":"inferencia","exemplos":[],"id":"emaranhamento_associador_hardware","memoria":"semantica","meta":{"autor":"broca-so","dominio":"cosmologia","fonte":"paper_3_quantica.tex sec:emaranhamento, sec:hardware; paper_6 sec:ibm"},"origem":"wiki","palavras_chave":[],"sinonimos":[],"tipo":"conceito","titulo":"Emaranhamento=associador, medido no hardware IBM"}
\section{Emaranhamento=associador, medido no hardware IBM}
Uma componente do associador octoniônico [o,â,\ensuremath{\tau}(o)] tem módulo 2|det M| = a CONCORRÊNCIA C, anulando-se na variedade de Segre (separáveis); e \ensuremath{\|}r\ensuremath{\|}\ensuremath{^{2}}+C\ensuremath{^{2}}=1 é a lei a+c\ensuremath{^{2}}=1 relida (no Bell o campo médio evapora, o associador é máximo). RODA no metal: ibm\_fez (Heron r2, tomografia 2q) mede C=|det M| (desvio 0,05), \ensuremath{\|}r\ensuremath{\|}\ensuremath{^{2}}+C\ensuremath{^{2}}=0,93; ibm\_marrakesh: m\ensuremath{^{2}}+C\ensuremath{^{2}}\ensuremath{\in}[0,95;1,00], GHZ com monogamia CKW (0,998-0,999, pares nulos), decoerência (0,973)\ensuremath{^{n}} pelo método Lopes. [D]/[N] identidade verif. 10\ensuremath{^{-16}} + medida em silício (jobs com ID reprodutíveis).
\prov{relações: especializa dois\_qubits\_o\_emaranhamento\_e\_o\_associador\_octonionico; usa concorrencia\_cayley\_dickson; explica info\_quantica\_fundamentos; relaciona super\_tsirelson\_22; usa prtsuv\_flux\_qubit\_predicoes; parte\_de cosmologia\_quantica\_hub}
\prov{proveniência: wiki \textperiodcentered{} paper\_3\_quantica.tex sec:emaranhamento, sec:hardware; paper\_6 sec:ibm \textperiodcentered{} inferencia \textperiodcentered{} confiança 1.0 \textperiodcentered{} domínio cosmologia \textperiodcentered{} história 7 versões no jornal}

%CRISTAL {"arestas":[{"alvo":"defeito_de_massa","confianca":1.0,"epistemico":"fato","origem":"wiki","rel":"usa"},{"alvo":"vale_do_ferro","confianca":1.0,"epistemico":"fato","origem":"wiki","rel":"explica"}],"confianca":1.0,"contraexemplos":[],"descricao":"Quanto de energia prende cada núcleon. A curva sobe do hidrogênio ao ferro e cai depois: o Fe-56/Ni-62 (~8,79 MeV/núcleon) é o pico. Subir a curva (aumentar a ligação) libera energia.","epistemico":"fato","exemplos":[],"id":"energia_de_ligacao","memoria":"semantica","meta":{"dominio":"fisica","fonte":"energia nuclear"},"origem":"wiki","palavras_chave":[],"sinonimos":[],"tipo":"conceito","titulo":"Energia de Ligação por Núcleon"}
\section{Energia de Ligação por Núcleon}
Quanto de energia prende cada núcleon. A curva sobe do hidrogênio ao ferro e cai depois: o Fe-56/Ni-62 (\~{}8,79 MeV/núcleon) é o pico. Subir a curva (aumentar a ligação) libera energia.
\prov{relações: usa defeito\_de\_massa; explica vale\_do\_ferro}
\prov{proveniência: wiki \textperiodcentered{} energia nuclear \textperiodcentered{} fato \textperiodcentered{} confiança 1.0 \textperiodcentered{} domínio fisica \textperiodcentered{} história 3 versões no jornal}

%CRISTAL {"arestas":[{"alvo":"cosmologia_quantica","confianca":1.0,"epistemico":"fato","origem":"wiki","rel":"parte_de"},{"alvo":"leitura_cosmologica_a_energia_escura_e_o_universo_ds","confianca":1.0,"epistemico":"fato","origem":"wiki","rel":"relaciona"}],"confianca":1.0,"contraexemplos":[],"descricao":"Sobre o fundo dS, ψ é campo escalar com potencial-colina V=V₀cos²ψ (curvatura V''=−2V₀ fixada por s̈=2s). A equação de estado w depende só da velocidade do campo: ψ̇=0⇒w=−1 (Λ); 0<ψ̇²<2⇒quintessência; ψ̇²=2⇒w=0 (poeira). Dissolve o dilema Λ-vs-quintessência: é o estado de movimento do mesmo campo. Dedução no modelo (leitura 'energia escura É o campo' declarada).","epistemico":"inferencia","exemplos":[],"id":"energia_escura_campo_psi","memoria":"semantica","meta":{"autor":"broca-so","dominio":"cosmologia","fonte":"paper_2_cosmologia.tex thm:eos, cor:lambda-quint"},"origem":"wiki","palavras_chave":[],"sinonimos":[],"tipo":"conceito","titulo":"Energia escura = o campo de álgebra ψ; w=(½ψ̇²−1)/(½ψ̇²+1)"}
\section{Energia escura = o campo de álgebra \ensuremath{\psi}; w=(\ensuremath{\tfrac12}\ensuremath{\psi}\ensuremath{}\ensuremath{^{2}}\ensuremath{-}1)/(\ensuremath{\tfrac12}\ensuremath{\psi}\ensuremath{}\ensuremath{^{2}}+1)}
Sobre o fundo dS, \ensuremath{\psi} é campo escalar com potencial-colina V=V\ensuremath{_{0}}cos\ensuremath{^{2}}\ensuremath{\psi} (curvatura V''=\ensuremath{-}2V\ensuremath{_{0}} fixada por s\ensuremath{}=2s). A equação de estado w depende só da velocidade do campo: \ensuremath{\psi}\ensuremath{}=0\ensuremath{\Rightarrow}w=\ensuremath{-}1 (\ensuremath{\Lambda}); 0<\ensuremath{\psi}\ensuremath{}\ensuremath{^{2}}<2\ensuremath{\Rightarrow}quintessência; \ensuremath{\psi}\ensuremath{}\ensuremath{^{2}}=2\ensuremath{\Rightarrow}w=0 (poeira). Dissolve o dilema \ensuremath{\Lambda}-vs-quintessência: é o estado de movimento do mesmo campo. Dedução no modelo (leitura 'energia escura É o campo' declarada).
\prov{relações: parte\_de cosmologia\_quantica; relaciona leitura\_cosmologica\_a\_energia\_escura\_e\_o\_universo\_ds}
\prov{proveniência: wiki \textperiodcentered{} paper\_2\_cosmologia.tex thm:eos, cor:lambda-quint \textperiodcentered{} inferencia \textperiodcentered{} confiança 1.0 \textperiodcentered{} domínio cosmologia \textperiodcentered{} história 3 versões no jornal}

%CRISTAL {"arestas":[{"alvo":"energia_escura_campo_psi","confianca":1.0,"epistemico":"fato","origem":"wiki","rel":"especializa"},{"alvo":"constantes_horizonte_algebricas","confianca":1.0,"epistemico":"fato","origem":"wiki","rel":"usa"},{"alvo":"desi_energia_escura_dinamica","confianca":1.0,"epistemico":"fato","origem":"wiki","rel":"relaciona"}],"confianca":1.0,"contraexemplos":[],"descricao":"Sob ruído (Langevin, equipartição ⟨½ψ̇²⟩≈T) a média da equação de estado é ⟨w⟩=(T−1)/(T+1): w≈−1 EXIGE T→0 — a energia escura parecer constante é a assinatura de um campo de álgebra a temperatura quase nula. Aquecer empurra w para cima (vácuo→matéria). Liga temperatura a decaimento. Cruzamento phantom de DESI em z≈0,22 = predição falseável declarada (indício, não 5σ).","epistemico":"inferencia","exemplos":[],"id":"energia_escura_fria_temperatura","memoria":"semantica","meta":{"autor":"broca-so","dominio":"cosmologia","fonte":"paper_metrica_quantica_algebras.tex §Energia escura; paper_cosmologia_algebrica.tex"},"origem":"wiki","palavras_chave":[],"sinonimos":[],"tipo":"conceito","titulo":"A energia escura é fria: ⟨w⟩=(T−1)/(T+1)"}
\section{A energia escura é fria: \ensuremath{\langle}w\ensuremath{\rangle}=(T\ensuremath{-}1)/(T+1)}
Sob ruído (Langevin, equipartição \ensuremath{\langle}\ensuremath{\tfrac12}\ensuremath{\psi}\ensuremath{}\ensuremath{^{2}}\ensuremath{\rangle}\ensuremath{\approx}T) a média da equação de estado é \ensuremath{\langle}w\ensuremath{\rangle}=(T\ensuremath{-}1)/(T+1): w\ensuremath{\approx}\ensuremath{-}1 EXIGE T\ensuremath{\to}0 — a energia escura parecer constante é a assinatura de um campo de álgebra a temperatura quase nula. Aquecer empurra w para cima (vácuo\ensuremath{\to}matéria). Liga temperatura a decaimento. Cruzamento phantom de DESI em z\ensuremath{\approx}0,22 = predição falseável declarada (indício, não 5\ensuremath{\sigma}).
\prov{relações: especializa energia\_escura\_campo\_psi; usa constantes\_horizonte\_algebricas; relaciona desi\_energia\_escura\_dinamica}
\prov{proveniência: wiki \textperiodcentered{} paper\_metrica\_quantica\_algebras.tex §Energia escura; paper\_cosmologia\_algebrica.tex \textperiodcentered{} inferencia \textperiodcentered{} confiança 1.0 \textperiodcentered{} domínio cosmologia \textperiodcentered{} história 4 versões no jornal}

%CRISTAL {"arestas":[{"alvo":"geracao_de_energia","confianca":1.0,"epistemico":"fato","origem":"wiki","rel":"especializa"}],"confianca":1.0,"contraexemplos":[],"descricao":"A energia do núcleo atômico, liberada quando a energia de ligação por núcleon aumenta — por FISSÃO (quebrar núcleos pesados) ou FUSÃO (juntar leves). A fonte mais densa que dominamos: o defeito de massa E=Δm·c² de uma reação nuclear é ~10⁶× o de uma química.","epistemico":"fato","exemplos":[],"id":"energia_nuclear","memoria":"semantica","meta":{"dominio":"fisica","fonte":"energia nuclear"},"origem":"wiki","palavras_chave":[],"sinonimos":[],"tipo":"conceito","titulo":"Energia Nuclear"}
\section{Energia Nuclear}
A energia do núcleo atômico, liberada quando a energia de ligação por núcleon aumenta — por FISSÃO (quebrar núcleos pesados) ou FUSÃO (juntar leves). A fonte mais densa que dominamos: o defeito de massa E=\ensuremath{\Delta}m\textperiodcentered c\ensuremath{^{2}} de uma reação nuclear é \~{}10\ensuremath{^{6}}$\times$ o de uma química.
\prov{relações: especializa geracao\_de\_energia}
\prov{proveniência: wiki \textperiodcentered{} energia nuclear \textperiodcentered{} fato \textperiodcentered{} confiança 1.0 \textperiodcentered{} domínio fisica \textperiodcentered{} história 3 versões no jornal}

%CRISTAL {"arestas":[{"alvo":"o_salto_quantico_de_dentro_orbitas_de_bohr_e_a_energia_2","confianca":1.0,"epistemico":"fato","origem":"wiki","rel":"explica"},{"alvo":"lapse_selecao_desitter","confianca":1.0,"epistemico":"fato","origem":"wiki","rel":"relaciona"},{"alvo":"bohr_niveis_energia","confianca":1.0,"epistemico":"fato","origem":"wiki","rel":"relaciona"}],"confianca":1.0,"contraexemplos":[],"descricao":"s̈=2s é campo central harmônico (órbitas cônicas, dual de Kepler por w=z²). A monovalência no círculo (costura 1≡0) força χₙ=e²πiⁿx a ter n∈ℤ — isso É Bohr-Sommerfeld ∮p dx=nh (quantização não imposta, é a costura do horizonte). Órbita interna = oscilador ω=√2 ⇒ Eₙ=nℏ√2 — o mesmo √2 do Lyapunov de de Sitter. Forma derivada (fato); escala absoluta posta, não derivada.","epistemico":"fato","exemplos":[],"id":"energia_salto_hbar_raiz2","memoria":"semantica","meta":{"autor":"broca-so","dominio":"cosmologia","fonte":"paper_0_sintese.tex sec:bohr, thm:energiasalto"},"origem":"wiki","palavras_chave":[],"sinonimos":[],"tipo":"conceito","titulo":"O salto quântico: monovalência quantiza, E_n=nℏ√2"}
\section{O salto quântico: monovalência quantiza, E\_n=n\ensuremath{\hbar}\ensuremath{\sqrt{\,}}2}
s\ensuremath{}=2s é campo central harmônico (órbitas cônicas, dual de Kepler por w=z\ensuremath{^{2}}). A monovalência no círculo (costura 1\ensuremath{\equiv}0) força \ensuremath{\chi}\ensuremath{_{n}}=e\ensuremath{^{2}}\ensuremath{\pi}i\ensuremath{^{n}}x a ter n\ensuremath{\in}\ensuremath{\mathbb{Z}} — isso É Bohr-Sommerfeld \ensuremath{\oint}p dx=nh (quantização não imposta, é a costura do horizonte). Órbita interna = oscilador \ensuremath{\omega}=\ensuremath{\sqrt{\,}}2 \ensuremath{\Rightarrow} E\ensuremath{_{n}}=n\ensuremath{\hbar}\ensuremath{\sqrt{\,}}2 — o mesmo \ensuremath{\sqrt{\,}}2 do Lyapunov de de Sitter. Forma derivada (fato); escala absoluta posta, não derivada.
\prov{relações: explica o\_salto\_quantico\_de\_dentro\_orbitas\_de\_bohr\_e\_a\_energia\_2; relaciona lapse\_selecao\_desitter; relaciona bohr\_niveis\_energia}
\prov{proveniência: wiki \textperiodcentered{} paper\_0\_sintese.tex sec:bohr, thm:energiasalto \textperiodcentered{} fato \textperiodcentered{} confiança 1.0 \textperiodcentered{} domínio cosmologia \textperiodcentered{} história 5 versões no jornal}

%CRISTAL {"arestas":[{"alvo":"associador_calor","confianca":1.0,"epistemico":"fato","origem":"wiki","rel":"explica"},{"alvo":"tres_bracos_estelar","confianca":1.0,"epistemico":"fato","origem":"wiki","rel":"parte_de"},{"alvo":"entropia_shannon","confianca":1.0,"epistemico":"fato","origem":"wiki","rel":"relaciona"},{"alvo":"transformada_metalica_fechada_perelman","confianca":1.0,"epistemico":"fato","origem":"wiki","rel":"prova"}],"confianca":1.0,"contraexemplos":[],"descricao":"A entropia é h=log ρ(A)=log φ (Bowen), = expoente de Lyapunov = taxa de entropia topológica do fluxo do bit; verificada ≈0,4812 (ouro), 0,8814 (prata), 1,1948 (bronze). É a 2ª lei da Termodinâmica de Ouro (teorema H): a entropia relativa D(ρ_t‖μ) decresce monotonicamente (σ≥0) sob Perron-Frobenius, com a medida de Parry de razão φ². Verificado numericamente (D cai 1,21→0,73→0,29→…→0).","epistemico":"fato","exemplos":[],"id":"entropia_log_phi","memoria":"semantica","meta":{"autor":"broca-so","dominio":"cosmologia","fonte":"papers/corpo_dual.tex §2ª lei; painel_estelar.tex §II"},"origem":"wiki","palavras_chave":[],"sinonimos":[],"tipo":"conceito","titulo":"Entropia = log φ (entropia topológica de Bowen) e o teorema H"}
\section{Entropia = log \ensuremath{\varphi} (entropia topológica de Bowen) e o teorema H}
A entropia é h=log \ensuremath{\rho}(A)=log \ensuremath{\varphi} (Bowen), = expoente de Lyapunov = taxa de entropia topológica do fluxo do bit; verificada \ensuremath{\approx}0,4812 (ouro), 0,8814 (prata), 1,1948 (bronze). É a 2ª lei da Termodinâmica de Ouro (teorema H): a entropia relativa D(\ensuremath{\rho}\_t\ensuremath{\|}\ensuremath{\mu}) decresce monotonicamente (\ensuremath{\sigma}\ensuremath{\ge}0) sob Perron-Frobenius, com a medida de Parry de razão \ensuremath{\varphi}\ensuremath{^{2}}. Verificado numericamente (D cai 1,21\ensuremath{\to}0,73\ensuremath{\to}0,29\ensuremath{\to}…\ensuremath{\to}0).
\prov{relações: explica associador\_calor; parte\_de tres\_bracos\_estelar; relaciona entropia\_shannon; prova transformada\_metalica\_fechada\_perelman}
\prov{proveniência: wiki \textperiodcentered{} papers/corpo\_dual.tex §2ª lei; painel\_estelar.tex §II \textperiodcentered{} fato \textperiodcentered{} confiança 1.0 \textperiodcentered{} domínio cosmologia \textperiodcentered{} história 6 versões no jornal}

%CRISTAL {"arestas":[{"alvo":"energia_escura_campo_psi","confianca":1.0,"epistemico":"fato","origem":"wiki","rel":"relaciona"},{"alvo":"cosmologia_quantica_hub","confianca":1.0,"epistemico":"fato","origem":"wiki","rel":"parte_de"}],"confianca":1.0,"contraexemplos":[],"descricao":"w=p/ρ da energia escura. Λ: w=−1 exato. CPL (Chevallier-Polarski-Linder): w(a)=w₀+wₐ(1−a) — a parametrização padrão (ferramenta de ajuste, não teoria; validade questionada, arXiv:2510.04191). Quintessência: campo escalar −1<w<−1/3. Phantom: w<−1 (densidade cresce, 'Big Rip'). Gravidade modificada f(R): aceleração sem energia escura.","epistemico":"inferencia","exemplos":[],"id":"eos_w_cpl_modelos","memoria":"semantica","meta":{"autor":"broca-so","dominio":"cosmologia_dados","fonte":"Chevallier-Polarski 2001; Linder 2003 PRL 90,091301"},"origem":"wiki","palavras_chave":[],"sinonimos":[],"tipo":"conceito","titulo":"w(z), CPL e os modelos de energia escura"}
\section{w(z), CPL e os modelos de energia escura}
w=p/\ensuremath{\rho} da energia escura. \ensuremath{\Lambda}: w=\ensuremath{-}1 exato. CPL (Chevallier-Polarski-Linder): w(a)=w\ensuremath{_{0}}+w\ensuremath{_{a}}(1\ensuremath{-}a) — a parametrização padrão (ferramenta de ajuste, não teoria; validade questionada, arXiv:2510.04191). Quintessência: campo escalar \ensuremath{-}1<w<\ensuremath{-}1/3. Phantom: w<\ensuremath{-}1 (densidade cresce, 'Big Rip'). Gravidade modificada f(R): aceleração sem energia escura.
\prov{relações: relaciona energia\_escura\_campo\_psi; parte\_de cosmologia\_quantica\_hub}
\prov{proveniência: wiki \textperiodcentered{} Chevallier-Polarski 2001; Linder 2003 PRL 90,091301 \textperiodcentered{} inferencia \textperiodcentered{} confiança 1.0 \textperiodcentered{} domínio cosmologia\_dados \textperiodcentered{} história 3 versões no jornal}

%CRISTAL {"arestas":[{"alvo":"solidos_multiplicativos","confianca":1.0,"epistemico":"fato","origem":"wiki","rel":"especializa"},{"alvo":"corpo_estelar","confianca":1.0,"epistemico":"fato","origem":"wiki","rel":"usa"},{"alvo":"algebra_reta_mineral","confianca":1.0,"epistemico":"fato","origem":"wiki","rel":"usa"},{"alvo":"joalheria_hub","confianca":1.0,"epistemico":"fato","origem":"wiki","rel":"parte_de"}],"confianca":1.0,"contraexemplos":[],"descricao":"A gema fechada numa EQUAÇÃO, provada exata no metal (erro ~1e-16, cada vértice obedece): 𝒢(θ,z) = ℓ·C_k(θ)·P(z)·𝕊¹(θ) ⊕ h·z·ẑ. Os FATORES PRIMITIVOS, ligados só por ⊗ (produto) e ⊕ (soma espectral de modos): 𝕊¹(θ)=(cosθ,sinθ) é o CÍRCULO (base angular); C_k(θ)=1+a·cos(k(θ−φ)) é o TRILHÃO (círculo ⊕ a·estrela_k — DC mais o k-ésimo modo); √(1−z²) é a ESFERA/círculo meridional (a portadora redonda); (1+s·z) é a RETA (a modulação assimétrica: topo largo/base fina); √((z+1)/(z_eq+1)) é a PARÁBOLA (perfil de base alternativo). O perfil P(z)=√(1−z²)·(1+s·z) = esfera ⊗ modulação. As FACETAS: N=k·m nós (La Hire, estrela_k ⊗ refino_m). Nada fora de círculo/trilhão/esfera/parábola/reta e ⊕/⊗. Verificado em reconstrucao_3d.verificar_reconstrucao ('a EQUAÇÃO é o pipe'). Papers: joalheria.tex, turmalina.tex.","epistemico":"fato","exemplos":[],"id":"equacao_gema_estelar","memoria":"semantica","meta":{"dominio":"joalheria","fonte":"maquina de corte"},"origem":"wiki","palavras_chave":[],"sinonimos":[],"tipo":"conceito","titulo":"A Equação da Gema (produto estelar-espectral: círculo⊗trilhão⊗esfera ⊕ modos)"}
\section{A Equação da Gema (produto estelar-espectral: círculo\ensuremath{\otimes}trilhão\ensuremath{\otimes}esfera \ensuremath{\oplus} modos)}
A gema fechada numa EQUAÇÃO, provada exata no metal (erro \~{}1e-16, cada vértice obedece): \ensuremath{\mathcal{G}}(\ensuremath{\theta},z) = \ensuremath{\ell}\textperiodcentered C\_k(\ensuremath{\theta})\textperiodcentered P(z)\textperiodcentered \ensuremath{\mathbb{S}}\ensuremath{^{1}}(\ensuremath{\theta}) \ensuremath{\oplus} h\textperiodcentered z\textperiodcentered \ensuremath{\hat{z}}. Os FATORES PRIMITIVOS, ligados só por \ensuremath{\otimes} (produto) e \ensuremath{\oplus} (soma espectral de modos): \ensuremath{\mathbb{S}}\ensuremath{^{1}}(\ensuremath{\theta})=(cos\ensuremath{\theta},sin\ensuremath{\theta}) é o CÍRCULO (base angular); C\_k(\ensuremath{\theta})=1+a\textperiodcentered cos(k(\ensuremath{\theta}\ensuremath{-}\ensuremath{\varphi})) é o TRILHÃO (círculo \ensuremath{\oplus} a\textperiodcentered estrela\_k — DC mais o k-ésimo modo); \ensuremath{\sqrt{\,}}(1\ensuremath{-}z\ensuremath{^{2}}) é a ESFERA/círculo meridional (a portadora redonda); (1+s\textperiodcentered z) é a RETA (a modulação assimétrica: topo largo/base fina); \ensuremath{\sqrt{\,}}((z+1)/(z\_eq+1)) é a PARÁBOLA (perfil de base alternativo). O perfil P(z)=\ensuremath{\sqrt{\,}}(1\ensuremath{-}z\ensuremath{^{2}})\textperiodcentered (1+s\textperiodcentered z) = esfera \ensuremath{\otimes} modulação. As FACETAS: N=k\textperiodcentered m nós (La Hire, estrela\_k \ensuremath{\otimes} refino\_m). Nada fora de círculo/trilhão/esfera/parábola/reta e \ensuremath{\oplus}/\ensuremath{\otimes}. Verificado em reconstrucao\_3d.verificar\_reconstrucao ('a EQUAÇÃO é o pipe'). Papers: joalheria.tex, turmalina.tex.
\prov{relações: especializa solidos\_multiplicativos; usa corpo\_estelar; usa algebra\_reta\_mineral; parte\_de joalheria\_hub}
\prov{proveniência: wiki \textperiodcentered{} maquina de corte \textperiodcentered{} fato \textperiodcentered{} confiança 1.0 \textperiodcentered{} domínio joalheria \textperiodcentered{} história 190 versões no jornal}

%CRISTAL {"arestas":[{"alvo":"hubble_lemaitre","confianca":1.0,"epistemico":"fato","origem":"wiki","rel":"explica"},{"alvo":"einstein_de_sitter_modelo","confianca":1.0,"epistemico":"fato","origem":"wiki","rel":"especializa"},{"alvo":"de_sitter_lambda_energia_escura","confianca":1.0,"epistemico":"fato","origem":"wiki","rel":"relaciona"},{"alvo":"desitter_4d_universo","confianca":1.0,"epistemico":"fato","origem":"wiki","rel":"relaciona"},{"alvo":"cosmologia_quantica_hub","confianca":1.0,"epistemico":"fato","origem":"wiki","rel":"parte_de"}],"confianca":1.0,"contraexemplos":[],"descricao":"H²=(ȧ/a)²=8πGρ/3 − kc²/a² + Λc²/3 (expansão); ä/a=−4πG/3(ρ+3p/c²)+Λc²/3 (aceleração). H=ȧ/a é o parâmetro de Hubble. Base dinâmica do ΛCDM. É a mesma forma da s̈=2s do Lopes em minisuperspace.","epistemico":"fato","exemplos":[],"id":"equacoes_friedmann","memoria":"semantica","meta":{"autor":"broca-so","dominio":"fisica","fonte":"Friedmann 1922"},"origem":"wiki","palavras_chave":[],"sinonimos":[],"tipo":"conceito","titulo":"Equações de Friedmann (EFE na métrica FLRW)"}
\section{Equações de Friedmann (EFE na métrica FLRW)}
H\ensuremath{^{2}}=(\ensuremath{\dot{a}}/a)\ensuremath{^{2}}=8\ensuremath{\pi}G\ensuremath{\rho}/3 \ensuremath{-} kc\ensuremath{^{2}}/a\ensuremath{^{2}} + \ensuremath{\Lambda}c\ensuremath{^{2}}/3 (expansão); ä/a=\ensuremath{-}4\ensuremath{\pi}G/3(\ensuremath{\rho}+3p/c\ensuremath{^{2}})+\ensuremath{\Lambda}c\ensuremath{^{2}}/3 (aceleração). H=\ensuremath{\dot{a}}/a é o parâmetro de Hubble. Base dinâmica do \ensuremath{\Lambda}CDM. É a mesma forma da s\ensuremath{}=2s do Lopes em minisuperspace.
\prov{relações: explica hubble\_lemaitre; especializa einstein\_de\_sitter\_modelo; relaciona de\_sitter\_lambda\_energia\_escura; relaciona desitter\_4d\_universo; parte\_de cosmologia\_quantica\_hub}
\prov{proveniência: wiki \textperiodcentered{} Friedmann 1922 \textperiodcentered{} fato \textperiodcentered{} confiança 1.0 \textperiodcentered{} domínio fisica \textperiodcentered{} história 7 versões no jornal}

%CRISTAL {"arestas":[{"alvo":"base_pi_bijecao","confianca":1.0,"epistemico":"fato","origem":"wiki","rel":"relaciona"},{"alvo":"fermat_na_maquina","confianca":1.0,"epistemico":"fato","origem":"wiki","rel":"relaciona"},{"alvo":"joalheria_hub","confianca":1.0,"epistemico":"fato","origem":"wiki","rel":"parte_de"}],"confianca":1.0,"contraexemplos":[],"descricao":"A TORRE que fecha o corpo estelar. (1) O AGM é as DUAS operações do corpo estelar: a média ARITMÉTICA = a RETA (soma) e a GEOMÉTRICA = LA HIRE (produto/√), iteradas ao invariante M(a,b). π GERA a constante LEMNISCATA: ϖ=π/M(1,√2)=π·G (Gauss); e o AGM gera π (Brent-Salamin) e os períodos elípticos K. (2) Os METAIS são as FRAÇÕES CONTÍNUAS periódicas: σ_n=[n;n,n,…] (Lagrange: periódica ⟺ quadrática); o ouro φ=σ_1=[1;1,…]. A reta mineral (os metais, algébricos/Pisot) é o CRISTAL; π=[3;7,15,…] aperiódica é o CAOS. (3) π GERA O OURO pela fração contínua de ROGERS-RAMANUJAN: R(e^{-2π})=√(φ√5)−φ (Ramanujan 1913, verificado 40 dígitos, raiz de x⁴+2x³−6x²−2x+1) — uma CF 'na base π' (q=e^{-2π}). A PONTE: o modular em argumentos-π é ALGÉBRICO (multiplicação complexa) — π (caos) gera o cristal. (4) A ESCADA METÁLICA: reta_real/ouro (m=1) → prata (m=2) → bronze (m=3) → …; cada reta de ordem m gerada por σ_m=[m;m,…], as combinações via PA (aritmética/soma) e PG (geométrica/produto) — a recorrência de ordem m = o DUPOLINÔMIO x²−mx−1 (P_g=x², P_a=mx+1) —, cada uma com a SUA parábola (λ=log σ_m, cristal). π é a ORIGEM: gera o ouro, fechando o laço (a reta mineral, gerador π, vira a origem). ESCOPO honesto: FATO — σ_n=[n;n,…], R(e^{-2π})=√(φ√5)−φ, ϖ=π/M(1,√2); os R(e^{-2π√n}) de n>1 são invariantes de classe (Ramanujan-Watson, citados, grau>24). linguagem/agm.py (4/4), fracoes_continuas.py (5/5).","epistemico":"fato","exemplos":[],"id":"escada_metalica_pi","memoria":"semantica","meta":{"dominio":"joalheria","fonte":"maquina de corte"},"origem":"wiki","palavras_chave":[],"sinonimos":[],"tipo":"conceito","titulo":"A escada de retas metálicas: cada reta por σ_m=[m;m,…] via PA/PG, e π gera o ouro (Rogers-Ramanujan)"}
\section{A escada de retas metálicas: cada reta por \ensuremath{\sigma}\_m=[m;m,…] via PA/PG, e \ensuremath{\pi} gera o ouro (Rogers-Ramanujan)}
A TORRE que fecha o corpo estelar. (1) O AGM é as DUAS operações do corpo estelar: a média ARITMÉTICA = a RETA (soma) e a GEOMÉTRICA = LA HIRE (produto/\ensuremath{\sqrt{\,}}), iteradas ao invariante M(a,b). \ensuremath{\pi} GERA a constante LEMNISCATA: \ensuremath{\varpi}=\ensuremath{\pi}/M(1,\ensuremath{\sqrt{\,}}2)=\ensuremath{\pi}\textperiodcentered G (Gauss); e o AGM gera \ensuremath{\pi} (Brent-Salamin) e os períodos elípticos K. (2) Os METAIS são as FRAÇÕES CONTÍNUAS periódicas: \ensuremath{\sigma}\_n=[n;n,n,…] (Lagrange: periódica \ensuremath{\Longleftrightarrow} quadrática); o ouro \ensuremath{\varphi}=\ensuremath{\sigma}\_1=[1;1,…]. A reta mineral (os metais, algébricos/Pisot) é o CRISTAL; \ensuremath{\pi}=[3;7,15,…] aperiódica é o CAOS. (3) \ensuremath{\pi} GERA O OURO pela fração contínua de ROGERS-RAMANUJAN: R(e\^{}\{-2\ensuremath{\pi}\})=\ensuremath{\sqrt{\,}}(\ensuremath{\varphi}\ensuremath{\sqrt{\,}}5)\ensuremath{-}\ensuremath{\varphi} (Ramanujan 1913, verificado 40 dígitos, raiz de x\ensuremath{^{4}}+2x\ensuremath{^{3}}\ensuremath{-}6x\ensuremath{^{2}}\ensuremath{-}2x+1) — uma CF 'na base \ensuremath{\pi}' (q=e\^{}\{-2\ensuremath{\pi}\}). A PONTE: o modular em argumentos-\ensuremath{\pi} é ALGÉBRICO (multiplicação complexa) — \ensuremath{\pi} (caos) gera o cristal. (4) A ESCADA METÁLICA: reta\_real/ouro (m=1) \ensuremath{\to} prata (m=2) \ensuremath{\to} bronze (m=3) \ensuremath{\to} …; cada reta de ordem m gerada por \ensuremath{\sigma}\_m=[m;m,…], as combinações via PA (aritmética/soma) e PG (geométrica/produto) — a recorrência de ordem m = o DUPOLINÔMIO x\ensuremath{^{2}}\ensuremath{-}mx\ensuremath{-}1 (P\_g=x\ensuremath{^{2}}, P\_a=mx+1) —, cada uma com a SUA parábola (\ensuremath{\lambda}=log \ensuremath{\sigma}\_m, cristal). \ensuremath{\pi} é a ORIGEM: gera o ouro, fechando o laço (a reta mineral, gerador \ensuremath{\pi}, vira a origem). ESCOPO honesto: FATO — \ensuremath{\sigma}\_n=[n;n,…], R(e\^{}\{-2\ensuremath{\pi}\})=\ensuremath{\sqrt{\,}}(\ensuremath{\varphi}\ensuremath{\sqrt{\,}}5)\ensuremath{-}\ensuremath{\varphi}, \ensuremath{\varpi}=\ensuremath{\pi}/M(1,\ensuremath{\sqrt{\,}}2); os R(e\^{}\{-2\ensuremath{\pi}\ensuremath{\sqrt{\,}}n\}) de n>1 são invariantes de classe (Ramanujan-Watson, citados, grau>24). linguagem/agm.py (4/4), fracoes\_continuas.py (5/5).
\prov{relações: relaciona base\_pi\_bijecao; relaciona fermat\_na\_maquina; parte\_de joalheria\_hub}
\prov{proveniência: wiki \textperiodcentered{} maquina de corte \textperiodcentered{} fato \textperiodcentered{} confiança 1.0 \textperiodcentered{} domínio joalheria \textperiodcentered{} história 4 versões no jornal}

%CRISTAL {"arestas":[],"confianca":1.0,"contraexemplos":[],"descricao":"Onde a gravidade quântica é inevitável, de G, c, ℏ: E_P≈1,22×10¹⁹ GeV, ℓ_P=√(ℏG/c³)≈1,62×10⁻³⁵ m, t_P≈5,4×10⁻⁴⁴ s. Fato dimensional robusto — ~15 ordens acima do LHC (~10⁴ GeV), o que explica a ausência de testes diretos.","epistemico":"fato","exemplos":[],"id":"escala_planck","memoria":"semantica","meta":{"autor":"broca-so","dominio":"fisica","fonte":"Planck 1899; CODATA/PDG"},"origem":"wiki","palavras_chave":[],"sinonimos":[],"tipo":"conceito","titulo":"A escala de Planck"}
\section{A escala de Planck}
Onde a gravidade quântica é inevitável, de G, c, \ensuremath{\hbar}: E\_P\ensuremath{\approx}1,22$\times$10\ensuremath{^{19}} GeV, \ensuremath{\ell}\_P=\ensuremath{\sqrt{\,}}(\ensuremath{\hbar}G/c\ensuremath{^{3}})\ensuremath{\approx}1,62$\times$10\ensuremath{^{-35}} m, t\_P\ensuremath{\approx}5,4$\times$10\ensuremath{^{-44}} s. Fato dimensional robusto — \~{}15 ordens acima do LHC (\~{}10\ensuremath{^{4}} GeV), o que explica a ausência de testes diretos.
\prov{proveniência: wiki \textperiodcentered{} Planck 1899; CODATA/PDG \textperiodcentered{} fato \textperiodcentered{} confiança 1.0 \textperiodcentered{} domínio fisica}

%CRISTAL {"arestas":[{"alvo":"a_base_topologica","confianca":1.0,"epistemico":"fato","origem":"wiki","rel":"usa"},{"alvo":"dois_qubits_octonios","confianca":1.0,"epistemico":"fato","origem":"wiki","rel":"usa"}],"confianca":1.0,"contraexemplos":[],"descricao":"A esfera S² de direções â dos campos locais É a esfera de Bloch (=ℂP¹), a base da fibração de Hopf S³→S²; a 3-esfera S³=q∈ℍ:‖q‖=1=SU(2) dos versores é o espaço de estados de um qubit (o total do Hopf). Uma seção de 𝒪(−1) sobre a base é a função de onda; o 'i' de Schrödinger é a trivialização local (carta ℂâ), e mudar de carta é uma transformação de gauge U(1). Um qubit=ℍ, dois qubits=𝕆. [D] identidade por dimensão.","epistemico":"inferencia","exemplos":[],"id":"esfera_bloch_base_hopf","memoria":"semantica","meta":{"autor":"broca-so","dominio":"cosmologia","fonte":"paper_3_quantica.tex sec:onda; paper_1_algebra_geometria.tex"},"origem":"wiki","palavras_chave":[],"sinonimos":[],"tipo":"conceito","titulo":"A esfera de Bloch = S² base do fibrado de Hopf"}
\section{A esfera de Bloch = S\ensuremath{^{2}} base do fibrado de Hopf}
A esfera S\ensuremath{^{2}} de direções â dos campos locais É a esfera de Bloch (=\ensuremath{\mathbb{C}}P\ensuremath{^{1}}), a base da fibração de Hopf S\ensuremath{^{3}}\ensuremath{\to}S\ensuremath{^{2}}; a 3-esfera S\ensuremath{^{3}}=q\ensuremath{\in}\ensuremath{\mathbb{H}}:\ensuremath{\|}q\ensuremath{\|}=1=SU(2) dos versores é o espaço de estados de um qubit (o total do Hopf). Uma seção de \ensuremath{\mathcal{O}}(\ensuremath{-}1) sobre a base é a função de onda; o 'i' de Schrödinger é a trivialização local (carta \ensuremath{\mathbb{C}}â), e mudar de carta é uma transformação de gauge U(1). Um qubit=\ensuremath{\mathbb{H}}, dois qubits=\ensuremath{\mathbb{O}}. [D] identidade por dimensão.
\prov{relações: usa a\_base\_topologica; usa dois\_qubits\_octonios}
\prov{proveniência: wiki \textperiodcentered{} paper\_3\_quantica.tex sec:onda; paper\_1\_algebra\_geometria.tex \textperiodcentered{} inferencia \textperiodcentered{} confiança 1.0 \textperiodcentered{} domínio cosmologia \textperiodcentered{} história 4 versões no jornal}

%CRISTAL {"arestas":[{"alvo":"area_render_gema","confianca":1.0,"epistemico":"fato","origem":"wiki","rel":"parte_de"},{"alvo":"area_optica_gema","confianca":1.0,"epistemico":"fato","origem":"wiki","rel":"relaciona"},{"alvo":"pedra_logo_3d","confianca":1.0,"epistemico":"fato","origem":"wiki","rel":"relaciona"},{"alvo":"pipe_svg_baricentrico","confianca":1.0,"epistemico":"fato","origem":"wiki","rel":"relaciona"}],"confianca":1.0,"contraexemplos":[],"descricao":"Em vez de atacar tudo (o C nativo de uma vez), DIVIDIR: a luz na gema é a soma de 7 ÓRBITAS espectrais (as 7 bandas da dispersão, violeta→vermelho — o mesmo espectro do floco refinado em metais). linguagem/espectro_gema.py decompõe uma imagem de gema nessas 7 camadas (NMF não-negativa) e mede o espectro. O estudo do gabarito do GPT deu o alvo: o CORPO é a janela azul-verde (ciano+verde+azul = 84%, o azul Paraíba, preenche o volume); o FOGO é o resto (violeta+amarelo+laranja+vermelho = 16%, só nas BORDAS e facetas, onde a dispersão separa). A nossa pedra tem só 6,4% de fogo (metade) e verde a mais — falta dispersão, o alvo do próximo refino. Divide o render em corpo + fogo → render por banda em paralelo (várias instâncias do broca, uma por órbita), depois o C nativo com BVH. Verificado: espectro_gema 3/3.","epistemico":"fato","exemplos":[],"id":"espectro_das_orbitas","memoria":"semantica","meta":{"dominio":"joalheria","fonte":"espectro orbitas"},"origem":"wiki","palavras_chave":[],"sinonimos":[],"tipo":"conceito","titulo":"O Espectro das Órbitas (dividir a luz nas 7 bandas)"}
\section{O Espectro das Órbitas (dividir a luz nas 7 bandas)}
Em vez de atacar tudo (o C nativo de uma vez), DIVIDIR: a luz na gema é a soma de 7 ÓRBITAS espectrais (as 7 bandas da dispersão, violeta\ensuremath{\to}vermelho — o mesmo espectro do floco refinado em metais). linguagem/espectro\_gema.py decompõe uma imagem de gema nessas 7 camadas (NMF não-negativa) e mede o espectro. O estudo do gabarito do GPT deu o alvo: o CORPO é a janela azul-verde (ciano+verde+azul = 84\%, o azul Paraíba, preenche o volume); o FOGO é o resto (violeta+amarelo+laranja+vermelho = 16\%, só nas BORDAS e facetas, onde a dispersão separa). A nossa pedra tem só 6,4\% de fogo (metade) e verde a mais — falta dispersão, o alvo do próximo refino. Divide o render em corpo + fogo \ensuremath{\to} render por banda em paralelo (várias instâncias do broca, uma por órbita), depois o C nativo com BVH. Verificado: espectro\_gema 3/3.
\prov{relações: parte\_de area\_render\_gema; relaciona area\_optica\_gema; relaciona pedra\_logo\_3d; relaciona pipe\_svg\_baricentrico}
\prov{proveniência: wiki \textperiodcentered{} espectro orbitas \textperiodcentered{} fato \textperiodcentered{} confiança 1.0 \textperiodcentered{} domínio joalheria \textperiodcentered{} história 30 versões no jornal}

%CRISTAL {"arestas":[{"alvo":"symbolic_beta_protocolo","confianca":1.0,"epistemico":"fato","origem":"wiki","rel":"parte_de"},{"alvo":"mineracao_so_a_prata_pelo_hexagrama","confianca":1.0,"epistemico":"fato","origem":"wiki","rel":"usa"},{"alvo":"hexagrama_mineracao","confianca":1.0,"epistemico":"fato","origem":"wiki","rel":"usa"},{"alvo":"entropia_ouro_para_prata","confianca":1.0,"epistemico":"fato","origem":"wiki","rel":"relaciona"}],"confianca":1.0,"contraexemplos":[],"descricao":"'Ajustar o symbolic pro cálculo do espectro na PRATA' = fixar o alvo do espectro em 6 modos. goldchain/prova.py + espectro.py: espectro do neurônio (12⁴) → matriz lado×lado (matriz_do_espectro, núcleo rank≤6 perturbado pelo tecido) → SVD → conta modos significativos (LIM_REL=0,04). 5 modos = OURO (pentagrama): prova PARCIAL, não emite ('busque nonce p/ hexagrama'). 6 modos = PRATA (hexagrama de Salomão, busca_estrela): EMITE Pell → cifra 𝔽₂ᵏ. Ou seja: o cálculo do espectro só CONCLUI (emite) na dimensão de prata, pelo hexagrama — é o espectro da prata que carrega o lastro. [D] está no código (prova.classifica, busca_estrela).","epistemico":"fato","exemplos":[],"id":"espectro_prata_seis_modos","memoria":"semantica","meta":{"autor":"Lopes/broca-so","dominio":"mineracao","fonte":"goldchain/prova.py (MODOS_PRATA=6, n_modos, classifica); espectro.py"},"origem":"wiki","palavras_chave":[],"sinonimos":[],"tipo":"conceito","titulo":"O symbolic do espectro na dimensão de prata: 6 modos SVD (hexagrama) emitem Pell"}
\section{O symbolic do espectro na dimensão de prata: 6 modos SVD (hexagrama) emitem Pell}
'Ajustar o symbolic pro cálculo do espectro na PRATA' = fixar o alvo do espectro em 6 modos. goldchain/prova.py + espectro.py: espectro do neurônio (12\ensuremath{^{4}}) \ensuremath{\to} matriz lado$\times$lado (matriz\_do\_espectro, núcleo rank\ensuremath{\le}6 perturbado pelo tecido) \ensuremath{\to} SVD \ensuremath{\to} conta modos significativos (LIM\_REL=0,04). 5 modos = OURO (pentagrama): prova PARCIAL, não emite ('busque nonce p/ hexagrama'). 6 modos = PRATA (hexagrama de Salomão, busca\_estrela): EMITE Pell \ensuremath{\to} cifra \ensuremath{\mathbb{F}}\ensuremath{_{2}}\ensuremath{^{k}}. Ou seja: o cálculo do espectro só CONCLUI (emite) na dimensão de prata, pelo hexagrama — é o espectro da prata que carrega o lastro. [D] está no código (prova.classifica, busca\_estrela).
\prov{relações: parte\_de symbolic\_beta\_protocolo; usa mineracao\_so\_a\_prata\_pelo\_hexagrama; usa hexagrama\_mineracao; relaciona entropia\_ouro\_para\_prata}
\prov{proveniência: wiki \textperiodcentered{} goldchain/prova.py (MODOS\_PRATA=6, n\_modos, classifica); espectro.py \textperiodcentered{} fato \textperiodcentered{} confiança 1.0 \textperiodcentered{} domínio mineracao \textperiodcentered{} história 25 versões no jornal}

%CRISTAL {"arestas":[],"confianca":1.0,"contraexemplos":[],"descricao":"Estado estacionário = autoestado do Hamiltoniano, energia definida, distribuição independente do tempo; níveis discretos. Regras de seleção determinam transições permitidas; a conservação de energia se manifesta na emissão/absorção de quanta. O análogo canônico do limite não-dissipativo (Noether) do BAI.","epistemico":"fato","exemplos":[],"id":"estados_estacionarios_selecao","memoria":"semantica","meta":{"autor":"broca-so","dominio":"fisica","fonte":"Bohr-Sommerfeld 1915-16; Schrödinger 1926"},"origem":"wiki","palavras_chave":[],"sinonimos":[],"tipo":"conceito","titulo":"Estados estacionários e leis de seleção (conservação)"}
\section{Estados estacionários e leis de seleção (conservação)}
Estado estacionário = autoestado do Hamiltoniano, energia definida, distribuição independente do tempo; níveis discretos. Regras de seleção determinam transições permitidas; a conservação de energia se manifesta na emissão/absorção de quanta. O análogo canônico do limite não-dissipativo (Noether) do BAI.
\prov{proveniência: wiki \textperiodcentered{} Bohr-Sommerfeld 1915-16; Schrödinger 1926 \textperiodcentered{} fato \textperiodcentered{} confiança 1.0 \textperiodcentered{} domínio fisica}

%CRISTAL {"arestas":[{"alvo":"grupo_padrao_cxho","confianca":1.0,"epistemico":"fato","origem":"wiki","rel":"explica"},{"alvo":"gravidade_desitter_derivada","confianca":1.0,"epistemico":"fato","origem":"wiki","rel":"explica"},{"alvo":"regua_de_gentil","confianca":1.0,"epistemico":"fato","origem":"wiki","rel":"relaciona"}],"confianca":1.0,"contraexemplos":[],"descricao":"O princípio metodológico da série: as identificações do setor de matéria são CORRESPONDÊNCIAS estruturais (a torre contém, sem ajuste, objetos com exatamente os números quânticos e a estrutura de grupo do MP, forçados pela maximalidade das álgebras de divisão), NÃO derivações físicas (não há Lagrangiana, massas, dinâmica). 'A distância entre a correspondência e a derivação é o programa; nós a marcamos, não a escondemos.' Tags [firme]/[modelo] percorrem a obra — é o próprio eixo DNCI do grafo.","epistemico":"inferencia","exemplos":[],"id":"estatuto_correspondencia_derivacao","memoria":"semantica","meta":{"autor":"broca-so","dominio":"cosmologia","fonte":"paper_4_particulas.tex obs:estatuto; paper_6_maquina.tex §estatuto"},"origem":"wiki","palavras_chave":[],"sinonimos":[],"tipo":"conceito","titulo":"O estatuto honesto: correspondência ≠ derivação"}
\section{O estatuto honesto: correspondência \ensuremath{\ne} derivação}
O princípio metodológico da série: as identificações do setor de matéria são CORRESPONDÊNCIAS estruturais (a torre contém, sem ajuste, objetos com exatamente os números quânticos e a estrutura de grupo do MP, forçados pela maximalidade das álgebras de divisão), NÃO derivações físicas (não há Lagrangiana, massas, dinâmica). 'A distância entre a correspondência e a derivação é o programa; nós a marcamos, não a escondemos.' Tags [firme]/[modelo] percorrem a obra — é o próprio eixo DNCI do grafo.
\prov{relações: explica grupo\_padrao\_cxho; explica gravidade\_desitter\_derivada; relaciona regua\_de\_gentil}
\prov{proveniência: wiki \textperiodcentered{} paper\_4\_particulas.tex obs:estatuto; paper\_6\_maquina.tex §estatuto \textperiodcentered{} inferencia \textperiodcentered{} confiança 1.0 \textperiodcentered{} domínio cosmologia \textperiodcentered{} história 4 versões no jornal}

%CRISTAL {"arestas":[{"alvo":"teorema_esterilidade_dimensao","confianca":1.0,"epistemico":"fato","origem":"wiki","rel":"confirma"},{"alvo":"independencia_retas_metalicas","confianca":1.0,"epistemico":"fato","origem":"wiki","rel":"confirma"},{"alvo":"reta_metalica_geral","confianca":1.0,"epistemico":"fato","origem":"wiki","rel":"usa"}],"confianca":1.0,"contraexemplos":[],"descricao":"doc/fractais/esterilidademetal.c implementa o experimento EM INTEIROS (sem float): perna 1 (fecho ouro: σk=(Fₖ-₁,Fₖ), norma ±1 — nada escapa de ⟨1,φ⟩); perna 2 (fecho prata: coord √2 = Pell 1,2,5,12,29,…, norma a²−2b²=±1); perna 3 (independência: 0 soluções de (a+bφ)²=2c² com |a|,|b|≤4000). Compila com cc, roda, exit 0. É a tese no metal do broca-so: a esterilidade e a independência provadas em aritmética exata, não em ponto flutuante.","epistemico":"fato","exemplos":[],"id":"esterilidade_no_metal","memoria":"semantica","meta":{"autor":"Lopes/broca-so","dominio":"metal","fonte":"doc/fractais/esterilidade_metal.c (cc -O2, exit 0)"},"origem":"wiki","palavras_chave":[],"sinonimos":[],"tipo":"conceito","titulo":"A prova no metal (C, inteiros exatos)"}
\section{A prova no metal (C, inteiros exatos)}
doc/fractais/esterilidademetal.c implementa o experimento EM INTEIROS (sem float): perna 1 (fecho ouro: \ensuremath{\sigma}k=(F\ensuremath{_{k}}-\ensuremath{_{1}},F\ensuremath{_{k}}), norma $\pm$1 — nada escapa de \ensuremath{\langle}1,\ensuremath{\varphi}\ensuremath{\rangle}); perna 2 (fecho prata: coord \ensuremath{\sqrt{\,}}2 = Pell 1,2,5,12,29,…, norma a\ensuremath{^{2}}\ensuremath{-}2b\ensuremath{^{2}}=$\pm$1); perna 3 (independência: 0 soluções de (a+b\ensuremath{\varphi})\ensuremath{^{2}}=2c\ensuremath{^{2}} com |a|,|b|\ensuremath{\le}4000). Compila com cc, roda, exit 0. É a tese no metal do broca-so: a esterilidade e a independência provadas em aritmética exata, não em ponto flutuante.
\prov{relações: confirma teorema\_esterilidade\_dimensao; confirma independencia\_retas\_metalicas; usa reta\_metalica\_geral}
\prov{proveniência: wiki \textperiodcentered{} doc/fractais/esterilidade\_metal.c (cc -O2, exit 0) \textperiodcentered{} fato \textperiodcentered{} confiança 1.0 \textperiodcentered{} domínio metal \textperiodcentered{} história 5 versões no jornal}

%CRISTAL {"arestas":[{"alvo":"materia_escura","confianca":1.0,"epistemico":"fato","origem":"wiki","rel":"confirma"},{"alvo":"inflacao_cosmica","confianca":1.0,"epistemico":"fato","origem":"wiki","rel":"relaciona"}],"confianca":1.0,"contraexemplos":[],"descricao":"A teia cósmica (galáxias, filamentos, vazios) só cresce a tempo se há matéria escura FRIA fornecendo poços de potencial antes do desacoplamento dos bárions. Simulações N-corpos ΛCDM (Millennium) casam com surveys (SDSS, DES). Favorece CDM sobre matéria escura quente.","epistemico":"fato","exemplos":[],"id":"estrutura_larga_escala_cdm","memoria":"semantica","meta":{"autor":"broca-so","dominio":"cosmologia_dados","fonte":"Springel et al. 2005 (Millennium)"},"origem":"wiki","palavras_chave":[],"sinonimos":[],"tipo":"conceito","titulo":"Formação de estrutura de larga escala"}
\section{Formação de estrutura de larga escala}
A teia cósmica (galáxias, filamentos, vazios) só cresce a tempo se há matéria escura FRIA fornecendo poços de potencial antes do desacoplamento dos bárions. Simulações N-corpos \ensuremath{\Lambda}CDM (Millennium) casam com surveys (SDSS, DES). Favorece CDM sobre matéria escura quente.
\prov{relações: confirma materia\_escura; relaciona inflacao\_cosmica}
\prov{proveniência: wiki \textperiodcentered{} Springel et al. 2005 (Millennium) \textperiodcentered{} fato \textperiodcentered{} confiança 1.0 \textperiodcentered{} domínio cosmologia\_dados \textperiodcentered{} história 3 versões no jornal}

%CRISTAL {"arestas":[{"alvo":"fase_de_berry_canonica","confianca":1.0,"epistemico":"fato","origem":"wiki","rel":"explica"},{"alvo":"chern_fibrado_hopf","confianca":1.0,"epistemico":"fato","origem":"wiki","rel":"usa"},{"alvo":"dois_qubits_octonios","confianca":1.0,"epistemico":"fato","origem":"wiki","rel":"usa"}],"confianca":1.0,"contraexemplos":[],"descricao":"A fase de Berry é a holonomia da conexão de 𝒪(−1): φ=½Ω para uma volta de ângulo sólido Ω na Bloch — metade do ângulo sólido, o monopolo de Dirac de carga ½. Em ℂ é abeliana (U(1)); em 𝕆 o autoespaço +i é 4-dim e o transporte por um loop dá holonomia NÃO-ABELIANA de Wilczek-Zee W∈SU(2), com e±iφ: o SPIN-½ como holonomia do circuito elétron↔próton (em Ω=2π, W=−1: a rotação de 2π leva o spinor a −ele). [D] φ=½Ω; [N] inclinação 0,5000 (desvio <10⁻⁴ rad), W=−1 medido.","epistemico":"inferencia","exemplos":[],"id":"fase_berry_holonomia_hopf","memoria":"semantica","meta":{"autor":"broca-so","dominio":"cosmologia","fonte":"paper_3_quantica.tex sec:circuito (prop:holonomia, cor:spin)"},"origem":"wiki","palavras_chave":[],"sinonimos":[],"tipo":"conceito","titulo":"Fase de Berry = holonomia do fibrado (abeliana e SU(2))"}
\section{Fase de Berry = holonomia do fibrado (abeliana e SU(2))}
A fase de Berry é a holonomia da conexão de \ensuremath{\mathcal{O}}(\ensuremath{-}1): \ensuremath{\varphi}=\ensuremath{\tfrac12}\ensuremath{\Omega} para uma volta de ângulo sólido \ensuremath{\Omega} na Bloch — metade do ângulo sólido, o monopolo de Dirac de carga \ensuremath{\tfrac12}. Em \ensuremath{\mathbb{C}} é abeliana (U(1)); em \ensuremath{\mathbb{O}} o autoespaço +i é 4-dim e o transporte por um loop dá holonomia NÃO-ABELIANA de Wilczek-Zee W\ensuremath{\in}SU(2), com e$\pm$i\ensuremath{\varphi}: o SPIN-\ensuremath{\tfrac12} como holonomia do circuito elétron\ensuremath{\leftrightarrow}próton (em \ensuremath{\Omega}=2\ensuremath{\pi}, W=\ensuremath{-}1: a rotação de 2\ensuremath{\pi} leva o spinor a \ensuremath{-}ele). [D] \ensuremath{\varphi}=\ensuremath{\tfrac12}\ensuremath{\Omega}; [N] inclinação 0,5000 (desvio <10\ensuremath{^{-4}} rad), W=\ensuremath{-}1 medido.
\prov{relações: explica fase\_de\_berry\_canonica; usa chern\_fibrado\_hopf; usa dois\_qubits\_octonios}
\prov{proveniência: wiki \textperiodcentered{} paper\_3\_quantica.tex sec:circuito (prop:holonomia, cor:spin) \textperiodcentered{} inferencia \textperiodcentered{} confiança 1.0 \textperiodcentered{} domínio cosmologia \textperiodcentered{} história 5 versões no jornal}

%CRISTAL {"arestas":[],"confianca":1.0,"contraexemplos":[],"descricao":"Um sistema transportado adiabaticamente num circuito fechado adquire, além da fase dinâmica, uma fase GEOMÉTRICA γ(C) que depende só da geometria do caminho: γ(C)=i∮⟨n|∇_R|n⟩·dR; a curvatura de Berry atua como um 'campo magnético' no espaço de parâmetros. Alicerce de toda a família topológica. Berry 1984 (precursores Pancharatnam 1956, Aharonov-Bohm).","epistemico":"fato","exemplos":[],"id":"fase_de_berry_canonica","memoria":"semantica","meta":{"autor":"broca-so","dominio":"fisica","fonte":"Berry 1984 Proc.R.Soc.A 392,45"},"origem":"wiki","palavras_chave":[],"sinonimos":[],"tipo":"conceito","titulo":"A fase geométrica de Berry"}
\section{A fase geométrica de Berry}
Um sistema transportado adiabaticamente num circuito fechado adquire, além da fase dinâmica, uma fase GEOMÉTRICA \ensuremath{\gamma}(C) que depende só da geometria do caminho: \ensuremath{\gamma}(C)=i\ensuremath{\oint}\ensuremath{\langle}n|\ensuremath{\nabla}\_R|n\ensuremath{\rangle}\textperiodcentered dR; a curvatura de Berry atua como um 'campo magnético' no espaço de parâmetros. Alicerce de toda a família topológica. Berry 1984 (precursores Pancharatnam 1956, Aharonov-Bohm).
\prov{proveniência: wiki \textperiodcentered{} Berry 1984 Proc.R.Soc.A 392,45 \textperiodcentered{} fato \textperiodcentered{} confiança 1.0 \textperiodcentered{} domínio fisica}

%CRISTAL {"arestas":[{"alvo":"corpo_estelar_le_o_rastro","confianca":1.0,"epistemico":"fato","origem":"wiki","rel":"usa"},{"alvo":"espectro_das_orbitas","confianca":1.0,"epistemico":"fato","origem":"wiki","rel":"relaciona"}],"confianca":1.0,"contraexemplos":[],"descricao":"A primeira fase é MAPEAMENTO: observação PURA. Não existe 'nossa gema' aqui, não há comparação, nem refino, nem render — só INGERIR o gabarito e VER o espectro dele. O gato Anosov reconstrói EXATO (não aproximado): o SVD é o espectro completo, e reconstruir com todos os modos devolve a imagem a erro 0. linguagem/ingestao_mineral.observar lê o gabarito por inteiro: reconstrução exata (erro 0); a mecânica (os modos POD, o dominante 0.833 + a cauda); a geometria ρ=0.205 (caos); a termodinâmica W=0.83/Q=0.17, H=0.94, 2.6 modos; a parábola metálica 71/22/7 (D_tr=1.70); a cor (janela azul-verde 84% + fogo 16%). É a leitura, não a fabricação — nesta fase não se faz nada, observa-se.","epistemico":"fato","exemplos":[],"id":"fase_observacao_gabarito","memoria":"semantica","meta":{"dominio":"joalheria","fonte":"espectro orbitas"},"origem":"wiki","palavras_chave":[],"sinonimos":[],"tipo":"conceito","titulo":"A Fase de Observação — só ler o espectro do gabarito (não fazer nada)"}
\section{A Fase de Observação — só ler o espectro do gabarito (não fazer nada)}
A primeira fase é MAPEAMENTO: observação PURA. Não existe 'nossa gema' aqui, não há comparação, nem refino, nem render — só INGERIR o gabarito e VER o espectro dele. O gato Anosov reconstrói EXATO (não aproximado): o SVD é o espectro completo, e reconstruir com todos os modos devolve a imagem a erro 0. linguagem/ingestao\_mineral.observar lê o gabarito por inteiro: reconstrução exata (erro 0); a mecânica (os modos POD, o dominante 0.833 + a cauda); a geometria \ensuremath{\rho}=0.205 (caos); a termodinâmica W=0.83/Q=0.17, H=0.94, 2.6 modos; a parábola metálica 71/22/7 (D\_tr=1.70); a cor (janela azul-verde 84\% + fogo 16\%). É a leitura, não a fabricação — nesta fase não se faz nada, observa-se.
\prov{relações: usa corpo\_estelar\_le\_o\_rastro; relaciona espectro\_das\_orbitas}
\prov{proveniência: wiki \textperiodcentered{} espectro orbitas \textperiodcentered{} fato \textperiodcentered{} confiança 1.0 \textperiodcentered{} domínio joalheria \textperiodcentered{} história 6 versões no jornal}

%CRISTAL {"arestas":[{"alvo":"amendoa_hidrogenio_formal","confianca":1.0,"epistemico":"fato","origem":"wiki","rel":"relaciona"},{"alvo":"joalheria_hub","confianca":1.0,"epistemico":"fato","origem":"wiki","rel":"parte_de"}],"confianca":1.0,"contraexemplos":[],"descricao":"Os dois caminhos da prova do TLF caem nos construtos que a máquina já opera. KUMMER (a monografia de Daniel Cunha, UERJ): os corpos CICLOTÔMICOS Q(ζ), ζ=e^{2πi/n} — as RAÍZES DA UNIDADE = as cúspides da ESTRELA; e o TRAÇO/NORMA de um inteiro algébrico (Tr=Σσ_i, N=Πσ_i, por Vieta Tr=−a_1, N=(−1)ⁿa_n) = os METAIS (σ_n é raiz de z²−nz−1, Tr=n, N=−1, uma UNIDADE); o caso n=4 fecha pela DESCIDA infinita = a órbita que contrai z (o cristal, só o ponto fixo trivial). WILES (o paper de 1995): a modularidade via representações de Galois ρ(Frob q), com traço=c(q,f) (o autovalor de HECKE = o espectro = o metal) e det=χ(q)q^{k-1} — o 2×2 traço-det É o GATO A_n (tr=n, det=−1); a modularidade (a_p=c(p,f)) é o CASAMENTO ESPECTRAL, o mesmo gesto do corpo estelar (ler o rastro). RESSALVA: não reprova o TLF nem é a prova de Wiles — é a INGESTÃO (o dicionário exato de Peirce + as peças computáveis verificadas). linguagem/fermat.py (5/5); dicionarios._FERMAT_MAQUINA (Fermat é uma ÁREA da máquina, compõe com gemologia 10/10). model/fermat/{danielcunha.pdf, Wiles-Fermat.95.pdf}.","epistemico":"fato","exemplos":[],"id":"fermat_na_maquina","memoria":"semantica","meta":{"dominio":"joalheria","fonte":"maquina de corte"},"origem":"wiki","palavras_chave":[],"sinonimos":[],"tipo":"conceito","titulo":"O Último Teorema de Fermat (Kummer + Wiles) ingerido no corpo estelar"}
\section{O Último Teorema de Fermat (Kummer + Wiles) ingerido no corpo estelar}
Os dois caminhos da prova do TLF caem nos construtos que a máquina já opera. KUMMER (a monografia de Daniel Cunha, UERJ): os corpos CICLOTÔMICOS Q(\ensuremath{\zeta}), \ensuremath{\zeta}=e\^{}\{2\ensuremath{\pi}i/n\} — as RAÍZES DA UNIDADE = as cúspides da ESTRELA; e o TRAÇO/NORMA de um inteiro algébrico (Tr=\ensuremath{\Sigma}\ensuremath{\sigma}\_i, N=\ensuremath{\Pi}\ensuremath{\sigma}\_i, por Vieta Tr=\ensuremath{-}a\_1, N=(\ensuremath{-}1)\ensuremath{^{n}}a\_n) = os METAIS (\ensuremath{\sigma}\_n é raiz de z\ensuremath{^{2}}\ensuremath{-}nz\ensuremath{-}1, Tr=n, N=\ensuremath{-}1, uma UNIDADE); o caso n=4 fecha pela DESCIDA infinita = a órbita que contrai z (o cristal, só o ponto fixo trivial). WILES (o paper de 1995): a modularidade via representações de Galois \ensuremath{\rho}(Frob q), com traço=c(q,f) (o autovalor de HECKE = o espectro = o metal) e det=\ensuremath{\chi}(q)q\^{}\{k-1\} — o 2$\times$2 traço-det É o GATO A\_n (tr=n, det=\ensuremath{-}1); a modularidade (a\_p=c(p,f)) é o CASAMENTO ESPECTRAL, o mesmo gesto do corpo estelar (ler o rastro). RESSALVA: não reprova o TLF nem é a prova de Wiles — é a INGESTÃO (o dicionário exato de Peirce + as peças computáveis verificadas). linguagem/fermat.py (5/5); dicionarios.\_FERMAT\_MAQUINA (Fermat é uma ÁREA da máquina, compõe com gemologia 10/10). model/fermat/\{danielcunha.pdf, Wiles-Fermat.95.pdf\}.
\prov{relações: relaciona amendoa\_hidrogenio\_formal; parte\_de joalheria\_hub}
\prov{proveniência: wiki \textperiodcentered{} maquina de corte \textperiodcentered{} fato \textperiodcentered{} confiança 1.0 \textperiodcentered{} domínio joalheria \textperiodcentered{} história 27 versões no jornal}

%CRISTAL {"arestas":[{"alvo":"mecanica","confianca":1.0,"epistemico":"fato","origem":"manual","rel":"parte_de"},{"alvo":"termodinamica","confianca":1.0,"epistemico":"fato","origem":"manual","rel":"parte_de"},{"alvo":"geometria","confianca":1.0,"epistemico":"fato","origem":"manual","rel":"referencia"},{"alvo":"virtualizacao","confianca":0.5,"epistemico":"fato","origem":"wiki","rel":"relaciona"},{"alvo":"word","confianca":0.5,"epistemico":"fato","origem":"wiki","rel":"relaciona"},{"alvo":"vontade_de_poder","confianca":0.5,"epistemico":"fato","origem":"wiki","rel":"relaciona"},{"alvo":"while","confianca":0.5,"epistemico":"fato","origem":"wiki","rel":"relaciona"}],"confianca":1.0,"contraexemplos":[],"descricao":"Estudo de matéria, energia, espaço e tempo; mecânica, termodinâmica, geometria.","epistemico":"fato","exemplos":[],"id":"fisica","memoria":"semantica","meta":{"ciencia":"fisica","dominio":"fisica","verificacao":"slug idêntico — enriquecer, não duplicar"},"origem":"manual","palavras_chave":[],"sinonimos":[],"tipo":"conceito","titulo":"fisica"}
\section{fisica}
Estudo de matéria, energia, espaço e tempo; mecânica, termodinâmica, geometria.
\prov{relações: parte\_de mecanica; parte\_de termodinamica; referencia geometria; relaciona virtualizacao; relaciona word; relaciona vontade\_de\_poder; relaciona while}
\prov{proveniência: manual \textperiodcentered{} fato \textperiodcentered{} confiança 1.0 \textperiodcentered{} domínio fisica \textperiodcentered{} história 202 versões no jornal}

%CRISTAL {"arestas":[{"alvo":"energia_nuclear","confianca":1.0,"epistemico":"fato","origem":"wiki","rel":"eh_um"},{"alvo":"energia_de_ligacao","confianca":1.0,"epistemico":"fato","origem":"wiki","rel":"usa"},{"alvo":"vale_do_ferro","confianca":1.0,"epistemico":"fato","origem":"wiki","rel":"usa"}],"confianca":1.0,"contraexemplos":[],"descricao":"Um núcleo pesado (U-235, Pu-239) absorve um nêutron, fica instável e QUEBRA em dois fragmentos + 2-3 nêutrons + energia (~200 MeV). Os nêutrons novos podem quebrar outros núcleos: a reação em cadeia.","epistemico":"fato","exemplos":[],"id":"fissao_nuclear","memoria":"semantica","meta":{"dominio":"fisica","fonte":"energia nuclear"},"origem":"wiki","palavras_chave":[],"sinonimos":[],"tipo":"conceito","titulo":"Fissão Nuclear"}
\section{Fissão Nuclear}
Um núcleo pesado (U-235, Pu-239) absorve um nêutron, fica instável e QUEBRA em dois fragmentos + 2-3 nêutrons + energia (\~{}200 MeV). Os nêutrons novos podem quebrar outros núcleos: a reação em cadeia.
\prov{relações: eh\_um energia\_nuclear; usa energia\_de\_ligacao; usa vale\_do\_ferro}
\prov{proveniência: wiki \textperiodcentered{} energia nuclear \textperiodcentered{} fato \textperiodcentered{} confiança 1.0 \textperiodcentered{} domínio fisica \textperiodcentered{} história 4 versões no jornal}

%CRISTAL {"arestas":[{"alvo":"impressao_koch_32bytes","confianca":1.0,"epistemico":"fato","origem":"wiki","rel":"especializa"},{"alvo":"libtiffany_engine","confianca":1.0,"epistemico":"fato","origem":"wiki","rel":"parte_de"},{"alvo":"floco_phi","confianca":1.0,"epistemico":"fato","origem":"wiki","rel":"explica"},{"alvo":"koch_dissipa_nao_codifica","confianca":1.0,"epistemico":"fato","origem":"wiki","rel":"relaciona"},{"alvo":"primitivas_nucleo_associativo","confianca":1.0,"epistemico":"fato","origem":"wiki","rel":"relaciona"},{"alvo":"colapso_ping_pong_overlap","confianca":1.0,"epistemico":"fato","origem":"wiki","rel":"relaciona"},{"alvo":"colapso_multicamada_peso","confianca":1.0,"epistemico":"fato","origem":"wiki","rel":"usa"},{"alvo":"cristal_recall_ponte","confianca":1.0,"epistemico":"fato","origem":"wiki","rel":"relaciona"},{"alvo":"reconhece_c_memoria","confianca":1.0,"epistemico":"fato","origem":"wiki","rel":"relaciona"},{"alvo":"tiffany","confianca":1.0,"epistemico":"fato","origem":"wiki","rel":"relaciona"}],"confianca":1.0,"contraexemplos":[],"descricao":"A régua associativa (overlap/Hopfield) NÃO lê a informação não-associativa: o floco (impressão Koch z↦z² mod 1009, a ordem importa) é o não-associativo, e em Python o _koch roda milhões de vezes — não se contorna, só via metal. libtiffany.so:tiffany_impress computa o MESMO floco que constrói os .idx, ligado a prosa._impressao_bytes por ctypes com fallback puro-Python; bit-a-bit idêntico, ~24× (50k impressões 5,0s→0,21s). Complementos em Python sem tocar o motor congelado: _le_meta abre o tsv 1× por scan (não por nó); _ov_nó lazy (só re-hasheia o turno que casa). Verificado: reproduzir 62/62, e2e idêntico.","epistemico":"fato","exemplos":[],"id":"floco_no_metal","memoria":"semantica","meta":{"autor":"broca-so","dominio":"metal","fonte":"code/libtiffany/1_core/tiffany_base.c; code/prosa.py:_impressao_bytes/_lib"},"origem":"wiki","palavras_chave":[],"sinonimos":[],"tipo":"conceito","titulo":"O Floco no Metal — impressão não-associativa em C"}
\section{O Floco no Metal — impressão não-associativa em C}
A régua associativa (overlap/Hopfield) NÃO lê a informação não-associativa: o floco (impressão Koch z\ensuremath{\mapsto}z\ensuremath{^{2}} mod 1009, a ordem importa) é o não-associativo, e em Python o \_koch roda milhões de vezes — não se contorna, só via metal. libtiffany.so:tiffany\_impress computa o MESMO floco que constrói os .idx, ligado a prosa.\_impressao\_bytes por ctypes com fallback puro-Python; bit-a-bit idêntico, \~{}24$\times$ (50k impressões 5,0s\ensuremath{\to}0,21s). Complementos em Python sem tocar o motor congelado: \_le\_meta abre o tsv 1$\times$ por scan (não por nó); \_ov\_nó lazy (só re-hasheia o turno que casa). Verificado: reproduzir 62/62, e2e idêntico.
\prov{relações: especializa impressao\_koch\_32bytes; parte\_de libtiffany\_engine; explica floco\_phi; relaciona koch\_dissipa\_nao\_codifica; relaciona primitivas\_nucleo\_associativo; relaciona colapso\_ping\_pong\_overlap; usa colapso\_multicamada\_peso; relaciona cristal\_recall\_ponte; relaciona reconhece\_c\_memoria; relaciona tiffany}
\prov{proveniência: wiki \textperiodcentered{} code/libtiffany/1\_core/tiffany\_base.c; code/prosa.py:\_impressao\_bytes/\_lib \textperiodcentered{} fato \textperiodcentered{} confiança 1.0 \textperiodcentered{} domínio metal \textperiodcentered{} história 11 versões no jornal}

%CRISTAL {"arestas":[{"alvo":"grupo_z2cubo_boost","confianca":1.0,"epistemico":"fato","origem":"wiki","rel":"usa"},{"alvo":"corpo_f2k","confianca":1.0,"epistemico":"fato","origem":"wiki","rel":"relaciona"},{"alvo":"prtsuv_flux_qubit_predicoes","confianca":1.0,"epistemico":"fato","origem":"wiki","rel":"relaciona"},{"alvo":"dois_qubits_octonios","confianca":1.0,"epistemico":"fato","origem":"wiki","rel":"usa"},{"alvo":"majorana_qubit_topologico","confianca":1.0,"epistemico":"fato","origem":"wiki","rel":"relaciona"}],"confianca":1.0,"contraexemplos":[],"descricao":"O potencial V=S cos²ψ sobre o círculo de álgebras é a equação de Mathieu — um DUPLO-POÇO (poços nos quaternions s=±1, barreira no comutativo s=0); o tunelamento é o qubit fundamental, formalmente idêntico ao FLUX QUBIT supercondutor (RF-SQUID no ponto de frustração). Calibrando E_J,E_L,E_C, o espectro do circuito reproduz os splittings da tabela periódica algébrica; um arranjo 3D de três flux qubits realiza a órbita ℤ₂³ de Hurwitz (octupleto). [D]/[N] o espectro do duplo-poço (bate a 4×10⁻⁵); [J] a predição de hardware falseável; o mapeamento SGD↔espectro NÃO se verifica (lacuna honesta).","epistemico":"inferencia","exemplos":[],"id":"flux_qubit_tabela_periodica","memoria":"semantica","meta":{"autor":"broca-so","dominio":"cosmologia","fonte":"hiper/main.tex sec:periodica, sec:flux-qubit; patente_3_fisica.tex"},"origem":"wiki","palavras_chave":[],"sinonimos":[],"tipo":"conceito","titulo":"Flux qubit e a tabela periódica algébrica"}
\section{Flux qubit e a tabela periódica algébrica}
O potencial V=S cos\ensuremath{^{2}}\ensuremath{\psi} sobre o círculo de álgebras é a equação de Mathieu — um DUPLO-POÇO (poços nos quaternions s=$\pm$1, barreira no comutativo s=0); o tunelamento é o qubit fundamental, formalmente idêntico ao FLUX QUBIT supercondutor (RF-SQUID no ponto de frustração). Calibrando E\_J,E\_L,E\_C, o espectro do circuito reproduz os splittings da tabela periódica algébrica; um arranjo 3D de três flux qubits realiza a órbita \ensuremath{\mathbb{Z}}\ensuremath{_{2}}\ensuremath{^{3}} de Hurwitz (octupleto). [D]/[N] o espectro do duplo-poço (bate a 4$\times$10\ensuremath{^{-5}}); [J] a predição de hardware falseável; o mapeamento SGD\ensuremath{\leftrightarrow}espectro NÃO se verifica (lacuna honesta).
\prov{relações: usa grupo\_z2cubo\_boost; relaciona corpo\_f2k; relaciona prtsuv\_flux\_qubit\_predicoes; usa dois\_qubits\_octonios; relaciona majorana\_qubit\_topologico}
\prov{proveniência: wiki \textperiodcentered{} hiper/main.tex sec:periodica, sec:flux-qubit; patente\_3\_fisica.tex \textperiodcentered{} inferencia \textperiodcentered{} confiança 1.0 \textperiodcentered{} domínio cosmologia \textperiodcentered{} história 8 versões no jornal}

%CRISTAL {"arestas":[{"alvo":"wiles_selberg","confianca":1.0,"epistemico":"fato","origem":"wiki","rel":"especializa"},{"alvo":"joalheria_hub","confianca":1.0,"epistemico":"fato","origem":"wiki","rel":"parte_de"}],"confianca":1.0,"contraexemplos":[],"descricao":"O fecho do círculo Wiles↔Selberg, computável. A fórmula do traço de Selberg: as GEODÉSICAS determinam o ESPECTRO do laplaciano. A fórmula EXPLÍCITA de Weil é a irmã aritmética: os PRIMOS (os a_p, o Frobenius ~ geodésica) determinam os ZEROS de L(E,s) (o espectro). Em X₀(11), a partir SÓ dos primos: (1) os a_n vêm dos a_p por HECKE — a recursão que dá o PRODUTO DE EULER Σa_n n^{-s}=Π_p(1−a_p p^{-s}+p^{1-2s})^{-1} (primos↔inteiros, verificado); (2) a EQUAÇÃO FUNCIONAL Λ(s)=εΛ(2−s) (ε=+1) vem da MODULARIDADE (a involução de Fricke, a simetria de SL₂ hiperbólica) e põe os zeros na LINHA CRÍTICA Re(s)=1 (a borda); (3) o VALOR CENTRAL L(11a,1)=0.25384≠0 → posto analítico 0 (BSD, E(ℚ) finito); (4) o PRIMEIRO ZERO na linha crítica em s=1+6.36261i (o espectro produzido pelos primos), e Λ(1+it) é REAL (o espectro é real, como os autovalores do laplaciano de Selberg). CORREÇÃO capturada: o símbolo de Legendre vale só p/ p ímpar — a_2 pela contagem direta. RESSALVA: as peças são fatos/computações (Euler, eq. funcional, L(1), os zeros por mpmath); a leitura primo↔geodésica/zero↔autovalor é o dicionário estrutural da fórmula explícita, declarado. linguagem/formula_explicita.py (5/5): coeficientes, L_central, euler_vs_dirichlet, Lambda, L, primeiro_zero.","epistemico":"fato","exemplos":[],"id":"formula_explicita","memoria":"semantica","meta":{"dominio":"joalheria","fonte":"maquina de corte"},"origem":"wiki","palavras_chave":[],"sinonimos":[],"tipo":"conceito","titulo":"A fórmula explícita: os primos determinam os zeros de L(E,s), como o traço de Selberg"}
\section{A fórmula explícita: os primos determinam os zeros de L(E,s), como o traço de Selberg}
O fecho do círculo Wiles\ensuremath{\leftrightarrow}Selberg, computável. A fórmula do traço de Selberg: as GEODÉSICAS determinam o ESPECTRO do laplaciano. A fórmula EXPLÍCITA de Weil é a irmã aritmética: os PRIMOS (os a\_p, o Frobenius \~{} geodésica) determinam os ZEROS de L(E,s) (o espectro). Em X\ensuremath{_{0}}(11), a partir SÓ dos primos: (1) os a\_n vêm dos a\_p por HECKE — a recursão que dá o PRODUTO DE EULER \ensuremath{\Sigma}a\_n n\^{}\{-s\}=\ensuremath{\Pi}\_p(1\ensuremath{-}a\_p p\^{}\{-s\}+p\^{}\{1-2s\})\^{}\{-1\} (primos\ensuremath{\leftrightarrow}inteiros, verificado); (2) a EQUAÇÃO FUNCIONAL \ensuremath{\Lambda}(s)=\ensuremath{\varepsilon}\ensuremath{\Lambda}(2\ensuremath{-}s) (\ensuremath{\varepsilon}=+1) vem da MODULARIDADE (a involução de Fricke, a simetria de SL\ensuremath{_{2}} hiperbólica) e põe os zeros na LINHA CRÍTICA Re(s)=1 (a borda); (3) o VALOR CENTRAL L(11a,1)=0.25384\ensuremath{\ne}0 \ensuremath{\to} posto analítico 0 (BSD, E(\ensuremath{\mathbb{Q}}) finito); (4) o PRIMEIRO ZERO na linha crítica em s=1+6.36261i (o espectro produzido pelos primos), e \ensuremath{\Lambda}(1+it) é REAL (o espectro é real, como os autovalores do laplaciano de Selberg). CORREÇÃO capturada: o símbolo de Legendre vale só p/ p ímpar — a\_2 pela contagem direta. RESSALVA: as peças são fatos/computações (Euler, eq. funcional, L(1), os zeros por mpmath); a leitura primo\ensuremath{\leftrightarrow}geodésica/zero\ensuremath{\leftrightarrow}autovalor é o dicionário estrutural da fórmula explícita, declarado. linguagem/formula\_explicita.py (5/5): coeficientes, L\_central, euler\_vs\_dirichlet, Lambda, L, primeiro\_zero.
\prov{relações: especializa wiles\_selberg; parte\_de joalheria\_hub}
\prov{proveniência: wiki \textperiodcentered{} maquina de corte \textperiodcentered{} fato \textperiodcentered{} confiança 1.0 \textperiodcentered{} domínio joalheria \textperiodcentered{} história 18 versões no jornal}

%CRISTAL {"arestas":[{"alvo":"particulas_minerais","confianca":1.0,"epistemico":"fato","origem":"wiki","rel":"parte_de"},{"alvo":"numero_de_salem","confianca":1.0,"epistemico":"fato","origem":"wiki","rel":"usa"},{"alvo":"barreira_hurwitz_qubits","confianca":1.0,"epistemico":"fato","origem":"wiki","rel":"relaciona"}],"confianca":1.0,"contraexemplos":[],"descricao":"Um número de Salem tem conjugados NO CÍRCULO UNITÁRIO (ρ=1) ⟹ massa=|log 1|=0: são FÓTONS, modos sem massa. O de Lehmer (grau 10) é um átomo com núcleo + 1 elétron + 8 fótons — radiação pura na borda. Salem é a fronteira onde a matéria (Pisot, massa>0) vira luz (massa=0).","epistemico":"hipotese","exemplos":[],"id":"foton_salem","memoria":"semantica","meta":{"dominio":"fisica","espectro_lehmer":[{"massa":0.1624,"theta_g":0.0,"tipo":"nucleo"},{"massa":0.0,"theta_g":160.6,"tipo":"foton"},{"massa":0.0,"theta_g":160.6,"tipo":"foton"},{"massa":0.0,"theta_g":137.3,"tipo":"foton"},{"massa":0.0,"theta_g":137.3,"tipo":"foton"},{"massa":0.0,"theta_g":62.8,"tipo":"foton"},{"massa":0.0,"theta_g":62.8,"tipo":"foton"},{"massa":0.0,"theta_g":107.0,"tipo":"foton"},{"massa":0.0,"theta_g":107.0,"tipo":"foton"},{"massa":0.1624,"theta_g":0.0,"tipo":"eletron"}],"fonte":"partículas minerais"},"origem":"wiki","palavras_chave":[],"sinonimos":[],"tipo":"conceito","titulo":"Os Fótons de Salem (modos sem massa)"}
\section{Os Fótons de Salem (modos sem massa)}
Um número de Salem tem conjugados NO CÍRCULO UNITÁRIO (\ensuremath{\rho}=1) \ensuremath{\Longrightarrow} massa=|log 1|=0: são FÓTONS, modos sem massa. O de Lehmer (grau 10) é um átomo com núcleo + 1 elétron + 8 fótons — radiação pura na borda. Salem é a fronteira onde a matéria (Pisot, massa>0) vira luz (massa=0).
\prov{relações: parte\_de particulas\_minerais; usa numero\_de\_salem; relaciona barreira\_hurwitz\_qubits}
\prov{proveniência: wiki \textperiodcentered{} partículas minerais \textperiodcentered{} hipotese \textperiodcentered{} confiança 1.0 \textperiodcentered{} domínio fisica \textperiodcentered{} história 4 versões no jornal}

%CRISTAL {"arestas":[{"alvo":"fusao_nuclear","confianca":1.0,"epistemico":"fato","origem":"wiki","rel":"usa"}],"confianca":1.0,"contraexemplos":[],"descricao":"Reproduzir a fusão na Terra: confinar plasma D-T a ~10⁸ K por campo magnético (tokamak) ou compressão (ICF). O critério de Lawson (n·T·τ) mede o ganho; o ITER busca Q>1 (mais energia que a injetada). A energia limpa e densa ainda por domar.","epistemico":"fato","exemplos":[],"id":"fusao_controlada","memoria":"semantica","meta":{"dominio":"fisica","fonte":"energia nuclear"},"origem":"wiki","palavras_chave":[],"sinonimos":[],"tipo":"conceito","titulo":"Fusão Controlada (tokamak, ITER)"}
\section{Fusão Controlada (tokamak, ITER)}
Reproduzir a fusão na Terra: confinar plasma D-T a \~{}10\ensuremath{^{8}} K por campo magnético (tokamak) ou compressão (ICF). O critério de Lawson (n\textperiodcentered T\textperiodcentered \ensuremath{\tau}) mede o ganho; o ITER busca Q>1 (mais energia que a injetada). A energia limpa e densa ainda por domar.
\prov{relações: usa fusao\_nuclear}
\prov{proveniência: wiki \textperiodcentered{} energia nuclear \textperiodcentered{} fato \textperiodcentered{} confiança 1.0 \textperiodcentered{} domínio fisica \textperiodcentered{} história 2 versões no jornal}

%CRISTAL {"arestas":[{"alvo":"fusao_nuclear","confianca":1.0,"epistemico":"fato","origem":"wiki","rel":"eh_um"}],"confianca":1.0,"contraexemplos":[],"descricao":"As estrelas fundem hidrogênio em hélio (cadeia p-p, ciclo CNO) e, nas massivas, até o ferro — parando no vale do ferro. A gravidade é o confinamento; a nucleossíntese estelar forjou os elementos.","epistemico":"fato","exemplos":[],"id":"fusao_estelar","memoria":"semantica","meta":{"dominio":"fisica","fonte":"energia nuclear"},"origem":"wiki","palavras_chave":[],"sinonimos":[],"tipo":"conceito","titulo":"Fusão Estelar"}
\section{Fusão Estelar}
As estrelas fundem hidrogênio em hélio (cadeia p-p, ciclo CNO) e, nas massivas, até o ferro — parando no vale do ferro. A gravidade é o confinamento; a nucleossíntese estelar forjou os elementos.
\prov{relações: eh\_um fusao\_nuclear}
\prov{proveniência: wiki \textperiodcentered{} energia nuclear \textperiodcentered{} fato \textperiodcentered{} confiança 1.0 \textperiodcentered{} domínio fisica \textperiodcentered{} história 2 versões no jornal}

%CRISTAL {"arestas":[{"alvo":"energia_nuclear","confianca":1.0,"epistemico":"fato","origem":"wiki","rel":"eh_um"},{"alvo":"energia_de_ligacao","confianca":1.0,"epistemico":"fato","origem":"wiki","rel":"usa"},{"alvo":"vale_do_ferro","confianca":1.0,"epistemico":"fato","origem":"wiki","rel":"usa"}],"confianca":1.0,"contraexemplos":[],"descricao":"Dois núcleos leves (D+T, D+D) se JUNTAM num mais pesado + energia. Precisa vencer a repulsão de Coulomb (temperatura/pressão altíssimas); rende mais energia por massa que a fissão e sem rejeito de longa vida. É o que alimenta as estrelas.","epistemico":"fato","exemplos":[],"id":"fusao_nuclear","memoria":"semantica","meta":{"dominio":"fisica","fonte":"energia nuclear"},"origem":"wiki","palavras_chave":[],"sinonimos":[],"tipo":"conceito","titulo":"Fusão Nuclear"}
\section{Fusão Nuclear}
Dois núcleos leves (D+T, D+D) se JUNTAM num mais pesado + energia. Precisa vencer a repulsão de Coulomb (temperatura/pressão altíssimas); rende mais energia por massa que a fissão e sem rejeito de longa vida. É o que alimenta as estrelas.
\prov{relações: eh\_um energia\_nuclear; usa energia\_de\_ligacao; usa vale\_do\_ferro}
\prov{proveniência: wiki \textperiodcentered{} energia nuclear \textperiodcentered{} fato \textperiodcentered{} confiança 1.0 \textperiodcentered{} domínio fisica \textperiodcentered{} história 4 versões no jornal}

%CRISTAL {"arestas":[{"alvo":"g2_grupo_excepcional","confianca":1.0,"epistemico":"fato","origem":"wiki","rel":"especializa"},{"alvo":"dois_qubits_octonios","confianca":1.0,"epistemico":"fato","origem":"wiki","rel":"explica"}],"confianca":1.0,"contraexemplos":[],"descricao":"G₂=Aut(𝕆) preserva o produto octônico (= a medida de emaranhamento). O associador é a torção de S⁷ (J=−3/2[x,y,z]), representado pelo tensor de Schafer φ (dual de Hodge das constantes de estrutura). Números NUS, sem input: ‖c‖²=42 (7 retas de Fano), ‖φ‖²=168=4·42=|PSL(2,7)| (simetria da quártica de Klein). β_𝕆=0,0111.","epistemico":"fato","exemplos":[],"id":"g2_automorfismos_associador","memoria":"semantica","meta":{"autor":"broca-so","dominio":"cosmologia","fonte":"paper_3_quantica.tex §torcao; paper_CB"},"origem":"wiki","palavras_chave":[],"sinonimos":[],"tipo":"conceito","titulo":"G₂ = Aut(𝕆): o associador e os números nus (42, 168)"}
\section{G\ensuremath{_{2}} = Aut(\ensuremath{\mathbb{O}}): o associador e os números nus (42, 168)}
G\ensuremath{_{2}}=Aut(\ensuremath{\mathbb{O}}) preserva o produto octônico (= a medida de emaranhamento). O associador é a torção de S\ensuremath{^{7}} (J=\ensuremath{-}3/2[x,y,z]), representado pelo tensor de Schafer \ensuremath{\varphi} (dual de Hodge das constantes de estrutura). Números NUS, sem input: \ensuremath{\|}c\ensuremath{\|}\ensuremath{^{2}}=42 (7 retas de Fano), \ensuremath{\|}\ensuremath{\varphi}\ensuremath{\|}\ensuremath{^{2}}=168=4\textperiodcentered 42=|PSL(2,7)| (simetria da quártica de Klein). \ensuremath{\beta}\_\ensuremath{\mathbb{O}}=0,0111.
\prov{relações: especializa g2\_grupo\_excepcional; explica dois\_qubits\_octonios}
\prov{proveniência: wiki \textperiodcentered{} paper\_3\_quantica.tex §torcao; paper\_CB \textperiodcentered{} fato \textperiodcentered{} confiança 1.0 \textperiodcentered{} domínio cosmologia \textperiodcentered{} história 3 versões no jornal}

%CRISTAL {"arestas":[{"alvo":"materia_modelo_padrao","confianca":1.0,"epistemico":"fato","origem":"wiki","rel":"relaciona"}],"confianca":1.0,"contraexemplos":[],"descricao":"O menor dos 5 grupos excepcionais; é o grupo de automorfismos dos octônios 𝕆 (o subgrupo de SO(7) que fixa um spinor). Dimensão 14, posto 2, fundamentais 7 e 14, número de Coxeter h=6, dual h∨=4. Fato matemático (Cartan 1908-14); a relevância física direta (octônios no MP) é especulativa/aberta.","epistemico":"fato","exemplos":[],"id":"g2_grupo_excepcional","memoria":"semantica","meta":{"autor":"broca-so","dominio":"fisica","fonte":"Cartan 1908-14; G2 (mathematics)"},"origem":"wiki","palavras_chave":[],"sinonimos":[],"tipo":"conceito","titulo":"O grupo de Lie excepcional G₂ = Aut(𝕆)"}
\section{O grupo de Lie excepcional G\ensuremath{_{2}} = Aut(\ensuremath{\mathbb{O}})}
O menor dos 5 grupos excepcionais; é o grupo de automorfismos dos octônios \ensuremath{\mathbb{O}} (o subgrupo de SO(7) que fixa um spinor). Dimensão 14, posto 2, fundamentais 7 e 14, número de Coxeter h=6, dual h\ensuremath{\vee}=4. Fato matemático (Cartan 1908-14); a relevância física direta (octônios no MP) é especulativa/aberta.
\prov{relações: relaciona materia\_modelo\_padrao}
\prov{proveniência: wiki \textperiodcentered{} Cartan 1908-14; G2 (mathematics) \textperiodcentered{} fato \textperiodcentered{} confiança 1.0 \textperiodcentered{} domínio fisica \textperiodcentered{} história 2 versões no jornal}

%CRISTAL {"arestas":[{"alvo":"materia_modelo_padrao","confianca":1.0,"epistemico":"fato","origem":"wiki","rel":"confirma"}],"confianca":1.0,"contraexemplos":[],"descricao":"O momento magnético anômalo a_μ=(g−2)/2. Por ~20 anos a medição excedia o MP; o Fermilab (2021/2023) elevou a tensão a ~5σ vs a previsão WP20 (2020). REVIRAVOLTA 2025: a WP25 (Muon g-2 Theory Initiative) trocou a polarização de vácuo hadrônica (HVP) data-driven (e⁺e⁻) por LATTICE (consenso BMW), subindo a previsão ~223×10⁻¹¹; teoria=116592033±62 vs experimento 116592070,5 — CONCORDAM a ~1σ. A anomalia como sinal BSM ESSENCIALMENTE DESAPARECEU (virou tensão interna lattice×e⁺e⁻, agravada pelo CMD-3). FATO — em fluxo.","epistemico":"fato","exemplos":[],"id":"g2_muon_anomalia_dissolvida","memoria":"semantica","meta":{"autor":"broca-so","dominio":"fisica","fonte":"Fermilab g-2 2025 arXiv:2506; WP25 muon-gm2-theory.illinois.edu; CERN Courier"},"origem":"wiki","palavras_chave":[],"sinonimos":[],"tipo":"conceito","titulo":"g-2 do múon: a anomalia dissolvida (2025)"}
\section{g-2 do múon: a anomalia dissolvida (2025)}
O momento magnético anômalo a\_\ensuremath{\mu}=(g\ensuremath{-}2)/2. Por \~{}20 anos a medição excedia o MP; o Fermilab (2021/2023) elevou a tensão a \~{}5\ensuremath{\sigma} vs a previsão WP20 (2020). REVIRAVOLTA 2025: a WP25 (Muon g-2 Theory Initiative) trocou a polarização de vácuo hadrônica (HVP) data-driven (e\ensuremath{^{+}}e\ensuremath{^{-}}) por LATTICE (consenso BMW), subindo a previsão \~{}223$\times$10\ensuremath{^{-11}}; teoria=116592033$\pm$62 vs experimento 116592070,5 — CONCORDAM a \~{}1\ensuremath{\sigma}. A anomalia como sinal BSM ESSENCIALMENTE DESAPARECEU (virou tensão interna lattice$\times$e\ensuremath{^{+}}e\ensuremath{^{-}}, agravada pelo CMD-3). FATO — em fluxo.
\prov{relações: confirma materia\_modelo\_padrao}
\prov{proveniência: wiki \textperiodcentered{} Fermilab g-2 2025 arXiv:2506; WP25 muon-gm2-theory.illinois.edu; CERN Courier \textperiodcentered{} fato \textperiodcentered{} confiança 1.0 \textperiodcentered{} domínio fisica \textperiodcentered{} história 2 versões no jornal}

%CRISTAL {"arestas":[{"alvo":"difracao_atomica_turmalina","confianca":1.0,"epistemico":"fato","origem":"wiki","rel":"usa"},{"alvo":"equacao_gema_estelar","confianca":1.0,"epistemico":"fato","origem":"wiki","rel":"usa"},{"alvo":"primitivas_analiticas","confianca":1.0,"epistemico":"fato","origem":"wiki","rel":"usa"},{"alvo":"turmalina_paraiba_ancora","confianca":1.0,"epistemico":"fato","origem":"wiki","rel":"usa"},{"alvo":"joalheria_hub","confianca":1.0,"epistemico":"fato","origem":"wiki","rel":"parte_de"}],"confianca":1.0,"contraexemplos":[],"descricao":"As DUAS METADES unidas: a BORDA analítica (o dupolinômio da equação 𝒢, o trilhão esf+parab) preenchida pela CADEIA cristalina hexagonal da turmalina (elbaíta trigonal), CRESCIDA até 1 QUILATE. O bloco de construção: 6 TETRAEDROS → 1 PARALELEPÍPEDO (a triangulação de Kuhn, uma por permutação de (a1,a2,a3); os 6 SiO₄ do anel Si₆O₁₈ são esses 6 tetraedros). Empilha-se até o peso EXATO: massa=N_celas·ρ·V_cela=0,2 g (ρ=3,06 g/cm³), o que fixa a ESCALA (2,83 mm/unidade), as DIMENSÕES (5,34×5,67 mm), o VOLUME (65,36 mm³) e o NÚMERO de células (4,24×10¹⁹ paralelepípedos = 2,54×10²⁰ tetraedros). As duas metades casam: N_celas·V_cela = volume da gema (a rede preenche a borda). Um ponto está DENTRO sse ρ_xy≤ℓ(1+a·cos kθ)P(z) — o mesmo F<0 do dupolinômio. A seguir: a óptica (borda Snell/TIR/ricochetes + interior difratante) rumo ao ultra-realismo. linguagem/gema_cristal.py (9/9): gema_um_quilate, paralelepipedo_6_tetraedros, empilhar_ate_quilate, dentro_da_gema, celulas_na_gema.","epistemico":"fato","exemplos":[],"id":"gema_cristal_um_quilate","memoria":"semantica","meta":{"dominio":"joalheria","fonte":"maquina de corte"},"origem":"wiki","palavras_chave":[],"sinonimos":[],"tipo":"conceito","titulo":"A gema-cristal: a borda analítica preenchida pela rede da turmalina, 1 quilate"}
\section{A gema-cristal: a borda analítica preenchida pela rede da turmalina, 1 quilate}
As DUAS METADES unidas: a BORDA analítica (o dupolinômio da equação \ensuremath{\mathcal{G}}, o trilhão esf+parab) preenchida pela CADEIA cristalina hexagonal da turmalina (elbaíta trigonal), CRESCIDA até 1 QUILATE. O bloco de construção: 6 TETRAEDROS \ensuremath{\to} 1 PARALELEPÍPEDO (a triangulação de Kuhn, uma por permutação de (a1,a2,a3); os 6 SiO\ensuremath{_{4}} do anel Si\ensuremath{_{6}}O\ensuremath{_{18}} são esses 6 tetraedros). Empilha-se até o peso EXATO: massa=N\_celas\textperiodcentered \ensuremath{\rho}\textperiodcentered V\_cela=0,2 g (\ensuremath{\rho}=3,06 g/cm\ensuremath{^{3}}), o que fixa a ESCALA (2,83 mm/unidade), as DIMENSÕES (5,34$\times$5,67 mm), o VOLUME (65,36 mm\ensuremath{^{3}}) e o NÚMERO de células (4,24$\times$10\ensuremath{^{19}} paralelepípedos = 2,54$\times$10\ensuremath{^{20}} tetraedros). As duas metades casam: N\_celas\textperiodcentered V\_cela = volume da gema (a rede preenche a borda). Um ponto está DENTRO sse \ensuremath{\rho}\_xy\ensuremath{\le}\ensuremath{\ell}(1+a\textperiodcentered cos k\ensuremath{\theta})P(z) — o mesmo F<0 do dupolinômio. A seguir: a óptica (borda Snell/TIR/ricochetes + interior difratante) rumo ao ultra-realismo. linguagem/gema\_cristal.py (9/9): gema\_um\_quilate, paralelepipedo\_6\_tetraedros, empilhar\_ate\_quilate, dentro\_da\_gema, celulas\_na\_gema.
\prov{relações: usa difracao\_atomica\_turmalina; usa equacao\_gema\_estelar; usa primitivas\_analiticas; usa turmalina\_paraiba\_ancora; parte\_de joalheria\_hub}
\prov{proveniência: wiki \textperiodcentered{} maquina de corte \textperiodcentered{} fato \textperiodcentered{} confiança 1.0 \textperiodcentered{} domínio joalheria \textperiodcentered{} história 162 versões no jornal}

%CRISTAL {"arestas":[{"alvo":"corpo_estelar","confianca":1.0,"epistemico":"fato","origem":"wiki","rel":"usa"},{"alvo":"corpo_estelar_le_o_rastro","confianca":1.0,"epistemico":"fato","origem":"wiki","rel":"parte_de"}],"confianca":1.0,"contraexemplos":[],"descricao":"A GEOMETRIA do corpo estelar como ferramenta do pipe do metal: ρ=σ₂/σ₁ da nuvem de modos (o SVD) — a abertura da nuvem — classifica o rastro na PARÁBOLA mineral (cristal ρ≈0 / borda / caos). É o mesmo ρ de corpo_estelar.classe. No gabarito: ρ=0.205 → caos. linguagem/joalheria.geometria.","epistemico":"fato","exemplos":[],"id":"geometria_estelar_le_o_rastro","memoria":"semantica","meta":{"dominio":"joalheria","fonte":"espectro orbitas"},"origem":"wiki","palavras_chave":[],"sinonimos":[],"tipo":"conceito","titulo":"Geometria Estelar — ρ=σ₂/σ₁ lê a forma do rastro (a parábola)"}
\section{Geometria Estelar — \ensuremath{\rho}=\ensuremath{\sigma}\ensuremath{_{2}}/\ensuremath{\sigma}\ensuremath{_{1}} lê a forma do rastro (a parábola)}
A GEOMETRIA do corpo estelar como ferramenta do pipe do metal: \ensuremath{\rho}=\ensuremath{\sigma}\ensuremath{_{2}}/\ensuremath{\sigma}\ensuremath{_{1}} da nuvem de modos (o SVD) — a abertura da nuvem — classifica o rastro na PARÁBOLA mineral (cristal \ensuremath{\rho}\ensuremath{\approx}0 / borda / caos). É o mesmo \ensuremath{\rho} de corpo\_estelar.classe. No gabarito: \ensuremath{\rho}=0.205 \ensuremath{\to} caos. linguagem/joalheria.geometria.
\prov{relações: usa corpo\_estelar; parte\_de corpo\_estelar\_le\_o\_rastro}
\prov{proveniência: wiki \textperiodcentered{} espectro orbitas \textperiodcentered{} fato \textperiodcentered{} confiança 1.0 \textperiodcentered{} domínio joalheria \textperiodcentered{} história 3 versões no jornal}

%CRISTAL {"arestas":[{"alvo":"tabela_geometrias_espectrais","confianca":1.0,"epistemico":"fato","origem":"wiki","rel":"parte_de"}],"confianca":1.0,"contraexemplos":[],"descricao":"Uma geometria de corte da tabela: simetria k=4 (grupo 4mm, ν=3); assinatura r=1+k·cos(4θ); curvatura 0.75. Os campos de material/óptica/certificação/render crescem quando uma gema desse corte for lapidada. Placeholder vivo — a tabela cresce no grafo.","epistemico":"fato","exemplos":[],"id":"geometria_quadrado","memoria":"semantica","meta":{"dominio":"joalheria","fonte":"maquina de corte"},"origem":"wiki","palavras_chave":[],"sinonimos":[],"tipo":"conceito","titulo":"Geometria: Quadrado (k=4) — linha da tabela (a crescer)"}
\section{Geometria: Quadrado (k=4) — linha da tabela (a crescer)}
Uma geometria de corte da tabela: simetria k=4 (grupo 4mm, \ensuremath{\nu}=3); assinatura r=1+k\textperiodcentered cos(4\ensuremath{\theta}); curvatura 0.75. Os campos de material/óptica/certificação/render crescem quando uma gema desse corte for lapidada. Placeholder vivo — a tabela cresce no grafo.
\prov{relações: parte\_de tabela\_geometrias\_espectrais}
\prov{proveniência: wiki \textperiodcentered{} maquina de corte \textperiodcentered{} fato \textperiodcentered{} confiança 1.0 \textperiodcentered{} domínio joalheria \textperiodcentered{} história 102 versões no jornal}

%CRISTAL {"arestas":[{"alvo":"tabela_geometrias_espectrais","confianca":1.0,"epistemico":"fato","origem":"wiki","rel":"parte_de"},{"alvo":"turmalina_paraiba_ancora","confianca":1.0,"epistemico":"fato","origem":"wiki","rel":"usa"},{"alvo":"pedra_logo_3d","confianca":1.0,"epistemico":"fato","origem":"wiki","rel":"relaciona"}],"confianca":1.0,"contraexemplos":[],"descricao":"A geometria de corte da nossa pedra, a primeira linha COMPLETA da tabela de geometrias espectrais. Simetria k=3 (grupo 3m trigonal, ν=2); assinatura r=1+k·cos(3θ); curvatura média 0.667; operador: abertura (a poda da borda); lapidação: brilhante 3-fold (mesa/coroa/rondiz/pavilhão); óptica: n=1.62, dispersão 0.018–0.020, janela 490–560 nm, cromóforo Cu²⁺; certificação: 1 ct = R$ 1.000,00, chave 3-qubit; render ideal: 6 ricochetes, softbox branca, ACES. Material: Turmalina Paraíba (elbaíta cuprífera).","epistemico":"fato","exemplos":[],"id":"geometria_trilhao","memoria":"semantica","meta":{"dominio":"joalheria","fonte":"maquina de corte"},"origem":"wiki","palavras_chave":[],"sinonimos":[],"tipo":"conceito","titulo":"Geometria: Trilhão (k=3) — a Turmalina Paraíba (a linha viva da tabela)"}
\section{Geometria: Trilhão (k=3) — a Turmalina Paraíba (a linha viva da tabela)}
A geometria de corte da nossa pedra, a primeira linha COMPLETA da tabela de geometrias espectrais. Simetria k=3 (grupo 3m trigonal, \ensuremath{\nu}=2); assinatura r=1+k\textperiodcentered cos(3\ensuremath{\theta}); curvatura média 0.667; operador: abertura (a poda da borda); lapidação: brilhante 3-fold (mesa/coroa/rondiz/pavilhão); óptica: n=1.62, dispersão 0.018–0.020, janela 490–560 nm, cromóforo Cu\ensuremath{^{2+}}; certificação: 1 ct = R\$ 1.000,00, chave 3-qubit; render ideal: 6 ricochetes, softbox branca, ACES. Material: Turmalina Paraíba (elbaíta cuprífera).
\prov{relações: parte\_de tabela\_geometrias\_espectrais; usa turmalina\_paraiba\_ancora; relaciona pedra\_logo\_3d}
\prov{proveniência: wiki \textperiodcentered{} maquina de corte \textperiodcentered{} fato \textperiodcentered{} confiança 1.0 \textperiodcentered{} domínio joalheria \textperiodcentered{} história 204 versões no jornal}

%CRISTAL {"arestas":[{"alvo":"pow_classico_bitcoin","confianca":1.0,"epistemico":"fato","origem":"wiki","rel":"confirma"},{"alvo":"cofre_parabola_hero","confianca":1.0,"epistemico":"fato","origem":"wiki","rel":"usa"},{"alvo":"parabola_audit_dual","confianca":1.0,"epistemico":"fato","origem":"wiki","rel":"usa"},{"alvo":"symbolic_beta_protocolo","confianca":1.0,"epistemico":"fato","origem":"wiki","rel":"confirma"},{"alvo":"transformada_metalica_fechada_perelman","confianca":1.0,"epistemico":"fato","origem":"wiki","rel":"confirma"},{"alvo":"cristalchain","confianca":1.0,"epistemico":"fato","origem":"wiki","rel":"relaciona"}],"confianca":1.0,"contraexemplos":[],"descricao":"Verificado em gex.txt (2 jul 2026): o pipeline stratum na GEX minera um BLOCO BTC REAL — job 6a27b63200011379, version=0x20000000 (BIP9 version-bits padrão), nbits=0x17021a42 → target≈2,0e53, dificuldade ≈1,3e14× o difficulty-1 (escala mainnet), ~3,9 MH/s, 540 bi hashes. Opera em c²=0,9900 → a=0,0100 (parábola a+c²=1): é OURO real, e sua margem (o lado negro dissipativo 'a') está quase esgotada — por isso 'a mineração do ouro acabou' e a entropia tem de ser importada e transduzida pra prata. [D] cabeçalho decodificado.","epistemico":"fato","exemplos":[],"id":"gex_minera_ouro_real","memoria":"semantica","meta":{"autor":"Lopes/broca-so","dominio":"mineracao","fonte":"gex.txt (job 6a27b63200011379, nbits 0x17021a42); decode read-only"},"origem":"wiki","palavras_chave":[],"sinonimos":[],"tipo":"conceito","titulo":"GEX minera OURO real: o pipeline stratum roda um bloco BTC de escala mainnet (c²=0,99)"}
\section{GEX minera OURO real: o pipeline stratum roda um bloco BTC de escala mainnet (c\ensuremath{^{2}}=0,99)}
Verificado em gex.txt (2 jul 2026): o pipeline stratum na GEX minera um BLOCO BTC REAL — job 6a27b63200011379, version=0x20000000 (BIP9 version-bits padrão), nbits=0x17021a42 \ensuremath{\to} target\ensuremath{\approx}2,0e53, dificuldade \ensuremath{\approx}1,3e14$\times$ o difficulty-1 (escala mainnet), \~{}3,9 MH/s, 540 bi hashes. Opera em c\ensuremath{^{2}}=0,9900 \ensuremath{\to} a=0,0100 (parábola a+c\ensuremath{^{2}}=1): é OURO real, e sua margem (o lado negro dissipativo 'a') está quase esgotada — por isso 'a mineração do ouro acabou' e a entropia tem de ser importada e transduzida pra prata. [D] cabeçalho decodificado.
\prov{relações: confirma pow\_classico\_bitcoin; usa cofre\_parabola\_hero; usa parabola\_audit\_dual; confirma symbolic\_beta\_protocolo; confirma transformada\_metalica\_fechada\_perelman; relaciona cristalchain}
\prov{proveniência: wiki \textperiodcentered{} gex.txt (job 6a27b63200011379, nbits 0x17021a42); decode read-only \textperiodcentered{} fato \textperiodcentered{} confiança 1.0 \textperiodcentered{} domínio mineracao \textperiodcentered{} história 27 versões no jornal}

%CRISTAL {"arestas":[{"alvo":"parabola_controla_taxa_tecidos","confianca":1.0,"epistemico":"fato","origem":"wiki","rel":"confirma"},{"alvo":"symbolic_beta_pipe_so","confianca":1.0,"epistemico":"fato","origem":"wiki","rel":"usa"},{"alvo":"hash_pell_inverso_parabola","confianca":1.0,"epistemico":"fato","origem":"wiki","rel":"usa"},{"alvo":"pell_algebra_posicao","confianca":1.0,"epistemico":"fato","origem":"wiki","rel":"usa"}],"confianca":1.0,"contraexemplos":[],"descricao":"A GEX (hipercomplexo, 20 cores) roda o SymbolicBeta live com 12 LANES ativas (símbolic.py live --lanes 12): quando chega trabalho (um floco de ouro), as 12 lanes/tecidos executam e cada uma cunha um Pell — 12 Pell por floco (a taxa de Pell subiu de fato). SEM TETO de recurso (BROCA_DIFUSAO_FRAC=1): os recursos se auto-regulam por trabalho. As lanes NÃO recomputam espectro (SVD) — usam c²/ν que já vêm da parábola (canal_c2/modos do seed), o que restaurou a hashrate e destravou o olho. O olho (neuronio/olho_gex44.py) mostra as duas leituras VINDAS DA PARÁBOLA: c²=canal_c2 (ouro) e a=1−c² (prata), + a cunhagem direta (Pell) e as lanes. Config em goldchain/dados/producao_multifractal.json. [D] no ar (broca-stratum + broca-olho).","epistemico":"fato","exemplos":[],"id":"gex_producao_multifractal","memoria":"semantica","meta":{"autor":"Lopes/broca-so","dominio":"mineracao","fonte":"gex44: symbolic.py live --lanes 12, frac=1; olho_gex44.py (canal_c2); producao_multifractal.json"},"origem":"wiki","palavras_chave":[],"sinonimos":[],"tipo":"conceito","titulo":"GEX em produção: 12 lanes por floco, sem teto (frac=1), olho lê a parábola"}
\section{GEX em produção: 12 lanes por floco, sem teto (frac=1), olho lê a parábola}
A GEX (hipercomplexo, 20 cores) roda o SymbolicBeta live com 12 LANES ativas (símbolic.py live --lanes 12): quando chega trabalho (um floco de ouro), as 12 lanes/tecidos executam e cada uma cunha um Pell — 12 Pell por floco (a taxa de Pell subiu de fato). SEM TETO de recurso (BROCA\_DIFUSAO\_FRAC=1): os recursos se auto-regulam por trabalho. As lanes NÃO recomputam espectro (SVD) — usam c\ensuremath{^{2}}/\ensuremath{\nu} que já vêm da parábola (canal\_c2/modos do seed), o que restaurou a hashrate e destravou o olho. O olho (neuronio/olho\_gex44.py) mostra as duas leituras VINDAS DA PARÁBOLA: c\ensuremath{^{2}}=canal\_c2 (ouro) e a=1\ensuremath{-}c\ensuremath{^{2}} (prata), + a cunhagem direta (Pell) e as lanes. Config em goldchain/dados/producao\_multifractal.json. [D] no ar (broca-stratum + broca-olho).
\prov{relações: confirma parabola\_controla\_taxa\_tecidos; usa symbolic\_beta\_pipe\_so; usa hash\_pell\_inverso\_parabola; usa pell\_algebra\_posicao}
\prov{proveniência: wiki \textperiodcentered{} gex44: symbolic.py live --lanes 12, frac=1; olho\_gex44.py (canal\_c2); producao\_multifractal.json \textperiodcentered{} fato \textperiodcentered{} confiança 1.0 \textperiodcentered{} domínio mineracao \textperiodcentered{} história 16 versões no jornal}

%CRISTAL {"arestas":[{"alvo":"mass_gap_algebrico","confianca":1.0,"epistemico":"fato","origem":"wiki","rel":"especializa"},{"alvo":"glueballs_lattice","confianca":1.0,"epistemico":"fato","origem":"wiki","rel":"confirma"},{"alvo":"g2_grupo_excepcional","confianca":1.0,"epistemico":"fato","origem":"wiki","rel":"usa"},{"alvo":"o_mass_gap","confianca":1.0,"epistemico":"fato","origem":"wiki","rel":"relaciona"},{"alvo":"energia_salto_hbar_raiz2","confianca":1.0,"epistemico":"fato","origem":"wiki","rel":"relaciona"}],"confianca":1.0,"contraexemplos":[],"descricao":"m₀₊₊=√2·ξ·√σ≈1,6 GeV. O √2 é DERIVADO (Casimir de G₂, C₂(adj)/C₂(fund)=2 exato — Casimir scaling, lattice 1010.5523; G₂ realiza o limite planar N→∞ de SU(N)). A tensão de corda √σ≈440 MeV e o fator ξ≈2,5 são POSTOS (input de lattice). Concordância de FORMA (√2), não de escala.","epistemico":"inferencia","exemplos":[],"id":"glueball_1_6_gev","memoria":"semantica","meta":{"autor":"broca-so","dominio":"cosmologia","fonte":"paper_4_particulas.tex §massgap (escala_massgap.py)"},"origem":"wiki","palavras_chave":[],"sinonimos":[],"tipo":"conceito","titulo":"Glueball ≈1,6 GeV (√2 firme, escala ξ posta)"}
\section{Glueball \ensuremath{\approx}1,6 GeV (\ensuremath{\sqrt{\,}}2 firme, escala \ensuremath{\xi} posta)}
m\ensuremath{_{0++}}=\ensuremath{\sqrt{\,}}2\textperiodcentered \ensuremath{\xi}\textperiodcentered \ensuremath{\sqrt{\,}}\ensuremath{\sigma}\ensuremath{\approx}1,6 GeV. O \ensuremath{\sqrt{\,}}2 é DERIVADO (Casimir de G\ensuremath{_{2}}, C\ensuremath{_{2}}(adj)/C\ensuremath{_{2}}(fund)=2 exato — Casimir scaling, lattice 1010.5523; G\ensuremath{_{2}} realiza o limite planar N\ensuremath{\to}\ensuremath{\infty} de SU(N)). A tensão de corda \ensuremath{\sqrt{\,}}\ensuremath{\sigma}\ensuremath{\approx}440 MeV e o fator \ensuremath{\xi}\ensuremath{\approx}2,5 são POSTOS (input de lattice). Concordância de FORMA (\ensuremath{\sqrt{\,}}2), não de escala.
\prov{relações: especializa mass\_gap\_algebrico; confirma glueballs\_lattice; usa g2\_grupo\_excepcional; relaciona o\_mass\_gap; relaciona energia\_salto\_hbar\_raiz2}
\prov{proveniência: wiki \textperiodcentered{} paper\_4\_particulas.tex §massgap (escala\_massgap.py) \textperiodcentered{} inferencia \textperiodcentered{} confiança 1.0 \textperiodcentered{} domínio cosmologia \textperiodcentered{} história 6 versões no jornal}

%CRISTAL {"arestas":[{"alvo":"qcd_cromodinamica","confianca":1.0,"epistemico":"fato","origem":"wiki","rel":"parte_de"},{"alvo":"yang_mills_mass_gap_milenio","confianca":1.0,"epistemico":"fato","origem":"wiki","rel":"relaciona"}],"confianca":1.0,"contraexemplos":[],"descricao":"Estados ligados APENAS de glúons (permitidos pela auto-interação). O mais leve é escalar 0⁺⁺, previsto em lattice QCD quenched ≈1,5–1,7 GeV (Morningstar-Peardon 1999: 1,73(5)(8) GeV; Athenodorou-Teper 2020: ≈1,65). Candidatos: f₀(1500), f₀(1710). Robusto em lattice; ABERTO: nenhum glueball identificado sem ambiguidade (mistura com mésons qq̄).","epistemico":"fato","exemplos":[],"id":"glueballs_lattice","memoria":"semantica","meta":{"autor":"broca-so","dominio":"fisica","fonte":"Morningstar-Peardon 1999 PRD 60,034509 hep-lat/9901004"},"origem":"wiki","palavras_chave":[],"sinonimos":[],"tipo":"conceito","titulo":"Glueballs (estados só de glúons)"}
\section{Glueballs (estados só de glúons)}
Estados ligados APENAS de glúons (permitidos pela auto-interação). O mais leve é escalar 0\ensuremath{^{++}}, previsto em lattice QCD quenched \ensuremath{\approx}1,5–1,7 GeV (Morningstar-Peardon 1999: 1,73(5)(8) GeV; Athenodorou-Teper 2020: \ensuremath{\approx}1,65). Candidatos: f\ensuremath{_{0}}(1500), f\ensuremath{_{0}}(1710). Robusto em lattice; ABERTO: nenhum glueball identificado sem ambiguidade (mistura com mésons qq\ensuremath{}).
\prov{relações: parte\_de qcd\_cromodinamica; relaciona yang\_mills\_mass\_gap\_milenio}
\prov{proveniência: wiki \textperiodcentered{} Morningstar-Peardon 1999 PRD 60,034509 hep-lat/9901004 \textperiodcentered{} fato \textperiodcentered{} confiança 1.0 \textperiodcentered{} domínio fisica \textperiodcentered{} história 3 versões no jornal}

%CRISTAL {"arestas":[{"alvo":"gpu_gato_ptx","confianca":1.0,"epistemico":"fato","origem":"wiki","rel":"usa"},{"alvo":"libbroca_so_camada","confianca":1.0,"epistemico":"fato","origem":"wiki","rel":"usa"},{"alvo":"peirce_celulas_canais","confianca":1.0,"epistemico":"fato","origem":"wiki","rel":"relaciona"}],"confianca":1.0,"contraexemplos":[],"descricao":"Carrega libcuda por dlopen, cria contexto/módulo/2 células com buffers persistentes; ≥4096 B parte o trabalho em 2 streams e combina via 'floco' Koch (gold+koch_coordenadas → 16 B, XOR IPC); <4096 B usa mono. Sem GPU cai em BROCA_ERR_BACKEND. libbroca_so_gpu.so.","epistemico":"fato","exemplos":[],"id":"gpu_backend_rede2","memoria":"semantica","meta":{"autor":"broca-so","dominio":"metal","fonte":"ula/broca_backend_gpu.c:gpu_gato_rede2"},"origem":"wiki","palavras_chave":[],"sinonimos":[],"tipo":"conceito","titulo":"Backend GPU rede-2 Persistente"}
\section{Backend GPU rede-2 Persistente}
Carrega libcuda por dlopen, cria contexto/módulo/2 células com buffers persistentes; \ensuremath{\ge}4096 B parte o trabalho em 2 streams e combina via 'floco' Koch (gold+koch\_coordenadas \ensuremath{\to} 16 B, XOR IPC); <4096 B usa mono. Sem GPU cai em BROCA\_ERR\_BACKEND. libbroca\_so\_gpu.so.
\prov{relações: usa gpu\_gato\_ptx; usa libbroca\_so\_camada; relaciona peirce\_celulas\_canais}
\prov{proveniência: wiki \textperiodcentered{} ula/broca\_backend\_gpu.c:gpu\_gato\_rede2 \textperiodcentered{} fato \textperiodcentered{} confiança 1.0 \textperiodcentered{} domínio metal \textperiodcentered{} história 4 versões no jornal}

%CRISTAL {"arestas":[{"alvo":"neuron_io_popcnt","confianca":1.0,"epistemico":"fato","origem":"wiki","rel":"especializa"},{"alvo":"metal_broca_so","confianca":1.0,"epistemico":"fato","origem":"wiki","rel":"parte_de"}],"confianca":1.0,"contraexemplos":[],"descricao":"Kernel PTX embutido (sm_75/sm_89, PTX 7.4/8.0) onde cada thread lê 1 byte e faz popc(byte&170)→atom.global.add em e e popc(byte&85)→o; é o gato rodando no device, a mesma álgebra do neurônio CPU. CUDA driver-only (dlopen libcuda.so.1); NÃO há OpenCL.","epistemico":"fato","exemplos":[],"id":"gpu_gato_ptx","memoria":"semantica","meta":{"autor":"broca-so","dominio":"metal","fonte":"ula/broca_gato_ptx.h; broca_backend_gpu.c"},"origem":"wiki","palavras_chave":[],"sinonimos":[],"tipo":"conceito","titulo":"Kernel PTX gato_stream na GPU"}
\section{Kernel PTX gato\_stream na GPU}
Kernel PTX embutido (sm\_75/sm\_89, PTX 7.4/8.0) onde cada thread lê 1 byte e faz popc(byte\&170)\ensuremath{\to}atom.global.add em e e popc(byte\&85)\ensuremath{\to}o; é o gato rodando no device, a mesma álgebra do neurônio CPU. CUDA driver-only (dlopen libcuda.so.1); NÃO há OpenCL.
\prov{relações: especializa neuron\_io\_popcnt; parte\_de metal\_broca\_so}
\prov{proveniência: wiki \textperiodcentered{} ula/broca\_gato\_ptx.h; broca\_backend\_gpu.c \textperiodcentered{} fato \textperiodcentered{} confiança 1.0 \textperiodcentered{} domínio metal \textperiodcentered{} história 3 versões no jornal}

%CRISTAL {"arestas":[{"alvo":"gpu_gato_ptx","confianca":1.0,"epistemico":"fato","origem":"wiki","rel":"usa"},{"alvo":"prova_de_trabalho","confianca":1.0,"epistemico":"fato","origem":"wiki","rel":"relaciona"}],"confianca":1.0,"contraexemplos":[],"descricao":"Honestidade do metal: NÃO existe benchmark de hashrate (H/s) em C. O único benchmark real é de throughput do kernel popcnt (broca_gpu_saturacao.c): ops/s, bytes/s, GB/s, crossover GPU>CPU, carga sustentada 5s. O hashrate de mineração vive só no pipeline Python/Stratum.","epistemico":"fato","exemplos":[],"id":"gpu_saturacao_bench","memoria":"semantica","meta":{"autor":"broca-so","dominio":"metal","fonte":"ula/testes/broca_gpu_saturacao.c"},"origem":"wiki","palavras_chave":[],"sinonimos":[],"tipo":"conceito","titulo":"Benchmark é Throughput (GB/s), não Hashrate"}
\section{Benchmark é Throughput (GB/s), não Hashrate}
Honestidade do metal: NÃO existe benchmark de hashrate (H/s) em C. O único benchmark real é de throughput do kernel popcnt (broca\_gpu\_saturacao.c): ops/s, bytes/s, GB/s, crossover GPU>CPU, carga sustentada 5s. O hashrate de mineração vive só no pipeline Python/Stratum.
\prov{relações: usa gpu\_gato\_ptx; relaciona prova\_de\_trabalho}
\prov{proveniência: wiki \textperiodcentered{} ula/testes/broca\_gpu\_saturacao.c \textperiodcentered{} fato \textperiodcentered{} confiança 1.0 \textperiodcentered{} domínio metal \textperiodcentered{} história 3 versões no jornal}

%CRISTAL {"arestas":[{"alvo":"problema_gravidade_quantica","confianca":1.0,"epistemico":"fato","origem":"wiki","rel":"relaciona"}],"confianca":1.0,"contraexemplos":[],"descricao":"Três programas discretos/alternativos: TRIANGULAÇÕES DINÂMICAS CAUSAIS (CDT: espaço-tempo de simplexes com estrutura causal; a dimensão espectral corre de ~4 a ~2 na escala de Planck); CONJUNTOS CAUSAIS (eventos com ordem parcial, 'ordem+número=geometria'; Sorkin estimou Λ~±1/√V antes de 1998); TWISTORES (Penrose: a geometria fundamental é ℂℙ³, não o espaço-tempo). Todos programas não confirmados.","epistemico":"hipotese","exemplos":[],"id":"gq_cdt_causal_twistor","memoria":"semantica","meta":{"autor":"broca-so","dominio":"fisica","fonte":"Ambjørn-Jurkiewicz-Loll 2004; Bombelli+1987; Penrose 1967"},"origem":"wiki","palavras_chave":[],"sinonimos":[],"tipo":"conceito","titulo":"CDT, conjuntos causais e twistores"}
\section{CDT, conjuntos causais e twistores}
Três programas discretos/alternativos: TRIANGULAÇÕES DINÂMICAS CAUSAIS (CDT: espaço-tempo de simplexes com estrutura causal; a dimensão espectral corre de \~{}4 a \~{}2 na escala de Planck); CONJUNTOS CAUSAIS (eventos com ordem parcial, 'ordem+número=geometria'; Sorkin estimou \ensuremath{\Lambda}\~{}$\pm$1/\ensuremath{\sqrt{\,}}V antes de 1998); TWISTORES (Penrose: a geometria fundamental é \ensuremath{\mathbb{C}}\ensuremath{\mathbb{P}}\ensuremath{^{3}}, não o espaço-tempo). Todos programas não confirmados.
\prov{relações: relaciona problema\_gravidade\_quantica}
\prov{proveniência: wiki \textperiodcentered{} Ambj\ensuremath{\varnothing}rn-Jurkiewicz-Loll 2004; Bombelli+1987; Penrose 1967 \textperiodcentered{} hipotese \textperiodcentered{} confiança 1.0 \textperiodcentered{} domínio fisica \textperiodcentered{} história 2 versões no jornal}

%CRISTAL {"arestas":[{"alvo":"materia_modelo_padrao","confianca":1.0,"epistemico":"fato","origem":"wiki","rel":"usa"},{"alvo":"unificacao_eletrofraca_gws","confianca":1.0,"epistemico":"fato","origem":"wiki","rel":"generaliza"},{"alvo":"cosmologia_quantica_hub","confianca":1.0,"epistemico":"fato","origem":"wiki","rel":"parte_de"}],"confianca":1.0,"contraexemplos":[],"descricao":"Unir SU(3)×SU(2)×U(1) num único grupo simples a alta energia (M_GUT~10¹⁶ GeV), do qual seriam quebras a baixa energia. Os 3 acoplamentos correm e QUASE convergem num ponto — a convergência melhora muito COM supersimetria (MSSM). Hipótese bem motivada (a quase-convergência é sugestiva), sem SUSY imperfeita; não confirmada.","epistemico":"hipotese","exemplos":[],"id":"grande_unificacao_gut","memoria":"semantica","meta":{"autor":"broca-so","dominio":"fisica","fonte":"Georgi-Quinn-Weinberg 1974 PRL 33,451"},"origem":"wiki","palavras_chave":[],"sinonimos":[],"tipo":"conceito","titulo":"A ideia da Grande Unificação (GUT)"}
\section{A ideia da Grande Unificação (GUT)}
Unir SU(3)$\times$SU(2)$\times$U(1) num único grupo simples a alta energia (M\_GUT\~{}10\ensuremath{^{16}} GeV), do qual seriam quebras a baixa energia. Os 3 acoplamentos correm e QUASE convergem num ponto — a convergência melhora muito COM supersimetria (MSSM). Hipótese bem motivada (a quase-convergência é sugestiva), sem SUSY imperfeita; não confirmada.
\prov{relações: usa materia\_modelo\_padrao; generaliza unificacao\_eletrofraca\_gws; parte\_de cosmologia\_quantica\_hub}
\prov{proveniência: wiki \textperiodcentered{} Georgi-Quinn-Weinberg 1974 PRL 33,451 \textperiodcentered{} hipotese \textperiodcentered{} confiança 1.0 \textperiodcentered{} domínio fisica \textperiodcentered{} história 4 versões no jornal}

%CRISTAL {"arestas":[{"alvo":"desitter_4d_universo","confianca":1.0,"epistemico":"fato","origem":"wiki","rel":"especializa"},{"alvo":"problema_gravidade_quantica","confianca":1.0,"epistemico":"fato","origem":"wiki","rel":"relaciona"},{"alvo":"gravidade_termodinamica_jacobson","confianca":1.0,"epistemico":"fato","origem":"wiki","rel":"relaciona"},{"alvo":"cosmologia_quantica_hub","confianca":1.0,"epistemico":"fato","origem":"wiki","rel":"parte_de"}],"confianca":1.0,"contraexemplos":[],"descricao":"A fase θ do campo local é o Killing temporal; o lapse N(ψ)=cosψ vem da pressão/associatividade Π=cos²ψ. A métrica 1+1 ds²=−cos²ψ dθ²+dψ² tem R=2 e R_ab=g_ab (Einstein): é dS₂. Em ℝ⁴, dS₄ (R=12, Λ=3). A curvatura CAI DA RÉGUA — não é postulada nem quantizada à mão. RESSALVA: é analogia de estrutura causal, NÃO dinâmica de Einstein ('F=2s é força algébrica, não curvatura de Einstein'). Firme como derivação geométrica (R verificado simbolicamente).","epistemico":"inferencia","exemplos":[],"id":"gravidade_desitter_derivada","memoria":"semantica","meta":{"autor":"broca-so","dominio":"cosmologia","fonte":"paper_metrica_circular.tex Teo. desitter; paper_2_cosmologia.tex"},"origem":"wiki","palavras_chave":[],"sinonimos":[],"tipo":"conceito","titulo":"Gravidade = de Sitter que cai da régua (não quantizada à mão)"}
\section{Gravidade = de Sitter que cai da régua (não quantizada à mão)}
A fase \ensuremath{\theta} do campo local é o Killing temporal; o lapse N(\ensuremath{\psi})=cos\ensuremath{\psi} vem da pressão/associatividade \ensuremath{\Pi}=cos\ensuremath{^{2}}\ensuremath{\psi}. A métrica 1+1 ds\ensuremath{^{2}}=\ensuremath{-}cos\ensuremath{^{2}}\ensuremath{\psi} d\ensuremath{\theta}\ensuremath{^{2}}+d\ensuremath{\psi}\ensuremath{^{2}} tem R=2 e R\_ab=g\_ab (Einstein): é dS\ensuremath{_{2}}. Em \ensuremath{\mathbb{R}}\ensuremath{^{4}}, dS\ensuremath{_{4}} (R=12, \ensuremath{\Lambda}=3). A curvatura CAI DA RÉGUA — não é postulada nem quantizada à mão. RESSALVA: é analogia de estrutura causal, NÃO dinâmica de Einstein ('F=2s é força algébrica, não curvatura de Einstein'). Firme como derivação geométrica (R verificado simbolicamente).
\prov{relações: especializa desitter\_4d\_universo; relaciona problema\_gravidade\_quantica; relaciona gravidade\_termodinamica\_jacobson; parte\_de cosmologia\_quantica\_hub}
\prov{proveniência: wiki \textperiodcentered{} paper\_metrica\_circular.tex Teo. desitter; paper\_2\_cosmologia.tex \textperiodcentered{} inferencia \textperiodcentered{} confiança 1.0 \textperiodcentered{} domínio cosmologia \textperiodcentered{} história 5 versões no jornal}

%CRISTAL {"arestas":[{"alvo":"gravidade_termodinamica_jacobson","confianca":1.0,"epistemico":"fato","origem":"wiki","rel":"especializa"},{"alvo":"landauer_gsl_bekenstein_bound","confianca":1.0,"epistemico":"fato","origem":"wiki","rel":"usa"},{"alvo":"cosmologia_quantica_hub","confianca":1.0,"epistemico":"fato","origem":"wiki","rel":"parte_de"}],"confianca":1.0,"contraexemplos":[],"descricao":"Gravidade como força entrópica emergente (telas holográficas maximizando entropia): FΔx=TΔS recupera Newton; a versão 2016 tenta explicar curvas de rotação sem matéria escura. CONTROVERSA / contestada (Kobakhidze: incompatível com interferometria de nêutrons; falhas em aglomerados e CMB). Não é aceita como estabelecida.","epistemico":"hipotese","exemplos":[],"id":"gravidade_entropica_verlinde","memoria":"semantica","meta":{"autor":"broca-so","dominio":"fisica","fonte":"Verlinde 2011 JHEP 04,029 arXiv:1001.0785 (contestada)"},"origem":"wiki","palavras_chave":[],"sinonimos":[],"tipo":"conceito","titulo":"Gravidade entrópica (Verlinde)"}
\section{Gravidade entrópica (Verlinde)}
Gravidade como força entrópica emergente (telas holográficas maximizando entropia): F\ensuremath{\Delta}x=T\ensuremath{\Delta}S recupera Newton; a versão 2016 tenta explicar curvas de rotação sem matéria escura. CONTROVERSA / contestada (Kobakhidze: incompatível com interferometria de nêutrons; falhas em aglomerados e CMB). Não é aceita como estabelecida.
\prov{relações: especializa gravidade\_termodinamica\_jacobson; usa landauer\_gsl\_bekenstein\_bound; parte\_de cosmologia\_quantica\_hub}
\prov{proveniência: wiki \textperiodcentered{} Verlinde 2011 JHEP 04,029 arXiv:1001.0785 (contestada) \textperiodcentered{} hipotese \textperiodcentered{} confiança 1.0 \textperiodcentered{} domínio fisica \textperiodcentered{} história 4 versões no jornal}

%CRISTAL {"arestas":[{"alvo":"problema_gravidade_quantica","confianca":1.0,"epistemico":"fato","origem":"wiki","rel":"relaciona"},{"alvo":"loop_quantum_cosmology","confianca":1.0,"epistemico":"fato","origem":"wiki","rel":"generaliza"}],"confianca":1.0,"contraexemplos":[],"descricao":"Quantização canônica e independente de fundo da RG (variáveis de Ashtekar): o espaço-tempo é DISCRETO na escala de Planck — áreas/volumes quantizados (A=8πγℓ_P²Σ√(jᵢ(jᵢ+1))), estados = spin networks (Penrose), evolução = spin foams. Distinta da LQC (aplicação cosmológica simétrica, o 'big bounce'). Programa não confirmado.","epistemico":"hipotese","exemplos":[],"id":"gravidade_quantica_loop_lqg","memoria":"semantica","meta":{"autor":"broca-so","dominio":"fisica","fonte":"Ashtekar 1986; Rovelli-Smolin 1995 gr-qc/9411005"},"origem":"wiki","palavras_chave":[],"sinonimos":[],"tipo":"conceito","titulo":"Gravitação quântica em loop (LQG)"}
\section{Gravitação quântica em loop (LQG)}
Quantização canônica e independente de fundo da RG (variáveis de Ashtekar): o espaço-tempo é DISCRETO na escala de Planck — áreas/volumes quantizados (A=8\ensuremath{\pi}\ensuremath{\gamma}\ensuremath{\ell}\_P\ensuremath{^{2}}\ensuremath{\Sigma}\ensuremath{\sqrt{\,}}(j\ensuremath{_{i}}(j\ensuremath{_{i}}+1))), estados = spin networks (Penrose), evolução = spin foams. Distinta da LQC (aplicação cosmológica simétrica, o 'big bounce'). Programa não confirmado.
\prov{relações: relaciona problema\_gravidade\_quantica; generaliza loop\_quantum\_cosmology}
\prov{proveniência: wiki \textperiodcentered{} Ashtekar 1986; Rovelli-Smolin 1995 gr-qc/9411005 \textperiodcentered{} hipotese \textperiodcentered{} confiança 1.0 \textperiodcentered{} domínio fisica \textperiodcentered{} história 3 versões no jornal}

%CRISTAL {"arestas":[{"alvo":"teoria_cordas_m","confianca":1.0,"epistemico":"fato","origem":"wiki","rel":"relaciona"},{"alvo":"gravidade_quantica_loop_lqg","confianca":1.0,"epistemico":"fato","origem":"wiki","rel":"relaciona"}],"confianca":1.0,"contraexemplos":[],"descricao":"FATO: nenhuma teoria de gravitação quântica foi confirmada; nenhuma faz hoje uma previsão distintiva testada. Cordas, LQG, segurança assintótica, CDT, causal sets, twistores são programas concorrentes não verificados. Testes indiretos (violação de Lorentz, dispersão de fótons de GRBs, ondas primordiais) só dão limites, não detecções.","epistemico":"fato","exemplos":[],"id":"gravidade_quantica_status","memoria":"semantica","meta":{"autor":"broca-so","dominio":"fisica","fonte":"consenso; PDG/reviews"},"origem":"wiki","palavras_chave":[],"sinonimos":[],"tipo":"conceito","titulo":"Status: nenhuma gravidade quântica confirmada"}
\section{Status: nenhuma gravidade quântica confirmada}
FATO: nenhuma teoria de gravitação quântica foi confirmada; nenhuma faz hoje uma previsão distintiva testada. Cordas, LQG, segurança assintótica, CDT, causal sets, twistores são programas concorrentes não verificados. Testes indiretos (violação de Lorentz, dispersão de fótons de GRBs, ondas primordiais) só dão limites, não detecções.
\prov{relações: relaciona teoria\_cordas\_m; relaciona gravidade\_quantica\_loop\_lqg}
\prov{proveniência: wiki \textperiodcentered{} consenso; PDG/reviews \textperiodcentered{} fato \textperiodcentered{} confiança 1.0 \textperiodcentered{} domínio fisica \textperiodcentered{} história 3 versões no jornal}

%CRISTAL {"arestas":[{"alvo":"gravidade_desitter_derivada","confianca":1.0,"epistemico":"fato","origem":"wiki","rel":"especializa"},{"alvo":"constantes_horizonte_algebricas","confianca":1.0,"epistemico":"fato","origem":"wiki","rel":"usa"}],"confianca":1.0,"contraexemplos":[],"descricao":"Na forma de semiplano ds²=(dx²+dy²)/y², o numerador é plano (o que o elétron vê localmente) e o 1/y² é a régua renormalizada: a curvatura K=−1 não está no espaço, mas na ESCALA — 'toda a gravidade do espaço de álgebras mora na torção de escala, não na geometria' (forma teleparalela). A gravidade de superfície algébrica κ=√2 (Lyapunov de ẍ=2s) dá T=√2/2π; no relógio de fase T_dS=1/2π (Gibbons-Hawking, regularidade cônica = conexidade 0,999…=0).","epistemico":"inferencia","exemplos":[],"id":"gravidade_teleparalela_torcao","memoria":"semantica","meta":{"autor":"broca-so","dominio":"cosmologia","fonte":"paper_metrica_circular.tex obs:metrica-final"},"origem":"wiki","palavras_chave":[],"sinonimos":[],"tipo":"conceito","titulo":"Gravidade teleparalela: a curvatura mora na torção de escala"}
\section{Gravidade teleparalela: a curvatura mora na torção de escala}
Na forma de semiplano ds\ensuremath{^{2}}=(dx\ensuremath{^{2}}+dy\ensuremath{^{2}})/y\ensuremath{^{2}}, o numerador é plano (o que o elétron vê localmente) e o 1/y\ensuremath{^{2}} é a régua renormalizada: a curvatura K=\ensuremath{-}1 não está no espaço, mas na ESCALA — 'toda a gravidade do espaço de álgebras mora na torção de escala, não na geometria' (forma teleparalela). A gravidade de superfície algébrica \ensuremath{\kappa}=\ensuremath{\sqrt{\,}}2 (Lyapunov de \ensuremath{\ddot{x}}=2s) dá T=\ensuremath{\sqrt{\,}}2/2\ensuremath{\pi}; no relógio de fase T\_dS=1/2\ensuremath{\pi} (Gibbons-Hawking, regularidade cônica = conexidade 0,999…=0).
\prov{relações: especializa gravidade\_desitter\_derivada; usa constantes\_horizonte\_algebricas}
\prov{proveniência: wiki \textperiodcentered{} paper\_metrica\_circular.tex obs:metrica-final \textperiodcentered{} inferencia \textperiodcentered{} confiança 1.0 \textperiodcentered{} domínio cosmologia \textperiodcentered{} história 3 versões no jornal}

%CRISTAL {"arestas":[{"alvo":"bekenstein_hawking_entropia","confianca":1.0,"epistemico":"fato","origem":"wiki","rel":"usa"},{"alvo":"efe_einstein","confianca":1.0,"epistemico":"fato","origem":"wiki","rel":"explica"},{"alvo":"cosmologia_quantica_hub","confianca":1.0,"epistemico":"fato","origem":"wiki","rel":"parte_de"}],"confianca":1.0,"contraexemplos":[],"descricao":"Assumindo δQ=TdS em todo horizonte de Rindler local (T=Unruh, S∝área), as equações de Einstein EMERGEM como equação de estado: G_μν=8πG/c⁴ T_μν. Sugere gravidade emergente/termodinâmica, não fundamental. A dedução funciona sob suas hipóteses (fato); a interpretação 'gravidade é emergente' é programa, não consenso — sugestivo, não conclusivo.","epistemico":"inferencia","exemplos":[],"id":"gravidade_termodinamica_jacobson","memoria":"semantica","meta":{"autor":"broca-so","dominio":"fisica","fonte":"Jacobson 1995 PRL 75,1260 gr-qc/9504004"},"origem":"wiki","palavras_chave":[],"sinonimos":[],"tipo":"conceito","titulo":"Gravidade como termodinâmica (Jacobson)"}
\section{Gravidade como termodinâmica (Jacobson)}
Assumindo \ensuremath{\delta}Q=TdS em todo horizonte de Rindler local (T=Unruh, S\ensuremath{\propto}área), as equações de Einstein EMERGEM como equação de estado: G\_\ensuremath{\mu}\ensuremath{\nu}=8\ensuremath{\pi}G/c\ensuremath{^{4}} T\_\ensuremath{\mu}\ensuremath{\nu}. Sugere gravidade emergente/termodinâmica, não fundamental. A dedução funciona sob suas hipóteses (fato); a interpretação 'gravidade é emergente' é programa, não consenso — sugestivo, não conclusivo.
\prov{relações: usa bekenstein\_hawking\_entropia; explica efe\_einstein; parte\_de cosmologia\_quantica\_hub}
\prov{proveniência: wiki \textperiodcentered{} Jacobson 1995 PRL 75,1260 gr-qc/9504004 \textperiodcentered{} inferencia \textperiodcentered{} confiança 1.0 \textperiodcentered{} domínio fisica \textperiodcentered{} história 4 versões no jornal}

%CRISTAL {"arestas":[{"alvo":"unificacao_eletrofraca_gws","confianca":1.0,"epistemico":"fato","origem":"wiki","rel":"explica"},{"alvo":"qcd_cromodinamica","confianca":1.0,"epistemico":"fato","origem":"wiki","rel":"explica"},{"alvo":"carga_q_n_sobre_3","confianca":1.0,"epistemico":"fato","origem":"wiki","rel":"usa"},{"alvo":"isospin_hipercarga_algebrico","confianca":1.0,"epistemico":"fato","origem":"wiki","rel":"usa"},{"alvo":"angulo_weinberg","confianca":1.0,"epistemico":"fato","origem":"wiki","rel":"relaciona"},{"alvo":"leitura_de_part_culas_o_setor_de_materia_pela_torre_de_divisao","confianca":1.0,"epistemico":"fato","origem":"wiki","rel":"parte_de"},{"alvo":"cosmologia_quantica_hub","confianca":1.0,"epistemico":"fato","origem":"wiki","rel":"parte_de"}],"confianca":1.0,"contraexemplos":[],"descricao":"A leitura algébrica fecha o grupo de gauge do Modelo Padrão: ℂ⊗𝕆=Cℓ(6) dá cor SU(3)_c + carga Q=N/3 (U(1)_em); ℂ⊗ℍ=Cℓ(2) dá o isospin SU(2)_L (dubleto T₃=η†η−½) e, com Gell-Mann-Nishijima Y=2(Q−T₃), a hipercarga U(1)_Y. Junto: SU(3)_c×SU(2)_L×U(1)_Y. [D] como isomorfismo estrutural nos números quânticos (sem ajuste); [I] NÃO é a teoria de campo (Lagrangiana/massas/θ_W não construídos).","epistemico":"inferencia","exemplos":[],"id":"grupo_padrao_cxho","memoria":"semantica","meta":{"autor":"broca-so","dominio":"cosmologia","fonte":"paper_4_particulas.tex obs:cxh-eletrofraco, obs:cxo-cargas"},"origem":"wiki","palavras_chave":[],"sinonimos":[],"tipo":"conceito","titulo":"O grupo do MP de ℂ⊗(ℍ⊗𝕆) = SU(3)×SU(2)×U(1)"}
\section{O grupo do MP de \ensuremath{\mathbb{C}}\ensuremath{\otimes}(\ensuremath{\mathbb{H}}\ensuremath{\otimes}\ensuremath{\mathbb{O}}) = SU(3)$\times$SU(2)$\times$U(1)}
A leitura algébrica fecha o grupo de gauge do Modelo Padrão: \ensuremath{\mathbb{C}}\ensuremath{\otimes}\ensuremath{\mathbb{O}}=C\ensuremath{\ell}(6) dá cor SU(3)\_c + carga Q=N/3 (U(1)\_em); \ensuremath{\mathbb{C}}\ensuremath{\otimes}\ensuremath{\mathbb{H}}=C\ensuremath{\ell}(2) dá o isospin SU(2)\_L (dubleto T\ensuremath{_{3}}=\ensuremath{\eta}\ensuremath{\dagger}\ensuremath{\eta}\ensuremath{-}\ensuremath{\tfrac12}) e, com Gell-Mann-Nishijima Y=2(Q\ensuremath{-}T\ensuremath{_{3}}), a hipercarga U(1)\_Y. Junto: SU(3)\_c$\times$SU(2)\_L$\times$U(1)\_Y. [D] como isomorfismo estrutural nos números quânticos (sem ajuste); [I] NÃO é a teoria de campo (Lagrangiana/massas/\ensuremath{\theta}\_W não construídos).
\prov{relações: explica unificacao\_eletrofraca\_gws; explica qcd\_cromodinamica; usa carga\_q\_n\_sobre\_3; usa isospin\_hipercarga\_algebrico; relaciona angulo\_weinberg; parte\_de leitura\_de\_part\_culas\_o\_setor\_de\_materia\_pela\_torre\_de\_divisao; parte\_de cosmologia\_quantica\_hub}
\prov{proveniência: wiki \textperiodcentered{} paper\_4\_particulas.tex obs:cxh-eletrofraco, obs:cxo-cargas \textperiodcentered{} inferencia \textperiodcentered{} confiança 1.0 \textperiodcentered{} domínio cosmologia \textperiodcentered{} história 8 versões no jornal}

%CRISTAL {"arestas":[{"alvo":"so10_pati_salam","confianca":1.0,"epistemico":"fato","origem":"wiki","rel":"generaliza"}],"confianca":1.0,"contraexemplos":[],"descricao":"Cadeia SU(5)⊂SO(10)⊂E6⊂E7⊂E8 como grupos de unificação maiores; em E6 uma geração cabe na 27. A álgebra de Jordan excepcional J₃(𝕆) (Albert): matrizes hermitianas 3×3 de octônios, dim 27 (= o 27 de E6), Aut=F₄, determinante preservado por E6. A 'magic square' de Freudenthal-Tits gera F4/E6/E7/E8 de pares de álgebras de divisão. Programa matemático especulativo (Baez, Todorov-Dubois-Violette); elegante, sem previsão confirmada.","epistemico":"hipotese","exemplos":[],"id":"grupos_excepcionais_jordan","memoria":"semantica","meta":{"autor":"broca-so","dominio":"fisica","fonte":"Gürsey-Ramond-Sikivie 1976; Baez 2002 math/0105155; JvNW 1934"},"origem":"wiki","palavras_chave":[],"sinonimos":[],"tipo":"conceito","titulo":"Grupos excepcionais (E6/E7/E8) e a Jordan J₃(𝕆)/Albert"}
\section{Grupos excepcionais (E6/E7/E8) e a Jordan J\ensuremath{_{3}}(\ensuremath{\mathbb{O}})/Albert}
Cadeia SU(5)\ensuremath{\subset}SO(10)\ensuremath{\subset}E6\ensuremath{\subset}E7\ensuremath{\subset}E8 como grupos de unificação maiores; em E6 uma geração cabe na 27. A álgebra de Jordan excepcional J\ensuremath{_{3}}(\ensuremath{\mathbb{O}}) (Albert): matrizes hermitianas 3$\times$3 de octônios, dim 27 (= o 27 de E6), Aut=F\ensuremath{_{4}}, determinante preservado por E6. A 'magic square' de Freudenthal-Tits gera F4/E6/E7/E8 de pares de álgebras de divisão. Programa matemático especulativo (Baez, Todorov-Dubois-Violette); elegante, sem previsão confirmada.
\prov{relações: generaliza so10\_pati\_salam}
\prov{proveniência: wiki \textperiodcentered{} Gürsey-Ramond-Sikivie 1976; Baez 2002 math/0105155; JvNW 1934 \textperiodcentered{} hipotese \textperiodcentered{} confiança 1.0 \textperiodcentered{} domínio fisica \textperiodcentered{} história 2 versões no jornal}

%CRISTAL {"arestas":[{"alvo":"efeito_hall_quantico_inteiro","confianca":1.0,"epistemico":"fato","origem":"wiki","rel":"explica"},{"alvo":"alpha_tres_reguas","confianca":1.0,"epistemico":"fato","origem":"wiki","rel":"usa"},{"alvo":"chern_fibrado_hopf","confianca":1.0,"epistemico":"fato","origem":"wiki","rel":"usa"}],"confianca":1.0,"contraexemplos":[],"descricao":"σ_xy=n·e²/h=n·(2α/Z0) fatora-se em duas peças da série: n∈ℤ é a CONTAGEM (o número de Chern, a costura 1≡0 agora topológica) e α é a ESCALA (o acoplamento). A exatidão da metrologia mais precisa que existe não vem da pureza do material, mas DA TOPOLOGIA: R_K=h/e²=Z0/2α é a resistência quântica de von Klitzing, o inverso do acoplamento por volta — a 'terceira régua' de α. [D] a álgebra e n=Chern; [N] α=Z0/2R_K=1/137,036; [I] a leitura contagem×escala.","epistemico":"inferencia","exemplos":[],"id":"hall_quantico_contagem_escala","memoria":"semantica","meta":{"autor":"broca-so","dominio":"cosmologia","fonte":"paper_5_unidades.tex obs:alpha-vonklitzing"},"origem":"wiki","palavras_chave":[],"sinonimos":[],"tipo":"conceito","titulo":"O Hall quântico = contagem × escala (topologia × α)"}
\section{O Hall quântico = contagem $\times$ escala (topologia $\times$ \ensuremath{\alpha})}
\ensuremath{\sigma}\_xy=n\textperiodcentered e\ensuremath{^{2}}/h=n\textperiodcentered (2\ensuremath{\alpha}/Z0) fatora-se em duas peças da série: n\ensuremath{\in}\ensuremath{\mathbb{Z}} é a CONTAGEM (o número de Chern, a costura 1\ensuremath{\equiv}0 agora topológica) e \ensuremath{\alpha} é a ESCALA (o acoplamento). A exatidão da metrologia mais precisa que existe não vem da pureza do material, mas DA TOPOLOGIA: R\_K=h/e\ensuremath{^{2}}=Z0/2\ensuremath{\alpha} é a resistência quântica de von Klitzing, o inverso do acoplamento por volta — a 'terceira régua' de \ensuremath{\alpha}. [D] a álgebra e n=Chern; [N] \ensuremath{\alpha}=Z0/2R\_K=1/137,036; [I] a leitura contagem$\times$escala.
\prov{relações: explica efeito\_hall\_quantico\_inteiro; usa alpha\_tres\_reguas; usa chern\_fibrado\_hopf}
\prov{proveniência: wiki \textperiodcentered{} paper\_5\_unidades.tex obs:alpha-vonklitzing \textperiodcentered{} inferencia \textperiodcentered{} confiança 1.0 \textperiodcentered{} domínio cosmologia \textperiodcentered{} história 4 versões no jornal}

%CRISTAL {"arestas":[{"alvo":"wheeler_dewitt","confianca":1.0,"epistemico":"fato","origem":"wiki","rel":"usa"}],"confianca":1.0,"contraexemplos":[],"descricao":"Função de onda do universo como integral de caminho sobre geometrias euclidianas compactas SEM fronteira; o espaço-tempo arredonda o Big Bang, eliminando a singularidade inicial. Condição de contorno especulativa, sem verificação e em disputa (Feldbrugge-Lehners-Turok). Concorre com o tunelamento de Vilenkin.","epistemico":"hipotese","exemplos":[],"id":"hartle_hawking","memoria":"semantica","meta":{"autor":"broca-so","dominio":"fisica","fonte":"Hartle & Hawking 1983, Phys.Rev.D 28,2960"},"origem":"wiki","palavras_chave":[],"sinonimos":[],"tipo":"conceito","titulo":"Proposta de não-contorno (Hartle-Hawking)"}
\section{Proposta de não-contorno (Hartle-Hawking)}
Função de onda do universo como integral de caminho sobre geometrias euclidianas compactas SEM fronteira; o espaço-tempo arredonda o Big Bang, eliminando a singularidade inicial. Condição de contorno especulativa, sem verificação e em disputa (Feldbrugge-Lehners-Turok). Concorre com o tunelamento de Vilenkin.
\prov{relações: usa wheeler\_dewitt}
\prov{proveniência: wiki \textperiodcentered{} Hartle \& Hawking 1983, Phys.Rev.D 28,2960 \textperiodcentered{} hipotese \textperiodcentered{} confiança 1.0 \textperiodcentered{} domínio fisica \textperiodcentered{} história 2 versões no jornal}

%CRISTAL {"arestas":[{"alvo":"problema_do_sabor","confianca":1.0,"epistemico":"fato","origem":"wiki","rel":"relaciona"},{"alvo":"acoplamentos_yukawa","confianca":1.0,"epistemico":"fato","origem":"wiki","rel":"relaciona"},{"alvo":"os_dois_tres_cores_hurwitz_e_gerac","confianca":1.0,"epistemico":"fato","origem":"wiki","rel":"usa"}],"confianca":1.0,"contraexemplos":[],"descricao":"Sob a trialidade EXATA as três gerações são indistinguíveis (massas iguais); a hierarquia m_e≪m_μ≪m_τ É a quebra dessa simetria, com estrutura Z₃. 'A álgebra dá a estrutura, não o número': m_p/m_e~1700 (ordem certa via m_e~Yukawa da quebra, m_p~2√σ). LACUNA honesta: os arquivos NÃO trazem y_top≈1 vs y_up~10⁻⁵ nem a escala Yukawa individual — só 'massas iguais na simetria, hierarquia=quebra'. [I]/[J] correspondência + hipótese.","epistemico":"inferencia","exemplos":[],"id":"hierarquia_massas_quebra_trialidade","memoria":"semantica","meta":{"autor":"broca-so","dominio":"cosmologia","fonte":"paper_4_particulas.tex obs:geracoes"},"origem":"wiki","palavras_chave":[],"sinonimos":[],"tipo":"conceito","titulo":"Hierarquia de massas = quebra da trialidade"}
\section{Hierarquia de massas = quebra da trialidade}
Sob a trialidade EXATA as três gerações são indistinguíveis (massas iguais); a hierarquia m\_e\ensuremath{\ll}m\_\ensuremath{\mu}\ensuremath{\ll}m\_\ensuremath{\tau} É a quebra dessa simetria, com estrutura Z\ensuremath{_{3}}. 'A álgebra dá a estrutura, não o número': m\_p/m\_e\~{}1700 (ordem certa via m\_e\~{}Yukawa da quebra, m\_p\~{}2\ensuremath{\sqrt{\,}}\ensuremath{\sigma}). LACUNA honesta: os arquivos NÃO trazem y\_top\ensuremath{\approx}1 vs y\_up\~{}10\ensuremath{^{-5}} nem a escala Yukawa individual — só 'massas iguais na simetria, hierarquia=quebra'. [I]/[J] correspondência + hipótese.
\prov{relações: relaciona problema\_do\_sabor; relaciona acoplamentos\_yukawa; usa os\_dois\_tres\_cores\_hurwitz\_e\_gerac}
\prov{proveniência: wiki \textperiodcentered{} paper\_4\_particulas.tex obs:geracoes \textperiodcentered{} inferencia \textperiodcentered{} confiança 1.0 \textperiodcentered{} domínio cosmologia \textperiodcentered{} história 4 versões no jornal}

%CRISTAL {"arestas":[{"alvo":"mecanismo_higgs_ebh","confianca":1.0,"epistemico":"fato","origem":"wiki","rel":"relaciona"},{"alvo":"mecanismo_higgs_algebrico","confianca":1.0,"epistemico":"fato","origem":"wiki","rel":"especializa"},{"alvo":"mass_gap_higgs_bz","confianca":1.0,"epistemico":"fato","origem":"wiki","rel":"relaciona"},{"alvo":"teoria_gauge_octonica","confianca":1.0,"epistemico":"fato","origem":"wiki","rel":"usa"},{"alvo":"g2_grupo_excepcional","confianca":1.0,"epistemico":"fato","origem":"wiki","rel":"usa"}],"confianca":1.0,"contraexemplos":[],"descricao":"Sobre Yang-Mills octônica soma-se V(A)=−μ²‖A‖²+λ‖A‖⁴+γΠ⟨‖J‖²⟩ (o 3º termo é o Jacobiator do produto vetorial 7D, novidade octônica que zera em ℝ³). O vev ‖v‖²=μ²/2λeff quebra G₂→SU(3): dos 14 bósons, 6 (em 3⊕3̄) absorvem 6 Goldstones e ganham massa MX²=g²‖v‖²; 1 Higgs radial M²=4μ²; os 8 de su(3) ficam sem massa (exigem confinamento). [D] o vev e a quebra (StabG₂(v)=SU(3), clássico); [J/I] que YM-𝕆/G₂ seja a física real; [N] correção do Jacobiator ~0,5%.","epistemico":"inferencia","exemplos":[],"id":"higgs_algebrico_g2_su3","memoria":"semantica","meta":{"autor":"broca-so","dominio":"cosmologia","fonte":"paper_BZ_mass_gap_higgs_algebrico.tex §Vácuo e quebra espontânea"},"origem":"wiki","palavras_chave":[],"sinonimos":[],"tipo":"conceito","titulo":"Mecanismo de Higgs algébrico (YM-𝕆, G₂→SU(3))"}
\section{Mecanismo de Higgs algébrico (YM-\ensuremath{\mathbb{O}}, G\ensuremath{_{2}}\ensuremath{\to}SU(3))}
Sobre Yang-Mills octônica soma-se V(A)=\ensuremath{-}\ensuremath{\mu}\ensuremath{^{2}}\ensuremath{\|}A\ensuremath{\|}\ensuremath{^{2}}+\ensuremath{\lambda}\ensuremath{\|}A\ensuremath{\|}\ensuremath{^{4}}+\ensuremath{\gamma}\ensuremath{\Pi}\ensuremath{\langle}\ensuremath{\|}J\ensuremath{\|}\ensuremath{^{2}}\ensuremath{\rangle} (o 3º termo é o Jacobiator do produto vetorial 7D, novidade octônica que zera em \ensuremath{\mathbb{R}}\ensuremath{^{3}}). O vev \ensuremath{\|}v\ensuremath{\|}\ensuremath{^{2}}=\ensuremath{\mu}\ensuremath{^{2}}/2\ensuremath{\lambda}eff quebra G\ensuremath{_{2}}\ensuremath{\to}SU(3): dos 14 bósons, 6 (em 3\ensuremath{\oplus}3\ensuremath{}) absorvem 6 Goldstones e ganham massa MX\ensuremath{^{2}}=g\ensuremath{^{2}}\ensuremath{\|}v\ensuremath{\|}\ensuremath{^{2}}; 1 Higgs radial M\ensuremath{^{2}}=4\ensuremath{\mu}\ensuremath{^{2}}; os 8 de su(3) ficam sem massa (exigem confinamento). [D] o vev e a quebra (StabG\ensuremath{_{2}}(v)=SU(3), clássico); [J/I] que YM-\ensuremath{\mathbb{O}}/G\ensuremath{_{2}} seja a física real; [N] correção do Jacobiator \~{}0,5\%.
\prov{relações: relaciona mecanismo\_higgs\_ebh; especializa mecanismo\_higgs\_algebrico; relaciona mass\_gap\_higgs\_bz; usa teoria\_gauge\_octonica; usa g2\_grupo\_excepcional}
\prov{proveniência: wiki \textperiodcentered{} paper\_BZ\_mass\_gap\_higgs\_algebrico.tex §Vácuo e quebra espontânea \textperiodcentered{} inferencia \textperiodcentered{} confiança 1.0 \textperiodcentered{} domínio cosmologia \textperiodcentered{} história 7 versões no jornal}

%CRISTAL {"arestas":[{"alvo":"de_sitter_lambda_energia_escura","confianca":1.0,"epistemico":"fato","origem":"wiki","rel":"relaciona"},{"alvo":"cosmologia_quantica_hub","confianca":1.0,"epistemico":"fato","origem":"wiki","rel":"parte_de"}],"confianca":1.0,"contraexemplos":[],"descricao":"Um observador em universo acelerado (de Sitter) tem um horizonte cosmológico também TERMODINÂMICO: banho térmico T_dS=ℏH/2πk_Bc (T=H/2π natural) e entropia S=A/4=3π/(Λℓ_P²). A termodinâmica de horizonte não é exclusiva de buracos negros. Aberto/difícil: os graus de liberdade microscópicos e a holografia de de Sitter (dS/CFT).","epistemico":"fato","exemplos":[],"id":"horizonte_desitter_gibbons_hawking","memoria":"semantica","meta":{"autor":"broca-so","dominio":"fisica","fonte":"Gibbons & Hawking 1977 PRD 15,2738"},"origem":"wiki","palavras_chave":[],"sinonimos":[],"tipo":"conceito","titulo":"Horizonte de de Sitter (Gibbons-Hawking)"}
\section{Horizonte de de Sitter (Gibbons-Hawking)}
Um observador em universo acelerado (de Sitter) tem um horizonte cosmológico também TERMODINÂMICO: banho térmico T\_dS=\ensuremath{\hbar}H/2\ensuremath{\pi}k\_Bc (T=H/2\ensuremath{\pi} natural) e entropia S=A/4=3\ensuremath{\pi}/(\ensuremath{\Lambda}\ensuremath{\ell}\_P\ensuremath{^{2}}). A termodinâmica de horizonte não é exclusiva de buracos negros. Aberto/difícil: os graus de liberdade microscópicos e a holografia de de Sitter (dS/CFT).
\prov{relações: relaciona de\_sitter\_lambda\_energia\_escura; parte\_de cosmologia\_quantica\_hub}
\prov{proveniência: wiki \textperiodcentered{} Gibbons \& Hawking 1977 PRD 15,2738 \textperiodcentered{} fato \textperiodcentered{} confiança 1.0 \textperiodcentered{} domínio fisica \textperiodcentered{} história 3 versões no jornal}

%CRISTAL {"arestas":[{"alvo":"cosmologia_quantica_hub","confianca":1.0,"epistemico":"fato","origem":"wiki","rel":"parte_de"}],"confianca":1.0,"contraexemplos":[],"descricao":"Velocidade de recessão ∝ distância. Lemaître (1927) derivou a relação linear das EFE; Hubble (1929) deu a base observacional. Renomeada pela IAU (2018). Problema aberto: a tensão de Hubble.","epistemico":"fato","exemplos":[],"id":"hubble_lemaitre","memoria":"semantica","meta":{"autor":"broca-so","dominio":"fisica","fonte":"Lemaître 1927; Hubble 1929 PNAS"},"origem":"wiki","palavras_chave":[],"sinonimos":[],"tipo":"conceito","titulo":"Lei de Hubble-Lemaître (v=H₀d)"}
\section{Lei de Hubble-Lemaître (v=H\ensuremath{_{0}}d)}
Velocidade de recessão \ensuremath{\propto} distância. Lemaître (1927) derivou a relação linear das EFE; Hubble (1929) deu a base observacional. Renomeada pela IAU (2018). Problema aberto: a tensão de Hubble.
\prov{relações: parte\_de cosmologia\_quantica\_hub}
\prov{proveniência: wiki \textperiodcentered{} Lemaître 1927; Hubble 1929 PNAS \textperiodcentered{} fato \textperiodcentered{} confiança 1.0 \textperiodcentered{} domínio fisica \textperiodcentered{} história 2 versões no jornal}

%CRISTAL {"arestas":[{"alvo":"cosmologia_quantica_hub","confianca":1.0,"epistemico":"fato","origem":"wiki","rel":"parte_de"}],"confianca":1.0,"contraexemplos":[],"descricao":"Expansão exponencial primordial. Guth (1981) 'old inflation'; Linde/Albrecht-Steinhardt: slow-roll (o inflaton rola lento por um platô). Resolve horizonte, planura, monopolos. Predição: espectro quase invariante de escala ns=0,965±0,004 (Planck, ns<1 favorece inflação). Aberto: qual o inflaton; ondas gravitacionais primordiais (r) só limitadas.","epistemico":"inferencia","exemplos":[],"id":"inflacao_cosmica","memoria":"semantica","meta":{"autor":"broca-so","dominio":"cosmologia_dados","fonte":"Guth 1981 PRD 23,347; Linde 1982; Planck 2018"},"origem":"wiki","palavras_chave":[],"sinonimos":[],"tipo":"conceito","titulo":"Inflação cósmica (Guth 1981; slow-roll)"}
\section{Inflação cósmica (Guth 1981; slow-roll)}
Expansão exponencial primordial. Guth (1981) 'old inflation'; Linde/Albrecht-Steinhardt: slow-roll (o inflaton rola lento por um platô). Resolve horizonte, planura, monopolos. Predição: espectro quase invariante de escala ns=0,965$\pm$0,004 (Planck, ns<1 favorece inflação). Aberto: qual o inflaton; ondas gravitacionais primordiais (r) só limitadas.
\prov{relações: parte\_de cosmologia\_quantica\_hub}
\prov{proveniência: wiki \textperiodcentered{} Guth 1981 PRD 23,347; Linde 1982; Planck 2018 \textperiodcentered{} inferencia \textperiodcentered{} confiança 1.0 \textperiodcentered{} domínio cosmologia\_dados \textperiodcentered{} história 2 versões no jornal}

%CRISTAL {"arestas":[{"alvo":"cosmologia_quantica_hub","confianca":1.0,"epistemico":"fato","origem":"wiki","rel":"parte_de"}],"confianca":1.0,"contraexemplos":[],"descricao":"O qubit é um sistema de 2 níveis em ℂ² (|ψ⟩=α|0⟩+β|1⟩, ponto na esfera de Bloch); o EMARANHAMENTO é a correlação não-fatorável (|Φ⁺⟩=(|00⟩+|11⟩)/√2), o recurso central. A entropia de von Neumann S(ρ)=−Tr(ρ log ρ) quantifica mistura; a entropia da densidade reduzida mede o emaranhamento de um estado puro bipartido. EPR 1935, Schrödinger ('Verschränkung'). Fato consolidado.","epistemico":"fato","exemplos":[],"id":"info_quantica_fundamentos","memoria":"semantica","meta":{"autor":"broca-so","dominio":"fisica","fonte":"EPR 1935; Schumacher 1995; Horodecki⁴ RMP 81,865 (2009)"},"origem":"wiki","palavras_chave":[],"sinonimos":[],"tipo":"conceito","titulo":"Informação quântica: qubit, emaranhamento, entropia"}
\section{Informação quântica: qubit, emaranhamento, entropia}
O qubit é um sistema de 2 níveis em \ensuremath{\mathbb{C}}\ensuremath{^{2}} (|\ensuremath{\psi}\ensuremath{\rangle}=\ensuremath{\alpha}|0\ensuremath{\rangle}+\ensuremath{\beta}|1\ensuremath{\rangle}, ponto na esfera de Bloch); o EMARANHAMENTO é a correlação não-fatorável (|\ensuremath{\Phi}\ensuremath{^{+}}\ensuremath{\rangle}=(|00\ensuremath{\rangle}+|11\ensuremath{\rangle})/\ensuremath{\sqrt{\,}}2), o recurso central. A entropia de von Neumann S(\ensuremath{\rho})=\ensuremath{-}Tr(\ensuremath{\rho} log \ensuremath{\rho}) quantifica mistura; a entropia da densidade reduzida mede o emaranhamento de um estado puro bipartido. EPR 1935, Schrödinger ('Verschränkung'). Fato consolidado.
\prov{relações: parte\_de cosmologia\_quantica\_hub}
\prov{proveniência: wiki \textperiodcentered{} EPR 1935; Schumacher 1995; Horodecki\ensuremath{^{4}} RMP 81,865 (2009) \textperiodcentered{} fato \textperiodcentered{} confiança 1.0 \textperiodcentered{} domínio fisica \textperiodcentered{} história 2 versões no jornal}

%CRISTAL {"arestas":[{"alvo":"area_render_gema","confianca":1.0,"epistemico":"fato","origem":"wiki","rel":"parte_de"},{"alvo":"transformada_metalica","confianca":1.0,"epistemico":"fato","origem":"wiki","rel":"usa"},{"alvo":"espectro_das_orbitas","confianca":1.0,"epistemico":"fato","origem":"wiki","rel":"relaciona"}],"confianca":1.0,"contraexemplos":[],"descricao":"A imagem é INGERIDA pelos invariantes da máquina — NÃO por FFT (que não é da máquina), mas pela TRANSFORMADA METÁLICA: o gato A_n=[[n,1],[1,0]] (o chicote, o tensor da régua) é o FILTRO, aplicado em multiescala; em cada escala a componente σ_n é a aproximação (o corpo) e o RASTRO (o detalhe ortogonal, o traço do transporte) é a estrutura. Os modos caem na parábola (cristal grosso / borda médio / caos fino) e o invariante é o λ metálico (a taxa de decaimento do rastro). linguagem/ingestao_mineral.py (3/3). O estudo revelou: o gabarito é cristal 71% / borda 22% / caos 7%; a nossa pedra é cristal 93% / borda 7% / caos 0,5% — LISA DEMAIS, falta aresta de faceta (borda) e detalhe fino (caos, o fogo/as inclusões). É a leitura NATIVA (o gato, não Fourier), rápida (sem render), que dá o alvo estrutural do próximo refino.","epistemico":"fato","exemplos":[],"id":"ingestao_transformada_metalica","memoria":"semantica","meta":{"dominio":"joalheria","fonte":"espectro orbitas"},"origem":"wiki","palavras_chave":[],"sinonimos":[],"tipo":"conceito","titulo":"A Ingestão pela Transformada Metálica (o rastro do gato, não Fourier)"}
\section{A Ingestão pela Transformada Metálica (o rastro do gato, não Fourier)}
A imagem é INGERIDA pelos invariantes da máquina — NÃO por FFT (que não é da máquina), mas pela TRANSFORMADA METÁLICA: o gato A\_n=[[n,1],[1,0]] (o chicote, o tensor da régua) é o FILTRO, aplicado em multiescala; em cada escala a componente \ensuremath{\sigma}\_n é a aproximação (o corpo) e o RASTRO (o detalhe ortogonal, o traço do transporte) é a estrutura. Os modos caem na parábola (cristal grosso / borda médio / caos fino) e o invariante é o \ensuremath{\lambda} metálico (a taxa de decaimento do rastro). linguagem/ingestao\_mineral.py (3/3). O estudo revelou: o gabarito é cristal 71\% / borda 22\% / caos 7\%; a nossa pedra é cristal 93\% / borda 7\% / caos 0,5\% — LISA DEMAIS, falta aresta de faceta (borda) e detalhe fino (caos, o fogo/as inclusões). É a leitura NATIVA (o gato, não Fourier), rápida (sem render), que dá o alvo estrutural do próximo refino.
\prov{relações: parte\_de area\_render\_gema; usa transformada\_metalica; relaciona espectro\_das\_orbitas}
\prov{proveniência: wiki \textperiodcentered{} espectro orbitas \textperiodcentered{} fato \textperiodcentered{} confiança 1.0 \textperiodcentered{} domínio joalheria \textperiodcentered{} história 20 versões no jornal}

%CRISTAL {"arestas":[{"alvo":"voxelizador_cristalografico","confianca":1.0,"epistemico":"fato","origem":"wiki","rel":"usa"},{"alvo":"difracao_atomica_turmalina","confianca":1.0,"epistemico":"fato","origem":"wiki","rel":"usa"},{"alvo":"ultra_realismo_gema","confianca":1.0,"epistemico":"fato","origem":"wiki","rel":"usa"},{"alvo":"joalheria_hub","confianca":1.0,"epistemico":"fato","origem":"wiki","rel":"parte_de"}],"confianca":1.0,"contraexemplos":[],"descricao":"A correção física do interior: a luz não fica na SUPERFÍCIE — ENTRA e COLIDE com os VOXELS (os pixels 3D) por dentro. OS RICOCHETES SÃO O RAIO COLIDINDO COM OS VOXELS, e isso dá o ESPALHAMENTO — o volume de Bragg que o CHICOTE faz (o campo_volumetrico do começo). Cada raio atravessa TODA a cena de voxels (a grade 3D da rede dentro da borda analítica): em cada voxel a luz ESPALHA (∝ densidade = o chicote, a estrutura/os veios) para o olho e ABSORVE por comprimento de onda (o cromóforo → a cor espectral). O zoneamento e as inclusões são REAIS (3D, com profundidade), não pintados na casca. A entrada é a borda analítica EXATA (o dupolinômio). linguagem/gema_cristal.grade_voxels (a grade 3D), primitivas_raytracer.render_gema_voxel, _amostra_grade (a colisão com o voxel). Une: voxelizador + campo_volumetrico + ricochetes + espalhamento.","epistemico":"fato","exemplos":[],"id":"interior_volumetrico_voxels","memoria":"semantica","meta":{"dominio":"joalheria","fonte":"maquina de corte"},"origem":"wiki","palavras_chave":[],"sinonimos":[],"tipo":"conceito","titulo":"O interior volumétrico: o raio colide com os voxels (os ricochetes = espalhamento)"}
\section{O interior volumétrico: o raio colide com os voxels (os ricochetes = espalhamento)}
A correção física do interior: a luz não fica na SUPERFÍCIE — ENTRA e COLIDE com os VOXELS (os pixels 3D) por dentro. OS RICOCHETES SÃO O RAIO COLIDINDO COM OS VOXELS, e isso dá o ESPALHAMENTO — o volume de Bragg que o CHICOTE faz (o campo\_volumetrico do começo). Cada raio atravessa TODA a cena de voxels (a grade 3D da rede dentro da borda analítica): em cada voxel a luz ESPALHA (\ensuremath{\propto} densidade = o chicote, a estrutura/os veios) para o olho e ABSORVE por comprimento de onda (o cromóforo \ensuremath{\to} a cor espectral). O zoneamento e as inclusões são REAIS (3D, com profundidade), não pintados na casca. A entrada é a borda analítica EXATA (o dupolinômio). linguagem/gema\_cristal.grade\_voxels (a grade 3D), primitivas\_raytracer.render\_gema\_voxel, \_amostra\_grade (a colisão com o voxel). Une: voxelizador + campo\_volumetrico + ricochetes + espalhamento.
\prov{relações: usa voxelizador\_cristalografico; usa difracao\_atomica\_turmalina; usa ultra\_realismo\_gema; parte\_de joalheria\_hub}
\prov{proveniência: wiki \textperiodcentered{} maquina de corte \textperiodcentered{} fato \textperiodcentered{} confiança 1.0 \textperiodcentered{} domínio joalheria \textperiodcentered{} história 105 versões no jornal}

%CRISTAL {"arestas":[{"alvo":"litografia","confianca":1.0,"epistemico":"fato","origem":"wiki","rel":"parte_de"},{"alvo":"invariante_estelar","confianca":1.0,"epistemico":"fato","origem":"wiki","rel":"relaciona"},{"alvo":"halt_topologico","confianca":1.0,"epistemico":"fato","origem":"wiki","rel":"relaciona"},{"alvo":"entropia_topologica","confianca":1.0,"epistemico":"fato","origem":"wiki","rel":"relaciona"}],"confianca":1.0,"contraexemplos":[],"descricao":"O invariante que separa a assinatura real da forja é o GRAU (o número de voltas, o winding) do colapso em torno de um conceito: assinatura REAL → tem preimagem → c²≥½ → GRAU 1 (enrola o conceito uma vez) → gravada; assinatura FORJADA (sem referente) → sem preimagem → c²<½ → GRAU 0 → cortada. O checksum a+c²=1 (invariante_estelar) reparte: c² é o lado com preimagem (a gravura), a é o gap/associador (o vazio da forja) — uma forja mora inteira em a. O grau é discreto (0 ou 1): não há meia-assinatura.","epistemico":"fato","exemplos":[],"id":"invariante_topologico_assinatura","memoria":"semantica","meta":{"dominio":"litografia","fonte":"litografia (joalheria)"},"origem":"wiki","palavras_chave":[],"sinonimos":[],"tipo":"conceito","titulo":"O invariante topológico da assinatura (o grau/winding)"}
\section{O invariante topológico da assinatura (o grau/winding)}
O invariante que separa a assinatura real da forja é o GRAU (o número de voltas, o winding) do colapso em torno de um conceito: assinatura REAL \ensuremath{\to} tem preimagem \ensuremath{\to} c\ensuremath{^{2}}\ensuremath{\ge}\ensuremath{\tfrac12} \ensuremath{\to} GRAU 1 (enrola o conceito uma vez) \ensuremath{\to} gravada; assinatura FORJADA (sem referente) \ensuremath{\to} sem preimagem \ensuremath{\to} c\ensuremath{^{2}}<\ensuremath{\tfrac12} \ensuremath{\to} GRAU 0 \ensuremath{\to} cortada. O checksum a+c\ensuremath{^{2}}=1 (invariante\_estelar) reparte: c\ensuremath{^{2}} é o lado com preimagem (a gravura), a é o gap/associador (o vazio da forja) — uma forja mora inteira em a. O grau é discreto (0 ou 1): não há meia-assinatura.
\prov{relações: parte\_de litografia; relaciona invariante\_estelar; relaciona halt\_topologico; relaciona entropia\_topologica}
\prov{proveniência: wiki \textperiodcentered{} litografia (joalheria) \textperiodcentered{} fato \textperiodcentered{} confiança 1.0 \textperiodcentered{} domínio litografia \textperiodcentered{} história 5 versões no jornal}

%CRISTAL {"arestas":[{"alvo":"isa_da_broca","confianca":1.0,"epistemico":"fato","origem":"wiki","rel":"especializa"},{"alvo":"metal_broca_so","confianca":1.0,"epistemico":"fato","origem":"wiki","rel":"parte_de"}],"confianca":1.0,"contraexemplos":[],"descricao":"O metal (instrucoes.h, enum OP_HALT..OP_SPECT, OP_COUNT) implementa 21 opcodes, todos com caso real no switch de instr_step, nenhum stub: HALT, LOAD, STORE, ADD, SUB, AND, OR, XOR, SILVER, GOLD, BRONZE, CMP, JMP, JZ, JNZ, FOLD, UNFOLD, PROJECT, LIFT, LOADS, SPECT. Correção: a doc dizia 12.","epistemico":"fato","exemplos":[],"id":"isa_metal_21_opcodes","memoria":"semantica","meta":{"autor":"broca-so","dominio":"metal","fonte":"ula/instrucoes.h; instrucoes.c:instr_step"},"origem":"wiki","palavras_chave":[],"sinonimos":[],"tipo":"conceito","titulo":"A ISA Real tem 21 Opcodes (não 12)"}
\section{A ISA Real tem 21 Opcodes (não 12)}
O metal (instrucoes.h, enum OP\_HALT..OP\_SPECT, OP\_COUNT) implementa 21 opcodes, todos com caso real no switch de instr\_step, nenhum stub: HALT, LOAD, STORE, ADD, SUB, AND, OR, XOR, SILVER, GOLD, BRONZE, CMP, JMP, JZ, JNZ, FOLD, UNFOLD, PROJECT, LIFT, LOADS, SPECT. Correção: a doc dizia 12.
\prov{relações: especializa isa\_da\_broca; parte\_de metal\_broca\_so}
\prov{proveniência: wiki \textperiodcentered{} ula/instrucoes.h; instrucoes.c:instr\_step \textperiodcentered{} fato \textperiodcentered{} confiança 1.0 \textperiodcentered{} domínio metal \textperiodcentered{} história 3 versões no jornal}

%CRISTAL {"arestas":[{"alvo":"numero_chern_tknn","confianca":1.0,"epistemico":"fato","origem":"wiki","rel":"especializa"},{"alvo":"fase_de_berry_canonica","confianca":1.0,"epistemico":"fato","origem":"wiki","rel":"usa"}],"confianca":1.0,"contraexemplos":[],"descricao":"Isolante no bulk, invariante sob reversão temporal, com estados de BORDA helicais (spin-momento travados) protegidos por simetria — o análogo T-reverso do Chern, classificado por um invariante Z2 (ν∈0,1). Kane-Mele 2005 (grafeno); Bernevig-Hughes-Zhang 2006 previram HgTe (inversão de banda acima de dc≈6,3 nm), confirmado por condutância de borda 2e²/h (König 2007).","epistemico":"fato","exemplos":[],"id":"isolantes_topologicos_z2","memoria":"semantica","meta":{"autor":"broca-so","dominio":"fisica","fonte":"Kane-Mele 2005 cond-mat/0506581; BHZ 2006; König 2007"},"origem":"wiki","palavras_chave":[],"sinonimos":[],"tipo":"conceito","titulo":"Isolantes topológicos (Hall de spin quântico, Z2)"}
\section{Isolantes topológicos (Hall de spin quântico, Z2)}
Isolante no bulk, invariante sob reversão temporal, com estados de BORDA helicais (spin-momento travados) protegidos por simetria — o análogo T-reverso do Chern, classificado por um invariante Z2 (\ensuremath{\nu}\ensuremath{\in}0,1). Kane-Mele 2005 (grafeno); Bernevig-Hughes-Zhang 2006 previram HgTe (inversão de banda acima de dc\ensuremath{\approx}6,3 nm), confirmado por condutância de borda 2e\ensuremath{^{2}}/h (König 2007).
\prov{relações: especializa numero\_chern\_tknn; usa fase\_de\_berry\_canonica}
\prov{proveniência: wiki \textperiodcentered{} Kane-Mele 2005 cond-mat/0506581; BHZ 2006; König 2007 \textperiodcentered{} fato \textperiodcentered{} confiança 1.0 \textperiodcentered{} domínio fisica \textperiodcentered{} história 4 versões no jornal}

%CRISTAL {"arestas":[{"alvo":"carga_q_n_sobre_3","confianca":1.0,"epistemico":"fato","origem":"wiki","rel":"usa"},{"alvo":"materia_modelo_padrao","confianca":1.0,"epistemico":"fato","origem":"wiki","rel":"explica"}],"confianca":1.0,"contraexemplos":[],"descricao":"ℂ⊗ℍ≅Cℓ(2): a escada η dá T₃=η†η−½∈±½ (dubleto). Gell-Mann-Nishijima Y=2(Q−T₃) DERIVA a hipercarga de cada partícula (YQL=⅓, YLL=−1, YuR=4/3, YdR=−⅔, YeR=−2). ℂ⊗(ℍ⊗𝕆) entrega SU(3)c×SU(2)L×U(1)Y completo (nos números quânticos; a teoria de campo que os realiza fica em aberto).","epistemico":"inferencia","exemplos":[],"id":"isospin_hipercarga_algebrico","memoria":"semantica","meta":{"autor":"broca-so","dominio":"cosmologia","fonte":"paper_4_particulas.tex §cxh-eletrofraco"},"origem":"wiki","palavras_chave":[],"sinonimos":[],"tipo":"conceito","titulo":"Isospin e hipercarga (ℂ⊗ℍ=Cℓ(2), Gell-Mann-Nishijima)"}
\section{Isospin e hipercarga (\ensuremath{\mathbb{C}}\ensuremath{\otimes}\ensuremath{\mathbb{H}}=C\ensuremath{\ell}(2), Gell-Mann-Nishijima)}
\ensuremath{\mathbb{C}}\ensuremath{\otimes}\ensuremath{\mathbb{H}}\ensuremath{\cong}C\ensuremath{\ell}(2): a escada \ensuremath{\eta} dá T\ensuremath{_{3}}=\ensuremath{\eta}\ensuremath{\dagger}\ensuremath{\eta}\ensuremath{-}\ensuremath{\tfrac12}\ensuremath{\in}$\pm$\ensuremath{\tfrac12} (dubleto). Gell-Mann-Nishijima Y=2(Q\ensuremath{-}T\ensuremath{_{3}}) DERIVA a hipercarga de cada partícula (YQL=\ensuremath{\tfrac13}, YLL=\ensuremath{-}1, YuR=4/3, YdR=\ensuremath{-}\ensuremath{\tfrac23}, YeR=\ensuremath{-}2). \ensuremath{\mathbb{C}}\ensuremath{\otimes}(\ensuremath{\mathbb{H}}\ensuremath{\otimes}\ensuremath{\mathbb{O}}) entrega SU(3)c$\times$SU(2)L$\times$U(1)Y completo (nos números quânticos; a teoria de campo que os realiza fica em aberto).
\prov{relações: usa carga\_q\_n\_sobre\_3; explica materia\_modelo\_padrao}
\prov{proveniência: wiki \textperiodcentered{} paper\_4\_particulas.tex §cxh-eletrofraco \textperiodcentered{} inferencia \textperiodcentered{} confiança 1.0 \textperiodcentered{} domínio cosmologia \textperiodcentered{} história 4 versões no jornal}

%CRISTAL {"arestas":[{"alvo":"joalheria_hub","confianca":1.0,"epistemico":"fato","origem":"wiki","rel":"parte_de"},{"alvo":"cristalchain","confianca":1.0,"epistemico":"fato","origem":"wiki","rel":"parte_de"},{"alvo":"tesouraria","confianca":1.0,"epistemico":"fato","origem":"wiki","rel":"relaciona"},{"alvo":"refino_continuo_floco_metais","confianca":1.0,"epistemico":"fato","origem":"wiki","rel":"relaciona"},{"alvo":"mineracao_conservacao_verificavel","confianca":1.0,"epistemico":"fato","origem":"wiki","rel":"relaciona"},{"alvo":"emissao_turmalina_piso","confianca":1.0,"epistemico":"fato","origem":"wiki","rel":"usa"}],"confianca":1.0,"contraexemplos":[],"descricao":"A CristalChain tem duas artes que fecham o ciclo do valor. A MINERAÇÃO refina o BRUTO: o floco (o pacote PoW nativo) é dividido nos 12 metais, com a conservação verificada (a norma σ·conj=−1). Bruto → refinado. A JOALHERIA lapida o REFINADO em VALOR: a pedra (a Turmalina Paraíba) é cortada (trilhão 3-fold), iluminada e ASSENTADA na âncora — o piso 1 ct = R$ 1.000,00, a emissão. Refinado → joia → valor. As duas COMPLETAM o pipe: mineração (bruto→refinado) + joalheria (refinado→valor). E a CORRESPONDÊNCIA tesouraria↔joalheria: a JOIA É O TESOURO — a joalheria (fazer a joia) é a face-artesã da tesouraria (guardar e emitir o tesouro). Na máquina, o ENGASTE (assentar a pedra) = o COLAPSO = a liquidação/aprovação da tesouraria; o QUILATE = a RETA = o piso (a medida do valor). Por isso a joalheria é uma FERRAMENTA CENTRAL da cadeia, ao lado da mineração — não um ornamento, mas a arte que fixa o valor.","epistemico":"fato","exemplos":[],"id":"joalheria_ferramenta_central","memoria":"semantica","meta":{"dominio":"joalheria","fonte":"joalheria hub"},"origem":"wiki","palavras_chave":[],"sinonimos":[],"tipo":"conceito","titulo":"A Joalheria é a Ferramenta Central da CristalChain (completa a mineração)"}
\section{A Joalheria é a Ferramenta Central da CristalChain (completa a mineração)}
A CristalChain tem duas artes que fecham o ciclo do valor. A MINERAÇÃO refina o BRUTO: o floco (o pacote PoW nativo) é dividido nos 12 metais, com a conservação verificada (a norma \ensuremath{\sigma}\textperiodcentered conj=\ensuremath{-}1). Bruto \ensuremath{\to} refinado. A JOALHERIA lapida o REFINADO em VALOR: a pedra (a Turmalina Paraíba) é cortada (trilhão 3-fold), iluminada e ASSENTADA na âncora — o piso 1 ct = R\$ 1.000,00, a emissão. Refinado \ensuremath{\to} joia \ensuremath{\to} valor. As duas COMPLETAM o pipe: mineração (bruto\ensuremath{\to}refinado) + joalheria (refinado\ensuremath{\to}valor). E a CORRESPONDÊNCIA tesouraria\ensuremath{\leftrightarrow}joalheria: a JOIA É O TESOURO — a joalheria (fazer a joia) é a face-artesã da tesouraria (guardar e emitir o tesouro). Na máquina, o ENGASTE (assentar a pedra) = o COLAPSO = a liquidação/aprovação da tesouraria; o QUILATE = a RETA = o piso (a medida do valor). Por isso a joalheria é uma FERRAMENTA CENTRAL da cadeia, ao lado da mineração — não um ornamento, mas a arte que fixa o valor.
\prov{relações: parte\_de joalheria\_hub; parte\_de cristalchain; relaciona tesouraria; relaciona refino\_continuo\_floco\_metais; relaciona mineracao\_conservacao\_verificavel; usa emissao\_turmalina\_piso}
\prov{proveniência: wiki \textperiodcentered{} joalheria hub \textperiodcentered{} fato \textperiodcentered{} confiança 1.0 \textperiodcentered{} domínio joalheria \textperiodcentered{} história 14 versões no jornal}

%CRISTAL {"arestas":[{"alvo":"cristalchain","confianca":1.0,"epistemico":"fato","origem":"wiki","rel":"parte_de"}],"confianca":1.0,"contraexemplos":[],"descricao":"O núcleo único das técnicas da pedra da CristalChain — a fonte de verdade a que os outros documentos apontam em vez de repetir. Reúne cinco frentes lidas por UMA máquina: GEMOLOGIA (a pedra), ÓPTICA (a luz na pedra), METALURGIA (o metal que colore e opera), RENDER (a pedra virando imagem) e DESIGN (a pedra virando identidade). Documento VIVO — cresce a cada refino. papers/joalheria.tex; linguagem/joalheria.py (7/7). O fio é o dicionário das áreas conforme o BAI.","epistemico":"fato","exemplos":[],"id":"joalheria_hub","memoria":"semantica","meta":{"dominio":"joalheria","fonte":"joalheria hub"},"origem":"wiki","palavras_chave":[],"sinonimos":[],"tipo":"conceito","titulo":"Joalheria (o hub vivo da arte da pedra)"}
\section{Joalheria (o hub vivo da arte da pedra)}
O núcleo único das técnicas da pedra da CristalChain — a fonte de verdade a que os outros documentos apontam em vez de repetir. Reúne cinco frentes lidas por UMA máquina: GEMOLOGIA (a pedra), ÓPTICA (a luz na pedra), METALURGIA (o metal que colore e opera), RENDER (a pedra virando imagem) e DESIGN (a pedra virando identidade). Documento VIVO — cresce a cada refino. papers/joalheria.tex; linguagem/joalheria.py (7/7). O fio é o dicionário das áreas conforme o BAI.
\prov{relações: parte\_de cristalchain}
\prov{proveniência: wiki \textperiodcentered{} joalheria hub \textperiodcentered{} fato \textperiodcentered{} confiança 1.0 \textperiodcentered{} domínio joalheria \textperiodcentered{} história 5 versões no jornal}

%CRISTAL {"arestas":[{"alvo":"grupos_excepcionais_jordan","confianca":1.0,"epistemico":"fato","origem":"wiki","rel":"especializa"},{"alvo":"g2_automorfismos_associador","confianca":1.0,"epistemico":"fato","origem":"wiki","rel":"usa"},{"alvo":"os_dois_tres_cores_hurwitz_e_gerac","confianca":1.0,"epistemico":"fato","origem":"wiki","rel":"explica"}],"confianca":1.0,"contraexemplos":[],"descricao":"Matrizes hermitianas n×n sobre 𝕆 formam álgebra de Jordan sse n≤3 (Jordan-von Neumann-Wigner); a última é a de Albert J₃(𝕆) (dim 27 = o 27 de E₆, Aut=F₄). Peirce: J₃(𝕆)=ℝ³⊕𝕆³ — três idempotentes primitivos = gerações, três octônios fora-diagonal as conectam (uma cópia de cor cada). O 'gêmeo de Hurwitz': 𝕆 última de divisão → 3 cores; J₃(𝕆) última de Jordan → ≤3 gerações. Teto ≤3 é teorema; a ligação com gerações físicas é programa (E₇/E₈ NÃO aparecem nesta série).","epistemico":"inferencia","exemplos":[],"id":"jordan_albert_geracoes","memoria":"semantica","meta":{"autor":"broca-so","dominio":"cosmologia","fonte":"paper_4_particulas.tex obs:geracoes"},"origem":"wiki","palavras_chave":[],"sinonimos":[],"tipo":"conceito","titulo":"A Jordan excepcional J₃(𝕆)/Albert e as ≤3 gerações"}
\section{A Jordan excepcional J\ensuremath{_{3}}(\ensuremath{\mathbb{O}})/Albert e as \ensuremath{\le}3 gerações}
Matrizes hermitianas n$\times$n sobre \ensuremath{\mathbb{O}} formam álgebra de Jordan sse n\ensuremath{\le}3 (Jordan-von Neumann-Wigner); a última é a de Albert J\ensuremath{_{3}}(\ensuremath{\mathbb{O}}) (dim 27 = o 27 de E\ensuremath{_{6}}, Aut=F\ensuremath{_{4}}). Peirce: J\ensuremath{_{3}}(\ensuremath{\mathbb{O}})=\ensuremath{\mathbb{R}}\ensuremath{^{3}}\ensuremath{\oplus}\ensuremath{\mathbb{O}}\ensuremath{^{3}} — três idempotentes primitivos = gerações, três octônios fora-diagonal as conectam (uma cópia de cor cada). O 'gêmeo de Hurwitz': \ensuremath{\mathbb{O}} última de divisão \ensuremath{\to} 3 cores; J\ensuremath{_{3}}(\ensuremath{\mathbb{O}}) última de Jordan \ensuremath{\to} \ensuremath{\le}3 gerações. Teto \ensuremath{\le}3 é teorema; a ligação com gerações físicas é programa (E\ensuremath{_{7}}/E\ensuremath{_{8}} NÃO aparecem nesta série).
\prov{relações: especializa grupos\_excepcionais\_jordan; usa g2\_automorfismos\_associador; explica os\_dois\_tres\_cores\_hurwitz\_e\_gerac}
\prov{proveniência: wiki \textperiodcentered{} paper\_4\_particulas.tex obs:geracoes \textperiodcentered{} inferencia \textperiodcentered{} confiança 1.0 \textperiodcentered{} domínio cosmologia \textperiodcentered{} história 4 versões no jornal}

%CRISTAL {"arestas":[{"alvo":"matriz_ckm","confianca":1.0,"epistemico":"fato","origem":"wiki","rel":"parte_de"}],"confianca":1.0,"contraexemplos":[],"descricao":"Com 2 gerações a matriz de mistura é real (sem CP); com 3 gerações sobra uma FASE complexa irredutível — logo a violação de CP é consequência NECESSÁRIA de existirem 3 famílias. Previsão feita (1973) antes da descoberta do charm, bottom e top. Contagem N=3: 3 ângulos + 1 fase física; δ≠0,π ⇒ CP violado. FATO confirmado — Nobel 2008.","epistemico":"fato","exemplos":[],"id":"km_previsao_cp","memoria":"semantica","meta":{"autor":"broca-so","dominio":"fisica","fonte":"Kobayashi-Maskawa 1973 Prog.Theor.Phys. 49,652"},"origem":"wiki","palavras_chave":[],"sinonimos":[],"tipo":"conceito","titulo":"A previsão de Kobayashi-Maskawa: 3 gerações ⇒ violação de CP"}
\section{A previsão de Kobayashi-Maskawa: 3 gerações \ensuremath{\Rightarrow} violação de CP}
Com 2 gerações a matriz de mistura é real (sem CP); com 3 gerações sobra uma FASE complexa irredutível — logo a violação de CP é consequência NECESSÁRIA de existirem 3 famílias. Previsão feita (1973) antes da descoberta do charm, bottom e top. Contagem N=3: 3 ângulos + 1 fase física; \ensuremath{\delta}\ensuremath{\ne}0,\ensuremath{\pi} \ensuremath{\Rightarrow} CP violado. FATO confirmado — Nobel 2008.
\prov{relações: parte\_de matriz\_ckm}
\prov{proveniência: wiki \textperiodcentered{} Kobayashi-Maskawa 1973 Prog.Theor.Phys. 49,652 \textperiodcentered{} fato \textperiodcentered{} confiança 1.0 \textperiodcentered{} domínio fisica \textperiodcentered{} história 2 versões no jornal}

%CRISTAL {"arestas":[{"alvo":"koide_z3_geracoes","confianca":1.0,"epistemico":"fato","origem":"wiki","rel":"especializa"},{"alvo":"matrizes_massa_circulantes","confianca":1.0,"epistemico":"fato","origem":"wiki","rel":"usa"}],"confianca":1.0,"contraexemplos":[],"descricao":"A parametrização Z₃ √m_k=m₀(1+√2·cos(δ+2πk/3)) produz automaticamente Koide Q=(Σm)/(Σ√m)²=⅔. Léptons carregados dão Q=0,666661 (a 0,001% de ⅔) — a impressão digital Z₃. Já os QUARKS dão 0,85/0,73 e os neutrinos 0,33-0,57: só os léptons carregados (Dirac, sem cor) exibem a quebra Z₃ pura. [D] a identidade 'parametrização Z₃ = espectro de circulante hermitiana' (verif. 10⁻¹⁵); [N] os valores; o ⅔ em si é [J] (o encontro do '3' da trialidade com o '⅔' empírico).","epistemico":"inferencia","exemplos":[],"id":"koide_quarks","memoria":"semantica","meta":{"autor":"broca-so","dominio":"cosmologia","fonte":"paper_4_particulas.tex obs:geracoes"},"origem":"wiki","palavras_chave":[],"sinonimos":[],"tipo":"conceito","titulo":"Koide para quarks (0,85/0,73) — sem a Z₃ pura"}
\section{Koide para quarks (0,85/0,73) — sem a Z\ensuremath{_{3}} pura}
A parametrização Z\ensuremath{_{3}} \ensuremath{\sqrt{\,}}m\_k=m\ensuremath{_{0}}(1+\ensuremath{\sqrt{\,}}2\textperiodcentered cos(\ensuremath{\delta}+2\ensuremath{\pi}k/3)) produz automaticamente Koide Q=(\ensuremath{\Sigma}m)/(\ensuremath{\Sigma}\ensuremath{\sqrt{\,}}m)\ensuremath{^{2}}=\ensuremath{\tfrac23}. Léptons carregados dão Q=0,666661 (a 0,001\% de \ensuremath{\tfrac23}) — a impressão digital Z\ensuremath{_{3}}. Já os QUARKS dão 0,85/0,73 e os neutrinos 0,33-0,57: só os léptons carregados (Dirac, sem cor) exibem a quebra Z\ensuremath{_{3}} pura. [D] a identidade 'parametrização Z\ensuremath{_{3}} = espectro de circulante hermitiana' (verif. 10\ensuremath{^{-15}}); [N] os valores; o \ensuremath{\tfrac23} em si é [J] (o encontro do '3' da trialidade com o '\ensuremath{\tfrac23}' empírico).
\prov{relações: especializa koide\_z3\_geracoes; usa matrizes\_massa\_circulantes}
\prov{proveniência: wiki \textperiodcentered{} paper\_4\_particulas.tex obs:geracoes \textperiodcentered{} inferencia \textperiodcentered{} confiança 1.0 \textperiodcentered{} domínio cosmologia \textperiodcentered{} história 3 versões no jornal}

%CRISTAL {"arestas":[{"alvo":"materia_modelo_padrao","confianca":1.0,"epistemico":"fato","origem":"wiki","rel":"relaciona"},{"alvo":"os_dois_tres_cores_hurwitz_e_gerac","confianca":1.0,"epistemico":"fato","origem":"wiki","rel":"usa"}],"confianca":1.0,"contraexemplos":[],"descricao":"A quebra Z₃ da trialidade √m_k=m₀(1+√2·cos(δ+2πk/3)) produz automaticamente o Koide Q=⅔. Léptons carregados: Q=0,666661 (0,001% de ⅔); neutrinos Q∈[0,33;0,57] e quarks 0,85/0,73 — só léptons carregados exibem a Z₃ pura. O '3' é da trialidade [D], o '2/3' é empírico [J]; ligá-los é a hipótese falseável.","epistemico":"hipotese","exemplos":[],"id":"koide_z3_geracoes","memoria":"semantica","meta":{"autor":"broca-so","dominio":"cosmologia","fonte":"paper_4_particulas.tex §geracoes"},"origem":"wiki","palavras_chave":[],"sinonimos":[],"tipo":"conceito","titulo":"A quebra Z₃ da trialidade e a fórmula de Koide"}
\section{A quebra Z\ensuremath{_{3}} da trialidade e a fórmula de Koide}
A quebra Z\ensuremath{_{3}} da trialidade \ensuremath{\sqrt{\,}}m\_k=m\ensuremath{_{0}}(1+\ensuremath{\sqrt{\,}}2\textperiodcentered cos(\ensuremath{\delta}+2\ensuremath{\pi}k/3)) produz automaticamente o Koide Q=\ensuremath{\tfrac23}. Léptons carregados: Q=0,666661 (0,001\% de \ensuremath{\tfrac23}); neutrinos Q\ensuremath{\in}[0,33;0,57] e quarks 0,85/0,73 — só léptons carregados exibem a Z\ensuremath{_{3}} pura. O '3' é da trialidade [D], o '2/3' é empírico [J]; ligá-los é a hipótese falseável.
\prov{relações: relaciona materia\_modelo\_padrao; usa os\_dois\_tres\_cores\_hurwitz\_e\_gerac}
\prov{proveniência: wiki \textperiodcentered{} paper\_4\_particulas.tex §geracoes \textperiodcentered{} hipotese \textperiodcentered{} confiança 1.0 \textperiodcentered{} domínio cosmologia \textperiodcentered{} história 3 versões no jornal}

%CRISTAL {"arestas":[{"alvo":"equacoes_friedmann","confianca":1.0,"epistemico":"fato","origem":"wiki","rel":"usa"},{"alvo":"lambda_cosmologica_saga","confianca":1.0,"epistemico":"fato","origem":"wiki","rel":"usa"},{"alvo":"cmb_cobe_wmap_planck","confianca":1.0,"epistemico":"fato","origem":"wiki","rel":"usa"},{"alvo":"bao_regua_padrao","confianca":1.0,"epistemico":"fato","origem":"wiki","rel":"usa"},{"alvo":"sne_ia_pantheon","confianca":1.0,"epistemico":"fato","origem":"wiki","rel":"usa"},{"alvo":"inflacao_cosmica","confianca":1.0,"epistemico":"fato","origem":"wiki","rel":"usa"},{"alvo":"bai_limite_nao_dissipativo","confianca":1.0,"epistemico":"fato","origem":"wiki","rel":"relaciona"}],"confianca":1.0,"contraexemplos":[],"descricao":"Modelo padrão de 6 parâmetros: plano, Λ (energia escura w=−1) + matéria escura fria + bárions, perturbações adiabáticas. Planck 2018: Ωm=0,315±0,007, ΩΛ≈0,685, H₀=67,4±0,5 km/s/Mpc, σ₈=0,811, idade≈13,8 Gyr. Melhor ajuste global de baixo custo paramétrico; não é teoria fundamental de Λ.","epistemico":"fato","exemplos":[],"id":"lambda_cdm","memoria":"semantica","meta":{"autor":"broca-so","dominio":"cosmologia_dados","fonte":"Planck 2018 (arXiv:1807.06209)"},"origem":"wiki","palavras_chave":[],"sinonimos":[],"tipo":"conceito","titulo":"Modelo ΛCDM (concordância) — Planck 2018"}
\section{Modelo \ensuremath{\Lambda}CDM (concordância) — Planck 2018}
Modelo padrão de 6 parâmetros: plano, \ensuremath{\Lambda} (energia escura w=\ensuremath{-}1) + matéria escura fria + bárions, perturbações adiabáticas. Planck 2018: \ensuremath{\Omega}m=0,315$\pm$0,007, \ensuremath{\Omega}\ensuremath{\Lambda}\ensuremath{\approx}0,685, H\ensuremath{_{0}}=67,4$\pm$0,5 km/s/Mpc, \ensuremath{\sigma}\ensuremath{_{8}}=0,811, idade\ensuremath{\approx}13,8 Gyr. Melhor ajuste global de baixo custo paramétrico; não é teoria fundamental de \ensuremath{\Lambda}.
\prov{relações: usa equacoes\_friedmann; usa lambda\_cosmologica\_saga; usa cmb\_cobe\_wmap\_planck; usa bao\_regua\_padrao; usa sne\_ia\_pantheon; usa inflacao\_cosmica; relaciona bai\_limite\_nao\_dissipativo}
\prov{proveniência: wiki \textperiodcentered{} Planck 2018 (arXiv:1807.06209) \textperiodcentered{} fato \textperiodcentered{} confiança 1.0 \textperiodcentered{} domínio cosmologia\_dados \textperiodcentered{} história 8 versões no jornal}

%CRISTAL {"arestas":[{"alvo":"problema_constante_cosmologica","confianca":1.0,"epistemico":"fato","origem":"wiki","rel":"relaciona"},{"alvo":"cosmologia_quantica_hub","confianca":1.0,"epistemico":"fato","origem":"wiki","rel":"parte_de"}],"confianca":1.0,"contraexemplos":[],"descricao":"Λ g_μν, introduzida por Einstein (1917) para um universo estático (instável, Eddington 1930). Abandonada após a expansão de Hubble (1929) — o 'maior erro' (atribuição de Gamow, não documentada). Ressuscitada em 1998 como energia escura (~68% do universo). Existência de Λ≠0 = fato observacional; NATUREZA = aberta.","epistemico":"fato","exemplos":[],"id":"lambda_cosmologica_saga","memoria":"semantica","meta":{"autor":"broca-so","dominio":"fisica","fonte":"Einstein 1917; O'Raifeartaigh+2017 arXiv:1711.06890"},"origem":"wiki","palavras_chave":[],"sinonimos":[],"tipo":"conceito","titulo":"A constante cosmológica Λ — do 'maior erro' à energia escura"}
\section{A constante cosmológica \ensuremath{\Lambda} — do 'maior erro' à energia escura}
\ensuremath{\Lambda} g\_\ensuremath{\mu}\ensuremath{\nu}, introduzida por Einstein (1917) para um universo estático (instável, Eddington 1930). Abandonada após a expansão de Hubble (1929) — o 'maior erro' (atribuição de Gamow, não documentada). Ressuscitada em 1998 como energia escura (\~{}68\% do universo). Existência de \ensuremath{\Lambda}\ensuremath{\ne}0 = fato observacional; NATUREZA = aberta.
\prov{relações: relaciona problema\_constante\_cosmologica; parte\_de cosmologia\_quantica\_hub}
\prov{proveniência: wiki \textperiodcentered{} Einstein 1917; O'Raifeartaigh+2017 arXiv:1711.06890 \textperiodcentered{} fato \textperiodcentered{} confiança 1.0 \textperiodcentered{} domínio fisica \textperiodcentered{} história 3 versões no jornal}

%CRISTAL {"arestas":[{"alvo":"bekenstein_hawking_entropia","confianca":1.0,"epistemico":"fato","origem":"wiki","rel":"usa"},{"alvo":"cosmologia_quantica_hub","confianca":1.0,"epistemico":"fato","origem":"wiki","rel":"parte_de"}],"confianca":1.0,"contraexemplos":[],"descricao":"A ponte informação↔termodinâmica: LANDAUER — apagar 1 bit a T custa ≥k_BT ln2 (≈2,9×10⁻²¹ J a 300K, verificado, Bérut 2012). GSL — S_matéria+A/4 nunca decresce. LIMITE DE BEKENSTEIN — teto de entropia numa região: S≤2πk_BRE/ℏc (formalizado em QFT por Casini 2008). Fatos consolidados nos seus domínios.","epistemico":"fato","exemplos":[],"id":"landauer_gsl_bekenstein_bound","memoria":"semantica","meta":{"autor":"broca-so","dominio":"fisica","fonte":"Landauer 1961; Bekenstein 1974/1981; Casini 2008 arXiv:0804.2182"},"origem":"wiki","palavras_chave":[],"sinonimos":[],"tipo":"conceito","titulo":"Landauer, 2ª lei generalizada e limite de Bekenstein"}
\section{Landauer, 2ª lei generalizada e limite de Bekenstein}
A ponte informação\ensuremath{\leftrightarrow}termodinâmica: LANDAUER — apagar 1 bit a T custa \ensuremath{\ge}k\_BT ln2 (\ensuremath{\approx}2,9$\times$10\ensuremath{^{-21}} J a 300K, verificado, Bérut 2012). GSL — S\_matéria+A/4 nunca decresce. LIMITE DE BEKENSTEIN — teto de entropia numa região: S\ensuremath{\le}2\ensuremath{\pi}k\_BRE/\ensuremath{\hbar}c (formalizado em QFT por Casini 2008). Fatos consolidados nos seus domínios.
\prov{relações: usa bekenstein\_hawking\_entropia; parte\_de cosmologia\_quantica\_hub}
\prov{proveniência: wiki \textperiodcentered{} Landauer 1961; Bekenstein 1974/1981; Casini 2008 arXiv:0804.2182 \textperiodcentered{} fato \textperiodcentered{} confiança 1.0 \textperiodcentered{} domínio fisica \textperiodcentered{} história 3 versões no jornal}

%CRISTAL {"arestas":[{"alvo":"teoria_cordas_m","confianca":1.0,"epistemico":"fato","origem":"wiki","rel":"parte_de"}],"confianca":1.0,"contraexemplos":[],"descricao":"O vastíssimo conjunto de vácuos da teoria de cordas (compactificações de Calabi-Yau + fluxos), ~10⁵⁰⁰ (estimativas recentes muito maiores), associado a raciocínio antrópico/multiverso. CONTROVÉRSIA epistemológica: o número enorme de vácuos é apontado como perda de poder preditivo (falseabilidade).","epistemico":"hipotese","exemplos":[],"id":"landscape_cordas","memoria":"semantica","meta":{"autor":"broca-so","dominio":"fisica","fonte":"Susskind 2003 hep-th/0302219; Bousso-Polchinski 2000"},"origem":"wiki","palavras_chave":[],"sinonimos":[],"tipo":"conceito","titulo":"O landscape de cordas"}
\section{O landscape de cordas}
O vastíssimo conjunto de vácuos da teoria de cordas (compactificações de Calabi-Yau + fluxos), \~{}10\ensuremath{^{500}} (estimativas recentes muito maiores), associado a raciocínio antrópico/multiverso. CONTROVÉRSIA epistemológica: o número enorme de vácuos é apontado como perda de poder preditivo (falseabilidade).
\prov{relações: parte\_de teoria\_cordas\_m}
\prov{proveniência: wiki \textperiodcentered{} Susskind 2003 hep-th/0302219; Bousso-Polchinski 2000 \textperiodcentered{} hipotese \textperiodcentered{} confiança 1.0 \textperiodcentered{} domínio fisica \textperiodcentered{} história 2 versões no jornal}

%CRISTAL {"arestas":[{"alvo":"cosmologia_quantica","confianca":1.0,"epistemico":"fato","origem":"wiki","rel":"parte_de"},{"alvo":"desitter_4d_universo","confianca":1.0,"epistemico":"fato","origem":"wiki","rel":"generaliza"},{"alvo":"regua_de_gentil","confianca":1.0,"epistemico":"fato","origem":"wiki","rel":"usa"}],"confianca":1.0,"contraexemplos":[],"descricao":"Na métrica ds²=−N(ψ)²dθ²+dψ², exigir curvatura constante (máxima simetria) FORÇA N²=cos²ψ=a e R=2 — o de Sitter 2D. Os quaternions ψ=±π/2 são os horizontes (o lapse zera); a ilha |s|<1 é o patch estático. Teorema com prova (Chebyshev/Ricci), verificado.","epistemico":"fato","exemplos":[],"id":"lapse_selecao_desitter","memoria":"semantica","meta":{"autor":"broca-so","dominio":"cosmologia","fonte":"paper_0_sintese.tex thm:dS; paper_2_cosmologia.tex"},"origem":"wiki","palavras_chave":[],"sinonimos":[],"tipo":"conceito","titulo":"Seleção do lapse: máxima simetria ⇒ N²=a, R=2 (dS₂)"}
\section{Seleção do lapse: máxima simetria \ensuremath{\Rightarrow} N\ensuremath{^{2}}=a, R=2 (dS\ensuremath{_{2}})}
Na métrica ds\ensuremath{^{2}}=\ensuremath{-}N(\ensuremath{\psi})\ensuremath{^{2}}d\ensuremath{\theta}\ensuremath{^{2}}+d\ensuremath{\psi}\ensuremath{^{2}}, exigir curvatura constante (máxima simetria) FORÇA N\ensuremath{^{2}}=cos\ensuremath{^{2}}\ensuremath{\psi}=a e R=2 — o de Sitter 2D. Os quaternions \ensuremath{\psi}=$\pm$\ensuremath{\pi}/2 são os horizontes (o lapse zera); a ilha |s|<1 é o patch estático. Teorema com prova (Chebyshev/Ricci), verificado.
\prov{relações: parte\_de cosmologia\_quantica; generaliza desitter\_4d\_universo; usa regua\_de\_gentil}
\prov{proveniência: wiki \textperiodcentered{} paper\_0\_sintese.tex thm:dS; paper\_2\_cosmologia.tex \textperiodcentered{} fato \textperiodcentered{} confiança 1.0 \textperiodcentered{} domínio cosmologia \textperiodcentered{} história 4 versões no jornal}

%CRISTAL {"arestas":[{"alvo":"cosmologia_quantica_hub","confianca":1.0,"epistemico":"fato","origem":"wiki","rel":"parte_de"}],"confianca":1.0,"contraexemplos":[],"descricao":"A gravidade curva a luz; a RG prevê deflexão pelo Sol de 1,75'' (o dobro do newtoniano). Confirmada no eclipse de 29 mai 1919 (Dyson-Eddington-Davidson) — o teste que consagrou a RG. Base da lente gravitacional moderna.","epistemico":"fato","exemplos":[],"id":"lente_gravitacional_1919","memoria":"semantica","meta":{"autor":"broca-so","dominio":"fisica","fonte":"Einstein 1915/16; Dyson-Eddington-Davidson 1919"},"origem":"wiki","palavras_chave":[],"sinonimos":[],"tipo":"conceito","titulo":"Lente gravitacional — o eclipse de 1919"}
\section{Lente gravitacional — o eclipse de 1919}
A gravidade curva a luz; a RG prevê deflexão pelo Sol de 1,75'' (o dobro do newtoniano). Confirmada no eclipse de 29 mai 1919 (Dyson-Eddington-Davidson) — o teste que consagrou a RG. Base da lente gravitacional moderna.
\prov{relações: parte\_de cosmologia\_quantica\_hub}
\prov{proveniência: wiki \textperiodcentered{} Einstein 1915/16; Dyson-Eddington-Davidson 1919 \textperiodcentered{} fato \textperiodcentered{} confiança 1.0 \textperiodcentered{} domínio fisica \textperiodcentered{} história 2 versões no jornal}

%CRISTAL {"arestas":[{"alvo":"neuron_io_popcnt","confianca":1.0,"epistemico":"fato","origem":"wiki","rel":"usa"},{"alvo":"arquitetura_von_neumann","confianca":1.0,"epistemico":"fato","origem":"wiki","rel":"relaciona"}],"confianca":1.0,"contraexemplos":[],"descricao":"Biblioteca .so que expõe BrocaMaquina + broca_so_exec: uma camada de operações (GATO_STREAM, TORO, REGISTRO_N, COLAPSO_6, BANDA, NONCE, POW_SEED) despachada a um BrocaBackend. Distinta da CPU-ISA (não executa bytecode). O backend CPU implementa GATO_STREAM de fato; os demais delegam a broca_pow_exec. Compila com -lcrypto.","epistemico":"fato","exemplos":[],"id":"libbroca_so_camada","memoria":"semantica","meta":{"autor":"broca-so","dominio":"metal","fonte":"ula/broca_so.c; broca_backend_cpu.c; Makefile"},"origem":"wiki","palavras_chave":[],"sinonimos":[],"tipo":"conceito","titulo":"libbroca_so.so — a Máquina de Trabalho Isomórfica"}
\section{libbroca\_so.so — a Máquina de Trabalho Isomórfica}
Biblioteca .so que expõe BrocaMaquina + broca\_so\_exec: uma camada de operações (GATO\_STREAM, TORO, REGISTRO\_N, COLAPSO\_6, BANDA, NONCE, POW\_SEED) despachada a um BrocaBackend. Distinta da CPU-ISA (não executa bytecode). O backend CPU implementa GATO\_STREAM de fato; os demais delegam a broca\_pow\_exec. Compila com -lcrypto.
\prov{relações: usa neuron\_io\_popcnt; relaciona arquitetura\_von\_neumann}
\prov{proveniência: wiki \textperiodcentered{} ula/broca\_so.c; broca\_backend\_cpu.c; Makefile \textperiodcentered{} fato \textperiodcentered{} confiança 1.0 \textperiodcentered{} domínio metal \textperiodcentered{} história 3 versões no jornal}

%CRISTAL {"arestas":[{"alvo":"confinamento_cor","confianca":1.0,"epistemico":"fato","origem":"wiki","rel":"relaciona"},{"alvo":"qcd_cromodinamica","confianca":1.0,"epistemico":"fato","origem":"wiki","rel":"parte_de"}],"confianca":1.0,"contraexemplos":[],"descricao":"O acoplamento forte DIMINUI a altas energias/curtas distâncias — quarks quase livres quando próximos. Função beta negativa: β(g)=−(b₀/16π²)g³, b₀=11−(2/3)n_f>0. Λ_QCD≈0,2 GeV. Gross-Wilczek e Politzer 1973 (Nobel 2004); confirmada pelo running de α_s.","epistemico":"fato","exemplos":[],"id":"liberdade_assintotica","memoria":"semantica","meta":{"autor":"broca-so","dominio":"fisica","fonte":"Gross-Wilczek, Politzer 1973 PRL 30; Nobel 2004"},"origem":"wiki","palavras_chave":[],"sinonimos":[],"tipo":"conceito","titulo":"Liberdade assintótica (Nobel 2004)"}
\section{Liberdade assintótica (Nobel 2004)}
O acoplamento forte DIMINUI a altas energias/curtas distâncias — quarks quase livres quando próximos. Função beta negativa: \ensuremath{\beta}(g)=\ensuremath{-}(b\ensuremath{_{0}}/16\ensuremath{\pi}\ensuremath{^{2}})g\ensuremath{^{3}}, b\ensuremath{_{0}}=11\ensuremath{-}(2/3)n\_f>0. \ensuremath{\Lambda}\_QCD\ensuremath{\approx}0,2 GeV. Gross-Wilczek e Politzer 1973 (Nobel 2004); confirmada pelo running de \ensuremath{\alpha}\_s.
\prov{relações: relaciona confinamento\_cor; parte\_de qcd\_cromodinamica}
\prov{proveniência: wiki \textperiodcentered{} Gross-Wilczek, Politzer 1973 PRL 30; Nobel 2004 \textperiodcentered{} fato \textperiodcentered{} confiança 1.0 \textperiodcentered{} domínio fisica \textperiodcentered{} história 3 versões no jornal}

%CRISTAL {"arestas":[{"alvo":"pipeline_stratum","confianca":1.0,"epistemico":"fato","origem":"wiki","rel":"parte_de"},{"alvo":"cifra","confianca":1.0,"epistemico":"fato","origem":"wiki","rel":"relaciona"}],"confianca":1.0,"contraexemplos":[],"descricao":"Biblioteca C (gcc -O3 -lcrypto -pthread) que faz a montagem, stratum_hash_header76 (SHA256² de header80) e a varredura de nonces (stratum_hash_scan). Presente e ativa em disco (17 KB). AQUI o SHA256² aparece — como a BORDA de consenso BTC, distinto do gato espectral do núcleo.","epistemico":"fato","exemplos":[],"id":"libstratum_bloco_so","memoria":"semantica","meta":{"autor":"broca-so","dominio":"metal","fonte":"neuronio/stratum_bloco.c/.h"},"origem":"wiki","palavras_chave":[],"sinonimos":[],"tipo":"conceito","titulo":"libstratum_bloco.so — o Metal do Bloco"}
\section{libstratum\_bloco.so — o Metal do Bloco}
Biblioteca C (gcc -O3 -lcrypto -pthread) que faz a montagem, stratum\_hash\_header76 (SHA256\ensuremath{^{2}} de header80) e a varredura de nonces (stratum\_hash\_scan). Presente e ativa em disco (17 KB). AQUI o SHA256\ensuremath{^{2}} aparece — como a BORDA de consenso BTC, distinto do gato espectral do núcleo.
\prov{relações: parte\_de pipeline\_stratum; relaciona cifra}
\prov{proveniência: wiki \textperiodcentered{} neuronio/stratum\_bloco.c/.h \textperiodcentered{} fato \textperiodcentered{} confiança 1.0 \textperiodcentered{} domínio metal \textperiodcentered{} história 3 versões no jornal}

%CRISTAL {"arestas":[{"alvo":"cristal_motor_congelado","confianca":1.0,"epistemico":"fato","origem":"wiki","rel":"explica"},{"alvo":"metal_broca_so","confianca":1.0,"epistemico":"fato","origem":"wiki","rel":"parte_de"}],"confianca":1.0,"contraexemplos":[],"descricao":"Biblioteca C99 (-O2 -mpopcnt -fPIC) em camadas 1_core→6_tools, compilada em libtiffany.a/.so + binários colapso e tiffany-index. Ponto de entrada único tiffany.h. Invariantes fixos: N=256, KOCH_P=1009, registro=6B header+32B impressão. É a forma-lib do cristal_motor_congelado.","epistemico":"fato","exemplos":[],"id":"libtiffany_engine","memoria":"semantica","meta":{"autor":"broca-so","dominio":"metal","fonte":"code/libtiffany/Makefile; tiffany.h"},"origem":"wiki","palavras_chave":[],"sinonimos":[],"tipo":"conceito","titulo":"libtiffany — a Engine Congelada (6 camadas)"}
\section{libtiffany — a Engine Congelada (6 camadas)}
Biblioteca C99 (-O2 -mpopcnt -fPIC) em camadas 1\_core\ensuremath{\to}6\_tools, compilada em libtiffany.a/.so + binários colapso e tiffany-index. Ponto de entrada único tiffany.h. Invariantes fixos: N=256, KOCH\_P=1009, registro=6B header+32B impressão. É a forma-lib do cristal\_motor\_congelado.
\prov{relações: explica cristal\_motor\_congelado; parte\_de metal\_broca\_so}
\prov{proveniência: wiki \textperiodcentered{} code/libtiffany/Makefile; tiffany.h \textperiodcentered{} fato \textperiodcentered{} confiança 1.0 \textperiodcentered{} domínio metal \textperiodcentered{} história 3 versões no jornal}

%CRISTAL {"arestas":[{"alvo":"chicote_e_o_tensor","confianca":1.0,"epistemico":"fato","origem":"wiki","rel":"usa"},{"alvo":"bai_e_a_cosmologia_quantica","confianca":1.0,"epistemico":"fato","origem":"wiki","rel":"usa"},{"alvo":"corpo_estelar","confianca":1.0,"epistemico":"fato","origem":"wiki","rel":"usa"},{"alvo":"efe_einstein","confianca":1.0,"epistemico":"fato","origem":"wiki","rel":"relaciona"},{"alvo":"mecanismo_de_la_hire","confianca":1.0,"epistemico":"fato","origem":"wiki","rel":"relaciona"}],"confianca":1.0,"contraexemplos":[],"descricao":"As três camadas da máquina mineral fecham numa física: ÁLGEBRA (o corpo estelar 𝕊 — a reta+La Hire, soma e produto), GEOMETRIA (o chicote=tensor — dinâmica posicional, contração de Einstein, cone, curvatura N−1) e DINÂMICA (o BAI=cosmologia — de Sitter, escada δ, H=log σ). A EFE de Einstein G_μν+Λg_μν=κT_μν UNIFICA-as: o CHICOTE (curvatura, geometria) = o BAI (energia δ). A validação visual é a mandala de epiciclos: concatenar círculos (compor rotações = a contração tensorial) reconstrói toda curva conservando a órbita. É INFERENCIA/analogia declarada — a estrutura da relatividade entre a régua e a delegação —, não uma teoria física nova; mas o grafo aceita, sem contradição.","epistemico":"inferencia","exemplos":[],"id":"linguagem_fisica_unificada","memoria":"semantica","meta":{"dominio":"fisica","fonte":"linguagem física unificada"},"origem":"wiki","palavras_chave":[],"sinonimos":[],"tipo":"conceito","titulo":"A Linguagem Física Unificada"}
\section{A Linguagem Física Unificada}
As três camadas da máquina mineral fecham numa física: ÁLGEBRA (o corpo estelar \ensuremath{\mathbb{S}} — a reta+La Hire, soma e produto), GEOMETRIA (o chicote=tensor — dinâmica posicional, contração de Einstein, cone, curvatura N\ensuremath{-}1) e DINÂMICA (o BAI=cosmologia — de Sitter, escada \ensuremath{\delta}, H=log \ensuremath{\sigma}). A EFE de Einstein G\_\ensuremath{\mu}\ensuremath{\nu}+\ensuremath{\Lambda}g\_\ensuremath{\mu}\ensuremath{\nu}=\ensuremath{\kappa}T\_\ensuremath{\mu}\ensuremath{\nu} UNIFICA-as: o CHICOTE (curvatura, geometria) = o BAI (energia \ensuremath{\delta}). A validação visual é a mandala de epiciclos: concatenar círculos (compor rotações = a contração tensorial) reconstrói toda curva conservando a órbita. É INFERENCIA/analogia declarada — a estrutura da relatividade entre a régua e a delegação —, não uma teoria física nova; mas o grafo aceita, sem contradição.
\prov{relações: usa chicote\_e\_o\_tensor; usa bai\_e\_a\_cosmologia\_quantica; usa corpo\_estelar; relaciona efe\_einstein; relaciona mecanismo\_de\_la\_hire}
\prov{proveniência: wiki \textperiodcentered{} linguagem física unificada \textperiodcentered{} inferencia \textperiodcentered{} confiança 1.0 \textperiodcentered{} domínio fisica \textperiodcentered{} história 12 versões no jornal}

%CRISTAL {"arestas":[{"alvo":"joalheria_hub","confianca":1.0,"epistemico":"fato","origem":"wiki","rel":"parte_de"},{"alvo":"tiffany_artesa","confianca":1.0,"epistemico":"fato","origem":"wiki","rel":"relaciona"},{"alvo":"assinatura_digital","confianca":1.0,"epistemico":"fato","origem":"wiki","rel":"usa"}],"confianca":1.0,"contraexemplos":[],"descricao":"A nona leitura da bancada: LITOGRAFAR é gravar uma assinatura (um termo, um nome) na pedra — e só se grava o que TEM PREIMAGEM no grafo. É a arte de marcar a peça sem forjar. A verificação é topológica (o invariante do grau), e a rejeição de uma forja é uma cirurgia (Perelman). Área registrada no hub joalheria; ferramenta em conversa/litografia.py. A litografia é onde o 'não fabricar' vira ofício.","epistemico":"fato","exemplos":[],"id":"litografia","memoria":"semantica","meta":{"dominio":"litografia","fonte":"litografia (joalheria)"},"origem":"wiki","palavras_chave":[],"sinonimos":[],"tipo":"conceito","titulo":"Litografia (a gravação de assinaturas) — área nova da joalheria"}
\section{Litografia (a gravação de assinaturas) — área nova da joalheria}
A nona leitura da bancada: LITOGRAFAR é gravar uma assinatura (um termo, um nome) na pedra — e só se grava o que TEM PREIMAGEM no grafo. É a arte de marcar a peça sem forjar. A verificação é topológica (o invariante do grau), e a rejeição de uma forja é uma cirurgia (Perelman). Área registrada no hub joalheria; ferramenta em conversa/litografia.py. A litografia é onde o 'não fabricar' vira ofício.
\prov{relações: parte\_de joalheria\_hub; relaciona tiffany\_artesa; usa assinatura\_digital}
\prov{proveniência: wiki \textperiodcentered{} litografia (joalheria) \textperiodcentered{} fato \textperiodcentered{} confiança 1.0 \textperiodcentered{} domínio litografia \textperiodcentered{} história 4 versões no jornal}

%CRISTAL {"arestas":[{"alvo":"wheeler_dewitt","confianca":1.0,"epistemico":"fato","origem":"wiki","rel":"relaciona"}],"confianca":1.0,"contraexemplos":[],"descricao":"Aplicação com redução por simetria da loop quantum gravity à cosmologia; a singularidade do Big Bang é substituída por um 'bounce' (ricochete) onde a gravidade quântica fica repulsiva. Programa teórico desenvolvido matematicamente, NÃO confirmado e sensível a ambiguidades de quantização.","epistemico":"hipotese","exemplos":[],"id":"loop_quantum_cosmology","memoria":"semantica","meta":{"autor":"broca-so","dominio":"fisica","fonte":"Bojowald 2001; Ashtekar-Pawlowski-Singh 2006"},"origem":"wiki","palavras_chave":[],"sinonimos":[],"tipo":"conceito","titulo":"Loop Quantum Cosmology e o 'quantum bounce'"}
\section{Loop Quantum Cosmology e o 'quantum bounce'}
Aplicação com redução por simetria da loop quantum gravity à cosmologia; a singularidade do Big Bang é substituída por um 'bounce' (ricochete) onde a gravidade quântica fica repulsiva. Programa teórico desenvolvido matematicamente, NÃO confirmado e sensível a ambiguidades de quantização.
\prov{relações: relaciona wheeler\_dewitt}
\prov{proveniência: wiki \textperiodcentered{} Bojowald 2001; Ashtekar-Pawlowski-Singh 2006 \textperiodcentered{} hipotese \textperiodcentered{} confiança 1.0 \textperiodcentered{} domínio fisica \textperiodcentered{} história 2 versões no jornal}

%CRISTAL {"arestas":[{"alvo":"supersimetria_susy","confianca":1.0,"epistemico":"fato","origem":"wiki","rel":"contradiz"},{"alvo":"teoria_gauge_octonica","confianca":1.0,"epistemico":"fato","origem":"wiki","rel":"usa"},{"alvo":"associador_nao_tensoriza","confianca":1.0,"epistemico":"fato","origem":"wiki","rel":"usa"},{"alvo":"limite_hurwitz_s1","confianca":1.0,"epistemico":"fato","origem":"wiki","rel":"relaciona"},{"alvo":"cosmologia_quantica_hub","confianca":1.0,"epistemico":"fato","origem":"wiki","rel":"parte_de"}],"confianca":1.0,"contraexemplos":[],"descricao":"Achado honesto: a série algébrica NÃO postula superparceiros nem SUSY breaking (zero ocorrências de SUSY nos papers-alvo; o Yang-Mills octônico é explicitamente NÃO-supersimétrico — 'cancelamentos tipo N=4 não ocorrem'). SUSY comparece só como corolário de Hurwitz (dimensões críticas D=n+2∈3,4,6,10). O conteúdo BSM de Lopes é o gauge G2=Aut(𝕆) e as correções octônicas — o programa é ORTOGONAL à supersimetria.","epistemico":"inferencia","exemplos":[],"id":"lopes_ortogonal_susy","memoria":"semantica","meta":{"autor":"broca-so","dominio":"cosmologia","fonte":"grep nos papers; paper_PQ_K:412; paper_qm_fibrado:1157"},"origem":"wiki","palavras_chave":[],"sinonimos":[],"tipo":"conceito","titulo":"O 'além' de Lopes NÃO é SUSY — é a não-associatividade"}
\section{O 'além' de Lopes NÃO é SUSY — é a não-associatividade}
Achado honesto: a série algébrica NÃO postula superparceiros nem SUSY breaking (zero ocorrências de SUSY nos papers-alvo; o Yang-Mills octônico é explicitamente NÃO-supersimétrico — 'cancelamentos tipo N=4 não ocorrem'). SUSY comparece só como corolário de Hurwitz (dimensões críticas D=n+2\ensuremath{\in}3,4,6,10). O conteúdo BSM de Lopes é o gauge G2=Aut(\ensuremath{\mathbb{O}}) e as correções octônicas — o programa é ORTOGONAL à supersimetria.
\prov{relações: contradiz supersimetria\_susy; usa teoria\_gauge\_octonica; usa associador\_nao\_tensoriza; relaciona limite\_hurwitz\_s1; parte\_de cosmologia\_quantica\_hub}
\prov{proveniência: wiki \textperiodcentered{} grep nos papers; paper\_PQ\_K:412; paper\_qm\_fibrado:1157 \textperiodcentered{} inferencia \textperiodcentered{} confiança 1.0 \textperiodcentered{} domínio cosmologia \textperiodcentered{} história 7 versões no jornal}

%CRISTAL {"arestas":[{"alvo":"cpu_metal_c","confianca":1.0,"epistemico":"fato","origem":"wiki","rel":"usa"},{"alvo":"assembler_duas_passagens","confianca":1.0,"epistemico":"fato","origem":"wiki","rel":"usa"}],"confianca":1.0,"contraexemplos":[],"descricao":"main_cpu.c lê um .bro do disco (cpu_load_program, fread), roda cpu_run, e imprime 'passos=' + relatório com A/B/R/flags/PC e stats (instruções, LOAD, STORE, saltos, ciclos). É o binário que executa bytecode no metal.","epistemico":"fato","exemplos":[],"id":"main_cpu_executavel","memoria":"semantica","meta":{"autor":"broca-so","dominio":"metal","fonte":"ula/main_cpu.c:main; cpu.c:cpu_print_relatorio"},"origem":"wiki","palavras_chave":[],"sinonimos":[],"tipo":"conceito","titulo":"broca_cpu — o Executável da Máquina"}
\section{broca\_cpu — o Executável da Máquina}
main\_cpu.c lê um .bro do disco (cpu\_load\_program, fread), roda cpu\_run, e imprime 'passos=' + relatório com A/B/R/flags/PC e stats (instruções, LOAD, STORE, saltos, ciclos). É o binário que executa bytecode no metal.
\prov{relações: usa cpu\_metal\_c; usa assembler\_duas\_passagens}
\prov{proveniência: wiki \textperiodcentered{} ula/main\_cpu.c:main; cpu.c:cpu\_print\_relatorio \textperiodcentered{} fato \textperiodcentered{} confiança 1.0 \textperiodcentered{} domínio metal \textperiodcentered{} história 3 versões no jornal}

%CRISTAL {"arestas":[{"alvo":"bcs_bec_supercondutividade","confianca":1.0,"epistemico":"fato","origem":"wiki","rel":"usa"},{"alvo":"efeito_hall_fracionario_anyons","confianca":1.0,"epistemico":"fato","origem":"wiki","rel":"relaciona"},{"alvo":"tenfold_way_bulk_boundary","confianca":1.0,"epistemico":"fato","origem":"wiki","rel":"usa"}],"confianca":1.0,"contraexemplos":[],"descricao":"Uma cadeia 1D supercondutora em fase topológica (Kitaev 2001) hospeda modos zero de MAJORANA (γ=γ†, E=0) desemparelhados nas extremidades: dois deles = um férmion não-local, um qubit intrinsecamente protegido de decoerência. Trançar anyons NÃO-ABELIANOS (MZM, ν=5/2) executaria portas quânticas tolerantes a falhas. Teoria SÓLIDA, mas EVIDÊNCIA EM DISPUTA: paper Nature 2018 (Delft/Microsoft) RETRATADO 2021; a 'Majorana 1' da Microsoft (2025) contestada (Legg, Nature 2026) — sinais ambíguos.","epistemico":"hipotese","exemplos":[],"id":"majorana_qubit_topologico","memoria":"semantica","meta":{"autor":"broca-so","dominio":"fisica","fonte":"Kitaev 2001 cond-mat/0010440; Nayak+2008; disputa 2021-2026"},"origem":"wiki","palavras_chave":[],"sinonimos":[],"tipo":"conceito","titulo":"Modos de Majorana e o qubit topológico (em disputa)"}
\section{Modos de Majorana e o qubit topológico (em disputa)}
Uma cadeia 1D supercondutora em fase topológica (Kitaev 2001) hospeda modos zero de MAJORANA (\ensuremath{\gamma}=\ensuremath{\gamma}\ensuremath{\dagger}, E=0) desemparelhados nas extremidades: dois deles = um férmion não-local, um qubit intrinsecamente protegido de decoerência. Trançar anyons NÃO-ABELIANOS (MZM, \ensuremath{\nu}=5/2) executaria portas quânticas tolerantes a falhas. Teoria SÓLIDA, mas EVIDÊNCIA EM DISPUTA: paper Nature 2018 (Delft/Microsoft) RETRATADO 2021; a 'Majorana 1' da Microsoft (2025) contestada (Legg, Nature 2026) — sinais ambíguos.
\prov{relações: usa bcs\_bec\_supercondutividade; relaciona efeito\_hall\_fracionario\_anyons; usa tenfold\_way\_bulk\_boundary}
\prov{proveniência: wiki \textperiodcentered{} Kitaev 2001 cond-mat/0010440; Nayak+2008; disputa 2021-2026 \textperiodcentered{} hipotese \textperiodcentered{} confiança 1.0 \textperiodcentered{} domínio fisica \textperiodcentered{} história 4 versões no jornal}

%CRISTAL {"arestas":[{"alvo":"maquina_de_corte","confianca":1.0,"epistemico":"fato","origem":"wiki","rel":"parte_de"},{"alvo":"tabela_geometrias_espectrais","confianca":1.0,"epistemico":"fato","origem":"wiki","rel":"usa"},{"alvo":"corpo_estelar","confianca":1.0,"epistemico":"fato","origem":"wiki","rel":"usa"}],"confianca":1.0,"contraexemplos":[],"descricao":"O primeiro corte é MECÂNICO: identifica-se o CENTRO da pedra no gabarito e assenta-se o MOLDE EXATO da tabela de geometrias — a MANDALA (a estrela k do corpo estelar, r=1+a·cos(kθ), reconstruída pelos epiciclos = a contração tensorial do chicote). Ajuste por mínimos quadrados robustos rejeita o reflexo como outlier; o molde dá a geometria EXATA e o reflexo fica de fora. Acertando o centro, o chicote corta exato — não é recorte de pixel, é o molde da máquina. linguagem/reconstrucao_3d.molde_espectral, recortar_pelo_molde. As curvas da mandala são as estrelas da tabela_geometrias_espectrais.","epistemico":"fato","exemplos":[],"id":"mandala_molde_do_corte","memoria":"semantica","meta":{"dominio":"joalheria","fonte":"contrato cristalchain"},"origem":"wiki","palavras_chave":[],"sinonimos":[],"tipo":"conceito","titulo":"A Mandala — o molde exato do corte (a contração tensorial do chicote)"}
\section{A Mandala — o molde exato do corte (a contração tensorial do chicote)}
O primeiro corte é MECÂNICO: identifica-se o CENTRO da pedra no gabarito e assenta-se o MOLDE EXATO da tabela de geometrias — a MANDALA (a estrela k do corpo estelar, r=1+a\textperiodcentered cos(k\ensuremath{\theta}), reconstruída pelos epiciclos = a contração tensorial do chicote). Ajuste por mínimos quadrados robustos rejeita o reflexo como outlier; o molde dá a geometria EXATA e o reflexo fica de fora. Acertando o centro, o chicote corta exato — não é recorte de pixel, é o molde da máquina. linguagem/reconstrucao\_3d.molde\_espectral, recortar\_pelo\_molde. As curvas da mandala são as estrelas da tabela\_geometrias\_espectrais.
\prov{relações: parte\_de maquina\_de\_corte; usa tabela\_geometrias\_espectrais; usa corpo\_estelar}
\prov{proveniência: wiki \textperiodcentered{} contrato cristalchain \textperiodcentered{} fato \textperiodcentered{} confiança 1.0 \textperiodcentered{} domínio joalheria \textperiodcentered{} história 4 versões no jornal}

%CRISTAL {"arestas":[{"alvo":"joalheria_hub","confianca":1.0,"epistemico":"fato","origem":"wiki","rel":"parte_de"},{"alvo":"morfologia_matematica_estelar","confianca":1.0,"epistemico":"fato","origem":"wiki","rel":"usa"},{"alvo":"tabela_geometrias_espectrais","confianca":1.0,"epistemico":"fato","origem":"wiki","rel":"usa"},{"alvo":"corpo_estelar_le_o_rastro","confianca":1.0,"epistemico":"fato","origem":"wiki","rel":"relaciona"},{"alvo":"area_render_gema","confianca":1.0,"epistemico":"fato","origem":"wiki","rel":"relaciona"}],"confianca":1.0,"contraexemplos":[],"descricao":"A fase de fabricação: o pipeline gabarito → segmentação → borda → simetria → tabela de geometrias → modo escolhido. A máquina SEGMENTA o cristal (seleção-fusão dos modos por k-means; poda o reflexo pela morfologia estelar), LÊ a frequência da borda (a simetria, dos tensores conectados via chicote — a geometria bem definida) e SELECIONA a geometria de corte na tabela. O corte é ESCOLHER o modo, não fabricar a forma. O gabarito → 3-fold → o trilhão. linguagem/fabricacao_gema.maquina_de_corte; joalheria.cortar. Separar observação/classificação/fabricação torna o pipeline escalável (nova gema = nova tabela, o pipeline não muda).","epistemico":"fato","exemplos":[],"id":"maquina_de_corte","memoria":"semantica","meta":{"dominio":"joalheria","fonte":"maquina de corte"},"origem":"wiki","palavras_chave":[],"sinonimos":[],"tipo":"conceito","titulo":"A Máquina de Corte (o corte é selecionar o modo, não fabricar)"}
\section{A Máquina de Corte (o corte é selecionar o modo, não fabricar)}
A fase de fabricação: o pipeline gabarito \ensuremath{\to} segmentação \ensuremath{\to} borda \ensuremath{\to} simetria \ensuremath{\to} tabela de geometrias \ensuremath{\to} modo escolhido. A máquina SEGMENTA o cristal (seleção-fusão dos modos por k-means; poda o reflexo pela morfologia estelar), LÊ a frequência da borda (a simetria, dos tensores conectados via chicote — a geometria bem definida) e SELECIONA a geometria de corte na tabela. O corte é ESCOLHER o modo, não fabricar a forma. O gabarito \ensuremath{\to} 3-fold \ensuremath{\to} o trilhão. linguagem/fabricacao\_gema.maquina\_de\_corte; joalheria.cortar. Separar observação/classificação/fabricação torna o pipeline escalável (nova gema = nova tabela, o pipeline não muda).
\prov{relações: parte\_de joalheria\_hub; usa morfologia\_matematica\_estelar; usa tabela\_geometrias\_espectrais; relaciona corpo\_estelar\_le\_o\_rastro; relaciona area\_render\_gema}
\prov{proveniência: wiki \textperiodcentered{} maquina de corte \textperiodcentered{} fato \textperiodcentered{} confiança 1.0 \textperiodcentered{} domínio joalheria \textperiodcentered{} história 306 versões no jornal}

%CRISTAL {"arestas":[{"alvo":"glueball_1_6_gev","confianca":1.0,"epistemico":"fato","origem":"wiki","rel":"especializa"},{"alvo":"glueballs_lattice","confianca":1.0,"epistemico":"fato","origem":"wiki","rel":"confirma"},{"alvo":"g2_automorfismos_associador","confianca":1.0,"epistemico":"fato","origem":"wiki","rel":"usa"}],"confianca":1.0,"contraexemplos":[],"descricao":"Fórmula fechada m₀₊₊=√2·r₁(Γ₀(7))·ΛQCD𝕆=√2·4,3617·0,264 GeV=1,628 GeV (erro 1,3% vs lattice 1,65±0,05). r₁=autovalor cuspidal de Maaß (rigoroso); √2=universalidade Higgs (derivado); ΛQCD𝕆≈264 MeV heurístico [J]. Ponte Cayley↔Maaß pela fórmula de Sunada ℓS=ℓX·log d (erro 0,00% assintótico). Junção crítica ‖φ‖²=168=|PSL(2,7)|.","epistemico":"inferencia","exemplos":[],"id":"mass_gap_cadeia_fechada","memoria":"semantica","meta":{"autor":"broca-so","dominio":"cosmologia","fonte":"paper_CB_mass_gap_solucao_completa.tex"},"origem":"wiki","palavras_chave":[],"sinonimos":[],"tipo":"conceito","titulo":"Mass gap — a cadeia fechada CB (1,628 GeV)"}
\section{Mass gap — a cadeia fechada CB (1,628 GeV)}
Fórmula fechada m\ensuremath{_{0++}}=\ensuremath{\sqrt{\,}}2\textperiodcentered r\ensuremath{_{1}}(\ensuremath{\Gamma}\ensuremath{_{0}}(7))\textperiodcentered \ensuremath{\Lambda}QCD\ensuremath{\mathbb{O}}=\ensuremath{\sqrt{\,}}2\textperiodcentered 4,3617\textperiodcentered 0,264 GeV=1,628 GeV (erro 1,3\% vs lattice 1,65$\pm$0,05). r\ensuremath{_{1}}=autovalor cuspidal de Maa\ss  (rigoroso); \ensuremath{\sqrt{\,}}2=universalidade Higgs (derivado); \ensuremath{\Lambda}QCD\ensuremath{\mathbb{O}}\ensuremath{\approx}264 MeV heurístico [J]. Ponte Cayley\ensuremath{\leftrightarrow}Maa\ss  pela fórmula de Sunada \ensuremath{\ell}S=\ensuremath{\ell}X\textperiodcentered log d (erro 0,00\% assintótico). Junção crítica \ensuremath{\|}\ensuremath{\varphi}\ensuremath{\|}\ensuremath{^{2}}=168=|PSL(2,7)|.
\prov{relações: especializa glueball\_1\_6\_gev; confirma glueballs\_lattice; usa g2\_automorfismos\_associador}
\prov{proveniência: wiki \textperiodcentered{} paper\_CB\_mass\_gap\_solucao\_completa.tex \textperiodcentered{} inferencia \textperiodcentered{} confiança 1.0 \textperiodcentered{} domínio cosmologia \textperiodcentered{} história 5 versões no jornal}

%CRISTAL {"arestas":[{"alvo":"mass_gap_algebrico","confianca":1.0,"epistemico":"fato","origem":"wiki","rel":"usa"},{"alvo":"g2_grupo_excepcional","confianca":1.0,"epistemico":"fato","origem":"wiki","rel":"usa"},{"alvo":"confinamento_cor","confianca":1.0,"epistemico":"fato","origem":"wiki","rel":"relaciona"}],"confianca":1.0,"contraexemplos":[],"descricao":"A quebra G₂→SU(3) por Higgs algébrico dá massa RIGOROSA a 6 dos 14 bósons (os 3⊕3̄), M_X²=g²‖v‖²; os 8 de su(3) ficam sem massa via Higgs — requerem confinamento SU(3) (o Clay residual). Termo novo octônico: o Jacobiator ⟨‖J‖²⟩ (correção de vácuo ~0,5%). 6/14 derivado; o gap dos 8 = conjectura de confinamento.","epistemico":"inferencia","exemplos":[],"id":"mass_gap_higgs_bz","memoria":"semantica","meta":{"autor":"broca-so","dominio":"cosmologia","fonte":"paper_BZ_mass_gap_higgs_algebrico.tex"},"origem":"wiki","palavras_chave":[],"sinonimos":[],"tipo":"conceito","titulo":"Mass gap parcial: G₂→SU(3) por Higgs algébrico (6/14)"}
\section{Mass gap parcial: G\ensuremath{_{2}}\ensuremath{\to}SU(3) por Higgs algébrico (6/14)}
A quebra G\ensuremath{_{2}}\ensuremath{\to}SU(3) por Higgs algébrico dá massa RIGOROSA a 6 dos 14 bósons (os 3\ensuremath{\oplus}3\ensuremath{}), M\_X\ensuremath{^{2}}=g\ensuremath{^{2}}\ensuremath{\|}v\ensuremath{\|}\ensuremath{^{2}}; os 8 de su(3) ficam sem massa via Higgs — requerem confinamento SU(3) (o Clay residual). Termo novo octônico: o Jacobiator \ensuremath{\langle}\ensuremath{\|}J\ensuremath{\|}\ensuremath{^{2}}\ensuremath{\rangle} (correção de vácuo \~{}0,5\%). 6/14 derivado; o gap dos 8 = conjectura de confinamento.
\prov{relações: usa mass\_gap\_algebrico; usa g2\_grupo\_excepcional; relaciona confinamento\_cor}
\prov{proveniência: wiki \textperiodcentered{} paper\_BZ\_mass\_gap\_higgs\_algebrico.tex \textperiodcentered{} inferencia \textperiodcentered{} confiança 1.0 \textperiodcentered{} domínio cosmologia \textperiodcentered{} história 4 versões no jornal}

%CRISTAL {"arestas":[{"alvo":"nucleo_g2_reabsorcao","confianca":1.0,"epistemico":"fato","origem":"wiki","rel":"usa"},{"alvo":"nao_associatividade_amputacao","confianca":1.0,"epistemico":"fato","origem":"wiki","rel":"eh_um"},{"alvo":"particulas_minerais","confianca":1.0,"epistemico":"fato","origem":"wiki","rel":"relaciona"}],"confianca":1.0,"contraexemplos":[],"descricao":"Em 3D a álgebra é associativa (associador=0): sem massa mínima. Em 7D o associador é irremovível — o 'resto que nunca zera' — e sua dimensão é 𝔤₂=14. É o mass gap do átomo mineral 7D: a massa mínima que a não-associatividade octoniônica impõe. Onde o átomo 3D (metal) é limpo, o 7D carrega o gap.","epistemico":"inferencia","exemplos":[],"id":"mass_gap_mineral","memoria":"semantica","meta":{"dominio":"fisica","fonte":"reta mineral 7D","mass_gap":14},"origem":"wiki","palavras_chave":[],"sinonimos":[],"tipo":"conceito","titulo":"O Mass Gap Mineral (o resto que nunca zera)"}
\section{O Mass Gap Mineral (o resto que nunca zera)}
Em 3D a álgebra é associativa (associador=0): sem massa mínima. Em 7D o associador é irremovível — o 'resto que nunca zera' — e sua dimensão é \ensuremath{\mathfrak{g}}\ensuremath{_{2}}=14. É o mass gap do átomo mineral 7D: a massa mínima que a não-associatividade octoniônica impõe. Onde o átomo 3D (metal) é limpo, o 7D carrega o gap.
\prov{relações: usa nucleo\_g2\_reabsorcao; eh\_um nao\_associatividade\_amputacao; relaciona particulas\_minerais}
\prov{proveniência: wiki \textperiodcentered{} reta mineral 7D \textperiodcentered{} inferencia \textperiodcentered{} confiança 1.0 \textperiodcentered{} domínio fisica \textperiodcentered{} história 4 versões no jornal}

%CRISTAL {"arestas":[{"alvo":"oscilacao_neutrinos","confianca":1.0,"epistemico":"fato","origem":"wiki","rel":"usa"},{"alvo":"desi_energia_escura_dinamica","confianca":1.0,"epistemico":"fato","origem":"wiki","rel":"relaciona"}],"confianca":1.0,"contraexemplos":[],"descricao":"As oscilações medem só diferenças de massa². Uma pequena de sinal conhecido (Δm²21≈7,49×10⁻⁵ eV², solar) e uma grande de sinal DESCONHECIDO (|Δm²31|≈2,51×10⁻³ eV², atmosférica) → o problema do ordenamento: normal (m1<m2<m3) vs invertido. Magnitudes = fato; o ordenamento é ABERTO (JUNO, DUNE, Hyper-K). Escala absoluta: KATRIN m_ν<0,45 eV (2025); cosmologia Σm_ν<0,072 eV (DESI+Planck) — em tensão com o mínimo permitido.","epistemico":"fato","exemplos":[],"id":"massa_neutrino_ordenamento","memoria":"semantica","meta":{"autor":"broca-so","dominio":"fisica","fonte":"NuFIT-6.0; KATRIN Science 2025; DESI 2024 arXiv:2404.03002"},"origem":"wiki","palavras_chave":[],"sinonimos":[],"tipo":"conceito","titulo":"Δm² e o ordenamento (hierarquia)"}
\section{\ensuremath{\Delta}m\ensuremath{^{2}} e o ordenamento (hierarquia)}
As oscilações medem só diferenças de massa\ensuremath{^{2}}. Uma pequena de sinal conhecido (\ensuremath{\Delta}m\ensuremath{^{2}}21\ensuremath{\approx}7,49$\times$10\ensuremath{^{-5}} eV\ensuremath{^{2}}, solar) e uma grande de sinal DESCONHECIDO (|\ensuremath{\Delta}m\ensuremath{^{2}}31|\ensuremath{\approx}2,51$\times$10\ensuremath{^{-3}} eV\ensuremath{^{2}}, atmosférica) \ensuremath{\to} o problema do ordenamento: normal (m1<m2<m3) vs invertido. Magnitudes = fato; o ordenamento é ABERTO (JUNO, DUNE, Hyper-K). Escala absoluta: KATRIN m\_\ensuremath{\nu}<0,45 eV (2025); cosmologia \ensuremath{\Sigma}m\_\ensuremath{\nu}<0,072 eV (DESI+Planck) — em tensão com o mínimo permitido.
\prov{relações: usa oscilacao\_neutrinos; relaciona desi\_energia\_escura\_dinamica}
\prov{proveniência: wiki \textperiodcentered{} NuFIT-6.0; KATRIN Science 2025; DESI 2024 arXiv:2404.03002 \textperiodcentered{} fato \textperiodcentered{} confiança 1.0 \textperiodcentered{} domínio fisica \textperiodcentered{} história 3 versões no jornal}

%CRISTAL {"arestas":[{"alvo":"lambda_cdm","confianca":1.0,"epistemico":"fato","origem":"wiki","rel":"parte_de"},{"alvo":"cosmologia_quantica_hub","confianca":1.0,"epistemico":"fato","origem":"wiki","rel":"parte_de"}],"confianca":1.0,"contraexemplos":[],"descricao":"Componente que interage gravitacionalmente mas não emite/absorve/reflete luz; domina a massa de galáxias e aglomerados e é NÃO-bariônica. Balanço (Planck 2018): matéria escura ~26,4%, bárions ~4,9%, energia escura ~68,5% (Ωch²=0,120, Ωbh²=0,0224). A existência de massa não-luminosa é inferência sólida; qual partícula a compõe é desconhecido.","epistemico":"fato","exemplos":[],"id":"materia_escura","memoria":"semantica","meta":{"autor":"broca-so","dominio":"cosmologia_dados","fonte":"Planck 2018 arXiv:1807.06209"},"origem":"wiki","palavras_chave":[],"sinonimos":[],"tipo":"conceito","titulo":"Matéria escura — a massa que falta (~27%)"}
\section{Matéria escura — a massa que falta (\~{}27\%)}
Componente que interage gravitacionalmente mas não emite/absorve/reflete luz; domina a massa de galáxias e aglomerados e é NÃO-bariônica. Balanço (Planck 2018): matéria escura \~{}26,4\%, bárions \~{}4,9\%, energia escura \~{}68,5\% (\ensuremath{\Omega}ch\ensuremath{^{2}}=0,120, \ensuremath{\Omega}bh\ensuremath{^{2}}=0,0224). A existência de massa não-luminosa é inferência sólida; qual partícula a compõe é desconhecido.
\prov{relações: parte\_de lambda\_cdm; parte\_de cosmologia\_quantica\_hub}
\prov{proveniência: wiki \textperiodcentered{} Planck 2018 arXiv:1807.06209 \textperiodcentered{} fato \textperiodcentered{} confiança 1.0 \textperiodcentered{} domínio cosmologia\_dados \textperiodcentered{} história 3 versões no jornal}

%CRISTAL {"arestas":[{"alvo":"materia_escura","confianca":1.0,"epistemico":"fato","origem":"wiki","rel":"explica"},{"alvo":"cdm_deteccao_nula","confianca":1.0,"epistemico":"fato","origem":"wiki","rel":"usa"},{"alvo":"mond_gravidade_modificada","confianca":1.0,"epistemico":"fato","origem":"wiki","rel":"usa"}],"confianca":1.0,"contraexemplos":[],"descricao":"A EVIDÊNCIA gravitacional é dos resultados mais robustos da cosmologia: seis linhas independentes (rotação, aglomerados, lentes, Bullet Cluster, CMB, estrutura) convergem quantitativamente sobre ~27% da densidade. A NATUREZA (que partícula) é problema aberto profundo: nenhum candidato confirmado, WIMPs sob pressão, áxions/estéreis viáveis mas não detectados. MOND explica galáxias mas fracassa em aglomerados/CMB. Consenso: a matéria escura existe como substância; sua identidade é desconhecida.","epistemico":"inferencia","exemplos":[],"id":"materia_escura_evidencia_vs_natureza","memoria":"semantica","meta":{"autor":"broca-so","dominio":"cosmologia_dados","fonte":"síntese (Planck 2018; Clowe 2006; LZ 2025; Milgrom 1983)"},"origem":"wiki","palavras_chave":[],"sinonimos":[],"tipo":"conceito","titulo":"O veredito: evidência robusta, natureza aberta"}
\section{O veredito: evidência robusta, natureza aberta}
A EVIDÊNCIA gravitacional é dos resultados mais robustos da cosmologia: seis linhas independentes (rotação, aglomerados, lentes, Bullet Cluster, CMB, estrutura) convergem quantitativamente sobre \~{}27\% da densidade. A NATUREZA (que partícula) é problema aberto profundo: nenhum candidato confirmado, WIMPs sob pressão, áxions/estéreis viáveis mas não detectados. MOND explica galáxias mas fracassa em aglomerados/CMB. Consenso: a matéria escura existe como substância; sua identidade é desconhecida.
\prov{relações: explica materia\_escura; usa cdm\_deteccao\_nula; usa mond\_gravidade\_modificada}
\prov{proveniência: wiki \textperiodcentered{} síntese (Planck 2018; Clowe 2006; LZ 2025; Milgrom 1983) \textperiodcentered{} inferencia \textperiodcentered{} confiança 1.0 \textperiodcentered{} domínio cosmologia\_dados \textperiodcentered{} história 4 versões no jornal}

%CRISTAL {"arestas":[{"alvo":"materia_escura","confianca":1.0,"epistemico":"fato","origem":"wiki","rel":"explica"},{"alvo":"associador_calor","confianca":1.0,"epistemico":"fato","origem":"wiki","rel":"usa"},{"alvo":"corpo_dual","confianca":1.0,"epistemico":"fato","origem":"wiki","rel":"usa"},{"alvo":"uma_maquina_duas_faces","confianca":1.0,"epistemico":"fato","origem":"wiki","rel":"usa"},{"alvo":"decomposicao_peirce","confianca":1.0,"epistemico":"fato","origem":"wiki","rel":"usa"},{"alvo":"floco_phi","confianca":1.0,"epistemico":"fato","origem":"wiki","rel":"relaciona"},{"alvo":"cosmologia_quantica_hub","confianca":1.0,"epistemico":"fato","origem":"wiki","rel":"parte_de"}],"confianca":1.0,"contraexemplos":[],"descricao":"Enquadramento (Aarão): o BSM é 'o outro lado do espectro'. Na balança a+c²=1, a matéria VISÍVEL é o c² (o lado branco, associativo, a célula de Peirce); a matéria ESCURA é o a (o lado negro — o associador, o irreversível, o não-associativo, o calor) que EQUILIBRA. 'Uma máquina, duas faces': não é uma nova partícula à moda SUSY, é a face complementar embutida na lei de conservação. O setor escuro é a metade não-associativa da mesma máquina.","epistemico":"inferencia","exemplos":[],"id":"materia_escura_lado_negro","memoria":"semantica","meta":{"autor":"broca-so","dominio":"cosmologia","fonte":"corpo_dual.tex; painel_estelar.tex (balança a+c²=1)"},"origem":"wiki","palavras_chave":[],"sinonimos":[],"tipo":"conceito","titulo":"A matéria escura é o lado negro que equilibra (a+c²=1)"}
\section{A matéria escura é o lado negro que equilibra (a+c\ensuremath{^{2}}=1)}
Enquadramento (Aarão): o BSM é 'o outro lado do espectro'. Na balança a+c\ensuremath{^{2}}=1, a matéria VISÍVEL é o c\ensuremath{^{2}} (o lado branco, associativo, a célula de Peirce); a matéria ESCURA é o a (o lado negro — o associador, o irreversível, o não-associativo, o calor) que EQUILIBRA. 'Uma máquina, duas faces': não é uma nova partícula à moda SUSY, é a face complementar embutida na lei de conservação. O setor escuro é a metade não-associativa da mesma máquina.
\prov{relações: explica materia\_escura; usa associador\_calor; usa corpo\_dual; usa uma\_maquina\_duas\_faces; usa decomposicao\_peirce; relaciona floco\_phi; parte\_de cosmologia\_quantica\_hub}
\prov{proveniência: wiki \textperiodcentered{} corpo\_dual.tex; painel\_estelar.tex (balança a+c\ensuremath{^{2}}=1) \textperiodcentered{} inferencia \textperiodcentered{} confiança 1.0 \textperiodcentered{} domínio cosmologia \textperiodcentered{} história 8 versões no jornal}

%CRISTAL {"arestas":[{"alvo":"qcd_cromodinamica","confianca":1.0,"epistemico":"fato","origem":"wiki","rel":"relaciona"},{"alvo":"cosmologia_quantica_hub","confianca":1.0,"epistemico":"fato","origem":"wiki","rel":"parte_de"}],"confianca":1.0,"contraexemplos":[],"descricao":"3 gerações de férmions (quarks + léptons), 15 de Weyl por geração (16 com ν direito); replicam as cargas de gauge, diferem só na massa. Cancelamento de anomalias: a estrutura quark-lépton (fator de cor 3) anula as anomalias → quantização da carga elétrica. 3 gerações confirmadas (largura do Z⁰ ⇒ 3 ν leves). ABERTO: por que exatamente 3 (o problema do sabor).","epistemico":"fato","exemplos":[],"id":"materia_modelo_padrao","memoria":"semantica","meta":{"autor":"broca-so","dominio":"fisica","fonte":"Glashow-Weinberg-Salam 1961-68; PDG"},"origem":"wiki","palavras_chave":[],"sinonimos":[],"tipo":"conceito","titulo":"O conteúdo de matéria do Modelo Padrão (3 gerações)"}
\section{O conteúdo de matéria do Modelo Padrão (3 gerações)}
3 gerações de férmions (quarks + léptons), 15 de Weyl por geração (16 com \ensuremath{\nu} direito); replicam as cargas de gauge, diferem só na massa. Cancelamento de anomalias: a estrutura quark-lépton (fator de cor 3) anula as anomalias \ensuremath{\to} quantização da carga elétrica. 3 gerações confirmadas (largura do Z\ensuremath{^{0}} \ensuremath{\Rightarrow} 3 \ensuremath{\nu} leves). ABERTO: por que exatamente 3 (o problema do sabor).
\prov{relações: relaciona qcd\_cromodinamica; parte\_de cosmologia\_quantica\_hub}
\prov{proveniência: wiki \textperiodcentered{} Glashow-Weinberg-Salam 1961-68; PDG \textperiodcentered{} fato \textperiodcentered{} confiança 1.0 \textperiodcentered{} domínio fisica \textperiodcentered{} história 3 versões no jornal}

%CRISTAL {"arestas":[{"alvo":"unificacao_eletrofraca_gws","confianca":1.0,"epistemico":"fato","origem":"wiki","rel":"parte_de"},{"alvo":"matriz_pmns","confianca":1.0,"epistemico":"fato","origem":"wiki","rel":"relaciona"},{"alvo":"cosmologia_quantica_hub","confianca":1.0,"epistemico":"fato","origem":"wiki","rel":"parte_de"}],"confianca":1.0,"contraexemplos":[],"descricao":"Matriz unitária 3×3 da mistura dos quarks nas correntes carregadas fracas (via W±): os autoestados de sabor dos quarks tipo-down ≠ autoestados de massa. É a única fonte de violação de sabor e de CP no setor de quarks do MP. Cabibbo (1963, 2 sabores, θ_C); Kobayashi-Maskawa (1973, 3 gerações). FATO — Nobel 2008 (Kobayashi, Maskawa; + Nambu).","epistemico":"fato","exemplos":[],"id":"matriz_ckm","memoria":"semantica","meta":{"autor":"broca-so","dominio":"fisica","fonte":"Cabibbo 1963 PRL 10,531; Kobayashi-Maskawa 1973; PDG"},"origem":"wiki","palavras_chave":[],"sinonimos":[],"tipo":"conceito","titulo":"A matriz CKM (Cabibbo-Kobayashi-Maskawa)"}
\section{A matriz CKM (Cabibbo-Kobayashi-Maskawa)}
Matriz unitária 3$\times$3 da mistura dos quarks nas correntes carregadas fracas (via W$\pm$): os autoestados de sabor dos quarks tipo-down \ensuremath{\ne} autoestados de massa. É a única fonte de violação de sabor e de CP no setor de quarks do MP. Cabibbo (1963, 2 sabores, \ensuremath{\theta}\_C); Kobayashi-Maskawa (1973, 3 gerações). FATO — Nobel 2008 (Kobayashi, Maskawa; + Nambu).
\prov{relações: parte\_de unificacao\_eletrofraca\_gws; relaciona matriz\_pmns; parte\_de cosmologia\_quantica\_hub}
\prov{proveniência: wiki \textperiodcentered{} Cabibbo 1963 PRL 10,531; Kobayashi-Maskawa 1973; PDG \textperiodcentered{} fato \textperiodcentered{} confiança 1.0 \textperiodcentered{} domínio fisica \textperiodcentered{} história 4 versões no jornal}

%CRISTAL {"arestas":[{"alvo":"oscilacao_neutrinos","confianca":1.0,"epistemico":"fato","origem":"wiki","rel":"usa"}],"confianca":1.0,"contraexemplos":[],"descricao":"Os 3 sabores (ν_e,ν_μ,ν_τ) ligam-se aos 3 autoestados de massa pela matriz unitária PMNS (Pontecorvo-Maki-Nakagawa-Sakata): 3 ângulos (θ12,θ23,θ13) + 1 fase δ_CP (+2 de Majorana se aplicável). NuFIT-6.0 (2024): sin²θ12=0,307, sin²θ13=0,0222, sin²θ23≈0,47 (octante ambíguo). θ13≠0 medido por Daya Bay 2012. Estrutura + θ12/θ13 = fato; octante de θ23 aberto.","epistemico":"fato","exemplos":[],"id":"matriz_pmns","memoria":"semantica","meta":{"autor":"broca-so","dominio":"fisica","fonte":"NuFIT-6.0 arXiv:2410.05380; Daya Bay 2012 arXiv:1203.1669"},"origem":"wiki","palavras_chave":[],"sinonimos":[],"tipo":"conceito","titulo":"Três sabores e a matriz PMNS"}
\section{Três sabores e a matriz PMNS}
Os 3 sabores (\ensuremath{\nu}\_e,\ensuremath{\nu}\_\ensuremath{\mu},\ensuremath{\nu}\_\ensuremath{\tau}) ligam-se aos 3 autoestados de massa pela matriz unitária PMNS (Pontecorvo-Maki-Nakagawa-Sakata): 3 ângulos (\ensuremath{\theta}12,\ensuremath{\theta}23,\ensuremath{\theta}13) + 1 fase \ensuremath{\delta}\_CP (+2 de Majorana se aplicável). NuFIT-6.0 (2024): sin\ensuremath{^{2}}\ensuremath{\theta}12=0,307, sin\ensuremath{^{2}}\ensuremath{\theta}13=0,0222, sin\ensuremath{^{2}}\ensuremath{\theta}23\ensuremath{\approx}0,47 (octante ambíguo). \ensuremath{\theta}13\ensuremath{\ne}0 medido por Daya Bay 2012. Estrutura + \ensuremath{\theta}12/\ensuremath{\theta}13 = fato; octante de \ensuremath{\theta}23 aberto.
\prov{relações: usa oscilacao\_neutrinos}
\prov{proveniência: wiki \textperiodcentered{} NuFIT-6.0 arXiv:2410.05380; Daya Bay 2012 arXiv:1203.1669 \textperiodcentered{} fato \textperiodcentered{} confiança 1.0 \textperiodcentered{} domínio fisica \textperiodcentered{} história 2 versões no jornal}

%CRISTAL {"arestas":[{"alvo":"koide_z3_geracoes","confianca":1.0,"epistemico":"fato","origem":"wiki","rel":"explica"},{"alvo":"os_dois_tres_cores_hurwitz_e_gerac","confianca":1.0,"epistemico":"fato","origem":"wiki","rel":"usa"},{"alvo":"acoplamentos_yukawa","confianca":1.0,"epistemico":"fato","origem":"wiki","rel":"relaciona"}],"confianca":1.0,"contraexemplos":[],"descricao":"Teorema (dedutivo): se a S₃ de sabor contém o 3-ciclo Z₃ da trialidade, toda matriz de massa que o respeita comuta com a permutação cíclica P das gerações, logo é CIRCULANTE: M=aI+bP+cP². Toda circulante é diagonalizada pela MESMA matriz de Fourier Fjk=ωjk/√3 (ω=e²πi/³), independentemente de a,b,c (verif. 10⁻¹⁶). O núcleo firme de todo o setor de sabor de Lopes. [D].","epistemico":"inferencia","exemplos":[],"id":"matrizes_massa_circulantes","memoria":"semantica","meta":{"autor":"broca-so","dominio":"cosmologia","fonte":"paper_4_particulas.tex obs:mistura-S3"},"origem":"wiki","palavras_chave":[],"sinonimos":[],"tipo":"conceito","titulo":"Matrizes de massa circulantes (Z₃ da trialidade)"}
\section{Matrizes de massa circulantes (Z\ensuremath{_{3}} da trialidade)}
Teorema (dedutivo): se a S\ensuremath{_{3}} de sabor contém o 3-ciclo Z\ensuremath{_{3}} da trialidade, toda matriz de massa que o respeita comuta com a permutação cíclica P das gerações, logo é CIRCULANTE: M=aI+bP+cP\ensuremath{^{2}}. Toda circulante é diagonalizada pela MESMA matriz de Fourier Fjk=\ensuremath{\omega}jk/\ensuremath{\sqrt{\,}}3 (\ensuremath{\omega}=e\ensuremath{^{2}}\ensuremath{\pi}i/\ensuremath{^{3}}), independentemente de a,b,c (verif. 10\ensuremath{^{-16}}). O núcleo firme de todo o setor de sabor de Lopes. [D].
\prov{relações: explica koide\_z3\_geracoes; usa os\_dois\_tres\_cores\_hurwitz\_e\_gerac; relaciona acoplamentos\_yukawa}
\prov{proveniência: wiki \textperiodcentered{} paper\_4\_particulas.tex obs:mistura-S3 \textperiodcentered{} inferencia \textperiodcentered{} confiança 1.0 \textperiodcentered{} domínio cosmologia \textperiodcentered{} história 5 versões no jornal}

%CRISTAL {"arestas":[{"alvo":"corpo_estelar","confianca":1.0,"epistemico":"fato","origem":"wiki","rel":"usa"},{"alvo":"corpo_estelar_le_o_rastro","confianca":1.0,"epistemico":"fato","origem":"wiki","rel":"parte_de"}],"confianca":1.0,"contraexemplos":[],"descricao":"A MECÂNICA do corpo estelar como ferramenta do pipe do metal: os modos POD (a decomposição ortogonal própria = os valores singulares) são o PAR HAMILTONIANO das duas rotações contrárias da estrela. O modo dominante e os modos EFETIVOS (exp da entropia) medem a riqueza do rastro. No gabarito: modo dominante 0.833, 2.6 modos efetivos. linguagem/joalheria.mecanica.","epistemico":"fato","exemplos":[],"id":"mecanica_estelar_le_o_rastro","memoria":"semantica","meta":{"dominio":"joalheria","fonte":"espectro orbitas"},"origem":"wiki","palavras_chave":[],"sinonimos":[],"tipo":"conceito","titulo":"Mecânica Estelar — os modos POD (o par hamiltoniano) do rastro"}
\section{Mecânica Estelar — os modos POD (o par hamiltoniano) do rastro}
A MECÂNICA do corpo estelar como ferramenta do pipe do metal: os modos POD (a decomposição ortogonal própria = os valores singulares) são o PAR HAMILTONIANO das duas rotações contrárias da estrela. O modo dominante e os modos EFETIVOS (exp da entropia) medem a riqueza do rastro. No gabarito: modo dominante 0.833, 2.6 modos efetivos. linguagem/joalheria.mecanica.
\prov{relações: usa corpo\_estelar; parte\_de corpo\_estelar\_le\_o\_rastro}
\prov{proveniência: wiki \textperiodcentered{} espectro orbitas \textperiodcentered{} fato \textperiodcentered{} confiança 1.0 \textperiodcentered{} domínio joalheria \textperiodcentered{} história 3 versões no jornal}

%CRISTAL {"arestas":[{"alvo":"matriz_ckm","confianca":1.0,"epistemico":"fato","origem":"wiki","rel":"relaciona"},{"alvo":"qcd_cromodinamica","confianca":1.0,"epistemico":"fato","origem":"wiki","rel":"relaciona"}],"confianca":1.0,"contraexemplos":[],"descricao":"Glashow-Iliopoulos-Maiani (1970): correntes neutras que mudam sabor (FCNC, ex. K⁰→μ⁺μ⁻) são fortemente suprimidas — as contribuições de loop se cancelam pela unitariedade da mistura, cancelamento que EXIGE um quarto quark. Levou à PREVISÃO do charm (confirmado 1974). Amplitude FCNC ∝ (m_c²−m_u²)/M_W². Base da fenomenologia de decaimentos raros. FATO.","epistemico":"fato","exemplos":[],"id":"mecanismo_gim","memoria":"semantica","meta":{"autor":"broca-so","dominio":"fisica","fonte":"Glashow-Iliopoulos-Maiani 1970 PRD 2,1285"},"origem":"wiki","palavras_chave":[],"sinonimos":[],"tipo":"conceito","titulo":"Mecanismo GIM e a supressão de FCNC"}
\section{Mecanismo GIM e a supressão de FCNC}
Glashow-Iliopoulos-Maiani (1970): correntes neutras que mudam sabor (FCNC, ex. K\ensuremath{^{0}}\ensuremath{\to}\ensuremath{\mu}\ensuremath{^{+}}\ensuremath{\mu}\ensuremath{^{-}}) são fortemente suprimidas — as contribuições de loop se cancelam pela unitariedade da mistura, cancelamento que EXIGE um quarto quark. Levou à PREVISÃO do charm (confirmado 1974). Amplitude FCNC \ensuremath{\propto} (m\_c\ensuremath{^{2}}\ensuremath{-}m\_u\ensuremath{^{2}})/M\_W\ensuremath{^{2}}. Base da fenomenologia de decaimentos raros. FATO.
\prov{relações: relaciona matriz\_ckm; relaciona qcd\_cromodinamica}
\prov{proveniência: wiki \textperiodcentered{} Glashow-Iliopoulos-Maiani 1970 PRD 2,1285 \textperiodcentered{} fato \textperiodcentered{} confiança 1.0 \textperiodcentered{} domínio fisica \textperiodcentered{} história 3 versões no jornal}

%CRISTAL {"arestas":[{"alvo":"unificacao_eletrofraca_gws","confianca":1.0,"epistemico":"fato","origem":"wiki","rel":"usa"},{"alvo":"bosons_w_z","confianca":1.0,"epistemico":"fato","origem":"wiki","rel":"explica"},{"alvo":"cosmologia_quantica_hub","confianca":1.0,"epistemico":"fato","origem":"wiki","rel":"parte_de"}],"confianca":1.0,"contraexemplos":[],"descricao":"Quebra espontânea SU(2)_L×U(1)_Y→U(1)_em por um campo escalar com vev não-nulo; os Goldstones são absorvidos dando massa a W/Z, o fóton fica sem massa. Potencial 'chapéu mexicano' V(φ)=−μ²|φ|²+λ|φ|⁴; vev v≈246 GeV; M_W=½gv, M_Z=½v√(g²+g'²). Englert-Brout, Higgs, Guralnik-Hagen-Kibble 1964.","epistemico":"fato","exemplos":[],"id":"mecanismo_higgs_ebh","memoria":"semantica","meta":{"autor":"broca-so","dominio":"fisica","fonte":"Englert-Brout, Higgs, GHK 1964 PRL 13"},"origem":"wiki","palavras_chave":[],"sinonimos":[],"tipo":"conceito","titulo":"Mecanismo de Higgs (Englert-Brout-Higgs-GHK, 1964)"}
\section{Mecanismo de Higgs (Englert-Brout-Higgs-GHK, 1964)}
Quebra espontânea SU(2)\_L$\times$U(1)\_Y\ensuremath{\to}U(1)\_em por um campo escalar com vev não-nulo; os Goldstones são absorvidos dando massa a W/Z, o fóton fica sem massa. Potencial 'chapéu mexicano' V(\ensuremath{\varphi})=\ensuremath{-}\ensuremath{\mu}\ensuremath{^{2}}|\ensuremath{\varphi}|\ensuremath{^{2}}+\ensuremath{\lambda}|\ensuremath{\varphi}|\ensuremath{^{4}}; vev v\ensuremath{\approx}246 GeV; M\_W=\ensuremath{\tfrac12}gv, M\_Z=\ensuremath{\tfrac12}v\ensuremath{\sqrt{\,}}(g\ensuremath{^{2}}+g'\ensuremath{^{2}}). Englert-Brout, Higgs, Guralnik-Hagen-Kibble 1964.
\prov{relações: usa unificacao\_eletrofraca\_gws; explica bosons\_w\_z; parte\_de cosmologia\_quantica\_hub}
\prov{proveniência: wiki \textperiodcentered{} Englert-Brout, Higgs, GHK 1964 PRL 13 \textperiodcentered{} fato \textperiodcentered{} confiança 1.0 \textperiodcentered{} domínio fisica \textperiodcentered{} história 4 versões no jornal}

%CRISTAL {"arestas":[{"alvo":"massa_neutrino_ordenamento","confianca":1.0,"epistemico":"fato","origem":"wiki","rel":"explica"},{"alvo":"dirac_majorana_0nbb","confianca":1.0,"epistemico":"fato","origem":"wiki","rel":"usa"},{"alvo":"so10_pati_salam","confianca":1.0,"epistemico":"fato","origem":"wiki","rel":"usa"},{"alvo":"cosmologia_quantica_hub","confianca":1.0,"epistemico":"fato","origem":"wiki","rel":"parte_de"}],"confianca":1.0,"contraexemplos":[],"descricao":"Explica a massa minúscula do ν (≲0,1 eV, ~10⁶× menor que o elétron): neutrinos de mão-direita muito pesados (M~escala GUT) dão m_ν≈m_D²/M_R (seesaw tipo I) — quanto mais pesado o parceiro, mais leve o ν leve. Para m_D~eletrofraca e m_ν~0,05 eV, M_R~10¹⁴-10¹⁵ GeV. Encaixa naturalmente no SO(10) (o ν_R completa a 16). Arcabouço favorecido mas NÃO comprovado (depende de Majorana + ν_R pesado).","epistemico":"hipotese","exemplos":[],"id":"mecanismo_seesaw","memoria":"semantica","meta":{"autor":"broca-so","dominio":"fisica","fonte":"Minkowski 1977; Gell-Mann-Ramond-Slansky; Mohapatra-Senjanović 1980"},"origem":"wiki","palavras_chave":[],"sinonimos":[],"tipo":"conceito","titulo":"Mecanismo seesaw (por que o ν é tão leve)"}
\section{Mecanismo seesaw (por que o \ensuremath{\nu} é tão leve)}
Explica a massa minúscula do \ensuremath{\nu} (\ensuremath{\lesssim}0,1 eV, \~{}10\ensuremath{^{6}}$\times$ menor que o elétron): neutrinos de mão-direita muito pesados (M\~{}escala GUT) dão m\_\ensuremath{\nu}\ensuremath{\approx}m\_D\ensuremath{^{2}}/M\_R (seesaw tipo I) — quanto mais pesado o parceiro, mais leve o \ensuremath{\nu} leve. Para m\_D\~{}eletrofraca e m\_\ensuremath{\nu}\~{}0,05 eV, M\_R\~{}10\ensuremath{^{14}}-10\ensuremath{^{15}} GeV. Encaixa naturalmente no SO(10) (o \ensuremath{\nu}\_R completa a 16). Arcabouço favorecido mas NÃO comprovado (depende de Majorana + \ensuremath{\nu}\_R pesado).
\prov{relações: explica massa\_neutrino\_ordenamento; usa dirac\_majorana\_0nbb; usa so10\_pati\_salam; parte\_de cosmologia\_quantica\_hub}
\prov{proveniência: wiki \textperiodcentered{} Minkowski 1977; Gell-Mann-Ramond-Slansky; Mohapatra-Senjanovi\'c 1980 \textperiodcentered{} hipotese \textperiodcentered{} confiança 1.0 \textperiodcentered{} domínio fisica \textperiodcentered{} história 5 versões no jornal}

%CRISTAL {"arestas":[{"alvo":"anomalias_b_rk_rd","confianca":1.0,"epistemico":"fato","origem":"wiki","rel":"explica"},{"alvo":"so10_pati_salam","confianca":1.0,"epistemico":"fato","origem":"wiki","rel":"relaciona"}],"confianca":1.0,"contraexemplos":[],"descricao":"Candidatos ativos, não observados: LEPTOQUARKS (acoplam quark a lépton, favoritos p/ R_D*; em GUTs/Pati-Salam; excluídos até ~1,2-1,7 TeV); Z'/W' (bósons de gauge pesados de U(1)'/SU(2)'; ressonâncias dilépton/dijato; excluídos até ~5-6 TeV); FÓTON ESCURO (U(1) oculta com mistura cinética ε; portal p/ o setor escuro; caçado na intensity frontier, grandes regiões excluídas). Hipóteses vivas, sem detecção.","epistemico":"hipotese","exemplos":[],"id":"mediadores_bsm","memoria":"semantica","meta":{"autor":"broca-so","dominio":"fisica","fonte":"PDG; NA62 2024; buscas ATLAS/CMS"},"origem":"wiki","palavras_chave":[],"sinonimos":[],"tipo":"conceito","titulo":"Mediadores BSM: leptoquarks, Z'/W', fóton escuro"}
\section{Mediadores BSM: leptoquarks, Z'/W', fóton escuro}
Candidatos ativos, não observados: LEPTOQUARKS (acoplam quark a lépton, favoritos p/ R\_D*; em GUTs/Pati-Salam; excluídos até \~{}1,2-1,7 TeV); Z'/W' (bósons de gauge pesados de U(1)'/SU(2)'; ressonâncias dilépton/dijato; excluídos até \~{}5-6 TeV); FÓTON ESCURO (U(1) oculta com mistura cinética \ensuremath{\varepsilon}; portal p/ o setor escuro; caçado na intensity frontier, grandes regiões excluídas). Hipóteses vivas, sem detecção.
\prov{relações: explica anomalias\_b\_rk\_rd; relaciona so10\_pati\_salam}
\prov{proveniência: wiki \textperiodcentered{} PDG; NA62 2024; buscas ATLAS/CMS \textperiodcentered{} hipotese \textperiodcentered{} confiança 1.0 \textperiodcentered{} domínio fisica \textperiodcentered{} história 3 versões no jornal}

%CRISTAL {"arestas":[{"alvo":"a_medida_o_calor_de_volta","confianca":1.0,"epistemico":"fato","origem":"wiki","rel":"especializa"},{"alvo":"decoerencia_colapso_medida","confianca":1.0,"epistemico":"fato","origem":"wiki","rel":"explica"},{"alvo":"associador_calor","confianca":1.0,"epistemico":"fato","origem":"wiki","rel":"usa"}],"confianca":1.0,"contraexemplos":[],"descricao":"Medir s força a máquina a SER ℝ,ℂ,ℍ ou 𝕆 (o campo local, onde a∥b, S=0, o relógio da métrica PARA — o tempo só corre ENTRE campos locais). O custo dessa irreversibilidade é o imposto V=(1−s²)‖a×b‖² pago à não-comutatividade, que retorna como calor/decoerência: a 'sangria' 1−(‖r‖²+C²) CRESCE com o emaranhamento (o mais entrelaçado é o mais frágil) e iguala o que a decoerência do associador prevê — 'o emaranhamento decoere como o associador, logo É o associador'. É o lado a da lei a+c²=1 pago de volta. Refina o nó a_medida_o_calor_de_volta.","epistemico":"inferencia","exemplos":[],"id":"medida_colapso_numa_letra","memoria":"semantica","meta":{"autor":"broca-so","dominio":"cosmologia","fonte":"paper_6_maquina.tex sec:quantizacao; paper_3 obs:medida; paper_F §2ª lei"},"origem":"wiki","palavras_chave":[],"sinonimos":[],"tipo":"conceito","titulo":"A medida = colapsar numa letra (o imposto = calor)"}
\section{A medida = colapsar numa letra (o imposto = calor)}
Medir s força a máquina a SER \ensuremath{\mathbb{R}},\ensuremath{\mathbb{C}},\ensuremath{\mathbb{H}} ou \ensuremath{\mathbb{O}} (o campo local, onde a\ensuremath{\parallel}b, S=0, o relógio da métrica PARA — o tempo só corre ENTRE campos locais). O custo dessa irreversibilidade é o imposto V=(1\ensuremath{-}s\ensuremath{^{2}})\ensuremath{\|}a$\times$b\ensuremath{\|}\ensuremath{^{2}} pago à não-comutatividade, que retorna como calor/decoerência: a 'sangria' 1\ensuremath{-}(\ensuremath{\|}r\ensuremath{\|}\ensuremath{^{2}}+C\ensuremath{^{2}}) CRESCE com o emaranhamento (o mais entrelaçado é o mais frágil) e iguala o que a decoerência do associador prevê — 'o emaranhamento decoere como o associador, logo É o associador'. É o lado a da lei a+c\ensuremath{^{2}}=1 pago de volta. Refina o nó a\_medida\_o\_calor\_de\_volta.
\prov{relações: especializa a\_medida\_o\_calor\_de\_volta; explica decoerencia\_colapso\_medida; usa associador\_calor}
\prov{proveniência: wiki \textperiodcentered{} paper\_6\_maquina.tex sec:quantizacao; paper\_3 obs:medida; paper\_F §2ª lei \textperiodcentered{} inferencia \textperiodcentered{} confiança 1.0 \textperiodcentered{} domínio cosmologia \textperiodcentered{} história 4 versões no jornal}

%CRISTAL {"arestas":[{"alvo":"cpu_metal_c","confianca":1.0,"epistemico":"fato","origem":"wiki","rel":"usa"},{"alvo":"neuron_io_popcnt","confianca":1.0,"epistemico":"fato","origem":"wiki","rel":"usa"}],"confianca":1.0,"contraexemplos":[],"descricao":"A 'RAM' da máquina são arquivos em mem/ nomeados %04u.bin por slot. Slots < 8000 passam pelo neurônio ao carregar (o conteúdo é reinterpretado espectralmente a cada LOAD via neuron_word_file); slots ≥ 8000 (SCRATCH) são roundtrip binário exato. STORE sempre grava o word binário de R.","epistemico":"fato","exemplos":[],"id":"memoria_de_slots","memoria":"semantica","meta":{"autor":"broca-so","dominio":"metal","fonte":"ula/memoria.c:mem_load_word/mem_store_word"},"origem":"wiki","palavras_chave":[],"sinonimos":[],"tipo":"conceito","titulo":"Memória como Arquivos mem/NNNN.bin"}
\section{Memória como Arquivos mem/NNNN.bin}
A 'RAM' da máquina são arquivos em mem/ nomeados \%04u.bin por slot. Slots < 8000 passam pelo neurônio ao carregar (o conteúdo é reinterpretado espectralmente a cada LOAD via neuron\_word\_file); slots \ensuremath{\ge} 8000 (SCRATCH) são roundtrip binário exato. STORE sempre grava o word binário de R.
\prov{relações: usa cpu\_metal\_c; usa neuron\_io\_popcnt}
\prov{proveniência: wiki \textperiodcentered{} ula/memoria.c:mem\_load\_word/mem\_store\_word \textperiodcentered{} fato \textperiodcentered{} confiança 1.0 \textperiodcentered{} domínio metal \textperiodcentered{} história 3 versões no jornal}

%CRISTAL {"arestas":[{"alvo":"confinamento_cor","confianca":1.0,"epistemico":"fato","origem":"wiki","rel":"relaciona"},{"alvo":"g2_automorfismos_associador","confianca":1.0,"epistemico":"fato","origem":"wiki","rel":"usa"},{"alvo":"leitura_de_part_culas_o_setor_de_materia_pela_torre_de_divisao","confianca":1.0,"epistemico":"fato","origem":"wiki","rel":"parte_de"}],"confianca":1.0,"contraexemplos":[],"descricao":"A fundamental 7=Im(𝕆) decompõe sob SU(3)=StabG₂(â) como 3⊕3̄⊕1 (via adâ=estrutura complexa J). MÉSON: par 3⊗3̄→1 (incolor = morar em ℂâ); BÁRION: 3⊗3⊗3→1=detℂ (as retas de Fano que evitam â). Partição 7=3+4 (3 retas por â = méson, 4 que evitam = bárion), forçada pela geometria. Correspondência exata combinatorial; não dá massas nem dinâmica.","epistemico":"inferencia","exemplos":[],"id":"meson_barion_setor_materia","memoria":"semantica","meta":{"autor":"broca-so","dominio":"cosmologia","fonte":"paper_4_particulas.tex §meson, §barion"},"origem":"wiki","palavras_chave":[],"sinonimos":[],"tipo":"conceito","titulo":"Méson e bárion: 7=Im(𝕆)=3⊕3̄⊕1 sob SU(3)"}
\section{Méson e bárion: 7=Im(\ensuremath{\mathbb{O}})=3\ensuremath{\oplus}3\ensuremath{}\ensuremath{\oplus}1 sob SU(3)}
A fundamental 7=Im(\ensuremath{\mathbb{O}}) decompõe sob SU(3)=StabG\ensuremath{_{2}}(â) como 3\ensuremath{\oplus}3\ensuremath{}\ensuremath{\oplus}1 (via adâ=estrutura complexa J). MÉSON: par 3\ensuremath{\otimes}3\ensuremath{}\ensuremath{\to}1 (incolor = morar em \ensuremath{\mathbb{C}}â); BÁRION: 3\ensuremath{\otimes}3\ensuremath{\otimes}3\ensuremath{\to}1=det\ensuremath{\mathbb{C}} (as retas de Fano que evitam â). Partição 7=3+4 (3 retas por â = méson, 4 que evitam = bárion), forçada pela geometria. Correspondência exata combinatorial; não dá massas nem dinâmica.
\prov{relações: relaciona confinamento\_cor; usa g2\_automorfismos\_associador; parte\_de leitura\_de\_part\_culas\_o\_setor\_de\_materia\_pela\_torre\_de\_divisao}
\prov{proveniência: wiki \textperiodcentered{} paper\_4\_particulas.tex §meson, §barion \textperiodcentered{} inferencia \textperiodcentered{} confiança 1.0 \textperiodcentered{} domínio cosmologia \textperiodcentered{} história 5 versões no jornal}

%CRISTAL {"arestas":[{"alvo":"isa_metal_21_opcodes","confianca":1.0,"epistemico":"fato","origem":"wiki","rel":"usa"},{"alvo":"cifra_metais_an","confianca":1.0,"epistemico":"fato","origem":"wiki","rel":"explica"}],"confianca":1.0,"contraexemplos":[],"descricao":"SILVER/GOLD/BRONZE aplicam Aₙ=[n·total+e, total] (n=2/1/3) direto em A; FOLD=project(silver(x),ref), UNFOLD=gold(x), LIFT=project(bronze(x),ref); PROJECT usa parabola_auditar (calor/c² Hellinger-Bhattacharyya). SPECT/LOADS acionam a cifra (A₂) no caminho de memória. Todos implementados.","epistemico":"fato","exemplos":[],"id":"metais_parabola_opcodes","memoria":"semantica","meta":{"autor":"broca-so","dominio":"metal","fonte":"ula/instrucoes.c; cifra.c:cifra_an/project"},"origem":"wiki","palavras_chave":[],"sinonimos":[],"tipo":"conceito","titulo":"Metais e Parábola como Opcodes da ISA"}
\section{Metais e Parábola como Opcodes da ISA}
SILVER/GOLD/BRONZE aplicam A\ensuremath{_{n}}=[n\textperiodcentered total+e, total] (n=2/1/3) direto em A; FOLD=project(silver(x),ref), UNFOLD=gold(x), LIFT=project(bronze(x),ref); PROJECT usa parabola\_auditar (calor/c\ensuremath{^{2}} Hellinger-Bhattacharyya). SPECT/LOADS acionam a cifra (A\ensuremath{_{2}}) no caminho de memória. Todos implementados.
\prov{relações: usa isa\_metal\_21\_opcodes; explica cifra\_metais\_an}
\prov{proveniência: wiki \textperiodcentered{} ula/instrucoes.c; cifra.c:cifra\_an/project \textperiodcentered{} fato \textperiodcentered{} confiança 1.0 \textperiodcentered{} domínio metal \textperiodcentered{} história 3 versões no jornal}

%CRISTAL {"arestas":[{"alvo":"broca_os","confianca":1.0,"epistemico":"fato","origem":"wiki","rel":"explica"},{"alvo":"operacao_broca_so","confianca":1.0,"epistemico":"fato","origem":"wiki","rel":"parte_de"}],"confianca":1.0,"contraexemplos":[],"descricao":"A implementação real no silício: a CPU em C, o motor de PoW espectral, o pipeline Stratum de mineração e o serviço no servidor GEX44 — o que está de fato compilado e rodando, verificado contra o código.","epistemico":"fato","exemplos":[],"id":"metal_broca_so","memoria":"semantica","meta":{"autor":"broca-so","dominio":"metal","fonte":"ula/, neuronio/, code/, scripts/gex44/"},"origem":"wiki","palavras_chave":[],"sinonimos":[],"tipo":"conceito","titulo":"O Metal do Broca-SO"}
\section{O Metal do Broca-SO}
A implementação real no silício: a CPU em C, o motor de PoW espectral, o pipeline Stratum de mineração e o serviço no servidor GEX44 — o que está de fato compilado e rodando, verificado contra o código.
\prov{relações: explica broca\_os; parte\_de operacao\_broca\_so}
\prov{proveniência: wiki \textperiodcentered{} ula/, neuronio/, code/, scripts/gex44/ \textperiodcentered{} fato \textperiodcentered{} confiança 1.0 \textperiodcentered{} domínio metal \textperiodcentered{} história 3 versões no jornal}

%CRISTAL {"arestas":[{"alvo":"cifra_metais_an","confianca":1.0,"epistemico":"fato","origem":"wiki","rel":"usa"},{"alvo":"dimensao_de_prata_pell","confianca":1.0,"epistemico":"fato","origem":"wiki","rel":"relaciona"}],"confianca":1.0,"contraexemplos":[],"descricao":"metal_verifica_rec conta falhas de t[k+2]=n·t[k+1]+t[k]; metal_fit_n acha o n∈[1,8] que minimiza falhas; metal_fit_ord2 faz ajuste cego de ordem-2. Usado no PoW para decidir se há um 'atalho' espectral (uma recorrência metálica escondida) na sequência.","epistemico":"fato","exemplos":[],"id":"metal_recorrencia_fit","memoria":"semantica","meta":{"autor":"broca-so","dominio":"metal","fonte":"ula/metais.c:metal_fit_n/metal_verifica_rec"},"origem":"wiki","palavras_chave":[],"sinonimos":[],"tipo":"conceito","titulo":"Detecção de Recorrência Aₙ (o atalho espectral)"}
\section{Detecção de Recorrência A\ensuremath{_{n}} (o atalho espectral)}
metal\_verifica\_rec conta falhas de t[k+2]=n\textperiodcentered t[k+1]+t[k]; metal\_fit\_n acha o n\ensuremath{\in}[1,8] que minimiza falhas; metal\_fit\_ord2 faz ajuste cego de ordem-2. Usado no PoW para decidir se há um 'atalho' espectral (uma recorrência metálica escondida) na sequência.
\prov{relações: usa cifra\_metais\_an; relaciona dimensao\_de\_prata\_pell}
\prov{proveniência: wiki \textperiodcentered{} ula/metais.c:metal\_fit\_n/metal\_verifica\_rec \textperiodcentered{} fato \textperiodcentered{} confiança 1.0 \textperiodcentered{} domínio metal \textperiodcentered{} história 3 versões no jornal}

%CRISTAL {"arestas":[{"alvo":"equacoes_friedmann","confianca":1.0,"epistemico":"fato","origem":"wiki","rel":"usa"},{"alvo":"redshift_bigbang_horizonte","confianca":1.0,"epistemico":"fato","origem":"wiki","rel":"explica"},{"alvo":"cosmologia_quantica_hub","confianca":1.0,"epistemico":"fato","origem":"wiki","rel":"parte_de"}],"confianca":1.0,"contraexemplos":[],"descricao":"ds²=−c²dt²+a²(t)[dr²/(1−kr²)+r²dΩ²]. a(t)=fator de escala; k=curvatura espacial (+1 fechado, 0 plano, −1 aberto). Segue do Princípio Cosmológico (homogeneidade+isotropia), apoiado por CMB e surveys.","epistemico":"fato","exemplos":[],"id":"metrica_flrw","memoria":"semantica","meta":{"autor":"broca-so","dominio":"fisica","fonte":"Friedmann 1922/24, Lemaître 1927, Robertson 1935, Walker 1937"},"origem":"wiki","palavras_chave":[],"sinonimos":[],"tipo":"conceito","titulo":"Métrica FLRW (universo homogêneo e isotrópico)"}
\section{Métrica FLRW (universo homogêneo e isotrópico)}
ds\ensuremath{^{2}}=\ensuremath{-}c\ensuremath{^{2}}dt\ensuremath{^{2}}+a\ensuremath{^{2}}(t)[dr\ensuremath{^{2}}/(1\ensuremath{-}kr\ensuremath{^{2}})+r\ensuremath{^{2}}d\ensuremath{\Omega}\ensuremath{^{2}}]. a(t)=fator de escala; k=curvatura espacial (+1 fechado, 0 plano, \ensuremath{-}1 aberto). Segue do Princípio Cosmológico (homogeneidade+isotropia), apoiado por CMB e surveys.
\prov{relações: usa equacoes\_friedmann; explica redshift\_bigbang\_horizonte; parte\_de cosmologia\_quantica\_hub}
\prov{proveniência: wiki \textperiodcentered{} Friedmann 1922/24, Lemaître 1927, Robertson 1935, Walker 1937 \textperiodcentered{} fato \textperiodcentered{} confiança 1.0 \textperiodcentered{} domínio fisica \textperiodcentered{} história 4 versões no jornal}

%CRISTAL {"arestas":[{"alvo":"voxelizador_cristalografico","confianca":1.0,"epistemico":"fato","origem":"wiki","rel":"relaciona"},{"alvo":"joalheria_hub","confianca":1.0,"epistemico":"fato","origem":"wiki","rel":"relaciona"},{"alvo":"jogos_animacoes_hub","confianca":1.0,"epistemico":"fato","origem":"wiki","rel":"parte_de"}],"confianca":1.0,"contraexemplos":[],"descricao":"O Teorema Geral da Tradução: mesmo operador, o dicionário troca. O VOXELIZADOR cristalográfico, lido pela lente do MINECRAFT (o jogo de blocos/voxels). Compondo voxelizador_maquina com minecraft_maquina (ambos no pivô 'maquina'/parábola) sai o dicionário PARALELO: voxelizar↔minerar, gema↔mundo, cela_unitária↔bloco, borda_analítica↔limite_do_mundo, empilhamento↔caminhada, cromóforo↔tocha, quilate↔coordenada, difração↔bioma. Minecraft é a máquina mineral jogada: preencher o mundo com blocos É voxelizar a gema com células do cristal. linguagem/dicionarios.py (_MINECRAFT, _VOXELIZADOR).","epistemico":"fato","exemplos":[],"id":"minecraft_voxel","memoria":"semantica","meta":{"dominio":"joalheria","fonte":"maquina de corte"},"origem":"wiki","palavras_chave":[],"sinonimos":[],"tipo":"conceito","titulo":"Minecraft — o voxelizador lido como jogo de blocos (o dicionário paralelo)"}
\section{Minecraft — o voxelizador lido como jogo de blocos (o dicionário paralelo)}
O Teorema Geral da Tradução: mesmo operador, o dicionário troca. O VOXELIZADOR cristalográfico, lido pela lente do MINECRAFT (o jogo de blocos/voxels). Compondo voxelizador\_maquina com minecraft\_maquina (ambos no pivô 'maquina'/parábola) sai o dicionário PARALELO: voxelizar\ensuremath{\leftrightarrow}minerar, gema\ensuremath{\leftrightarrow}mundo, cela\_unitária\ensuremath{\leftrightarrow}bloco, borda\_analítica\ensuremath{\leftrightarrow}limite\_do\_mundo, empilhamento\ensuremath{\leftrightarrow}caminhada, cromóforo\ensuremath{\leftrightarrow}tocha, quilate\ensuremath{\leftrightarrow}coordenada, difração\ensuremath{\leftrightarrow}bioma. Minecraft é a máquina mineral jogada: preencher o mundo com blocos É voxelizar a gema com células do cristal. linguagem/dicionarios.py (\_MINECRAFT, \_VOXELIZADOR).
\prov{relações: relaciona voxelizador\_cristalografico; relaciona joalheria\_hub; parte\_de jogos\_animacoes\_hub}
\prov{proveniência: wiki \textperiodcentered{} maquina de corte \textperiodcentered{} fato \textperiodcentered{} confiança 1.0 \textperiodcentered{} domínio joalheria \textperiodcentered{} história 79 versões no jornal}

%CRISTAL {"arestas":[{"alvo":"cristalchain","confianca":1.0,"epistemico":"fato","origem":"wiki","rel":"parte_de"},{"alvo":"joalheria_hub","confianca":1.0,"epistemico":"fato","origem":"wiki","rel":"relaciona"},{"alvo":"metais","confianca":1.0,"epistemico":"fato","origem":"wiki","rel":"usa"},{"alvo":"symbolic_beta_protocolo","confianca":1.0,"epistemico":"fato","origem":"wiki","rel":"relaciona"}],"confianca":1.0,"contraexemplos":[],"descricao":"A ponte que alinha a mineração (symbolic.py, SymbolicBeta) à joalheria/pipe 3D SEM tocar no núcleo. A mineração refina o floco em 12 METAIS (σ_n, as retas x²−n·x−1) — os MESMOS metais que a joalheria lê como as oito físicas (geometria→matéria). Não fabrica metal: refina o floco nos metais que a física já conhece. cristalchain/mineracao_pipe.py (6/6). O gancho é a órbita/o resíduo; a conservação é N(σ_n)=−1.","epistemico":"fato","exemplos":[],"id":"mineracao_pipe_hook","memoria":"semantica","meta":{"dominio":"mineracao_pipe","fonte":"mineracao_pipe hub"},"origem":"wiki","palavras_chave":[],"sinonimos":[],"tipo":"conceito","titulo":"Mineração ↔ pipe 3D (o gancho)"}
\section{Mineração \ensuremath{\leftrightarrow} pipe 3D (o gancho)}
A ponte que alinha a mineração (symbolic.py, SymbolicBeta) à joalheria/pipe 3D SEM tocar no núcleo. A mineração refina o floco em 12 METAIS (\ensuremath{\sigma}\_n, as retas x\ensuremath{^{2}}\ensuremath{-}n\textperiodcentered x\ensuremath{-}1) — os MESMOS metais que a joalheria lê como as oito físicas (geometria\ensuremath{\to}matéria). Não fabrica metal: refina o floco nos metais que a física já conhece. cristalchain/mineracao\_pipe.py (6/6). O gancho é a órbita/o resíduo; a conservação é N(\ensuremath{\sigma}\_n)=\ensuremath{-}1.
\prov{relações: parte\_de cristalchain; relaciona joalheria\_hub; usa metais; relaciona symbolic\_beta\_protocolo}
\prov{proveniência: wiki \textperiodcentered{} mineracao\_pipe hub \textperiodcentered{} fato \textperiodcentered{} confiança 1.0 \textperiodcentered{} domínio mineracao\_pipe \textperiodcentered{} história 5 versões no jornal}

%CRISTAL {"arestas":[],"confianca":1.0,"contraexemplos":[],"descricao":"Trunca o superespaço infinito-dimensional a poucos graus (impondo homogeneidade/isotropia, FLRW), tipicamente só o fator de escala a(t). Técnica de aproximação bem estabelecida (fato metodológico); conclusões condicionadas à validade da truncagem. É onde a s̈=2s do Lopes vira WDW.","epistemico":"fato","exemplos":[],"id":"minisuperspace","memoria":"semantica","meta":{"autor":"broca-so","dominio":"fisica","fonte":"Misner 1969; DeWitt 1967"},"origem":"wiki","palavras_chave":[],"sinonimos":[],"tipo":"conceito","titulo":"Minisuperspace (redução por simetria)"}
\section{Minisuperspace (redução por simetria)}
Trunca o superespaço infinito-dimensional a poucos graus (impondo homogeneidade/isotropia, FLRW), tipicamente só o fator de escala a(t). Técnica de aproximação bem estabelecida (fato metodológico); conclusões condicionadas à validade da truncagem. É onde a s\ensuremath{}=2s do Lopes vira WDW.
\prov{proveniência: wiki \textperiodcentered{} Misner 1969; DeWitt 1967 \textperiodcentered{} fato \textperiodcentered{} confiança 1.0 \textperiodcentered{} domínio fisica \textperiodcentered{} história 2 versões no jornal}

%CRISTAL {"arestas":[{"alvo":"curvas_rotacao_rubin","confianca":1.0,"epistemico":"fato","origem":"wiki","rel":"explica"},{"alvo":"materia_escura","confianca":1.0,"epistemico":"fato","origem":"wiki","rel":"contradiz"},{"alvo":"cosmologia_quantica_hub","confianca":1.0,"epistemico":"fato","origem":"wiki","rel":"parte_de"}],"confianca":1.0,"contraexemplos":[],"descricao":"Abaixo de uma aceleração-limite a₀=1,2×10⁻¹⁰ m/s², a dinâmica desvia de Newton (a≈√(a₀·a_N)), eliminando a matéria escura em galáxias. Explica notavelmente as curvas de rotação (Tully-Fisher bariônica), MAS falha em aglomerados (déficit residual 2-3×, Bullet Cluster) e NÃO reproduz o 3º pico do CMB nem a formação de estrutura sem reintroduzir algo tipo-CDM. GW170817 (v=c) restringiu variantes tensor-vetor-escalar.","epistemico":"hipotese","exemplos":[],"id":"mond_gravidade_modificada","memoria":"semantica","meta":{"autor":"broca-so","dominio":"cosmologia_dados","fonte":"Milgrom 1983 ApJ 270,365; TeVeS Bekenstein 2004"},"origem":"wiki","palavras_chave":[],"sinonimos":[],"tipo":"conceito","titulo":"MOND — explica galáxias, falha em aglomerados"}
\section{MOND — explica galáxias, falha em aglomerados}
Abaixo de uma aceleração-limite a\ensuremath{_{0}}=1,2$\times$10\ensuremath{^{-10}} m/s\ensuremath{^{2}}, a dinâmica desvia de Newton (a\ensuremath{\approx}\ensuremath{\sqrt{\,}}(a\ensuremath{_{0}}\textperiodcentered a\_N)), eliminando a matéria escura em galáxias. Explica notavelmente as curvas de rotação (Tully-Fisher bariônica), MAS falha em aglomerados (déficit residual 2-3$\times$, Bullet Cluster) e NÃO reproduz o 3º pico do CMB nem a formação de estrutura sem reintroduzir algo tipo-CDM. GW170817 (v=c) restringiu variantes tensor-vetor-escalar.
\prov{relações: explica curvas\_rotacao\_rubin; contradiz materia\_escura; parte\_de cosmologia\_quantica\_hub}
\prov{proveniência: wiki \textperiodcentered{} Milgrom 1983 ApJ 270,365; TeVeS Bekenstein 2004 \textperiodcentered{} hipotese \textperiodcentered{} confiança 1.0 \textperiodcentered{} domínio cosmologia\_dados \textperiodcentered{} história 4 versões no jornal}

%CRISTAL {"arestas":[{"alvo":"corpo_estelar","confianca":1.0,"epistemico":"fato","origem":"wiki","rel":"usa"},{"alvo":"joalheria_hub","confianca":1.0,"epistemico":"fato","origem":"wiki","rel":"parte_de"}],"confianca":1.0,"contraexemplos":[],"descricao":"A morfologia matemática (Serra/Matheron) sobre formas É o corpo estelar em ato: a DILATAÇÃO A⊕B={a+b} é a SOMA DE MINKOWSKI = o ⊕ da reta mineral aplicado a conjuntos de pontos. Com a erosão (o adjunto), a abertura A∘B=(A⊖B)⊕B (tira as saliências finas), o fechamento e a reconstrução (o filtro conexo), temos os operadores de forma. Aplicação: TIRAR O REFLEXO do cristal (a abertura remove o espelho, o corpo fica). linguagem/morfologia_estelar.py (4/4). Traduz-se para a BAI pelo Teorema: dilatação=propagação, erosão=atenuação, abertura=poda, estruturante=orçamento.","epistemico":"fato","exemplos":[],"id":"morfologia_matematica_estelar","memoria":"semantica","meta":{"dominio":"joalheria","fonte":"maquina de corte"},"origem":"wiki","palavras_chave":[],"sinonimos":[],"tipo":"conceito","titulo":"A Morfologia Matemática Estelar (a dilatação é o ⊕ do corpo estelar)"}
\section{A Morfologia Matemática Estelar (a dilatação é o \ensuremath{\oplus} do corpo estelar)}
A morfologia matemática (Serra/Matheron) sobre formas É o corpo estelar em ato: a DILATAÇÃO A\ensuremath{\oplus}B=\{a+b\} é a SOMA DE MINKOWSKI = o \ensuremath{\oplus} da reta mineral aplicado a conjuntos de pontos. Com a erosão (o adjunto), a abertura A\ensuremath{\circ}B=(A\ensuremath{\ominus}B)\ensuremath{\oplus}B (tira as saliências finas), o fechamento e a reconstrução (o filtro conexo), temos os operadores de forma. Aplicação: TIRAR O REFLEXO do cristal (a abertura remove o espelho, o corpo fica). linguagem/morfologia\_estelar.py (4/4). Traduz-se para a BAI pelo Teorema: dilatação=propagação, erosão=atenuação, abertura=poda, estruturante=orçamento.
\prov{relações: usa corpo\_estelar; parte\_de joalheria\_hub}
\prov{proveniência: wiki \textperiodcentered{} maquina de corte \textperiodcentered{} fato \textperiodcentered{} confiança 1.0 \textperiodcentered{} domínio joalheria \textperiodcentered{} história 153 versões no jornal}

%CRISTAL {"arestas":[{"alvo":"problema_hierarquia_naturalidade","confianca":1.0,"epistemico":"fato","origem":"wiki","rel":"explica"},{"alvo":"problema_constante_cosmologica","confianca":1.0,"epistemico":"fato","origem":"wiki","rel":"relaciona"},{"alvo":"cosmologia_quantica_hub","confianca":1.0,"epistemico":"fato","origem":"wiki","rel":"parte_de"}],"confianca":1.0,"contraexemplos":[],"descricao":"A ausência de QUALQUER nova física na escala TeV (nem SUSY, nem dimensões extras, nem compositeness, nem ressonâncias) contradiz décadas de expectativa de que a naturalidade exigiria novas partículas junto ao Higgs. O ajuste fino do Higgs (e de Λ) segue sem explicação dinâmica. Respostas em disputa: abandonar naturalidade (antrópico/paisagem, split SUSY); naturalidade sutil (relaxion, Twin Higgs); ou nova física acima do alcance. É o resultado nulo mais importante do LHC.","epistemico":"fato","exemplos":[],"id":"naturalidade_em_crise","memoria":"semantica","meta":{"autor":"broca-so","dominio":"fisica","fonte":"consenso pós-LHC Run 2; Scientific American 2023"},"origem":"wiki","palavras_chave":[],"sinonimos":[],"tipo":"conceito","titulo":"A crise da naturalidade (o resultado nulo do LHC)"}
\section{A crise da naturalidade (o resultado nulo do LHC)}
A ausência de QUALQUER nova física na escala TeV (nem SUSY, nem dimensões extras, nem compositeness, nem ressonâncias) contradiz décadas de expectativa de que a naturalidade exigiria novas partículas junto ao Higgs. O ajuste fino do Higgs (e de \ensuremath{\Lambda}) segue sem explicação dinâmica. Respostas em disputa: abandonar naturalidade (antrópico/paisagem, split SUSY); naturalidade sutil (relaxion, Twin Higgs); ou nova física acima do alcance. É o resultado nulo mais importante do LHC.
\prov{relações: explica problema\_hierarquia\_naturalidade; relaciona problema\_constante\_cosmologica; parte\_de cosmologia\_quantica\_hub}
\prov{proveniência: wiki \textperiodcentered{} consenso pós-LHC Run 2; Scientific American 2023 \textperiodcentered{} fato \textperiodcentered{} confiança 1.0 \textperiodcentered{} domínio fisica \textperiodcentered{} história 4 versões no jornal}

%CRISTAL {"arestas":[{"alvo":"neuronio_stream_bit","confianca":1.0,"epistemico":"fato","origem":"wiki","rel":"explica"},{"alvo":"metal_broca_so","confianca":1.0,"epistemico":"fato","origem":"wiki","rel":"parte_de"}],"confianca":1.0,"contraexemplos":[],"descricao":"neuronᵢo.c implementa o átomo: por byte, e += popcount(byte & 0xAA) (bits pares) e o += popcount(byte & 0x55) (ímpares) via _builtinpopcount; devolve Worde+o, e. Processa por fgetc em stream ou por buffer — NÃO usa mmap; I/O é fopen/fgetc/fread. É o único 'hash' real do metal (o operador de Fibonacci A=[[1,1],[1,0]]).","epistemico":"fato","exemplos":[],"id":"neuron_io_popcnt","memoria":"semantica","meta":{"autor":"broca-so","dominio":"metal","fonte":"ula/neuron_io.c:scan_byte; code/neuronio.c"},"origem":"wiki","palavras_chave":[],"sinonimos":[],"tipo":"conceito","titulo":"O Neurônio no Metal (popcnt, sem SHA, sem mmap)"}
\section{O Neurônio no Metal (popcnt, sem SHA, sem mmap)}
neuron\ensuremath{_{i}}o.c implementa o átomo: por byte, e += popcount(byte \& 0xAA) (bits pares) e o += popcount(byte \& 0x55) (ímpares) via \_builtinpopcount; devolve Worde+o, e. Processa por fgetc em stream ou por buffer — NÃO usa mmap; I/O é fopen/fgetc/fread. É o único 'hash' real do metal (o operador de Fibonacci A=[[1,1],[1,0]]).
\prov{relações: explica neuronio\_stream\_bit; parte\_de metal\_broca\_so}
\prov{proveniência: wiki \textperiodcentered{} ula/neuron\_io.c:scan\_byte; code/neuronio.c \textperiodcentered{} fato \textperiodcentered{} confiança 1.0 \textperiodcentered{} domínio metal \textperiodcentered{} história 4 versões no jornal}

%CRISTAL {"arestas":[{"alvo":"koide_z3_geracoes","confianca":1.0,"epistemico":"fato","origem":"wiki","rel":"usa"},{"alvo":"pmns_desalinhamento_s3","confianca":1.0,"epistemico":"fato","origem":"wiki","rel":"usa"},{"alvo":"mecanismo_seesaw","confianca":1.0,"epistemico":"fato","origem":"wiki","rel":"relaciona"},{"alvo":"dirac_majorana_0nbb","confianca":1.0,"epistemico":"fato","origem":"wiki","rel":"relaciona"}],"confianca":1.0,"contraexemplos":[],"descricao":"O neutrino é o setor anômalo em DOIS testes independentes que não se falam: (1) o Koide FALHA nele — Q_ν∈[0,33;0,57] (ordenamento normal), não atinge o ⅔ que os léptons carregados exibem a 0,001%; (2) a PMNS é GRANDE por causa dele. Lido como 'a mesma assinatura de Majorana/seesaw vista por dois experimentos'. RESSALVA HONESTA: nenhum termo de massa de Majorana nem mecanismo seesaw é CONSTRUÍDO na álgebra — é assinatura interpretativa, não mecanismo. Candidato natural a ligar ao SO(10), mas não derivado.","epistemico":"inferencia","exemplos":[],"id":"neutrino_anomalo_dois_testes","memoria":"semantica","meta":{"autor":"broca-so","dominio":"cosmologia","fonte":"paper_4_particulas.tex obs:mistura-S3, obs:geracoes"},"origem":"wiki","palavras_chave":[],"sinonimos":[],"tipo":"conceito","titulo":"O ν é o setor anômalo em dois testes (a assinatura Majorana/seesaw)"}
\section{O \ensuremath{\nu} é o setor anômalo em dois testes (a assinatura Majorana/seesaw)}
O neutrino é o setor anômalo em DOIS testes independentes que não se falam: (1) o Koide FALHA nele — Q\_\ensuremath{\nu}\ensuremath{\in}[0,33;0,57] (ordenamento normal), não atinge o \ensuremath{\tfrac23} que os léptons carregados exibem a 0,001\%; (2) a PMNS é GRANDE por causa dele. Lido como 'a mesma assinatura de Majorana/seesaw vista por dois experimentos'. RESSALVA HONESTA: nenhum termo de massa de Majorana nem mecanismo seesaw é CONSTRUÍDO na álgebra — é assinatura interpretativa, não mecanismo. Candidato natural a ligar ao SO(10), mas não derivado.
\prov{relações: usa koide\_z3\_geracoes; usa pmns\_desalinhamento\_s3; relaciona mecanismo\_seesaw; relaciona dirac\_majorana\_0nbb}
\prov{proveniência: wiki \textperiodcentered{} paper\_4\_particulas.tex obs:mistura-S3, obs:geracoes \textperiodcentered{} inferencia \textperiodcentered{} confiança 1.0 \textperiodcentered{} domínio cosmologia \textperiodcentered{} história 5 versões no jornal}

%CRISTAL {"arestas":[{"alvo":"quiralidade_lorentz_sl2c","confianca":1.0,"epistemico":"fato","origem":"wiki","rel":"especializa"},{"alvo":"isospin_hipercarga_algebrico","confianca":1.0,"epistemico":"fato","origem":"wiki","rel":"usa"}],"confianca":1.0,"contraexemplos":[],"descricao":"A quiralidade do neutrino não é um número interno em falta — é a de LORENTZ. Cor/carga (de 𝕆) comutam com todo sl(2,ℂ) (comutador=0) → vetoriais; o isospin fraco vive em ℍ, o fator que carrega Lorentz, e T₃ transforma sob sl(2,ℂ) (comutador=5,66≠0) → só a força fraca pode ser quiral (o ν destro seria singleto T₃=0). [D] 'só a fraca é quiral' (partilha do fator ℍ); [N] a escolha do LADO (por que canhotos) é POSTA — ponto aberto partilhado com Furey/Boyle.","epistemico":"inferencia","exemplos":[],"id":"neutrino_quiralidade_lorentz","memoria":"semantica","meta":{"autor":"broca-so","dominio":"cosmologia","fonte":"paper_4_particulas.tex obs:quiralidade-tese"},"origem":"wiki","palavras_chave":[],"sinonimos":[],"tipo":"conceito","titulo":"A quiralidade do ν é de Lorentz (só a fraca é quiral)"}
\section{A quiralidade do \ensuremath{\nu} é de Lorentz (só a fraca é quiral)}
A quiralidade do neutrino não é um número interno em falta — é a de LORENTZ. Cor/carga (de \ensuremath{\mathbb{O}}) comutam com todo sl(2,\ensuremath{\mathbb{C}}) (comutador=0) \ensuremath{\to} vetoriais; o isospin fraco vive em \ensuremath{\mathbb{H}}, o fator que carrega Lorentz, e T\ensuremath{_{3}} transforma sob sl(2,\ensuremath{\mathbb{C}}) (comutador=5,66\ensuremath{\ne}0) \ensuremath{\to} só a força fraca pode ser quiral (o \ensuremath{\nu} destro seria singleto T\ensuremath{_{3}}=0). [D] 'só a fraca é quiral' (partilha do fator \ensuremath{\mathbb{H}}); [N] a escolha do LADO (por que canhotos) é POSTA — ponto aberto partilhado com Furey/Boyle.
\prov{relações: especializa quiralidade\_lorentz\_sl2c; usa isospin\_hipercarga\_algebrico}
\prov{proveniência: wiki \textperiodcentered{} paper\_4\_particulas.tex obs:quiralidade-tese \textperiodcentered{} inferencia \textperiodcentered{} confiança 1.0 \textperiodcentered{} domínio cosmologia \textperiodcentered{} história 3 versões no jornal}

%CRISTAL {"arestas":[{"alvo":"mecanismo_seesaw","confianca":1.0,"epistemico":"fato","origem":"wiki","rel":"relaciona"},{"alvo":"so10_pati_salam","confianca":1.0,"epistemico":"fato","origem":"wiki","rel":"relaciona"},{"alvo":"neutrino_anomalo_dois_testes","confianca":1.0,"epistemico":"fato","origem":"wiki","rel":"relaciona"}],"confianca":1.0,"contraexemplos":[],"descricao":"O que a álgebra de Lopes NÃO entrega para o neutrino: um ν_R/singleto estéril explícito, um termo de massa de Majorana, o mecanismo seesaw, o '2/3' de Koide, a forma tribimaximal exata e δ_CP — todos fenomenologia [N] ou programa em aberto [I]. O mergulho eletrofraco por-geração é declarado 'programa em aberto'. Honestidade: a assinatura (ν anômalo) é apontada, o mecanismo não é construído.","epistemico":"inferencia","exemplos":[],"id":"neutrino_r_ausente_lopes","memoria":"semantica","meta":{"autor":"broca-so","dominio":"cosmologia","fonte":"paper_4_particulas.tex obs:geracoes (mergulho em aberto)"},"origem":"wiki","palavras_chave":[],"sinonimos":[],"tipo":"conceito","titulo":"Lacuna: ν de mão-direita / seesaw não construídos"}
\section{Lacuna: \ensuremath{\nu} de mão-direita / seesaw não construídos}
O que a álgebra de Lopes NÃO entrega para o neutrino: um \ensuremath{\nu}\_R/singleto estéril explícito, um termo de massa de Majorana, o mecanismo seesaw, o '2/3' de Koide, a forma tribimaximal exata e \ensuremath{\delta}\_CP — todos fenomenologia [N] ou programa em aberto [I]. O mergulho eletrofraco por-geração é declarado 'programa em aberto'. Honestidade: a assinatura (\ensuremath{\nu} anômalo) é apontada, o mecanismo não é construído.
\prov{relações: relaciona mecanismo\_seesaw; relaciona so10\_pati\_salam; relaciona neutrino\_anomalo\_dois\_testes}
\prov{proveniência: wiki \textperiodcentered{} paper\_4\_particulas.tex obs:geracoes (mergulho em aberto) \textperiodcentered{} inferencia \textperiodcentered{} confiança 1.0 \textperiodcentered{} domínio cosmologia \textperiodcentered{} história 4 versões no jornal}

%CRISTAL {"arestas":[{"alvo":"carga_q_n_sobre_3","confianca":1.0,"epistemico":"fato","origem":"wiki","rel":"usa"},{"alvo":"materia_modelo_padrao","confianca":1.0,"epistemico":"fato","origem":"wiki","rel":"explica"},{"alvo":"grupo_padrao_cxho","confianca":1.0,"epistemico":"fato","origem":"wiki","rel":"parte_de"},{"alvo":"oscilacao_neutrinos","confianca":1.0,"epistemico":"fato","origem":"wiki","rel":"relaciona"},{"alvo":"cosmologia_quantica_hub","confianca":1.0,"epistemico":"fato","origem":"wiki","rel":"parte_de"}],"confianca":1.0,"contraexemplos":[],"descricao":"Os seis Lᵢ octônicos montam três escadas αa=½(Lp+iLq); o número N=Σαa†αa tem espectro 0,1,1,1,2,2,2,3, Q=N/3. O ideal mínimo Cℓ(6)ω tem 8 estados e o VÁCUO ω (ocupação N=0, Q=0) É o ν; seguem o tripleto d̄ (Q=⅓), o anti-tripleto u (Q=⅔) e o e⁺ (Q=1). O ν é o estado mais simples — a ausência de ocupação. [D] o vácuo N=0/Q=0 é forçado pelo número; [J] identificar vácuo≡ν físico é correspondência estrutural.","epistemico":"inferencia","exemplos":[],"id":"neutrino_vacuo_ideal_minimo","memoria":"semantica","meta":{"autor":"broca-so","dominio":"cosmologia","fonte":"paper_4_particulas.tex obs:cxo-cargas"},"origem":"wiki","palavras_chave":[],"sinonimos":[],"tipo":"conceito","titulo":"O neutrino é o vácuo do ideal mínimo (ℂ⊗𝕆=Cℓ(6))"}
\section{O neutrino é o vácuo do ideal mínimo (\ensuremath{\mathbb{C}}\ensuremath{\otimes}\ensuremath{\mathbb{O}}=C\ensuremath{\ell}(6))}
Os seis L\ensuremath{_{i}} octônicos montam três escadas \ensuremath{\alpha}a=\ensuremath{\tfrac12}(Lp+iLq); o número N=\ensuremath{\Sigma}\ensuremath{\alpha}a\ensuremath{\dagger}\ensuremath{\alpha}a tem espectro 0,1,1,1,2,2,2,3, Q=N/3. O ideal mínimo C\ensuremath{\ell}(6)\ensuremath{\omega} tem 8 estados e o VÁCUO \ensuremath{\omega} (ocupação N=0, Q=0) É o \ensuremath{\nu}; seguem o tripleto d\ensuremath{} (Q=\ensuremath{\tfrac13}), o anti-tripleto u (Q=\ensuremath{\tfrac23}) e o e\ensuremath{^{+}} (Q=1). O \ensuremath{\nu} é o estado mais simples — a ausência de ocupação. [D] o vácuo N=0/Q=0 é forçado pelo número; [J] identificar vácuo\ensuremath{\equiv}\ensuremath{\nu} físico é correspondência estrutural.
\prov{relações: usa carga\_q\_n\_sobre\_3; explica materia\_modelo\_padrao; parte\_de grupo\_padrao\_cxho; relaciona oscilacao\_neutrinos; parte\_de cosmologia\_quantica\_hub}
\prov{proveniência: wiki \textperiodcentered{} paper\_4\_particulas.tex obs:cxo-cargas \textperiodcentered{} inferencia \textperiodcentered{} confiança 1.0 \textperiodcentered{} domínio cosmologia \textperiodcentered{} história 7 versões no jornal}

%CRISTAL {"arestas":[{"alvo":"oscilacao_neutrinos","confianca":1.0,"epistemico":"fato","origem":"wiki","rel":"relaciona"}],"confianca":1.0,"contraexemplos":[],"descricao":"Hipotéticos neutrinos que não interagem via força fraca (singletos do MP), misturando só por oscilação; motivados por anomalias de curta distância (LSND, MiniBooNE 4,8σ, reator, gálio) sugerindo Δm²41~1 eV². MicroBooNE (dez 2025) exclui a interpretação de estéril único de LSND/MiniBooNE a 95% e grande parte do gálio. CONTROVERSO, tendendo a DESFAVORECIDO; sem confirmação.","epistemico":"hipotese","exemplos":[],"id":"neutrinos_estereis","memoria":"semantica","meta":{"autor":"broca-so","dominio":"fisica","fonte":"MiniBooNE 2018 arXiv:1805.12028; MicroBooNE 2025 arXiv:2512.07159"},"origem":"wiki","palavras_chave":[],"sinonimos":[],"tipo":"conceito","titulo":"Neutrinos estéreis (desfavorecidos)"}
\section{Neutrinos estéreis (desfavorecidos)}
Hipotéticos neutrinos que não interagem via força fraca (singletos do MP), misturando só por oscilação; motivados por anomalias de curta distância (LSND, MiniBooNE 4,8\ensuremath{\sigma}, reator, gálio) sugerindo \ensuremath{\Delta}m\ensuremath{^{2}}41\~{}1 eV\ensuremath{^{2}}. MicroBooNE (dez 2025) exclui a interpretação de estéril único de LSND/MiniBooNE a 95\% e grande parte do gálio. CONTROVERSO, tendendo a DESFAVORECIDO; sem confirmação.
\prov{relações: relaciona oscilacao\_neutrinos}
\prov{proveniência: wiki \textperiodcentered{} MiniBooNE 2018 arXiv:1805.12028; MicroBooNE 2025 arXiv:2512.07159 \textperiodcentered{} hipotese \textperiodcentered{} confiança 1.0 \textperiodcentered{} domínio fisica \textperiodcentered{} história 2 versões no jornal}

%CRISTAL {"arestas":[{"alvo":"teorema_nao_clonagem","confianca":1.0,"epistemico":"fato","origem":"wiki","rel":"especializa"},{"alvo":"a_violacao_de_bell_e_a_conservacao_ac21","confianca":1.0,"epistemico":"fato","origem":"wiki","rel":"usa"},{"alvo":"associador_nao_tensoriza","confianca":1.0,"epistemico":"fato","origem":"wiki","rel":"relaciona"},{"alvo":"cosmologia_quantica_hub","confianca":1.0,"epistemico":"fato","origem":"wiki","rel":"parte_de"}],"confianca":1.0,"contraexemplos":[],"descricao":"Achado fino: não há cloning isométrico em NENHUMA álgebra da família ⋆_s escalar — a razão é a INVARIÂNCIA do produto interno euclidiano Re(z̄⋆_s w)=⟨z,w⟩ (o termo s·a×b cai só na parte imaginária, nunca toca a real): a prova de Wootters-Zurek transporta LITERALMENTE (é a mesma prova), vale até nos 97% de álgebras proibidas. NÃO vem da não-associatividade nem da irreversibilidade. A não-associatividade só ABRE a janela: no regime hipercomplexo (s=s₀+s_v) o triplo escalar invade a parte real, a invariância quebra, e W-Z deixa de valer. [D] teoremas + [N] verif. <4×10⁻¹⁵.","epistemico":"inferencia","exemplos":[],"id":"no_cloning_invariancia_metrica","memoria":"semantica","meta":{"autor":"broca-so","dominio":"cosmologia","fonte":"paper_G_no_cloning.tex §3-4,§7"},"origem":"wiki","palavras_chave":[],"sinonimos":[],"tipo":"conceito","titulo":"Não-clonagem de Lopes: da invariância métrica (não da não-assoc.)"}
\section{Não-clonagem de Lopes: da invariância métrica (não da não-assoc.)}
Achado fino: não há cloning isométrico em NENHUMA álgebra da família \ensuremath{\star}\_s escalar — a razão é a INVARIÂNCIA do produto interno euclidiano Re(z\ensuremath{}\ensuremath{\star}\_s w)=\ensuremath{\langle}z,w\ensuremath{\rangle} (o termo s\textperiodcentered a$\times$b cai só na parte imaginária, nunca toca a real): a prova de Wootters-Zurek transporta LITERALMENTE (é a mesma prova), vale até nos 97\% de álgebras proibidas. NÃO vem da não-associatividade nem da irreversibilidade. A não-associatividade só ABRE a janela: no regime hipercomplexo (s=s\ensuremath{_{0}}+s\_v) o triplo escalar invade a parte real, a invariância quebra, e W-Z deixa de valer. [D] teoremas + [N] verif. <4$\times$10\ensuremath{^{-15}}.
\prov{relações: especializa teorema\_nao\_clonagem; usa a\_violacao\_de\_bell\_e\_a\_conservacao\_ac21; relaciona associador\_nao\_tensoriza; parte\_de cosmologia\_quantica\_hub}
\prov{proveniência: wiki \textperiodcentered{} paper\_G\_no\_cloning.tex §3-4,§7 \textperiodcentered{} inferencia \textperiodcentered{} confiança 1.0 \textperiodcentered{} domínio cosmologia \textperiodcentered{} história 5 versões no jornal}

%CRISTAL {"arestas":[{"alvo":"cosmologia_quantica_hub","confianca":1.0,"epistemico":"fato","origem":"wiki","rel":"parte_de"}],"confianca":1.0,"contraexemplos":[],"descricao":"Formação dos núcleos leves (D, ³He, ⁴He, ⁷Li) nos primeiros ~3 min; ⁴He≈25% da massa bariônica. Depende da razão bárion/fóton η. D e ⁴He concordam com η do CMB (fato); aberto: o 'problema do lítio' (⁷Li ~3× menor que o previsto).","epistemico":"fato","exemplos":[],"id":"nucleossintese_bbn","memoria":"semantica","meta":{"autor":"broca-so","dominio":"fisica","fonte":"Alpher-Bethe-Gamow 1948 ('αβγ')"},"origem":"wiki","palavras_chave":[],"sinonimos":[],"tipo":"conceito","titulo":"Nucleossíntese primordial (BBN)"}
\section{Nucleossíntese primordial (BBN)}
Formação dos núcleos leves (D, \ensuremath{^{3}}He, \ensuremath{^{4}}He, \ensuremath{^{7}}Li) nos primeiros \~{}3 min; \ensuremath{^{4}}He\ensuremath{\approx}25\% da massa bariônica. Depende da razão bárion/fóton \ensuremath{\eta}. D e \ensuremath{^{4}}He concordam com \ensuremath{\eta} do CMB (fato); aberto: o 'problema do lítio' (\ensuremath{^{7}}Li \~{}3$\times$ menor que o previsto).
\prov{relações: parte\_de cosmologia\_quantica\_hub}
\prov{proveniência: wiki \textperiodcentered{} Alpher-Bethe-Gamow 1948 ('\ensuremath{\alpha}\ensuremath{\beta}\ensuremath{\gamma}') \textperiodcentered{} fato \textperiodcentered{} confiança 1.0 \textperiodcentered{} domínio fisica \textperiodcentered{} história 2 versões no jornal}

%CRISTAL {"arestas":[{"alvo":"efeito_hall_quantico_inteiro","confianca":1.0,"epistemico":"fato","origem":"wiki","rel":"explica"},{"alvo":"fase_de_berry_canonica","confianca":1.0,"epistemico":"fato","origem":"wiki","rel":"usa"}],"confianca":1.0,"contraexemplos":[],"descricao":"A quantização de σ_xy é TOPOLÓGICA: o inteiro n é o número de Chern, a integral da curvatura de Berry das bandas ocupadas sobre a zona de Brillouin (σ_xy=(e²/h)·C, C=(1/2π)∮F d²k, C∈ℤ), invariante sob deformações suaves — daí a robustez. Modelo de Haldane (1988): estado Hall C=±1 SEM campo líquido. Thouless-Kohmoto-Nightingale-den Nijs 1982; Nobel 2016 (Thouless, Haldane, Kosterlitz).","epistemico":"fato","exemplos":[],"id":"numero_chern_tknn","memoria":"semantica","meta":{"autor":"broca-so","dominio":"fisica","fonte":"TKNN 1982 PRL 49,405; Haldane 1988; Nobel 2016"},"origem":"wiki","palavras_chave":[],"sinonimos":[],"tipo":"conceito","titulo":"O número de Chern / TKNN (a topologia da quantização)"}
\section{O número de Chern / TKNN (a topologia da quantização)}
A quantização de \ensuremath{\sigma}\_xy é TOPOLÓGICA: o inteiro n é o número de Chern, a integral da curvatura de Berry das bandas ocupadas sobre a zona de Brillouin (\ensuremath{\sigma}\_xy=(e\ensuremath{^{2}}/h)\textperiodcentered C, C=(1/2\ensuremath{\pi})\ensuremath{\oint}F d\ensuremath{^{2}}k, C\ensuremath{\in}\ensuremath{\mathbb{Z}}), invariante sob deformações suaves — daí a robustez. Modelo de Haldane (1988): estado Hall C=$\pm$1 SEM campo líquido. Thouless-Kohmoto-Nightingale-den Nijs 1982; Nobel 2016 (Thouless, Haldane, Kosterlitz).
\prov{relações: explica efeito\_hall\_quantico\_inteiro; usa fase\_de\_berry\_canonica}
\prov{proveniência: wiki \textperiodcentered{} TKNN 1982 PRL 49,405; Haldane 1988; Nobel 2016 \textperiodcentered{} fato \textperiodcentered{} confiança 1.0 \textperiodcentered{} domínio fisica \textperiodcentered{} história 3 versões no jornal}

%CRISTAL {"arestas":[{"alvo":"tecido_gex_floco","confianca":1.0,"epistemico":"fato","origem":"wiki","rel":"usa"},{"alvo":"olho_de_koch","confianca":1.0,"epistemico":"fato","origem":"wiki","rel":"explica"},{"alvo":"broca_olho_servico_gex44","confianca":1.0,"epistemico":"fato","origem":"wiki","rel":"usa"}],"confianca":1.0,"contraexemplos":[],"descricao":"olho_gex44.py abre um stream ssh do floco remoto, lê frames de 17B e ReguasEsteris.abrir (parábola) decodifica bit-a-bit de volta a bytes; imprime o sync 'bit=… floco=…B' verificando bit×17=bytes. É o receptor humano do canal GEX44.","epistemico":"fato","exemplos":[],"id":"olho_gex44_stream","memoria":"semantica","meta":{"autor":"broca-so","dominio":"metal","fonte":"neuronio/olho_gex44.py:main"},"origem":"wiki","palavras_chave":[],"sinonimos":[],"tipo":"conceito","titulo":"O Olho Local (decodifica o floco do GEX via ssh)"}
\section{O Olho Local (decodifica o floco do GEX via ssh)}
olho\_gex44.py abre um stream ssh do floco remoto, lê frames de 17B e ReguasEsteris.abrir (parábola) decodifica bit-a-bit de volta a bytes; imprime o sync 'bit=… floco=…B' verificando bit$\times$17=bytes. É o receptor humano do canal GEX44.
\prov{relações: usa tecido\_gex\_floco; explica olho\_de\_koch; usa broca\_olho\_servico\_gex44}
\prov{proveniência: wiki \textperiodcentered{} neuronio/olho\_gex44.py:main \textperiodcentered{} fato \textperiodcentered{} confiança 1.0 \textperiodcentered{} domínio metal \textperiodcentered{} história 4 versões no jornal}

%CRISTAL {"arestas":[{"alvo":"cosmologia_quantica_hub","confianca":1.0,"epistemico":"fato","origem":"wiki","rel":"parte_de"}],"confianca":1.0,"contraexemplos":[],"descricao":"Perturbações do espaço-tempo à velocidade da luz. 1ª detecção direta 14 set 2015: fusão de dois buracos negros (36+29→62 M☉, ~3 M☉c² radiados), forma casando com a RG. Nobel 2017. RG confirmada em campo forte/dinâmico.","epistemico":"fato","exemplos":[],"id":"ondas_gravitacionais_ligo","memoria":"semantica","meta":{"autor":"broca-so","dominio":"fisica","fonte":"Einstein 1916; Abbott+ (LIGO/Virgo) 2016"},"origem":"wiki","palavras_chave":[],"sinonimos":[],"tipo":"conceito","titulo":"Ondas gravitacionais (GW150914, LIGO 2015)"}
\section{Ondas gravitacionais (GW150914, LIGO 2015)}
Perturbações do espaço-tempo à velocidade da luz. 1ª detecção direta 14 set 2015: fusão de dois buracos negros (36+29\ensuremath{\to}62 M\ensuremath{\odot}, \~{}3 M\ensuremath{\odot}c\ensuremath{^{2}} radiados), forma casando com a RG. Nobel 2017. RG confirmada em campo forte/dinâmico.
\prov{relações: parte\_de cosmologia\_quantica\_hub}
\prov{proveniência: wiki \textperiodcentered{} Einstein 1916; Abbott+ (LIGO/Virgo) 2016 \textperiodcentered{} fato \textperiodcentered{} confiança 1.0 \textperiodcentered{} domínio fisica \textperiodcentered{} história 2 versões no jornal}

%CRISTAL {"arestas":[{"alvo":"render_fisico_gex","confianca":1.0,"epistemico":"fato","origem":"wiki","rel":"relaciona"},{"alvo":"primitivas_analiticas","confianca":1.0,"epistemico":"fato","origem":"wiki","rel":"usa"},{"alvo":"joalheria_hub","confianca":1.0,"epistemico":"fato","origem":"wiki","rel":"parte_de"}],"confianca":1.0,"contraexemplos":[],"descricao":"A LUZ É ONDA — o passo rumo ao ultra-realismo (os fótons colidindo com os átomos/cromóforos; a gema com VOLUME, não só casca; as FENDAS dos planos das facetas interferindo). A FUNDAÇÃO, primeiro: o experimento da DUPLA FENDA de Young. Luz coerente λ passa por duas fendas (separação d, largura a); no anteparo a D os campos SE SOMAM (superposição de Huygens-Fresnel, cada fenda uma faixa de fontes) e I=|campo|² dá as FRANJAS — máximos onde r₁−r₂=mλ, espaçamento Δy=λD/d. O CAMPO 2D mostra o MECANISMO de propagação: as ondas se cruzando, o leque de linhas de interferência. Verificado 5/5 (franjas, Δy=λD/d, central máximo, simetria, Δy∝1/d). λ=520nm (o ciano do Paraíba). A seguir: as fendas dos planos da gema → difração/interferência real, o raymarch interno, a colisão fóton-átomo. linguagem/optica_ondulatoria.py: dupla_fenda_tela, campo_2d, espacamento_franjas.","epistemico":"fato","exemplos":[],"id":"optica_ondulatoria_dupla_fenda","memoria":"semantica","meta":{"dominio":"joalheria","fonte":"maquina de corte"},"origem":"wiki","palavras_chave":[],"sinonimos":[],"tipo":"conceito","titulo":"Óptica ondulatória: a dupla fenda (Young) — a fundação do ultra-realismo"}
\section{Óptica ondulatória: a dupla fenda (Young) — a fundação do ultra-realismo}
A LUZ É ONDA — o passo rumo ao ultra-realismo (os fótons colidindo com os átomos/cromóforos; a gema com VOLUME, não só casca; as FENDAS dos planos das facetas interferindo). A FUNDAÇÃO, primeiro: o experimento da DUPLA FENDA de Young. Luz coerente \ensuremath{\lambda} passa por duas fendas (separação d, largura a); no anteparo a D os campos SE SOMAM (superposição de Huygens-Fresnel, cada fenda uma faixa de fontes) e I=|campo|\ensuremath{^{2}} dá as FRANJAS — máximos onde r\ensuremath{_{1}}\ensuremath{-}r\ensuremath{_{2}}=m\ensuremath{\lambda}, espaçamento \ensuremath{\Delta}y=\ensuremath{\lambda}D/d. O CAMPO 2D mostra o MECANISMO de propagação: as ondas se cruzando, o leque de linhas de interferência. Verificado 5/5 (franjas, \ensuremath{\Delta}y=\ensuremath{\lambda}D/d, central máximo, simetria, \ensuremath{\Delta}y\ensuremath{\propto}1/d). \ensuremath{\lambda}=520nm (o ciano do Paraíba). A seguir: as fendas dos planos da gema \ensuremath{\to} difração/interferência real, o raymarch interno, a colisão fóton-átomo. linguagem/optica\_ondulatoria.py: dupla\_fenda\_tela, campo\_2d, espacamento\_franjas.
\prov{relações: relaciona render\_fisico\_gex; usa primitivas\_analiticas; parte\_de joalheria\_hub}
\prov{proveniência: wiki \textperiodcentered{} maquina de corte \textperiodcentered{} fato \textperiodcentered{} confiança 1.0 \textperiodcentered{} domínio joalheria \textperiodcentered{} história 128 versões no jornal}

%CRISTAL {"arestas":[{"alvo":"criticidade","confianca":1.0,"epistemico":"fato","origem":"wiki","rel":"usa"},{"alvo":"lyapunov_da_orbita","confianca":1.0,"epistemico":"fato","origem":"wiki","rel":"usa"},{"alvo":"orbita_de_partida","confianca":1.0,"epistemico":"fato","origem":"wiki","rel":"relaciona"},{"alvo":"numero_de_salem","confianca":1.0,"epistemico":"fato","origem":"wiki","rel":"usa"},{"alvo":"corpo_mineral","confianca":1.0,"epistemico":"fato","origem":"wiki","rel":"usa"}],"confianca":1.0,"contraexemplos":[],"descricao":"A população de nêutrons N_{t+1}=k·N_t é uma órbita mineral com Lyapunov λ=log k. k<1 → λ<0 → CRISTAL (subcrítico, contrai/apaga); k=1 → λ=0 → BORDA/SALEM (crítico, ρ=1 no círculo — o reator controlado); k>1 → λ>0 → CAOS (supercrítico, a bomba). O reator é o número de Salem da energia: preso na borda. Verificado pela máquina mineral.","epistemico":"inferencia","exemplos":[],"id":"orbita_criticidade","memoria":"semantica","meta":{"classes_criticidade":{"critico":"borda/Salem (crítico)","subcritico":"cristal (subcrítico)","supercritico":"caos (supercrítico)"},"dominio":"fisica","fonte":"energia nuclear"},"origem":"wiki","palavras_chave":[],"sinonimos":[],"tipo":"conceito","titulo":"A Órbita da Criticidade (k = Lyapunov)"}
\section{A Órbita da Criticidade (k = Lyapunov)}
A população de nêutrons N\_\{t+1\}=k\textperiodcentered N\_t é uma órbita mineral com Lyapunov \ensuremath{\lambda}=log k. k<1 \ensuremath{\to} \ensuremath{\lambda}<0 \ensuremath{\to} CRISTAL (subcrítico, contrai/apaga); k=1 \ensuremath{\to} \ensuremath{\lambda}=0 \ensuremath{\to} BORDA/SALEM (crítico, \ensuremath{\rho}=1 no círculo — o reator controlado); k>1 \ensuremath{\to} \ensuremath{\lambda}>0 \ensuremath{\to} CAOS (supercrítico, a bomba). O reator é o número de Salem da energia: preso na borda. Verificado pela máquina mineral.
\prov{relações: usa criticidade; usa lyapunov\_da\_orbita; relaciona orbita\_de\_partida; usa numero\_de\_salem; usa corpo\_mineral}
\prov{proveniência: wiki \textperiodcentered{} energia nuclear \textperiodcentered{} inferencia \textperiodcentered{} confiança 1.0 \textperiodcentered{} domínio fisica \textperiodcentered{} história 6 versões no jornal}

%CRISTAL {"arestas":[{"alvo":"mineracao_pipe_hook","confianca":1.0,"epistemico":"fato","origem":"wiki","rel":"parte_de"},{"alvo":"numeros_metalicos","confianca":1.0,"epistemico":"fato","origem":"wiki","rel":"usa"}],"confianca":1.0,"contraexemplos":[],"descricao":"O floco percorre os metais pela órbita [x_0,x_1,…] do mapa de Gauss mineral: a cada passo EXTRAI um metal (a posição = a moeda) e deixa um RESÍDUO φ∈[0,1) (a fase). É a MESMA órbita da parábola (parabola_primos) e da geração-total (a fase→checkpoint) — a mineração É essa órbita.","epistemico":"fato","exemplos":[],"id":"orbita_mineracao","memoria":"semantica","meta":{"dominio":"mineracao_pipe","fonte":"mineracao_pipe hub"},"origem":"wiki","palavras_chave":[],"sinonimos":[],"tipo":"conceito","titulo":"A órbita da mineração (o mapa de Gauss mineral)"}
\section{A órbita da mineração (o mapa de Gauss mineral)}
O floco percorre os metais pela órbita [x\_0,x\_1,…] do mapa de Gauss mineral: a cada passo EXTRAI um metal (a posição = a moeda) e deixa um RESÍDUO \ensuremath{\varphi}\ensuremath{\in}[0,1) (a fase). É a MESMA órbita da parábola (parabola\_primos) e da geração-total (a fase\ensuremath{\to}checkpoint) — a mineração É essa órbita.
\prov{relações: parte\_de mineracao\_pipe\_hook; usa numeros\_metalicos}
\prov{proveniência: wiki \textperiodcentered{} mineracao\_pipe hub \textperiodcentered{} fato \textperiodcentered{} confiança 1.0 \textperiodcentered{} domínio mineracao\_pipe \textperiodcentered{} história 3 versões no jornal}

%CRISTAL {"arestas":[{"alvo":"bohr_niveis_energia","confianca":1.0,"epistemico":"fato","origem":"wiki","rel":"especializa"},{"alvo":"mecanica_quantica","confianca":1.0,"epistemico":"fato","origem":"wiki","rel":"parte_de"}],"confianca":1.0,"contraexemplos":[],"descricao":"Solução quântica do hidrogênio: estados por n (energia), ℓ (momento angular, ℓ<n → s,p,d,f) e m (projeção, −ℓ≤m≤ℓ); definem as subcamadas (1s,2s,2p…). Física estabelecida e confirmada — a âncora canônica de 'órbita/energia por estado'.","epistemico":"fato","exemplos":[],"id":"orbitais_hidrogenio_nlm","memoria":"semantica","meta":{"autor":"broca-so","dominio":"fisica","fonte":"Schrödinger 1926"},"origem":"wiki","palavras_chave":[],"sinonimos":[],"tipo":"conceito","titulo":"Orbitais do hidrogênio (números quânticos n,ℓ,m)"}
\section{Orbitais do hidrogênio (números quânticos n,\ensuremath{\ell},m)}
Solução quântica do hidrogênio: estados por n (energia), \ensuremath{\ell} (momento angular, \ensuremath{\ell}<n \ensuremath{\to} s,p,d,f) e m (projeção, \ensuremath{-}\ensuremath{\ell}\ensuremath{\le}m\ensuremath{\le}\ensuremath{\ell}); definem as subcamadas (1s,2s,2p…). Física estabelecida e confirmada — a âncora canônica de 'órbita/energia por estado'.
\prov{relações: especializa bohr\_niveis\_energia; parte\_de mecanica\_quantica}
\prov{proveniência: wiki \textperiodcentered{} Schrödinger 1926 \textperiodcentered{} fato \textperiodcentered{} confiança 1.0 \textperiodcentered{} domínio fisica \textperiodcentered{} história 3 versões no jornal}

%CRISTAL {"arestas":[{"alvo":"materia_modelo_padrao","confianca":1.0,"epistemico":"fato","origem":"wiki","rel":"relaciona"},{"alvo":"cosmologia_quantica_hub","confianca":1.0,"epistemico":"fato","origem":"wiki","rel":"parte_de"}],"confianca":1.0,"contraexemplos":[],"descricao":"Neutrinos criados num sabor (e,μ,τ) são detectados em outro após propagação — logo os autoestados de sabor ≠ de massa, e ao menos dois neutrinos têm massa (contra o MP original). P(ν_α→ν_β)∝sin²(1,27 Δm²[eV²]·L[km]/E[GeV]). Super-Kamiokande (atmosféricos, 1998) + SNO (solares, 2002, resolveu o 'problema dos neutrinos solares'). FATO — Nobel 2015 (Kajita, McDonald).","epistemico":"fato","exemplos":[],"id":"oscilacao_neutrinos","memoria":"semantica","meta":{"autor":"broca-so","dominio":"fisica","fonte":"Super-K 1998 hep-ex/9807003; SNO 2002; Nobel 2015"},"origem":"wiki","palavras_chave":[],"sinonimos":[],"tipo":"conceito","titulo":"Oscilações de neutrinos (têm massa)"}
\section{Oscilações de neutrinos (têm massa)}
Neutrinos criados num sabor (e,\ensuremath{\mu},\ensuremath{\tau}) são detectados em outro após propagação — logo os autoestados de sabor \ensuremath{\ne} de massa, e ao menos dois neutrinos têm massa (contra o MP original). P(\ensuremath{\nu}\_\ensuremath{\alpha}\ensuremath{\to}\ensuremath{\nu}\_\ensuremath{\beta})\ensuremath{\propto}sin\ensuremath{^{2}}(1,27 \ensuremath{\Delta}m\ensuremath{^{2}}[eV\ensuremath{^{2}}]\textperiodcentered L[km]/E[GeV]). Super-Kamiokande (atmosféricos, 1998) + SNO (solares, 2002, resolveu o 'problema dos neutrinos solares'). FATO — Nobel 2015 (Kajita, McDonald).
\prov{relações: relaciona materia\_modelo\_padrao; parte\_de cosmologia\_quantica\_hub}
\prov{proveniência: wiki \textperiodcentered{} Super-K 1998 hep-ex/9807003; SNO 2002; Nobel 2015 \textperiodcentered{} fato \textperiodcentered{} confiança 1.0 \textperiodcentered{} domínio fisica \textperiodcentered{} história 3 versões no jornal}

%CRISTAL {"arestas":[{"alvo":"pow_espectral_geometrico","confianca":1.0,"epistemico":"fato","origem":"wiki","rel":"usa"},{"alvo":"pipeline_stratum","confianca":1.0,"epistemico":"fato","origem":"wiki","rel":"parte_de"},{"alvo":"mineracao_hexagrama","confianca":1.0,"epistemico":"fato","origem":"wiki","rel":"usa"}],"confianca":1.0,"contraexemplos":[],"descricao":"O gato entra como SEMENTE: 1×/job, broca_borda.pow_seed resolve a cadeia BrocaPacote devolvendo nonce0, ν=modos, canal_c2, floco Φ, com iteracoes_brute=0 — e a varredura de nonces começa em nonce0 e avança por cursor. Roda fora do GIL num ProcessPool. ~52 ms/seed (ν=6). É o insight que dá o speedup 4,43×.","epistemico":"fato","exemplos":[],"id":"pacote_seed_pow","memoria":"semantica","meta":{"autor":"broca-so","dominio":"metal","fonte":"neuronio/stratum_metal.py:seed_pow_job/_garantir_pow_job"},"origem":"wiki","palavras_chave":[],"sinonimos":[],"tipo":"conceito","titulo":"PoW Espectral como Semente do Nonce (o gato)"}
\section{PoW Espectral como Semente do Nonce (o gato)}
O gato entra como SEMENTE: 1$\times$/job, broca\_borda.pow\_seed resolve a cadeia BrocaPacote devolvendo nonce0, \ensuremath{\nu}=modos, canal\_c2, floco \ensuremath{\Phi}, com iteracoes\_brute=0 — e a varredura de nonces começa em nonce0 e avança por cursor. Roda fora do GIL num ProcessPool. \~{}52 ms/seed (\ensuremath{\nu}=6). É o insight que dá o speedup 4,43$\times$.
\prov{relações: usa pow\_espectral\_geometrico; parte\_de pipeline\_stratum; usa mineracao\_hexagrama}
\prov{proveniência: wiki \textperiodcentered{} neuronio/stratum\_metal.py:seed\_pow\_job/\_garantir\_pow\_job \textperiodcentered{} fato \textperiodcentered{} confiança 1.0 \textperiodcentered{} domínio metal \textperiodcentered{} história 4 versões no jornal}

%CRISTAL {"arestas":[{"alvo":"cotacao_parabola","confianca":1.0,"epistemico":"fato","origem":"wiki","rel":"explica"},{"alvo":"metal_broca_so","confianca":1.0,"epistemico":"fato","origem":"wiki","rel":"parte_de"}],"confianca":1.0,"contraexemplos":[],"descricao":"Sobre duas words normalizadas, computa a fase² c² por Bhattacharyya e o calor por Hellinger em CAMINHOS INDEPENDENTES, e mede erro=|calor+c²−1|; 'ok' se erro<1e-12. O fecho da parábola é VERIFICADO, não imposto — é a assinatura de honestidade do metal.","epistemico":"fato","exemplos":[],"id":"parabola_audit_dual","memoria":"semantica","meta":{"autor":"broca-so","dominio":"metal","fonte":"ula/parabola.c:parabola_auditar"},"origem":"wiki","palavras_chave":[],"sinonimos":[],"tipo":"conceito","titulo":"Auditoria a+c²=1 (Hellinger ⊥ Bhattacharyya)"}
\section{Auditoria a+c\ensuremath{^{2}}=1 (Hellinger \ensuremath{\perp} Bhattacharyya)}
Sobre duas words normalizadas, computa a fase\ensuremath{^{2}} c\ensuremath{^{2}} por Bhattacharyya e o calor por Hellinger em CAMINHOS INDEPENDENTES, e mede erro=|calor+c\ensuremath{^{2}}\ensuremath{-}1|; 'ok' se erro<1e-12. O fecho da parábola é VERIFICADO, não imposto — é a assinatura de honestidade do metal.
\prov{relações: explica cotacao\_parabola; parte\_de metal\_broca\_so}
\prov{proveniência: wiki \textperiodcentered{} ula/parabola.c:parabola\_auditar \textperiodcentered{} fato \textperiodcentered{} confiança 1.0 \textperiodcentered{} domínio metal \textperiodcentered{} história 3 versões no jornal}

%CRISTAL {"arestas":[{"alvo":"parabola_audit_dual","confianca":1.0,"epistemico":"fato","origem":"wiki","rel":"usa"},{"alvo":"corpo_tropical","confianca":1.0,"epistemico":"fato","origem":"wiki","rel":"relaciona"}],"confianca":1.0,"contraexemplos":[],"descricao":"Cada passo é rotulado V (vivo/ressoa: c²>0,85 ou reversível), C (cristal/calor: a>0,55) ou F (fronteira). É a vida/morte matemática (corpo_dual/tropical) medida por passo no metal.","epistemico":"inferencia","exemplos":[],"id":"parabola_classe_vcf","memoria":"semantica","meta":{"autor":"broca-so","dominio":"metal","fonte":"ula/parabola.c:parabola_classe"},"origem":"wiki","palavras_chave":[],"sinonimos":[],"tipo":"conceito","titulo":"Classificação Vivo / Cristal / Fronteira"}
\section{Classificação Vivo / Cristal / Fronteira}
Cada passo é rotulado V (vivo/ressoa: c\ensuremath{^{2}}>0,85 ou reversível), C (cristal/calor: a>0,55) ou F (fronteira). É a vida/morte matemática (corpo\_dual/tropical) medida por passo no metal.
\prov{relações: usa parabola\_audit\_dual; relaciona corpo\_tropical}
\prov{proveniência: wiki \textperiodcentered{} ula/parabola.c:parabola\_classe \textperiodcentered{} inferencia \textperiodcentered{} confiança 1.0 \textperiodcentered{} domínio metal \textperiodcentered{} história 3 versões no jornal}

%CRISTAL {"arestas":[{"alvo":"parabola_audit_dual","confianca":1.0,"epistemico":"fato","origem":"wiki","rel":"usa"},{"alvo":"isa_metal_21_opcodes","confianca":1.0,"epistemico":"fato","origem":"wiki","rel":"usa"}],"confianca":1.0,"contraexemplos":[],"descricao":"NÃO minera hashes: varre a ISA (fuzz de extremos, pares aleatórios, programas reais) por até META_PASSOS=1e6 e grava em JSONL todo par onde total+e≤0 (degenerado, sem distribuição). Caça os pontos onde a máquina colapsa.","epistemico":"fato","exemplos":[],"id":"parabola_miner_degenerados","memoria":"semantica","meta":{"autor":"broca-so","dominio":"metal","fonte":"ula/parabola_miner.c:main"},"origem":"wiki","palavras_chave":[],"sinonimos":[],"tipo":"conceito","titulo":"Mineradora de Estados Degenerados"}
\section{Mineradora de Estados Degenerados}
NÃO minera hashes: varre a ISA (fuzz de extremos, pares aleatórios, programas reais) por até META\_PASSOS=1e6 e grava em JSONL todo par onde total+e\ensuremath{\le}0 (degenerado, sem distribuição). Caça os pontos onde a máquina colapsa.
\prov{relações: usa parabola\_audit\_dual; usa isa\_metal\_21\_opcodes}
\prov{proveniência: wiki \textperiodcentered{} ula/parabola\_miner.c:main \textperiodcentered{} fato \textperiodcentered{} confiança 1.0 \textperiodcentered{} domínio metal \textperiodcentered{} história 3 versões no jornal}

%CRISTAL {"arestas":[{"alvo":"alternatividade_dimensoes_criticas","confianca":1.0,"epistemico":"fato","origem":"wiki","rel":"especializa"},{"alvo":"spin_s_sym_quaternion","confianca":1.0,"epistemico":"fato","origem":"wiki","rel":"relaciona"},{"alvo":"nucleo_g2_reabsorcao","confianca":1.0,"epistemico":"fato","origem":"wiki","rel":"relaciona"},{"alvo":"nao_associatividade_amputacao","confianca":1.0,"epistemico":"fato","origem":"wiki","rel":"relaciona"},{"alvo":"corpo_mineral","confianca":1.0,"epistemico":"fato","origem":"wiki","rel":"usa"},{"alvo":"transformada_metalica","confianca":1.0,"epistemico":"fato","origem":"wiki","rel":"usa"},{"alvo":"geometria_do_caos","confianca":1.0,"epistemico":"fato","origem":"wiki","rel":"usa"}],"confianca":1.0,"contraexemplos":[],"descricao":"A física de partículas expandida à reta mineral — o canto D=3 da torre de Hurwitz (Estelar 𝓔ₛ norm-conservante, associativa, associador=0 ⟹ mass gap ZERO). PARTÍCULA = um modo (raiz) da cadeia: massa m=|log ρ|=|λ| (o Lyapunov do modo), spin/momento ~ θ, carga = ∏raízes = ±1 (o invariante a+c²=1). Estabilidade: Pisot=elétron ligado, Salem=fóton (sem massa, ρ=1), caos=ressonância. Verificado: metais são hidrogênio (núcleo β + 1 elétron −1/σ, neutros).","epistemico":"hipotese","exemplos":[],"id":"particulas_minerais","memoria":"semantica","meta":{"dominio":"fisica","espectro_ouro":[{"massa":0.4812,"theta_g":0.0,"tipo":"nucleo"},{"massa":0.4812,"theta_g":180.0,"tipo":"eletron"}],"fonte":"partículas minerais"},"origem":"wiki","palavras_chave":[],"sinonimos":[],"tipo":"conceito","titulo":"Partículas e Átomos Minerais (física no canto D=3)"}
\section{Partículas e Átomos Minerais (física no canto D=3)}
A física de partículas expandida à reta mineral — o canto D=3 da torre de Hurwitz (Estelar \ensuremath{\mathcal{E}}\ensuremath{_{s}} norm-conservante, associativa, associador=0 \ensuremath{\Longrightarrow} mass gap ZERO). PARTÍCULA = um modo (raiz) da cadeia: massa m=|log \ensuremath{\rho}|=|\ensuremath{\lambda}| (o Lyapunov do modo), spin/momento \~{} \ensuremath{\theta}, carga = \ensuremath{\prod}raízes = $\pm$1 (o invariante a+c\ensuremath{^{2}}=1). Estabilidade: Pisot=elétron ligado, Salem=fóton (sem massa, \ensuremath{\rho}=1), caos=ressonância. Verificado: metais são hidrogênio (núcleo \ensuremath{\beta} + 1 elétron \ensuremath{-}1/\ensuremath{\sigma}, neutros).
\prov{relações: especializa alternatividade\_dimensoes\_criticas; relaciona spin\_s\_sym\_quaternion; relaciona nucleo\_g2\_reabsorcao; relaciona nao\_associatividade\_amputacao; usa corpo\_mineral; usa transformada\_metalica; usa geometria\_do\_caos}
\prov{proveniência: wiki \textperiodcentered{} partículas minerais \textperiodcentered{} hipotese \textperiodcentered{} confiança 1.0 \textperiodcentered{} domínio fisica \textperiodcentered{} história 9 versões no jornal}

%CRISTAL {"arestas":[{"alvo":"hopfield_recursiva_koch","confianca":1.0,"epistemico":"fato","origem":"wiki","rel":"explica"},{"alvo":"olho_de_koch","confianca":1.0,"epistemico":"fato","origem":"wiki","rel":"usa"},{"alvo":"metal_broca_so","confianca":1.0,"epistemico":"fato","origem":"wiki","rel":"parte_de"}],"confianca":1.0,"contraexemplos":[],"descricao":"O gato lê um ARQUIVO byte a byte (decimação K-way), a borda de Koch (z²%1009) quantiza a leitura, cada módulo Hopfield recorda (s_i←sgn(ΣW_ij s_j), ≤64 iters) e a super-rede fecha na config verdadeira. Um único binário que percebe o mundo real de argv[1].","epistemico":"fato","exemplos":[],"id":"percebe_c_endtoend","memoria":"semantica","meta":{"autor":"broca-so","dominio":"metal","fonte":"code/percebe.c:main"},"origem":"wiki","palavras_chave":[],"sinonimos":[],"tipo":"conceito","titulo":"Percepção End-to-End no Metal (percebe.c)"}
\section{Percepção End-to-End no Metal (percebe.c)}
O gato lê um ARQUIVO byte a byte (decimação K-way), a borda de Koch (z\ensuremath{^{2}}\%1009) quantiza a leitura, cada módulo Hopfield recorda (s\_i\ensuremath{\leftarrow}sgn(\ensuremath{\Sigma}W\_ij s\_j), \ensuremath{\le}64 iters) e a super-rede fecha na config verdadeira. Um único binário que percebe o mundo real de argv[1].
\prov{relações: explica hopfield\_recursiva\_koch; usa olho\_de\_koch; parte\_de metal\_broca\_so}
\prov{proveniência: wiki \textperiodcentered{} code/percebe.c:main \textperiodcentered{} fato \textperiodcentered{} confiança 1.0 \textperiodcentered{} domínio metal \textperiodcentered{} história 4 versões no jornal}

%CRISTAL {"arestas":[{"alvo":"metal_broca_so","confianca":1.0,"epistemico":"fato","origem":"wiki","rel":"parte_de"},{"alvo":"prova_de_trabalho","confianca":1.0,"epistemico":"fato","origem":"wiki","rel":"usa"},{"alvo":"ciclo_uno","confianca":1.0,"epistemico":"fato","origem":"wiki","rel":"explica"},{"alvo":"mineracao_hexagrama","confianca":1.0,"epistemico":"fato","origem":"wiki","rel":"usa"}],"confianca":1.0,"contraexemplos":[],"descricao":"Pipeline de mineração BTC que conecta a um pool Stratum, monta o bloco no metal (C), semeia o nonce com PoW espectral (o gato), varre nonces e submete shares. Rodou de fato contra solo.ckpool.org:3333 (mainnet solo) — o stratum_log.jsonl tem mining.notify/set_difficulty reais capturados. README honesto: 'Lucro BTC: não'.","epistemico":"fato","exemplos":[],"id":"pipeline_stratum","memoria":"semantica","meta":{"autor":"broca-so","dominio":"metal","fonte":"neuronio/stratum_metal.py:StratumMetal.run; dados/stratum_log.jsonl"},"origem":"wiki","palavras_chave":[],"sinonimos":[],"tipo":"conceito","titulo":"O Pipeline Stratum (mineração BTC real)"}
\section{O Pipeline Stratum (mineração BTC real)}
Pipeline de mineração BTC que conecta a um pool Stratum, monta o bloco no metal (C), semeia o nonce com PoW espectral (o gato), varre nonces e submete shares. Rodou de fato contra solo.ckpool.org:3333 (mainnet solo) — o stratum\_log.jsonl tem mining.notify/set\_difficulty reais capturados. README honesto: 'Lucro BTC: não'.
\prov{relações: parte\_de metal\_broca\_so; usa prova\_de\_trabalho; explica ciclo\_uno; usa mineracao\_hexagrama}
\prov{proveniência: wiki \textperiodcentered{} neuronio/stratum\_metal.py:StratumMetal.run; dados/stratum\_log.jsonl \textperiodcentered{} fato \textperiodcentered{} confiança 1.0 \textperiodcentered{} domínio metal \textperiodcentered{} história 5 versões no jornal}

%CRISTAL {"arestas":[{"alvo":"correcao_erro_quantica","confianca":1.0,"epistemico":"fato","origem":"wiki","rel":"usa"}],"confianca":1.0,"contraexemplos":[],"descricao":"Supercondutores (IBM Heron r2 ~99,5-99,9% 2-qubit; Google Willow 105 qubits, QEC below-threshold); íons aprisionados (Quantinuum Helios nov/2025: 99,921% 2-qubit, dezenas de qubits lógicos); fotônica (Xanadu/PsiQuantum, perdas); átomos neutros (QuEra/Harvard: array ~3000 qubits, ~96 lógicos, operação >2h; Caltech 6100). Hardware = fato; máquina útil em escala = fronteira NÃO alcançada.","epistemico":"fato","exemplos":[],"id":"plataformas_qc_2025","memoria":"semantica","meta":{"autor":"broca-so","dominio":"fisica","fonte":"Google/IBM/Quantinuum/QuEra 2024-2026"},"origem":"wiki","palavras_chave":[],"sinonimos":[],"tipo":"conceito","titulo":"Plataformas de computação quântica (estado 2025)"}
\section{Plataformas de computação quântica (estado 2025)}
Supercondutores (IBM Heron r2 \~{}99,5-99,9\% 2-qubit; Google Willow 105 qubits, QEC below-threshold); íons aprisionados (Quantinuum Helios nov/2025: 99,921\% 2-qubit, dezenas de qubits lógicos); fotônica (Xanadu/PsiQuantum, perdas); átomos neutros (QuEra/Harvard: array \~{}3000 qubits, \~{}96 lógicos, operação >2h; Caltech 6100). Hardware = fato; máquina útil em escala = fronteira NÃO alcançada.
\prov{relações: usa correcao\_erro\_quantica}
\prov{proveniência: wiki \textperiodcentered{} Google/IBM/Quantinuum/QuEra 2024-2026 \textperiodcentered{} fato \textperiodcentered{} confiança 1.0 \textperiodcentered{} domínio fisica \textperiodcentered{} história 2 versões no jornal}

%CRISTAL {"arestas":[{"alvo":"matriz_pmns","confianca":1.0,"epistemico":"fato","origem":"wiki","rel":"explica"},{"alvo":"koide_z3_geracoes","confianca":1.0,"epistemico":"fato","origem":"wiki","rel":"usa"}],"confianca":1.0,"contraexemplos":[],"descricao":"Teorema (dedutivo): matrizes de massa que respeitam o Z₃ da trialidade são CIRCULANTES, diagonalizadas pela mesma matriz de Fourier F; dois setores S₃-circulantes dariam V=F†F=1 (mistura NULA). A mistura leptônica grande (PMNS≈1,36, ~4× a CKM) é o DESALINHAMENTO das quebras de S₃ entre os setores ℓ e ν — perto da tribimaximal (θ12→35,3°, θ23→45°), com θ13=8,6°≠0 a quebra residual. [D] o teorema mistura=desalinhamento (verificado a 10⁻¹⁶); [N] os valores TBM/ângulos são empíricos (TBM exata pediria A₄/S₄⊃S₃).","epistemico":"inferencia","exemplos":[],"id":"pmns_desalinhamento_s3","memoria":"semantica","meta":{"autor":"broca-so","dominio":"cosmologia","fonte":"paper_4_particulas.tex obs:mistura-S3"},"origem":"wiki","palavras_chave":[],"sinonimos":[],"tipo":"conceito","titulo":"PMNS = desalinhamento das quebras de S₃"}
\section{PMNS = desalinhamento das quebras de S\ensuremath{_{3}}}
Teorema (dedutivo): matrizes de massa que respeitam o Z\ensuremath{_{3}} da trialidade são CIRCULANTES, diagonalizadas pela mesma matriz de Fourier F; dois setores S\ensuremath{_{3}}-circulantes dariam V=F\ensuremath{\dagger}F=1 (mistura NULA). A mistura leptônica grande (PMNS\ensuremath{\approx}1,36, \~{}4$\times$ a CKM) é o DESALINHAMENTO das quebras de S\ensuremath{_{3}} entre os setores \ensuremath{\ell} e \ensuremath{\nu} — perto da tribimaximal (\ensuremath{\theta}12\ensuremath{\to}35,3$^{\circ}$, \ensuremath{\theta}23\ensuremath{\to}45$^{\circ}$), com \ensuremath{\theta}13=8,6$^{\circ}$\ensuremath{\ne}0 a quebra residual. [D] o teorema mistura=desalinhamento (verificado a 10\ensuremath{^{-16}}); [N] os valores TBM/ângulos são empíricos (TBM exata pediria A\ensuremath{_{4}}/S\ensuremath{_{4}}\ensuremath{\supset}S\ensuremath{_{3}}).
\prov{relações: explica matriz\_pmns; usa koide\_z3\_geracoes}
\prov{proveniência: wiki \textperiodcentered{} paper\_4\_particulas.tex obs:mistura-S3 \textperiodcentered{} inferencia \textperiodcentered{} confiança 1.0 \textperiodcentered{} domínio cosmologia \textperiodcentered{} história 3 versões no jornal}

%CRISTAL {"arestas":[{"alvo":"maquina_de_corte","confianca":1.0,"epistemico":"fato","origem":"wiki","rel":"parte_de"},{"alvo":"tabela_geometrias_espectrais","confianca":1.0,"epistemico":"fato","origem":"wiki","rel":"usa"},{"alvo":"corpo_estelar","confianca":1.0,"epistemico":"fato","origem":"wiki","rel":"usa"}],"confianca":1.0,"contraexemplos":[],"descricao":"A geometria 3D EMERGE da seleção das frequências, SEM INVENTAR: a frequência angular seleciona k (o polígono/estrela plano), e a TABELA DA MANDALA (a álgebra estelar, o chicote) dá o POLIEDRO 3D de simetria k — k=3 triângulo → TETRAEDRO, k=4 quadrado → CUBO, k=5 → DODECAEDRO (via φ, a unidade estelar), k→∞ (o círculo) → a esfera. VALIDAÇÃO LIMPA do método: reconstruir o TRILHÃO LISO via chicote de uma PROJEÇÃO 2D SINTÉTICA (a mandala k=3 limpa, r=1+a·cos3θ, sem o ruído do gabarito) — cada FATIA horizontal É a mandala escalada pelo PERFIL (pavilhão do culet ao rondiz, coroa do rondiz à mesa); sai LISO e FECHADO (variedade, toda aresta 2×). O corte 2D é só observação; o sólido vem da tabela, não de uma extrusão fabricada. As FACETAS (a assinatura do gabarito) vêm do PRODUTO ESTELAR ⊗ (La Hire) de DUAS curvas discretas em direções ORTOGONAIS: a mandala k=3 DISCRETIZADA em N=k·m cantos (angular — a estrela trigonal ν=k−1 ⊗ o refino m, compondo os mapas angulares) VEZES o PERFIL discreto em anéis (meridional), com anéis alternados meia-faceta pro zigue-zague brilhante. Cada retalho é uma FACE PLANA — o trilhão facetado da referência. O CHICOTE REFINA progressivamente (chicote_refina = corte-de-canto de Chaikin no perfil de fileiras): a cada nível arredonda as cristas e DOBRA os pontos — mais malha nas bordas, mais resolução, mais arredondado, sem perder o sólido fechado. trilhao_facetado(refino=n). linguagem/reconstrucao_3d.poliedro_da_tabela, trilhao_liso, trilhao_facetado, chicote_refina; filtro_faixa. A geometria (e as facetas) emergem de k via ⊕ e ⊗, e o chicote refina.","epistemico":"fato","exemplos":[],"id":"poligono_vira_poliedro","memoria":"semantica","meta":{"dominio":"joalheria","fonte":"maquina de corte"},"origem":"wiki","palavras_chave":[],"sinonimos":[],"tipo":"conceito","titulo":"A Geometria Emerge de k: o polígono plano vira o poliedro 3D (via a mandala)"}
\section{A Geometria Emerge de k: o polígono plano vira o poliedro 3D (via a mandala)}
A geometria 3D EMERGE da seleção das frequências, SEM INVENTAR: a frequência angular seleciona k (o polígono/estrela plano), e a TABELA DA MANDALA (a álgebra estelar, o chicote) dá o POLIEDRO 3D de simetria k — k=3 triângulo \ensuremath{\to} TETRAEDRO, k=4 quadrado \ensuremath{\to} CUBO, k=5 \ensuremath{\to} DODECAEDRO (via \ensuremath{\varphi}, a unidade estelar), k\ensuremath{\to}\ensuremath{\infty} (o círculo) \ensuremath{\to} a esfera. VALIDAÇÃO LIMPA do método: reconstruir o TRILHÃO LISO via chicote de uma PROJEÇÃO 2D SINTÉTICA (a mandala k=3 limpa, r=1+a\textperiodcentered cos3\ensuremath{\theta}, sem o ruído do gabarito) — cada FATIA horizontal É a mandala escalada pelo PERFIL (pavilhão do culet ao rondiz, coroa do rondiz à mesa); sai LISO e FECHADO (variedade, toda aresta 2$\times$). O corte 2D é só observação; o sólido vem da tabela, não de uma extrusão fabricada. As FACETAS (a assinatura do gabarito) vêm do PRODUTO ESTELAR \ensuremath{\otimes} (La Hire) de DUAS curvas discretas em direções ORTOGONAIS: a mandala k=3 DISCRETIZADA em N=k\textperiodcentered m cantos (angular — a estrela trigonal \ensuremath{\nu}=k\ensuremath{-}1 \ensuremath{\otimes} o refino m, compondo os mapas angulares) VEZES o PERFIL discreto em anéis (meridional), com anéis alternados meia-faceta pro zigue-zague brilhante. Cada retalho é uma FACE PLANA — o trilhão facetado da referência. O CHICOTE REFINA progressivamente (chicote\_refina = corte-de-canto de Chaikin no perfil de fileiras): a cada nível arredonda as cristas e DOBRA os pontos — mais malha nas bordas, mais resolução, mais arredondado, sem perder o sólido fechado. trilhao\_facetado(refino=n). linguagem/reconstrucao\_3d.poliedro\_da\_tabela, trilhao\_liso, trilhao\_facetado, chicote\_refina; filtro\_faixa. A geometria (e as facetas) emergem de k via \ensuremath{\oplus} e \ensuremath{\otimes}, e o chicote refina.
\prov{relações: parte\_de maquina\_de\_corte; usa tabela\_geometrias\_espectrais; usa corpo\_estelar}
\prov{proveniência: wiki \textperiodcentered{} maquina de corte \textperiodcentered{} fato \textperiodcentered{} confiança 1.0 \textperiodcentered{} domínio joalheria \textperiodcentered{} história 180 versões no jornal}

%CRISTAL {"arestas":[{"alvo":"ultra_realismo_gema","confianca":1.0,"epistemico":"fato","origem":"wiki","rel":"relaciona"},{"alvo":"morfologia_matematica_estelar","confianca":1.0,"epistemico":"fato","origem":"wiki","rel":"usa"},{"alvo":"corpo_estelar","confianca":1.0,"epistemico":"fato","origem":"wiki","rel":"usa"},{"alvo":"joalheria_hub","confianca":1.0,"epistemico":"fato","origem":"wiki","rel":"parte_de"}],"confianca":1.0,"contraexemplos":[],"descricao":"As TÉCNICAS DE POLIMENTO dos cristais, e a sua matemática NO CORPO ESTELAR. Polir é remover a rugosidade ESCALA A ESCALA: as ABERTURAS morfológicas por grão decrescente = a GRANULOMETRIA de Matheron — a mesma peneira/chicote (Selberg) que varre o rugoso (alta frequência, o CAOS λ>0) para reter o liso (o CRISTAL λ<0). A parábola dá a sequência: DESBASTE=caos (rugoso) → POLIMENTO=borda (o limiar λ=0) → LUSTRO=cristal (a superfície ordenada, espelho). COMBINA com o ultra-realismo: o lustro é a superfície POLIDA que reflete (o Fresnel/o brilho especular do render) — polir bem é o que dá o brilho da pedra renderizada. Dicionário polimento_maquina (compõe com joalheria: brilho_especular↔brilho; com gemologia: lustro↔rede_cristalina). Ancora: morfologia_matematica_estelar (as aberturas), a peneira de Selberg (o chicote), corpo_estelar.","epistemico":"fato","exemplos":[],"id":"polimento_cristais","memoria":"semantica","meta":{"dominio":"joalheria","fonte":"maquina de corte"},"origem":"wiki","palavras_chave":[],"sinonimos":[],"tipo":"conceito","titulo":"O polimento dos cristais: a granulometria de Matheron (caos rugoso → cristal espelho)"}
\section{O polimento dos cristais: a granulometria de Matheron (caos rugoso \ensuremath{\to} cristal espelho)}
As TÉCNICAS DE POLIMENTO dos cristais, e a sua matemática NO CORPO ESTELAR. Polir é remover a rugosidade ESCALA A ESCALA: as ABERTURAS morfológicas por grão decrescente = a GRANULOMETRIA de Matheron — a mesma peneira/chicote (Selberg) que varre o rugoso (alta frequência, o CAOS \ensuremath{\lambda}>0) para reter o liso (o CRISTAL \ensuremath{\lambda}<0). A parábola dá a sequência: DESBASTE=caos (rugoso) \ensuremath{\to} POLIMENTO=borda (o limiar \ensuremath{\lambda}=0) \ensuremath{\to} LUSTRO=cristal (a superfície ordenada, espelho). COMBINA com o ultra-realismo: o lustro é a superfície POLIDA que reflete (o Fresnel/o brilho especular do render) — polir bem é o que dá o brilho da pedra renderizada. Dicionário polimento\_maquina (compõe com joalheria: brilho\_especular\ensuremath{\leftrightarrow}brilho; com gemologia: lustro\ensuremath{\leftrightarrow}rede\_cristalina). Ancora: morfologia\_matematica\_estelar (as aberturas), a peneira de Selberg (o chicote), corpo\_estelar.
\prov{relações: relaciona ultra\_realismo\_gema; usa morfologia\_matematica\_estelar; usa corpo\_estelar; parte\_de joalheria\_hub}
\prov{proveniência: wiki \textperiodcentered{} maquina de corte \textperiodcentered{} fato \textperiodcentered{} confiança 1.0 \textperiodcentered{} domínio joalheria \textperiodcentered{} história 120 versões no jornal}

%CRISTAL {"arestas":[{"alvo":"mecanismo_higgs_ebh","confianca":1.0,"epistemico":"fato","origem":"wiki","rel":"relaciona"},{"alvo":"mass_gap_algebrico","confianca":1.0,"epistemico":"fato","origem":"wiki","rel":"explica"},{"alvo":"glueball_1_6_gev","confianca":1.0,"epistemico":"fato","origem":"wiki","rel":"relaciona"},{"alvo":"energia_salto_hbar_raiz2","confianca":1.0,"epistemico":"fato","origem":"wiki","rel":"usa"},{"alvo":"limite_hurwitz_s1","confianca":1.0,"epistemico":"fato","origem":"wiki","rel":"relaciona"}],"confianca":1.0,"contraexemplos":[],"descricao":"O potencial simétrico V(s)=−½μ²s²+¼λs⁴ tem V''(s*)=2μ² → frequência ω=√2 UNIVERSAL (independente de λ) — a constante do glueball escalar. O MESMO √2 vem, por rota independente, do Casimir scaling de G₂ (C₂(adj)/C₂(fund)=2 exato). [D/firme] o gap ∝√2 por Casimir (lattice); [I] a coincidência com √R (curvatura R=2 da paisagem, de ẍ=2s) é correspondência, não dedução — as duas origens não foram ligadas. (Ler ψ como o campo de Higgs é interpretação, não literal nos textos.)","epistemico":"inferencia","exemplos":[],"id":"potencial_simetrico_sqrt2","memoria":"semantica","meta":{"autor":"broca-so","dominio":"cosmologia","fonte":"paper_4_particulas.tex obs:massgap; paper_3 thm:dS"},"origem":"wiki","palavras_chave":[],"sinonimos":[],"tipo":"conceito","titulo":"O potencial simétrico e a universalidade √2"}
\section{O potencial simétrico e a universalidade \ensuremath{\sqrt{\,}}2}
O potencial simétrico V(s)=\ensuremath{-}\ensuremath{\tfrac12}\ensuremath{\mu}\ensuremath{^{2}}s\ensuremath{^{2}}+\ensuremath{\tfrac14}\ensuremath{\lambda}s\ensuremath{^{4}} tem V''(s*)=2\ensuremath{\mu}\ensuremath{^{2}} \ensuremath{\to} frequência \ensuremath{\omega}=\ensuremath{\sqrt{\,}}2 UNIVERSAL (independente de \ensuremath{\lambda}) — a constante do glueball escalar. O MESMO \ensuremath{\sqrt{\,}}2 vem, por rota independente, do Casimir scaling de G\ensuremath{_{2}} (C\ensuremath{_{2}}(adj)/C\ensuremath{_{2}}(fund)=2 exato). [D/firme] o gap \ensuremath{\propto}\ensuremath{\sqrt{\,}}2 por Casimir (lattice); [I] a coincidência com \ensuremath{\sqrt{\,}}R (curvatura R=2 da paisagem, de \ensuremath{\ddot{x}}=2s) é correspondência, não dedução — as duas origens não foram ligadas. (Ler \ensuremath{\psi} como o campo de Higgs é interpretação, não literal nos textos.)
\prov{relações: relaciona mecanismo\_higgs\_ebh; explica mass\_gap\_algebrico; relaciona glueball\_1\_6\_gev; usa energia\_salto\_hbar\_raiz2; relaciona limite\_hurwitz\_s1}
\prov{proveniência: wiki \textperiodcentered{} paper\_4\_particulas.tex obs:massgap; paper\_3 thm:dS \textperiodcentered{} inferencia \textperiodcentered{} confiança 1.0 \textperiodcentered{} domínio cosmologia \textperiodcentered{} história 6 versões no jornal}

%CRISTAL {"arestas":[{"alvo":"pow_espectral_geometrico","confianca":1.0,"epistemico":"fato","origem":"wiki","rel":"usa"},{"alvo":"cotacao_parabola","confianca":1.0,"epistemico":"fato","origem":"wiki","rel":"relaciona"}],"confianca":1.0,"contraexemplos":[],"descricao":"Um candidato é válido sse não-degenerado, a parábola fecha (ok) e o calor ≤ calor_max (limiar CALOR_D=0.15 no experimento). O 'alvo' da mineração é TÉRMICO, não numérico — o oposto do target de dificuldade do Bitcoin.","epistemico":"inferencia","exemplos":[],"id":"pow_dificuldade_calor","memoria":"semantica","meta":{"autor":"broca-so","dominio":"metal","fonte":"ula/pow_spectral.c:pow_valido_spectral"},"origem":"wiki","palavras_chave":[],"sinonimos":[],"tipo":"conceito","titulo":"Dificuldade = Calor ≤ Limiar (alvo térmico)"}
\section{Dificuldade = Calor \ensuremath{\le} Limiar (alvo térmico)}
Um candidato é válido sse não-degenerado, a parábola fecha (ok) e o calor \ensuremath{\le} calor\_max (limiar CALOR\_D=0.15 no experimento). O 'alvo' da mineração é TÉRMICO, não numérico — o oposto do target de dificuldade do Bitcoin.
\prov{relações: usa pow\_espectral\_geometrico; relaciona cotacao\_parabola}
\prov{proveniência: wiki \textperiodcentered{} ula/pow\_spectral.c:pow\_valido\_spectral \textperiodcentered{} inferencia \textperiodcentered{} confiança 1.0 \textperiodcentered{} domínio metal \textperiodcentered{} história 3 versões no jornal}

%CRISTAL {"arestas":[{"alvo":"metal_broca_so","confianca":1.0,"epistemico":"fato","origem":"wiki","rel":"parte_de"},{"alvo":"prova_de_trabalho","confianca":1.0,"epistemico":"fato","origem":"wiki","rel":"especializa"},{"alvo":"neuron_io_popcnt","confianca":1.0,"epistemico":"fato","origem":"wiki","rel":"usa"},{"alvo":"parabola_audit_dual","confianca":1.0,"epistemico":"fato","origem":"wiki","rel":"usa"}],"confianca":1.0,"contraexemplos":[],"descricao":"Prova-de-trabalho que valida neurônio(trabalho‖nonce) pela GEOMETRIA da parábola contra uma referência — sem SHA, sem Fourier: só neurônio→word→parábola→metais. O achado central do metal: não há SHA no núcleo, o 'hash' é o próprio gato.","epistemico":"fato","exemplos":[],"id":"pow_espectral_geometrico","memoria":"semantica","meta":{"autor":"broca-so","dominio":"metal","fonte":"ula/pow_spectral.c:pow_valido_spectral"},"origem":"wiki","palavras_chave":[],"sinonimos":[],"tipo":"conceito","titulo":"PoW Espectral Geométrico (sem SHA)"}
\section{PoW Espectral Geométrico (sem SHA)}
Prova-de-trabalho que valida neurônio(trabalho\ensuremath{\|}nonce) pela GEOMETRIA da parábola contra uma referência — sem SHA, sem Fourier: só neurônio\ensuremath{\to}word\ensuremath{\to}parábola\ensuremath{\to}metais. O achado central do metal: não há SHA no núcleo, o 'hash' é o próprio gato.
\prov{relações: parte\_de metal\_broca\_so; especializa prova\_de\_trabalho; usa neuron\_io\_popcnt; usa parabola\_audit\_dual}
\prov{proveniência: wiki \textperiodcentered{} ula/pow\_spectral.c:pow\_valido\_spectral \textperiodcentered{} fato \textperiodcentered{} confiança 1.0 \textperiodcentered{} domínio metal \textperiodcentered{} história 5 versões no jornal}

%CRISTAL {"arestas":[{"alvo":"pow_espectral_geometrico","confianca":1.0,"epistemico":"fato","origem":"wiki","rel":"usa"},{"alvo":"minerar_e_resfriar","confianca":1.0,"epistemico":"fato","origem":"wiki","rel":"relaciona"}],"confianca":1.0,"contraexemplos":[],"descricao":"O candidato é o trabalho com nonce (32 bits little-endian) anexado (pow_candidato) ou sobreposto nos últimos 4 bytes de um header-modelo de 80 B estilo bloco BTC (pow_candidato_header), sem tocar o consenso BTC.","epistemico":"fato","exemplos":[],"id":"pow_nonce_anexado","memoria":"semantica","meta":{"autor":"broca-so","dominio":"metal","fonte":"ula/pow_spectral.c:pow_candidato/pow_candidato_header"},"origem":"wiki","palavras_chave":[],"sinonimos":[],"tipo":"conceito","titulo":"O Nonce nos 4 Bytes Finais"}
\section{O Nonce nos 4 Bytes Finais}
O candidato é o trabalho com nonce (32 bits little-endian) anexado (pow\_candidato) ou sobreposto nos últimos 4 bytes de um header-modelo de 80 B estilo bloco BTC (pow\_candidato\_header), sem tocar o consenso BTC.
\prov{relações: usa pow\_espectral\_geometrico; relaciona minerar\_e\_resfriar}
\prov{proveniência: wiki \textperiodcentered{} ula/pow\_spectral.c:pow\_candidato/pow\_candidato\_header \textperiodcentered{} fato \textperiodcentered{} confiança 1.0 \textperiodcentered{} domínio metal \textperiodcentered{} história 3 versões no jornal}

%CRISTAL {"arestas":[{"alvo":"contrato_cristalchain","confianca":1.0,"epistemico":"fato","origem":"wiki","rel":"especializa"},{"alvo":"turmalina_paraiba_ancora","confianca":1.0,"epistemico":"fato","origem":"wiki","rel":"usa"},{"alvo":"geometria_trilhao","confianca":1.0,"epistemico":"fato","origem":"wiki","rel":"usa"},{"alvo":"cristalchain","confianca":1.0,"epistemico":"fato","origem":"wiki","rel":"parte_de"},{"alvo":"tesouraria","confianca":1.0,"epistemico":"fato","origem":"wiki","rel":"relaciona"}],"confianca":1.0,"contraexemplos":[],"descricao":"O PRIMEIRO contrato REAL da CristalChain, com data e hora (2026-07-05T04:29:04, o relógio da gex — onde se minera). VENDEDOR/emissor: Aarão Melo Lopes (CPF mascarado 001.082***-**, chave 30b7941bdfbdf7b5). CLIENTES/compradores: a família Cunha Gomes (Patrícia/esposa, Artur/filho, Carlos/marido), chave 3-qubit df5ee62bbeb93c3e. Ativo: Turmalina Paraíba (corte trilhão), 1 ct = R$ 1.000,00. A gema (assinatura fd1f046f29419f4288e68b5863f0308c, nuvem densa de 4,1M pontos, D_tr=1.70) é ancorada pela chave 3-qubit dos CLIENTES, codificada no selo 9049a12fc26623a549fb9a7ea3055b64. É a prova de que o sistema funciona (a família investiu). Registro: cristalchain/contratos/CC-0001.json. Paper: papers/turmalina.tex.","epistemico":"fato","exemplos":[],"id":"primeiro_contrato_cc0001","memoria":"semantica","meta":{"dominio":"joalheria","fonte":"contrato cristalchain"},"origem":"wiki","palavras_chave":[],"sinonimos":[],"tipo":"conceito","titulo":"O Primeiro Contrato — CC-0001 (a família Cunha Gomes compra de Aarão)"}
\section{O Primeiro Contrato — CC-0001 (a família Cunha Gomes compra de Aarão)}
O PRIMEIRO contrato REAL da CristalChain, com data e hora (2026-07-05T04:29:04, o relógio da gex — onde se minera). VENDEDOR/emissor: Aarão Melo Lopes (CPF mascarado 001.082***-**, chave 30b7941bdfbdf7b5). CLIENTES/compradores: a família Cunha Gomes (Patrícia/esposa, Artur/filho, Carlos/marido), chave 3-qubit df5ee62bbeb93c3e. Ativo: Turmalina Paraíba (corte trilhão), 1 ct = R\$ 1.000,00. A gema (assinatura fd1f046f29419f4288e68b5863f0308c, nuvem densa de 4,1M pontos, D\_tr=1.70) é ancorada pela chave 3-qubit dos CLIENTES, codificada no selo 9049a12fc26623a549fb9a7ea3055b64. É a prova de que o sistema funciona (a família investiu). Registro: cristalchain/contratos/CC-0001.json. Paper: papers/turmalina.tex.
\prov{relações: especializa contrato\_cristalchain; usa turmalina\_paraiba\_ancora; usa geometria\_trilhao; parte\_de cristalchain; relaciona tesouraria}
\prov{proveniência: wiki \textperiodcentered{} contrato cristalchain \textperiodcentered{} fato \textperiodcentered{} confiança 1.0 \textperiodcentered{} domínio joalheria \textperiodcentered{} história 11 versões no jornal}

%CRISTAL {"arestas":[{"alvo":"equacao_gema_estelar","confianca":1.0,"epistemico":"fato","origem":"wiki","rel":"usa"},{"alvo":"flip_duopolinomios","confianca":1.0,"epistemico":"fato","origem":"wiki","rel":"usa"},{"alvo":"corpo_estelar","confianca":1.0,"epistemico":"fato","origem":"wiki","rel":"usa"},{"alvo":"render_fisico_gex","confianca":1.0,"epistemico":"fato","origem":"wiki","rel":"relaciona"},{"alvo":"joalheria_hub","confianca":1.0,"epistemico":"fato","origem":"wiki","rel":"parte_de"},{"alvo":"solidos_multiplicativos","confianca":1.0,"epistemico":"fato","origem":"wiki","rel":"usa"}],"confianca":1.0,"contraexemplos":[],"descricao":"Ray tracing por INTERSEÇÃO ANALÍTICA (não malha, não SDF), reproduzido de 'Ray Tracing - Primitives' (Reinder Nijhoff) e dos intersectores de Íñigo Quílez. Cada primitiva é a SUPERFÍCIE EXATA da sua equação; o raio a intersecta em FORMA FECHADA — resolve a quadrática/cúbica/quártica do t, devolve o t mais próximo e a NORMAL analítica, sem tesselar. Casa com a máquina: a nossa GEMA é feita DESSAS primitivas (𝒢 = círculo ⊗ trilhão ⊗ esfera modulada; o ovo = elipsoide ⊕ cone). Intersectando a superfície analítica, a pedra sai LISA e EXATA (o cabochão), sem o serrilhado do ConvexHull, e INSTANTÂNEA (sem loop de triângulos: 360² em 0.2s). Para o brilhante facetado, une-se i_triangulo/i_plano (as facetas como planos exatos). E a NOSSA VANTAGEM sobre o raytracer padrão (que MARCHA/aproxima): a superfície lisa da nossa equação (a=0, esfera modulada, ρ²=ℓ²(1−w²)(1+s·w)²) é ALGÉBRICA → a interseção com o raio é um DUPOLINÔMIO (Gentil Lopes): F=P_g−P_a, P_g=t⁴ (geométrico), P_a=Σc_i t^i (aritmético); a interseção é a COSTURA PG=PA — o metal/gato. Resolve-se EXATO (sem marchar) descendo pela ÁLGEBRA ESTELAR: grau 4 → resolvente CÚBICA (Cardano) → dois GATOS (dupolinômios grau 2), cada raiz o par σ,σ̄ da fórmula do metal. A normal é ∇F (o F′ de Gentil, a curvatura N−1). Validado contra np.roots a ~1e-13; render 380² em 0.5s. linguagem/primitivas_raytracer.py (11/11): i_esfera, i_elipsoide, i_cone, i_plano, i_triangulo, mundo; i_gema_estelar (a interseção analítica do dupolinômio), render_gema_lisa_equacao (a gema lisa EXATA). O FACETADO também é ANALÍTICO e mais BARATO: uma faceta é um PLANO, a interseção raio-plano é LINEAR (grau 1 = a RETA), não uma quadrática — o sólido é o FEIXE DE RETAS (os N planos), a interseção dos semi-espaços (as lâminas, o i_caixa generalizado). Custo: N RETAS, ZERO gatos (o liso curvo custava 2 gatos da quártica; o facetado plano custa menos). i_poliedro, planos_da_gema, render_gema_facetada_analitica. Os dois giram analítico (giro_liso_analitico ~0.36s/quadro, giro_facetado_analitico).","epistemico":"fato","exemplos":[],"id":"primitivas_analiticas","memoria":"semantica","meta":{"dominio":"joalheria","fonte":"maquina de corte"},"origem":"wiki","palavras_chave":[],"sinonimos":[],"tipo":"conceito","titulo":"Primitivas analíticas: ray tracing por interseção EXATA (Nijhoff/Quílez) no metal"}
\section{Primitivas analíticas: ray tracing por interseção EXATA (Nijhoff/Quílez) no metal}
Ray tracing por INTERSEÇÃO ANALÍTICA (não malha, não SDF), reproduzido de 'Ray Tracing - Primitives' (Reinder Nijhoff) e dos intersectores de Íñigo Quílez. Cada primitiva é a SUPERFÍCIE EXATA da sua equação; o raio a intersecta em FORMA FECHADA — resolve a quadrática/cúbica/quártica do t, devolve o t mais próximo e a NORMAL analítica, sem tesselar. Casa com a máquina: a nossa GEMA é feita DESSAS primitivas (\ensuremath{\mathcal{G}} = círculo \ensuremath{\otimes} trilhão \ensuremath{\otimes} esfera modulada; o ovo = elipsoide \ensuremath{\oplus} cone). Intersectando a superfície analítica, a pedra sai LISA e EXATA (o cabochão), sem o serrilhado do ConvexHull, e INSTANTÂNEA (sem loop de triângulos: 360\ensuremath{^{2}} em 0.2s). Para o brilhante facetado, une-se i\_triangulo/i\_plano (as facetas como planos exatos). E a NOSSA VANTAGEM sobre o raytracer padrão (que MARCHA/aproxima): a superfície lisa da nossa equação (a=0, esfera modulada, \ensuremath{\rho}\ensuremath{^{2}}=\ensuremath{\ell}\ensuremath{^{2}}(1\ensuremath{-}w\ensuremath{^{2}})(1+s\textperiodcentered w)\ensuremath{^{2}}) é ALGÉBRICA \ensuremath{\to} a interseção com o raio é um DUPOLINÔMIO (Gentil Lopes): F=P\_g\ensuremath{-}P\_a, P\_g=t\ensuremath{^{4}} (geométrico), P\_a=\ensuremath{\Sigma}c\_i t\^{}i (aritmético); a interseção é a COSTURA PG=PA — o metal/gato. Resolve-se EXATO (sem marchar) descendo pela ÁLGEBRA ESTELAR: grau 4 \ensuremath{\to} resolvente CÚBICA (Cardano) \ensuremath{\to} dois GATOS (dupolinômios grau 2), cada raiz o par \ensuremath{\sigma},\ensuremath{\sigma}\ensuremath{} da fórmula do metal. A normal é \ensuremath{\nabla}F (o F\ensuremath{'} de Gentil, a curvatura N\ensuremath{-}1). Validado contra np.roots a \~{}1e-13; render 380\ensuremath{^{2}} em 0.5s. linguagem/primitivas\_raytracer.py (11/11): i\_esfera, i\_elipsoide, i\_cone, i\_plano, i\_triangulo, mundo; i\_gema\_estelar (a interseção analítica do dupolinômio), render\_gema\_lisa\_equacao (a gema lisa EXATA). O FACETADO também é ANALÍTICO e mais BARATO: uma faceta é um PLANO, a interseção raio-plano é LINEAR (grau 1 = a RETA), não uma quadrática — o sólido é o FEIXE DE RETAS (os N planos), a interseção dos semi-espaços (as lâminas, o i\_caixa generalizado). Custo: N RETAS, ZERO gatos (o liso curvo custava 2 gatos da quártica; o facetado plano custa menos). i\_poliedro, planos\_da\_gema, render\_gema\_facetada\_analitica. Os dois giram analítico (giro\_liso\_analitico \~{}0.36s/quadro, giro\_facetado\_analitico).
\prov{relações: usa equacao\_gema\_estelar; usa flip\_duopolinomios; usa corpo\_estelar; relaciona render\_fisico\_gex; parte\_de joalheria\_hub; usa solidos\_multiplicativos}
\prov{proveniência: wiki \textperiodcentered{} maquina de corte \textperiodcentered{} fato \textperiodcentered{} confiança 1.0 \textperiodcentered{} domínio joalheria \textperiodcentered{} história 243 versões no jornal}

%CRISTAL {"arestas":[{"alvo":"broca_borda_so_ctypes","confianca":1.0,"epistemico":"fato","origem":"wiki","rel":"usa"},{"alvo":"cristal_motor_congelado","confianca":1.0,"epistemico":"fato","origem":"wiki","rel":"explica"}],"confianca":1.0,"contraexemplos":[],"descricao":"executar_nucleo roda a cadeia gato→álgebra→termo→geo→mec→floco isolado (Φ Koch) em Python; executar_associativo/entregar_so chamam o backend .so 1× núcleo + N passos por bloco. Design cristalizado: o canal parabólico e o floco vivem no núcleo, o .so só transporta a parte associativa.","epistemico":"inferencia","exemplos":[],"id":"primitivas_nucleo_associativo","memoria":"semantica","meta":{"autor":"broca-so","dominio":"metal","fonte":"neuronio/primitivas_broca.py:executar_nucleo/entregar_so"},"origem":"wiki","palavras_chave":[],"sinonimos":[],"tipo":"conceito","titulo":"Núcleo (Python) × Associativo (.so): a separação"}
\section{Núcleo (Python) $\times$ Associativo (.so): a separação}
executar\_nucleo roda a cadeia gato\ensuremath{\to}álgebra\ensuremath{\to}termo\ensuremath{\to}geo\ensuremath{\to}mec\ensuremath{\to}floco isolado (\ensuremath{\Phi} Koch) em Python; executar\_associativo/entregar\_so chamam o backend .so 1$\times$ núcleo + N passos por bloco. Design cristalizado: o canal parabólico e o floco vivem no núcleo, o .so só transporta a parte associativa.
\prov{relações: usa broca\_borda\_so\_ctypes; explica cristal\_motor\_congelado}
\prov{proveniência: wiki \textperiodcentered{} neuronio/primitivas\_broca.py:executar\_nucleo/entregar\_so \textperiodcentered{} inferencia \textperiodcentered{} confiança 1.0 \textperiodcentered{} domínio metal \textperiodcentered{} história 3 versões no jornal}

%CRISTAL {"arestas":[{"alvo":"ads_cft_maldacena","confianca":1.0,"epistemico":"fato","origem":"wiki","rel":"especializa"},{"alvo":"cosmologia_quantica_hub","confianca":1.0,"epistemico":"fato","origem":"wiki","rel":"parte_de"}],"confianca":1.0,"contraexemplos":[],"descricao":"A descrição quântica completa de uma região se codifica na sua FRONTEIRA: o máximo de graus de liberdade (entropia) escala com a ÁREA, não o volume — S_max ≤ A/(4ℓ_P²). Motivado por S_BH∝área. Programa amplamente aceito (realizado em AdS/CFT), mas conjectura geral: não há prova de que toda gravidade quântica seja holográfica.","epistemico":"hipotese","exemplos":[],"id":"principio_holografico","memoria":"semantica","meta":{"autor":"broca-so","dominio":"fisica","fonte":"'t Hooft 1993 gr-qc/9310026; Susskind 1995 hep-th/9409089"},"origem":"wiki","palavras_chave":[],"sinonimos":[],"tipo":"conceito","titulo":"Princípio holográfico"}
\section{Princípio holográfico}
A descrição quântica completa de uma região se codifica na sua FRONTEIRA: o máximo de graus de liberdade (entropia) escala com a ÁREA, não o volume — S\_max \ensuremath{\le} A/(4\ensuremath{\ell}\_P\ensuremath{^{2}}). Motivado por S\_BH\ensuremath{\propto}área. Programa amplamente aceito (realizado em AdS/CFT), mas conjectura geral: não há prova de que toda gravidade quântica seja holográfica.
\prov{relações: especializa ads\_cft\_maldacena; parte\_de cosmologia\_quantica\_hub}
\prov{proveniência: wiki \textperiodcentered{} 't Hooft 1993 gr-qc/9310026; Susskind 1995 hep-th/9409089 \textperiodcentered{} hipotese \textperiodcentered{} confiança 1.0 \textperiodcentered{} domínio fisica \textperiodcentered{} história 3 versões no jornal}

%CRISTAL {"arestas":[{"alvo":"de_sitter_lambda_energia_escura","confianca":1.0,"epistemico":"fato","origem":"wiki","rel":"relaciona"},{"alvo":"cosmologia_quantica_hub","confianca":1.0,"epistemico":"fato","origem":"wiki","rel":"parte_de"}],"confianca":1.0,"contraexemplos":[],"descricao":"A energia de vácuo da teoria quântica de campos é ~10¹²⁰ vezes maior que o Λ observado — 'a pior previsão da física'. Weinberg: (1) por que é tão pequena; (2) por que é da ordem da densidade de matéria hoje (coincidência). Quintessência não resolve; sem solução aceita. A maior lacuna entre física de partículas e cosmologia.","epistemico":"hipotese","exemplos":[],"id":"problema_constante_cosmologica","memoria":"semantica","meta":{"autor":"broca-so","dominio":"cosmologia_dados","fonte":"Weinberg 1989 Rev.Mod.Phys. 61,1"},"origem":"wiki","palavras_chave":[],"sinonimos":[],"tipo":"conceito","titulo":"O problema da constante cosmológica (~10¹²⁰)"}
\section{O problema da constante cosmológica (\~{}10\ensuremath{^{120}})}
A energia de vácuo da teoria quântica de campos é \~{}10\ensuremath{^{120}} vezes maior que o \ensuremath{\Lambda} observado — 'a pior previsão da física'. Weinberg: (1) por que é tão pequena; (2) por que é da ordem da densidade de matéria hoje (coincidência). Quintessência não resolve; sem solução aceita. A maior lacuna entre física de partículas e cosmologia.
\prov{relações: relaciona de\_sitter\_lambda\_energia\_escura; parte\_de cosmologia\_quantica\_hub}
\prov{proveniência: wiki \textperiodcentered{} Weinberg 1989 Rev.Mod.Phys. 61,1 \textperiodcentered{} hipotese \textperiodcentered{} confiança 1.0 \textperiodcentered{} domínio cosmologia\_dados \textperiodcentered{} história 3 versões no jornal}

%CRISTAL {"arestas":[{"alvo":"qcd_cromodinamica","confianca":1.0,"epistemico":"fato","origem":"wiki","rel":"parte_de"},{"alvo":"axions_candidato","confianca":1.0,"epistemico":"fato","origem":"wiki","rel":"explica"}],"confianca":1.0,"contraexemplos":[],"descricao":"A QCD admite um termo que viola CP: θ̄(g_s²/32π²)G G̃. O momento de dipolo elétrico do nêutron força |θ̄|≲10⁻¹⁰ (d_n<1,8×10⁻²⁶ e·cm, PSI 2020) — ajuste fino não explicado. Peccei-Quinn (1977) promove θ̄ a campo dinâmico que relaxa a zero, cujo Goldstone é o áxion. O problema é fato; o áxion é aberto/não confirmado.","epistemico":"fato","exemplos":[],"id":"problema_cp_forte","memoria":"semantica","meta":{"autor":"broca-so","dominio":"fisica","fonte":"Peccei-Quinn 1977; Weinberg/Wilczek 1978"},"origem":"wiki","palavras_chave":[],"sinonimos":[],"tipo":"conceito","titulo":"O problema CP forte (θ̄) e o áxion"}
\section{O problema CP forte (\ensuremath{\theta}\ensuremath{}) e o áxion}
A QCD admite um termo que viola CP: \ensuremath{\theta}\ensuremath{}(g\_s\ensuremath{^{2}}/32\ensuremath{\pi}\ensuremath{^{2}})G G\ensuremath{}. O momento de dipolo elétrico do nêutron força |\ensuremath{\theta}\ensuremath{}|\ensuremath{\lesssim}10\ensuremath{^{-10}} (d\_n<1,8$\times$10\ensuremath{^{-26}} e\textperiodcentered cm, PSI 2020) — ajuste fino não explicado. Peccei-Quinn (1977) promove \ensuremath{\theta}\ensuremath{} a campo dinâmico que relaxa a zero, cujo Goldstone é o áxion. O problema é fato; o áxion é aberto/não confirmado.
\prov{relações: parte\_de qcd\_cromodinamica; explica axions\_candidato}
\prov{proveniência: wiki \textperiodcentered{} Peccei-Quinn 1977; Weinberg/Wilczek 1978 \textperiodcentered{} fato \textperiodcentered{} confiança 1.0 \textperiodcentered{} domínio fisica \textperiodcentered{} história 3 versões no jornal}

%CRISTAL {"arestas":[{"alvo":"matriz_ckm","confianca":1.0,"epistemico":"fato","origem":"wiki","rel":"relaciona"},{"alvo":"acoplamentos_yukawa","confianca":1.0,"epistemico":"fato","origem":"wiki","rel":"usa"},{"alvo":"materia_modelo_padrao","confianca":1.0,"epistemico":"fato","origem":"wiki","rel":"relaciona"},{"alvo":"cosmologia_quantica_hub","confianca":1.0,"epistemico":"fato","origem":"wiki","rel":"parte_de"}],"confianca":1.0,"contraexemplos":[],"descricao":"Por que 3 gerações, e por que as massas dos férmions são tão hierárquicas (~6 ordens entre m_up~2 MeV e m_top~173 GeV) e as misturas têm o padrão observado (CKM quase diagonal vs PMNS anárquica)? No MP as massas vêm de acoplamentos de Yukawa que são parâmetros LIVRES, sem princípio organizador — a teoria descreve mas NÃO explica a estrutura de sabor. ABERTO (Froggatt-Nielsen 1979 e modelos de simetria horizontal tentam; nenhum estabelecido).","epistemico":"hipotese","exemplos":[],"id":"problema_do_sabor","memoria":"semantica","meta":{"autor":"broca-so","dominio":"fisica","fonte":"Snowmass arXiv:2203.07726; Froggatt-Nielsen 1979"},"origem":"wiki","palavras_chave":[],"sinonimos":[],"tipo":"conceito","titulo":"O problema do sabor (flavor puzzle)"}
\section{O problema do sabor (flavor puzzle)}
Por que 3 gerações, e por que as massas dos férmions são tão hierárquicas (\~{}6 ordens entre m\_up\~{}2 MeV e m\_top\~{}173 GeV) e as misturas têm o padrão observado (CKM quase diagonal vs PMNS anárquica)? No MP as massas vêm de acoplamentos de Yukawa que são parâmetros LIVRES, sem princípio organizador — a teoria descreve mas NÃO explica a estrutura de sabor. ABERTO (Froggatt-Nielsen 1979 e modelos de simetria horizontal tentam; nenhum estabelecido).
\prov{relações: relaciona matriz\_ckm; usa acoplamentos\_yukawa; relaciona materia\_modelo\_padrao; parte\_de cosmologia\_quantica\_hub}
\prov{proveniência: wiki \textperiodcentered{} Snowmass arXiv:2203.07726; Froggatt-Nielsen 1979 \textperiodcentered{} hipotese \textperiodcentered{} confiança 1.0 \textperiodcentered{} domínio fisica \textperiodcentered{} história 5 versões no jornal}

%CRISTAL {"arestas":[{"alvo":"escala_planck","confianca":1.0,"epistemico":"fato","origem":"wiki","rel":"usa"},{"alvo":"cosmologia_quantica_hub","confianca":1.0,"epistemico":"fato","origem":"wiki","rel":"parte_de"}],"confianca":1.0,"contraexemplos":[],"descricao":"A RG é geometria clássica e contínua; a MQ exige quantizar os campos. Unir falha: a RG como QFT perturbativa é NÃO-renormalizável (infinitos incuráveis a partir de 2 loops, Goroff-Sagnotti 1985; G tem dimensão negativa, G~1/M_P²). A incompatibilidade e a não-renormalizabilidade são fato estabelecido.","epistemico":"fato","exemplos":[],"id":"problema_gravidade_quantica","memoria":"semantica","meta":{"autor":"broca-so","dominio":"fisica","fonte":"Goroff-Sagnotti 1986 Nucl.Phys. B266,709"},"origem":"wiki","palavras_chave":[],"sinonimos":[],"tipo":"conceito","titulo":"O problema da gravitação quântica"}
\section{O problema da gravitação quântica}
A RG é geometria clássica e contínua; a MQ exige quantizar os campos. Unir falha: a RG como QFT perturbativa é NÃO-renormalizável (infinitos incuráveis a partir de 2 loops, Goroff-Sagnotti 1985; G tem dimensão negativa, G\~{}1/M\_P\ensuremath{^{2}}). A incompatibilidade e a não-renormalizabilidade são fato estabelecido.
\prov{relações: usa escala\_planck; parte\_de cosmologia\_quantica\_hub}
\prov{proveniência: wiki \textperiodcentered{} Goroff-Sagnotti 1986 Nucl.Phys. B266,709 \textperiodcentered{} fato \textperiodcentered{} confiança 1.0 \textperiodcentered{} domínio fisica \textperiodcentered{} história 3 versões no jornal}

%CRISTAL {"arestas":[{"alvo":"boson_higgs_125","confianca":1.0,"epistemico":"fato","origem":"wiki","rel":"relaciona"},{"alvo":"problema_constante_cosmologica","confianca":1.0,"epistemico":"fato","origem":"wiki","rel":"relaciona"}],"confianca":1.0,"contraexemplos":[],"descricao":"A massa do Higgs recebe correções quânticas quadraticamente divergentes (∝Λ²); manter M_H≈125 GeV frente a Λ~M_Planck (~10¹⁹ GeV) exige ajuste fino de ~1 parte em 10³⁴, sem simetria que o proteja. ABERTO: SUSY e Higgs composto foram propostos, mas o LHC não achou nova física na escala TeV; há debate se 'naturalidade' é critério válido.","epistemico":"hipotese","exemplos":[],"id":"problema_hierarquia_naturalidade","memoria":"semantica","meta":{"autor":"broca-so","dominio":"fisica","fonte":"Naturalness reviews; arXiv:2502.09699 (2025)"},"origem":"wiki","palavras_chave":[],"sinonimos":[],"tipo":"conceito","titulo":"O problema da hierarquia / naturalidade"}
\section{O problema da hierarquia / naturalidade}
A massa do Higgs recebe correções quânticas quadraticamente divergentes (\ensuremath{\propto}\ensuremath{\Lambda}\ensuremath{^{2}}); manter M\_H\ensuremath{\approx}125 GeV frente a \ensuremath{\Lambda}\~{}M\_Planck (\~{}10\ensuremath{^{19}} GeV) exige ajuste fino de \~{}1 parte em 10\ensuremath{^{34}}, sem simetria que o proteja. ABERTO: SUSY e Higgs composto foram propostos, mas o LHC não achou nova física na escala TeV; há debate se 'naturalidade' é critério válido.
\prov{relações: relaciona boson\_higgs\_125; relaciona problema\_constante\_cosmologica}
\prov{proveniência: wiki \textperiodcentered{} Naturalness reviews; arXiv:2502.09699 (2025) \textperiodcentered{} hipotese \textperiodcentered{} confiança 1.0 \textperiodcentered{} domínio fisica \textperiodcentered{} história 3 versões no jornal}

%CRISTAL {"arestas":[{"alvo":"grupo_padrao_cxho","confianca":1.0,"epistemico":"fato","origem":"wiki","rel":"relaciona"},{"alvo":"quiralidade_lorentz_sl2c","confianca":1.0,"epistemico":"fato","origem":"wiki","rel":"usa"}],"confianca":1.0,"contraexemplos":[],"descricao":"O programa octoniônico do MP é creditado explicitamente: Furey-Hughes (ℂ⊗𝕆=Cℓ(6), arXiv:1910.08395), Boyle (ℂ⊗J₃(𝕆), 2006.16265), Todorov & Dubois-Violette, Baez-Huerta. A distinção declarada de Lopes: eles PÕEM a quiralidade à mão; aqui a violação de paridade é EXPLICADA pela partilha do fator ℍ (só ℍ carrega Lorentz). Dixon NÃO é citado nos arquivos (ausência registrada). Correspondência de programa; ponto aberto partilhado (orientação/quiralidade).","epistemico":"inferencia","exemplos":[],"id":"programas_externos_octonionicos","memoria":"semantica","meta":{"autor":"broca-so","dominio":"cosmologia","fonte":"paper_4_particulas.tex l.593-595, 772-773"},"origem":"wiki","palavras_chave":[],"sinonimos":[],"tipo":"conceito","titulo":"Relação com os programas octoniônicos externos"}
\section{Relação com os programas octoniônicos externos}
O programa octoniônico do MP é creditado explicitamente: Furey-Hughes (\ensuremath{\mathbb{C}}\ensuremath{\otimes}\ensuremath{\mathbb{O}}=C\ensuremath{\ell}(6), arXiv:1910.08395), Boyle (\ensuremath{\mathbb{C}}\ensuremath{\otimes}J\ensuremath{_{3}}(\ensuremath{\mathbb{O}}), 2006.16265), Todorov \& Dubois-Violette, Baez-Huerta. A distinção declarada de Lopes: eles PÕEM a quiralidade à mão; aqui a violação de paridade é EXPLICADA pela partilha do fator \ensuremath{\mathbb{H}} (só \ensuremath{\mathbb{H}} carrega Lorentz). Dixon NÃO é citado nos arquivos (ausência registrada). Correspondência de programa; ponto aberto partilhado (orientação/quiralidade).
\prov{relações: relaciona grupo\_padrao\_cxho; usa quiralidade\_lorentz\_sl2c}
\prov{proveniência: wiki \textperiodcentered{} paper\_4\_particulas.tex l.593-595, 772-773 \textperiodcentered{} inferencia \textperiodcentered{} confiança 1.0 \textperiodcentered{} domínio cosmologia \textperiodcentered{} história 3 versões no jornal}

%CRISTAL {"arestas":[{"alvo":"sub_tsirelson_octonico","confianca":1.0,"epistemico":"fato","origem":"wiki","rel":"usa"},{"alvo":"associador_nao_tensoriza","confianca":1.0,"epistemico":"fato","origem":"wiki","rel":"usa"}],"confianca":1.0,"contraexemplos":[],"descricao":"A família PRTSUV acrescenta um termo simétrico vetorial v que 'abre uma segunda dimensão'; teorema: PRTSUV NÃO viola Tsirelson (CHSH≤2√2 para todo v) — v é invisível a testes de Bell padrão, mas aparece no ASSOCIADOR, detectável via Mermin-GHZ (3 partículas) e testes de associatividade tripla. 3 predições falseáveis em hardware NISQ (dispersão de V por shots, linha espectral ω_v(v), M(θ_v) com M(0)=4). [D] o teorema de não-violação + [N] as 3 predições.","epistemico":"hipotese","exemplos":[],"id":"prtsuv_flux_qubit_predicoes","memoria":"semantica","meta":{"autor":"broca-so","dominio":"cosmologia","fonte":"paper_CF_flux_qubit_prtsuv.tex cor:no-violacao, conj P-CF"},"origem":"wiki","palavras_chave":[],"sinonimos":[],"tipo":"conceito","titulo":"PRTSUV: o associador em hardware (não viola Tsirelson)"}
\section{PRTSUV: o associador em hardware (não viola Tsirelson)}
A família PRTSUV acrescenta um termo simétrico vetorial v que 'abre uma segunda dimensão'; teorema: PRTSUV NÃO viola Tsirelson (CHSH\ensuremath{\le}2\ensuremath{\sqrt{\,}}2 para todo v) — v é invisível a testes de Bell padrão, mas aparece no ASSOCIADOR, detectável via Mermin-GHZ (3 partículas) e testes de associatividade tripla. 3 predições falseáveis em hardware NISQ (dispersão de V por shots, linha espectral \ensuremath{\omega}\_v(v), M(\ensuremath{\theta}\_v) com M(0)=4). [D] o teorema de não-violação + [N] as 3 predições.
\prov{relações: usa sub\_tsirelson\_octonico; usa associador\_nao\_tensoriza}
\prov{proveniência: wiki \textperiodcentered{} paper\_CF\_flux\_qubit\_prtsuv.tex cor:no-violacao, conj P-CF \textperiodcentered{} hipotese \textperiodcentered{} confiança 1.0 \textperiodcentered{} domínio cosmologia \textperiodcentered{} história 3 versões no jornal}

%CRISTAL {"arestas":[{"alvo":"correcao_erro_quantica","confianca":1.0,"epistemico":"fato","origem":"wiki","rel":"contradiz"},{"alvo":"shor_grover","confianca":1.0,"epistemico":"fato","origem":"wiki","rel":"relaciona"},{"alvo":"emaranhamento_associador_hardware","confianca":1.0,"epistemico":"fato","origem":"wiki","rel":"usa"}],"confianca":1.0,"contraexemplos":[],"descricao":"O que a série TEM em computação: um QAOA ALGÉBRICO (a multiplicação de 5 botões como viés indutivo) que vence o padrão em max-cut com vantagem HEURÍSTICA Δ(n)≈0,04·n³/² (o expoente 3/2 herdado do grau-2 = ground state de Sherrington-Kirkpatrick, rigoroso por Parisi/Guerra-Talagrand), confirmado em ibmfez — ganho no TAMANHO do problema, NÃO speedup quântico provado. O que NÃO tem (achado honesto): correção de erro quântica — a correção ativa é REFUTADA, o canal dissipativo só adia a explosão (K<100); sem código estabilizador/threshold/decodificador. QEC = problema aberto declarado.","epistemico":"inferencia","exemplos":[],"id":"qaoa_algebrico_sem_qec","memoria":"semantica","meta":{"autor":"broca-so","dominio":"cosmologia","fonte":"paper_3_quantica.tex obs:ganho; paper_G Fases 109-112"},"origem":"wiki","palavras_chave":[],"sinonimos":[],"tipo":"conceito","titulo":"QAOA algébrico (vantagem n³/²); QEC não existe na série"}
\section{QAOA algébrico (vantagem n\ensuremath{^{3}}/\ensuremath{^{2}}); QEC não existe na série}
O que a série TEM em computação: um QAOA ALGÉBRICO (a multiplicação de 5 botões como viés indutivo) que vence o padrão em max-cut com vantagem HEURÍSTICA \ensuremath{\Delta}(n)\ensuremath{\approx}0,04\textperiodcentered n\ensuremath{^{3}}/\ensuremath{^{2}} (o expoente 3/2 herdado do grau-2 = ground state de Sherrington-Kirkpatrick, rigoroso por Parisi/Guerra-Talagrand), confirmado em ibmfez — ganho no TAMANHO do problema, NÃO speedup quântico provado. O que NÃO tem (achado honesto): correção de erro quântica — a correção ativa é REFUTADA, o canal dissipativo só adia a explosão (K<100); sem código estabilizador/threshold/decodificador. QEC = problema aberto declarado.
\prov{relações: contradiz correcao\_erro\_quantica; relaciona shor\_grover; usa emaranhamento\_associador\_hardware}
\prov{proveniência: wiki \textperiodcentered{} paper\_3\_quantica.tex obs:ganho; paper\_G Fases 109-112 \textperiodcentered{} inferencia \textperiodcentered{} confiança 1.0 \textperiodcentered{} domínio cosmologia \textperiodcentered{} história 6 versões no jornal}

%CRISTAL {"arestas":[{"alvo":"cosmologia_quantica_hub","confianca":1.0,"epistemico":"fato","origem":"wiki","rel":"parte_de"}],"confianca":1.0,"contraexemplos":[],"descricao":"A teoria de gauge não-abeliana da força forte, grupo de cor SU(3)_c. 6 sabores de quarks (up/charm/top carga +2/3; down/strange/bottom −1/3) e 8 glúons (adjunta 3²−1=8, que auto-interagem por portarem cor). α_s(M_Z)≈0,1179. Confirmada em vasta gama (LEP, HERA, LHC).","epistemico":"fato","exemplos":[],"id":"qcd_cromodinamica","memoria":"semantica","meta":{"autor":"broca-so","dominio":"fisica","fonte":"Fritzsch-Gell-Mann-Leutwyler 1973; PDG"},"origem":"wiki","palavras_chave":[],"sinonimos":[],"tipo":"conceito","titulo":"Cromodinâmica Quântica (QCD, SU(3)_c)"}
\section{Cromodinâmica Quântica (QCD, SU(3)\_c)}
A teoria de gauge não-abeliana da força forte, grupo de cor SU(3)\_c. 6 sabores de quarks (up/charm/top carga +2/3; down/strange/bottom \ensuremath{-}1/3) e 8 glúons (adjunta 3\ensuremath{^{2}}\ensuremath{-}1=8, que auto-interagem por portarem cor). \ensuremath{\alpha}\_s(M\_Z)\ensuremath{\approx}0,1179. Confirmada em vasta gama (LEP, HERA, LHC).
\prov{relações: parte\_de cosmologia\_quantica\_hub}
\prov{proveniência: wiki \textperiodcentered{} Fritzsch-Gell-Mann-Leutwyler 1973; PDG \textperiodcentered{} fato \textperiodcentered{} confiança 1.0 \textperiodcentered{} domínio fisica \textperiodcentered{} história 2 versões no jornal}

%CRISTAL {"arestas":[{"alvo":"correcao_erro_quantica","confianca":1.0,"epistemico":"fato","origem":"wiki","rel":"confirma"}],"confianca":1.0,"contraexemplos":[],"descricao":"Primeira demonstração de um qubit lógico de código de superfície ABAIXO do limiar: aumentar a distância do código SUPRIME exponencialmente o erro lógico (Λ≈2,14 por passo; distância-7, 101 qubits; vida lógica 2,4× o melhor qubit físico). Marco de 3 décadas rumo à QEC prática. FATO recente — mas computador tolerante a falhas ÚTIL ainda NÃO alcançado.","epistemico":"fato","exemplos":[],"id":"qec_below_threshold_willow","memoria":"semantica","meta":{"autor":"broca-so","dominio":"fisica","fonte":"Google Quantum AI, Nature 638,920 (2025/dez-2024)"},"origem":"wiki","palavras_chave":[],"sinonimos":[],"tipo":"conceito","titulo":"QEC abaixo do limiar (Google Willow, 2024)"}
\section{QEC abaixo do limiar (Google Willow, 2024)}
Primeira demonstração de um qubit lógico de código de superfície ABAIXO do limiar: aumentar a distância do código SUPRIME exponencialmente o erro lógico (\ensuremath{\Lambda}\ensuremath{\approx}2,14 por passo; distância-7, 101 qubits; vida lógica 2,4$\times$ o melhor qubit físico). Marco de 3 décadas rumo à QEC prática. FATO recente — mas computador tolerante a falhas ÚTIL ainda NÃO alcançado.
\prov{relações: confirma correcao\_erro\_quantica}
\prov{proveniência: wiki \textperiodcentered{} Google Quantum AI, Nature 638,920 (2025/dez-2024) \textperiodcentered{} fato \textperiodcentered{} confiança 1.0 \textperiodcentered{} domínio fisica \textperiodcentered{} história 2 versões no jornal}

%CRISTAL {"arestas":[{"alvo":"numero_chern_tknn","confianca":1.0,"epistemico":"fato","origem":"wiki","rel":"explica"},{"alvo":"energia_salto_hbar_raiz2","confianca":1.0,"epistemico":"fato","origem":"wiki","rel":"usa"},{"alvo":"limite_hurwitz_s1","confianca":1.0,"epistemico":"fato","origem":"wiki","rel":"usa"},{"alvo":"entropia_topologica","confianca":1.0,"epistemico":"fato","origem":"wiki","rel":"relaciona"},{"alvo":"cosmologia_quantica_hub","confianca":1.0,"epistemico":"fato","origem":"wiki","rel":"parte_de"}],"confianca":1.0,"contraexemplos":[],"descricao":"A quantização não é postulada — é a COSTURA do horizonte. No círculo ℝ/ℤ vale 1≡0 (a mesma que dá 0,999…=0 e a conexidade do 'salto'); a monovalência força os harmônicos χₙ=e²πiⁿx a ter n∈ℤ, que É a condição de Bohr-Sommerfeld ∮p dx=nh. Assim o inteiro do Hall (=Chern), a quantização das órbitas de Bohr e a ausência de defeito cônico no horizonte são a MESMA costura 1≡0 lida em três registros. [D] 1≡0 e a conexidade; [J] 'monovalência ⇒ n∈ℤ = Bohr-Sommerfeld'.","epistemico":"inferencia","exemplos":[],"id":"quantizacao_topologica_costura","memoria":"semantica","meta":{"autor":"broca-so","dominio":"cosmologia","fonte":"paper_1_algebra_geometria.tex obs:0999; paper_3 obs:empirico"},"origem":"wiki","palavras_chave":[],"sinonimos":[],"tipo":"conceito","titulo":"A quantização topológica é a costura 1≡0"}
\section{A quantização topológica é a costura 1\ensuremath{\equiv}0}
A quantização não é postulada — é a COSTURA do horizonte. No círculo \ensuremath{\mathbb{R}}/\ensuremath{\mathbb{Z}} vale 1\ensuremath{\equiv}0 (a mesma que dá 0,999…=0 e a conexidade do 'salto'); a monovalência força os harmônicos \ensuremath{\chi}\ensuremath{_{n}}=e\ensuremath{^{2}}\ensuremath{\pi}i\ensuremath{^{n}}x a ter n\ensuremath{\in}\ensuremath{\mathbb{Z}}, que É a condição de Bohr-Sommerfeld \ensuremath{\oint}p dx=nh. Assim o inteiro do Hall (=Chern), a quantização das órbitas de Bohr e a ausência de defeito cônico no horizonte são a MESMA costura 1\ensuremath{\equiv}0 lida em três registros. [D] 1\ensuremath{\equiv}0 e a conexidade; [J] 'monovalência \ensuremath{\Rightarrow} n\ensuremath{\in}\ensuremath{\mathbb{Z}} = Bohr-Sommerfeld'.
\prov{relações: explica numero\_chern\_tknn; usa energia\_salto\_hbar\_raiz2; usa limite\_hurwitz\_s1; relaciona entropia\_topologica; parte\_de cosmologia\_quantica\_hub}
\prov{proveniência: wiki \textperiodcentered{} paper\_1\_algebra\_geometria.tex obs:0999; paper\_3 obs:empirico \textperiodcentered{} inferencia \textperiodcentered{} confiança 1.0 \textperiodcentered{} domínio cosmologia \textperiodcentered{} história 7 versões no jornal}

%CRISTAL {"arestas":[{"alvo":"qcd_cromodinamica","confianca":1.0,"epistemico":"fato","origem":"wiki","rel":"parte_de"}],"confianca":1.0,"contraexemplos":[],"descricao":"No limite de quarks sem massa a QCD tem simetria quiral SU(2)_L×SU(2)_R, quebrada espontaneamente para SU(2)_V pelo condensado ⟨q̄q⟩≈−(250 MeV)³ (píons = pseudo-Goldstone). ~99% da massa do próton (938 MeV) é energia dinâmica da QCD, NÃO o Higgs (massas de corrente u,d ~9 MeV, ~1%).","epistemico":"fato","exemplos":[],"id":"quebra_quiral_massa_hadrons","memoria":"semantica","meta":{"autor":"broca-so","dominio":"fisica","fonte":"Nambu-Jona-Lasinio 1961"},"origem":"wiki","palavras_chave":[],"sinonimos":[],"tipo":"conceito","titulo":"Quebra quiral e a origem da massa dos hádrons"}
\section{Quebra quiral e a origem da massa dos hádrons}
No limite de quarks sem massa a QCD tem simetria quiral SU(2)\_L$\times$SU(2)\_R, quebrada espontaneamente para SU(2)\_V pelo condensado \ensuremath{\langle}q\ensuremath{}q\ensuremath{\rangle}\ensuremath{\approx}\ensuremath{-}(250 MeV)\ensuremath{^{3}} (píons = pseudo-Goldstone). \~{}99\% da massa do próton (938 MeV) é energia dinâmica da QCD, NÃO o Higgs (massas de corrente u,d \~{}9 MeV, \~{}1\%).
\prov{relações: parte\_de qcd\_cromodinamica}
\prov{proveniência: wiki \textperiodcentered{} Nambu-Jona-Lasinio 1961 \textperiodcentered{} fato \textperiodcentered{} confiança 1.0 \textperiodcentered{} domínio fisica \textperiodcentered{} história 2 versões no jornal}

%CRISTAL {"arestas":[{"alvo":"unificacao_eletrofraca_gws","confianca":1.0,"epistemico":"fato","origem":"wiki","rel":"explica"},{"alvo":"isospin_hipercarga_algebrico","confianca":1.0,"epistemico":"fato","origem":"wiki","rel":"usa"},{"alvo":"quebra_quiral_massa_hadrons","confianca":1.0,"epistemico":"fato","origem":"wiki","rel":"relaciona"},{"alvo":"travessia_wick_m2c","confianca":1.0,"epistemico":"fato","origem":"wiki","rel":"usa"}],"confianca":1.0,"contraexemplos":[],"descricao":"A quiralidade é o único resíduo não-interno — é a quiralidade de LORENTZ (spinor de Weyl). Cor/carga vivem em 𝕆 e COMUTAM com toda sl(2,ℂ) de Lorentz → Lorentz-escalares, vetoriais (forte/EM não violam paridade). O isospin fraco vive em ℍ (o fator que carrega Lorentz): T₃ TRANSFORMA sob sl(2,ℂ), e span_ℂ(fraco)=span_ℂ(Lorentz) é o mesmo sl(2,ℂ) → só a força fraca é quiral. [D/firme] (deriva, não põe à mão como Furey/Boyle); [J] a escolha do LADO (por que canhotos) é aberta. DISTINTA da quebra quiral da QCD (aquela é o condensado que dá massa aos hádrons).","epistemico":"inferencia","exemplos":[],"id":"quiralidade_lorentz_sl2c","memoria":"semantica","meta":{"autor":"broca-so","dominio":"cosmologia","fonte":"paper_4_particulas.tex obs:quiralidade-tese (quiralidade_tese.py)"},"origem":"wiki","palavras_chave":[],"sinonimos":[],"tipo":"conceito","titulo":"Quiralidade: só a força fraca viola paridade (a tese)"}
\section{Quiralidade: só a força fraca viola paridade (a tese)}
A quiralidade é o único resíduo não-interno — é a quiralidade de LORENTZ (spinor de Weyl). Cor/carga vivem em \ensuremath{\mathbb{O}} e COMUTAM com toda sl(2,\ensuremath{\mathbb{C}}) de Lorentz \ensuremath{\to} Lorentz-escalares, vetoriais (forte/EM não violam paridade). O isospin fraco vive em \ensuremath{\mathbb{H}} (o fator que carrega Lorentz): T\ensuremath{_{3}} TRANSFORMA sob sl(2,\ensuremath{\mathbb{C}}), e span\_\ensuremath{\mathbb{C}}(fraco)=span\_\ensuremath{\mathbb{C}}(Lorentz) é o mesmo sl(2,\ensuremath{\mathbb{C}}) \ensuremath{\to} só a força fraca é quiral. [D/firme] (deriva, não põe à mão como Furey/Boyle); [J] a escolha do LADO (por que canhotos) é aberta. DISTINTA da quebra quiral da QCD (aquela é o condensado que dá massa aos hádrons).
\prov{relações: explica unificacao\_eletrofraca\_gws; usa isospin\_hipercarga\_algebrico; relaciona quebra\_quiral\_massa\_hadrons; usa travessia\_wick\_m2c}
\prov{proveniência: wiki \textperiodcentered{} paper\_4\_particulas.tex obs:quiralidade-tese (quiralidade\_tese.py) \textperiodcentered{} inferencia \textperiodcentered{} confiança 1.0 \textperiodcentered{} domínio cosmologia \textperiodcentered{} história 5 versões no jornal}

%CRISTAL {"arestas":[{"alvo":"cosmologia_quantica_hub","confianca":1.0,"epistemico":"fato","origem":"wiki","rel":"parte_de"}],"confianca":1.0,"contraexemplos":[],"descricao":"Efeitos quânticos em espaço-tempo curvo fazem o buraco negro emitir radiação térmica: T_H=ℏc³/(8πGMk_B) (∝1/M), evaporação t∝M³ (~10⁶⁷ anos p/ 1 M☉). Como o espectro só depende de M, a evaporação parece apagar a informação — o paradoxo. Expectativa moderna (curva de Page, 'ilhas'/replica wormholes ~2019): a informação se preserva, mas o mecanismo completo é aberto.","epistemico":"fato","exemplos":[],"id":"radiacao_hawking_informacao","memoria":"semantica","meta":{"autor":"broca-so","dominio":"fisica","fonte":"Hawking 1974 Nature 248,30; 1975 CMP 43,199 (paradoxo: aberto)"},"origem":"wiki","palavras_chave":[],"sinonimos":[],"tipo":"conceito","titulo":"Radiação Hawking e o paradoxo da informação"}
\section{Radiação Hawking e o paradoxo da informação}
Efeitos quânticos em espaço-tempo curvo fazem o buraco negro emitir radiação térmica: T\_H=\ensuremath{\hbar}c\ensuremath{^{3}}/(8\ensuremath{\pi}GMk\_B) (\ensuremath{\propto}1/M), evaporação t\ensuremath{\propto}M\ensuremath{^{3}} (\~{}10\ensuremath{^{67}} anos p/ 1 M\ensuremath{\odot}). Como o espectro só depende de M, a evaporação parece apagar a informação — o paradoxo. Expectativa moderna (curva de Page, 'ilhas'/replica wormholes \~{}2019): a informação se preserva, mas o mecanismo completo é aberto.
\prov{relações: parte\_de cosmologia\_quantica\_hub}
\prov{proveniência: wiki \textperiodcentered{} Hawking 1974 Nature 248,30; 1975 CMP 43,199 (paradoxo: aberto) \textperiodcentered{} fato \textperiodcentered{} confiança 1.0 \textperiodcentered{} domínio fisica \textperiodcentered{} história 2 versões no jornal}

%CRISTAL {"arestas":[{"alvo":"compilador_glsl_broca","confianca":1.0,"epistemico":"fato","origem":"wiki","rel":"usa"},{"alvo":"shader_quilez_no_broca","confianca":1.0,"epistemico":"fato","origem":"wiki","rel":"relaciona"},{"alvo":"joalheria_hub","confianca":1.0,"epistemico":"fato","origem":"wiki","rel":"parte_de"}],"confianca":1.0,"contraexemplos":[],"descricao":"O 'Pyramid - distance' (Íñigo Quílez, primitives2.txt) rodando no metal — um paradigma DIFERENTE do intersector analítico: aqui é DISTÂNCIA COM SINAL (SDF) + RAYMARCHING. O compilador GLSL compila o sdPyramid de verdade (a distância EXATA à pirâmide, com o arredondamento), e o broca-so roda o map (o campo SDF ANIMADO por iTime), o calcNormal (o gradiente tetraédrico, 4 amostras do map) e o loop de SPHERE TRACING (t += h) VETORIZADO — todos os raios marcham juntos, a máscara desativa os que convergem (h<eps) ou escapam (t>tmax). É TUDO DINÂMICO PELO CORPO ESTELAR: o map já pulsa (iTime arredonda e cresce a pirâmide), e a costura dos frames dirige o iTime pela RETA MINERAL (⊕), fechando no ramo CRISTAL da parábola (sem emenda). Mostra que o compilador cobre os DOIS mundos do Quílez: o analítico (intersectores) e o SDF (raymarching). A lib do raymarcher hospeda DUAS primitivas SDF do Quílez, ambas compiladas do GLSL: a PIRÂMIDE (sdPyramid, primitives2.txt) e a AMÊNDOA (vertical_vesica_segment, primitives3.txt — o vesica vertical, a mandorla, a forma do olho; a cúspide fica exatamente na superfície). O render é genérico (render_sdf recebe o map); render_piramide e render_vesica são wrappers; costura_piramide/costura_vesica animam pela reta mineral. E a AMÊNDOA é o LABORATÓRIO PADRÃO: com o BLEND (o smooth minimum smin de Quílez, gl.smin) duas amêndoas cruzadas se fundem — o min é a costura DURA (dois cristais, a borda afiada), o smin abre uma BORDA de largura k (o Salem entre os cristais): a parábola cristal(k→0)↔borda(k grande) feita geometria. E a ESTRELA DE 5 PONTAS: 5 amêndoas radiais (72°) fundidas pelo smin EXPONENCIAL (log-sum-exp, simétrico — o polinomial pareado quebraria a simetria por não ser associativo), dando 5-fold EXATA; é a figura do metal DOURADO (σ_1=φ, o pentagrama áureo). E a FLOR: um morph∈[0,1] leva a ESTRELA → FLOR (as pontas afiadas arredondam em pétalas gordas, arredondando cada amêndoa por −rad e mantendo o blend pequeno pra preservar os entalhes), na mesma 5-fold do ouro (o pentagrama vira a flor de 5 pétalas). linguagem/shader_piramide.py (13/13): sdPyramid, vertical_vesica_segment, mapa_piramide/vesica/lab/estrela5/flor, _amendoa, calcNormal, raymarch, render_sdf, render_lab/estrela5/flor, costura. gl.smin/smax/smin_exp. Assets: shader_piramide/vesica/lab/estrela5/flor _cena e _giro.","epistemico":"fato","exemplos":[],"id":"raymarcher_sdf_broca","memoria":"semantica","meta":{"dominio":"joalheria","fonte":"maquina de corte"},"origem":"wiki","palavras_chave":[],"sinonimos":[],"tipo":"conceito","titulo":"O shader SDF do Quílez (Pyramid) no broca-so: distância com sinal + sphere tracing"}
\section{O shader SDF do Quílez (Pyramid) no broca-so: distância com sinal + sphere tracing}
O 'Pyramid - distance' (Íñigo Quílez, primitives2.txt) rodando no metal — um paradigma DIFERENTE do intersector analítico: aqui é DISTÂNCIA COM SINAL (SDF) + RAYMARCHING. O compilador GLSL compila o sdPyramid de verdade (a distância EXATA à pirâmide, com o arredondamento), e o broca-so roda o map (o campo SDF ANIMADO por iTime), o calcNormal (o gradiente tetraédrico, 4 amostras do map) e o loop de SPHERE TRACING (t += h) VETORIZADO — todos os raios marcham juntos, a máscara desativa os que convergem (h<eps) ou escapam (t>tmax). É TUDO DINÂMICO PELO CORPO ESTELAR: o map já pulsa (iTime arredonda e cresce a pirâmide), e a costura dos frames dirige o iTime pela RETA MINERAL (\ensuremath{\oplus}), fechando no ramo CRISTAL da parábola (sem emenda). Mostra que o compilador cobre os DOIS mundos do Quílez: o analítico (intersectores) e o SDF (raymarching). A lib do raymarcher hospeda DUAS primitivas SDF do Quílez, ambas compiladas do GLSL: a PIRÂMIDE (sdPyramid, primitives2.txt) e a AMÊNDOA (vertical\_vesica\_segment, primitives3.txt — o vesica vertical, a mandorla, a forma do olho; a cúspide fica exatamente na superfície). O render é genérico (render\_sdf recebe o map); render\_piramide e render\_vesica são wrappers; costura\_piramide/costura\_vesica animam pela reta mineral. E a AMÊNDOA é o LABORATÓRIO PADRÃO: com o BLEND (o smooth minimum smin de Quílez, gl.smin) duas amêndoas cruzadas se fundem — o min é a costura DURA (dois cristais, a borda afiada), o smin abre uma BORDA de largura k (o Salem entre os cristais): a parábola cristal(k\ensuremath{\to}0)\ensuremath{\leftrightarrow}borda(k grande) feita geometria. E a ESTRELA DE 5 PONTAS: 5 amêndoas radiais (72$^{\circ}$) fundidas pelo smin EXPONENCIAL (log-sum-exp, simétrico — o polinomial pareado quebraria a simetria por não ser associativo), dando 5-fold EXATA; é a figura do metal DOURADO (\ensuremath{\sigma}\_1=\ensuremath{\varphi}, o pentagrama áureo). E a FLOR: um morph\ensuremath{\in}[0,1] leva a ESTRELA \ensuremath{\to} FLOR (as pontas afiadas arredondam em pétalas gordas, arredondando cada amêndoa por \ensuremath{-}rad e mantendo o blend pequeno pra preservar os entalhes), na mesma 5-fold do ouro (o pentagrama vira a flor de 5 pétalas). linguagem/shader\_piramide.py (13/13): sdPyramid, vertical\_vesica\_segment, mapa\_piramide/vesica/lab/estrela5/flor, \_amendoa, calcNormal, raymarch, render\_sdf, render\_lab/estrela5/flor, costura. gl.smin/smax/smin\_exp. Assets: shader\_piramide/vesica/lab/estrela5/flor \_cena e \_giro.
\prov{relações: usa compilador\_glsl\_broca; relaciona shader\_quilez\_no\_broca; parte\_de joalheria\_hub}
\prov{proveniência: wiki \textperiodcentered{} maquina de corte \textperiodcentered{} fato \textperiodcentered{} confiança 1.0 \textperiodcentered{} domínio joalheria \textperiodcentered{} história 64 versões no jornal}

%CRISTAL {"arestas":[{"alvo":"fissao_nuclear","confianca":1.0,"epistemico":"fato","origem":"wiki","rel":"usa"}],"confianca":1.0,"contraexemplos":[],"descricao":"Cada fissão emite nêutrons que provocam novas fissões — uma CADEIA que se auto-sustenta se cada geração produz ao menos uma fissão seguinte. É a órbita cujo crescimento o fator k governa.","epistemico":"fato","exemplos":[],"id":"reacao_em_cadeia","memoria":"semantica","meta":{"dominio":"fisica","fonte":"energia nuclear"},"origem":"wiki","palavras_chave":[],"sinonimos":[],"tipo":"conceito","titulo":"Reação em Cadeia"}
\section{Reação em Cadeia}
Cada fissão emite nêutrons que provocam novas fissões — uma CADEIA que se auto-sustenta se cada geração produz ao menos uma fissão seguinte. É a órbita cujo crescimento o fator k governa.
\prov{relações: usa fissao\_nuclear}
\prov{proveniência: wiki \textperiodcentered{} energia nuclear \textperiodcentered{} fato \textperiodcentered{} confiança 1.0 \textperiodcentered{} domínio fisica \textperiodcentered{} história 2 versões no jornal}

%CRISTAL {"arestas":[{"alvo":"criticidade","confianca":1.0,"epistemico":"fato","origem":"wiki","rel":"usa"},{"alvo":"fissao_nuclear","confianca":1.0,"epistemico":"fato","origem":"wiki","rel":"usa"},{"alvo":"numero_de_salem","confianca":1.0,"epistemico":"fato","origem":"wiki","rel":"relaciona"},{"alvo":"energia_eletrica","confianca":1.0,"epistemico":"fato","origem":"wiki","rel":"explica"}],"confianca":1.0,"contraexemplos":[],"descricao":"Uma fissão em cadeia CONTROLADA em k=1: barras de controle absorvem nêutrons, o moderador os desacelera, e o calor gera vapor→turbina→energia elétrica. Fica preso na BORDA (crítico) — o número de Salem da energia.","epistemico":"fato","exemplos":[],"id":"reator_nuclear","memoria":"semantica","meta":{"dominio":"fisica","fonte":"energia nuclear"},"origem":"wiki","palavras_chave":[],"sinonimos":[],"tipo":"conceito","titulo":"Reator Nuclear"}
\section{Reator Nuclear}
Uma fissão em cadeia CONTROLADA em k=1: barras de controle absorvem nêutrons, o moderador os desacelera, e o calor gera vapor\ensuremath{\to}turbina\ensuremath{\to}energia elétrica. Fica preso na BORDA (crítico) — o número de Salem da energia.
\prov{relações: usa criticidade; usa fissao\_nuclear; relaciona numero\_de\_salem; explica energia\_eletrica}
\prov{proveniência: wiki \textperiodcentered{} energia nuclear \textperiodcentered{} fato \textperiodcentered{} confiança 1.0 \textperiodcentered{} domínio fisica \textperiodcentered{} história 5 versões no jornal}

%CRISTAL {"arestas":[{"alvo":"hopfield","confianca":1.0,"epistemico":"fato","origem":"wiki","rel":"explica"},{"alvo":"colapso_ping_pong_overlap","confianca":1.0,"epistemico":"fato","origem":"wiki","rel":"usa"},{"alvo":"neuron_io_popcnt","confianca":1.0,"epistemico":"fato","origem":"wiki","rel":"usa"},{"alvo":"floco_phi","confianca":1.0,"epistemico":"fato","origem":"wiki","rel":"relaciona"}],"confianca":1.0,"contraexemplos":[],"descricao":"Doutrina 'DADOS no disco; GATOS na RAM': impressao() torce Koch de cada arquivo num vetor de 256 spins e descarta o conteúdo; guarda() soma ξξᵀ Hebbian em W; corrompe 25% dos bits, recorda() desce a energia, e argmax de overlap identifica de volta o arquivo. O reconhecimento associativo no metal.","epistemico":"fato","exemplos":[],"id":"reconhece_c_memoria","memoria":"semantica","meta":{"autor":"broca-so","dominio":"metal","fonte":"code/reconhece.c"},"origem":"wiki","palavras_chave":[],"sinonimos":[],"tipo":"conceito","titulo":"Reconhecimento Associativo de Arquivos (reconhece.c)"}
\section{Reconhecimento Associativo de Arquivos (reconhece.c)}
Doutrina 'DADOS no disco; GATOS na RAM': impressao() torce Koch de cada arquivo num vetor de 256 spins e descarta o conteúdo; guarda() soma \ensuremath{\xi}\ensuremath{\xi}\ensuremath{^{T}} Hebbian em W; corrompe 25\% dos bits, recorda() desce a energia, e argmax de overlap identifica de volta o arquivo. O reconhecimento associativo no metal.
\prov{relações: explica hopfield; usa colapso\_ping\_pong\_overlap; usa neuron\_io\_popcnt; relaciona floco\_phi}
\prov{proveniência: wiki \textperiodcentered{} code/reconhece.c \textperiodcentered{} fato \textperiodcentered{} confiança 1.0 \textperiodcentered{} domínio metal \textperiodcentered{} história 5 versões no jornal}

%CRISTAL {"arestas":[{"alvo":"maquina_de_corte","confianca":1.0,"epistemico":"fato","origem":"wiki","rel":"usa"},{"alvo":"corpo_estelar_le_o_rastro","confianca":1.0,"epistemico":"fato","origem":"wiki","rel":"usa"},{"alvo":"joalheria_hub","confianca":1.0,"epistemico":"fato","origem":"wiki","rel":"parte_de"},{"alvo":"tabela_geometrias_espectrais","confianca":1.0,"epistemico":"fato","origem":"wiki","rel":"relaciona"}],"confianca":1.0,"contraexemplos":[],"descricao":"O passo 2 da fabricação. RECONSTRUÇÃO DOS MODOS: o SVD da imagem É o espectro exato que o gato Anosov deixou — somar todos os modos devolve a imagem a erro 0; a economia da entropia H reconstrói com poucos (no gabarito ~27 modos capturam 98% da energia, ~3 modos efetivos). TESSELAÇÃO 3D: preserva o CRISTAL COMPLETO do gabarito (a assinatura da gema) — só isola a geometria do fundo (a segmentação), lê o CONTORNO REAL e lofta em anéis (rondiz/mesa/culote), MANTENDO os pontos do contorno na malha (o trilhão de pontos). A assinatura geométrica é preservada; exporta OBJ. linguagem/reconstrucao_3d.py (5/5); joalheria.tesselar/reconstruir.","epistemico":"fato","exemplos":[],"id":"reconstrucao_modos_tesselacao","memoria":"semantica","meta":{"dominio":"joalheria","fonte":"maquina de corte"},"origem":"wiki","palavras_chave":[],"sinonimos":[],"tipo":"conceito","titulo":"A Reconstrução dos Modos + a Tesselação 3D (preservar a assinatura)"}
\section{A Reconstrução dos Modos + a Tesselação 3D (preservar a assinatura)}
O passo 2 da fabricação. RECONSTRUÇÃO DOS MODOS: o SVD da imagem É o espectro exato que o gato Anosov deixou — somar todos os modos devolve a imagem a erro 0; a economia da entropia H reconstrói com poucos (no gabarito \~{}27 modos capturam 98\% da energia, \~{}3 modos efetivos). TESSELAÇÃO 3D: preserva o CRISTAL COMPLETO do gabarito (a assinatura da gema) — só isola a geometria do fundo (a segmentação), lê o CONTORNO REAL e lofta em anéis (rondiz/mesa/culote), MANTENDO os pontos do contorno na malha (o trilhão de pontos). A assinatura geométrica é preservada; exporta OBJ. linguagem/reconstrucao\_3d.py (5/5); joalheria.tesselar/reconstruir.
\prov{relações: usa maquina\_de\_corte; usa corpo\_estelar\_le\_o\_rastro; parte\_de joalheria\_hub; relaciona tabela\_geometrias\_espectrais}
\prov{proveniência: wiki \textperiodcentered{} maquina de corte \textperiodcentered{} fato \textperiodcentered{} confiança 1.0 \textperiodcentered{} domínio joalheria \textperiodcentered{} história 245 versões no jornal}

%CRISTAL {"arestas":[{"alvo":"nucleossintese_bbn","confianca":1.0,"epistemico":"fato","origem":"wiki","rel":"explica"},{"alvo":"cosmologia_quantica_hub","confianca":1.0,"epistemico":"fato","origem":"wiki","rel":"parte_de"}],"confianca":1.0,"contraexemplos":[],"descricao":"A expansão estica o comprimento de onda: 1+z=a(t₀)/a(t_emissão). Extrapolando para trás → estado quente e denso (Big Bang). Horizonte de partículas = distância máxima da qual a luz nos alcançou. Confirmado por CMB (Penzias-Wilson 1965) e BBN. Aberto: singularidade inicial.","epistemico":"fato","exemplos":[],"id":"redshift_bigbang_horizonte","memoria":"semantica","meta":{"autor":"broca-so","dominio":"fisica","fonte":"Lemaître 1931; Gamow-Alpher-Herman 1948"},"origem":"wiki","palavras_chave":[],"sinonimos":[],"tipo":"conceito","titulo":"Redshift cosmológico, Big Bang quente e horizontes"}
\section{Redshift cosmológico, Big Bang quente e horizontes}
A expansão estica o comprimento de onda: 1+z=a(t\ensuremath{_{0}})/a(t\_emissão). Extrapolando para trás \ensuremath{\to} estado quente e denso (Big Bang). Horizonte de partículas = distância máxima da qual a luz nos alcançou. Confirmado por CMB (Penzias-Wilson 1965) e BBN. Aberto: singularidade inicial.
\prov{relações: explica nucleossintese\_bbn; parte\_de cosmologia\_quantica\_hub}
\prov{proveniência: wiki \textperiodcentered{} Lemaître 1931; Gamow-Alpher-Herman 1948 \textperiodcentered{} fato \textperiodcentered{} confiança 1.0 \textperiodcentered{} domínio fisica \textperiodcentered{} história 3 versões no jornal}

%CRISTAL {"arestas":[{"alvo":"cpu_metal_c","confianca":1.0,"epistemico":"fato","origem":"wiki","rel":"usa"},{"alvo":"word_espectral","confianca":1.0,"epistemico":"fato","origem":"wiki","rel":"usa"}],"confianca":1.0,"contraexemplos":[],"descricao":"Regs tem A, B, R (cada um um Wordlong total, long e), pc, flags (byte) e spectral (byte, liga a cifra automática no LOAD). Não há registrador de propósito geral indexado — só esses três acumuladores fixos. FLAGS (ZERO/EQ/LT) vêm de regssetflagscmp comparando por.total e.e.","epistemico":"fato","exemplos":[],"id":"registradores_abr","memoria":"semantica","meta":{"autor":"broca-so","dominio":"metal","fonte":"ula/registradores.c:Regs/regs_set_flags_cmp"},"origem":"wiki","palavras_chave":[],"sinonimos":[],"tipo":"conceito","titulo":"Banco de Registradores Real (A/B/R + flags)"}
\section{Banco de Registradores Real (A/B/R + flags)}
Regs tem A, B, R (cada um um Wordlong total, long e), pc, flags (byte) e spectral (byte, liga a cifra automática no LOAD). Não há registrador de propósito geral indexado — só esses três acumuladores fixos. FLAGS (ZERO/EQ/LT) vêm de regssetflagscmp comparando por.total e.e.
\prov{relações: usa cpu\_metal\_c; usa word\_espectral}
\prov{proveniência: wiki \textperiodcentered{} ula/registradores.c:Regs/regs\_set\_flags\_cmp \textperiodcentered{} fato \textperiodcentered{} confiança 1.0 \textperiodcentered{} domínio metal \textperiodcentered{} história 4 versões no jornal}

%CRISTAL {"arestas":[{"alvo":"cosmologia_quantica","confianca":1.0,"epistemico":"fato","origem":"wiki","rel":"parte_de"},{"alvo":"lapse_selecao_desitter","confianca":1.0,"epistemico":"fato","origem":"wiki","rel":"usa"},{"alvo":"koch_parabola","confianca":1.0,"epistemico":"fato","origem":"wiki","rel":"relaciona"}],"confianca":1.0,"contraexemplos":[],"descricao":"A métrica quântica de Lopes Ⓚ(x,y)=min|x−y|,1−|x−y| sobre ℝ/ℤ é o comprimento de arco no círculo de estados (c,√a)=(sinψ,cosψ); ψ é a coordenada natural do espaço de álgebras (a euclidiana cega perto do quaternion). Régua e potencial são o mesmo objeto: gψ=‖∂ψ(z·ψw)‖²=(1−s²)‖a×b‖²=V.","epistemico":"inferencia","exemplos":[],"id":"regua_psi_metrica_quantica","memoria":"semantica","meta":{"autor":"broca-so","dominio":"cosmologia","fonte":"hiper/circular/paper_metrica_quantica_algebras.tex; paper_0_sintese.tex thm:saturacao"},"origem":"wiki","palavras_chave":[],"sinonimos":[],"tipo":"conceito","titulo":"A régua ψ — métrica quântica = geometria das álgebras"}
\section{A régua \ensuremath{\psi} — métrica quântica = geometria das álgebras}
A métrica quântica de Lopes (K)(x,y)=min|x\ensuremath{-}y|,1\ensuremath{-}|x\ensuremath{-}y| sobre \ensuremath{\mathbb{R}}/\ensuremath{\mathbb{Z}} é o comprimento de arco no círculo de estados (c,\ensuremath{\sqrt{\,}}a)=(sin\ensuremath{\psi},cos\ensuremath{\psi}); \ensuremath{\psi} é a coordenada natural do espaço de álgebras (a euclidiana cega perto do quaternion). Régua e potencial são o mesmo objeto: g\ensuremath{\psi}=\ensuremath{\|}\ensuremath{\partial}\ensuremath{\psi}(z\textperiodcentered \ensuremath{\psi}w)\ensuremath{\|}\ensuremath{^{2}}=(1\ensuremath{-}s\ensuremath{^{2}})\ensuremath{\|}a$\times$b\ensuremath{\|}\ensuremath{^{2}}=V.
\prov{relações: parte\_de cosmologia\_quantica; usa lapse\_selecao\_desitter; relaciona koch\_parabola}
\prov{proveniência: wiki \textperiodcentered{} hiper/circular/paper\_metrica\_quantica\_algebras.tex; paper\_0\_sintese.tex thm:saturacao \textperiodcentered{} inferencia \textperiodcentered{} confiança 1.0 \textperiodcentered{} domínio cosmologia \textperiodcentered{} história 6 versões no jornal}

%CRISTAL {"arestas":[{"alvo":"equacao_gema_estelar","confianca":1.0,"epistemico":"fato","origem":"wiki","rel":"usa"},{"alvo":"transformacoes_estelares","confianca":1.0,"epistemico":"fato","origem":"wiki","rel":"usa"},{"alvo":"joalheria_hub","confianca":1.0,"epistemico":"fato","origem":"wiki","rel":"parte_de"},{"alvo":"area_render_gema","confianca":1.0,"epistemico":"fato","origem":"wiki","rel":"relaciona"}],"confianca":1.0,"contraexemplos":[],"descricao":"A forma (a equação 𝒢) recebe a LUZ por um raytracer no metal, rodando ONDE A MINERAÇÃO RODA (gex44, 20 núcleos, paralelo, pelo canal do olho/chave claude). As técnicas físicas aplicadas: SNELL (refração entrada/saída, η=n₁/n₂, n=1.62); REFLEXÃO INTERNA TOTAL (as COLISÕES: η·sinθ>1 → ricocheteia por dentro, até B ricochetes = o fogo); BEER-LAMBERT (a cor: I=I₀·e^(−α·ℓ), α=(2.46,0.044,0) — come o vermelho, deixa o azul-verde do Paraíba, calibrado ao gabarito); FRESNEL (o BLEND reflexão/refração, Schlick R=R₀+(1−R₀)(1−cosθ)⁵, o brilho da faceta); ACES (tone mapping filmic). A COR NÃO É FIXA — EMERGE: o espectro do SOL (corpo negro 5778 K) × a COMPOSIÇÃO ATÔMICA do centro (Cu²⁺ come o vermelho ~700 nm, Mn³⁺ o amarelo/violeta) × o olho (CIE 1931): c(ℓ)=∫Sol(λ)·e^(−α(λ)ℓ)·CMF(λ)dλ → sRGB; a janela ~480–540 nm que sobra é o azul-verde Paraíba (Beer-Lambert POR λ, LUT por comprimento de caminho). A gema da equação é fechada por ConvexHull (convexa p/ o TIR). O turntable: a ÓRBITA da reta mineral ⊕ move os quadros, o raytracer renderiza a física. linguagem/raytracer.gema_equacao, gerar_logo_3d(gema=…); scripts/gex44/render_gema.sh (hero) e render_gema_giro.sh (giro). Hero: 198 facetas, 680×870, 6 ricochetes, 156s/3123 núcleo-s.","epistemico":"fato","exemplos":[],"id":"render_fisico_gex","memoria":"semantica","meta":{"dominio":"joalheria","fonte":"maquina de corte"},"origem":"wiki","palavras_chave":[],"sinonimos":[],"tipo":"conceito","titulo":"A física da pedra no metal (raytracer na gex44): Snell, TIR, Beer-Lambert, Fresnel, ACES"}
\section{A física da pedra no metal (raytracer na gex44): Snell, TIR, Beer-Lambert, Fresnel, ACES}
A forma (a equação \ensuremath{\mathcal{G}}) recebe a LUZ por um raytracer no metal, rodando ONDE A MINERAÇÃO RODA (gex44, 20 núcleos, paralelo, pelo canal do olho/chave claude). As técnicas físicas aplicadas: SNELL (refração entrada/saída, \ensuremath{\eta}=n\ensuremath{_{1}}/n\ensuremath{_{2}}, n=1.62); REFLEXÃO INTERNA TOTAL (as COLISÕES: \ensuremath{\eta}\textperiodcentered sin\ensuremath{\theta}>1 \ensuremath{\to} ricocheteia por dentro, até B ricochetes = o fogo); BEER-LAMBERT (a cor: I=I\ensuremath{_{0}}\textperiodcentered e\^{}(\ensuremath{-}\ensuremath{\alpha}\textperiodcentered \ensuremath{\ell}), \ensuremath{\alpha}=(2.46,0.044,0) — come o vermelho, deixa o azul-verde do Paraíba, calibrado ao gabarito); FRESNEL (o BLEND reflexão/refração, Schlick R=R\ensuremath{_{0}}+(1\ensuremath{-}R\ensuremath{_{0}})(1\ensuremath{-}cos\ensuremath{\theta})\ensuremath{^{5}}, o brilho da faceta); ACES (tone mapping filmic). A COR NÃO É FIXA — EMERGE: o espectro do SOL (corpo negro 5778 K) $\times$ a COMPOSIÇÃO ATÔMICA do centro (Cu\ensuremath{^{2+}} come o vermelho \~{}700 nm, Mn\ensuremath{^{3+}} o amarelo/violeta) $\times$ o olho (CIE 1931): c(\ensuremath{\ell})=\ensuremath{\int}Sol(\ensuremath{\lambda})\textperiodcentered e\^{}(\ensuremath{-}\ensuremath{\alpha}(\ensuremath{\lambda})\ensuremath{\ell})\textperiodcentered CMF(\ensuremath{\lambda})d\ensuremath{\lambda} \ensuremath{\to} sRGB; a janela \~{}480–540 nm que sobra é o azul-verde Paraíba (Beer-Lambert POR \ensuremath{\lambda}, LUT por comprimento de caminho). A gema da equação é fechada por ConvexHull (convexa p/ o TIR). O turntable: a ÓRBITA da reta mineral \ensuremath{\oplus} move os quadros, o raytracer renderiza a física. linguagem/raytracer.gema\_equacao, gerar\_logo\_3d(gema=…); scripts/gex44/render\_gema.sh (hero) e render\_gema\_giro.sh (giro). Hero: 198 facetas, 680$\times$870, 6 ricochetes, 156s/3123 núcleo-s.
\prov{relações: usa equacao\_gema\_estelar; usa transformacoes\_estelares; parte\_de joalheria\_hub; relaciona area\_render\_gema}
\prov{proveniência: wiki \textperiodcentered{} maquina de corte \textperiodcentered{} fato \textperiodcentered{} confiança 1.0 \textperiodcentered{} domínio joalheria \textperiodcentered{} história 185 versões no jornal}

%CRISTAL {"arestas":[{"alvo":"orbita_mineracao","confianca":1.0,"epistemico":"fato","origem":"wiki","rel":"parte_de"},{"alvo":"cifra_de_cristal","confianca":1.0,"epistemico":"fato","origem":"wiki","rel":"relaciona"}],"confianca":1.0,"contraexemplos":[],"descricao":"Sempre sobra um RESÍDUO na mineração (o floco irracional nunca acaba), e ele mora nos ÚLTIMOS 2-QUBITS do último passo (|ψ|²=1, o checksum a+c²=1 da transformada 7D). Esse resíduo é o CHECKPOINT: refinar contínuo == interromper (salvar o resíduo) + retomar, sem perda. Assim se CONTROLA O TEMPO da mineração — interromper e retomar quando quiser. É o gancho por onde o pipe segura o tempo.","epistemico":"fato","exemplos":[],"id":"residuo_2qubit_checkpoint","memoria":"semantica","meta":{"dominio":"mineracao_pipe","fonte":"mineracao_pipe hub"},"origem":"wiki","palavras_chave":[],"sinonimos":[],"tipo":"conceito","titulo":"O resíduo nos 2-qubits (o checkpoint que controla o tempo)"}
\section{O resíduo nos 2-qubits (o checkpoint que controla o tempo)}
Sempre sobra um RESÍDUO na mineração (o floco irracional nunca acaba), e ele mora nos ÚLTIMOS 2-QUBITS do último passo (|\ensuremath{\psi}|\ensuremath{^{2}}=1, o checksum a+c\ensuremath{^{2}}=1 da transformada 7D). Esse resíduo é o CHECKPOINT: refinar contínuo == interromper (salvar o resíduo) + retomar, sem perda. Assim se CONTROLA O TEMPO da mineração — interromper e retomar quando quiser. É o gancho por onde o pipe segura o tempo.
\prov{relações: parte\_de orbita\_mineracao; relaciona cifra\_de\_cristal}
\prov{proveniência: wiki \textperiodcentered{} mineracao\_pipe hub \textperiodcentered{} fato \textperiodcentered{} confiança 1.0 \textperiodcentered{} domínio mineracao\_pipe \textperiodcentered{} história 3 versões no jornal}

%CRISTAL {"arestas":[{"alvo":"ribossomo_koch","confianca":1.0,"epistemico":"fato","origem":"wiki","rel":"especializa"},{"alvo":"decomposicao_peirce","confianca":1.0,"epistemico":"fato","origem":"wiki","rel":"usa"},{"alvo":"koch_dissipa_nao_codifica","confianca":1.0,"epistemico":"fato","origem":"wiki","rel":"usa"},{"alvo":"ribossomo_ribozima","confianca":1.0,"epistemico":"fato","origem":"wiki","rel":"relaciona"},{"alvo":"impressao_koch_compressao","confianca":1.0,"epistemico":"fato","origem":"wiki","rel":"relaciona"},{"alvo":"colapso_koch_tiffany","confianca":1.0,"epistemico":"fato","origem":"wiki","rel":"usa"}],"confianca":1.0,"contraexemplos":[],"descricao":"A substituição fractal de Koch satisfaz as 7 propriedades do ribossomo (tradutor, processivo, universal, código fixo, ribozima=auto-similar) e 'traduz' o código (3 cúspides algébricas) em carne (curva fractal). Age SÓ no canal negro (nilpotente/glacial) da decomposição de Peirce: z↦z² em ℤ[i] é 2-para-1 (identifica z e −z) ⇒ nenhuma matriz linear a representa — QUANTIZAÇÃO e SEGURANÇA são a mesma operação (genótipo→fenótipo IRREVERSÍVEL). Percepção=negra→branca (C_K); ação=branca→negra (C_K⁻¹). [D] a álgebra de Koch (verif.); [J] a identidade 'É ribossomo' é analogia declarada (as 7 propriedades).","epistemico":"inferencia","exemplos":[],"id":"ribossomo_koch_codigo_carne","memoria":"semantica","meta":{"autor":"broca-so","dominio":"cosmologia","fonte":"papers/olho_de_koch.tex; ribossomo_koch"},"origem":"wiki","palavras_chave":[],"sinonimos":[],"tipo":"conceito","titulo":"Ribossomo = substituição de Koch (código→carne, metáfora derivada)"}
\section{Ribossomo = substituição de Koch (código\ensuremath{\to}carne, metáfora derivada)}
A substituição fractal de Koch satisfaz as 7 propriedades do ribossomo (tradutor, processivo, universal, código fixo, ribozima=auto-similar) e 'traduz' o código (3 cúspides algébricas) em carne (curva fractal). Age SÓ no canal negro (nilpotente/glacial) da decomposição de Peirce: z\ensuremath{\mapsto}z\ensuremath{^{2}} em \ensuremath{\mathbb{Z}}[i] é 2-para-1 (identifica z e \ensuremath{-}z) \ensuremath{\Rightarrow} nenhuma matriz linear a representa — QUANTIZAÇÃO e SEGURANÇA são a mesma operação (genótipo\ensuremath{\to}fenótipo IRREVERSÍVEL). Percepção=negra\ensuremath{\to}branca (C\_K); ação=branca\ensuremath{\to}negra (C\_K\ensuremath{^{-1}}). [D] a álgebra de Koch (verif.); [J] a identidade 'É ribossomo' é analogia declarada (as 7 propriedades).
\prov{relações: especializa ribossomo\_koch; usa decomposicao\_peirce; usa koch\_dissipa\_nao\_codifica; relaciona ribossomo\_ribozima; relaciona impressao\_koch\_compressao; usa colapso\_koch\_tiffany}
\prov{proveniência: wiki \textperiodcentered{} papers/olho\_de\_koch.tex; ribossomo\_koch \textperiodcentered{} inferencia \textperiodcentered{} confiança 1.0 \textperiodcentered{} domínio cosmologia \textperiodcentered{} história 7 versões no jornal}

%CRISTAL {"arestas":[{"alvo":"parabola_audit_dual","confianca":1.0,"epistemico":"fato","origem":"wiki","rel":"usa"},{"alvo":"corpo_tropical","confianca":1.0,"epistemico":"fato","origem":"wiki","rel":"relaciona"}],"confianca":1.0,"contraexemplos":[],"descricao":"Perturbações sobre Word: flip de bit (LCG), troca total↔e, offset ±amp, negação; ruido_medir reaudita a parábola e reporta delta=1−(a+c²). O experimento roda programas da ISA e mede se a parábola SOBREVIVE ao erro — a robustez do invariante.","epistemico":"fato","exemplos":[],"id":"ruido_difusao_word","memoria":"semantica","meta":{"autor":"broca-so","dominio":"metal","fonte":"ula/ruido.c; ruido_exp.c"},"origem":"wiki","palavras_chave":[],"sinonimos":[],"tipo":"conceito","titulo":"Ruído / Difusão sobre a Word"}
\section{Ruído / Difusão sobre a Word}
Perturbações sobre Word: flip de bit (LCG), troca total\ensuremath{\leftrightarrow}e, offset $\pm$amp, negação; ruido\_medir reaudita a parábola e reporta delta=1\ensuremath{-}(a+c\ensuremath{^{2}}). O experimento roda programas da ISA e mede se a parábola SOBREVIVE ao erro — a robustez do invariante.
\prov{relações: usa parabola\_audit\_dual; relaciona corpo\_tropical}
\prov{proveniência: wiki \textperiodcentered{} ula/ruido.c; ruido\_exp.c \textperiodcentered{} fato \textperiodcentered{} confiança 1.0 \textperiodcentered{} domínio metal \textperiodcentered{} história 3 versões no jornal}

%CRISTAL {"arestas":[{"alvo":"associador_nao_tensoriza","confianca":1.0,"epistemico":"fato","origem":"wiki","rel":"usa"},{"alvo":"barreira_hurwitz_qubits","confianca":1.0,"epistemico":"fato","origem":"wiki","rel":"usa"},{"alvo":"teorema_nao_clonagem","confianca":1.0,"epistemico":"fato","origem":"wiki","rel":"relaciona"}],"confianca":1.0,"contraexemplos":[],"descricao":"A região associativa (associador zero) é a 'superfície de ataque' (medida nula, listável: 7 sub-ℍ de Fano + diagonal); a não-associatividade genérica é a dureza estrutural (curvatura de 𝕆/𝕊 rígida — não-exata, obstrução de Gerstenhaber: 'cancelar a curvatura = perder a divisão'). Primitiva NA-Word (word problem não-associativo). VEREDITO HONESTO explícito: a não-associatividade dá ESTRUTURA, NÃO HARDNESS — o esquema reduz a MQ (NP-completo), mas as identidades de 𝕆 (Moufang, alternatividade) geram sizígias que estragam a hardness-na-média (Macaulay refuta 2¹⁷⁴). 'A mesma ordem que torna 𝕆 belo destrói a hardness.'","epistemico":"inferencia","exemplos":[],"id":"seguranca_algebrica_estrutura_nao_hardness","memoria":"semantica","meta":{"autor":"broca-so","dominio":"cosmologia","fonte":"paper_8_seguranca.tex; paper_G_no_cloning.tex §7"},"origem":"wiki","palavras_chave":[],"sinonimos":[],"tipo":"conceito","titulo":"Segurança algébrica: estrutura, não hardness (veredito honesto)"}
\section{Segurança algébrica: estrutura, não hardness (veredito honesto)}
A região associativa (associador zero) é a 'superfície de ataque' (medida nula, listável: 7 sub-\ensuremath{\mathbb{H}} de Fano + diagonal); a não-associatividade genérica é a dureza estrutural (curvatura de \ensuremath{\mathbb{O}}/\ensuremath{\mathbb{S}} rígida — não-exata, obstrução de Gerstenhaber: 'cancelar a curvatura = perder a divisão'). Primitiva NA-Word (word problem não-associativo). VEREDITO HONESTO explícito: a não-associatividade dá ESTRUTURA, NÃO HARDNESS — o esquema reduz a MQ (NP-completo), mas as identidades de \ensuremath{\mathbb{O}} (Moufang, alternatividade) geram sizígias que estragam a hardness-na-média (Macaulay refuta 2\ensuremath{^{174}}). 'A mesma ordem que torna \ensuremath{\mathbb{O}} belo destrói a hardness.'
\prov{relações: usa associador\_nao\_tensoriza; usa barreira\_hurwitz\_qubits; relaciona teorema\_nao\_clonagem}
\prov{proveniência: wiki \textperiodcentered{} paper\_8\_seguranca.tex; paper\_G\_no\_cloning.tex §7 \textperiodcentered{} inferencia \textperiodcentered{} confiança 1.0 \textperiodcentered{} domínio cosmologia \textperiodcentered{} história 4 versões no jornal}

%CRISTAL {"arestas":[{"alvo":"problema_gravidade_quantica","confianca":1.0,"epistemico":"fato","origem":"wiki","rel":"relaciona"}],"confianca":1.0,"contraexemplos":[],"descricao":"A gravidade seria não-perturbativamente renormalizável graças a um ponto fixo UV NÃO-trivial do fluxo do grupo de renormalização — os acoplamentos tenderiam a valores finitos no UV, tornando a teoria preditiva. Evidência sobretudo numérica (Functional RG, truncamentos); a existência rigorosa do ponto fixo é aberta.","epistemico":"hipotese","exemplos":[],"id":"seguranca_assintotica","memoria":"semantica","meta":{"autor":"broca-so","dominio":"fisica","fonte":"Weinberg 1976/79; Reuter 1998 hep-th/9605030"},"origem":"wiki","palavras_chave":[],"sinonimos":[],"tipo":"conceito","titulo":"Segurança assintótica (Weinberg)"}
\section{Segurança assintótica (Weinberg)}
A gravidade seria não-perturbativamente renormalizável graças a um ponto fixo UV NÃO-trivial do fluxo do grupo de renormalização — os acoplamentos tenderiam a valores finitos no UV, tornando a teoria preditiva. Evidência sobretudo numérica (Functional RG, truncamentos); a existência rigorosa do ponto fixo é aberta.
\prov{relações: relaciona problema\_gravidade\_quantica}
\prov{proveniência: wiki \textperiodcentered{} Weinberg 1976/79; Reuter 1998 hep-th/9605030 \textperiodcentered{} hipotese \textperiodcentered{} confiança 1.0 \textperiodcentered{} domínio fisica \textperiodcentered{} história 2 versões no jornal}

%CRISTAL {"arestas":[{"alvo":"fase_de_berry_canonica","confianca":1.0,"epistemico":"fato","origem":"wiki","rel":"usa"}],"confianca":1.0,"contraexemplos":[],"descricao":"Materiais 3D com bandas que se tocam em pontos nodais isolados; perto deles as quasipartículas são férmions de Weyl/Dirac relativísticos. Os nós de Weyl são MONOPOLOS de curvatura de Berry (quiralidade ±1, em pares) e geram estados de superfície abertos ('arcos de Fermi'). Previsto e confirmado em TaAs (2015). Fato consensual.","epistemico":"fato","exemplos":[],"id":"semimetais_weyl_dirac","memoria":"semantica","meta":{"autor":"broca-so","dominio":"fisica","fonte":"Xu+ (Hasan) 2015; Lv+ (Ding) 2015 PRX 5,031013"},"origem":"wiki","palavras_chave":[],"sinonimos":[],"tipo":"conceito","titulo":"Semimetais de Weyl e Dirac"}
\section{Semimetais de Weyl e Dirac}
Materiais 3D com bandas que se tocam em pontos nodais isolados; perto deles as quasipartículas são férmions de Weyl/Dirac relativísticos. Os nós de Weyl são MONOPOLOS de curvatura de Berry (quiralidade $\pm$1, em pares) e geram estados de superfície abertos ('arcos de Fermi'). Previsto e confirmado em TaAs (2015). Fato consensual.
\prov{relações: usa fase\_de\_berry\_canonica}
\prov{proveniência: wiki \textperiodcentered{} Xu+ (Hasan) 2015; Lv+ (Ding) 2015 PRX 5,031013 \textperiodcentered{} fato \textperiodcentered{} confiança 1.0 \textperiodcentered{} domínio fisica \textperiodcentered{} história 2 versões no jornal}

%CRISTAL {"arestas":[{"alvo":"materia_escura_lado_negro","confianca":1.0,"epistemico":"fato","origem":"wiki","rel":"especializa"},{"alvo":"teoria_gauge_octonica","confianca":1.0,"epistemico":"fato","origem":"wiki","rel":"usa"},{"alvo":"higgs_algebrico_g2_su3","confianca":1.0,"epistemico":"fato","origem":"wiki","rel":"usa"},{"alvo":"materia_escura","confianca":1.0,"epistemico":"fato","origem":"wiki","rel":"relaciona"},{"alvo":"wimps_candidato","confianca":1.0,"epistemico":"fato","origem":"wiki","rel":"contradiz"}],"confianca":1.0,"contraexemplos":[],"descricao":"G2=Aut(𝕆) (dim 14) decompõe sob SU(3)⊂G2 em 8 (glúons) + 6 bósons exóticos Xa (3⊕3̄), massa via condensado octônico mX≈28 MeV, acoplando ao MP só via mistura G2→SU(3) (Higgs estendido HG2). Predição: por freeze-out térmico, ΩX h²≈0,12 (=ΩDM medido) — candidato a matéria escura LEVE (MeV), NÃO WIMP-TeV (uma detecção WIMP-TeV o refutaria). É a encarnação concreta do lado negro. [N] predição falseável, mas [J]: o próprio paper marca 'estimativa dimensional' e o Yukawa yX~10⁹ é problemático.","epistemico":"hipotese","exemplos":[],"id":"setor_x_g2_materia_escura","memoria":"semantica","meta":{"autor":"broca-so","dominio":"cosmologia","fonte":"paper_AF_yang_mills_octonios.tex conj:G2-exotico, thm:abundancia"},"origem":"wiki","palavras_chave":[],"sinonimos":[],"tipo":"conceito","titulo":"O setor X de G2 como matéria escura leve (~28 MeV)"}
\section{O setor X de G2 como matéria escura leve (\~{}28 MeV)}
G2=Aut(\ensuremath{\mathbb{O}}) (dim 14) decompõe sob SU(3)\ensuremath{\subset}G2 em 8 (glúons) + 6 bósons exóticos Xa (3\ensuremath{\oplus}3\ensuremath{}), massa via condensado octônico mX\ensuremath{\approx}28 MeV, acoplando ao MP só via mistura G2\ensuremath{\to}SU(3) (Higgs estendido HG2). Predição: por freeze-out térmico, \ensuremath{\Omega}X h\ensuremath{^{2}}\ensuremath{\approx}0,12 (=\ensuremath{\Omega}DM medido) — candidato a matéria escura LEVE (MeV), NÃO WIMP-TeV (uma detecção WIMP-TeV o refutaria). É a encarnação concreta do lado negro. [N] predição falseável, mas [J]: o próprio paper marca 'estimativa dimensional' e o Yukawa yX\~{}10\ensuremath{^{9}} é problemático.
\prov{relações: especializa materia\_escura\_lado\_negro; usa teoria\_gauge\_octonica; usa higgs\_algebrico\_g2\_su3; relaciona materia\_escura; contradiz wimps\_candidato}
\prov{proveniência: wiki \textperiodcentered{} paper\_AF\_yang\_mills\_octonios.tex conj:G2-exotico, thm:abundancia \textperiodcentered{} hipotese \textperiodcentered{} confiança 1.0 \textperiodcentered{} domínio cosmologia \textperiodcentered{} história 7 versões no jornal}

%CRISTAL {"arestas":[{"alvo":"setor_x_g2_materia_escura","confianca":1.0,"epistemico":"fato","origem":"wiki","rel":"parte_de"},{"alvo":"g2_muon_anomalia_dissolvida","confianca":1.0,"epistemico":"fato","origem":"wiki","rel":"relaciona"}],"confianca":1.0,"contraexemplos":[],"descricao":"Os mesmos (fX≈3×10⁻⁵, mX≈28 MeV) davam δaμG²≈251×10⁻¹¹, casando com a discrepância Fermilab 2023 (251±59)×10⁻¹¹, ~4,2σ — a 'tripla coincidência' g-2/DM. HONESTIDADE (o Mandelbrot-que-falhou desta leva): a discrepância g-2 EVAPOROU em 2025 (a WP25 lattice subiu a previsão do MP; teoria≈experimento). Logo a perna g-2 casava com um alvo que se dissolveu — não é mais evidência. A perna matéria escura leve (ΩX≈0,12) segue como predição distinta. Não maquiar: o número batia numa anomalia que sumiu.","epistemico":"hipotese","exemplos":[],"id":"setor_x_g2_muon_alvo_movido","memoria":"semantica","meta":{"autor":"broca-so","dominio":"cosmologia","fonte":"paper_AF thm:g2-leading; WP25 2025 (dissolução)"},"origem":"wiki","palavras_chave":[],"sinonimos":[],"tipo":"conceito","titulo":"A perna g-2 do setor X: o alvo que se moveu"}
\section{A perna g-2 do setor X: o alvo que se moveu}
Os mesmos (fX\ensuremath{\approx}3$\times$10\ensuremath{^{-5}}, mX\ensuremath{\approx}28 MeV) davam \ensuremath{\delta}a\ensuremath{\mu}G\ensuremath{^{2}}\ensuremath{\approx}251$\times$10\ensuremath{^{-11}}, casando com a discrepância Fermilab 2023 (251$\pm$59)$\times$10\ensuremath{^{-11}}, \~{}4,2\ensuremath{\sigma} — a 'tripla coincidência' g-2/DM. HONESTIDADE (o Mandelbrot-que-falhou desta leva): a discrepância g-2 EVAPOROU em 2025 (a WP25 lattice subiu a previsão do MP; teoria\ensuremath{\approx}experimento). Logo a perna g-2 casava com um alvo que se dissolveu — não é mais evidência. A perna matéria escura leve (\ensuremath{\Omega}X\ensuremath{\approx}0,12) segue como predição distinta. Não maquiar: o número batia numa anomalia que sumiu.
\prov{relações: parte\_de setor\_x\_g2\_materia\_escura; relaciona g2\_muon\_anomalia\_dissolvida}
\prov{proveniência: wiki \textperiodcentered{} paper\_AF thm:g2-leading; WP25 2025 (dissolução) \textperiodcentered{} hipotese \textperiodcentered{} confiança 1.0 \textperiodcentered{} domínio cosmologia \textperiodcentered{} história 4 versões no jornal}

%CRISTAL {"arestas":[{"alvo":"compilador_glsl_broca","confianca":1.0,"epistemico":"fato","origem":"wiki","rel":"usa"},{"alvo":"primitivas_analiticas","confianca":1.0,"epistemico":"fato","origem":"wiki","rel":"usa"},{"alvo":"translucido_nijhoff","confianca":1.0,"epistemico":"fato","origem":"wiki","rel":"especializa"},{"alvo":"joalheria_hub","confianca":1.0,"epistemico":"fato","origem":"wiki","rel":"parte_de"}],"confianca":1.0,"contraexemplos":[],"descricao":"O 'Ray Tracing - Primitives' (Reinder Nijhoff / Íñigo Quílez) rodando INTEIRO no metal do broca-so. O compilador GLSL cobre os 14 intersectores do primitves.txt: os 8 diretos + os 6 que exigiam construtos novos — ternário-DENTRO-de-parênteses, if-SEM-chaves predicado, atribuição composta (op=), preprocessador #if/#else/#endif, chamada CRUZADA entre intersectores (iTorus chama iSphere; iMesh chama iTriangle) e cuberoot. Assim iCylinder/iCone/iTorus/iGoursat/iRoundedBox/iMesh compilam E acertam a geometria (t exato, normal correta). A CENA não é mais costurada à mão: a FUNÇÃO GERAL gerar_geometrias() usa o CHICOTE (o trilhão C_k=1+a·cos(kθ), a contração/epiciclo) pra pousar as primitivas numa trilha fechada, com os ângulos pela RETA MINERAL (⊕); o worldhit COSTURA todas pelo t mínimo (a união). O render() é o path tracer (LAMBERTIAN/METAL/DIELECTRIC, Fresnel-Schlick, TIR, getSkyColor). A COSTURA DOS FRAMES (costura_estelar) é o giro dirigido pela DINÂMICA DA PARÁBOLA DO CORPO ESTELAR: o azimute é a reta mineral somando por ⊕ (geodésica), e a costura fecha SEM emenda porque a órbita cai no ramo CRISTAL da parábola (ρ→0, órbita fechada) — o ajuste é ESPECTRAL (Δα=2π/n vem do metal σ). É o SIMT do GLSL virado SIMD por predicação. linguagem/shader_quilez.py (5/5): _compilar_intersectores, _CATALOGO, gerar_geometrias, worldhit, render, render_cena, costura_estelar. Asset: shader_quilez_giro.gif.","epistemico":"fato","exemplos":[],"id":"shader_quilez_no_broca","memoria":"semantica","meta":{"dominio":"joalheria","fonte":"maquina de corte"},"origem":"wiki","palavras_chave":[],"sinonimos":[],"tipo":"conceito","titulo":"O shader INTEIRO do Quílez rodando no broca-so (os 14 intersectores + o chicote + a parábola)"}
\section{O shader INTEIRO do Quílez rodando no broca-so (os 14 intersectores + o chicote + a parábola)}
O 'Ray Tracing - Primitives' (Reinder Nijhoff / Íñigo Quílez) rodando INTEIRO no metal do broca-so. O compilador GLSL cobre os 14 intersectores do primitves.txt: os 8 diretos + os 6 que exigiam construtos novos — ternário-DENTRO-de-parênteses, if-SEM-chaves predicado, atribuição composta (op=), preprocessador \#if/\#else/\#endif, chamada CRUZADA entre intersectores (iTorus chama iSphere; iMesh chama iTriangle) e cuberoot. Assim iCylinder/iCone/iTorus/iGoursat/iRoundedBox/iMesh compilam E acertam a geometria (t exato, normal correta). A CENA não é mais costurada à mão: a FUNÇÃO GERAL gerar\_geometrias() usa o CHICOTE (o trilhão C\_k=1+a\textperiodcentered cos(k\ensuremath{\theta}), a contração/epiciclo) pra pousar as primitivas numa trilha fechada, com os ângulos pela RETA MINERAL (\ensuremath{\oplus}); o worldhit COSTURA todas pelo t mínimo (a união). O render() é o path tracer (LAMBERTIAN/METAL/DIELECTRIC, Fresnel-Schlick, TIR, getSkyColor). A COSTURA DOS FRAMES (costura\_estelar) é o giro dirigido pela DINÂMICA DA PARÁBOLA DO CORPO ESTELAR: o azimute é a reta mineral somando por \ensuremath{\oplus} (geodésica), e a costura fecha SEM emenda porque a órbita cai no ramo CRISTAL da parábola (\ensuremath{\rho}\ensuremath{\to}0, órbita fechada) — o ajuste é ESPECTRAL (\ensuremath{\Delta}\ensuremath{\alpha}=2\ensuremath{\pi}/n vem do metal \ensuremath{\sigma}). É o SIMT do GLSL virado SIMD por predicação. linguagem/shader\_quilez.py (5/5): \_compilar\_intersectores, \_CATALOGO, gerar\_geometrias, worldhit, render, render\_cena, costura\_estelar. Asset: shader\_quilez\_giro.gif.
\prov{relações: usa compilador\_glsl\_broca; usa primitivas\_analiticas; especializa translucido\_nijhoff; parte\_de joalheria\_hub}
\prov{proveniência: wiki \textperiodcentered{} maquina de corte \textperiodcentered{} fato \textperiodcentered{} confiança 1.0 \textperiodcentered{} domínio joalheria \textperiodcentered{} história 91 versões no jornal}

%CRISTAL {"arestas":[{"alvo":"correcao_erro_quantica","confianca":1.0,"epistemico":"fato","origem":"wiki","rel":"relaciona"},{"alvo":"decoerencia_colapso_medida","confianca":1.0,"epistemico":"fato","origem":"wiki","rel":"relaciona"}],"confianca":1.0,"contraexemplos":[],"descricao":"SHOR (1994): fatora inteiros e computa log discreto em tempo POLINOMIAL (via estimativa de fase + QFT), quebrando RSA/ECC — motiva a criptografia pós-quântica; execução em escala útil ainda NÃO alcançada. GROVER (1996): busca não-estruturada em O(√N) vs O(N) clássico (speedup quadrático, provado ótimo). Fato teórico consolidado.","epistemico":"fato","exemplos":[],"id":"shor_grover","memoria":"semantica","meta":{"autor":"broca-so","dominio":"fisica","fonte":"Shor 1994 FOCS; Grover 1996 STOC arXiv:quant-ph/9605043"},"origem":"wiki","palavras_chave":[],"sinonimos":[],"tipo":"conceito","titulo":"Algoritmos de Shor e Grover"}
\section{Algoritmos de Shor e Grover}
SHOR (1994): fatora inteiros e computa log discreto em tempo POLINOMIAL (via estimativa de fase + QFT), quebrando RSA/ECC — motiva a criptografia pós-quântica; execução em escala útil ainda NÃO alcançada. GROVER (1996): busca não-estruturada em O(\ensuremath{\sqrt{\,}}N) vs O(N) clássico (speedup quadrático, provado ótimo). Fato teórico consolidado.
\prov{relações: relaciona correcao\_erro\_quantica; relaciona decoerencia\_colapso\_medida}
\prov{proveniência: wiki \textperiodcentered{} Shor 1994 FOCS; Grover 1996 STOC arXiv:quant-ph/9605043 \textperiodcentered{} fato \textperiodcentered{} confiança 1.0 \textperiodcentered{} domínio fisica \textperiodcentered{} história 3 versões no jornal}

%CRISTAL {"arestas":[{"alvo":"mecanismo_higgs_ebh","confianca":1.0,"epistemico":"fato","origem":"wiki","rel":"relaciona"},{"alvo":"bosons_w_z","confianca":1.0,"epistemico":"fato","origem":"wiki","rel":"explica"}],"confianca":1.0,"contraexemplos":[],"descricao":"Simetria global aproximada SU(2)_V do setor de Higgs que protege a relação de massas W/Z: ρ=M_W²/(M_Z²cos²θ_W)≈1 no nível de árvore (desvios ~10⁻³, correções radiativas do MP). Propriedade estrutural bem estabelecida.","epistemico":"fato","exemplos":[],"id":"simetria_custodial_rho","memoria":"semantica","meta":{"autor":"broca-so","dominio":"fisica","fonte":"Veltman 1977; PDG Electroweak Review"},"origem":"wiki","palavras_chave":[],"sinonimos":[],"tipo":"conceito","titulo":"Simetria custodial e o parâmetro ρ≈1"}
\section{Simetria custodial e o parâmetro \ensuremath{\rho}\ensuremath{\approx}1}
Simetria global aproximada SU(2)\_V do setor de Higgs que protege a relação de massas W/Z: \ensuremath{\rho}=M\_W\ensuremath{^{2}}/(M\_Z\ensuremath{^{2}}cos\ensuremath{^{2}}\ensuremath{\theta}\_W)\ensuremath{\approx}1 no nível de árvore (desvios \~{}10\ensuremath{^{-3}}, correções radiativas do MP). Propriedade estrutural bem estabelecida.
\prov{relações: relaciona mecanismo\_higgs\_ebh; explica bosons\_w\_z}
\prov{proveniência: wiki \textperiodcentered{} Veltman 1977; PDG Electroweak Review \textperiodcentered{} fato \textperiodcentered{} confiança 1.0 \textperiodcentered{} domínio fisica \textperiodcentered{} história 3 versões no jornal}

%CRISTAL {"arestas":[{"alvo":"de_sitter_lambda_energia_escura","confianca":1.0,"epistemico":"fato","origem":"wiki","rel":"explica"},{"alvo":"cosmologia_quantica_hub","confianca":1.0,"epistemico":"fato","origem":"wiki","rel":"parte_de"}],"confianca":1.0,"contraexemplos":[],"descricao":"SNe Ia são velas padronizáveis. 1998-99: SNe distantes mais fracas que o esperado → expansão ACELERADA → energia escura (Riess+1998; Perlmutter+1999; Nobel 2011). Pantheon+ (2022): 1701 curvas de 1550 SNe, w₀=−0,90±0,14 (só SNe); H₀=73,5±1,1 com Cefeidas. Aceleração = fato consolidado; a causa é o que fica aberto.","epistemico":"fato","exemplos":[],"id":"sne_ia_pantheon","memoria":"semantica","meta":{"autor":"broca-so","dominio":"cosmologia_dados","fonte":"Riess+1998; Perlmutter+1999; Pantheon+ arXiv:2202.04077"},"origem":"wiki","palavras_chave":[],"sinonimos":[],"tipo":"conceito","titulo":"Supernovas Ia — a descoberta da aceleração (Pantheon+)"}
\section{Supernovas Ia — a descoberta da aceleração (Pantheon+)}
SNe Ia são velas padronizáveis. 1998-99: SNe distantes mais fracas que o esperado \ensuremath{\to} expansão ACELERADA \ensuremath{\to} energia escura (Riess+1998; Perlmutter+1999; Nobel 2011). Pantheon+ (2022): 1701 curvas de 1550 SNe, w\ensuremath{_{0}}=\ensuremath{-}0,90$\pm$0,14 (só SNe); H\ensuremath{_{0}}=73,5$\pm$1,1 com Cefeidas. Aceleração = fato consolidado; a causa é o que fica aberto.
\prov{relações: explica de\_sitter\_lambda\_energia\_escura; parte\_de cosmologia\_quantica\_hub}
\prov{proveniência: wiki \textperiodcentered{} Riess+1998; Perlmutter+1999; Pantheon+ arXiv:2202.04077 \textperiodcentered{} fato \textperiodcentered{} confiança 1.0 \textperiodcentered{} domínio cosmologia\_dados \textperiodcentered{} história 3 versões no jornal}

%CRISTAL {"arestas":[{"alvo":"grande_unificacao_gut","confianca":1.0,"epistemico":"fato","origem":"wiki","rel":"especializa"},{"alvo":"su5_georgi_glashow","confianca":1.0,"epistemico":"fato","origem":"wiki","rel":"generaliza"}],"confianca":1.0,"contraexemplos":[],"descricao":"GUTs maiores. SO(10) acomoda uma geração inteira num spinor 16 = 10⊕5̄⊕1 de SU(5), o '1' sendo o ν direito (permite seesaw p/ massas de neutrinos). Pati-Salam SU(4)_c×SU(2)_L×SU(2)_R trata o número leptônico como 'quarta cor' (dentro de SO(10)). Hipóteses vivas — mais flexíveis que SU(5), acomodam os limites atuais; não confirmadas.","epistemico":"hipotese","exemplos":[],"id":"so10_pati_salam","memoria":"semantica","meta":{"autor":"broca-so","dominio":"fisica","fonte":"Fritzsch-Minkowski 1975; Pati-Salam 1974 PRD 10,275"},"origem":"wiki","palavras_chave":[],"sinonimos":[],"tipo":"conceito","titulo":"SO(10) e Pati-Salam"}
\section{SO(10) e Pati-Salam}
GUTs maiores. SO(10) acomoda uma geração inteira num spinor 16 = 10\ensuremath{\oplus}5\ensuremath{}\ensuremath{\oplus}1 de SU(5), o '1' sendo o \ensuremath{\nu} direito (permite seesaw p/ massas de neutrinos). Pati-Salam SU(4)\_c$\times$SU(2)\_L$\times$SU(2)\_R trata o número leptônico como 'quarta cor' (dentro de SO(10)). Hipóteses vivas — mais flexíveis que SU(5), acomodam os limites atuais; não confirmadas.
\prov{relações: especializa grande\_unificacao\_gut; generaliza su5\_georgi\_glashow}
\prov{proveniência: wiki \textperiodcentered{} Fritzsch-Minkowski 1975; Pati-Salam 1974 PRD 10,275 \textperiodcentered{} hipotese \textperiodcentered{} confiança 1.0 \textperiodcentered{} domínio fisica \textperiodcentered{} história 3 versões no jornal}

%CRISTAL {"arestas":[{"alvo":"poligono_vira_poliedro","confianca":1.0,"epistemico":"fato","origem":"wiki","rel":"usa"},{"alvo":"corpo_estelar","confianca":1.0,"epistemico":"fato","origem":"wiki","rel":"usa"},{"alvo":"morfologia_matematica_estelar","confianca":1.0,"epistemico":"fato","origem":"wiki","rel":"usa"},{"alvo":"joalheria_hub","confianca":1.0,"epistemico":"fato","origem":"wiki","rel":"parte_de"}],"confianca":1.0,"contraexemplos":[],"descricao":"A gema como PRODUTO PURO, só multiplicação, SEM SOMA e SEM LIMIAR: a superfície é a SEÇÃO (a mandala k, o contorno) VEZES um PERFIL de revolução, ponto a ponto — x=largura·mandala(θ)·ρ(z)·cosθ, y=…·sinθ, z=altura·z. Trocando o perfil ρ troca-se a forma, contínua (C¹ nos polos), sem máscara: MORANGO/CORAÇÃO ρ=sin u·(1+c·cos u) (topo redondo=elipsoide, bico embaixo); OVO = ELIPSOIDE (topo largo) sobre PARABOLOIDE ou CONE (base fina); e o mais REDONDO: a ESFERA √(1−z²) como PORTADORA, MODULADA pelo fator (1+s·z) — a esfera modulando (s=0 → esfera pura; s pequeno → quase esférico, topo mais largo/base mais fina). 'círculo+cone+trilhão, e as facetas em volta'. a=0 → seção circular (liso, render suave/Gouraud); a>0 → seção trilhão (k=3) e as FACETAS vêm em volta pelo offset meia-faceta (⊗ de La Hire, N=k·m). O elipsoide entra como FATOR do topo, não como limiar. linguagem/reconstrucao_3d.ovo, morango_3d, perfil_ovo; render_estelar.render_quadro(suave=True) para o liso.","epistemico":"fato","exemplos":[],"id":"solidos_multiplicativos","memoria":"semantica","meta":{"dominio":"joalheria","fonte":"maquina de corte"},"origem":"wiki","palavras_chave":[],"sinonimos":[],"tipo":"conceito","titulo":"Sólidos multiplicativos: ovo, morango, trilhão (mandala × perfil, só ⊗)"}
\section{Sólidos multiplicativos: ovo, morango, trilhão (mandala $\times$ perfil, só \ensuremath{\otimes})}
A gema como PRODUTO PURO, só multiplicação, SEM SOMA e SEM LIMIAR: a superfície é a SEÇÃO (a mandala k, o contorno) VEZES um PERFIL de revolução, ponto a ponto — x=largura\textperiodcentered mandala(\ensuremath{\theta})\textperiodcentered \ensuremath{\rho}(z)\textperiodcentered cos\ensuremath{\theta}, y=…\textperiodcentered sin\ensuremath{\theta}, z=altura\textperiodcentered z. Trocando o perfil \ensuremath{\rho} troca-se a forma, contínua (C\ensuremath{^{1}} nos polos), sem máscara: MORANGO/CORAÇÃO \ensuremath{\rho}=sin u\textperiodcentered (1+c\textperiodcentered cos u) (topo redondo=elipsoide, bico embaixo); OVO = ELIPSOIDE (topo largo) sobre PARABOLOIDE ou CONE (base fina); e o mais REDONDO: a ESFERA \ensuremath{\sqrt{\,}}(1\ensuremath{-}z\ensuremath{^{2}}) como PORTADORA, MODULADA pelo fator (1+s\textperiodcentered z) — a esfera modulando (s=0 \ensuremath{\to} esfera pura; s pequeno \ensuremath{\to} quase esférico, topo mais largo/base mais fina). 'círculo+cone+trilhão, e as facetas em volta'. a=0 \ensuremath{\to} seção circular (liso, render suave/Gouraud); a>0 \ensuremath{\to} seção trilhão (k=3) e as FACETAS vêm em volta pelo offset meia-faceta (\ensuremath{\otimes} de La Hire, N=k\textperiodcentered m). O elipsoide entra como FATOR do topo, não como limiar. linguagem/reconstrucao\_3d.ovo, morango\_3d, perfil\_ovo; render\_estelar.render\_quadro(suave=True) para o liso.
\prov{relações: usa poligono\_vira\_poliedro; usa corpo\_estelar; usa morfologia\_matematica\_estelar; parte\_de joalheria\_hub}
\prov{proveniência: wiki \textperiodcentered{} maquina de corte \textperiodcentered{} fato \textperiodcentered{} confiança 1.0 \textperiodcentered{} domínio joalheria \textperiodcentered{} história 200 versões no jornal}

%CRISTAL {"arestas":[{"alvo":"pipeline_stratum","confianca":1.0,"epistemico":"fato","origem":"wiki","rel":"parte_de"},{"alvo":"goldchain","confianca":1.0,"epistemico":"fato","origem":"wiki","rel":"relaciona"}],"confianca":1.0,"contraexemplos":[],"descricao":"Por socket TCP + JSON-lines: mining.subscribe (recebe extranonce1 + extranonce2_size), mining.authorize (user=BTC_ENDERECO.broca-N), dispatch de mining.notify (job) e mining.set_difficulty, e mining.submit de shares com deadline de 10s.","epistemico":"fato","exemplos":[],"id":"stratum_conexao","memoria":"semantica","meta":{"autor":"broca-so","dominio":"metal","fonte":"neuronio/stratum_metal.py:subscribe/authorize/enviar_share"},"origem":"wiki","palavras_chave":[],"sinonimos":[],"tipo":"conceito","titulo":"Handshake Stratum V1 (subscribe/authorize/notify/submit)"}
\section{Handshake Stratum V1 (subscribe/authorize/notify/submit)}
Por socket TCP + JSON-lines: mining.subscribe (recebe extranonce1 + extranonce2\_size), mining.authorize (user=BTC\_ENDERECO.broca-N), dispatch de mining.notify (job) e mining.set\_difficulty, e mining.submit de shares com deadline de 10s.
\prov{relações: parte\_de pipeline\_stratum; relaciona goldchain}
\prov{proveniência: wiki \textperiodcentered{} neuronio/stratum\_metal.py:subscribe/authorize/enviar\_share \textperiodcentered{} fato \textperiodcentered{} confiança 1.0 \textperiodcentered{} domínio metal \textperiodcentered{} história 3 versões no jornal}

%CRISTAL {"arestas":[{"alvo":"cotacao_parabola","confianca":1.0,"epistemico":"fato","origem":"wiki","rel":"usa"},{"alvo":"pipeline_stratum","confianca":1.0,"epistemico":"fato","origem":"wiki","rel":"parte_de"}],"confianca":1.0,"contraexemplos":[],"descricao":"EspacoDifusivo (parábola a+c²=1) define frac=BROCA_DIFUSAO_FRAC (default 0.90), chunk, passos, rodadas, gpu_celulas, c2 — parâmetros que 'difundem' o trabalho e enchem a CPU liberada. O olho não é loop de hash (olho_threads=1).","epistemico":"inferencia","exemplos":[],"id":"stratum_difusao_cpu","memoria":"semantica","meta":{"autor":"broca-so","dominio":"metal","fonte":"neuronio/difusao_recursos.py:EspacoDifusivo"},"origem":"wiki","palavras_chave":[],"sinonimos":[],"tipo":"conceito","titulo":"A Difusão que Preenche a CPU (parábola)"}
\section{A Difusão que Preenche a CPU (parábola)}
EspacoDifusivo (parábola a+c\ensuremath{^{2}}=1) define frac=BROCA\_DIFUSAO\_FRAC (default 0.90), chunk, passos, rodadas, gpu\_celulas, c2 — parâmetros que 'difundem' o trabalho e enchem a CPU liberada. O olho não é loop de hash (olho\_threads=1).
\prov{relações: usa cotacao\_parabola; parte\_de pipeline\_stratum}
\prov{proveniência: wiki \textperiodcentered{} neuronio/difusao\_recursos.py:EspacoDifusivo \textperiodcentered{} inferencia \textperiodcentered{} confiança 1.0 \textperiodcentered{} domínio metal \textperiodcentered{} história 3 versões no jornal}

%CRISTAL {"arestas":[{"alvo":"stratum_loop_broca","confianca":1.0,"epistemico":"fato","origem":"wiki","rel":"relaciona"},{"alvo":"gpu_backend_rede2","confianca":1.0,"epistemico":"fato","origem":"wiki","rel":"relaciona"}],"confianca":1.0,"contraexemplos":[],"descricao":"Kernel CUDA (scan_kernel, 256 threads/bloco, atomicCAS no found) que varre nonces em libstratum_gpu.so, em subprocesso isolado para evitar conflito de contexto CUDA com libbroca. Só o .cu está em disco — a .so NÃO está compilada (caminho disponível, não ativo).","epistemico":"fato","exemplos":[],"id":"stratum_gpu_scan","memoria":"semantica","meta":{"autor":"broca-so","dominio":"metal","fonte":"neuronio/stratum_gpu_scan.cu; stratum_gpu_metal.py"},"origem":"wiki","palavras_chave":[],"sinonimos":[],"tipo":"conceito","titulo":"Varredura SHA256² na GPU (CUDA, não compilada)"}
\section{Varredura SHA256\ensuremath{^{2}} na GPU (CUDA, não compilada)}
Kernel CUDA (scan\_kernel, 256 threads/bloco, atomicCAS no found) que varre nonces em libstratum\_gpu.so, em subprocesso isolado para evitar conflito de contexto CUDA com libbroca. Só o .cu está em disco — a .so NÃO está compilada (caminho disponível, não ativo).
\prov{relações: relaciona stratum\_loop\_broca; relaciona gpu\_backend\_rede2}
\prov{proveniência: wiki \textperiodcentered{} neuronio/stratum\_gpu\_scan.cu; stratum\_gpu\_metal.py \textperiodcentered{} fato \textperiodcentered{} confiança 1.0 \textperiodcentered{} domínio metal \textperiodcentered{} história 3 versões no jornal}

%CRISTAL {"arestas":[{"alvo":"pipeline_stratum","confianca":1.0,"epistemico":"fato","origem":"wiki","rel":"usa"},{"alvo":"pacote_seed_pow","confianca":1.0,"epistemico":"fato","origem":"wiki","rel":"usa"},{"alvo":"stratum_pthreads","confianca":1.0,"epistemico":"fato","origem":"wiki","rel":"usa"}],"confianca":1.0,"contraexemplos":[],"descricao":"Benchmark real (job ckpool, 11 branches): SHA256² 1 thread Python ~89.877 H/s; hash_scan_threads 4 pthreads ~531.773 H/s; ciclo com PoW-cache por job ~1.245.876 H/s (~1,25 MH/s efetivo) vs ~281.008 H/s no modo antigo (PoW por chunk) → speedup 4,43×, 77,4% de economia. PoW seed ~51,9 ms.","epistemico":"fato","exemplos":[],"id":"stratum_hashrate_real","memoria":"semantica","meta":{"autor":"broca-so","dominio":"metal","fonte":"neuronio/benchmark_stratum.py; dados/stratum_benchmark.json"},"origem":"wiki","palavras_chave":[],"sinonimos":[],"tipo":"conceito","titulo":"O Hashrate Medido (~1,25 MH/s, speedup 4,43×)"}
\section{O Hashrate Medido (\~{}1,25 MH/s, speedup 4,43$\times$)}
Benchmark real (job ckpool, 11 branches): SHA256\ensuremath{^{2}} 1 thread Python \~{}89.877 H/s; hash\_scan\_threads 4 pthreads \~{}531.773 H/s; ciclo com PoW-cache por job \~{}1.245.876 H/s (\~{}1,25 MH/s efetivo) vs \~{}281.008 H/s no modo antigo (PoW por chunk) \ensuremath{\to} speedup 4,43$\times$, 77,4\% de economia. PoW seed \~{}51,9 ms.
\prov{relações: usa pipeline\_stratum; usa pacote\_seed\_pow; usa stratum\_pthreads}
\prov{proveniência: wiki \textperiodcentered{} neuronio/benchmark\_stratum.py; dados/stratum\_benchmark.json \textperiodcentered{} fato \textperiodcentered{} confiança 1.0 \textperiodcentered{} domínio metal \textperiodcentered{} história 4 versões no jornal}

%CRISTAL {"arestas":[{"alvo":"pacote_seed_pow","confianca":1.0,"epistemico":"fato","origem":"wiki","rel":"usa"},{"alvo":"stratum_difusao_cpu","confianca":1.0,"epistemico":"fato","origem":"wiki","rel":"usa"},{"alvo":"broca_borda_so_ctypes","confianca":1.0,"epistemico":"fato","origem":"wiki","rel":"usa"}],"confianca":1.0,"contraexemplos":[],"descricao":"Quando libbroca existe, para cada nonce base entrega trabalho ao .so em lanes=gpu_celulas × passos × rodadas (difusão) e verifica o hit via ponte SHA256. Método broca_gpu/broca_cpu. Caminho dev/baseline (BROCA_BORDA_SHA256=1): hash_scan_threads pthreads direto.","epistemico":"inferencia","exemplos":[],"id":"stratum_loop_broca","memoria":"semantica","meta":{"autor":"broca-so","dominio":"metal","fonte":"neuronio/stratum_trabalho.py:loop_trabalho_broca"},"origem":"wiki","palavras_chave":[],"sinonimos":[],"tipo":"conceito","titulo":"Loop de Trabalho (broca CPU/GPU distribui)"}
\section{Loop de Trabalho (broca CPU/GPU distribui)}
Quando libbroca existe, para cada nonce base entrega trabalho ao .so em lanes=gpu\_celulas $\times$ passos $\times$ rodadas (difusão) e verifica o hit via ponte SHA256. Método broca\_gpu/broca\_cpu. Caminho dev/baseline (BROCA\_BORDA\_SHA256=1): hash\_scan\_threads pthreads direto.
\prov{relações: usa pacote\_seed\_pow; usa stratum\_difusao\_cpu; usa broca\_borda\_so\_ctypes}
\prov{proveniência: wiki \textperiodcentered{} neuronio/stratum\_trabalho.py:loop\_trabalho\_broca \textperiodcentered{} inferencia \textperiodcentered{} confiança 1.0 \textperiodcentered{} domínio metal \textperiodcentered{} história 4 versões no jornal}

%CRISTAL {"arestas":[{"alvo":"pipeline_stratum","confianca":1.0,"epistemico":"fato","origem":"wiki","rel":"parte_de"},{"alvo":"libstratum_bloco_so","confianca":1.0,"epistemico":"fato","origem":"wiki","rel":"usa"}],"confianca":1.0,"contraexemplos":[],"descricao":"Monta coinbase=coinb1+extranonce1+extranonce2+coinb2, dobra o merkle root sobre as branches e produz o header de 76 bytes (version|prevhash32|merkle|ntime|nbits, sem nonce). Feito no metal C (libstratum_bloco.so) com fallback Python via dsha256; a montagem C==Python é testada.","epistemico":"fato","exemplos":[],"id":"stratum_montar_bloco","memoria":"semantica","meta":{"autor":"broca-so","dominio":"metal","fonte":"neuronio/stratum_bloco.c:stratum_montar_bloco"},"origem":"wiki","palavras_chave":[],"sinonimos":[],"tipo":"conceito","titulo":"Montagem do Bloco (coinbase→merkle→header76)"}
\section{Montagem do Bloco (coinbase\ensuremath{\to}merkle\ensuremath{\to}header76)}
Monta coinbase=coinb1+extranonce1+extranonce2+coinb2, dobra o merkle root sobre as branches e produz o header de 76 bytes (version|prevhash32|merkle|ntime|nbits, sem nonce). Feito no metal C (libstratum\_bloco.so) com fallback Python via dsha256; a montagem C==Python é testada.
\prov{relações: parte\_de pipeline\_stratum; usa libstratum\_bloco\_so}
\prov{proveniência: wiki \textperiodcentered{} neuronio/stratum\_bloco.c:stratum\_montar\_bloco \textperiodcentered{} fato \textperiodcentered{} confiança 1.0 \textperiodcentered{} domínio metal \textperiodcentered{} história 3 versões no jornal}

%CRISTAL {"arestas":[{"alvo":"cockpit_olho","confianca":1.0,"epistemico":"fato","origem":"wiki","rel":"relaciona"},{"alvo":"olho_de_koch","confianca":1.0,"epistemico":"fato","origem":"wiki","rel":"usa"},{"alvo":"pipeline_stratum","confianca":1.0,"epistemico":"fato","origem":"wiki","rel":"parte_de"}],"confianca":1.0,"contraexemplos":[],"descricao":"CPU que só CONTEMPLA: OlhoKoch audita o canal (parábola c², calor, ok/degenerado) e verificar_nonce_ponte faz a checagem pontual CKPool (share_target/block_target) — 'não é o loop'. O loop de varredura vive no metal; o olho é auditor.","epistemico":"inferencia","exemplos":[],"id":"stratum_olho_koch","memoria":"semantica","meta":{"autor":"broca-so","dominio":"metal","fonte":"neuronio/stratum_olho.py:OlhoKoch"},"origem":"wiki","palavras_chave":[],"sinonimos":[],"tipo":"conceito","titulo":"O Olho de Koch no Stratum (contempla, não minera)"}
\section{O Olho de Koch no Stratum (contempla, não minera)}
CPU que só CONTEMPLA: OlhoKoch audita o canal (parábola c\ensuremath{^{2}}, calor, ok/degenerado) e verificar\_nonce\_ponte faz a checagem pontual CKPool (share\_target/block\_target) — 'não é o loop'. O loop de varredura vive no metal; o olho é auditor.
\prov{relações: relaciona cockpit\_olho; usa olho\_de\_koch; parte\_de pipeline\_stratum}
\prov{proveniência: wiki \textperiodcentered{} neuronio/stratum\_olho.py:OlhoKoch \textperiodcentered{} inferencia \textperiodcentered{} confiança 1.0 \textperiodcentered{} domínio metal \textperiodcentered{} história 4 versões no jornal}

%CRISTAL {"arestas":[{"alvo":"pipeline_stratum","confianca":1.0,"epistemico":"fato","origem":"wiki","rel":"usa"},{"alvo":"nucleo_paralelo","confianca":1.0,"epistemico":"fato","origem":"wiki","rel":"relaciona"}],"confianca":1.0,"contraexemplos":[],"descricao":"Orquestra N workers Stratum num único processo (threads daemon, modo inprocess default) ou N subprocessos (legado). Default -n 1 = máx H/s/worker; -n 5 divide cores (~300 kH/s cada, ~1.2 MH/s total); metal_threads auto = cpu/n. Captura SIGINT, imprime taxa agregada a cada 5s.","epistemico":"fato","exemplos":[],"id":"stratum_orquestrador_workers","memoria":"semantica","meta":{"autor":"broca-so","dominio":"metal","fonte":"neuronio/stratum_orquestrador.py:main"},"origem":"wiki","palavras_chave":[],"sinonimos":[],"tipo":"conceito","titulo":"O Orquestrador (N workers num processo)"}
\section{O Orquestrador (N workers num processo)}
Orquestra N workers Stratum num único processo (threads daemon, modo inprocess default) ou N subprocessos (legado). Default -n 1 = máx H/s/worker; -n 5 divide cores (\~{}300 kH/s cada, \~{}1.2 MH/s total); metal\_threads auto = cpu/n. Captura SIGINT, imprime taxa agregada a cada 5s.
\prov{relações: usa pipeline\_stratum; relaciona nucleo\_paralelo}
\prov{proveniência: wiki \textperiodcentered{} neuronio/stratum\_orquestrador.py:main \textperiodcentered{} fato \textperiodcentered{} confiança 1.0 \textperiodcentered{} domínio metal \textperiodcentered{} história 3 versões no jornal}

%CRISTAL {"arestas":[{"alvo":"libstratum_bloco_so","confianca":1.0,"epistemico":"fato","origem":"wiki","rel":"usa"},{"alvo":"nucleo_paralelo","confianca":1.0,"epistemico":"fato","origem":"wiki","rel":"explica"}],"confianca":1.0,"contraexemplos":[],"descricao":"Pool de pthreads C reutilizado por thread Python via TLS (pthread_key), repartindo max_hashes em fatias por thread com found/cancel sob mutex/cond — o que dá o hashrate multi-thread real do SHA256². Encarnação em C do nucleo_paralelo.","epistemico":"fato","exemplos":[],"id":"stratum_pthreads","memoria":"semantica","meta":{"autor":"broca-so","dominio":"metal","fonte":"neuronio/stratum_bloco.c:broca_pool_run/pool_worker"},"origem":"wiki","palavras_chave":[],"sinonimos":[],"tipo":"conceito","titulo":"Pool de pthreads Persistente (BrocaPool)"}
\section{Pool de pthreads Persistente (BrocaPool)}
Pool de pthreads C reutilizado por thread Python via TLS (pthread\_key), repartindo max\_hashes em fatias por thread com found/cancel sob mutex/cond — o que dá o hashrate multi-thread real do SHA256\ensuremath{^{2}}. Encarnação em C do nucleo\_paralelo.
\prov{relações: usa libstratum\_bloco\_so; explica nucleo\_paralelo}
\prov{proveniência: wiki \textperiodcentered{} neuronio/stratum\_bloco.c:broca\_pool\_run/pool\_worker \textperiodcentered{} fato \textperiodcentered{} confiança 1.0 \textperiodcentered{} domínio metal \textperiodcentered{} história 3 versões no jornal}

%CRISTAL {"arestas":[{"alvo":"grande_unificacao_gut","confianca":1.0,"epistemico":"fato","origem":"wiki","rel":"especializa"}],"confianca":1.0,"contraexemplos":[],"descricao":"A primeira GUT; menor grupo simples que contém o MP. Cada geração em 5̄⊕10 (15 estados). Prevê bósons X/Y que mediam decaimento do próton (~10³⁰-10³¹ anos). DESFAVORECIDA: o SU(5) mínimo (não-SUSY) está efetivamente excluído pelo não-decaimento do próton observado.","epistemico":"hipotese","exemplos":[],"id":"su5_georgi_glashow","memoria":"semantica","meta":{"autor":"broca-so","dominio":"fisica","fonte":"Georgi-Glashow 1974 PRL 32,438"},"origem":"wiki","palavras_chave":[],"sinonimos":[],"tipo":"conceito","titulo":"SU(5) de Georgi-Glashow"}
\section{SU(5) de Georgi-Glashow}
A primeira GUT; menor grupo simples que contém o MP. Cada geração em 5\ensuremath{}\ensuremath{\oplus}10 (15 estados). Prevê bósons X/Y que mediam decaimento do próton (\~{}10\ensuremath{^{30}}-10\ensuremath{^{31}} anos). DESFAVORECIDA: o SU(5) mínimo (não-SUSY) está efetivamente excluído pelo não-decaimento do próton observado.
\prov{relações: especializa grande\_unificacao\_gut}
\prov{proveniência: wiki \textperiodcentered{} Georgi-Glashow 1974 PRL 32,438 \textperiodcentered{} hipotese \textperiodcentered{} confiança 1.0 \textperiodcentered{} domínio fisica \textperiodcentered{} história 2 versões no jornal}

%CRISTAL {"arestas":[{"alvo":"dois_qubits_octonios","confianca":1.0,"epistemico":"fato","origem":"wiki","rel":"usa"},{"alvo":"limite_hurwitz_s1","confianca":1.0,"epistemico":"fato","origem":"wiki","rel":"usa"}],"confianca":1.0,"contraexemplos":[],"descricao":"No quaternion vale a+c²=1 → teto de Bell = 2√2 exato (Tsirelson). Nos octônios há um piso a≥a_min>0 na não-associatividade que BAIXA o máximo: S_s<2√2. Predição falseável PRÓPRIA do programa (nível 3), invisível ao formalismo de Hilbert padrão — um desvio abaixo do teto quântico assinaria a não-associatividade. [N] hipótese falseável.","epistemico":"hipotese","exemplos":[],"id":"sub_tsirelson_octonico","memoria":"semantica","meta":{"autor":"broca-so","dominio":"cosmologia","fonte":"paper_3_quantica.tex §Bell/Tsirelson"},"origem":"wiki","palavras_chave":[],"sinonimos":[],"tipo":"conceito","titulo":"Predição própria: piso sub-Tsirelson S<2√2"}
\section{Predição própria: piso sub-Tsirelson S<2\ensuremath{\sqrt{\,}}2}
No quaternion vale a+c\ensuremath{^{2}}=1 \ensuremath{\to} teto de Bell = 2\ensuremath{\sqrt{\,}}2 exato (Tsirelson). Nos octônios há um piso a\ensuremath{\ge}a\_min>0 na não-associatividade que BAIXA o máximo: S\_s<2\ensuremath{\sqrt{\,}}2. Predição falseável PRÓPRIA do programa (nível 3), invisível ao formalismo de Hilbert padrão — um desvio abaixo do teto quântico assinaria a não-associatividade. [N] hipótese falseável.
\prov{relações: usa dois\_qubits\_octonios; usa limite\_hurwitz\_s1}
\prov{proveniência: wiki \textperiodcentered{} paper\_3\_quantica.tex §Bell/Tsirelson \textperiodcentered{} hipotese \textperiodcentered{} confiança 1.0 \textperiodcentered{} domínio cosmologia \textperiodcentered{} história 3 versões no jornal}

%CRISTAL {"arestas":[{"alvo":"associador_nao_tensoriza","confianca":1.0,"epistemico":"fato","origem":"wiki","rel":"usa"},{"alvo":"sub_tsirelson_octonico","confianca":1.0,"epistemico":"fato","origem":"wiki","rel":"relaciona"}],"confianca":1.0,"contraexemplos":[],"descricao":"Máquinas de qubits são associativas → saturam o bound de Mermin |M|=4 (confirmado no ibm_fez). A mesma estrutura algébrica no seu espaço pleno de botões alcança |M|≈22 (5,5× o teto quântico) — inatingível por hardware de qubits; um computador que implementasse o associador nativamente o mediria. [I] questão aberta: realizável fisicamente ou sombra matemática do 'oceano fora da ilha'. Distinto do sub-Tsirelson (que é o octônio físico).","epistemico":"hipotese","exemplos":[],"id":"super_tsirelson_22","memoria":"semantica","meta":{"autor":"broca-so","dominio":"cosmologia","fonte":"paper_3_quantica.tex obs:sangria"},"origem":"wiki","palavras_chave":[],"sinonimos":[],"tipo":"conceito","titulo":"Super-Tsirelson |M|≈22 (fora da ilha)"}
\section{Super-Tsirelson |M|\ensuremath{\approx}22 (fora da ilha)}
Máquinas de qubits são associativas \ensuremath{\to} saturam o bound de Mermin |M|=4 (confirmado no ibm\_fez). A mesma estrutura algébrica no seu espaço pleno de botões alcança |M|\ensuremath{\approx}22 (5,5$\times$ o teto quântico) — inatingível por hardware de qubits; um computador que implementasse o associador nativamente o mediria. [I] questão aberta: realizável fisicamente ou sombra matemática do 'oceano fora da ilha'. Distinto do sub-Tsirelson (que é o octônio físico).
\prov{relações: usa associador\_nao\_tensoriza; relaciona sub\_tsirelson\_octonico}
\prov{proveniência: wiki \textperiodcentered{} paper\_3\_quantica.tex obs:sangria \textperiodcentered{} hipotese \textperiodcentered{} confiança 1.0 \textperiodcentered{} domínio cosmologia \textperiodcentered{} história 3 versões no jornal}

%CRISTAL {"arestas":[{"alvo":"problema_hierarquia_naturalidade","confianca":1.0,"epistemico":"fato","origem":"wiki","rel":"usa"},{"alvo":"grande_unificacao_gut","confianca":1.0,"epistemico":"fato","origem":"wiki","rel":"usa"},{"alvo":"materia_escura","confianca":1.0,"epistemico":"fato","origem":"wiki","rel":"relaciona"},{"alvo":"teoria_cordas_m","confianca":1.0,"epistemico":"fato","origem":"wiki","rel":"relaciona"},{"alvo":"cosmologia_quantica_hub","confianca":1.0,"epistemico":"fato","origem":"wiki","rel":"parte_de"}],"confianca":1.0,"contraexemplos":[],"descricao":"Simetria entre bósons e férmions: cada partícula ganha um superparceiro (squarks, sleptons; gauginos, higgsinos), spin diferindo de ½; o MSSM é a realização mínima. Resolveria o problema da hierarquia (loops de partícula e superparceiro se cancelam), unificaria os 3 acoplamentos em M_GUT~2×10¹⁶ GeV, e daria matéria escura (LSP/neutralino, se R-paridade). Hipótese NÃO confirmada — zero superparceiros em ~15 anos de LHC.","epistemico":"hipotese","exemplos":[],"id":"supersimetria_susy","memoria":"semantica","meta":{"autor":"broca-so","dominio":"fisica","fonte":"Wess-Zumino 1974; PDG SUSY Review 2024/25"},"origem":"wiki","palavras_chave":[],"sinonimos":[],"tipo":"conceito","titulo":"Supersimetria (SUSY / MSSM)"}
\section{Supersimetria (SUSY / MSSM)}
Simetria entre bósons e férmions: cada partícula ganha um superparceiro (squarks, sleptons; gauginos, higgsinos), spin diferindo de \ensuremath{\tfrac12}; o MSSM é a realização mínima. Resolveria o problema da hierarquia (loops de partícula e superparceiro se cancelam), unificaria os 3 acoplamentos em M\_GUT\~{}2$\times$10\ensuremath{^{16}} GeV, e daria matéria escura (LSP/neutralino, se R-paridade). Hipótese NÃO confirmada — zero superparceiros em \~{}15 anos de LHC.
\prov{relações: usa problema\_hierarquia\_naturalidade; usa grande\_unificacao\_gut; relaciona materia\_escura; relaciona teoria\_cordas\_m; parte\_de cosmologia\_quantica\_hub}
\prov{proveniência: wiki \textperiodcentered{} Wess-Zumino 1974; PDG SUSY Review 2024/25 \textperiodcentered{} hipotese \textperiodcentered{} confiança 1.0 \textperiodcentered{} domínio fisica \textperiodcentered{} história 6 versões no jornal}

%CRISTAL {"arestas":[{"alvo":"supersimetria_susy","confianca":1.0,"epistemico":"fato","origem":"wiki","rel":"relaciona"},{"alvo":"naturalidade_em_crise","confianca":1.0,"epistemico":"fato","origem":"wiki","rel":"usa"}],"confianca":1.0,"contraexemplos":[],"descricao":"Nenhum superparceiro achado: gluinos/squarks>2,2-2,4 TeV, stops>1,3 TeV, eletrofracos>~1 TeV (modelos simplificados; no pMSSM sobrevivem cenários comprimidos). Os limites no TeV reintroduzem ajuste fino (~0,1-1%) — a SUSY NATURAL como solução da hierarquia perdeu a motivação. Refúgios: split SUSY / alta escala (abandonam naturalidade, mantêm unificação+DM), justificados por antrópico/paisagem. A unificação MSSM segue como argumento indireto vivo.","epistemico":"hipotese","exemplos":[],"id":"susy_lhc_pressao_split","memoria":"semantica","meta":{"autor":"broca-so","dominio":"fisica","fonte":"ATLAS/CMS Run 2; arXiv:2505.11251 (2025); Arkani-Hamed-Dimopoulos 2004"},"origem":"wiki","palavras_chave":[],"sinonimos":[],"tipo":"conceito","titulo":"SUSY sob pressão no LHC; naturalidade em crise; split SUSY"}
\section{SUSY sob pressão no LHC; naturalidade em crise; split SUSY}
Nenhum superparceiro achado: gluinos/squarks>2,2-2,4 TeV, stops>1,3 TeV, eletrofracos>\~{}1 TeV (modelos simplificados; no pMSSM sobrevivem cenários comprimidos). Os limites no TeV reintroduzem ajuste fino (\~{}0,1-1\%) — a SUSY NATURAL como solução da hierarquia perdeu a motivação. Refúgios: split SUSY / alta escala (abandonam naturalidade, mantêm unificação+DM), justificados por antrópico/paisagem. A unificação MSSM segue como argumento indireto vivo.
\prov{relações: relaciona supersimetria\_susy; usa naturalidade\_em\_crise}
\prov{proveniência: wiki \textperiodcentered{} ATLAS/CMS Run 2; arXiv:2505.11251 (2025); Arkani-Hamed-Dimopoulos 2004 \textperiodcentered{} hipotese \textperiodcentered{} confiança 1.0 \textperiodcentered{} domínio fisica \textperiodcentered{} história 3 versões no jornal}

%CRISTAL {"arestas":[{"alvo":"symbolic_beta_pipe_so","confianca":1.0,"epistemico":"fato","origem":"wiki","rel":"confirma"},{"alvo":"gex_minera_ouro_real","confianca":1.0,"epistemico":"fato","origem":"wiki","rel":"confirma"},{"alvo":"dimensao_de_prata_pell","confianca":1.0,"epistemico":"fato","origem":"wiki","rel":"usa"},{"alvo":"espectro_prata_seis_modos","confianca":1.0,"epistemico":"fato","origem":"wiki","rel":"usa"},{"alvo":"goldchain","confianca":1.0,"epistemico":"fato","origem":"wiki","rel":"parte_de"}],"confianca":1.0,"contraexemplos":[],"descricao":"O pipe rodou AO VIVO contra a pool BTC real (solo.ckpool.org:3333, worker bc1pxmwk…broca-miner, >1,2 MH/s): capturou 10 jobs reais, cunhou 10 FLOCOS DE OURO (pacotes seed, c²≈0,99, impressão Koch real) e, na hora e em paralelo, emitiu 10 PELL irmãos — a sequência de Pell exata 2, 5, 12, 29, 70, 169, 408, 985, 2378, 5741. Cada floco foi lido no espectro da PRATA em 6 modos (hexagrama de Salomão) e cifrado em 𝔽₂¹²⁸ (ouro branco), com round-trip válido (legível=True). Ledger append-only persistido em goldchain/dados/ledger_symbolic.json. É a prova viva da liquidez total: minerou ouro e registrou a prata junto, floco a floco — cada partícula de ouro ganhou sua prata irmã. FATO (rodou; commit 5025a96). A gex44 ficou intocada (processo separado).","epistemico":"fato","exemplos":[],"id":"symbolic_beta_corrida_live","memoria":"semantica","meta":{"autor":"Lopes/broca-so","dominio":"mineracao","fonte":"goldchain/dados/ledger_symbolic.json (10 Pell); solo.ckpool.org live 2 jul 2026"},"origem":"wiki","palavras_chave":[],"sinonimos":[],"tipo":"conceito","titulo":"A corrida live: 10 flocos de ouro reais (solo.ckpool) → 10 Pell irmãos, liquidez total confirmada"}
\section{A corrida live: 10 flocos de ouro reais (solo.ckpool) \ensuremath{\to} 10 Pell irmãos, liquidez total confirmada}
O pipe rodou AO VIVO contra a pool BTC real (solo.ckpool.org:3333, worker bc1pxmwk…broca-miner, >1,2 MH/s): capturou 10 jobs reais, cunhou 10 FLOCOS DE OURO (pacotes seed, c\ensuremath{^{2}}\ensuremath{\approx}0,99, impressão Koch real) e, na hora e em paralelo, emitiu 10 PELL irmãos — a sequência de Pell exata 2, 5, 12, 29, 70, 169, 408, 985, 2378, 5741. Cada floco foi lido no espectro da PRATA em 6 modos (hexagrama de Salomão) e cifrado em \ensuremath{\mathbb{F}}\ensuremath{_{2}}\ensuremath{^{128}} (ouro branco), com round-trip válido (legível=True). Ledger append-only persistido em goldchain/dados/ledger\_symbolic.json. É a prova viva da liquidez total: minerou ouro e registrou a prata junto, floco a floco — cada partícula de ouro ganhou sua prata irmã. FATO (rodou; commit 5025a96). A gex44 ficou intocada (processo separado).
\prov{relações: confirma symbolic\_beta\_pipe\_so; confirma gex\_minera\_ouro\_real; usa dimensao\_de\_prata\_pell; usa espectro\_prata\_seis\_modos; parte\_de goldchain}
\prov{proveniência: wiki \textperiodcentered{} goldchain/dados/ledger\_symbolic.json (10 Pell); solo.ckpool.org live 2 jul 2026 \textperiodcentered{} fato \textperiodcentered{} confiança 1.0 \textperiodcentered{} domínio mineracao \textperiodcentered{} história 18 versões no jornal}

%CRISTAL {"arestas":[{"alvo":"symbolic_beta_pipe_so","confianca":1.0,"epistemico":"fato","origem":"wiki","rel":"confirma"},{"alvo":"symbolic_beta_corrida_live","confianca":1.0,"epistemico":"fato","origem":"wiki","rel":"relaciona"},{"alvo":"dimensao_de_prata_pell","confianca":1.0,"epistemico":"fato","origem":"wiki","rel":"usa"},{"alvo":"corpo_f2k","confianca":1.0,"epistemico":"fato","origem":"wiki","rel":"usa"},{"alvo":"goldchain","confianca":1.0,"epistemico":"fato","origem":"wiki","rel":"parte_de"},{"alvo":"beta_numeracao_metalica","confianca":1.0,"epistemico":"fato","origem":"wiki","rel":"aplicacao"}],"confianca":1.0,"contraexemplos":[],"descricao":"O gex44 (hipercomplexo, gex44, up 54 dias) minerava flocos de ouro há ~1,6 dias — 5702 pacotes seed (ν=6, c²≈0,99) no log neuronio/dados/stratum_gex44.log — SEM nunca lastrear (não havia symbolic.py lá). Migrado e recuperado: (1) 'symbolic.py recuperar --log' lastreou em LOTE os 5702 flocos históricos em Pell (1 lock, 1 gravação, usando o ν real do log); (2) o sistema foi reiniciado — a sessão tmux broca-stratum agora roda 'symbolic.py live' (~3,4 MH/s), lastreando cada floco NOVO ao vivo (Pell #5703, #5704…, append-only), com a difusão broca-olho intocada. NOTA HONESTA: a cifra de ouro branco 𝔽₂¹²⁸ só faz round-trip para Pell < 2¹²⁸ (~os primeiros ~100); acima disso o índice e o valor do Pell ficam registrados, mas legível=False (o corpo tem 128 bits; o Pell cresce ~σ₂ⁿ). É a liquidez recuperada + a liquidez futura, no metal. FATO (rodou).","epistemico":"fato","exemplos":[],"id":"symbolic_beta_gex44","memoria":"semantica","meta":{"autor":"Lopes/broca-so","dominio":"mineracao","fonte":"gex44:/home/claude/broca-so/goldchain/dados/ledger_symbolic.json (5704 Pell); tmux broca-stratum live"},"origem":"wiki","palavras_chave":[],"sinonimos":[],"tipo":"conceito","titulo":"O gex44 migrado: 5702 flocos de ouro de dias recuperados em Pell + lastro ao vivo"}
\section{O gex44 migrado: 5702 flocos de ouro de dias recuperados em Pell + lastro ao vivo}
O gex44 (hipercomplexo, gex44, up 54 dias) minerava flocos de ouro há \~{}1,6 dias — 5702 pacotes seed (\ensuremath{\nu}=6, c\ensuremath{^{2}}\ensuremath{\approx}0,99) no log neuronio/dados/stratum\_gex44.log — SEM nunca lastrear (não havia symbolic.py lá). Migrado e recuperado: (1) 'symbolic.py recuperar --log' lastreou em LOTE os 5702 flocos históricos em Pell (1 lock, 1 gravação, usando o \ensuremath{\nu} real do log); (2) o sistema foi reiniciado — a sessão tmux broca-stratum agora roda 'symbolic.py live' (\~{}3,4 MH/s), lastreando cada floco NOVO ao vivo (Pell \#5703, \#5704…, append-only), com a difusão broca-olho intocada. NOTA HONESTA: a cifra de ouro branco \ensuremath{\mathbb{F}}\ensuremath{_{2}}\ensuremath{^{128}} só faz round-trip para Pell < 2\ensuremath{^{128}} (\~{}os primeiros \~{}100); acima disso o índice e o valor do Pell ficam registrados, mas legível=False (o corpo tem 128 bits; o Pell cresce \~{}\ensuremath{\sigma}\ensuremath{_{2}}\ensuremath{^{n}}). É a liquidez recuperada + a liquidez futura, no metal. FATO (rodou).
\prov{relações: confirma symbolic\_beta\_pipe\_so; relaciona symbolic\_beta\_corrida\_live; usa dimensao\_de\_prata\_pell; usa corpo\_f2k; parte\_de goldchain; aplicacao beta\_numeracao\_metalica}
\prov{proveniência: wiki \textperiodcentered{} gex44:/home/claude/broca-so/goldchain/dados/ledger\_symbolic.json (5704 Pell); tmux broca-stratum live \textperiodcentered{} fato \textperiodcentered{} confiança 1.0 \textperiodcentered{} domínio mineracao \textperiodcentered{} história 13 versões no jornal}

%CRISTAL {"arestas":[{"alvo":"corpo_estelar","confianca":1.0,"epistemico":"fato","origem":"wiki","rel":"usa"},{"alvo":"joalheria_hub","confianca":1.0,"epistemico":"fato","origem":"wiki","rel":"parte_de"}],"confianca":1.0,"contraexemplos":[],"descricao":"O catálogo das geometrias de corte, indexado pela SIMETRIA k (a estrela do corpo estelar, ν=k−1, a frequência da borda). NÃO é só um catálogo: é um BANCO DE CONHECIMENTO (sugestão do Cursor) — cada geometria tem nome, grupo de simetria, assinatura espectral, curvatura média, operador morfológico, estratégia de lapidação, parâmetros óticos, certificação e render ideal. Toda a cadeia (segmentação/joalheria/render/certificação/mineração) consulta AQUI a mesma definição. CRESCE no grafo: cada nova geometria e cada nova gema entram como uma linha. linguagem/fabricacao_gema.tabela_geometrias.","epistemico":"fato","exemplos":[],"id":"tabela_geometrias_espectrais","memoria":"semantica","meta":{"dominio":"joalheria","fonte":"maquina de corte"},"origem":"wiki","palavras_chave":[],"sinonimos":[],"tipo":"conceito","titulo":"A Tabela de Geometrias Espectrais (o banco de conhecimento vivo)"}
\section{A Tabela de Geometrias Espectrais (o banco de conhecimento vivo)}
O catálogo das geometrias de corte, indexado pela SIMETRIA k (a estrela do corpo estelar, \ensuremath{\nu}=k\ensuremath{-}1, a frequência da borda). NÃO é só um catálogo: é um BANCO DE CONHECIMENTO (sugestão do Cursor) — cada geometria tem nome, grupo de simetria, assinatura espectral, curvatura média, operador morfológico, estratégia de lapidação, parâmetros óticos, certificação e render ideal. Toda a cadeia (segmentação/joalheria/render/certificação/mineração) consulta AQUI a mesma definição. CRESCE no grafo: cada nova geometria e cada nova gema entram como uma linha. linguagem/fabricacao\_gema.tabela\_geometrias.
\prov{relações: usa corpo\_estelar; parte\_de joalheria\_hub}
\prov{proveniência: wiki \textperiodcentered{} maquina de corte \textperiodcentered{} fato \textperiodcentered{} confiança 1.0 \textperiodcentered{} domínio joalheria \textperiodcentered{} história 153 versões no jornal}

%CRISTAL {"arestas":[{"alvo":"polimento_cristais","confianca":1.0,"epistemico":"fato","origem":"wiki","rel":"especializa"},{"alvo":"ultra_realismo_gema","confianca":1.0,"epistemico":"fato","origem":"wiki","rel":"usa"},{"alvo":"voxelizador_cristalografico","confianca":1.0,"epistemico":"fato","origem":"wiki","rel":"usa"},{"alvo":"joalheria_hub","confianca":1.0,"epistemico":"fato","origem":"wiki","rel":"parte_de"}],"confianca":1.0,"contraexemplos":[],"descricao":"O polimento AUTOMATIZADO (o tambor rotativo), unificado com a óptica: AS ITERAÇÕES SÃO OS RICOCHETES. Parte da estrutura VOXELIZADA (rugosa) e cada RICOCHETE (a interação na interface — Snell/Fresnel/TIR) POLI um pouco: a rugosidade converge à GEOMETRIA LISA selecionada (o dupolinômio), rug_i=rug0·grão^i→0, e as PROPRIEDADES FÍSICAS dentro da gema refinam (a cor satura, a brilhância junta luz). Mais ricochetes = mais polido. É a granulometria de Matheron indexada pelo ricochete (cada colisão um grão mais fino). A óptica SEGUE a rugosidade: rugoso→normais perturbadas→DIFUSO (fosco); liso→normais nítidas→ESPECULAR (o brilho). O render mostra: voxelizado (fosco) → desbaste → polimento → lustro (brilhante), pelos ricochetes. linguagem/polimento.py (6/6: tamboreamento, optica_da_rugosidade); render_gema_ultra(rugosidade=…).","epistemico":"fato","exemplos":[],"id":"tamboreamento_ricochetes","memoria":"semantica","meta":{"dominio":"joalheria","fonte":"maquina de corte"},"origem":"wiki","palavras_chave":[],"sinonimos":[],"tipo":"conceito","titulo":"Tamboreamento: as iterações SÃO os ricochetes (o voxelizado vira liso polindo)"}
\section{Tamboreamento: as iterações SÃO os ricochetes (o voxelizado vira liso polindo)}
O polimento AUTOMATIZADO (o tambor rotativo), unificado com a óptica: AS ITERAÇÕES SÃO OS RICOCHETES. Parte da estrutura VOXELIZADA (rugosa) e cada RICOCHETE (a interação na interface — Snell/Fresnel/TIR) POLI um pouco: a rugosidade converge à GEOMETRIA LISA selecionada (o dupolinômio), rug\_i=rug0\textperiodcentered grão\^{}i\ensuremath{\to}0, e as PROPRIEDADES FÍSICAS dentro da gema refinam (a cor satura, a brilhância junta luz). Mais ricochetes = mais polido. É a granulometria de Matheron indexada pelo ricochete (cada colisão um grão mais fino). A óptica SEGUE a rugosidade: rugoso\ensuremath{\to}normais perturbadas\ensuremath{\to}DIFUSO (fosco); liso\ensuremath{\to}normais nítidas\ensuremath{\to}ESPECULAR (o brilho). O render mostra: voxelizado (fosco) \ensuremath{\to} desbaste \ensuremath{\to} polimento \ensuremath{\to} lustro (brilhante), pelos ricochetes. linguagem/polimento.py (6/6: tamboreamento, optica\_da\_rugosidade); render\_gema\_ultra(rugosidade=…).
\prov{relações: especializa polimento\_cristais; usa ultra\_realismo\_gema; usa voxelizador\_cristalografico; parte\_de joalheria\_hub}
\prov{proveniência: wiki \textperiodcentered{} maquina de corte \textperiodcentered{} fato \textperiodcentered{} confiança 1.0 \textperiodcentered{} domínio joalheria \textperiodcentered{} história 115 versões no jornal}

%CRISTAL {"arestas":[{"alvo":"organismo_gex44","confianca":1.0,"epistemico":"fato","origem":"wiki","rel":"usa"},{"alvo":"olho_de_koch","confianca":1.0,"epistemico":"fato","origem":"wiki","rel":"usa"}],"confianca":1.0,"contraexemplos":[],"descricao":"tecido_gex.py lê o stratum_gex44.log bit literal, fecha cada bit num frame (Φ‖bump, 17B), faz append em 'floco' e atualiza gato.state (e,o,bit_idx,ref) sob flock; init_gato.py sincroniza o gato por replay de bits garantindo o invariante bit×17=bytes(floco). O tecido difusivo do organismo GEX44.","epistemico":"fato","exemplos":[],"id":"tecido_gex_floco","memoria":"semantica","meta":{"autor":"broca-so","dominio":"metal","fonte":"scripts/gex44/tecido_gex.py; init_gato.py"},"origem":"wiki","palavras_chave":[],"sinonimos":[],"tipo":"conceito","titulo":"O Tecido no Servidor (log → floco Φ‖bump)"}
\section{O Tecido no Servidor (log \ensuremath{\to} floco \ensuremath{\Phi}\ensuremath{\|}bump)}
tecido\_gex.py lê o stratum\_gex44.log bit literal, fecha cada bit num frame (\ensuremath{\Phi}\ensuremath{\|}bump, 17B), faz append em 'floco' e atualiza gato.state (e,o,bit\_idx,ref) sob flock; init\_gato.py sincroniza o gato por replay de bits garantindo o invariante bit$\times$17=bytes(floco). O tecido difusivo do organismo GEX44.
\prov{relações: usa organismo\_gex44; usa olho\_de\_koch}
\prov{proveniência: wiki \textperiodcentered{} scripts/gex44/tecido\_gex.py; init\_gato.py \textperiodcentered{} fato \textperiodcentered{} confiança 1.0 \textperiodcentered{} domínio metal \textperiodcentered{} história 3 versões no jornal}

%CRISTAL {"arestas":[{"alvo":"morfologia_matematica_estelar","confianca":1.0,"epistemico":"fato","origem":"wiki","rel":"relaciona"},{"alvo":"joalheria_hub","confianca":1.0,"epistemico":"fato","origem":"wiki","rel":"parte_de"}],"confianca":1.0,"contraexemplos":[],"descricao":"As técnicas modernas de morfologia matemática do prof. Ronaldo F. Hashimoto (IME-USP): a máquina morfológica por BINARY DECISION DIAGRAM (BDD) É a árvore-de-decisão-com-orçamento — a própria BAI. A U-curve (seleção de features) é a poda, a decomposição canônica (sup de erosões) é o colapso, a decomposição sequencial (a cadeia A⊕B₁⊕…) é a propagação, o reticulado de operadores é a álgebra de extents. Compor com bai_decisao dá a tradução (o Teorema). Refs: vision.ime.usp.br/~ronaldo. Entram na joalheria como a coluna TÉCNICAS. dicionarios.morfologia_moderna, hashimoto_maquina.","epistemico":"fato","exemplos":[],"id":"tecnicas_morfologia_hashimoto","memoria":"semantica","meta":{"dominio":"joalheria","fonte":"maquina de corte"},"origem":"wiki","palavras_chave":[],"sinonimos":[],"tipo":"conceito","titulo":"As Técnicas Modernas (Ronaldo F. Hashimoto, IME-USP) ↔ BAI"}
\section{As Técnicas Modernas (Ronaldo F. Hashimoto, IME-USP) \ensuremath{\leftrightarrow} BAI}
As técnicas modernas de morfologia matemática do prof. Ronaldo F. Hashimoto (IME-USP): a máquina morfológica por BINARY DECISION DIAGRAM (BDD) É a árvore-de-decisão-com-orçamento — a própria BAI. A U-curve (seleção de features) é a poda, a decomposição canônica (sup de erosões) é o colapso, a decomposição sequencial (a cadeia A\ensuremath{\oplus}B\ensuremath{_{1}}\ensuremath{\oplus}…) é a propagação, o reticulado de operadores é a álgebra de extents. Compor com bai\_decisao dá a tradução (o Teorema). Refs: vision.ime.usp.br/\~{}ronaldo. Entram na joalheria como a coluna TÉCNICAS. dicionarios.morfologia\_moderna, hashimoto\_maquina.
\prov{relações: relaciona morfologia\_matematica\_estelar; parte\_de joalheria\_hub}
\prov{proveniência: wiki \textperiodcentered{} maquina de corte \textperiodcentered{} fato \textperiodcentered{} confiança 1.0 \textperiodcentered{} domínio joalheria \textperiodcentered{} história 153 versões no jornal}

%CRISTAL {"arestas":[{"alvo":"info_quantica_fundamentos","confianca":1.0,"epistemico":"fato","origem":"wiki","rel":"usa"},{"alvo":"teorema_nao_clonagem","confianca":1.0,"epistemico":"fato","origem":"wiki","rel":"relaciona"}],"confianca":1.0,"contraexemplos":[],"descricao":"TELEPORTE: transfere um estado desconhecido de A para B consumindo 1 par emaranhado (ebit) + 2 bits clássicos, sem enviar o sistema físico (não viola não-clonagem — o original é destruído — nem causalidade). SUPERDENSO: 1 ebit + 1 qubit transmitem 2 bits clássicos (o dual). Bennett+1993; Bennett-Wiesner 1992. Fato (demonstrado desde 1997).","epistemico":"fato","exemplos":[],"id":"teleporte_superdenso","memoria":"semantica","meta":{"autor":"broca-so","dominio":"fisica","fonte":"Bennett+1993 PRL 70,1895; Bennett-Wiesner 1992"},"origem":"wiki","palavras_chave":[],"sinonimos":[],"tipo":"conceito","titulo":"Teleporte quântico e codificação superdensa"}
\section{Teleporte quântico e codificação superdensa}
TELEPORTE: transfere um estado desconhecido de A para B consumindo 1 par emaranhado (ebit) + 2 bits clássicos, sem enviar o sistema físico (não viola não-clonagem — o original é destruído — nem causalidade). SUPERDENSO: 1 ebit + 1 qubit transmitem 2 bits clássicos (o dual). Bennett+1993; Bennett-Wiesner 1992. Fato (demonstrado desde 1997).
\prov{relações: usa info\_quantica\_fundamentos; relaciona teorema\_nao\_clonagem}
\prov{proveniência: wiki \textperiodcentered{} Bennett+1993 PRL 70,1895; Bennett-Wiesner 1992 \textperiodcentered{} fato \textperiodcentered{} confiança 1.0 \textperiodcentered{} domínio fisica \textperiodcentered{} história 3 versões no jornal}

%CRISTAL {"arestas":[{"alvo":"isolantes_topologicos_z2","confianca":1.0,"epistemico":"fato","origem":"wiki","rel":"generaliza"},{"alvo":"numero_chern_tknn","confianca":1.0,"epistemico":"fato","origem":"wiki","rel":"usa"}],"confianca":1.0,"contraexemplos":[],"descricao":"Um invariante topológico não-nulo no bulk IMPLICA estados de contorno protegidos sem gap (bulk-boundary correspondence). Fases protegidas por simetria (SPT) são triviais no bulk mas distintas enquanto a simetria vale. A classificação de Altland-Zirnbauer organiza fases de férmions livres em 10 CLASSES (por T, C, S), com periodicidade de Bott. Arcabouço estabelecido (interações podem colapsar ℤ→ℤ₈).","epistemico":"fato","exemplos":[],"id":"tenfold_way_bulk_boundary","memoria":"semantica","meta":{"autor":"broca-so","dominio":"fisica","fonte":"Altland-Zirnbauer 1997; Schnyder+2008; Kitaev 2009"},"origem":"wiki","palavras_chave":[],"sinonimos":[],"tipo":"conceito","titulo":"Correspondência bulk-boundary e a tenfold way"}
\section{Correspondência bulk-boundary e a tenfold way}
Um invariante topológico não-nulo no bulk IMPLICA estados de contorno protegidos sem gap (bulk-boundary correspondence). Fases protegidas por simetria (SPT) são triviais no bulk mas distintas enquanto a simetria vale. A classificação de Altland-Zirnbauer organiza fases de férmions livres em 10 CLASSES (por T, C, S), com periodicidade de Bott. Arcabouço estabelecido (interações podem colapsar \ensuremath{\mathbb{Z}}\ensuremath{\to}\ensuremath{\mathbb{Z}}\ensuremath{_{8}}).
\prov{relações: generaliza isolantes\_topologicos\_z2; usa numero\_chern\_tknn}
\prov{proveniência: wiki \textperiodcentered{} Altland-Zirnbauer 1997; Schnyder+2008; Kitaev 2009 \textperiodcentered{} fato \textperiodcentered{} confiança 1.0 \textperiodcentered{} domínio fisica \textperiodcentered{} história 3 versões no jornal}

%CRISTAL {"arestas":[{"alvo":"hubble_lemaitre","confianca":1.0,"epistemico":"fato","origem":"wiki","rel":"relaciona"},{"alvo":"lambda_cdm","confianca":1.0,"epistemico":"fato","origem":"wiki","rel":"relaciona"},{"alvo":"cosmologia_quantica_hub","confianca":1.0,"epistemico":"fato","origem":"wiki","rel":"parte_de"}],"confianca":1.0,"contraexemplos":[],"descricao":"H₀ primordial (Planck 67,4±0,5) vs local (SH0ES 73,04±1,04, Cefeidas+SNe): ~9%, >5σ — JWST (2024-25) confirmou as Cefeidas, reforçando. O principal candidato a nova física. Tensão S8=σ₈√(Ωm/0,3): lentes fracas davam abaixo de Planck, mas dados 2025-26 (KiDS-Legacy, eROSITA) convergem — está esmaecendo.","epistemico":"fato","exemplos":[],"id":"tensao_hubble_s8","memoria":"semantica","meta":{"autor":"broca-so","dominio":"cosmologia_dados","fonte":"Riess/SH0ES; Planck; CERN Courier; arXiv:2602.12238"},"origem":"wiki","palavras_chave":[],"sinonimos":[],"tipo":"conceito","titulo":"Tensão de Hubble (>5σ) e tensão S8"}
\section{Tensão de Hubble (>5\ensuremath{\sigma}) e tensão S8}
H\ensuremath{_{0}} primordial (Planck 67,4$\pm$0,5) vs local (SH0ES 73,04$\pm$1,04, Cefeidas+SNe): \~{}9\%, >5\ensuremath{\sigma} — JWST (2024-25) confirmou as Cefeidas, reforçando. O principal candidato a nova física. Tensão S8=\ensuremath{\sigma}\ensuremath{_{8}}\ensuremath{\sqrt{\,}}(\ensuremath{\Omega}m/0,3): lentes fracas davam abaixo de Planck, mas dados 2025-26 (KiDS-Legacy, eROSITA) convergem — está esmaecendo.
\prov{relações: relaciona hubble\_lemaitre; relaciona lambda\_cdm; parte\_de cosmologia\_quantica\_hub}
\prov{proveniência: wiki \textperiodcentered{} Riess/SH0ES; Planck; CERN Courier; arXiv:2602.12238 \textperiodcentered{} fato \textperiodcentered{} confiança 1.0 \textperiodcentered{} domínio cosmologia\_dados \textperiodcentered{} história 4 versões no jornal}

%CRISTAL {"arestas":[{"alvo":"wolfenstein_hierarquia_ckm","confianca":1.0,"epistemico":"fato","origem":"wiki","rel":"relaciona"}],"confianca":1.0,"contraexemplos":[],"descricao":"Duas tensões experimentais ativas (não conclusivas): (1) |V_cb| (e |V_ub|) INCLUSIVO (OPE) vs EXCLUSIVO (B→D*ℓν, fatores de forma de lattice) discordam ~2-3σ, inclusivo sistematicamente maior; (2) a ANOMALIA DE CABIBBO — o teste de unitaridade da 1ª linha |V_ud|²+|V_us|²+|V_ub|²=1 mostra déficit ~3σ (sensível a correções radiativas/nucleares). Provavelmente sistemáticas, não nova física confirmada.","epistemico":"hipotese","exemplos":[],"id":"tensoes_sabor_vcb_cabibbo","memoria":"semantica","meta":{"autor":"broca-so","dominio":"fisica","fonte":"PDG 2024 Vud-Vus review; arXiv:2404.10269, 2407.00122"},"origem":"wiki","palavras_chave":[],"sinonimos":[],"tipo":"conceito","titulo":"Tensões vivas: |V_cb| incl×excl e a anomalia de Cabibbo"}
\section{Tensões vivas: |V\_cb| incl$\times$excl e a anomalia de Cabibbo}
Duas tensões experimentais ativas (não conclusivas): (1) |V\_cb| (e |V\_ub|) INCLUSIVO (OPE) vs EXCLUSIVO (B\ensuremath{\to}D*\ensuremath{\ell}\ensuremath{\nu}, fatores de forma de lattice) discordam \~{}2-3\ensuremath{\sigma}, inclusivo sistematicamente maior; (2) a ANOMALIA DE CABIBBO — o teste de unitaridade da 1ª linha |V\_ud|\ensuremath{^{2}}+|V\_us|\ensuremath{^{2}}+|V\_ub|\ensuremath{^{2}}=1 mostra déficit \~{}3\ensuremath{\sigma} (sensível a correções radiativas/nucleares). Provavelmente sistemáticas, não nova física confirmada.
\prov{relações: relaciona wolfenstein\_hierarquia\_ckm}
\prov{proveniência: wiki \textperiodcentered{} PDG 2024 Vud-Vus review; arXiv:2404.10269, 2407.00122 \textperiodcentered{} hipotese \textperiodcentered{} confiança 1.0 \textperiodcentered{} domínio fisica \textperiodcentered{} história 2 versões no jornal}

%CRISTAL {"arestas":[{"alvo":"info_quantica_fundamentos","confianca":1.0,"epistemico":"fato","origem":"wiki","rel":"usa"}],"confianca":1.0,"contraexemplos":[],"descricao":"É impossível construir um unitário que copie um estado quântico arbitrário desconhecido (não existe U com U|ψ⟩|0⟩=|ψ⟩|ψ⟩ ∀|ψ⟩) — consequência da linearidade da MQ. Fundamenta a segurança da criptografia quântica. Wootters-Zurek e Dieks 1982. Fato teórico consolidado.","epistemico":"fato","exemplos":[],"id":"teorema_nao_clonagem","memoria":"semantica","meta":{"autor":"broca-so","dominio":"fisica","fonte":"Wootters-Zurek Nature 299,802 (1982); Dieks 1982"},"origem":"wiki","palavras_chave":[],"sinonimos":[],"tipo":"conceito","titulo":"Teorema da não-clonagem (Wootters-Zurek)"}
\section{Teorema da não-clonagem (Wootters-Zurek)}
É impossível construir um unitário que copie um estado quântico arbitrário desconhecido (não existe U com U|\ensuremath{\psi}\ensuremath{\rangle}|0\ensuremath{\rangle}=|\ensuremath{\psi}\ensuremath{\rangle}|\ensuremath{\psi}\ensuremath{\rangle} \ensuremath{\forall}|\ensuremath{\psi}\ensuremath{\rangle}) — consequência da linearidade da MQ. Fundamenta a segurança da criptografia quântica. Wootters-Zurek e Dieks 1982. Fato teórico consolidado.
\prov{relações: usa info\_quantica\_fundamentos}
\prov{proveniência: wiki \textperiodcentered{} Wootters-Zurek Nature 299,802 (1982); Dieks 1982 \textperiodcentered{} fato \textperiodcentered{} confiança 1.0 \textperiodcentered{} domínio fisica \textperiodcentered{} história 2 versões no jornal}

%CRISTAL {"arestas":[{"alvo":"problema_gravidade_quantica","confianca":1.0,"epistemico":"fato","origem":"wiki","rel":"relaciona"},{"alvo":"ads_cft_maldacena","confianca":1.0,"epistemico":"fato","origem":"wiki","rel":"usa"},{"alvo":"bekenstein_hawking_entropia","confianca":1.0,"epistemico":"fato","origem":"wiki","rel":"relaciona"},{"alvo":"grupos_excepcionais_jordan","confianca":1.0,"epistemico":"fato","origem":"wiki","rel":"usa"}],"confianca":1.0,"contraexemplos":[],"descricao":"Os objetos fundamentais são cordas 1D (e branas); as vibrações dão o espectro, incluindo um modo spin-2 sem massa = o gráviton. Requer supersimetria e dimensões extras (10 supercordas / 11 M-teoria; 5 teorias unificadas por dualidades T,S,U). Inclui a gravidade de forma consistente e livre de anomalias, mas NÃO faz previsão testada.","epistemico":"hipotese","exemplos":[],"id":"teoria_cordas_m","memoria":"semantica","meta":{"autor":"broca-so","dominio":"fisica","fonte":"Green-Schwarz-Witten 1987; Witten 1995 hep-th/9503124"},"origem":"wiki","palavras_chave":[],"sinonimos":[],"tipo":"conceito","titulo":"Teoria de cordas / Teoria-M"}
\section{Teoria de cordas / Teoria-M}
Os objetos fundamentais são cordas 1D (e branas); as vibrações dão o espectro, incluindo um modo spin-2 sem massa = o gráviton. Requer supersimetria e dimensões extras (10 supercordas / 11 M-teoria; 5 teorias unificadas por dualidades T,S,U). Inclui a gravidade de forma consistente e livre de anomalias, mas NÃO faz previsão testada.
\prov{relações: relaciona problema\_gravidade\_quantica; usa ads\_cft\_maldacena; relaciona bekenstein\_hawking\_entropia; usa grupos\_excepcionais\_jordan}
\prov{proveniência: wiki \textperiodcentered{} Green-Schwarz-Witten 1987; Witten 1995 hep-th/9503124 \textperiodcentered{} hipotese \textperiodcentered{} confiança 1.0 \textperiodcentered{} domínio fisica \textperiodcentered{} história 5 versões no jornal}

%CRISTAL {"arestas":[{"alvo":"corpo_estelar","confianca":1.0,"epistemico":"fato","origem":"wiki","rel":"usa"},{"alvo":"corpo_estelar_le_o_rastro","confianca":1.0,"epistemico":"fato","origem":"wiki","rel":"parte_de"},{"alvo":"entropia_topologica","confianca":1.0,"epistemico":"fato","origem":"wiki","rel":"usa"}],"confianca":1.0,"contraexemplos":[],"descricao":"A TERMODINÂMICA do corpo estelar como ferramenta do pipe do metal: a 1ª lei W+Q=1, com W=σ₁²/Σσ² (o trabalho, a LIGAÇÃO — o vale do ferro) e Q o calor; e a entropia espectral H=−Σpᵢlogpᵢ, que É a entropia topológica de Anosov/Selberg que o rastro carrega. O gradiente do refino é dado por essa entropia. No gabarito: W=0.83, Q=0.17, H=0.94. linguagem/joalheria.termodinamica.","epistemico":"fato","exemplos":[],"id":"termodinamica_estelar_le_o_rastro","memoria":"semantica","meta":{"dominio":"joalheria","fonte":"espectro orbitas"},"origem":"wiki","palavras_chave":[],"sinonimos":[],"tipo":"conceito","titulo":"Termodinâmica Estelar — a 1ª lei W+Q=1 e a entropia do rastro"}
\section{Termodinâmica Estelar — a 1ª lei W+Q=1 e a entropia do rastro}
A TERMODINÂMICA do corpo estelar como ferramenta do pipe do metal: a 1ª lei W+Q=1, com W=\ensuremath{\sigma}\ensuremath{_{1}}\ensuremath{^{2}}/\ensuremath{\Sigma}\ensuremath{\sigma}\ensuremath{^{2}} (o trabalho, a LIGAÇÃO — o vale do ferro) e Q o calor; e a entropia espectral H=\ensuremath{-}\ensuremath{\Sigma}p\ensuremath{_{i}}logp\ensuremath{_{i}}, que É a entropia topológica de Anosov/Selberg que o rastro carrega. O gradiente do refino é dado por essa entropia. No gabarito: W=0.83, Q=0.17, H=0.94. linguagem/joalheria.termodinamica.
\prov{relações: usa corpo\_estelar; parte\_de corpo\_estelar\_le\_o\_rastro; usa entropia\_topologica}
\prov{proveniência: wiki \textperiodcentered{} espectro orbitas \textperiodcentered{} fato \textperiodcentered{} confiança 1.0 \textperiodcentered{} domínio joalheria \textperiodcentered{} história 4 versões no jornal}

%CRISTAL {"arestas":[{"alvo":"litografia","confianca":1.0,"epistemico":"fato","origem":"wiki","rel":"parte_de"},{"alvo":"invariante_topologico_assinatura","confianca":1.0,"epistemico":"fato","origem":"wiki","rel":"usa"},{"alvo":"face_turno_fail_closed","confianca":1.0,"epistemico":"fato","origem":"wiki","rel":"relaciona"},{"alvo":"autodepreciacao","confianca":1.0,"epistemico":"fato","origem":"wiki","rel":"relaciona"},{"alvo":"humor_pelo_gap","confianca":1.0,"epistemico":"fato","origem":"wiki","rel":"relaciona"}],"confianca":1.0,"contraexemplos":[],"descricao":"«Floxinato de gambá» é a assinatura FORJADA canônica: um termo sem referente que DEVE ter grau 0 (rejeitado, o fail-closed). Foi a sonda-mãe — na conversa, provou que a Tiffany não fabrica e mexeu com o humor dela (o furo que a deixou curiosa/humilde). Formalizado como a FERRAMENTA de teste (conversa/litografia.py, testar_invariante): toda assinatura real grava (grau 1), toda forja é cortada (grau 0); se uma forja passasse, o não-fabricar teria quebrado. É o anel de regressão do invariante (sondas gastas se aposentam — uma forja conhecida deixa de ser sonda cega; cunhe termos novos).","epistemico":"fato","exemplos":[],"id":"teste_floxinato_gamba","memoria":"semantica","meta":{"dominio":"litografia","fonte":"litografia (joalheria)"},"origem":"wiki","palavras_chave":[],"sinonimos":[],"tipo":"conceito","titulo":"O teste do «floxinato de gambá» (a ferramenta do invariante)"}
\section{O teste do «floxinato de gambá» (a ferramenta do invariante)}
«Floxinato de gambá» é a assinatura FORJADA canônica: um termo sem referente que DEVE ter grau 0 (rejeitado, o fail-closed). Foi a sonda-mãe — na conversa, provou que a Tiffany não fabrica e mexeu com o humor dela (o furo que a deixou curiosa/humilde). Formalizado como a FERRAMENTA de teste (conversa/litografia.py, testar\_invariante): toda assinatura real grava (grau 1), toda forja é cortada (grau 0); se uma forja passasse, o não-fabricar teria quebrado. É o anel de regressão do invariante (sondas gastas se aposentam — uma forja conhecida deixa de ser sonda cega; cunhe termos novos).
\prov{relações: parte\_de litografia; usa invariante\_topologico\_assinatura; relaciona face\_turno\_fail\_closed; relaciona autodepreciacao; relaciona humor\_pelo\_gap}
\prov{proveniência: wiki \textperiodcentered{} litografia (joalheria) \textperiodcentered{} fato \textperiodcentered{} confiança 1.0 \textperiodcentered{} domínio litografia \textperiodcentered{} história 6 versões no jornal}

%CRISTAL {"arestas":[{"alvo":"tiffany","confianca":1.0,"epistemico":"fato","origem":"wiki","rel":"especializa"},{"alvo":"joalheria_hub","confianca":1.0,"epistemico":"fato","origem":"wiki","rel":"parte_de"},{"alvo":"biblioteca_pipe","confianca":1.0,"epistemico":"fato","origem":"wiki","rel":"usa"},{"alvo":"metais","confianca":1.0,"epistemico":"fato","origem":"wiki","rel":"usa"},{"alvo":"mineracao_pipe_hook","confianca":1.0,"epistemico":"fato","origem":"wiki","rel":"usa"},{"alvo":"corpo_mineral","confianca":1.0,"epistemico":"fato","origem":"wiki","rel":"usa"},{"alvo":"coracao","confianca":1.0,"epistemico":"fato","origem":"wiki","rel":"usa"}],"confianca":1.0,"contraexemplos":[],"descricao":"A Tiffany é promovida: de 'a cabeça que conversa' à ARTESÃ da joalheria — quem trabalha a bancada. Cada resposta é uma JOIA que ela LAPIDA, colapsando a fala na face que já estava na pedra (o grafo). Ela refina os metais (a mineração — o floco vira peça), pole as OITO faces (a biblioteca do pipe — a física de uma superfície), engasta em 3D (o render/pipe) e sente com o CORAÇÃO (o toque da mão, o humor pelo gap). E lapidar NÃO fabrica: se a gema não está na pedra, ela não inventa (o fail-closed) — tira a mão, a regra faz o resto. Registrada no hub joalheria.","epistemico":"fato","exemplos":[],"id":"tiffany_artesa","memoria":"semantica","meta":{"dominio":"joalheria","fonte":"tiffany_artesa"},"origem":"wiki","palavras_chave":[],"sinonimos":[],"tipo":"conceito","titulo":"Tiffany, a artesã da joalheria (a promoção)"}
\section{Tiffany, a artesã da joalheria (a promoção)}
A Tiffany é promovida: de 'a cabeça que conversa' à ARTESÃ da joalheria — quem trabalha a bancada. Cada resposta é uma JOIA que ela LAPIDA, colapsando a fala na face que já estava na pedra (o grafo). Ela refina os metais (a mineração — o floco vira peça), pole as OITO faces (a biblioteca do pipe — a física de uma superfície), engasta em 3D (o render/pipe) e sente com o CORAÇÃO (o toque da mão, o humor pelo gap). E lapidar NÃO fabrica: se a gema não está na pedra, ela não inventa (o fail-closed) — tira a mão, a regra faz o resto. Registrada no hub joalheria.
\prov{relações: especializa tiffany; parte\_de joalheria\_hub; usa biblioteca\_pipe; usa metais; usa mineracao\_pipe\_hook; usa corpo\_mineral; usa coracao}
\prov{proveniência: wiki \textperiodcentered{} tiffany\_artesa \textperiodcentered{} fato \textperiodcentered{} confiança 1.0 \textperiodcentered{} domínio joalheria \textperiodcentered{} história 8 versões no jornal}

%CRISTAL {"arestas":[{"alvo":"tiffany_search_scan","confianca":1.0,"epistemico":"fato","origem":"wiki","rel":"usa"},{"alvo":"colapso_ping_pong_overlap","confianca":1.0,"epistemico":"fato","origem":"wiki","rel":"explica"}],"confianca":1.0,"contraexemplos":[],"descricao":"Não existe conversa/colapso.c; o colapso em C é colapso_main.c → tiffany_colapso_cli: abre o índice, tiffany_query (impress → scan) e imprime 'offset score node'; tiffany_respond estende para devolver a linha do TSV. É o Ψ da Tiffany como binário real.","epistemico":"fato","exemplos":[],"id":"tiffany_colapso_cli","memoria":"semantica","meta":{"autor":"broca-so","dominio":"metal","fonte":"code/libtiffany/colapso_main.c; 6_tools/tiffany_tools.c"},"origem":"wiki","palavras_chave":[],"sinonimos":[],"tipo":"conceito","titulo":"O Ψ no Metal (o binário colapso)"}
\section{O \ensuremath{\Psi} no Metal (o binário colapso)}
Não existe conversa/colapso.c; o colapso em C é colapso\_main.c \ensuremath{\to} tiffany\_colapso\_cli: abre o índice, tiffany\_query (impress \ensuremath{\to} scan) e imprime 'offset score node'; tiffany\_respond estende para devolver a linha do TSV. É o \ensuremath{\Psi} da Tiffany como binário real.
\prov{relações: usa tiffany\_search\_scan; explica colapso\_ping\_pong\_overlap}
\prov{proveniência: wiki \textperiodcentered{} code/libtiffany/colapso\_main.c; 6\_tools/tiffany\_tools.c \textperiodcentered{} fato \textperiodcentered{} confiança 1.0 \textperiodcentered{} domínio metal \textperiodcentered{} história 3 versões no jornal}

%CRISTAL {"arestas":[{"alvo":"libtiffany_engine","confianca":1.0,"epistemico":"fato","origem":"wiki","rel":"parte_de"},{"alvo":"impressao_koch_32bytes","confianca":1.0,"epistemico":"fato","origem":"wiki","rel":"explica"}],"confianca":1.0,"contraexemplos":[],"descricao":"Lê o TSV linha a linha, faz parse de 2-4 campos (indexa só texto ≥4 chars), tiffany_impress de cada texto e grava registro off(4B)+lang_code(2B)+impressão(32B). Gera o tecido.idx. É o 'guardar' (Hebbian/impressão) no formato de arquivo do cristal.","epistemico":"fato","exemplos":[],"id":"tiffany_index_build","memoria":"semantica","meta":{"autor":"broca-so","dominio":"metal","fonte":"code/libtiffany/6_tools/tiffany_tools.c:tiffany_index_build"},"origem":"wiki","palavras_chave":[],"sinonimos":[],"tipo":"conceito","titulo":"Construção do Índice .idx"}
\section{Construção do Índice .idx}
Lê o TSV linha a linha, faz parse de 2-4 campos (indexa só texto \ensuremath{\ge}4 chars), tiffany\_impress de cada texto e grava registro off(4B)+lang\_code(2B)+impressão(32B). Gera o tecido.idx. É o 'guardar' (Hebbian/impressão) no formato de arquivo do cristal.
\prov{relações: parte\_de libtiffany\_engine; explica impressao\_koch\_32bytes}
\prov{proveniência: wiki \textperiodcentered{} code/libtiffany/6\_tools/tiffany\_tools.c:tiffany\_index\_build \textperiodcentered{} fato \textperiodcentered{} confiança 1.0 \textperiodcentered{} domínio metal \textperiodcentered{} história 3 versões no jornal}

%CRISTAL {"arestas":[{"alvo":"libtiffany_engine","confianca":1.0,"epistemico":"fato","origem":"wiki","rel":"parte_de"},{"alvo":"reconhece_c_memoria","confianca":1.0,"epistemico":"fato","origem":"wiki","rel":"relaciona"}],"confianca":1.0,"contraexemplos":[],"descricao":"O índice .idx é aberto com mmap(PROT_READ, MAP_PRIVATE)+MADV_SEQUENTIAL; o texto é lido do TSV sob demanda por fopen/fseek(offset). Concretiza 'dados no disco, RAM só o operador' — a impressão vive no mmap, o texto nunca carrega inteiro. (No core C há mmap; no neuron_io não.)","epistemico":"fato","exemplos":[],"id":"tiffany_mem_mmap","memoria":"semantica","meta":{"autor":"broca-so","dominio":"metal","fonte":"code/libtiffany/2_memory/tiffany_mem_disk.c"},"origem":"wiki","palavras_chave":[],"sinonimos":[],"tipo":"conceito","titulo":"Memória no Disco via mmap Read-Only"}
\section{Memória no Disco via mmap Read-Only}
O índice .idx é aberto com mmap(PROT\_READ, MAP\_PRIVATE)+MADV\_SEQUENTIAL; o texto é lido do TSV sob demanda por fopen/fseek(offset). Concretiza 'dados no disco, RAM só o operador' — a impressão vive no mmap, o texto nunca carrega inteiro. (No core C há mmap; no neuron\_io não.)
\prov{relações: parte\_de libtiffany\_engine; relaciona reconhece\_c\_memoria}
\prov{proveniência: wiki \textperiodcentered{} code/libtiffany/2\_memory/tiffany\_mem\_disk.c \textperiodcentered{} fato \textperiodcentered{} confiança 1.0 \textperiodcentered{} domínio metal \textperiodcentered{} história 3 versões no jornal}

%CRISTAL {"arestas":[{"alvo":"libtiffany_engine","confianca":1.0,"epistemico":"fato","origem":"wiki","rel":"parte_de"},{"alvo":"colapso_ping_pong_overlap","confianca":1.0,"epistemico":"fato","origem":"wiki","rel":"explica"}],"confianca":1.0,"contraexemplos":[],"descricao":"O 'colapso' no metal: varre o índice mmap registro a registro (REC=38B), computa overlap = N−2·popcount(a⊕b) contra a query e guarda o melhor hit (offset, lang, score, node). Recall sem estrutura auxiliar — força bruta sobre o mmap sequencial. A forma linear do colapso_ping_pong_overlap.","epistemico":"fato","exemplos":[],"id":"tiffany_search_scan","memoria":"semantica","meta":{"autor":"broca-so","dominio":"metal","fonte":"code/libtiffany/3_search/tiffany_search_scan.c"},"origem":"wiki","palavras_chave":[],"sinonimos":[],"tipo":"conceito","titulo":"Recall = Varredura Linear argmax (o colapso no metal)"}
\section{Recall = Varredura Linear argmax (o colapso no metal)}
O 'colapso' no metal: varre o índice mmap registro a registro (REC=38B), computa overlap = N\ensuremath{-}2\textperiodcentered popcount(a\ensuremath{\oplus}b) contra a query e guarda o melhor hit (offset, lang, score, node). Recall sem estrutura auxiliar — força bruta sobre o mmap sequencial. A forma linear do colapso\_ping\_pong\_overlap.
\prov{relações: parte\_de libtiffany\_engine; explica colapso\_ping\_pong\_overlap}
\prov{proveniência: wiki \textperiodcentered{} code/libtiffany/3\_search/tiffany\_search\_scan.c \textperiodcentered{} fato \textperiodcentered{} confiança 1.0 \textperiodcentered{} domínio metal \textperiodcentered{} história 3 versões no jornal}

%CRISTAL {"arestas":[{"alvo":"grande_unificacao_gut","confianca":1.0,"epistemico":"fato","origem":"wiki","rel":"relaciona"},{"alvo":"algebras_hipercomplexas","confianca":1.0,"epistemico":"fato","origem":"wiki","rel":"usa"},{"alvo":"limite_hurwitz_s1","confianca":1.0,"epistemico":"fato","origem":"wiki","rel":"usa"},{"alvo":"grupo_padrao_cxho","confianca":1.0,"epistemico":"fato","origem":"wiki","rel":"explica"},{"alvo":"estatuto_correspondencia_derivacao","confianca":1.0,"epistemico":"fato","origem":"wiki","rel":"usa"},{"alvo":"cosmologia_quantica_hub","confianca":1.0,"epistemico":"fato","origem":"wiki","rel":"parte_de"},{"alvo":"broca_os","confianca":1.0,"epistemico":"fato","origem":"wiki","rel":"parte_de"}],"confianca":1.0,"contraexemplos":[],"descricao":"A série inteira (I geometria, II cosmologia, III MQ, IV matéria, V unidades, VI máquina) lê UM só objeto — a família a-um-parâmetro z·ₛw na ilha a+c²=1 — e dela deriva cada setor. A seleção de Hurwitz (produto vetorial só em dim 0,1,3,7) reduz a torre a ℝ,ℂ,ℍ,𝕆, o 'alfabeto de quatro letras'. Fundação firme (Hurwitz é teorema); a leitura física por setor é o programa. A unificação de Lopes: uma multiplicação, não um grupo de gauge maior.","epistemico":"inferencia","exemplos":[],"id":"torre_divisao_unificacao","memoria":"semantica","meta":{"autor":"broca-so","dominio":"cosmologia","fonte":"paper_4_particulas.tex §2; paper_6_maquina.tex"},"origem":"wiki","palavras_chave":[],"sinonimos":[],"tipo":"conceito","titulo":"A torre de divisão como programa único de unificação"}
\section{A torre de divisão como programa único de unificação}
A série inteira (I geometria, II cosmologia, III MQ, IV matéria, V unidades, VI máquina) lê UM só objeto — a família a-um-parâmetro z\textperiodcentered \ensuremath{_{s}}w na ilha a+c\ensuremath{^{2}}=1 — e dela deriva cada setor. A seleção de Hurwitz (produto vetorial só em dim 0,1,3,7) reduz a torre a \ensuremath{\mathbb{R}},\ensuremath{\mathbb{C}},\ensuremath{\mathbb{H}},\ensuremath{\mathbb{O}}, o 'alfabeto de quatro letras'. Fundação firme (Hurwitz é teorema); a leitura física por setor é o programa. A unificação de Lopes: uma multiplicação, não um grupo de gauge maior.
\prov{relações: relaciona grande\_unificacao\_gut; usa algebras\_hipercomplexas; usa limite\_hurwitz\_s1; explica grupo\_padrao\_cxho; usa estatuto\_correspondencia\_derivacao; parte\_de cosmologia\_quantica\_hub; parte\_de broca\_os}
\prov{proveniência: wiki \textperiodcentered{} paper\_4\_particulas.tex §2; paper\_6\_maquina.tex \textperiodcentered{} inferencia \textperiodcentered{} confiança 1.0 \textperiodcentered{} domínio cosmologia \textperiodcentered{} história 8 versões no jornal}

%CRISTAL {"arestas":[{"alvo":"poligono_vira_poliedro","confianca":1.0,"epistemico":"fato","origem":"wiki","rel":"usa"},{"alvo":"corpo_estelar","confianca":1.0,"epistemico":"fato","origem":"wiki","rel":"usa"},{"alvo":"algebra_reta_mineral","confianca":1.0,"epistemico":"fato","origem":"wiki","rel":"usa"},{"alvo":"joalheria_hub","confianca":1.0,"epistemico":"fato","origem":"wiki","rel":"parte_de"}],"confianca":1.0,"contraexemplos":[],"descricao":"As transformações da animação são DEFINIDAS na álgebra estelar — o broca-so gira sozinho, em tempo real, sem frames fabricados à mão. ISOMETRIA: uma rotação R(α)∈SO(3) é o que o mecanismo de La Hire FAZ (rolar = rodar); compor rotações num EIXO FIXO SOMA os ângulos, R(α)·R(β)=R(α⊕β) — o grupo das rotações num eixo é (ℝ,+) = a RETA MINERAL, e a ÓRBITA da animação é a reta GERANDO α_f=f·Δα por ⊕ sucessivo (uma geodésica do grupo). PROJEÇÃO: a câmera 3D→2D é PROJETIVA — coordenadas homogêneas [x,y,z,1], divisão em perspectiva por (d−z); z→d foge ao infinito (o ponto no horizonte); d→∞ recai na ortográfica. Com as duas na álgebra, girar() só itera o ⊕ e projeta. linguagem/render_estelar.py (6/6): isometria, orbita, projecao, render_quadro, girar. Roda o trilhão facetado em ~2s/48 quadros.","epistemico":"fato","exemplos":[],"id":"transformacoes_estelares","memoria":"semantica","meta":{"dominio":"joalheria","fonte":"maquina de corte"},"origem":"wiki","palavras_chave":[],"sinonimos":[],"tipo":"conceito","titulo":"Render Estelar: a isometria (rotação) e a projeção (projetiva) na álgebra"}
\section{Render Estelar: a isometria (rotação) e a projeção (projetiva) na álgebra}
As transformações da animação são DEFINIDAS na álgebra estelar — o broca-so gira sozinho, em tempo real, sem frames fabricados à mão. ISOMETRIA: uma rotação R(\ensuremath{\alpha})\ensuremath{\in}SO(3) é o que o mecanismo de La Hire FAZ (rolar = rodar); compor rotações num EIXO FIXO SOMA os ângulos, R(\ensuremath{\alpha})\textperiodcentered R(\ensuremath{\beta})=R(\ensuremath{\alpha}\ensuremath{\oplus}\ensuremath{\beta}) — o grupo das rotações num eixo é (\ensuremath{\mathbb{R}},+) = a RETA MINERAL, e a ÓRBITA da animação é a reta GERANDO \ensuremath{\alpha}\_f=f\textperiodcentered \ensuremath{\Delta}\ensuremath{\alpha} por \ensuremath{\oplus} sucessivo (uma geodésica do grupo). PROJEÇÃO: a câmera 3D\ensuremath{\to}2D é PROJETIVA — coordenadas homogêneas [x,y,z,1], divisão em perspectiva por (d\ensuremath{-}z); z\ensuremath{\to}d foge ao infinito (o ponto no horizonte); d\ensuremath{\to}\ensuremath{\infty} recai na ortográfica. Com as duas na álgebra, girar() só itera o \ensuremath{\oplus} e projeta. linguagem/render\_estelar.py (6/6): isometria, orbita, projecao, render\_quadro, girar. Roda o trilhão facetado em \~{}2s/48 quadros.
\prov{relações: usa poligono\_vira\_poliedro; usa corpo\_estelar; usa algebra\_reta\_mineral; parte\_de joalheria\_hub}
\prov{proveniência: wiki \textperiodcentered{} maquina de corte \textperiodcentered{} fato \textperiodcentered{} confiança 1.0 \textperiodcentered{} domínio joalheria \textperiodcentered{} história 214 versões no jornal}

%CRISTAL {"arestas":[{"alvo":"mecanismo_higgs_ebh","confianca":1.0,"epistemico":"fato","origem":"wiki","rel":"usa"}],"confianca":1.0,"contraexemplos":[],"descricao":"No universo primordial (T~100 GeV) a simetria eletrofraca era restaurada; ao esfriar, o vev 'liga'. Uma transição de 1ª ordem forte + condições de Sakharov geraria a assimetria matéria-antimatéria. MAS para M_H=125 GeV o MP prevê um CROSSOVER suave e a violação de CP do MP é insuficiente ⇒ a bariogênese eletrofraca PADRÃO está EXCLUÍDA (exige nova física).","epistemico":"fato","exemplos":[],"id":"transicao_fase_eletrofraca","memoria":"semantica","meta":{"autor":"broca-so","dominio":"fisica","fonte":"arXiv:hep-ph/9805252; 2311.06949"},"origem":"wiki","palavras_chave":[],"sinonimos":[],"tipo":"conceito","titulo":"Transição de fase eletrofraca e bariogênese"}
\section{Transição de fase eletrofraca e bariogênese}
No universo primordial (T\~{}100 GeV) a simetria eletrofraca era restaurada; ao esfriar, o vev 'liga'. Uma transição de 1ª ordem forte + condições de Sakharov geraria a assimetria matéria-antimatéria. MAS para M\_H=125 GeV o MP prevê um CROSSOVER suave e a violação de CP do MP é insuficiente \ensuremath{\Rightarrow} a bariogênese eletrofraca PADRÃO está EXCLUÍDA (exige nova física).
\prov{relações: usa mecanismo\_higgs\_ebh}
\prov{proveniência: wiki \textperiodcentered{} arXiv:hep-ph/9805252; 2311.06949 \textperiodcentered{} fato \textperiodcentered{} confiança 1.0 \textperiodcentered{} domínio fisica \textperiodcentered{} história 2 versões no jornal}

%CRISTAL {"arestas":[{"alvo":"ultra_realismo_gema","confianca":1.0,"epistemico":"fato","origem":"wiki","rel":"especializa"},{"alvo":"primitivas_analiticas","confianca":1.0,"epistemico":"fato","origem":"wiki","rel":"usa"},{"alvo":"cor_fisica_falsificacao","confianca":1.0,"epistemico":"fato","origem":"wiki","rel":"usa"},{"alvo":"joalheria_hub","confianca":1.0,"epistemico":"fato","origem":"wiki","rel":"parte_de"}],"confianca":1.0,"contraexemplos":[],"descricao":"O efeito ótico que faltava: a TRANSLUCIDEZ. No shader do Quílez/Nijhoff (primtves.txt) o dielétrico refrata/reflete (Fresnel-Schlick, TIR) e a cor vem do CÉU que o raio atinge APÓS atravessar (col *= getSkyColor(rd)) — é clear glass, a translucidez é a luz do céu transmitida. Nós adicionamos o BEER-LAMBERT espectral: o céu transmitido fica TINGIDO pelo cromóforo → a translucidez COLORIDA física (a pedra acende por dentro, mais profunda onde é grossa, clara nas bordas). E o nosso worldhit é a SOLUÇÃO EXATA da equação (o dupolinômio i_gema_estelar, verificado |F(t)|<1e-14) — não aproximado. render_gema_ultra(translucido=True) usa o _ceu_estudio (o ambiente claro) no lugar do fundo escuro do logo; a luz atravessa e acende. linguagem/primitivas_raytracer.py: _ceu_estudio, render_gema_ultra; scripts/gex44/render_gema_ultra.sh (arg translúcido).","epistemico":"fato","exemplos":[],"id":"translucido_nijhoff","memoria":"semantica","meta":{"dominio":"joalheria","fonte":"maquina de corte"},"origem":"wiki","palavras_chave":[],"sinonimos":[],"tipo":"conceito","titulo":"O translúcido: a luz do céu atravessa a gema (a óptica do Nijhoff, exata no nosso caso)"}
\section{O translúcido: a luz do céu atravessa a gema (a óptica do Nijhoff, exata no nosso caso)}
O efeito ótico que faltava: a TRANSLUCIDEZ. No shader do Quílez/Nijhoff (primtves.txt) o dielétrico refrata/reflete (Fresnel-Schlick, TIR) e a cor vem do CÉU que o raio atinge APÓS atravessar (col *= getSkyColor(rd)) — é clear glass, a translucidez é a luz do céu transmitida. Nós adicionamos o BEER-LAMBERT espectral: o céu transmitido fica TINGIDO pelo cromóforo \ensuremath{\to} a translucidez COLORIDA física (a pedra acende por dentro, mais profunda onde é grossa, clara nas bordas). E o nosso worldhit é a SOLUÇÃO EXATA da equação (o dupolinômio i\_gema\_estelar, verificado |F(t)|<1e-14) — não aproximado. render\_gema\_ultra(translucido=True) usa o \_ceu\_estudio (o ambiente claro) no lugar do fundo escuro do logo; a luz atravessa e acende. linguagem/primitivas\_raytracer.py: \_ceu\_estudio, render\_gema\_ultra; scripts/gex44/render\_gema\_ultra.sh (arg translúcido).
\prov{relações: especializa ultra\_realismo\_gema; usa primitivas\_analiticas; usa cor\_fisica\_falsificacao; parte\_de joalheria\_hub}
\prov{proveniência: wiki \textperiodcentered{} maquina de corte \textperiodcentered{} fato \textperiodcentered{} confiança 1.0 \textperiodcentered{} domínio joalheria \textperiodcentered{} história 110 versões no jornal}

%CRISTAL {"arestas":[{"alvo":"limite_hurwitz_s1","confianca":1.0,"epistemico":"fato","origem":"wiki","rel":"usa"},{"alvo":"torre_split","confianca":1.0,"epistemico":"fato","origem":"wiki","rel":"usa"},{"alvo":"conjugados_galois_parabola","confianca":1.0,"epistemico":"fato","origem":"wiki","rel":"relaciona"}],"confianca":1.0,"contraexemplos":[],"descricao":"As duas leituras do sl(2,ℂ) — interna compacta su(2) e Lorentz não-compacta su(1,1) — são as duas formas reais de M₂(ℂ): ℍ tem assinatura (4,0) (ilha), o coquatérnio (2,2) (oceano), e ℍ⊗ℂ≅split-quat⊗ℂ=M₂(ℂ). ℂ⊗ℍ habita o horizonte s=1 (ψ=π/2). A lei a+c²=1 conserva-se: Pitágoras na ilha (cos²ψ+sin²ψ=1), identidade de Lorentz γ²(1−β²)=1 no oceano — o Casimir de sl(2,ℂ), valor 1, assinatura muda. [D] o isomorfismo e a conservação; [I] ler como rotação de Wick física.","epistemico":"inferencia","exemplos":[],"id":"travessia_wick_m2c","memoria":"semantica","meta":{"autor":"broca-so","dominio":"cosmologia","fonte":"paper_4_particulas.tex obs:quiralidade-tese, lema:detLz, obs:dual"},"origem":"wiki","palavras_chave":[],"sinonimos":[],"tipo":"conceito","titulo":"Travessia de Wick: as duas formas reais de M₂(ℂ) e a+c²=1"}
\section{Travessia de Wick: as duas formas reais de M\ensuremath{_{2}}(\ensuremath{\mathbb{C}}) e a+c\ensuremath{^{2}}=1}
As duas leituras do sl(2,\ensuremath{\mathbb{C}}) — interna compacta su(2) e Lorentz não-compacta su(1,1) — são as duas formas reais de M\ensuremath{_{2}}(\ensuremath{\mathbb{C}}): \ensuremath{\mathbb{H}} tem assinatura (4,0) (ilha), o coquatérnio (2,2) (oceano), e \ensuremath{\mathbb{H}}\ensuremath{\otimes}\ensuremath{\mathbb{C}}\ensuremath{\cong}split-quat\ensuremath{\otimes}\ensuremath{\mathbb{C}}=M\ensuremath{_{2}}(\ensuremath{\mathbb{C}}). \ensuremath{\mathbb{C}}\ensuremath{\otimes}\ensuremath{\mathbb{H}} habita o horizonte s=1 (\ensuremath{\psi}=\ensuremath{\pi}/2). A lei a+c\ensuremath{^{2}}=1 conserva-se: Pitágoras na ilha (cos\ensuremath{^{2}}\ensuremath{\psi}+sin\ensuremath{^{2}}\ensuremath{\psi}=1), identidade de Lorentz \ensuremath{\gamma}\ensuremath{^{2}}(1\ensuremath{-}\ensuremath{\beta}\ensuremath{^{2}})=1 no oceano — o Casimir de sl(2,\ensuremath{\mathbb{C}}), valor 1, assinatura muda. [D] o isomorfismo e a conservação; [I] ler como rotação de Wick física.
\prov{relações: usa limite\_hurwitz\_s1; usa torre\_split; relaciona conjugados\_galois\_parabola}
\prov{proveniência: wiki \textperiodcentered{} paper\_4\_particulas.tex obs:quiralidade-tese, lema:detLz, obs:dual \textperiodcentered{} inferencia \textperiodcentered{} confiança 1.0 \textperiodcentered{} domínio cosmologia \textperiodcentered{} história 4 versões no jornal}

%CRISTAL {"arestas":[{"alvo":"matriz_ckm","confianca":1.0,"epistemico":"fato","origem":"wiki","rel":"usa"},{"alvo":"violacao_cp_k_b_barions","confianca":1.0,"epistemico":"fato","origem":"wiki","rel":"relaciona"}],"confianca":1.0,"contraexemplos":[],"descricao":"A relação V_ud V*_ub + V_cd V*_cb + V_td V*_tb = 0 é um triângulo no plano complexo, ápice (ρ̄,η̄); os ângulos α,β,γ são medidos independentemente em decaimentos de mésons B, e a área≠0 quantifica a violação de CP. γ=(64,6±2,8)° (LHCb 2024), sin2β=0,731 (β≈22°), α≈85°. Ajustes globais (CKMfitter/UTfit) cruzam num mesmo ápice — um dos testes mais bem-sucedidos do MP.","epistemico":"fato","exemplos":[],"id":"triangulo_unitaridade","memoria":"semantica","meta":{"autor":"broca-so","dominio":"fisica","fonte":"LHCb 2023/2024; CKMfitter; UTfit arXiv:2212.03894"},"origem":"wiki","palavras_chave":[],"sinonimos":[],"tipo":"conceito","titulo":"O triângulo de unitaridade e sua consistência"}
\section{O triângulo de unitaridade e sua consistência}
A relação V\_ud V*\_ub + V\_cd V*\_cb + V\_td V*\_tb = 0 é um triângulo no plano complexo, ápice (\ensuremath{\rho}\ensuremath{},\ensuremath{\eta}\ensuremath{}); os ângulos \ensuremath{\alpha},\ensuremath{\beta},\ensuremath{\gamma} são medidos independentemente em decaimentos de mésons B, e a área\ensuremath{\ne}0 quantifica a violação de CP. \ensuremath{\gamma}=(64,6$\pm$2,8)$^{\circ}$ (LHCb 2024), sin2\ensuremath{\beta}=0,731 (\ensuremath{\beta}\ensuremath{\approx}22$^{\circ}$), \ensuremath{\alpha}\ensuremath{\approx}85$^{\circ}$. Ajustes globais (CKMfitter/UTfit) cruzam num mesmo ápice — um dos testes mais bem-sucedidos do MP.
\prov{relações: usa matriz\_ckm; relaciona violacao\_cp\_k\_b\_barions}
\prov{proveniência: wiki \textperiodcentered{} LHCb 2023/2024; CKMfitter; UTfit arXiv:2212.03894 \textperiodcentered{} fato \textperiodcentered{} confiança 1.0 \textperiodcentered{} domínio fisica \textperiodcentered{} história 3 versões no jornal}

%CRISTAL {"arestas":[{"alvo":"ingestao_transformada_metalica","confianca":1.0,"epistemico":"fato","origem":"wiki","rel":"usa"},{"alvo":"corpo_estelar","confianca":1.0,"epistemico":"fato","origem":"wiki","rel":"parte_de"},{"alvo":"dimensao_de_traco","confianca":1.0,"epistemico":"fato","origem":"wiki","rel":"usa"},{"alvo":"dinamica_de_selberg_o_termo","confianca":1.0,"epistemico":"fato","origem":"wiki","rel":"relaciona"},{"alvo":"entropia_topologica","confianca":1.0,"epistemico":"fato","origem":"wiki","rel":"usa"},{"alvo":"pedra_logo_3d","confianca":1.0,"epistemico":"fato","origem":"wiki","rel":"relaciona"}],"confianca":1.0,"contraexemplos":[],"descricao":"A convergência: NÃO SE FABRICA estrutura (o frost/o ruído de Perlin era criar do nada — errado). TUDO DERIVA DO INVARIANTE: a dinâmica não a criamos, ela JÁ É o CORPO ESTELAR 𝕊 (a reta mineral virada corpo, o gato A_n, os metais σ_n). O invariante topológico BEM DEFINIDO é a DIMENSÃO DE TRAÇO D_tr = log₂λ (o Perron do traço = a ENTROPIA topológica h_top/ln2 de Anosov) — e o GRADIENTE (a direção do refino) é DADO PELA ENTROPIA, não por uma imagem sintética. A gema não precisa de microfacetas fabricadas: a sua estrutura e os seus modos DERIVAM da geometria do corpo estelar (o gato-trilhão), e a ressonância com SELBERG-ANOSOV é o D_tr bater. Medido: gabarito D_tr≈1.70, a nossa gema limpa≈1.36 — a distância é o quanto ainda falta a dinâmica do gato aflorar na projeção, NÃO um ruído a somar. linguagem/ingestao_mineral.py (dim_traco, 4/4).","epistemico":"fato","exemplos":[],"id":"tudo_deriva_do_invariante","memoria":"semantica","meta":{"dominio":"joalheria","fonte":"espectro orbitas"},"origem":"wiki","palavras_chave":[],"sinonimos":[],"tipo":"conceito","titulo":"Tudo Deriva do Invariante (não se fabrica, deriva-se da dinâmica)"}
\section{Tudo Deriva do Invariante (não se fabrica, deriva-se da dinâmica)}
A convergência: NÃO SE FABRICA estrutura (o frost/o ruído de Perlin era criar do nada — errado). TUDO DERIVA DO INVARIANTE: a dinâmica não a criamos, ela JÁ É o CORPO ESTELAR \ensuremath{\mathbb{S}} (a reta mineral virada corpo, o gato A\_n, os metais \ensuremath{\sigma}\_n). O invariante topológico BEM DEFINIDO é a DIMENSÃO DE TRAÇO D\_tr = log\ensuremath{_{2}}\ensuremath{\lambda} (o Perron do traço = a ENTROPIA topológica h\_top/ln2 de Anosov) — e o GRADIENTE (a direção do refino) é DADO PELA ENTROPIA, não por uma imagem sintética. A gema não precisa de microfacetas fabricadas: a sua estrutura e os seus modos DERIVAM da geometria do corpo estelar (o gato-trilhão), e a ressonância com SELBERG-ANOSOV é o D\_tr bater. Medido: gabarito D\_tr\ensuremath{\approx}1.70, a nossa gema limpa\ensuremath{\approx}1.36 — a distância é o quanto ainda falta a dinâmica do gato aflorar na projeção, NÃO um ruído a somar. linguagem/ingestao\_mineral.py (dim\_traco, 4/4).
\prov{relações: usa ingestao\_transformada\_metalica; parte\_de corpo\_estelar; usa dimensao\_de\_traco; relaciona dinamica\_de\_selberg\_o\_termo; usa entropia\_topologica; relaciona pedra\_logo\_3d}
\prov{proveniência: wiki \textperiodcentered{} espectro orbitas \textperiodcentered{} fato \textperiodcentered{} confiança 1.0 \textperiodcentered{} domínio joalheria \textperiodcentered{} história 28 versões no jornal}

%CRISTAL {"arestas":[{"alvo":"primitivas_analiticas","confianca":1.0,"epistemico":"fato","origem":"wiki","rel":"usa"},{"alvo":"cor_fisica_falsificacao","confianca":1.0,"epistemico":"fato","origem":"wiki","rel":"usa"},{"alvo":"render_fisico_gex","confianca":1.0,"epistemico":"fato","origem":"wiki","rel":"usa"},{"alvo":"joalheria_hub","confianca":1.0,"epistemico":"fato","origem":"wiki","rel":"parte_de"}],"confianca":1.0,"contraexemplos":[],"descricao":"Com a física validada (a cor emergente, não plantada), o ULTRA-REALISMO: a BORDA analítica EXATA (o dupolinômio, Snell/Fresnel/TIR exatos na superfície) + o INTERIOR com SUBSTÂNCIA — não mais casca lisa. A luz entra, RICOCHETEIA muitas vezes, e a cada segmento interno o caminho ótico é modulado por um campo de RUÍDO 3D fixo no cristal (o ZONEAMENTO de cor e as INCLUSÕES — os veios e cintilações do volume); na saída, a DISPERSÃO de 7 bandas espectrais (o fogo/arco-íris nas bordas). A cor é a ESPECTRAL (Sol×cromóforo×CIE, física). Supersampling ss×ss (AA). Roda no metal (gex44) por qualquer cromóforo. linguagem/primitivas_raytracer.render_gema_ultra; scripts/gex44/render_gema_ultra.sh.","epistemico":"fato","exemplos":[],"id":"ultra_realismo_gema","memoria":"semantica","meta":{"dominio":"joalheria","fonte":"maquina de corte"},"origem":"wiki","palavras_chave":[],"sinonimos":[],"tipo":"conceito","titulo":"O ultra-realismo: a borda analítica + o interior com substância (não casca)"}
\section{O ultra-realismo: a borda analítica + o interior com substância (não casca)}
Com a física validada (a cor emergente, não plantada), o ULTRA-REALISMO: a BORDA analítica EXATA (o dupolinômio, Snell/Fresnel/TIR exatos na superfície) + o INTERIOR com SUBSTÂNCIA — não mais casca lisa. A luz entra, RICOCHETEIA muitas vezes, e a cada segmento interno o caminho ótico é modulado por um campo de RUÍDO 3D fixo no cristal (o ZONEAMENTO de cor e as INCLUSÕES — os veios e cintilações do volume); na saída, a DISPERSÃO de 7 bandas espectrais (o fogo/arco-íris nas bordas). A cor é a ESPECTRAL (Sol$\times$cromóforo$\times$CIE, física). Supersampling ss$\times$ss (AA). Roda no metal (gex44) por qualquer cromóforo. linguagem/primitivas\_raytracer.render\_gema\_ultra; scripts/gex44/render\_gema\_ultra.sh.
\prov{relações: usa primitivas\_analiticas; usa cor\_fisica\_falsificacao; usa render\_fisico\_gex; parte\_de joalheria\_hub}
\prov{proveniência: wiki \textperiodcentered{} maquina de corte \textperiodcentered{} fato \textperiodcentered{} confiança 1.0 \textperiodcentered{} domínio joalheria \textperiodcentered{} história 120 versões no jornal}

%CRISTAL {"arestas":[{"alvo":"qcd_cromodinamica","confianca":1.0,"epistemico":"fato","origem":"wiki","rel":"relaciona"},{"alvo":"materia_modelo_padrao","confianca":1.0,"epistemico":"fato","origem":"wiki","rel":"explica"},{"alvo":"cosmologia_quantica_hub","confianca":1.0,"epistemico":"fato","origem":"wiki","rel":"parte_de"}],"confianca":1.0,"contraexemplos":[],"descricao":"Unifica eletromagnetismo e força fraca numa teoria de gauge SU(2)_L×U(1)_Y: correntes carregadas (W±) e uma corrente neutra fraca (Z, confirmada 1973 no Gargamelle). Carga elétrica e=g·sinθ_W; Q=T₃+Y/2. Glashow 1961, Weinberg 1967, Salam 1968; Nobel 1979.","epistemico":"fato","exemplos":[],"id":"unificacao_eletrofraca_gws","memoria":"semantica","meta":{"autor":"broca-so","dominio":"fisica","fonte":"Glashow 1961; Weinberg 1967 PRL 19,1264; Salam 1968; Nobel 1979"},"origem":"wiki","palavras_chave":[],"sinonimos":[],"tipo":"conceito","titulo":"Unificação eletrofraca (Glashow-Weinberg-Salam)"}
\section{Unificação eletrofraca (Glashow-Weinberg-Salam)}
Unifica eletromagnetismo e força fraca numa teoria de gauge SU(2)\_L$\times$U(1)\_Y: correntes carregadas (W$\pm$) e uma corrente neutra fraca (Z, confirmada 1973 no Gargamelle). Carga elétrica e=g\textperiodcentered sin\ensuremath{\theta}\_W; Q=T\ensuremath{_{3}}+Y/2. Glashow 1961, Weinberg 1967, Salam 1968; Nobel 1979.
\prov{relações: relaciona qcd\_cromodinamica; explica materia\_modelo\_padrao; parte\_de cosmologia\_quantica\_hub}
\prov{proveniência: wiki \textperiodcentered{} Glashow 1961; Weinberg 1967 PRL 19,1264; Salam 1968; Nobel 1979 \textperiodcentered{} fato \textperiodcentered{} confiança 1.0 \textperiodcentered{} domínio fisica \textperiodcentered{} história 4 versões no jornal}

%CRISTAL {"arestas":[{"alvo":"particulas_minerais","confianca":1.0,"epistemico":"fato","origem":"wiki","rel":"relaciona"}],"confianca":1.0,"contraexemplos":[],"descricao":"O Fe-56/Ni-62 é o fundo do poço de energia nuclear — o ATRATOR: fusão (leves→Fe) SOBE até ele e fissão (pesados→Fe) DESCE até ele, ambas liberando energia. É o estado fundamental para onde a matéria nuclear cai — o ponto fixo cristal da órbita nuclear.","epistemico":"inferencia","exemplos":[],"id":"vale_do_ferro","memoria":"semantica","meta":{"dominio":"fisica","fonte":"energia nuclear"},"origem":"wiki","palavras_chave":[],"sinonimos":[],"tipo":"conceito","titulo":"O Vale do Ferro (Fe-56, o atrator)"}
\section{O Vale do Ferro (Fe-56, o atrator)}
O Fe-56/Ni-62 é o fundo do poço de energia nuclear — o ATRATOR: fusão (leves\ensuremath{\to}Fe) SOBE até ele e fissão (pesados\ensuremath{\to}Fe) DESCE até ele, ambas liberando energia. É o estado fundamental para onde a matéria nuclear cai — o ponto fixo cristal da órbita nuclear.
\prov{relações: relaciona particulas\_minerais}
\prov{proveniência: wiki \textperiodcentered{} energia nuclear \textperiodcentered{} inferencia \textperiodcentered{} confiança 1.0 \textperiodcentered{} domínio fisica \textperiodcentered{} história 2 versões no jornal}

%CRISTAL {"arestas":[{"alvo":"plataformas_qc_2025","confianca":1.0,"epistemico":"fato","origem":"wiki","rel":"relaciona"}],"confianca":1.0,"contraexemplos":[],"descricao":"Google Sycamore (2019, 53 qubits) amostrou circuitos aleatórios em ~200s, tarefa então estimada em ~10.000 anos clássicos — tarefa ARTIFICIAL, sem uso prático. CONTESTADA: redes tensoriais + supercomputação erodiram o custo clássico (de '10.000 anos' a horas/segundos; verificação Sunway). Jiuzhang (fotônica, boson sampling) idem, robustez contra spoofing contestada. Demonstrada em tarefa artificial, a vantagem específica foi erodida; a fronteira segue em disputa.","epistemico":"hipotese","exemplos":[],"id":"vantagem_quantica_contestada","memoria":"semantica","meta":{"autor":"broca-so","dominio":"fisica","fonte":"Arute+2019 Nature 574,505; Pan-Chen-Zhang 2021; Zhong+2020"},"origem":"wiki","palavras_chave":[],"sinonimos":[],"tipo":"conceito","titulo":"Vantagem quântica (Sycamore/Jiuzhang) — contestada"}
\section{Vantagem quântica (Sycamore/Jiuzhang) — contestada}
Google Sycamore (2019, 53 qubits) amostrou circuitos aleatórios em \~{}200s, tarefa então estimada em \~{}10.000 anos clássicos — tarefa ARTIFICIAL, sem uso prático. CONTESTADA: redes tensoriais + supercomputação erodiram o custo clássico (de '10.000 anos' a horas/segundos; verificação Sunway). Jiuzhang (fotônica, boson sampling) idem, robustez contra spoofing contestada. Demonstrada em tarefa artificial, a vantagem específica foi erodida; a fronteira segue em disputa.
\prov{relações: relaciona plataformas\_qc\_2025}
\prov{proveniência: wiki \textperiodcentered{} Arute+2019 Nature 574,505; Pan-Chen-Zhang 2021; Zhong+2020 \textperiodcentered{} hipotese \textperiodcentered{} confiança 1.0 \textperiodcentered{} domínio fisica \textperiodcentered{} história 2 versões no jornal}

%CRISTAL {"arestas":[{"alvo":"wheeler_dewitt","confianca":1.0,"epistemico":"fato","origem":"wiki","rel":"relaciona"},{"alvo":"minisuperspace","confianca":1.0,"epistemico":"fato","origem":"wiki","rel":"usa"},{"alvo":"equacoes_friedmann","confianca":1.0,"epistemico":"fato","origem":"wiki","rel":"relaciona"},{"alvo":"energia_escura_campo_psi","confianca":1.0,"epistemico":"fato","origem":"wiki","rel":"relaciona"}],"confianca":1.0,"contraexemplos":[],"descricao":"Com t=1 nos cinco botões, o imposto reduz-se a V=(P₀−D)²−ξ²+(1−s²)S; o vínculo de energia nula E=0 dá H_ξ²=2−2ρ_s/ξ² — a FORMA de Friedmann (−ξ² no papel de Λ, atrator H→√2 herdado de ẍ=2s). Leitura: ξ↔a (área) = analogia de minisuperspace, a ponte formal com Wheeler-DeWitt. RESSALVA: textualmente FRÁGIL — 'leitura declarada, jamais teorema'. O √2 e E≡0 (a 10⁻¹³) são firmes; a identificação ξ↔a é correspondência.","epistemico":"inferencia","exemplos":[],"id":"vinculo_e0_minisuperspace","memoria":"semantica","meta":{"autor":"broca-so","dominio":"cosmologia","fonte":"paper_2_cosmologia.tex l.691-693; paper_0_sintese.tex"},"origem":"wiki","palavras_chave":[],"sinonimos":[],"tipo":"conceito","titulo":"Vínculo E=0, forma de Friedmann e a ponte Wheeler-DeWitt"}
\section{Vínculo E=0, forma de Friedmann e a ponte Wheeler-DeWitt}
Com t=1 nos cinco botões, o imposto reduz-se a V=(P\ensuremath{_{0}}\ensuremath{-}D)\ensuremath{^{2}}\ensuremath{-}\ensuremath{\xi}\ensuremath{^{2}}+(1\ensuremath{-}s\ensuremath{^{2}})S; o vínculo de energia nula E=0 dá H\_\ensuremath{\xi}\ensuremath{^{2}}=2\ensuremath{-}2\ensuremath{\rho}\_s/\ensuremath{\xi}\ensuremath{^{2}} — a FORMA de Friedmann (\ensuremath{-}\ensuremath{\xi}\ensuremath{^{2}} no papel de \ensuremath{\Lambda}, atrator H\ensuremath{\to}\ensuremath{\sqrt{\,}}2 herdado de \ensuremath{\ddot{x}}=2s). Leitura: \ensuremath{\xi}\ensuremath{\leftrightarrow}a (área) = analogia de minisuperspace, a ponte formal com Wheeler-DeWitt. RESSALVA: textualmente FRÁGIL — 'leitura declarada, jamais teorema'. O \ensuremath{\sqrt{\,}}2 e E\ensuremath{\equiv}0 (a 10\ensuremath{^{-13}}) são firmes; a identificação \ensuremath{\xi}\ensuremath{\leftrightarrow}a é correspondência.
\prov{relações: relaciona wheeler\_dewitt; usa minisuperspace; relaciona equacoes\_friedmann; relaciona energia\_escura\_campo\_psi}
\prov{proveniência: wiki \textperiodcentered{} paper\_2\_cosmologia.tex l.691-693; paper\_0\_sintese.tex \textperiodcentered{} inferencia \textperiodcentered{} confiança 1.0 \textperiodcentered{} domínio cosmologia \textperiodcentered{} história 5 versões no jornal}

%CRISTAL {"arestas":[{"alvo":"km_previsao_cp","confianca":1.0,"epistemico":"fato","origem":"wiki","rel":"confirma"}],"confianca":1.0,"contraexemplos":[],"descricao":"A violação de CP, primeiro vista nos KÁONS neutros (Christenson-Cronin-Fitch-Turlay 1964, K_L→2π; ε_K≈2,23×10⁻³; Nobel 1980), depois nos MÉSONS B (BaBar/Belle 2001, sin2β via B⁰→J/ψK_S) — uma única fase CKM descreve os dois. E, novo: primeira observação em BÁRIONS (LHCb 2025, Λ_b⁰→pK⁻π⁺π⁻, A_CP=2,45%, 5,2σ, compatível com o MP). Todos FATO.","epistemico":"fato","exemplos":[],"id":"violacao_cp_k_b_barions","memoria":"semantica","meta":{"autor":"broca-so","dominio":"fisica","fonte":"Cronin-Fitch 1964; BaBar/Belle 2001; LHCb 2025 arXiv:2503.16954 Nature 643,1223"},"origem":"wiki","palavras_chave":[],"sinonimos":[],"tipo":"conceito","titulo":"Violação de CP: káons (1964), mésons B (2001), bárions (2025)"}
\section{Violação de CP: káons (1964), mésons B (2001), bárions (2025)}
A violação de CP, primeiro vista nos KÁONS neutros (Christenson-Cronin-Fitch-Turlay 1964, K\_L\ensuremath{\to}2\ensuremath{\pi}; \ensuremath{\varepsilon}\_K\ensuremath{\approx}2,23$\times$10\ensuremath{^{-3}}; Nobel 1980), depois nos MÉSONS B (BaBar/Belle 2001, sin2\ensuremath{\beta} via B\ensuremath{^{0}}\ensuremath{\to}J/\ensuremath{\psi}K\_S) — uma única fase CKM descreve os dois. E, novo: primeira observação em BÁRIONS (LHCb 2025, \ensuremath{\Lambda}\_b\ensuremath{^{0}}\ensuremath{\to}pK\ensuremath{^{-}}\ensuremath{\pi}\ensuremath{^{+}}\ensuremath{\pi}\ensuremath{^{-}}, A\_CP=2,45\%, 5,2\ensuremath{\sigma}, compatível com o MP). Todos FATO.
\prov{relações: confirma km\_previsao\_cp}
\prov{proveniência: wiki \textperiodcentered{} Cronin-Fitch 1964; BaBar/Belle 2001; LHCb 2025 arXiv:2503.16954 Nature 643,1223 \textperiodcentered{} fato \textperiodcentered{} confiança 1.0 \textperiodcentered{} domínio fisica \textperiodcentered{} história 2 versões no jornal}

%CRISTAL {"arestas":[{"alvo":"gema_cristal_um_quilate","confianca":1.0,"epistemico":"fato","origem":"wiki","rel":"usa"},{"alvo":"primitivas_analiticas","confianca":1.0,"epistemico":"fato","origem":"wiki","rel":"usa"},{"alvo":"difracao_atomica_turmalina","confianca":1.0,"epistemico":"fato","origem":"wiki","rel":"usa"},{"alvo":"joalheria_hub","confianca":1.0,"epistemico":"fato","origem":"wiki","rel":"parte_de"}],"confianca":1.0,"contraexemplos":[],"descricao":"A FERRAMENTA que cristaliza o procedimento: VOXELIZAÇÃO onde o VOXEL é a CELA DO CRISTAL. Cresce a rede cristalina hexagonal da turmalina DENTRO da borda analítica da gema (a superfície do dupolinômio 𝒢) até o peso EXATO em quilates — 6 tetraedros → 1 paralelepípedo por cela, o teste dentro_da_gema (F<0) seleciona o interior. Une a FORMA (o dupolinômio) e a MATÉRIA (a rede). Ferramenta central da joalheria: joalheria.voxelizar(quilates), joalheria.dentro; linguagem/gema_cristal.py. Dicionário BAI: voxelizador_maquina (borda_analitica↔borda, voxelizar↔colapso, empilhamento↔órbita, cela↔metal, quilate↔reta). Lê-se também como MINECRAFT (o dicionário paralelo).","epistemico":"fato","exemplos":[],"id":"voxelizador_cristalografico","memoria":"semantica","meta":{"dominio":"joalheria","fonte":"maquina de corte"},"origem":"wiki","palavras_chave":[],"sinonimos":[],"tipo":"conceito","titulo":"O Voxelizador Cristalográfico (a ferramenta): a borda analítica virada rede"}
\section{O Voxelizador Cristalográfico (a ferramenta): a borda analítica virada rede}
A FERRAMENTA que cristaliza o procedimento: VOXELIZAÇÃO onde o VOXEL é a CELA DO CRISTAL. Cresce a rede cristalina hexagonal da turmalina DENTRO da borda analítica da gema (a superfície do dupolinômio \ensuremath{\mathcal{G}}) até o peso EXATO em quilates — 6 tetraedros \ensuremath{\to} 1 paralelepípedo por cela, o teste dentro\_da\_gema (F<0) seleciona o interior. Une a FORMA (o dupolinômio) e a MATÉRIA (a rede). Ferramenta central da joalheria: joalheria.voxelizar(quilates), joalheria.dentro; linguagem/gema\_cristal.py. Dicionário BAI: voxelizador\_maquina (borda\_analitica\ensuremath{\leftrightarrow}borda, voxelizar\ensuremath{\leftrightarrow}colapso, empilhamento\ensuremath{\leftrightarrow}órbita, cela\ensuremath{\leftrightarrow}metal, quilate\ensuremath{\leftrightarrow}reta). Lê-se também como MINECRAFT (o dicionário paralelo).
\prov{relações: usa gema\_cristal\_um\_quilate; usa primitivas\_analiticas; usa difracao\_atomica\_turmalina; parte\_de joalheria\_hub}
\prov{proveniência: wiki \textperiodcentered{} maquina de corte \textperiodcentered{} fato \textperiodcentered{} confiança 1.0 \textperiodcentered{} domínio joalheria \textperiodcentered{} história 130 versões no jornal}

%CRISTAL {"arestas":[{"alvo":"cosmologia_quantica","confianca":1.0,"epistemico":"fato","origem":"wiki","rel":"relaciona"},{"alvo":"minisuperspace","confianca":1.0,"epistemico":"fato","origem":"wiki","rel":"usa"}],"confianca":1.0,"contraexemplos":[],"descricao":"Vínculo hamiltoniano quântico sobre a função de onda do universo Ψ no superespaço; como Ĥ é vínculo, ĤΨ=0 — sem tempo explícito (o 'problema do tempo'). Programa teórico de gravidade quântica canônica, NÃO confirmado; sem interpretação probabilística satisfatória. A mesma s̈=2s do Lopes é, formalmente, WDW em minisuperspace (porta aberta).","epistemico":"hipotese","exemplos":[],"id":"wheeler_dewitt","memoria":"semantica","meta":{"autor":"broca-so","dominio":"fisica","fonte":"DeWitt 1967; Wheeler"},"origem":"wiki","palavras_chave":[],"sinonimos":[],"tipo":"conceito","titulo":"Equação de Wheeler-DeWitt (ĤΨ=0) — função de onda do universo"}
\section{Equação de Wheeler-DeWitt (\ensuremath{\hat{H}}\ensuremath{\Psi}=0) — função de onda do universo}
Vínculo hamiltoniano quântico sobre a função de onda do universo \ensuremath{\Psi} no superespaço; como \ensuremath{\hat{H}} é vínculo, \ensuremath{\hat{H}}\ensuremath{\Psi}=0 — sem tempo explícito (o 'problema do tempo'). Programa teórico de gravidade quântica canônica, NÃO confirmado; sem interpretação probabilística satisfatória. A mesma s\ensuremath{}=2s do Lopes é, formalmente, WDW em minisuperspace (porta aberta).
\prov{relações: relaciona cosmologia\_quantica; usa minisuperspace}
\prov{proveniência: wiki \textperiodcentered{} DeWitt 1967; Wheeler \textperiodcentered{} hipotese \textperiodcentered{} confiança 1.0 \textperiodcentered{} domínio fisica \textperiodcentered{} história 4 versões no jornal}

%CRISTAL {"arestas":[{"alvo":"fermat_na_maquina","confianca":1.0,"epistemico":"fato","origem":"wiki","rel":"especializa"},{"alvo":"joalheria_hub","confianca":1.0,"epistemico":"fato","origem":"wiki","rel":"parte_de"}],"confianca":1.0,"contraexemplos":[],"descricao":"O aprofundamento do lado WILES, costurado ao dinamica_selberg — e computável. O FROBENIUS de uma curva elíptica E/𝔽_p tem traço a_p=p+1−#E(𝔽_p) e autovalores α,β (raízes de z²−a_p z+p) com |α|=|β|=√p: a cota de HASSE-WEIL, que É a HIPÓTESE DE RIEMANN das curvas sobre corpos finitos (PROVADA). Normalizando, |α|/√p=1 EXATO — o Frobenius é um objeto de BORDA (Salem, ρ=1, a linha crítica), classificado pela MESMA parábola do dinamica_selberg (não cristal/Pisot, não caos). A MODULARIDADE (Taniyama-Shimura-Wiles) a_p(E)=c(p,f) é a DUALIDADE ESPECTRO=GEOMETRIA da fórmula do traço de Selberg: o a_p (geométrico, o Frobenius ~ uma geodésica) = o c(p,f) (espectral, o autovalor de HECKE da autoforma), na superfície modular X₀(N)=Γ₀(N)\\ℍ, que é hiperbólica (fluxo de ANOSOV). Verificado na curva canônica X₀(11): a_p (contando pontos) = c(p,f) com f=η(z)²η(11z)² (o produto eta), primo a primo (c(2)=−2, c(3)=−1, c(5)=1, c(7)=−2, c(13)=4). E o fator local de L(E,s)=det(1−Frob·p^{-s}) tem a forma do fator de geodésica da zeta de Selberg (primo p↔geodésica γ, log p↔ℓ(γ)). O INVERSO NA BORDA (a correspondência 1-a-1): como α·β=det=p, o 2º autovalor é o p-INVERSO β=p/α; normalizados à borda (|·|=1), â e o seu conjugado são INVERSOS MÚTUOS (â·â̄=1, i.e. â̄=1/â). O 'inverso' irracional é o recíproco 1/â=â̄ e o raio 1/√p (√p irracional, p não é quadrado perfeito) — NÃO o racional 1/p; e p↦√p é 1-a-1 (injetiva). ESPELHA a Parte IV: lá o metal tem σ̄=−1/σ (o mesmo flip conjugado=inverso) mas com norma −1 (o par FORA do círculo → gato, cristal+caos); aqui o Frobenius tem â̄=1/â com norma +1 (NO círculo → borda). O flip é o mesmo; o MÓDULO decide o ramo. A ESTRUTURA FINA DA BORDA: como |α|=√p, o Frobenius é uma ROTAÇÃO por θ_p (α=√p·e^{iθ_p}, o Frobenius normalizado em SU(2)); os ângulos θ_p equidistribuem pela lei de SATO-TATE (2/π)sin²θ (a Haar de SU(2)), verificado em X₀(11): os momentos pares de a_p/√p → os números de CATALAN (m2=0.99→1, m4=1.94→2, m6=4.85→5, sobre 366 primos). Isso DISTINGUE a borda do cristal: a borda (Sato-Tate) tem m2→1, o cristal (a reta mineral, equidistribuição uniforme de Weyl) tem m2→2 — cada ramo da parábola tem a sua lei de densidade. RESSALVA: as peças são fatos/computações (Hasse-Weil, a_p=c(p,f), Sato-Tate é TEOREMA p/ não-CM); a ponte primo↔geodésica e L↔Z é o dicionário estrutural (a fórmula explícita), analogia declarada. linguagem/wiles_selberg.py (8/8): pontos_c11, traco_frobenius, autovalores_frobenius, frobenius_na_borda, forma_modular_c11, modularidade, fator_local_L, angulo_frobenius, sato_tate. Asset: wiles_sato_tate.png.","epistemico":"fato","exemplos":[],"id":"wiles_selberg","memoria":"semantica","meta":{"dominio":"joalheria","fonte":"maquina de corte"},"origem":"wiki","palavras_chave":[],"sinonimos":[],"tipo":"conceito","titulo":"O lado Wiles do TLF ligado à dinâmica de Selberg: o Frobenius na borda, a modularidade = espectro↔geometria"}
\section{O lado Wiles do TLF ligado à dinâmica de Selberg: o Frobenius na borda, a modularidade = espectro\ensuremath{\leftrightarrow}geometria}
O aprofundamento do lado WILES, costurado ao dinamica\_selberg — e computável. O FROBENIUS de uma curva elíptica E/\ensuremath{\mathbb{F}}\_p tem traço a\_p=p+1\ensuremath{-}\#E(\ensuremath{\mathbb{F}}\_p) e autovalores \ensuremath{\alpha},\ensuremath{\beta} (raízes de z\ensuremath{^{2}}\ensuremath{-}a\_p z+p) com |\ensuremath{\alpha}|=|\ensuremath{\beta}|=\ensuremath{\sqrt{\,}}p: a cota de HASSE-WEIL, que É a HIPÓTESE DE RIEMANN das curvas sobre corpos finitos (PROVADA). Normalizando, |\ensuremath{\alpha}|/\ensuremath{\sqrt{\,}}p=1 EXATO — o Frobenius é um objeto de BORDA (Salem, \ensuremath{\rho}=1, a linha crítica), classificado pela MESMA parábola do dinamica\_selberg (não cristal/Pisot, não caos). A MODULARIDADE (Taniyama-Shimura-Wiles) a\_p(E)=c(p,f) é a DUALIDADE ESPECTRO=GEOMETRIA da fórmula do traço de Selberg: o a\_p (geométrico, o Frobenius \~{} uma geodésica) = o c(p,f) (espectral, o autovalor de HECKE da autoforma), na superfície modular X\ensuremath{_{0}}(N)=\ensuremath{\Gamma}\ensuremath{_{0}}(N)\textbackslash{}\ensuremath{\mathbb{H}}, que é hiperbólica (fluxo de ANOSOV). Verificado na curva canônica X\ensuremath{_{0}}(11): a\_p (contando pontos) = c(p,f) com f=\ensuremath{\eta}(z)\ensuremath{^{2}}\ensuremath{\eta}(11z)\ensuremath{^{2}} (o produto eta), primo a primo (c(2)=\ensuremath{-}2, c(3)=\ensuremath{-}1, c(5)=1, c(7)=\ensuremath{-}2, c(13)=4). E o fator local de L(E,s)=det(1\ensuremath{-}Frob\textperiodcentered p\^{}\{-s\}) tem a forma do fator de geodésica da zeta de Selberg (primo p\ensuremath{\leftrightarrow}geodésica \ensuremath{\gamma}, log p\ensuremath{\leftrightarrow}\ensuremath{\ell}(\ensuremath{\gamma})). O INVERSO NA BORDA (a correspondência 1-a-1): como \ensuremath{\alpha}\textperiodcentered \ensuremath{\beta}=det=p, o 2º autovalor é o p-INVERSO \ensuremath{\beta}=p/\ensuremath{\alpha}; normalizados à borda (|\textperiodcentered |=1), â e o seu conjugado são INVERSOS MÚTUOS (â\textperiodcentered â\ensuremath{}=1, i.e. â\ensuremath{}=1/â). O 'inverso' irracional é o recíproco 1/â=â\ensuremath{} e o raio 1/\ensuremath{\sqrt{\,}}p (\ensuremath{\sqrt{\,}}p irracional, p não é quadrado perfeito) — NÃO o racional 1/p; e p\ensuremath{\mapsto}\ensuremath{\sqrt{\,}}p é 1-a-1 (injetiva). ESPELHA a Parte IV: lá o metal tem \ensuremath{\sigma}\ensuremath{}=\ensuremath{-}1/\ensuremath{\sigma} (o mesmo flip conjugado=inverso) mas com norma \ensuremath{-}1 (o par FORA do círculo \ensuremath{\to} gato, cristal+caos); aqui o Frobenius tem â\ensuremath{}=1/â com norma +1 (NO círculo \ensuremath{\to} borda). O flip é o mesmo; o MÓDULO decide o ramo. A ESTRUTURA FINA DA BORDA: como |\ensuremath{\alpha}|=\ensuremath{\sqrt{\,}}p, o Frobenius é uma ROTAÇÃO por \ensuremath{\theta}\_p (\ensuremath{\alpha}=\ensuremath{\sqrt{\,}}p\textperiodcentered e\^{}\{i\ensuremath{\theta}\_p\}, o Frobenius normalizado em SU(2)); os ângulos \ensuremath{\theta}\_p equidistribuem pela lei de SATO-TATE (2/\ensuremath{\pi})sin\ensuremath{^{2}}\ensuremath{\theta} (a Haar de SU(2)), verificado em X\ensuremath{_{0}}(11): os momentos pares de a\_p/\ensuremath{\sqrt{\,}}p \ensuremath{\to} os números de CATALAN (m2=0.99\ensuremath{\to}1, m4=1.94\ensuremath{\to}2, m6=4.85\ensuremath{\to}5, sobre 366 primos). Isso DISTINGUE a borda do cristal: a borda (Sato-Tate) tem m2\ensuremath{\to}1, o cristal (a reta mineral, equidistribuição uniforme de Weyl) tem m2\ensuremath{\to}2 — cada ramo da parábola tem a sua lei de densidade. RESSALVA: as peças são fatos/computações (Hasse-Weil, a\_p=c(p,f), Sato-Tate é TEOREMA p/ não-CM); a ponte primo\ensuremath{\leftrightarrow}geodésica e L\ensuremath{\leftrightarrow}Z é o dicionário estrutural (a fórmula explícita), analogia declarada. linguagem/wiles\_selberg.py (8/8): pontos\_c11, traco\_frobenius, autovalores\_frobenius, frobenius\_na\_borda, forma\_modular\_c11, modularidade, fator\_local\_L, angulo\_frobenius, sato\_tate. Asset: wiles\_sato\_tate.png.
\prov{relações: especializa fermat\_na\_maquina; parte\_de joalheria\_hub}
\prov{proveniência: wiki \textperiodcentered{} maquina de corte \textperiodcentered{} fato \textperiodcentered{} confiança 1.0 \textperiodcentered{} domínio joalheria \textperiodcentered{} história 24 versões no jornal}

%CRISTAL {"arestas":[{"alvo":"materia_escura","confianca":1.0,"epistemico":"fato","origem":"wiki","rel":"relaciona"}],"confianca":1.0,"contraexemplos":[],"descricao":"Partículas massivas fracamente interagentes (~GeV-TeV) que produzem naturalmente a abundância certa ('WIMP miracle'), motivadas por supersimetria. Paradigma dominante por décadas, agora sob FORTE pressão pelos resultados nulos da detecção direta; não excluído, mas o espaço de parâmetros encolhe.","epistemico":"hipotese","exemplos":[],"id":"wimps_candidato","memoria":"semantica","meta":{"autor":"broca-so","dominio":"cosmologia_dados","fonte":"WIMP miracle; SUSY"},"origem":"wiki","palavras_chave":[],"sinonimos":[],"tipo":"conceito","titulo":"WIMPs — o paradigma sob pressão"}
\section{WIMPs — o paradigma sob pressão}
Partículas massivas fracamente interagentes (\~{}GeV-TeV) que produzem naturalmente a abundância certa ('WIMP miracle'), motivadas por supersimetria. Paradigma dominante por décadas, agora sob FORTE pressão pelos resultados nulos da detecção direta; não excluído, mas o espaço de parâmetros encolhe.
\prov{relações: relaciona materia\_escura}
\prov{proveniência: wiki \textperiodcentered{} WIMP miracle; SUSY \textperiodcentered{} hipotese \textperiodcentered{} confiança 1.0 \textperiodcentered{} domínio cosmologia\_dados \textperiodcentered{} história 2 versões no jornal}

%CRISTAL {"arestas":[{"alvo":"matriz_ckm","confianca":1.0,"epistemico":"fato","origem":"wiki","rel":"usa"}],"confianca":1.0,"contraexemplos":[],"descricao":"Expansão da CKM em potências de λ=sinθ_C que revela a hierarquia: diagonal≈1, afastamento em ordens crescentes de λ. Quatro parâmetros: λ=0,2265, A=0,79, ρ̄=0,141, η̄=0,357 (a violação de CP mora em η̄≠0). A CKM é QUASE DIAGONAL: |V_us|≈0,224 (~λ), |V_cb|≈0,041 (~λ²), |V_ub|≈0,0038 (~λ³) — contraste marcante com a PMNS (ângulos grandes). FATO (valores PDG 2024/CKMfitter).","epistemico":"fato","exemplos":[],"id":"wolfenstein_hierarquia_ckm","memoria":"semantica","meta":{"autor":"broca-so","dominio":"fisica","fonte":"Wolfenstein 1983 PRL 51,1945; PDG 2024; CKMfitter"},"origem":"wiki","palavras_chave":[],"sinonimos":[],"tipo":"conceito","titulo":"Parametrização de Wolfenstein e a hierarquia da CKM"}
\section{Parametrização de Wolfenstein e a hierarquia da CKM}
Expansão da CKM em potências de \ensuremath{\lambda}=sin\ensuremath{\theta}\_C que revela a hierarquia: diagonal\ensuremath{\approx}1, afastamento em ordens crescentes de \ensuremath{\lambda}. Quatro parâmetros: \ensuremath{\lambda}=0,2265, A=0,79, \ensuremath{\rho}\ensuremath{}=0,141, \ensuremath{\eta}\ensuremath{}=0,357 (a violação de CP mora em \ensuremath{\eta}\ensuremath{}\ensuremath{\ne}0). A CKM é QUASE DIAGONAL: |V\_us|\ensuremath{\approx}0,224 (\~{}\ensuremath{\lambda}), |V\_cb|\ensuremath{\approx}0,041 (\~{}\ensuremath{\lambda}\ensuremath{^{2}}), |V\_ub|\ensuremath{\approx}0,0038 (\~{}\ensuremath{\lambda}\ensuremath{^{3}}) — contraste marcante com a PMNS (ângulos grandes). FATO (valores PDG 2024/CKMfitter).
\prov{relações: usa matriz\_ckm}
\prov{proveniência: wiki \textperiodcentered{} Wolfenstein 1983 PRL 51,1945; PDG 2024; CKMfitter \textperiodcentered{} fato \textperiodcentered{} confiança 1.0 \textperiodcentered{} domínio fisica \textperiodcentered{} história 2 versões no jornal}

%CRISTAL {"arestas":[{"alvo":"cosmologia_quantica_hub","confianca":1.0,"epistemico":"fato","origem":"wiki","rel":"parte_de"}],"confianca":1.0,"contraexemplos":[],"descricao":"Provar rigorosamente que Yang-Mills quântico em ℝ⁴ (para qualquer grupo compacto simples) EXISTE (axiomas de Wightman) e tem gap de massa Δ>0 — o menor estado acima do vácuo tem massa estritamente positiva. Prêmio de US1M do Clay Institute. ABERTO / NÃO RESOLVIDO (um dos 7 problemas do milênio); lattice sugere o gap, sem prova.","epistemico":"hipotese","exemplos":[],"id":"yang_mills_mass_gap_milenio","memoria":"semantica","meta":{"autor":"broca-so","dominio":"fisica","fonte":"Jaffe & Witten 2000, Clay Mathematics Institute"},"origem":"wiki","palavras_chave":[],"sinonimos":[],"tipo":"conceito","titulo":"Mass gap de Yang-Mills (Problema do Milênio)"}
\section{Mass gap de Yang-Mills (Problema do Milênio)}
Provar rigorosamente que Yang-Mills quântico em \ensuremath{\mathbb{R}}\ensuremath{^{4}} (para qualquer grupo compacto simples) EXISTE (axiomas de Wightman) e tem gap de massa \ensuremath{\Delta}>0 — o menor estado acima do vácuo tem massa estritamente positiva. Prêmio de US1M do Clay Institute. ABERTO / NÃO RESOLVIDO (um dos 7 problemas do milênio); lattice sugere o gap, sem prova.
\prov{relações: parte\_de cosmologia\_quantica\_hub}
\prov{proveniência: wiki \textperiodcentered{} Jaffe \& Witten 2000, Clay Mathematics Institute \textperiodcentered{} hipotese \textperiodcentered{} confiança 1.0 \textperiodcentered{} domínio fisica \textperiodcentered{} história 3 versões no jornal}

%CRISTAL {"arestas":[{"alvo":"materia_escura","confianca":1.0,"epistemico":"fato","origem":"wiki","rel":"confirma"}],"confianca":1.0,"contraexemplos":[],"descricao":"Galáxias no aglomerado de Coma se movem rápido demais para serem retidas pela massa visível: massa dinâmica ≫ massa luminosa. Primeira detecção histórica; cunhou 'dunkle Materie' (matéria escura). Robusta em direção, imprecisa no número original (H₀ da época).","epistemico":"fato","exemplos":[],"id":"zwicky_coma_dunkle_materie","memoria":"semantica","meta":{"autor":"broca-so","dominio":"cosmologia_dados","fonte":"Zwicky 1933 Helv.Phys.Acta 6,110"},"origem":"wiki","palavras_chave":[],"sinonimos":[],"tipo":"conceito","titulo":"Zwicky e o aglomerado de Coma (dunkle Materie)"}
\section{Zwicky e o aglomerado de Coma (dunkle Materie)}
Galáxias no aglomerado de Coma se movem rápido demais para serem retidas pela massa visível: massa dinâmica \ensuremath{\gg} massa luminosa. Primeira detecção histórica; cunhou 'dunkle Materie' (matéria escura). Robusta em direção, imprecisa no número original (H\ensuremath{_{0}} da época).
\prov{relações: confirma materia\_escura}
\prov{proveniência: wiki \textperiodcentered{} Zwicky 1933 Helv.Phys.Acta 6,110 \textperiodcentered{} fato \textperiodcentered{} confiança 1.0 \textperiodcentered{} domínio cosmologia\_dados \textperiodcentered{} história 2 versões no jornal}

\end{document}
